\section{Small models and exact poset minimality}
\label{sec:C}

\subsection{The independent enumeration and its exact scope}
\begin{proposition}\label{prop:C-census}
On a set of $n$ labelled elements with identity fixed at $0$, the numbers
of monoid structures for $1\le n\le7$ are
\[
 1,\ 2,\ 11,\ 156,\ 4122,\ 208672,\ 18507440.
\]
For $n\le6$ the numbers up to isomorphism and opposite are
$1,2,6,27,156,1373$. Every integral $(H_1,H_2,H_3)$ multiplicity
in \texttt{OWN\_RESULTS.md}, Statement 6.1(a)--(c), is reproduced.
The connected category counts, with object identities labelled and
nonidentity arrows sorted by source and target, are as follows.
\begin{center}
\begin{tabular}{rrrrr}\toprule
Arrows & 2 objects & 3 objects & 4 objects & 5 objects\\\midrule
3&2&0&0&0\\
4&11&0&0&0\\
5&76&6&0&0\\
6&977&54&0&0\\
7&22798&498&32&0\\
8&1016844&6324&406&0\\
9&\text{not enumerated}&133543&4468&250\\\bottomrule
\end{tabular}
\end{center}
All connected categories with at most seven arrows are covered. The
additional eight- and nine-arrow counts have only the scopes displayed;
the two-object nine-arrow cell is not claimed. Cases with
$N<2k-1$ are zero because a connected graph on $k$ objects needs at least
$k-1$ nonidentity arrows.
\end{proposition}
\begin{proof}
The independent programs \texttt{c\_enum.cpp} and
\texttt{c\_categories.cpp} use constraint propagation. Source and target
sets fix the allowed values of each composition entry; unit entries are
fixed first. Whenever $ab$ and $bc$ are known, associativity identifies
the entries $(ab)c$ and $a(bc)$: a known entry forces the other, and
unequal known entries reject the branch. A remaining unknown entry is
then assigned each allowed value. Every completion satisfies every
associativity equation, and each possible table follows exactly one
branch, so this proves both soundness and completeness. The programs do
not import or reuse the supplied enumeration code.

All homology calculations use independently constructed normalized nerve
boundary matrices and exact integer Smith normal forms. If
$d_{q+1}d_q$ is read with the standard chain convention as
$d_qd_{q+1}=0$, the lattice $\ker d_q$ is saturated in $C_q$.
It follows that the nonunit Smith factors of $d_{q+1}$ computed in
$C_q$ are precisely the torsion factors in
$\ker d_q/\operatorname{im}d_{q+1}$. The free rank is
$\operatorname{rank}C_q-\operatorname{rank}d_q-
\operatorname{rank}d_{q+1}$. Relabelling caches identify only explicitly
verified isomorphisms or opposite categories. The full tables, counts,
homology multiplicities and reduction certificates are in the accompanying
TSV and JSON outputs. The seven-arrow monoid search keeps every class
with nontrivial universal group and every class not reducible by a corner
$eMe=eM$ or $eMe=Me$; the total labelled count is taken before this filter.
The integral computations reproduce every listed earlier multiplicity,
so there is no numerical disagreement with R7.
\end{proof}

\begin{theorem}\label{thm:C-RP2-unique}
Up to isomorphism, exactly one category with three objects and nine arrows
has classifying space of type $\RP^2$. It is $R_2$: the objects are
$0,1,2$, each $\operatorname{Hom}(i,j)$ for $i<j$ is $C_2$, all
endomorphisms are identities, and composition is addition in $C_2$.
Its classifying space is homeomorphic to $\RP^2$.
\end{theorem}
\begin{proof}
Among all $133543$ sorted labelled three-object nine-arrow tables, the
independent calculation finds $39$ classes up to isomorphism and opposite
with $H_1=\Z/2$. Exactly one also has $H_2=H_3=0$; its table is
the displayed rule for $R_2$. This filter is only a necessary condition;
the following geometric argument proves sufficiency. The opposite of
$R_2$ is isomorphic to $R_2$ by reversing the three object levels,
so the single isomorphism-and-opposite class is a single isomorphism class.

Let $P$ have six objects $(i,\epsilon)$, $i\in\{0,1,2\}$ and
$\epsilon\in C_2$, with a unique arrow whenever the first coordinate
strictly increases. The order complex is the join of three two-point
sets, the boundary of the octahedron, hence $S^2$. Simultaneously changing
all three binary coordinates is its antipodal involution. The quotient
category is $R_2$, with an arrow labelled by the difference of its endpoint
coordinates. Lifting a chain after choosing its initial binary coordinate
is unique, so the nerve quotient is exactly $NR_2$. Geometric realization
therefore gives the homeomorphism $BR_2\cong S^2/(x\sim-x)=\RP^2$.
\end{proof}

\subsection{Four seven-element monoids of type $S^3$}
\begin{theorem}\label{thm:C-S3-four}
There are exactly four isomorphism classes of categories with seven arrows
and classifying space of type $S^3$. All are monoids. They are $P,P^{\op}$,
$Q,Q^{\op}$, where $P$ and $Q$ have the following tables, with $0$ the
identity:
\[
 P=\begin{pmatrix}
0&1&2&3&4&5&6\\
1&1&1&3&3&1&1\\
2&2&2&4&4&2&2\\
3&1&1&3&3&1&1\\
4&2&2&4&4&2&2\\
5&1&2&3&4&5&5\\
6&1&2&3&4&6&6
\end{pmatrix},\qquad
 Q=\begin{pmatrix}
0&1&2&3&4&5&6\\
1&2&2&2&4&4&1\\
2&2&2&2&4&4&2\\
3&3&3&3&5&5&3\\
4&2&2&2&4&4&2\\
5&3&3&3&5&5&3\\
6&2&2&3&4&5&6
\end{pmatrix}.
\]
\end{theorem}
\begin{proof}
The census excludes every category with at least two objects and at most
seven arrows by $(H_1,H_2,H_3)\ne(0,0,\Z)$, and excludes every monoid
of order at most six. Any reducible order-seven monoid has the homotopy
type of a smaller corner, so cannot contribute. Among the 25 remaining
corner-irreducible classes, exactly these two tables have the required
triple. Their opposite tables account for four ordinary isomorphism classes:
the finite isomorphism check tries all $6!$ identity-fixing bijections and
finds neither table is self-opposite. The first table is a band, whereas
the second is not, so no cross identification is possible.

We prove the all-degree assertion by finite projective resolutions.
For $P$ use the names $a=x_{00}$, $b=x_{01}$, $c=x_{10}$,
$d=x_{11}$ for its rectangular band and $q,r$ for its remaining
nonidentity elements. They act on the band by
$x_{ij}q=x_{ij}r=x_{i0}$ and $qx_{ij}=rx_{ij}=x_{ij}$,
and satisfy $q^2=qr=q$, $r^2=rq=r$. Over $R=\Z P$ the resolution is
\[
 0\longrightarrow Rb\oplus Rq
 \xrightarrow{(u,v)\mapsto(0,u+v)}Ra\oplus R
 \xrightarrow{(u,v)\mapsto u+v(r-q)}Rq
 \xrightarrow{w\mapsto w(a-c)}Ra
 \xrightarrow{\epsilon}\Z\longrightarrow0.
\]
Here $Ra=\Z\{a,c\}$, $Rb=\Z\{b,d\}$ and
$Rq=\Z\{a,c,q,r\}$. The first nonaugmentation kernel is
$Ra\oplus\Z(r-q)$, and the next is $0\oplus\Z(P\setminus\{1\})$,
whose basis partitions as the displayed $Rb\oplus Rq$.
These observations prove exactness integrally and injectivity at the left.
After tensoring with the trivial module, the three nonaugmentation
maps are respectively
$\left(\begin{smallmatrix}0&0\\1&1\end{smallmatrix}\right)$,
$(1\;0)$ and $0$. Thus the reduced homology is $\Z$ in degree three only.

For $Q$, write $y=1,a=2,c=3,b=4,d=5,q=6$ in the numeric table;
the identity is still the numeric element $0$. In $R=\Z Q$ put
$h=q-a$. Then $h^2=h$ and $Rh=\Z\{q-a,y-a\}$.
There is a resolution
\[
 0\longrightarrow Rb\oplus Rh
 \xrightarrow{(ub,v)\mapsto u(b-a)+v}R
 \xrightarrow{w\mapsto wy}Rq
 \xrightarrow{w\mapsto w(a-c)}Ra
 \xrightarrow{\epsilon}\Z\longrightarrow0.
\]
The map $Rb\to R$ is well defined by $b\mapsto b-a$,
$d\mapsto d-c$, as direct multiplication checks. The kernel of
$Rq\to Ra$ is $\Z\{a,c,y\}$, exactly the image of right
multiplication by $y$. Its kernel in $R$ is the direct sum
$\Z\{b-a,d-c\}\oplus\Z\{q-a,y-a\}$, proving exactness.
The module $Rh$ has zero coinvariants, since its generating idempotent
has augmentation zero. Tensoring therefore gives
$0\to\Z\xrightarrow0\Z\xrightarrow1\Z\xrightarrow0\Z$,
with the leftmost $\Z$ in degree three. Both monoids have trivial
universal group: all elements of $P$ are idempotent; in $Q$ the band
and $q$ are idempotent and $ya=a$ kills $y$. The homological Whitehead
theorem now identifies both spaces with $S^3$. Opposite categories have
homeomorphic classifying spaces.
\end{proof}

\subsection{The exact minimum for posets of type $\RP^2$}
\begin{lemma}[A torsion-support estimate]\label{lem:C-torsion-support}
Let $K$ be a finite simplicial complex and suppose
$H_1(K;\Z)$ has nonzero $p$-torsion for a prime $p$.
There is a connected two-dimensional subcomplex $L$ with
$v$ vertices and $e$ edges such that $e\ge3v-3$, and the induced
full subcomplex on these vertices still has $p$-torsion in $H_1$.
\end{lemma}
\begin{proof}
The integral coefficient sequence for multiplication by $p$ supplies a
mod-$p$ two-cycle $z$ whose integral Bockstein has nonzero image in
$H_1(K;\Z)$. Choose such $z$ with the smallest possible number of
triangles in its support, and let $L$ consist of those triangles and
their faces. Any nonzero two-cycle $w$ supported on a proper subset
of these triangles is impossible: if its Bockstein survives in $K$,
it contradicts minimality directly; otherwise subtract a scalar
multiple of $w$ from $z$ that cancels at least one triangle, keeping
the same nonzero Bockstein image. Thus the kernel of the boundary
on the triangle support is one-dimensional. Indeed two independent
cycles on that support have a nonzero linear combination with strictly
smaller support. Hence $H_2(L;\F_p)=\F_p$.

The same minimality makes $L$ connected: a component supporting a
nonzero Bockstein would otherwise have smaller support. The Bockstein
$H_2(L;\F_p)\to H_1(L;\Z)[p]$ is nonzero, hence injective.
The universal coefficient sequence then gives
$H_2(L;\Z)\otimes\F_p=0$. The integral group $H_2(L;\Z)$ is
free, since $L$ is two-dimensional, and so it is zero. Therefore
$\chi(L)=1-\operatorname{rank}H_1(L;\Z)\le1$.

Every edge of $L$ belongs to at least two supporting triangles: the
boundary coefficient of a triangle with nonzero coefficient could not
cancel at an edge contained in only that triangle. If $f$ is the
number of triangles, $3f\ge2e$. Consequently
\[
 1\ge\chi(L)=v-e+f\ge v-e/3,
 \qquad e\ge3v-3.
\]
Finally the Bockstein class maps nontrivially through the induced full
subcomplex on the support vertices, because its further image in $K$
is nonzero. That full subcomplex therefore has nonzero $p$-torsion in
its first integral homology.
\end{proof}

\begin{theorem}\label{thm:C-poset49}
The least number of arrows, including identities, of a finite poset with
classifying space of type $\RP^2$ is exactly $49$. More generally every
finite poset whose first integral homology has torsion has at least
$49$ arrows.
\end{theorem}
\begin{proof}
Apply Lemma~\ref{lem:C-torsion-support} to the order complex of the poset
$P$. Its induced full subcomplex on the $v$ support vertices is the
order complex of the induced subposet. Cianci--Ottina,
\cite[Theorem 5.7]{cianciottina2015}, states that a finite $T_0$ space
with at most twelve points has free fundamental group and therefore
torsion-free first integral homology. A finite poset is a finite $T_0$
space, and its order complex computes that homology; hence $v\ge13$.
If $P$ has $N$ points, its arrows are its $N$ identities together with
all strict comparable pairs. Every supporting edge is such a pair, so
\[
 |\operatorname{Mor}P|\ge N+e\ge v+3v-3\ge49.
\]

For the upper bound, take the face poset of the regular cell decomposition
$S^2/C_2$ obtained from the octahedron under the antipodal action.
There are three vertex orbits, six edge orbits and four triangle orbits.
Each edge has two distinct endpoints; each triangle has three distinct
edges and three distinct vertices. Its face poset therefore has
$3+6+4=13$ identities and $12+12+12=36$ strict incidences, for a total
of $49$. The barycentric subdivision of this regular cell structure
is the order complex of that face poset and is homeomorphic to $\RP^2$.
All incidences, transitivity and homology are independently checked in
\texttt{code/c\_models.py}. This proves the bound is attained.
\end{proof}

\paragraph{Step-by-step check of the earlier bound R8.}
First, a poset minimizing arrows can have no comparability vertex of
degree zero or one: connectedness rules out the former, and the latter
is a removable beat point. Second, the opened Corollary 5.8 of
\cite{cianciottina2015} gives at least thirteen vertices, and its
Theorem 5.10 identifies the thirteen-point $\RP^2$ models; their incidence
count is the $49$ just computed. Third, a vertex-minimal full subcomplex
with nonzero mod-$2$ second homology has nonzero first homology in every
vertex link, by the deletion--star Mayer--Vietoris sequence. Fourth,
its links are flag complexes; nonzero first homology requires at least
four vertices, and the support itself requires at least six vertices
(five vertices of degree at least four would form a simplex).
Finally summing degrees gives $e\ge N+6$, and $N\ge14$ gives
$N+e\ge34$. Every step of the supplied proof is correct. The
Bockstein-support argument above sharpens its support size and edge
estimate, settling the former interval $[34,49]$.

\begin{proposition}\label{prop:C-rose}
The least number of arrows of a category of type $S^1\vee S^1$ is five.
\end{proposition}
\begin{proof}
Two objects with three parallel arrows between them have five arrows
including the two identities; the nerve is the graph with two vertices
and three edges, of type $S^1\vee S^1$. The independent census through
four arrows has no category with first homology $\Z^2$; one-object
categories have finite abelianized universal group and cannot provide
an exception. Thus five is minimal.
\end{proof}

\subsection{The exact arrow minimum for the torus}
\begin{lemma}\label{lem:C-small-sumset}
For finite nonempty subsets $A,B$ of a torsion-free abelian group,
$|A+B|\geq |A|+|B|-1$. If $|A|=|B|=3$ and $|A+B|=5$,
then, after separate translations, $A=B=\{0,x,2x\}$ for some
nonzero $x$.
\end{lemma}
\begin{proof}
The subgroup generated by the two finite sets is a finitely generated
torsion-free abelian group, hence admits a translation-invariant total
order. List each set increasingly. Increasing through all elements of
$A$ while fixing the first element of $B$, and then through the remaining
elements of $B$ while fixing the last element of $A$, gives the first
inequality. For the equality assertion translate the least elements to
zero and write $A=\{0,x,y\}$ and $B=\{0,u,v\}$, where
$0<x<y$ and $0<u<v$. Each of the following strictly increasing lists
must be the complete ordered list of the same five sums:
\[
\begin{gathered}
 0,\ x,\ y,\ y+u,\ y+v,\\
 0,\ u,\ x+u,\ y+u,\ y+v,\\
 0,\ x,\ x+u,\ x+v,\ y+v.
\end{gathered}
\]
Comparison gives $u=x$, $y=2x$, and $v=2x$.
\end{proof}

\begin{theorem}\label{thm:C-torus15}
The least number of arrows, counting identities, of a finite category
whose classifying space has the homotopy type of $S^1\times S^1$
is exactly $15$.
\end{theorem}
\begin{proof}
\emph{The upper bound.}
There are three objects $0,1,2$, identity endomorphisms only, and
\[
\begin{aligned}
 \operatorname{Hom}(0,1)&=\{0,x,y\},&
 \operatorname{Hom}(1,2)&=\{0,x,y\},\\
 \operatorname{Hom}(0,2)&=\{0,x,y,2x,x+y,2y\}.&&
\end{aligned}
\]
Here $x,y$ are the standard basis of $\Z^2$ and composition is addition.
The copies of a label in different hom-sets are distinct arrows.
The only nondegenerate simplices above dimension one are the nine
triangles indexed by $(a,b)\in\{0,x,y\}^2$.
The three triangles $(0,0),(x,x),(y,y)$ have the respective free edges
$0,2x,2y$ in $\operatorname{Hom}(0,2)$. Collapse these three pairs.
There remain three vertices, nine edges, and six triangles.
Every edge belongs to two triangles and each vertex link is
$K_{3,3}$ with a perfect matching removed, a six-cycle.
Thus the remaining space is a connected closed surface.
Writing $0,x,y$ as $0,1,2$, give the triangles $(0,1),(1,2),(2,0)$
coefficient $+1$ and their reversals coefficient $-1$.
Their boundaries cancel, so the surface is orientable. Its Euler
characteristic is $3-9+6=0$, hence it is a torus.
This proves a homotopy equivalence for the original category; the
original classifying space itself includes the three collapsed flaps.

\emph{Passage to an acyclic category for the lower bound.}
Suppose $BC\simeq T^2$. Map $C$ to its fundamental groupoid and
take its image $D$. The image has finitely many arrows and no more
arrows than $C$. Every endomorphism has trivial image: this follows
either from the finite-dimensional theorem of Section~\ref{sec:A},
or directly because its image is a finite submonoid of the
torsion-free group $\Z^2$. Replace $D$ by a skeleton.
It is now acyclic and is a subcategory of a groupoid whose isotropy
groups are $\Z^2$. The map $BC\to BD$ induces an isomorphism
of fundamental groups. To check the last assertion without any
faithfulness assumption on $C$, use the universal groupoid property:
$C\to D\to\Pi C$ gives $\Pi C\to\Pi D\to\Pi C$ with
identity composite, and the other composite is also the identity
because every arrow of $D$ is the image of an arrow of $C$.

Choose $\alpha,\beta\in H^1(BD;\Q)$ pulled back from the standard
generators of $H^1(B\Z^2;\Q)$. The composite
$BC\to BD\to B\Z^2$ is a homotopy equivalence, since its source
and target are $K(\Z^2,1)$ spaces and it induces a fundamental-group
isomorphism. Consequently there is a rational two-cycle $z$ in
$ND$ with $\langle\alpha\beta,z\rangle\ne0$.
Choose such a cycle with the smallest possible number of supporting
triangles and let $L$ consist of its support and all faces.
The support is connected. Moreover $H_2(L;\Q)$ is one-dimensional:
if there were two independent cycles, a nonzero linear combination
$w$ would have zero pairing with $\alpha\beta$, and subtracting
a suitable multiple of $w$ from $z$ would cancel a triangle while
preserving the nonzero pairing. The restrictions of $\alpha,\beta$
to $L$ are independent, because the square of a rational class
of degree one is zero. Thus $b_1(L)\ge2$ and
$\chi(L)=2-b_1(L)\le0$.

Let $v,e,f$ be its numbers of vertices, edges, and triangles.
Each triangle has three distinct edges, since $D$ is acyclic.
Every supporting edge is in at least two supporting triangles, by
the nonzero coefficients of the cycle. Hence
\[
 3f\ge2e,\qquad 0\ge v-e+f\ge v-e/3,
 \qquad e\ge3v.
\]
In particular $C$ has at least $v+e\ge4v$ arrows.
If $C$ had at most fourteen arrows, the support would have exactly
three vertices, ordered $0<1<2$.

\emph{The three-object obstruction.}
In the full subcategory of $D$ on these three objects, choose
one arrow in each consecutive hom-set to identify the three objects
with a base object of the ambient groupoid. Its hom-sets are then
finite subsets $A,B,C_0$ of $\Z^2$, with $0\in A,B$,
$A+B\subseteq C_0$, and composition is addition.
If $|A|=1$, the fundamental group of this three-object category is
free, by eliminating its triangular relations. If $A=\{0,a\}$
with $a\ne0$, the same conclusion follows explicitly as follows.
After contracting the two chosen edges of a maximal tree, eliminate
the generator for $b\in B$ in favor of the composite edge labelled
$b\in C_0$. The remaining triangular relations have the form
\[
 c_{b+a}=s\,c_b\quad(b\in B),\qquad c_0=1,
\]
up to the harmless reversal of the convention for path products;
$s$ is the generator of the second arrow in $A$.
The graph with vertices $C_0$ and edges $b\to b+a$ has at most
one incoming and one outgoing edge at each vertex, and has no cycle:
a cycle would imply $na=0$ for some positive integer $n$.
It is therefore a disjoint union of finite paths. Each relation
eliminates one vertex generator along a path; the condition $c_0=1$
eliminates the remaining generator of the path containing $0$.
There are no remaining relations, so the fundamental group is free.
The same reasoning applies when $|B|\le2$.
Every map from such a category's classifying space to $B\Z^2$
factors, up to homotopy, through the classifying space of its free
fundamental group. It therefore induces zero in degree-two rational
homology, contrary to the cycle $z$.

It follows that $|A|,|B|\ge3$. Lemma~\ref{lem:C-small-sumset}
and the fourteen-arrow bound give
\[
 14\ge 3+|A|+|B|+|C_0|
 \ge 2(|A|+|B|)+2\ge14.
\]
All inequalities are equalities. Thus $|A|=|B|=3$ and
$C_0=A+B$ has five elements. The lemma shows that, after changing
the two chosen transport arrows, all labels lie in a single cyclic
subgroup of $\Z^2$. The map of this entire three-object category to
$B\Z^2$ factors through that cyclic subgroup, and again has zero
degree-two homology image. This final contradiction proves the
lower bound.
\end{proof}

\begin{remark}\label{rem:C-torus-certificate}
The independent program \texttt{c\_torus15.py} writes the complete
$15\times15$ typed composition table, checks all $66$ composable
triples, checks the three free faces and all vertex links, verifies
the oriented cycle, and computes the integral boundary Smith form.
As a separate check on the only small additive lemma used above,
it enumerates all $21147$ set partitions of the nine formal sums
of two three-element sets. Of the $18002$ partitions with at most
five blocks, $17772$ directly identify elements in a row or column,
and $212$ force such an identification by their rational linear
relations. The remaining $18$ all leave a one-dimensional quotient
of the four-dimensional vector space of the formal summand generators.
This finite calculation is not needed for the ordered-group proof,
but independently verifies its precise three-by-three case.
The literature search recorded in the notes found categorical torus
constructions, but no arrow-minimality statement used here; no priority
claim is made.
\end{remark}


\subsection{The four-sphere minimum and explicit three-dimensional models}
\begin{theorem}\label{thm:C-S4-nine}
The minimum number of arrows of a finite category of type $S^4$ is $9$.
\end{theorem}
\begin{proof}
The nine-element monoid and complete finite projective resolution in
Proposition~\ref{prop:B-S4-nine} prove the upper bound.
For monoids of smaller order, repeatedly apply the one-sided corner
reductions used in the census. Among all the irreducible monoids of
orders two through seven there is no table with
$H_1=H_2=H_3=0$; the lists contain respectively $1,1,2,2,7,25$
isomorphism-or-opposite classes, and their exact integral homology
is recorded in \texttt{c\_monoid\_homology.json}.
For a remaining nongroup irreducible monoid of order eight, the
minimal ideal is a Rees matrix semigroup $I\times H\times J$.
If $|I|=1$ or $|J|=1$, its maximal subgroup is a reducing corner,
contradicting irreducibility. Thus $|I|,|J|\ge2$.
Since the identity lies outside this proper ideal, its size is at
most seven. Consequently $H=1$ and, up to opposition, its dimensions
are $2\times2$ or $2\times3$.
The independently generated constrained tables give $243$ and $8$
irreducible isomorphism-or-opposite classes, respectively.
Here the raw completion counts with the band fixed in its stated
labels are $22746$ and $56$; these are not counts of all labelled
order-eight monoids.

The programs \texttt{c\_order8\_filter.py} and
\texttt{c\_order8\_survivors.py} exclude $244$ tables by mod-two
homology in degrees one through three. Each of the seven remaining
tables has $H_4(-;\F_2)=0$, also excluding $S^4$.
Groups are excluded by their nontrivial universal groups and
Section~\ref{sec:A}. This is a complete structural enumeration of
the order-eight irreducible cases, not an inference from a randomly
chosen family.

For categories with several objects the complete eight-arrow
enumeration has the counts in Proposition~\ref{prop:C-census}.
After exact isomorphism and opposition checks, the only candidate
with mod-two Betti numbers $(1,0,0,0,1)$ through degree four is
the categorical join of the five-element sphere monoid $P_2$ and
the two-element group $C_2$. A categorical join adds one arrow
from each object of the first category to each object of the second,
with its forced compositions. Thus this category has
$5+2+1=8$ arrows. Its nerve realizes the topological join, so
\[
 B(P_2*C_2)\simeq S^2*\RP^\infty
 \simeq\Sigma^3\RP^\infty.
\]
Its fourth integral homology is $\Z/2$, not $\Z$; the independent
integer Smith calculation in \texttt{c\_categories8\_residual.py}
also checks this directly.

This exact eight-arrow test covers all smaller multiobject categories
as well. To see the reduction, one may add a new identity at any
one object $x$, leaving its old identity as an idempotent $e$ that
acts as the identity on all old arrows. This adds exactly one arrow.
In the Karoubi envelope the full subcategory on $(x,e)$ and the other
original objects is the old category. The functor sending $x$ to
$(x,e)$ and every old arrow to itself is naturally related to the
canonical inclusion by the component $e$ at $x$ and identities
elsewhere. These functors prove that the added arrow preserves
the classifying-space homotopy type. Iterate to reach eight arrows.
A connected category with at most eight arrows has at most four
objects, because it needs at least $k-1$ nonidentity arrows for
$k$ objects. Hence every case reached is in the displayed census.
\end{proof}

\begin{proposition}\label{prop:C-lens-uppers}
There are explicit finite categories with four objects and $16$
arrows of type $\RP^3$, and with four objects and $20$ arrows of
type $L(3,1)$.
\end{proposition}
\begin{proof}
For the first use objects $0<1<2<3$, take
$\operatorname{Hom}(i,j)=C_2$ for every $i<j$, and compose by
addition; all endomorphisms are identities. Its covering poset
has objects $(i,g)$, $g\in C_2$, and a unique arrow for every
strict increase in $i$. Its order complex is the join of four
two-point sets, the cross-polytope boundary $S^3$.
Simultaneously adding $1$ to the binary coordinates is the antipodal
action. The orbit category is the stated category and each nerve
chain lifts uniquely after its first vertex is chosen. Therefore
its classifying space is homeomorphic to $S^3/C_2=\RP^3$.
The arrow count is $4+2\binom42=16$.

For the second again use objects $0<1<2<3$. In $C_3$ set
\[
 \operatorname{Hom}(0,1)=\operatorname{Hom}(2,3)=\{0,1\},
 \qquad \operatorname{Hom}(i,j)=C_3
 \quad\text{for the four other pairs }i<j,
\]
and compose by addition. The covering poset on $(i,g)$ has an
arrow precisely when the difference of its group coordinates lies
in the corresponding displayed hom-set. The subposet on levels
$0,1$ is the face poset of a triangular circle, as is the subposet
on levels $2,3$; every object of the first is less than every
object of the second. Its order complex is consequently
$S^1*S^1\cong S^3$. The deck generator rotates both circles
by the same $2\pi/3$ angle. In the usual join identification with
the unit sphere of $\mathbb C^2$, this is
$(z_1,z_2)\mapsto(\zeta_3z_1,\zeta_3z_2)$.
Its quotient is the lens space $L(3,1)$, with this explicit
convention for the lens action. The arrow count is
$4+2+2+4\cdot3=20$.

The program \texttt{c\_lens\_models.py} writes both complete typed
composition tables, checks every unit and associativity equation,
constructs and checks both covering posets and their unique lifts,
and computes integral homology through their full dimension three.
The geometric covering identifications prove the asserted homotopy
types; the homology check is an independent certificate.
The lower bounds for these two models are addressed separately;
the present statement asserts only the two explicit upper bounds.
\end{proof}


\subsection{Exact arrow minima for the three-dimensional lens examples}

\begin{lemma}[Reduction to labelled acyclic categories]
\label{C:lens-labelled-reduction}
Suppose that $C$ is a connected finite category of finite CW type with
$\pi_1(BC)\cong C_p$, where $p=2$ or $3$. There is a connected acyclic
category $D$ with at most as many arrows as $C$, with objects admitting
a linear order, and with
\[
 \operatorname{Hom}_D(i,j)=S_{ij}\subseteq C_p\quad(i<j),
 \qquad S_{ij}+S_{jk}\subseteq S_{ik}.
\]
Its only endomorphisms are identities, composition is addition of
labels, and $BC\to BD\to BC_p$ induces the classifying map of the
universal cover of $BC$.
\end{lemma}
\begin{proof}
Take the image category of $C$ in its fundamental groupoid as in
Lemma~\ref{A:groupoid-image}. Theorem~\ref{A:bounded-categories} makes
every endomorphism image an identity. Pass to a skeleton; the resulting
category $D$ is acyclic. Indeed arrows in both directions in a groupoid
are mutually inverse when all endomorphism images are identities, so
their endpoints were already isomorphic. Both taking the image and
passing to a skeleton weakly decrease the number of arrows.

Choose a base object and paths in the fundamental groupoid to each
object. They identify each ambient groupoid hom set with $C_p$ and
composition with addition. Since $D$ is a subcategory of this groupoid,
its hom sets are subsets, with the indicated closure property. The
factorization is the localization functor on classifying spaces, so
its composite is the universal-cover classifying map. A linear
extension of the object order gives the displayed indexing.
\end{proof}



\begin{lemma}[The exhaustive small-arrow calculation]
\label{C:lens-exhaustion}
Every connected acyclic $C_2$-labelled category with at most $15$ arrows
has zero pullback of $u^3\in H^3(BC_2;\mathbb F_2)$, where
$u\in H^1(BC_2;\mathbb F_2)$ is the nonzero class.
\end{lemma}
\begin{proof}
We describe a finite exhaustive calculation, including the reduction
that makes its search complete. A nonzero degree-three class requires
a chain of four objects. If there are $n$ objects and $m$ nonempty
strict hom sets, that chain contributes six comparisons on four
vertices. Connectedness requires at least $n-4$ further comparisons
to connect the other vertices. Hence
\[
 m\geq n+2,\qquad \#\operatorname{Mor}(D)\geq n+m\geq2n+2.
\]
Thus the $15$-arrow search needs only $4\leq n\leq6$. Categories
without a four-object chain have no nondegenerate three-simplices and
need no calculation.

For each $n$, the program \texttt{a\_lens\_lower.py} generates every
strict upper triangular transitive relation with at most $15-n$
comparisons. It does so by adjoining the vertices in order: the
predecessor set of each new vertex is every possible down-set of the
earlier relation. This generates every underlying poset exactly once
for its chosen linear order. It discards disconnected relations and
those without a four-object chain.

For a remaining relation, process its pairs $(i,j)$ in increasing
$j-i$. Every pair used to form a two-step composite has then already
been processed. If their assigned hom sets are $S_{ik},S_{kj}$, the
new hom set is allowed to be every nonempty subset of $C_2$ containing
\[
 \bigcup_{i<k<j}(S_{ik}+S_{kj}).
\]
The total cardinality of the hom sets is bounded by $15-n$. These
conditions are exactly the identity-free composition closure conditions;
associativity follows from addition in $C_2$.

An object gauge changes arrow labels by $a\mapsto a+c_j-c_i$.
When a processed comparison joins two components of the graph of
previously processed comparisons, choose the gauge constant on each
component so that the new hom set is the standard representative of
its translation orbit. Previous hom sets are unchanged. For $C_2$ the
representatives are $\{0\}$ and $C_2$; for later reference, the same
argument for $C_3$ has representatives $\{0\},\{0,1\},C_3$.
For either prime every nonempty subset of a fixed size lies in the
same translation orbit. The comparisons used in this normalization
form a spanning forest. At every other comparison the program allows
all subsets satisfying the closure condition. Consequently every
labelled category is gauge isomorphic to at least one category tested;
the normalization removes no possible cohomology class.

The calculation uses the normalized finite nerve. A $d$-simplex is an
increasing sequence of $d+1$ objects together with labels on its $d$
successive arrows. Its inner faces add the adjacent labels modulo $p$.
For labels in $\{0,\ldots,p-1\}$ the integral-lift formula for the
Bockstein is
\[
 b(a,b)=\left\lfloor\frac{a+b}{p}\right\rfloor\pmod p.
\]
For $p=2$, this is $ab$, and the degree-three cocycle representing
$u\beta u=u^3$ is $a\lfloor(b+c)/2\rfloor=abc$.
Exact Gaussian elimination first tests whether $b$ is a coboundary.
If it is, the degree-three product is zero. Otherwise the same exact
calculation tests whether $abc$ lies in the image of the coboundary
$C^2\to C^3$. No degree cutoff approximation is involved: both
cochain spaces and their full differential are finite and enumerated.

The complete output is
\[
\begin{array}{c|r|r|r|r}
 n&\text{underlying posets}&\text{categories tested}
   &\beta u\ne0&u^3\ne0\\\hline
 4&1&13&5&0\\
 5&22&171&0&0\\
 6&212&426&0&0
\end{array}
\]
The full code and result are supplied in
\path{code/a_lens_lower.py} and
\path{code/outputs/a_RP3_lower.json}; the companion output records
the counts and a SHA-256 digest of every tested labelled category and
its two decisions. As a positive control, the four-object category
with every strict hom set equal to $C_2$ has nonzero $u^3$, as expected
from its $\mathbb RP^3$ covering model in
Proposition~\ref{prop:C-lens-uppers}.
\end{proof}

\begin{theorem}[The exact minimum for $\mathbb RP^3$]
\label{C:RP3-minimum}
The least number of arrows in a finite category whose classifying space
is homotopy equivalent to $\mathbb RP^3$ is $16$.
\end{theorem}
\begin{proof}
The universal-cover classifying map
$\mathbb RP^3\to BC_2=\mathbb RP^\infty$ pulls $u^3$ back nontrivially.
For any category of this type, Lemma~\ref{C:lens-labelled-reduction}
factors this map through an acyclic labelled category with no more
arrows. Lemma~\ref{C:lens-exhaustion} excludes $15$ or fewer arrows.
Proposition~\ref{prop:C-lens-uppers} supplies the $16$-arrow category,
including its complete composition table and antipodal $S^3$ covering
poset.
\end{proof}

\begin{theorem}[The exact minimum for $L(3,1)$]
\label{C:L31-minimum}
The least number of arrows in a finite category whose classifying space
is homotopy equivalent to $L(3,1)$ is $20$.
\end{theorem}
\begin{proof}
Let $u\in H^1(BC_3;\mathbb F_3)$ be a generator. Its pullback to
$L(3,1)$ is nonzero. The Bockstein $\beta u$ is also nonzero: in the
coefficient sequence
$0\to\mathbb Z/3\to\mathbb Z/9\to\mathbb Z/3\to0$, the map
\[
 H^1(L(3,1);\mathbb Z/9)\longrightarrow
 H^1(L(3,1);\mathbb Z/3)
\]
is zero, because every homomorphism $C_3\to\mathbb Z/9$ takes values
in the subgroup $3\mathbb Z/9$. Both $H^1$ and $H^2$ over
$\mathbb F_3$ are one-dimensional. Poincar\'e duality for the oriented
closed three-manifold $L(3,1)$ therefore gives $u\beta u\ne0$ in
degree three.

Apply Lemma~\ref{C:lens-labelled-reduction}. A hypothetical category
with at most $19$ arrows gives a labelled acyclic category with that
bound and nonzero $u\beta u$. The same chain and connectedness count
as in Lemma~\ref{C:lens-exhaustion} gives $4\leq n\leq8$.
Run the exhaustive algorithm proved complete there with $p=3$ and
arrow budget $19$. All nonempty subsets of $C_3$ of a fixed cardinality
are translates, so its gauge-normalization proof applies unchanged.
Exact elimination over $\mathbb F_3$ tests the cochains
\[
 b(a,b)=\left\lfloor\frac{a+b}{3}\right\rfloor,
 \qquad c(a,b,c)=a\left\lfloor\frac{b+c}{3}\right\rfloor.
\]
The full result is
\[
\begin{array}{c|r|r|r|r}
 n&\text{underlying posets}&\text{categories tested}
   &\beta u\ne0&u\beta u\ne0\\\hline
 4&1&76&26&0\\
 5&22&4001&183&0\\
 6&539&54346&0&0\\
 7&9890&170653&0&0\\
 8&37905&97739&0&0
\end{array}
\]
The supplied output \path{code/outputs/a_L31_lower.json} records
all counts, the zero-counterexample result, and the deterministic
enumeration digest. The same program verifies nonzero $u\beta u$ on
the $20$-arrow positive control from
Proposition~\ref{prop:C-lens-uppers}. Thus $19$ arrows are impossible,
and that proposition supplies a $20$-arrow category together with its
explicit $C_3$ covering poset homeomorphic to $S^3$.
\end{proof}


\subsection{The category lower bound for $\mathbb CP^2$}

\begin{lemma}[Two-object cup-product reductions]
\label{C:two-object-cup-pruning}
Let $C$ be a finite category with two objects and endomorphism monoids
$M,N$. Suppose first that all cross arrows go in one direction, and
write $S$ for that nonempty cross-biset.
\begin{enumerate}
\item If $|S|=1$, every positive-degree integral cohomology product of
      $BC$ vanishes.
\item If $|S|=2$ and one of $BM,BN$ is contractible, every product of
      two positive-degree integral classes of $BC$ is annihilated by two.
      If that vertex space is only rationally acyclic, every such
      rational product vanishes.
\item If both $BM$ and $BN$ are rationally acyclic, every positive-degree
      rational product of $BC$ vanishes, for arbitrary $S$.
\end{enumerate}
If arrows instead go in both directions and one endomorphism monoid is
trivial, $BC$ is homotopy equivalent to the classifying space of the
other endomorphism monoid. If each of the two cross-hom sets has one
element, $BC$ is contractible.
\end{lemma}
\begin{proof}
The collage homotopy-pushout description expresses $BC$ as
\[
 BM\ \cup^h_{E_S}\ BN,
\]
where $E_S$ is the two-sided double bar on the cross-biset. A singleton
biset has $E_S\simeq BM\times BN$, with the two projections, giving
$BC\simeq BM*BN$. A join of connected based spaces is homotopy
equivalent to a suspension, so its positive-degree products vanish.
If both vertex spaces are rationally acyclic, their maps to points give
a rational homology equivalence from the displayed homotopy pushout
to the unreduced suspension of $E_S$, by Mayer--Vietoris. The induced
rational cohomology ring isomorphism proves the third assertion.

For the second assertion, identify $S$ with the endpoints of the
interval $I=[0,1]$. Every self-map of the two-point set extends uniquely
to an affine map of $I$. The commuting left and right monoid actions
therefore extend to commuting interval actions, each of which is the
identity, reflection, or a constant endpoint map. The corresponding
double bar $E_I$ is homotopy equivalent over $BM\times BN$ to that
product, because $I$ is contractible. Assume $BN$ is contractible.
Collapsing that vertex and using the equivalence $E_I\to BM$ gives
\[
 BC\simeq\operatorname{hocofib}(E_S\to BM)
       \simeq\operatorname{hocofib}(E_S\to E_I).
\]
The last space is the reduced double bar with fiber
$I/\partial I=S^1$, with its basepoint bar collapsed. These comparisons
can be made as ordinary quotient comparisons: at each bar level the
endpoint inclusion is a CW cofibration, and all degeneracies are
cofibrant inclusions.

On the circle the monoid actions are identity, complex conjugation,
or the constant map to the basepoint. They all commute strictly with
$z\mapsto z^2$. This supplies a self-map $F$ of the quotient double
bar. In its reduced cellular chain complex the internal circle complex
is $\mathbb Z_\chi$ concentrated in degree one, where each monoid
element acts by $0$, $1$, or $-1$. The circle power map acts by
multiplication by two on this complex. Hence its map on the total
reduced bar chain complex is literally multiplication by two in every
degree, and $F^*$ is multiplication by two on every positive-degree
integral cohomology group. For positive classes $a,b$, multiplicativity
now gives
\[
 2ab=F^*(ab)=F^*(a)F^*(b)=4ab.
\]
Thus $2ab=0$.

If $BN$ is only rationally acyclic, the same two comparison maps are
rational homology equivalences: use Mayer--Vietoris for the first and
K\"unneth followed by the cofiber sequence for the second. Their maps
on rational cohomology are ring isomorphisms. The reduced circle-bar
argument therefore still gives vanishing of every positive rational
product. No acyclicity assertion from a finite range of homology is
used here.

For arrows in both directions and a trivial endomorphism monoid at
the second object, any pair of opposite cross arrows composes to its
identity. That object is a retract of the first. The first object is
therefore Cauchy dense, and its full-subcategory inclusion induces the
claimed equivalence on classifying spaces.

Finally suppose both cross-hom sets are singletons, with arrows
$i:1\to0$ and $r:0\to1$. Their composites $ir$ and $ri$ are zero
elements in $M,N$, respectively: multiplication by any endomorphism
leaves the singleton cross arrow unchanged. Split the idempotent $ir$
in the Karoubi category. The new object has exactly one arrow to and
from either old object, and only its identity as endomorphism; the
cross arrows identify its alternative presentation by $ri$.
It is consequently a zero object. Adding a split idempotent preserves
the classifying-space type, which proves contractibility.
\end{proof}

\begin{lemma}[Small vertex monoids in a triangular category]
\label{C:small-triangular-vertices}
No triangular two-object category with at most ten arrows has
classifying space $\mathbb CP^2$.
\end{lemma}
\begin{proof}
Write $m,n$ for its two endomorphism orders and $s\geq1$ for the
number of cross arrows. Thus $m+n+s\leq10$. Every monoid of order
at most four has finite-group homotopy type by
Theorem~\ref{E:all-six-types}, so its classifying space is rationally
acyclic. Lemma~\ref{C:two-object-cup-pruning} disposes of $s=1$,
of $s=2$ (when $\min(m,n)\leq4$), and of every case with both
$m,n\leq4$. Up to interchanging the objects, the remaining maximal
profiles are
\[
 (m,n,s)=(6,1,3),\quad(5,1,4),\quad(5,2,3).
\]
We prove statements valid for arbitrary cross-biset size in these
three endomorphism ranges.

First suppose $n=1$, so the category is the homotopy cofiber of
$E_S=B(S\rtimes M)\to BM$. If $H^2(BM;\mathbb Q)=0$, every
degree-two class of that cofiber lifts from its cohomology relative
to $BM$. The quotient by $BM$ is a suspension of $E_S$, so the
product of any two such lifted relative classes, and hence of the
original classes, is zero. The same argument applies when $BN$ is
merely rationally acyclic, by the rational cofiber comparison in
Lemma~\ref{C:two-object-cup-pruning}.

The complete order-six classification leaves only the following
monoids with nonzero rational $H^2$: the five-element sphere monoid,
its fresh-identity extension, the additional six-element sphere band,
the two regular $C_2$-band monoids of types
$S^2\vee\Sigma^2\mathbb RP^\infty$ and
$\mathbb RP^\infty/\mathbb RP^1$, and the six-element Sweeney monoid.
For all but the last, the trivial rational module has projective
dimension at most two. The sphere monoid and the additional band have
the explicit length-two resolutions already checked in
Theorem~\ref{E:all-six-types}. A fresh identity splits off an additional
rational ring factor and preserves that resolution. For either regular
group-band monoid, put $R=\mathbb QM$ and
$A=R/\mathbb QK=\mathbb QC_2$. The ideal $\mathbb QK$ is projective,
so $\operatorname{pd}_R A\leq1$. The row augmentation module $V$ is
an $A$-module and is projective over the semisimple ring $A$.
Hence $\operatorname{pd}_R V\leq1$, and the exact sequence
$0\to V\to Re\to\mathbb Q\to0$ gives
$\operatorname{pd}_R\mathbb Q\leq2$. The same statements hold for
opposites; the sphere bands and their fresh-identity extension admit
the row-column interchange, and the regular-band argument is symmetric.

For any finite right $M$-set $S$, the action category has homology
\[
 H_q(E_S;\mathbb Q)
  =\operatorname{Tor}^{\mathbb QM}_q(\mathbb QS,\mathbb Q).
\]
Thus $H_3(E_S;\mathbb Q)=0$ in these cases. Since
$H_4(BM;\mathbb Q)=0$, the cofiber also has zero fourth rational
homology and cannot be $\mathbb CP^2$.

For the Sweeney monoid use the complete resolution in
Theorem~\ref{thm:E-wedge}. Its modules $P_q$ are $P\oplus R$ for
every $q\geq2$, and its differentials satisfy
$d_{q+2}=d_q$ for $q\geq3$. The cofiber chain complex can be calculated by the algebraic
mapping cone of the natural augmentation map
\[
 \mathbb QS\otimes_R P_*\longrightarrow
 \mathbb Q\otimes_R P_*.
\]
Its chain modules and differentials are exactly two-periodic in the
range needed to compute $H_q$ for $q\geq4$. Consequently
$H_4(BC;\mathbb Q)\cong H_6(BC;\mathbb Q)$ for every finite $S$.
This excludes $\mathbb CP^2$, independently of any homology cutoff.

Finally consider $m\leq5$, $n=2$. The rational algebra of either
two-element monoid, $C_2$ or $\{1,0\}$, is semisimple. If the
larger vertex has zero rational $H^2$, the relative-cup argument above
applies. Otherwise that monoid is the five-element sphere monoid.
The double-bar edge space is the homotopy-orbit space of the
$(N,M)$-biset, and its homology is computed over
$\mathbb QN\otimes\mathbb QM^{\mathrm{op}}$. Tensoring the
length-two sphere resolution with the projective trivial
$\mathbb QN$-module gives projective dimension at most two for the
trivial module of this product algebra. The edge space again has
zero third rational homology; the two vertex spaces have zero fourth
rational homology. Mayer--Vietoris for the collage then gives
$H_4(BC;\mathbb Q)=0$. This treats every remaining profile.
\end{proof}

\begin{lemma}[Three-object star reductions]
\label{C:star-cup-pruning}
Suppose a three-object category has cross arrows only on two legs of
a star, both pointing toward its center or both pointing away from it.
If both cross-hom sets are singletons, all positive cohomology products
vanish. If one leg is a singleton with a trivial endomorphism monoid
at its leaf, and the other leaf has rationally acyclic classifying
space, all positive rational products vanish.
\end{lemma}
\begin{proof}
Write $X$ for the center's classifying space and $Y,Z$ for those of
the leaves. The nerve is the union of the two leg collages along $X$;
there are no composable pairs of cross arrows. For singleton legs
these collages are $X*Y$ and $X*Z$. Both inclusions of $X$ are
null-homotopic, by choosing a point in the respective nonempty leaf
space. Their homotopy pushout is therefore
\[
 (X*Y)\vee(X*Z)\vee\Sigma X.
\]
Each summand is a suspension, proving the first assertion.
For the second, the singleton leg ending at a point is a cone on $X$.
Collapsing this contractible collage leaves the homotopy pushout
$*\leftarrow E\to Z$, where $E$ is the other edge's double bar.
If $Z$ is rationally acyclic, its map to a point makes this pushout
rationally homology equivalent to the unreduced suspension of $E$,
with a rational cohomology ring isomorphism. This proves the claim.
\end{proof}

\begin{lemma}[Triangular categories with eleven arrows]
\label{C:triangular-eleven}
No triangular two-object category with at most eleven arrows has
classifying space $\mathbb CP^2$.
\end{lemma}
\begin{proof}
The preceding reductions leave only the new maximal profiles
\[
 (m,n,s)=(7,1,3),(6,2,3),(5,3,3).
\]
Cases with a smaller vertex of
order one or two and a larger vertex of order at most six or five,
respectively, were treated for arbitrary $s$ in
Lemma~\ref{C:small-triangular-vertices}.
For $(6,2,3)$ the same proof extends: the two-element monoid algebra
is semisimple over $\mathbb Q$, so tensoring with its projective
trivial module preserves the projective-dimension bound or the
Sweeney resolution's exact two-periodic tail. If the six-element
vertex has zero rational $H^2$, the relative-cup argument applies.

For $(7,1,3)$ we exhaust the finite three-point actions, using a small
complete structural list of target monoids. A seven-element monoid
with nonzero rational $H^2$ must have at least two minimal rows and
columns: otherwise the one-sided corner equivalence reduces it to a
finite group. Its minimal group is then trivial by the size bound,
and its minimal band is $2\times2$, $2\times3$, or $3\times2$.
The latter two cases are the band with an identity adjoined.
In the $2\times2$ case there are exactly two further nonidentity
elements. Each is determined on the band by a row transformation and
a column transformation of a two-point set. By the degree-two
balanced tensor formula, their two characters must agree for rational
$H^2$ to be nonzero. There are six such pairs of transformations.
The four products of the two extra elements may be any table entries
whose row and column translations equal the corresponding composite
translations. The program tries every such choice and checks all
$7^3$ associativity equations.

This yields $229$ normalized tables and $41$ ordinary isomorphism
classes, with exact row and column permutations and interchange of
the two extra elements used to recognize isomorphisms. Adding the
two nonsquare bands gives $43$ target monoids. For each, the program
chooses a minimal generating set and assigns every possible
transformation of a three-point set to each generator. It extends
these assignments along fixed generator words and checks all monoid
multiplication equations. Exact state permutations remove only
isomorphic actions. This exhausts $572$ right-action cases; taking
opposite monoids, already in the list, also covers the other arrow
orientation. The normalized category cup computation gives $194$
zero squares, $203$ failures of the required second Betti number,
and $173$ failures of the first Betti number. The two remaining
nonzero-square cases both have third mod-two Betti number one.
The complete generation and certificates are
\texttt{a\_CP2\_11\_triangular.py} and its JSON output.

For $(5,3,3)$, the five-element vertex must be the sphere monoid:
every other five-element monoid has zero rational $H^2$, so is
excluded by the relative-cup argument. All three-element vertex
monoids are rationally acyclic. The complete order-three census
contains seven ordinary isomorphism classes after opposites are
restored. The program \texttt{a\_CP2\_11\_bisets.py} enumerates all
three-point right actions of the sphere monoid and all left actions
of each three-element monoid, and keeps exactly the commuting pairs.
For completeness, after fixing the right action up to state
permutation it includes all conjugates of each left action, so it
does not discard a possible relative positioning of the actions.
There are $277$ resulting commuting cases. Of these, $54$ have
$b_1=0,b_2=1$, and every one has an explicit degree-three primitive
for the square. The other $223$ fail $b_1$ or $b_2$.
This completes the three profiles and the proof.
\end{proof}

\begin{proposition}[No $\mathbb CP^2$ category with at most eleven arrows]
\label{C:CP2-lower10}
\label{C:CP2-lower11}
\label{C:CP2-lower12}
Every finite category whose classifying space has the homotopy type of
$\mathbb CP^2$ has at least twelve arrows.
\end{proposition}
\begin{proof}
The one-object case is covered by the stronger monoid exclusion in
Proposition~\ref{prop:B-fourdim-lower14}. For a connected category with
$n$ arrows and $k\geq2$ objects, at least $k-1$ nonidentity arrows are
needed to connect the objects, so $k\leq(n+1)/2$. The complete
multiobject enumeration from Section~\ref{sec:C} therefore supplies
all cases through eight arrows, and all nine-arrow cases with
$3\leq k\leq5$.

For nine arrows and two objects, the additional program
\texttt{a\_categories\_cp2.cpp} uses the same complete sorted-hom-set
enumeration and associativity propagation as
\texttt{c\_categories.cpp}. It discards a hom-cardinality profile only
when it guarantees an initial or terminal object. Such a category is
contractible. At every remaining associative table, it computes the
exact normalized-nerve coboundary ranks over $\mathbb F_2$ and retains
the table precisely when $b_1=0$ and $b_2=1$. All potentially relevant
tables are therefore retained. The computation examines $8\,964\,544$
connected labelled tables after the initial/terminal-profile exclusion,
using $264\,734\,681$ propagation nodes, and retains $399\,484$ tables.
There are $776$ retained classes up to isomorphism and opposition.
As an independent check on this added filter, running it at eight arrows
produces exactly the same set of $7132$ candidate tables as the previously
completed eight-arrow census, not merely the same count.

The same complete engine at ten arrows retains respectively
$309\,960$, $6048$, and $480$ labelled tables for $k=3,4,5$.
It examines respectively $2\,591\,868$, $51\,356$, and $3900$
associative tables after the initial/terminal-profile exclusion.
These runs are complete; the two-object ten-arrow case is treated
separately below using the preceding structural lemmas.

For every retained category, \texttt{a\_category\_cups.py} constructs
the full normalized nerve in degrees through four, finds a generator
$x\in H^2(-;\mathbb F_2)$, and tests $x\smile x$ against the full image
of the degree-three coboundary. Every zero square is accompanied by an
explicit degree-three cochain primitive, checked on every normalized
four-simplex. Exact arrow relabelings and opposition alone are used to
reuse a computation; no hash invariant is treated as an isomorphism
test. The relevant counts are
\[
\begin{array}{c|c|r|r|r}
 n&k&b_1=0,\ b_2=1\text{ classes}&x^2=0&x^2\ne0\\\hline
 8&2&72&72&0\\
 8&3&23&23&0\\
 8&4&1&1&0\\
 9&2&776&773&3\\
 9&3&150&149&1\\
 9&4&8&8&0\\
 9&5&0&0&0\\
 10&3&1238&1233&5\\
 10&4&72&72&0\\
 10&5&3&3&0\\
 11&2&1&1&0\\
 11&3&2870&2842&28\\
 11&4&545&544&1\\
 11&5&22&22&0\\
 11&6&0&0&0
\end{array}
\]
All smaller multiobject cases have zero square whenever $b_1=0,b_2=1$;
their individual counts are recorded in
\texttt{a\_categories\_CP2\_small\_summary.json}.

At nine arrows, each of the three two-object residual categories has
$H^3(-;\mathbb F_2)\cong\mathbb F_2$, by the exact ranks of the full
degree-two and degree-three coboundaries. The single three-object
residual category has $H_3(-;\mathbb Z)\cong\mathbb Z/2$, and also
nonzero third mod-two homology. At ten arrows the five three-object
residual categories have third mod-two Betti numbers
$2,2,1,2,1$, respectively. Thus none of the residual categories
has the homology of $\mathbb CP^2$. Every other category in the table was excluded
by $b_1$, $b_2$, or the fact that the square of the degree-two generator
of $H^*(\mathbb CP^2;\mathbb F_2)$ is nonzero.

For two objects and ten arrows,
Lemma~\ref{C:small-triangular-vertices} removes every triangular
profile. Lemma~\ref{C:two-object-cup-pruning} removes profiles having
an endomorphism monoid of order one, by reduction to a smaller
monoid, and profiles with singleton hom sets in both directions,
by contractibility. Object exchange and category opposition allow
the remaining profiles to be normalized by
\[
 |\operatorname{End}(0)|\geq|\operatorname{End}(1)|,
 \qquad |\operatorname{Hom}(0,1)|\geq
        |\operatorname{Hom}(1,0)|.
\]
These are independent operations on the two cardinality pairs;
opposition does not change endomorphism cardinalities. The complete
program \texttt{a\_categories\_cp2\_strong.cpp} thus removes $106$
profiles by this symmetry and $39$ by the proved structural
reductions. Exactly $11$ connected profiles remain. Exhaustive
associativity propagation visits $1\,872\,751$ nodes and yields
$19\,983$ associative labelled tables. None has $b_1=0,b_2=1$.
The full output is
\texttt{a\_categories\_10\_2\_CP2\_strong.out}; its candidate
table file has only the header because the candidate set is empty.
This completes every ten-arrow case.

At eleven arrows, the program \texttt{a\_categories11\_cp2.cpp}
uses the same exhaustive propagation, with the proved reductions of
Lemmas~\ref{C:two-object-cup-pruning}, \ref{C:star-cup-pruning},
and \ref{C:triangular-eleven}. For two objects it removes $145$
profiles by object exchange and opposition and $48$ by the structural
lemmas. The $17$ remaining connected profiles produce $1\,002\,072$
associative labelled tables after $133\,327\,715$ propagation nodes.
Exactly $24$ have $b_1=0,b_2=1$; they are one isomorphism/opposition
class and have zero square. For three objects the star lemma removes
$384$ profiles, in addition to the initial/terminal-profile exclusion.
The complete counts at each object number are
\[
\begin{array}{c|r|r|r}
 k&\text{associative labelled tables}&\text{propagation nodes}
    &b_1=0,b_2=1\text{ labelled tables}\\\hline
 2&1\,002\,072&133\,327\,715&24\\
 3&28\,311\,846&626\,523\,753&1\,727\,652\\
 4&1\,030\,548&6\,178\,468&104\,016\\
 5&50\,470&132\,350&5040\\
 6&2580&2580&0
\end{array}
\]
These counts refer to the remaining profiles, so they do not silently
count structurally excluded categories as generated tables. The exact
cup-square calculation gives the eleven-arrow rows in the preceding
table. Every one of the $28$ three-object residuals has nonzero third
mod-two homology (its dimension is one, two, or three); the one
four-object residual has third mod-two Betti number one. The program
\texttt{a\_category\_residual\_ranks.py} records all normalized chain
dimensions and boundary ranks establishing these assertions. Thus
every eleven-arrow category is excluded, which proves the bound.

All complete candidate tables, normalized cocycles, explicit primitives,
and residual rank certificates are supplied in the
\texttt{a\_categories\_*\_CP2\_cups.json} files. Preliminary
nine- and ten-arrow streams marked \texttt{partial} were superseded
by the complete filtered searches and are not used for any
completeness claim.
\end{proof}

\subsection{A structural obstruction for every act of a monoid of order at most seven}

\begin{lemma}[The homology of a triangular act cone]
\label{C:act-cone-Tor}
Let $M$ be a finite monoid, let $S$ be a nonempty finite right
$M$-set, and let $\mathcal C(M,S)$ have two objects, endomorphism
monoids $M$ and $\{1\}$, and $\Hom(0,1)=S$, with composition given
by the action. Put $R=\mathbb QM$ and
$L=\ker(\mathbb QS\xrightarrow{\epsilon}\mathbb Q)$.
For every $n\geq1$ there is a natural isomorphism
\[
 \widetilde H_n(B\mathcal C(M,S);\mathbb Q)
 \cong\operatorname{Tor}^{R}_{n-1}(L,\mathbb Q).
\]
The same conclusion, with $L$ replaced by its coinvariants for the
other vertex monoid, holds for a triangular biset category with
endomorphism monoids $M,N$, provided the trivial left
$\mathbb QN^{\mathrm{op}}$-module is projective.
\end{lemma}
\begin{proof}
The mixed simplices in the nerve give the homotopy pushout
\[
 BM\ \longleftarrow\ E_S\ \longrightarrow\ *,
 \qquad E_S=B(S,M,*).
\]
Choose a projective resolution $P_\bullet\to\mathbb Q$ over $R$.
The edge space and the vertex space are computed, naturally with
respect to their map, by
$\mathbb QS\otimes_R P_\bullet$ and
$\mathbb Q\otimes_R P_\bullet$.
Since each $P_q$ is projective, tensoring the exact sequence
$0\to L\to\mathbb QS\to\mathbb Q\to0$ gives a short exact
sequence of complexes. The mapping cone of its last map is
quasi-isomorphic to $(L\otimes_R P_\bullet)[1]$.
The degree-zero augmentation contributes the one unreduced
$\mathbb Q$ in the classifying space and no positive homology.
This proves the displayed formula, including $n=1$.

For a biset, use $R\otimes\mathbb QN^{\mathrm{op}}$ and the
two-sided bar. Projectivity of the trivial module for the second
factor makes its coinvariant functor exact. Thus a projective
resolution for the first factor tensored with that trivial module
is a projective resolution for the product algebra, and taking
the second coinvariants gives the stated formula. In particular
$BN$ is rationally acyclic. Replacing its vertex in the double
mapping cylinder by a point is a rational homology equivalence,
as follows directly from the homotopy-pushout exact sequence.
\end{proof}

\begin{proposition}[No projective-plane act cone over at most seven elements]
\label{C:seven-act-cones}
For every finite monoid $M$ with $|M|\leq7$ and every nonempty
finite right $M$-set $S$, the category $\mathcal C(M,S)$ does not
have classifying space of type $\mathbb CP^2$.
There is no bound on $|S|$ in this assertion. The same exclusion
holds for a triangular biset category with other endomorphism
monoid $N$ whenever the trivial left
$\mathbb QN^{\mathrm{op}}$-module is projective.
\end{proposition}
\begin{proof}
If $H^2(BM;\mathbb Q)=0$, every degree-two class in the
triangular category lifts to cohomology relative to $BM$.
After replacing the rationally acyclic other vertex by a point,
the quotient by $BM$ is a suspension. Relative products therefore
vanish, and so do the products of these degree-two classes in
absolute cohomology. This excludes $\mathbb CP^2$.

For $|M|\leq6$ and $H^2(BM;\mathbb Q)\ne0$, the proof of
Lemma~\ref{C:small-triangular-vertices} supplies either a projective
resolution of the trivial rational module of length two, or an
exact resolution with modules periodic from degree two and
differentials periodic from degree three. That proof does not
bound the size of the act. We now verify the same sufficient
obstruction for every seven-element target, with one additional
syzygy pattern.

The classification of the $43$ seven-element monoids with
nonzero rational $H^2$ is the complete translation-table
enumeration in Lemma~\ref{C:triangular-eleven}. In particular it
includes reducible monoids as well as the irreducible cores; no
homotopy equivalence of monoids is silently used to transfer an
arbitrary act. Number these tables $0,\ldots,42$ in the
deterministic order of \texttt{a\_CP2\_11\_triangular.targets()}.
The exact resolution computation gives the following partition:
\[
\begin{array}{c|r|l}
 \text{property}&\text{count}&\text{table indices}\\\hline
 \operatorname{pd}_{\mathbb QM}\mathbb Q\leq2&27&
 3,4,5,6,8,9,10,11,14,15,16,17,18,19,22,24,\\
 &&27,28,29,32,33,34,35,36,37,41,42\\
 \text{two-periodic resolution tail}&11&
 1,2,7,13,21,23,26,31,38,39,40\\
 \Omega^3\mathbb Q\cong(\Omega^2\mathbb Q)^{\oplus2}&4&0,12,25,30\\
 \text{central }C_2\text{ splitting}&1&20
\end{array}
\]
Here is what is checked, so that these are all-degree certificates.
The first $27$ resolutions are integral resolutions by summands
$\mathbb ZM e$ and $\mathbb ZM(e-f)$, with
$e^2=e$, $f^2=f$, $ef=fe=f$. Every differential is equivariant,
every image is saturated over $\mathbb Z$, and the successive
rank sums give exactness. The final degree-two differential is
injective. The separate program
\texttt{b\_verify\_resolutions.py} checks these facts independently
of the resolution search.

For the $11$ periodic cases, the projectives in every degree from
two onward are the same explicit module. The matrices satisfy
$d_3=d_5$, and $d_4,d_5$ form an exact two-periodic complex:
both products are zero, both images are saturated, and their ranks
sum to the rank of the common projective. Thus the displayed
finite initial segment continues to an exact resolution with
$d_{q+2}=d_q$ for every $q\geq3$.
For every right module $L$ it follows that
\[
 \operatorname{Tor}_{q+2}^{\mathbb QM}(L,\mathbb Q)
 \cong\operatorname{Tor}_{q}^{\mathbb QM}(L,\mathbb Q)
 \qquad(q\geq3).
\]

For each of the four replication cases, the program
\texttt{a\_seven\_cone\_syzygies.py} separately verifies the
projective initial segment through $P_2$, including saturation
and exactness at $P_0,P_1$. It computes the actual kernels
$\Omega^2\mathbb Q$ and $\Omega^3\mathbb Q$, of dimensions
four and eight, with the action of every monoid element.
It then supplies an $8$ by $8$ rational intertwining matrix
from two copies of the former module to the latter. The
determinants, in index order $0,12,25,30$, are respectively
$-2,-6,-2,8$. Every intertwining equation is verified exactly;
all matrices and module actions are in the JSON certificate.
Dimension shifting therefore gives, for every right module $L$,
\[
 \operatorname{Tor}_{q+1}^{\mathbb QM}(L,\mathbb Q)
 \cong
 \operatorname{Tor}_{q}^{\mathbb QM}(L,\mathbb Q)^{\oplus2}
 \qquad(q\geq3).
\]

Finally, table $20$ has a central unit $g=6$ of order two which
acts as the identity on every nonunit. Over $\mathbb Q$ the
central idempotents $(1+g)/2$ and $(1-g)/2$ split its algebra
as the product of the six-element Sweeney monoid algebra and
$\mathbb Q$. The first factor has basis its new identity,
the four band elements, and the remaining nonunit; multiplication
is exactly the six-element table. The trivial module is supported
on this first factor. Its Tor groups with an arbitrary module
therefore have the same periodicity just proved for that monoid.

In every case, Lemma~\ref{C:act-cone-Tor} shows that nonzero
fourth rational homology in the triangular category either is
impossible or persists in arbitrarily high degrees: the periodic
cases have $H_{n+2}\cong H_n$ for $n\geq4$, and the replication
cases have $H_{n+1}\cong H_n^{\oplus2}$ for $n\geq4$.
This is incompatible with the rational homology of
$\mathbb CP^2$. Exactness of the second-vertex coinvariants
proves the biset version by the same Tor formulas.

The program \texttt{a\_seven\_cone\_targets.py} produces the
complete resolution records and the partition. The replication matrices
are in \texttt{a\_seven\_cone\_syzygies.json}. The computation
does not treat failure of a resolution search to terminate as
proof of an infinite tail: each used tail has the displayed
exact periodic or module-isomorphism certificate.
\end{proof}

\subsection{The obstruction for every act of an eight-element monoid}

\begin{lemma}[A stable retract propagates Tor]
\label{C:stable-retract-Tor}
Let $R$ be a finite-dimensional algebra over a field, let $Q$ be
a finite-dimensional left module, and choose projective syzygies
$U=\Omega^2Q$ and $V=\Omega^{2+s}Q$, where $s\geq1$.
Suppose there are homomorphisms
\[
 F_i:U\longrightarrow V,\quad G_i:V\longrightarrow U,
 \qquad
 A_j:U\longrightarrow P_j,\quad B_j:P_j\longrightarrow U,
\]
with every $P_j$ projective, such that
\[
 1_U=\sum_i G_iF_i+\sum_j B_jA_j.                         \tag{\ref{C:stable-retract-Tor}.1}
\]
Then, for every right module $L$ and every $n\geq3$,
\[
 \operatorname{Tor}_n^R(L,Q)\ne0
 \quad\Longrightarrow\quad
 \operatorname{Tor}_{n+s}^R(L,Q)\ne0.
\]
If the first sum is empty, $U$ is projective and all these
degree-$n$ groups are zero.
\end{lemma}
\begin{proof}
Apply the additive functor $\operatorname{Tor}_q^R(L,-)$ for
$q\geq1$. Every composite through a $P_j$ becomes zero.
Consequently $\operatorname{Tor}_q^R(L,V)=0$ forces the identity
of $\operatorname{Tor}_q^R(L,U)$ to be zero. Dimension shifting,
with $q=n-2$, gives the implication. If the first sum is empty,
the identity of $U$ factors through $\bigoplus_jP_j$, so $U$
is a projective direct summand of that module.
\end{proof}

\begin{lemma}[A derived local Frobenius comparison]
\label{C:derived-Frobenius-Tor}
Let $R$ be a finite-dimensional $k$-algebra and let $A$ be a
quotient algebra which is local and Frobenius, with unique simple
module $k$ and $\operatorname{pd}_R A\leq1$ on the left.
Suppose a finite-dimensional left $R$-module $Q$ fits into
\[
 0\longrightarrow k\longrightarrow P\longrightarrow Q
 \longrightarrow0,
\]
where $P$ is projective and the first module is obtained by
restriction from the simple $A$-module. For every finite-dimensional
right $R$-module $L$, eventual vanishing of
$\operatorname{Tor}_n^R(L,Q)$ implies its vanishing for every
$n\geq3$.
\end{lemma}
\begin{proof}
We first record the needed fact about complexes over $A$.
Let $C$ be a derived $A$-module with finite-dimensional
cohomology concentrated in degrees zero and one. If
\[
 \operatorname{Ext}_A^n(k,C)=0\qquad\text{for all sufficiently large }n,
                                                               \tag{\ref{C:derived-Frobenius-Tor}.1}\label{eq:C-Frobenius-eventual-Ext}
\]
then these groups vanish for every $n\geq2$.

To prove this fact, choose a bounded-below injective resolution
$I^\bullet$ of $C$, starting in degree zero, with finite-dimensional
terms. Such a resolution is obtained by resolving a bounded
finite-dimensional representative, or by successive injective
envelopes. Since $A$ is local and Frobenius, its finite-dimensional
injectives are finite sums of the unique indecomposable injective
$A$. We may make the complex minimal directly. A unit matrix
entry in a differential can be isolated by invertible row and
column operations, splitting off a contractible complex
$A\xrightarrow{1}A$. The neighboring differential components
are zero by the complex equations. Remove these summands
successively by degree. Only finitely many can be removed at
each degree, so the resulting bounded-below complex is defined
degree by degree, remains an injective resolution of $C$, and
has every differential entry in $J=\operatorname{rad}(A)$.

These differentials kill the socles of the injective terms:
the socle of a local Frobenius algebra is annihilated on both
sides by its radical.
Thus the complex $\operatorname{Hom}_A(k,I^\bullet)$ has zero
differential. Each nonzero injective term has a nonzero socle,
so \eqref{eq:C-Frobenius-eventual-Ext} forces $I^n=0$ for all
sufficiently large $n$.
If the last nonzero term were $I^N$ with $N>1$, the equality
$H^N(C)=0$ would make $I^{N-1}\to I^N$ surjective. Since
$I^N$ is also projective, that surjection would split off
another contractible summand, contrary to minimality.
Equivalently, a matrix with entries in $J$ cannot surject onto
a nonzero free module. Therefore $I^n=0$ for every $n\geq2$,
proving the fact.

Now put $N=\operatorname{Hom}_k(L,k)$ and form the full derived
right adjoint
\[
 C=\operatorname{RHom}_R(A,N).
\]
The hypothesis on the left projective dimension of $A$ implies
that $H^j(C)=0$ outside $j=0,1$; its cohomology is finite-dimensional.
Restriction of modules along $R\to A$ is exact and has right
adjoint $\operatorname{Hom}_R(A,-)$. Applying this adjunction
to an injective resolution of $N$ gives
\[
 \operatorname{RHom}_R(k,N)
 \simeq
 \operatorname{RHom}_A\bigl(k,\operatorname{RHom}_R(A,N)\bigr).
                                                               \tag{\ref{C:derived-Frobenius-Tor}.2}\label{eq:C-Frobenius-derived-adjunction}
\]
Indeed, the right adjoint sends injective $R$-modules to
injective $A$-modules, so this is an equality of the appropriate
derived Hom computations. No ordinary change-of-rings Tor
formula is being assumed.

Duality of finite-dimensional vector spaces and the displayed
short exact sequence give
\[
 \operatorname{Hom}_k\bigl(\operatorname{Tor}_n^R(L,Q),k\bigr)
 \cong\operatorname{Ext}_R^n(Q,N)
 \cong\operatorname{Ext}_R^{n-1}(k,N)\qquad(n\geq2).
\]
Eventual vanishing therefore gives
\eqref{eq:C-Frobenius-eventual-Ext} for $C$ by
\eqref{eq:C-Frobenius-derived-adjunction}.
The proved minimal-injective-complex fact makes
$\operatorname{Ext}_R^m(k,N)=0$ for every $m\geq2$.
The same dimension shift and faithful vector-space duality
give $\operatorname{Tor}_n^R(L,Q)=0$ for all $n\geq3$.
\end{proof}

\begin{proposition}[No projective-plane act cone over eight elements]
\label{C:eight-act-cones}
For every finite monoid $M$ of order at most eight and every
nonempty finite right $M$-set $S$, the category $\mathcal C(M,S)$
of Lemma~\ref{C:act-cone-Tor} does not have classifying space
homotopy equivalent to $\mathbb CP^2$.
There is no bound on the cardinality of $S$.
The same exclusion holds for a triangular biset category with
other endomorphism monoid $N$ whenever the trivial left
$\mathbb QN^{\mathrm{op}}$-module is projective.
\end{proposition}
\begin{proof}
Orders at most seven are covered by
Proposition~\ref{C:seven-act-cones}. If
$H^2(BM;\mathbb Q)=0$, the relative-suspension argument in
that proposition makes every product of degree-two classes
zero. It remains to consider eight-element monoids with
nonzero rational second cohomology.

We use the complete own translation-table enumeration
\path{code/c_eight_H2_targets.cpp}, whose construction is
spelled out in the proof of Lemma~\ref{C:twelve-triangular}.
The minimal ideal must be a rectangular band of shape
$2\times2$, $2\times3$, or $3\times2$: fewer rows or columns
give group type, and a nontrivial local group with two rows
and columns would already require at least nine elements.
For the $2\times2$ shape, nonzero rational $H_2$ forces the
row and column augmentation characters of every extra element
to agree. The complete list of compatible translations and
all associative products gives 492 ordinary isomorphism
classes. The other two shapes give twelve classes each.
This produces the 516 tables indexed $0,\ldots,515$ in
\path{code/outputs/c_eight_H2_targets.tsv}; both orientations
are present. Exact rational balanced-module ranks show that
the six indices
\[
 496,498,499,506,512,513
\]
have $H_2(BM;\mathbb Q)=0$. Retaining them as harmless extra
entries, rather than dropping them, gives the following exact
partition of the entire 516-table list:
\[
\begin{array}{l|r}
 \text{certificate}&\text{number of tables}\\ \hline
 \text{original resolution of length at most two}&261\\
 \text{original two-periodic resolution tail}&157\\
 \text{split cancellation in an original length-three resolution}&5\\
 \text{additional projective second syzygy}&9\\
 \text{stable retract with recurrence step one}&65\\
 \text{stable retract with recurrence step two}&15\\
 \text{derived local Frobenius comparison}&4
\end{array}                                                     \tag{\ref{C:eight-act-cones}.1}\label{eq:C-eight-partition}
\]
The complete index-by-index partition, with the hash of every
underlying resolution certificate, is
\path{code/outputs/a_eight_cone_complete_partition.json}.
All resolutions and syzygies in this proof are for the actual
algebra $R=\mathbb QM$ and its trivial module $\mathbb Q$.
No equivalence of monoid classifying spaces is used to transfer
an arbitrary act.

For the first 261 cases, the program
\path{code/b_verify_resolutions.py} independently verifies the
complete integral projective resolution: each projective is
$\mathbb ZM e$ or $\mathbb ZM(e-f)$ for commuting idempotents
with $ef=fe=f$, each differential is equivariant, every image
is saturated, and the successive rank sums prove exactness.
The final differential is injective. In the 157 periodic cases
the same checker verifies the initial segment, equality of the
projectives from degree two onward, and the two zero matrix
products and rank sums that close the two-periodic tail.
Thus, for every right module $L$, either
$\operatorname{Tor}_n^R(L,\mathbb Q)=0$ for $n\geq3$, or
\[
 \operatorname{Tor}_{n+2}^R(L,\mathbb Q)
 \cong\operatorname{Tor}_n^R(L,\mathbb Q)\qquad(n\geq3).
\]
The exact indices, tables, verified ranks, and certificate
hashes are recorded in
\path{code/outputs/c_eight_cone_resolutions.json}.

The five split cancellations occur at indices
$182,307,382,407,437$. The last projective is $R\varepsilon$,
where $\varepsilon=g-f$, and its differential is literally
inclusion into a summand $R\eta$ of $P_2$, where $\eta=1-h$.
In the same index order, the triples $(g,f,h)$ are
\[
 (6,1,7),\ (7,2,6),\ (5,1,6),\ (6,3,7),\ (7,4,6).
\]
Direct multiplication gives
$\varepsilon^2=\varepsilon$, $\eta^2=\eta$, and
$\varepsilon\eta=\eta\varepsilon=\varepsilon$.
Hence $R\eta=R\varepsilon\oplus R(\eta-\varepsilon)$.
Cancel the contractible last summand to obtain a resolution
of length at most two. These integral ring equations, the
literal generator image, and the original complete resolutions
are independently checked by
\path{code/c_eight_split_cancellations.py}, with all five
records in its JSON output.

For the remaining syzygy calculations we use actual kernels
in the saved projective initial segments. The search
\path{code/a_eight_cone_syzygies.py} computes their full
$M$-actions and expresses $1_{\Omega^2\mathbb Q}$ as in
Lemma~\ref{C:stable-retract-Tor}, with all projective factors
equal to the free regular module $R$.
For the nine indices
\[
 32,137,371,450,459,463,464,497,501,
\]
the identity factors entirely through free modules, proving
that $\Omega^2\mathbb Q$ is projective.
For fifteen indices the later syzygy is $\Omega^4\mathbb Q$:
\[
 14,30,52,125,135,143,165,249,261,295,345,358,369,401,433.
\]
For the other 65 successful stable identities it is
$\Omega^3\mathbb Q$. Each certificate stores every intertwining
matrix, its rational coefficient, the two complete module
actions, and the exact identity obtained by summing the
composites. The independent verifier
\path{code/a_eight_cone_verify.py} imports neither the resolution
generator nor the syzygy search. It reconstructs each
projective as the image of its actual idempotent in $\mathbb QM$,
checks the kernel bases and exactness, and verifies every
intertwining equation and all 89 identities over $\mathbb Q$.
Consequently nonzero $\operatorname{Tor}_3^R(L,\mathbb Q)$
persists in arbitrarily high degrees in all 80 recurrence cases.

The four remaining indices are $9,10,16,18$.
Write $m_j$ for the element indexed $j$, so $m_0$ is the
global identity. Here $K=\{m_1,m_2,m_3,m_4\}$ is the minimal
$2\times2$ band, and
the quotient $A=\mathbb QM/\mathbb QK$ has basis $1,x,y,c$
with $c=m_5$, $x=m_6$, and $y=m_7$ in the table coordinates.
Its multiplication is exactly that of
Lemma~\ref{B:binary-pencil-Tor}, with respective matrices
\[
 \begin{pmatrix}0&1\\1&0\end{pmatrix},\quad
 \begin{pmatrix}0&1\\1&1\end{pmatrix},\quad
 \begin{pmatrix}1&0\\0&1\end{pmatrix},\quad
 \begin{pmatrix}1&0\\1&1\end{pmatrix}.
\]
Each is invertible, so the coefficient-of-$c$ pairing proves
that $A$ is local Frobenius. The ideal $\mathbb QK$ is the
direct sum of its two column modules, $Rm_1$ and $Rm_2$.
Both are projective, so
$\operatorname{pd}_R A\leq1$.
The augmentation $Rm_1\to\mathbb Q$ has kernel spanned by
$m_3-m_1$. Every nonidentity element kills this vector, whereas
the global identity, indexed zero, fixes it. Thus
$\Omega_R\mathbb Q$ is precisely the simple $A$-module.
All these assertions are direct multiplication-table equations
and are separately checked for each of the four records in
\path{code/outputs/a_eight_cone_verified.json}.
Lemma~\ref{C:derived-Frobenius-Tor} therefore says that eventual
Tor vanishing forces $\operatorname{Tor}_3^R(L,\mathbb Q)=0$
in these last cases as well.

The rows of \eqref{eq:C-eight-partition} now show, for every
finite-dimensional right
$R$-module $L$, that nonzero degree-three Tor is incompatible
with bounded Tor. Apply Lemma~\ref{C:act-cone-Tor} with
$L=\ker(\mathbb QS\to\mathbb Q)$. A projective-plane act cone
would have nonzero rational fourth homology and zero homology
in every higher degree, which is impossible.
The biset statement follows by the exact second-vertex
coinvariant argument in that lemma.
\end{proof}

\subsection{An all-order obstruction for two minimal rows}

\begin{lemma}[Finite affine acts on a line]\label{C:finite-line-act-cups}
Let $M$ be a monoid and let $T\subset\mathbb Q$ be a nonempty finite
right $M$-set whose action maps have the form
\[
 t\longmapsto\chi_m t+\beta_m,
 \qquad \chi_m\in\{0,1,-1\}.
\]
Every positive-degree rational cohomology product of the act-cone
category $\mathcal C(M,T)$ is zero.
\end{lemma}
\begin{proof}
For a singleton $T$, the edge bar is $BM$ and its projection to
$BM$ is the identity, so the cone is contractible.  Suppose $T$
has at least two elements, and put $I=[\min T,\max T]$.
An affine map of slope $1$ preserving a finite subset of the line
must be the identity: a nonzero translation cannot permute that
finite set.  A map of slope $-1$ is a bijection of $T$, so its
translation term is necessarily $\min T+\max T$.  Thus all such
maps are the same reflection of $I$.  The remaining maps are
constants at points of $T$.  The given action consequently extends
to an actual affine right action on the contractible interval $I$.
Its multiplication identities hold on $I$, since they hold as
affine identities on the at least two points of $T$.

The levelwise weak equivalence from the interval bar to $BM$ gives
\[
 B\mathcal C(M,T)
 \simeq\operatorname{hocofib}\bigl(B(T,M,*)\longrightarrow BM\bigr)
 \simeq\operatorname{hocofib}\bigl(B(T,M,*)
                                  \longrightarrow B(I,M,*)\bigr).
                                                               \tag{\ref{C:finite-line-act-cups}.1}
\]
The last cofiber is the pointed bar with fiber $I/T$, with its
basepoint bar collapsed.  This follows level by level from the
cofibration $T\subset I$; realization preserves this cofiber.
Writing the elements of $T$ in increasing order identifies $I/T$
with a wedge of $|T|-1$ circles, one for each consecutive interval.
Each action map is the identity, the common reflection, or the
constant pointed map.  The reflection permutes these circle
summands and reverses their orientations.

For an integer $d\geq2$, take the degree-$d$ circle map on every
summand, using its linear interval parameter modulo its endpoints.
It commutes strictly with the reflection: on the circle parameter,
the two expressions are $d(1-t)$ and $1-dt$, equal modulo one.
It also commutes with every constant map and the identity.
It therefore induces a self-map $F_d$ of the pointed bar in
(\ref{C:finite-line-act-cups}.1).

The reduced cellular chain complex of $I/T$ is concentrated in
degree one, with a free generator for each consecutive interval.
All action maps are cellular, and $F_d$ is multiplication by $d$
on this entire reduced chain complex.  The normalized pointed-bar
cellular complex thus has the same literal multiplication-by-$d$
map in every degree.  Hence $F_d^*$ is multiplication by $d$ on
every positive-degree rational cohomology group.  For positive
classes $u,v$, naturality of the cup product gives
\[
 d\,uv=F_d^*(uv)=F_d^*(u)F_d^*(v)=d^2uv.
\]
Since $d^2-d\ne0$ in $\mathbb Q$, this proves $uv=0$.
\end{proof}

\begin{theorem}[Two minimal rows exclude every act-cone square]
\label{C:two-row-act-cups}
Let $M$ be a finite monoid whose minimal Rees matrix ideal has at
most two rows.  For every nonempty finite right $M$-set $S$, all
products
\[
 H^2(B\mathcal C(M,S);\mathbb Q)\otimes
 H^2(B\mathcal C(M,S);\mathbb Q)
 \longrightarrow H^4(B\mathcal C(M,S);\mathbb Q)
\]
are zero.  In particular no such act-cone category has the homotopy
type of $\mathbb{CP}^2$.  There is no restriction on the order of
$M$, the size of $S$, or the maximal group of the minimal ideal.
The opposite assertion holds for left acts and at most two minimal
columns.
\end{theorem}
\begin{proof}
Put $R=\mathbb QM$.  Choose a minimal idempotent $e$ and write
$H=eMe$ for its maximal group.  The averaging element
\[
 \bar e=|H|^{-1}\sum_{h\in H}h
\]
is an idempotent of $R$.  The projective left module $P=R\bar e$
has one basis vector for each row of the minimal ideal: $Me$ is
the disjoint union of the corresponding free right $H$-orbits,
and multiplying by $\bar e$ averages each orbit.  Left
multiplication by every $m\in M$ sends each such basis vector to
another row basis vector.  Its augmentation $P\to\mathbb Q$
therefore sends each basis vector to $1$.

If there is only one row, this augmentation is an isomorphism.
The trivial module is projective, and
Lemma~\ref{C:act-cone-Tor} gives $H^2(B\mathcal C(M,S);\mathbb Q)=0$.
Suppose there are two rows.  The kernel of the augmentation is
the one-dimensional row-difference module, giving an exact sequence
\[
 0\longrightarrow\mathbb Q_\chi\longrightarrow P
      \longrightarrow\mathbb Q\longrightarrow0,
 \qquad \chi:M\longrightarrow\{0,1,-1\}.             \tag{\ref{C:two-row-act-cups}.1}
\]
Here $\chi_m$ is the action of the two-point row transformation on
its augmentation line: it is $1$ for the identity, $-1$ for the
interchange, and $0$ for a constant.  In particular $\chi$ is an
actual multiplicative character, including its zero values.

Let $L_S=\ker(\mathbb QS\to\mathbb Q)$, a right $R$-module,
and write $D L_S=\operatorname{Hom}_{\mathbb Q}(L_S,\mathbb Q)$
with the contragredient left action.  The natural act-cone formula
and finite-dimensional vector-space duality give
\[
 H^2(B\mathcal C(M,S);\mathbb Q)
 \cong D\operatorname{Tor}_1^R(L_S,\mathbb Q)
 \cong\operatorname{Ext}_R^1(\mathbb Q,D L_S).        \tag{\ref{C:two-row-act-cups}.2}
\]
The second isomorphism follows by dualizing the complex obtained
from a projective resolution of the left module $\mathbb Q$;
it is natural in every map of right acts $S\to T$.
Applying $\operatorname{Hom}_R(-,D L_S)$ to
(\ref{C:two-row-act-cups}.1), and using projectivity of $P$, gives
a surjection
\[
 \operatorname{Hom}_R(\mathbb Q_\chi,D L_S)
 \longrightarrow H^2(B\mathcal C(M,S);\mathbb Q).     \tag{\ref{C:two-row-act-cups}.3}
\]
Thus every degree-two class has a representative consisting of a
linear functional $\ell:L_S\to\mathbb Q$ satisfying
\[
 \ell(v m)=\chi_m\ell(v)\qquad(v\in L_S, m\in M).  \tag{\ref{C:two-row-act-cups}.4}
\]
No assertion that $L_S$ itself is annihilated by the minimal ideal
is required.

Fix $s_0\in S$ and define
\[
 q(s)=\ell([s]-[s_0]),\qquad T=q(S)\subset\mathbb Q.
\]
For every $s\in S$ and $m\in M$, equation
(\ref{C:two-row-act-cups}.4) gives
\[
 q(sm)=\chi_m q(s)+q(s_0m).                           \tag{\ref{C:two-row-act-cups}.5}
\]
Consequently $q:S\to T$ is a map of finite right $M$-sets when
$T$ is given these affine action maps.  This is well defined even
when $q$ is not injective.  The functional
\[
 \ell_T\left(\sum_{t\in T}a_t[t]\right)
       =\sum_{t\in T}a_t t,
 \qquad \sum_{t\in T}a_t=0,
\]
is $\chi$-equivariant by (\ref{C:two-row-act-cups}.5), since the
translation term is multiplied by $\sum a_t=0$.
Moreover $\ell=\ell_T\circ L(q)$.

Let $x$ be the original degree-two class, and let $y$ be the
class for $\mathcal C(M,T)$ obtained from $\ell_T$ by the same
connecting homomorphism (\ref{C:two-row-act-cups}.3).
Naturality of that connecting map and of
(\ref{C:two-row-act-cups}.2) proves
\[
 x=(B\mathcal C(M,q))^*y.
\]
The target act is affine on a finite rational subset of a line.
Lemma~\ref{C:finite-line-act-cups} gives $y^2=0$, and hence $x^2=0$.
This holds for every $x$.  Finally rational polarization,
$(x+z)^2-x^2-z^2=2xz$, proves that every product of two degree-two
classes is zero.  Reversing the category interchanges right acts
with left acts and rows with columns, proving the opposite statement.
\end{proof}

\begin{corollary}\label{C:two-row-biset-cups}
The same exclusion applies to a triangular biset category whose
first endomorphism monoid has at most two minimal rows and whose
other endomorphism monoid $N$ has projective trivial rational module
on the side used for its coinvariants.
\end{corollary}
\begin{proof}
Take the actual finite orbit act $S/N$.  Its quotient map from the
cross-biset is equivariant for the first monoid and defines a functor
from the biset category to $\mathcal C(M,S/N)$.  In every degree of
the first monoid's bar, exact $N$-coinvariants identify the second
bar rationally with the orbit set.  Thus the map of edge bars is a
rational homology equivalence.  The vertex map $BN\to *$ is also
a rational homology equivalence.  The mixed-nerve homotopy-pushout
comparison gives a rational homology equivalence of the complete
classifying spaces, and therefore an isomorphism of their rational
cohomology rings.  Theorem~\ref{C:two-row-act-cups} applies to the
target.  This is an actual quotient-functor comparison; no finite-type
property of an individual vertex or edge space is inferred.
\end{proof}

% New structural comparison, separate from any blanket lower bound.
\subsection{Removing an opposite corner with projective coinvariants}

\begin{lemma}[Projective trivial modules and the minimal column]
\label{C:projective-coinvariants-minimal-column}
Let $N$ be a finite monoid with minimal Rees ideal $I\times G\times J$.
The trivial right $\mathbb QN$-module is projective if and only if
$|J|=1$. In this case a minimal idempotent $e$ satisfies
$eN=eNe=G$, and
\[
 q_N:N\longrightarrow G,\qquad n\longmapsto en
\]
is a unital monoid homomorphism, with $e$ the identity of $G$.
Over a field of characteristic $p$, the same trivial right module
is projective if and only if $|J|=1$ and $p\nmid|G|$.
\end{lemma}
\begin{proof}
A splitting of the right augmentation is an element $a\in kN$
with augmentation one and $an=a$ for every $n\in N$.
For each minimal idempotent $f$ it follows that $a=af\in kNf$.
Distinct minimal columns give disjoint basis supports for these
left ideals. Their intersection is zero, contradicting the
augmentation of $a$. Thus there is only one minimal column.

For such an ideal, $eN$ is the group $eNe$.
Consequently $(en)e=en$ and
$(en)(em)=enm$, proving the homomorphism assertion.
If $|G|$ is invertible in $k$, the element
$|G|^{-1}\sum_{g\in G}g$ is right invariant under $N$, since
right multiplication by $n$ acts on this sum as right multiplication
by $en\in G$. Its augmentation is one, giving the splitting.

Conversely, any invariant element belongs to $kK$, since $a=ae$.
Its coefficients on each fixed-row $G$-orbit are constant under
right multiplication by $G$. Its augmentation is therefore a
sum of multiples of $|G|$. In characteristic $p$ dividing $|G|$
this cannot be one.
\end{proof}

\begin{lemma}[The equivariant action-groupoid comparison]
\label{C:opposite-corner-action-groupoid}
Let $N$ satisfy the preceding lemma, and let $S$ be a left
$N$-set with a commuting right action by a monoid $M$.
Let $S\mathbin{/\mkern-6mu/}N$ denote its left action category.
There is an actual right-$M$-equivariant functor
\[
 p:S\mathbin{/\mkern-6mu/}N
       \longrightarrow eS\mathbin{/\mkern-6mu/}G,
 \qquad s\longmapsto es,\quad n\longmapsto en,
\]
which induces a weak equivalence of classifying spaces.
Moreover $N\backslash S$ is the finite right $M$-set
$T=G\backslash eS$.
\end{lemma}
\begin{proof}
The action category has an arrow labelled $n$ from $s$ to $ns$.
The displayed functor preserves its target because
$(en)(es)=ens$, and preserves compositions because $q_N$ is a
homomorphism. It commutes strictly with the right $M$-action.

For $t\in eS$, an object of $(t\downarrow p)$ is a pair $(s,g)$
with $g\in G$ and $es=gt$. The pair $(t,e)$ is one such object.
There is a morphism from $(s,g)$ to $(t,e)$ labelled $g^{-1}\in N$:
indeed $g^{-1}s=g^{-1}es=t$ and $e g^{-1}g=e$.
For an arrow labelled $n$ from $(s,g)$ to $(s',g')$, the comma
equation is $g'=(en)g$. Hence
\[
 (g')^{-1}n=(g')^{-1}(en)=g^{-1}.
\]
These arrows give a natural transformation from the identity of
the comma category to its constant functor at $(t,e)$.
The component at $(t,e)$ need not be the identity of $N$;
this is immaterial to the contraction of its nerve.
Thus every comma category is contractible, and Quillen's theorem A
proves the weak equivalence.

Every $s$ is related to $es$ under the $N$-orbit equivalence.
For $n\in N$, the states $ens=(en)(es)$ and $es$ lie in the same
$G$-orbit. This gives inverse identifications of the two orbit
sets, compatible with their right $M$-actions.
\end{proof}


\begin{theorem}[Finite-target and good-prime corner removal]
\label{C:finite-target-opposite-corner}
Let $C$ be the triangular two-object category with endomorphism
monoids $M,N$, cross-biset $S$ with left $N$-action and right
$M$-action, and no reverse cross arrows.
Assume the trivial right $\mathbb QN$-module is projective.
Put $T=N\backslash S$ and let
\[
 Q:C\longrightarrow C(M,T)
\]
be the actual quotient functor, equal to the identity on $M$,
sending $N$ to the identity at the other object and $S$ to $T$.
Then:
\begin{enumerate}
\item $BQ$ induces an isomorphism of cohomology rings over
$\mathbb Q$, and over $\mathbb F_p$ whenever $p\nmid|G|$.
\item For every connected finite CW complex $X$, precomposition is
a weak equivalence of derived unbased mapping spaces
\[
 \operatorname{Map}(BC(M,T),X)
       \longrightarrow \operatorname{Map}(BC,X).
\]
\item If $BC$ has finite CW homotopy type, there is a homotopy
left inverse $h:BC(M,T)\to BC$ to $BQ$.
\end{enumerate}
The size of the act category is $|M|+|T|+1\leq|M|+|S|+1$.
No assertion that it has finite CW type or the same integral
homology as $BC$ is included.
\end{theorem}
\begin{proof}
Use the mixed-nerve homotopy-pushout description
\[
 BC\simeq\operatorname{hocolim}
 \bigl(BM\longleftarrow E\longrightarrow BN\bigr),
 \qquad
 E=B\bigl(B(S\mathbin{/\mkern-6mu/}N),M,*\bigr).
\]
Lemma~\ref{C:opposite-corner-action-groupoid} is an actual
$M$-equivariant weak equivalence. Applied to the one-point left
$N$-set, the same lemma gives $BN\simeq BG$ through $Bq_N$.
Thus this is equivalently the span with edge
$B(B(eS\mathbin{/\mkern-6mu/}G),M,*)$ and other vertex $BG$.

The groupoid $eS\mathbin{/\mkern-6mu/}G$ has one component for each
$t\in T$. Each component is equivalent to $BG_t$ for a finite
stabilizer subgroup $G_t\leq G$. The map to the component set is
a rational homology equivalence and an $\mathbb F_p$ homology
equivalence when $p\nmid|G|$. Applying the first-quadrant bar
homology spectral sequence to the remaining $M$-bar preserves
these equivalences. The same assertion holds for $BG\to *$.
The map of the two homotopy-pushout spans consequently proves (1);
because the map is the actual functor $Q$, it is an isomorphism
of rings, not merely of graded vector spaces.

For (2) use the full mapping-space form of Miller's theorem A
\cite[Theorem A, p.~39]{miller1984}, with its corrected proof
\cite{miller1985}. For every finite group $H$ and finite CW
complex $X$ the based mapping space $\operatorname{Map}_*(BH,X)$
is weakly contractible. Evaluation and its constant-map section
therefore give
\[
 X\xrightarrow{\simeq}\operatorname{Map}(BH,X).
\]
This is the full mapping-space assertion, not just the
fundamental-group consequence used elsewhere.

The canonical constant-map inclusion
\[
 \operatorname{Map}(T,X)\longrightarrow
 \operatorname{Map}\bigl(B(eS\mathbin{/\mkern-6mu/}G),X\bigr)
\]
is therefore a weak equivalence. It is natural for the actual
right $M$-action; no equivariant choice of component basepoints or
of contractions is needed. Mapping out of a homotopy colimit is
a homotopy limit of these mapping spaces. Equivalently, use the
simplicial $M$-bar and its derived totalization. Levelwise there
are only copies of the preceding weak equivalence, so its
homotopy-limit comparison is a weak equivalence. The corresponding
constant-map inclusion for $BG\to *$ is also a weak equivalence.
Mapping the two pushout spans into $X$ and taking their homotopy
pullbacks proves (2).

Choose a finite CW model $X$ of $BC$. Surjectivity on components
of the mapping-space equivalence lifts the chosen equivalence
$BC\to X$ through $BQ$, up to homotopy. Composing that lift with
an inverse equivalence $X\to BC$ gives (3).
\end{proof}


\begin{corollary}[The precise projective-plane consequences]
\label{C:finite-target-corner-CP2}
In the setting of the preceding theorem, suppose $BC\simeq\mathbb CP^2$.
Write $Y=BC(M,T)$ and let $u$ be an integral generator of
$H^2(BC;\mathbb Z)$. Then $Y$ is simply connected and has the
rational and good-prime cohomology rings of $\mathbb CP^2$.
There is a class $v\in H^2(Y;\mathbb Z)$ with
\[
 (BQ)^*v=u,\qquad (BQ)^*(v^2)=u^2.
\]
The group $H^2(Y;\mathbb Z)$ is infinite cyclic, $v$ is a generator,
and $v^2$ is primitive in the free quotient of $H^4(Y;\mathbb Z)$.
For every finite group $H$ and every map $a:BH\to Y$,
the composite $h a:BH\to BC$ is null-homotopic, and
$a^*v=0$ integrally.
\end{corollary}
\begin{proof}
The quotient functor is surjective on every hom set and fixes the
two objects, so it is surjective on fundamental groups. Since
$BC$ is simply connected, so is $Y$.
Part (1) of the theorem gives its asserted rational and
good-prime cohomology.

Take $v=h^*u$. The homotopy equality $hBQ\simeq1$ proves the
displayed identities integrally. A finite category has finitely
generated homology in each degree. Since $H_1(Y;\mathbb Z)=0$,
the universal coefficient theorem identifies $H^2(Y;\mathbb Z)$
with the dual of $H_2(Y;\mathbb Z)$; its rational rank is one.
It is therefore infinite cyclic. The split surjection $(BQ)^*$
makes $v$ its generator. The same split surjection in degree four,
whose rational rank is one, makes $v^2$ primitive modulo torsion.
Finally Miller's full theorem applies to $h a$ with finite target
$BC\simeq\mathbb CP^2$. Its nullity gives
$a^*v=(h a)^*u=0$.
\end{proof}

\begin{corollary}[A trivial minimal group gives an actual equivalence]
\label{C:trivial-group-opposite-corner}
In Theorem~\ref{C:finite-target-opposite-corner}, suppose the minimal
local group of $N$ is trivial. Then the actual quotient
$BC\to BC(M,N\backslash S)$ is a homotopy equivalence.
Consequently, if $BC\simeq\mathbb CP^2$, then
\[
 |C|\ \geq\ 21+|N|+|S|-|N\backslash S|.
\]
In particular a nontrivial such opposite corner cannot occur in a
projective-plane category with at most $22$ arrows.
\end{corollary}
\begin{proof}
The groupoid in Lemma~\ref{C:opposite-corner-action-groupoid} is
now the discrete set $eS=N\backslash S$, and $BN$ is contractible.
Thus the actual map of mixed-nerve pushout spans in the proof of
Theorem~\ref{C:finite-target-opposite-corner} is an objectwise
homotopy equivalence. This proves the first assertion without a
finite-target or coefficient comparison.
The exact act-cone minimum in
Theorem~\ref{C:act-cone-minimum-twenty-two} gives
$|M|+|N\backslash S|+1\geq22$. Subtracting this size from
$|C|=|M|+|N|+|S|$ gives the stated inequality.
\end{proof}

\begin{lemma}[Finite-target detection after an orbit quotient]
\label{C:finite-target-orbit-wedge}
Let $X=B\mathcal C(M,S)$ satisfy
Lemma~\ref{C:two-column-projective-orbit-kernel}, and suppose $X$
is simply connected. Suppose a map $h:X\to F$ to a simply
connected finite CW complex induces an isomorphism on integral
second cohomology. Write $Y=B\mathcal C(M,\bar S)$ and let $o$
be the number of core orbits. Then there is a map
\[
 Y\longrightarrow F\vee\bigvee_{o-1}S^2
\]
inducing an isomorphism on integral second cohomology.
In particular every integral second-cohomology class on $Y$
restricts to zero along every map $BH\to Y$ from a finite group.
\end{lemma}
\begin{proof}
Corollary~\ref{C:two-column-orbit-wedge} supplies a homotopy
equivalence $X\vee\bigvee_{o-1}S^2\to Y$ extending the actual
quotient on $X$. Compose its homotopy inverse with
$h\vee1:X\vee\bigvee_{o-1}S^2\to F\vee\bigvee_{o-1}S^2$.
The wedge decomposition and the hypothesis on $h^*$ prove the
cohomology assertion. The target is a finite CW complex, so
Miller's theorem makes every composite from $BH$ null-homotopic.
Since every second-cohomology class on $Y$ pulls back from this
target, all the stated restrictions vanish.
\end{proof}

\begin{remark}[Scope of the smaller act cone]
\label{C:finite-target-corner-scope}
Corollary~\ref{C:finite-target-corner-CP2} gives a finite-source
nullity condition for
every local-group restriction of the distinguished class $v$.
It does not identify the integral homotopy type of $Y$.
In particular a lower bound that assumes $Y$ itself has finite
CW type cannot be applied merely from this comparison.
Universal rational cup primitives and certified field cup
obstructions do apply through (1); any imported finite-type hull
argument must separately check that its use of finiteness is
exactly the local-group nullity supplied by the corollary.
\end{remark}




\subsection{Split row modules exclude rank-one act-cone squares}

\begin{proposition}\label{C:split-row-act-cups}
Let $M$ be a finite monoid, and let $V$ be the rational augmentation
module of the permutation action on the rows of its minimal ideal.
Suppose $V$ is a direct sum of one-dimensional rational $M$-modules.
For every nonempty finite right $M$-set $S$, the space
$H^2(B\mathcal C(M,S);\mathbb Q)$ is spanned by classes whose
squares are zero.  In particular, if this cohomology space is
one-dimensional, its generator has square zero.  Thus no such
act-cone category realizes $\mathbb{CP}^2$, for any size of $S$.
\end{proposition}

\begin{proof}
As in Theorem~\ref{C:two-row-act-cups}, averaging a maximal group in
the minimal ideal gives a projective module $P$ and an exact sequence
\[
 0\longrightarrow V\longrightarrow P\longrightarrow\mathbb Q
                       \longrightarrow0.
\]
Write $V=\bigoplus_i\mathbb Q_{\chi_i}$, with projections $\pi_i$.
Each rational character takes values in $\{0,1,-1\}$: its finite
multiplicative image in $\mathbb Q$ has no other nonzero elements.
Put $L_S=\ker(\mathbb QS\to\mathbb Q)$.  The same natural
Tor--Ext identification and connecting homomorphism give a surjection
\[
 \bigoplus_i\operatorname{Hom}_{\mathbb QM}
       (\mathbb Q_{\chi_i},D L_S)
 \cong\operatorname{Hom}_{\mathbb QM}(V,D L_S)
 \longrightarrow H^2(B\mathcal C(M,S);\mathbb Q).
\]
It suffices to consider a class contributed by one summand.  Its
representative is a scalar functional $\ell:L_S\to\mathbb Q$
with $\ell(vm)=\chi_i(m)\ell(v)$, extended to a map on $V$ by
the projection $\pi_i$.

Fix $s_0\in S$, put $q(s)=\ell([s]-[s_0])$, and let $T=q(S)$.
Exactly as in the two-row proof,
\[
 q(sm)=\chi_i(m)q(s)+q(s_0m).
\]
Thus $q:S\to T$ is a map of finite affine acts on a rational line.
The corresponding functional on $L_T$, extended by the same
projection $\pi_i$, defines a class $y$ for $\mathcal C(M,T)$.
Naturality of the connecting homomorphism shows that the original
class is $q^*y$.  Lemma~\ref{C:finite-line-act-cups} gives $y^2=0$.
Hence every class contributed by a single summand has zero square,
and these classes span $H^2$.

If $H^2$ is one-dimensional and nonzero, at least one of these
spanning classes is nonzero.  It is a scalar multiple of any
generator, so the generator's square is zero over $\mathbb Q$.
This contradicts the cohomology ring of $\mathbb{CP}^2$.
The proposition does not claim that mixed products of different
spanning classes vanish when $\dim H^2>1$.
\end{proof}

\begin{proposition}\label{C:nine-three-row-all-acts}
For each of the seven remaining three-row order-nine targets
\[
8000,\quad8005,\quad8025,\quad8030,\quad8073,\quad8105,\quad8129
\]
in the complete target list,
no act-cone category $\mathcal C(M,S)$ has type $\mathbb{CP}^2$,
regardless of the cardinality of the finite right act $S$.
\end{proposition}

\begin{proof}
The literal tables and exact character-line decompositions are stored
in \path{code/outputs/a_split_row_characters.json}.  In every table
the minimal idempotent has label $1$, and its minimal left ideal is
the three-element row set $\{1,3,5\}$.  On the augmentation basis
$[3]-[1],[5]-[1]$, every actor matrix is zero, the identity, or one
of the following rank-one idempotents:
\[
 \begin{pmatrix}0&0\\0&1\end{pmatrix},\qquad
 \begin{pmatrix}1&1\\0&0\end{pmatrix},\qquad
 \begin{pmatrix}0&-1\\0&1\end{pmatrix}.
\]
Only one of these three matrices occurs in any individual table.
Its kernel and image are invariant character lines.  The checker
rebuilds every row action from the multiplication table, verifies
the two line eigenvectors and every character multiplication, and
records the exact rational decomposition.  Apply
Proposition~\ref{C:split-row-act-cups}.

As a separate finite control, all $2214$ five-point act orbits on
these seven targets were exhaustively enumerated.  Every connected
act fails the exact rational homology required for $\mathbb{CP}^2$.
The per-target counts are $337,297,337,297,320,337,289$; the complete
action and tensor certificates are in
\path{code/outputs/a_nine_five_point_CP2_summary.json} and its seven
referenced files.  This finite control is not needed for the
all-cardinality conclusion.
\end{proof}

\subsection{Every act cone over a monoid of order at most nine}

\begin{theorem}\label{C:nine-act-cones}
For every finite monoid $M$ with $|M|\leq9$ and every nonempty
finite right $M$-set $S$, the category $\mathcal C(M,S)$ does not
have classifying space of type $\mathbb{CP}^2$.  There is no bound
on the cardinality of $S$.  The same exclusion holds for triangular
biset categories whose other endomorphism monoid has projective
trivial rational module on the side used for coinvariants.
\end{theorem}

\begin{proof}
Orders at most eight are covered by
Proposition~\ref{C:eight-act-cones}.  For order nine, if
$H^2(BM;\mathbb Q)=0$, every degree-two class of the act cone lifts
relative to $BM$.  The quotient by $BM$ is a suspension, so all
such products vanish, as in
Proposition~\ref{C:seven-act-cones}.  We may therefore assume this
vertex cohomology is nonzero.

Write the minimal ideal in Rees form with $r$ rows, $s$ columns,
and maximal group $H$.  If $r=1$ or $s=1$, the monoid has finite-group
classifying-space type and its positive rational cohomology is zero.
If $r\leq2$, Theorem~\ref{C:two-row-act-cups} applies directly.
The remaining size inequality is $rs|H|+1\leq9$ with $r\geq3$
and $s\geq2$.  It forces $H=1$, and leaves the shapes $3\times2$
or $4\times2$.  In the latter case the monoid is exactly the
rectangular band with an identity adjoined.  If $K$ is its band,
the ideal $\mathbb QK$ is projective on the left, and
$\mathbb QM/\mathbb QK=\mathbb Q$.  The row-augmentation kernel
in the projective cover $\mathbb QM e\to\mathbb Q$ is a direct
sum of copies of this quotient.  Thus
$\operatorname{pd}_{\mathbb QM}\mathbb Q\leq2$; the act-cone Tor
formula excludes degree-four homology.

For shape $3\times2$, the complete translation-table enumeration
has $191$ oriented order-nine tables, including reducible targets.
The exact, independently verified partition is
\[
\begin{array}{l|r}
 \text{all-act obstruction}&\text{number of tables}\\\hline
 H^2(BM;\mathbb Q)=0&62\\
 \operatorname{pd}_{\mathbb QM}\mathbb Q\leq2&73\\
 \text{exact periodic projective tail}&23\\
 \Omega^2\mathbb Q\text{ projective}&4\\
 \text{stable retract into a later syzygy}&22\\
 \text{split rational row augmentation}&7
\end{array}
\]
These counts sum to $191$.  The first two lines give the obstructions
already described.  The periodic and stable-retract lines force
nonzero degree-four act-cone homology to recur in higher degrees;
use the exact tensor formula of Lemma~\ref{C:act-cone-Tor} and
Lemma~\ref{C:stable-retract-Tor}.  A projective second syzygy is
another projective-dimension-two certificate.  Every module action,
projective differential, exact kernel, periodicity identification,
and stable identity is checked in the standalone verifier with
summary \path{code/outputs/b_nine_small_acts_verified.json}; the
source records retain the literal tables and exact certificate
locations.  The seven final tables are precisely those in
Proposition~\ref{C:nine-three-row-all-acts}, which excludes all their
finite acts by the split-character argument.  No transfer of an
arbitrary act along a mere homotopy equivalence of monoids is used.

For the biset assertion, take the actual finite orbit act of the
other monoid.  Exact rational coinvariants make the quotient functor
a rational homology equivalence on the edge bar and on the second
vertex, hence on the complete homotopy pushout.  It induces an
isomorphism of rational cohomology rings.  This is the same
comparison as Corollary~\ref{C:two-row-biset-cups}, now applying
the full order-nine act-cone exclusion to the quotient.
\end{proof}

\subsection{Replacing a minimal ideal by the states of an act}

\begin{theorem}[Constant minimal-ideal act compression]
\label{C:constant-minimal-act-compression}
Let $M$ be a finite monoid with proper minimal ideal $K$, and let
$I$ be the set of its minimal right ideals (the row set of its
Rees matrix presentation).  Let $S$ be a nonempty finite right
$M$-set.  Suppose that every $k\in K$ acts on $S$ as a constant
map.  There is an explicitly defined finite monoid $N$ with
\[
 |N|=|M|-|K|+|I|\,|S|
 \quad\hbox{and}\quad
 BN\simeq B\mathcal C(M,S).
\]
Here $\mathcal C(M,S)$ has objects $a,b$, endomorphism monoids
$M,\{1_b\}$, hom-set $\Hom(a,b)=S$, and no arrows from $b$ to $a$.
The conclusion allows nontrivial maximal groups in $K$.
\end{theorem}
\begin{proof}
Write $c_k\in S$ for the value of the constant action of $k$.
Right actions give $c_{kl}=c_l$ for $k,l\in K$.  In a completely
simple Rees matrix ideal, left multiplication by elements of $K$
varies the row and group coordinate arbitrarily within any fixed
column.  Consequently $c_k$ depends only on the column of $k$.
Let $\rho:K\to I$ be its row map.  The map
\[
 \eta:K\longrightarrow I\times S,
 \qquad \eta(k)=(\rho(k),c_k)
\]
is a homomorphism to the rectangular band
$(i,s)(j,t)=(i,t)$.  Left multiplication by $m\in M$ induces a
map $\sigma_m:I\to I$.  The constant-action identities imply
\[
 \eta(mk)=(\sigma_m\rho(k),c_k),\qquad
 \eta(km)=(\rho(k),c_km).                         \tag{\ref{C:constant-minimal-act-compression}.1}
\]

Set $N=(M\setminus K)\sqcup(I\times S)$.  Products of two old
elements are their product in $M$ if it lies outside $K$, and its
$\eta$-image otherwise.  Define all other products by
\[
 m(i,s)=(\sigma_m(i),s),\qquad
 (i,s)m=(i,sm),\qquad
 (i,s)(j,t)=(i,t).
\]
Equations (\ref{C:constant-minimal-act-compression}.1) verify every
associativity equation that meets the new band; the remaining
equations are those in $M$.  The old identity remains an identity.
Thus $N$ is a monoid and the identity outside $K$, together with
$\eta$ on $K$, defines a unital homomorphism $q:M\to N$.

Fix $i_0\in I,s_0\in S$ and put $e=(i_0,s_0)\in N$.
Let $D$ be the full subcategory of the Karoubi envelope of $N$
on the idempotents $1,e$.  Its four hom-sets are
\begin{align*}
 \End_D(1)&=N, & \End_D(e)&=\{e\},\\
 \Hom_D(1,e)&=eN=\{(i_0,s):s\in S\}, &
 \Hom_D(e,1)&=Ne=\{(i,s_0):i\in I\}.
\end{align*}
The object $e$ is a retract of $1$, so $BD\simeq BN$.
There is a functor $F:\mathcal C(M,S)\to D$ sending $a$ to $1$,
$b$ to $e$, $m$ to $q(m)$, and $s$ to $(i_0,s)$.
We verify Quillen's Theorem A directly for this functor.

The comma category $(F\downarrow e)$ has terminal object
$(b,e)$.  Indeed an object over $a$ is an element $(i_0,s)$,
and the unique arrow to $(b,e)$ is its corresponding $s\in S$.
For the other comma category $E=(F\downarrow1)$, the objects
over $a$ are the elements $n\in N$, and the objects over $b$
are symbols $*_i$, one for each $(i,s_0)\in Ne$.  Its arrows
among the old objects are
\[
 n\xrightarrow{\ m\ }n'\quad\Longleftrightarrow\quad n'q(m)=n.
                                                               \tag{\ref{C:constant-minimal-act-compression}.2}
\]
There is one arrow $n\to *_i$ exactly when $n=(i,s)$ for some
$s\in S$, and there are no other nonidentity arrows incident
from a new object.

For a right $M$-set $X$, let $A(X)$ denote the opposite of its
action category: its arrows are $x\to y$ labelled $m$ with
$ym=x$.  Choosing the final object and the arrow labels gives
a natural bijection $N_qA(X)=X\times M^q$.  In particular this
nerve construction preserves colimits of right $M$-sets and
preserves monomorphisms.  As right $M$-sets the construction of
$N$ is the pushout
\[
 N=M\sqcup_K(I\times S).                         \tag{\ref{C:constant-minimal-act-compression}.3}
\]
The full row category $A(\{i\}\times S)$ is a sieve in $A(N)$:
an arrow with target in that row has its source in the same row.
The description (\ref{C:constant-minimal-act-compression}.2)
therefore identifies the nerve of $E$ with the nerve of $A(N)$
with a terminal cone attached to each row nerve.  There are no
additional mixed simplices: a chain ending at $*_i$ has all its
previous objects in row $i$.  Using
(\ref{C:constant-minimal-act-compression}.3) gives the actual
pushout of realizations
\[
 BE=BA(M)\ \sqcup_{BA(K)}\
       \coprod_{i\in I}\operatorname{Cone}(BA(S)).                \tag{\ref{C:constant-minimal-act-compression}.4}
\]
The map $BA(K)\hookrightarrow BA(M)$ is a simplicial
subcomplex inclusion and hence a cofibration.  Thus this ordinary
pushout computes its homotopy pushout.  Contracting each cone to
its vertex consequently gives
\[
 BE\simeq BA(M)\sqcup^h_{BA(K)} I.
\]
The category $A(M)$ has terminal object $1$.  Also $K$ is the
disjoint union of the minimal right ideals $e_iM$, with one
idempotent $e_i$ for each $i\in I$.  Each $BA(e_iM)$ is
contractible.  More explicitly, the arrows
$x\to e_i$ labelled $x$ form a natural transformation from
the identity of $A(e_iM)$ to the constant functor at $e_i$:
for an arrow $x\xrightarrow{m}y$ its naturality equation is
$ym=x$.  The row projection $BA(K)\to I$ is therefore a
homotopy equivalence.  It follows that $BE$ is contractible.

Both comma categories of $F$ are contractible, so Quillen's
Theorem A gives $B\mathcal C(M,S)\simeq BD\simeq BN$.
\end{proof}

\begin{corollary}[A numerical obstruction from the monoid minimum]
\label{C:constant-minimal-act-size}
Under the hypotheses of
Theorem~\ref{C:constant-minimal-act-compression}, if
$B\mathcal C(M,S)\simeq\mathbb CP^2$, then
\[
 |M|-|K|+|I|\,|S|\geq23.
\]
\end{corollary}
\begin{proof}
Apply the exact monoid minimum in Theorem~\ref{B:CP2-minimum23}
to the monoid just constructed.
\end{proof}

\begin{remark}[The constant-action hypothesis is essential to the formula]
\label{C:constant-minimal-act-scope}
The theorem does not infer constant action from the number of
states.  In particular a nontrivial maximal group in $K$ can act
nontrivially on a small orbit.  Such an action must be treated
separately before applying this compression.  If $K=M$ is a
group and its identity acts constantly, $S$ is a singleton and
the act cone is contractible; this degenerate case needs no
replacement construction.
\end{remark}

\begin{lemma}[The constant-action quotient detects degree two]
\label{C:constant-action-quotient-H2}
Let $M$ be any finite monoid with minimal ideal $K$, and let $S$
be a nonempty finite right $M$-set.  Let $\bar S$ be the quotient
of $S$ by the equivalence relation generated by
\[
 sk\sim s'k\qquad(s,s'\in S, k\in K).
\]
Then $\bar S$ is a right $M$-set on which every element of $K$
acts constantly.  The induced functor of act cones gives a
surjection
\[
 H^2(B\mathcal C(M,\bar S);\mathbb Q)
 \longrightarrow H^2(B\mathcal C(M,S);\mathbb Q).
\]
If $K$ is a rectangular band, the same surjection holds with
integral coefficients.
In particular, if $|S|\leq3$ and some element of $K$ does not
act constantly on $S$, all degree-two rational products on
$B\mathcal C(M,S)$ vanish.  No restriction on the maximal group
of $K$ is imposed.
\end{lemma}
\begin{proof}
The relation is $M$-invariant because $km\in K$ for $k\in K$,
$m\in M$.  Put $R=\mathbb QM$, $L=\widetilde{\mathbb QS}$,
and $\bar L=\widetilde{\mathbb Q\bar S}$.  Its relation span
is exactly
\[
 \ker(L\longrightarrow\bar L)=LK:
 \qquad ([s]-[s'])k=[sk]-[s'k].
\]
Here taking the transitive closure introduces no additional linear
relations, since the difference between two points in an equivalence
class is the sum of differences along a generating chain.

Choose a minimal idempotent and average its maximal group, as in
Theorem~\ref{C:two-row-act-cups}.  This gives a projective left
$R$-module $P\cong\mathbb QI$ and an exact sequence
\[
 0\longrightarrow V\longrightarrow P\longrightarrow\mathbb Q
       \longrightarrow0,
 \qquad V=\widetilde{\mathbb QI},\quad KV=0.
\]
For the rational dual $D$, every map $V\to DL$ takes values in
the annihilator of $LK$: for $v\in V$, $\ell\in L$ and $k\in K$,
equivariance gives $f(v)(\ell k)=(kf(v))(\ell)=f(kv)(\ell)=0$.
Thus the inclusion $D\bar L\hookrightarrow DL$ induces an
isomorphism
\[
 \Hom_R(V,D\bar L)\xrightarrow{\ \cong\ }\Hom_R(V,DL).
\]
The natural act-cone bar and the above projective presentation
identify its degree-two cohomology with
\[
 \operatorname{Ext}^1_R(\mathbb Q,DL)
  =\operatorname{coker}\bigl(\Hom_R(P,DL)
                         \longrightarrow\Hom_R(V,DL)\bigr).
\]
Applying the same formula to $\bar L$ proves the claimed
surjection: the numerator is unchanged and the denominator can
only increase on passing to $DL$.

For a minimal band the same argument works over $\mathbb Z$.
Here $P=\mathbb ZI$ is projective without averaging, and
$L/LK=\widetilde{\mathbb Z\bar S}$ is a free abelian group:
the subgroup $LK$ is exactly the lattice generated by differences
inside equivalence classes.  Its integral dual is therefore the
annihilator of $LK$, and the same Hom-isomorphism and connecting
cokernels prove the integral assertion.

If a $K$-actor is nonconstant and $|S|\leq3$, the quotient has
at most two points.  Every action on a two-point set is affine
on $\{0,1\}$ with slope in $\{0,1,-1\}$, and a one-point act
has contractible cone.  Lemma~\ref{C:finite-line-act-cups} makes
all positive-degree products of the quotient cone zero.
The cohomology surjection and cup naturality prove the last claim.
\end{proof}

\begin{lemma}[A three-by-six band with five outside elements]
\label{C:three-by-six-five-outside-CP2}
A finite monoid with rectangular-band minimal ideal of shape
$3\times6$ or $6\times3$, and exactly five elements outside that
ideal, cannot have classifying space $\mathbb CP^2$.
\end{lemma}
\begin{proof}
After opposition the band--Tor modules have dimensions
$\dim V=2$, $\dim W=5$, over the five-dimensional contracted
monoid algebra $A$, with Tor dimensions $1,0,1,0,\ldots$.
A local $A$ is excluded by Proposition~\ref{B:local-quotient-CP2}.
The nonlocal classification in
Proposition~\ref{B:five-nonlocal-exceptions} leaves just the line
and Nakayama algebras of Lemma~\ref{B:five-exceptional-targets}.

For the zero-composite line algebra, its displayed resolutions
and $\dim V=2$, $\operatorname{Tor}_1=0$, $\dim\operatorname{Tor}_2=1$
force
\[
 V=S_1\oplus S_3,\qquad
 W=T_1\oplus T_3^{\oplus s}\oplus(e_1A)^{\oplus p}
                           \oplus(e_2A)^{\oplus q}.
\]
The dimension and tensor constraints are
$s+2p+2q=4$ and $1+s+p=1$.  Hence $s=p=0,q=2$.
Only the $T_1$ summand contributes to the balanced tensor product,
so the unique degree-two form has matrix rank one.  Its square
vanishes by Lemma~\ref{B:rank-one-square}.

For the Nakayama algebra, the full Tor table in the same lemma
forces $V=L$ and
$W=T_2\oplus(e_1A)^{\oplus p}\oplus(e_2A)^{\oplus q}$.
The dimension equation is $3p+2q=4$, with unique solution
$p=0,q=2$.  Its tensor dimension is $p+q=2$, contradicting the
required value one.  These cases are exhaustive and prove the claim.
\end{proof}

\begin{theorem}[The exact three-state act-cone minimum]
\label{C:three-state-cone-minimum22}
If a three-state act cone $\mathcal C(M,S)$ has classifying space
$\mathbb CP^2$, then it has at least twenty-two arrows.  This bound
is attained by the two models of Theorem~\ref{C:CP2-upper22}.
\end{theorem}
\begin{proof}
Suppose $m=|M|\leq17$, so the cone has at most twenty-one arrows.
Its nonzero degree-two square requires $H^2(BM;\mathbb Q)\ne0$;
otherwise every degree-two class is relative to $BM$ and products
vanish.  Write $|K|=rsh$, where $r,s$ are the numbers of rows and
columns and $h$ is the maximal-group order.  The rational balanced
presentation requires $r,s\geq2$, and
Theorem~\ref{C:two-row-act-cups} requires $r\geq3$.

Lemma~\ref{C:constant-action-quotient-H2} now forces every element
of $K$ to act constantly.  The compression is therefore an actual
homotopy-equivalent monoid, of order
\[
 |N|=m-rsh+3r.
\]
If $h\geq2$, then $sh\geq4$, and $|N|\leq m-r<23$,
contradicting Theorem~\ref{B:CP2-minimum23}.  If $h=1$ and
$s\geq3$, then $|N|\leq m<23$, with the same contradiction.
It remains that $h=1,s=2$.  Put $q=m-2r$, the number of elements
outside the minimal band.  Compression leaves precisely $q$
elements outside its new $r\times3$ band.  The exact monoid
minimum and Corollary~\ref{B:five-outside-minimal-band} require
\[
 m+r\geq23,\qquad q=m-2r\geq5.
\]
Together with $m\leq17$ these inequalities force
$m=17$, $r=6$, $q=5$.  The compressed monoid then has a
$6\times3$ band and five outside elements, excluded by
Lemma~\ref{C:three-by-six-five-outside-CP2}.
This proves the lower bound.  The models of
Theorem~\ref{C:CP2-upper22} have an eighteen-element endomorphism
monoid, three states and one sink identity, and give equality.
\end{proof}

\begin{lemma}[At most four states over a minimal band]
\label{C:four-state-band-dichotomy}
Let $M$ have rectangular-band minimal ideal $K=I\times J$, and
let $S$ be a nonempty finite right $M$-set with $|S|\leq4$.
If $H_1(B\mathcal C(M,S);\mathbb Q)=0$, then either every
degree-two rational cup product on this act cone is zero, or
every element of $K$ acts constantly on $S$.
\end{lemma}
\begin{proof}
First suppose that some $z\in S$ is fixed by every element of
$K$.  Adjoining a fresh identity to $K$ gives a unital monoid
$K^+\to M$ and a functor $K^+\to S\rtimes M$, taking its object
to $z$.  By the natural balanced-tensor presentation in
Lemma~\ref{lem:B-band-Tor}, the induced map
\[
 H_2(BK^+;\mathbb Q)
   =\widetilde{\mathbb QJ}\otimes_{\mathbb Q}
      \widetilde{\mathbb QI}
   \longrightarrow H_2(BM;\mathbb Q)
\]
is surjective: passing from $K^+$ to $M$ only adds balancing
relations.  Hence $H^2(BM;\mathbb Q)\to H^2(B(S\rtimes M);\mathbb Q)$
is injective.  The cofiber sequence for the act cone implies that
every class in its $H^2$ lifts to relative cohomology with respect
to $BM$.  The relative quotient is a suspension of the edge
space, so its positive-degree cup products vanish.  Their images
give the asserted zero products on the act cone.

Assume now that no point is fixed by all of $K$.  The set $SK$
is the disjoint union of its right $K$-orbits.  Each such orbit
is $M$-invariant, since $(xk)m=x(km)$ and $km\in K$.
It has at least two points, since a singleton orbit would be a
$K$-fixed point.  If there are at least two orbits, $|S|\leq4$
forces exactly two two-point orbits and $S=SK$.  The action
category is then disconnected.  This is impossible: the
degree-zero and degree-one part of the act-cone exact sequence,
using $H_1(BM;\mathbb Q)=0$, identifies
$H_1(B\mathcal C(M,S);\mathbb Q)$ with the reduced $H_0$ of
that action category.

There is therefore a single right $K$-orbit in $SK$.
Every $k\in K$ acts constantly on this orbit.  Indeed choose a
point $y=xk_0$ in the orbit and write every point as
$y\ell=xk_0\ell$.  The band product $k_0\ell k$ has the row
of $k_0$ and the column of $k$, independently of $\ell$.
Thus $(y\ell)k$ is independent of the point $y\ell$.
Finally, for an arbitrary $x\in S$, idempotence gives
$xk=(xk)k$.  The point $xk$ lies in $SK$, so this is the same
constant value for every $x$.  Thus every $k$ acts constantly
on all of $S$.
\end{proof}

\begin{corollary}[Four-state cones through seventeen arrows]
\label{C:four-state-cones-seventeen}
If $|S|=4$ and $|M|\leq12$, then
$B\mathcal C(M,S)\not\simeq\mathbb CP^2$.
Thus no four-state act cone with at most seventeen arrows
models $\mathbb CP^2$.  Also, if $|S|=3$, $|M|\leq15$, and
the minimal ideal of $M$ is a band, the same exclusion holds.
\end{corollary}
\begin{proof}
A possible model must have $H^2(BM;\mathbb Q)\ne0$, since
otherwise every degree-two cone class is relative and has zero
products.  Write the minimal-ideal shape as $I\times H\times J$,
with $r=|I|,s=|J|$.  The rational group-collapse and balanced-tensor
formulas of Lemmas~\ref{lem:B-collapse-minimal-group} and
\ref{lem:B-band-Tor} require $r,s\geq2$.
Theorem~\ref{C:two-row-act-cups} further requires $r\geq3$.
When $|M|\leq12$, a nontrivial $H$ would already require
$rs|H|+1\geq13$ elements.  Thus $K$ is a band.

Lemma~\ref{C:four-state-band-dichotomy} now forces constant
$K$-action.  For four states the compressed monoid has order
\[
 |N|=m-rs+4r\leq m+2r
 \leq m+2\left\lfloor\frac{m-1}{2}\right\rfloor\leq22
 \qquad(m=|M|\leq12).
\]
This contradicts Corollary~\ref{C:constant-minimal-act-size}.
For three states and a band ideal, the same argument gives
\[
 |N|=m-rs+3r\leq m+r
 \leq m+\left\lfloor\frac{m-1}{2}\right\rfloor\leq22
 \qquad(m\leq15).
\]
The band qualification in this last assertion is explicit;
no inference of constant action is made for a nontrivial
minimal group.
\end{proof}

\begin{theorem}[A three-object triangle as a band extension]
\label{C:triangle-band-compression}
Let $M$ be a finite monoid, $S$ a nonempty finite left $M$-set,
and $T$ a nonempty finite right $M$-set.  Let $C$ have objects
$0,1,2$, endomorphism monoids $\{1_0\},M,\{1_2\}$, and
\[
 \Hom_C(0,1)=S,\qquad \Hom_C(1,2)=T,\qquad
 \Hom_C(0,2)=\{w\},
\]
with no other arrows between distinct objects.  Every composite
$t s$ equals $w$.  Then $BC\simeq BN$, where
\[
 N=M\sqcup(S\times T),\qquad |N|=|M|+|S|\,|T|,
\]
and multiplication extends that of $M$ by
\[
 (s,t)(s',t')=(s,t'),\qquad
 m(s,t)=(ms,t),\qquad (s,t)m=(s,tm).
\]
\end{theorem}
\begin{proof}
The two action laws verify associativity of $N$ directly.
Choose $s_0\in S,t_0\in T$ and put $e=(s_0,t_0)$.
In the full subcategory $D\subset\operatorname{Kar}(N)$ on
$1,e$, the corner $eNe$ is a singleton,
$\Hom_D(e,1)=Ne\cong S$, and $\Hom_D(1,e)=eN\cong T$.
Define $F:C\to D$ by sending $0,2$ to $e$ and $1$ to $1$,
using the inclusion $M\hookrightarrow N$, sending
$s\mapsto(s,t_0)$ and $t\mapsto(s_0,t)$, and sending $w$ to $e$.
The composite of the latter two band elements is $e$, so this
is a functor.  Again $BD\simeq BN$ because $e$ is a retract of $1$.

The comma category $(F\downarrow e)$ has terminal object
$(2,e)$.  To analyze $E=(F\downarrow1)$, denote its objects
over $0$ and $2$ by $0_s,2_s$ for $s\in S$, and its objects
over $1$ by $1_n$ for $n\in N$.  The arrows among the latter
objects are labelled $m\in M$ with $n'm=n$, as in the previous
proof.  An arrow $0_s\to1_n$ is labelled $a\in S$ with
$n a=s$, where $M$ acts on $S$ by the given action and
$(u,t)\in S\times T$ acts constantly with value $u$.
An arrow $1_n\to2_s$ exists precisely when $n=(s,t)$;
it is then uniquely labelled $t$.  Finally there is one arrow
$0_s\to2_{s'}$ if and only if $s=s'$, labelled $w$.

Let $E_0$ be the full subcategory on all $0_s$ and the
$1_m$ with $m\in M$.  It has terminal object $1_1$:
the unique arrow $1_m\to1_1$ is labelled $m$, and the unique
arrow $0_s\to1_1$ is labelled $s$.
Let $E_1$ be the full subcategory on all $0_s,2_s$ and all
$1_{(s,t)}$.  It is the disjoint union, over $s\in S$, of
categories with terminal object $2_s$.
Their intersection is the discrete category of the objects $0_s$.

There are no arrows between an old label $1_m$ and a band label
$1_{(s,t)}$: right multiplication by elements of $M$ preserves
the two summands of $N$.  The shared objects $0_s$ have no
nonidentity incoming arrows.  Every composable chain in $E$
therefore lies entirely in $E_0$ or entirely in $E_1$, and its
nerve is the actual pushout
\[
 NE=NE_0\sqcup_S NE_1.
\]
Both inclusions of $S$ are simplicial cofibrations.  Moreover
$S\to BE_1$, selecting $0_s$ in its contractible component, is
a homotopy equivalence.  The pushout is consequently homotopy
equivalent to $BE_0$, which is contractible.  Quillen's Theorem A
now gives $BC\simeq BD\simeq BN$.
\end{proof}

\begin{corollary}\label{C:triangle-band-CP2-size}
A category in Theorem~\ref{C:triangle-band-compression} with
$BC\simeq\mathbb CP^2$ satisfies $|M|+|S|\,|T|\geq23$.
In particular the fifteen-arrow profile
\[
 (|\End(0)|,|\End(1)|,|\End(2)|,
   |\Hom(0,1)|,|\Hom(1,2)|,|\Hom(0,2)|)
       =(1,8,1,2,2,1)
\]
is impossible: its compressed monoid would have only twelve
elements.
\end{corollary}
\begin{proof}
Apply Theorem~\ref{B:CP2-minimum23} to the compressed monoid.
\end{proof}

\subsection{Two cross arrows and a degree-two acyclic corner}

\begin{lemma}\label{C:two-state-biset-H2zero}
Let $M,N$ be finite monoids and $S$ a two-element $(N,M)$-biset.
If $H^2(BM;\mathbb Q)=0$ or $H^2(BN;\mathbb Q)=0$, all products
of degree-two rational classes on its two-object collage vanish.
This hypothesis does not require either vertex to be rationally acyclic.
\end{lemma}
\begin{proof}
By opposition suppose $H^2(BM;\mathbb Q)=0$, and regard $S$ as
a left $N$-set and right $M$-set.  Every finite monoid has zero
first rational cohomology.  We examine the image of the left action
in the four-element transformation monoid of $S$.

If both constant maps occur, the commuting right action must be
the identity on both points.  The edge double bar is then
$BM\times E_L$, where $E_L$ is the left-action bar of $N$ on $S$.
Collapsing $M$ to the trivial monoid gives an actual collage functor
to the ordinary two-state left-act cone over $N$.  K\"unneth gives
isomorphisms in degrees one and two between $E_L$ and
$BM\times E_L$, because $H^1(BM;\mathbb Q)=H^2(BM;\mathbb Q)=0$.
The vertex degree-two groups also agree.  Naturality of the two
Mayer--Vietoris sequences consequently makes pullback on $H^2$ an
isomorphism.  Every positive-degree product on the target act cone
vanishes by Lemma~\ref{C:finite-line-act-cups}, applied to the
opposite monoid, so the source degree-two products vanish.

If exactly one constant occurs, no swap can occur in the left
image, since composing a swap with that constant would give the
other constant.  Its value is therefore fixed by every element
of $N$.  Commutation also makes it fixed by every element of $M$.
The corresponding fixed-state bar is an actual section
$BM\times BN\to E_S$.  Thus the direct sum of the two vertex
degree-two groups injects into the edge group.

If no constant occurs and the left action is trivial, the edge
is $E_R\times BN$.  Pullback from $BN$ is injective, and it is
the entire vertex degree-two group since $H^2(BM;\mathbb Q)=0$.
Finally, if the left image contains the swap and no constant,
commutation forces the right image to consist of permutations
too.  The edge is then the two-sheeted covering of $BM\times BN$
associated to this permutation action.  Transfer makes rational
pullback injective, so again the vertex degree-two group injects
into the edge group.

In each of the last three cases, Mayer--Vietoris makes the
restriction of every collage degree-two class to its two vertex
spaces zero.  Every such class lifts to relative cohomology;
the relative quotient of the double mapping cylinder is
$\Sigma(E_{S,+})$.  Its positive-degree products vanish.
Their images in absolute cohomology give the required conclusion.
\end{proof}

\subsection{The universal obstruction for two cross arrows}

\begin{lemma}[Cancelling a trivially acting vertex in degree two]
\label{C:trivial-biset-side-H2}
Let $S$ be a nonempty finite $(N,M)$-biset, with trivial left
$N$-action.  Collapsing $N$ to the trivial monoid gives an actual
functor from its two-object collage to the ordinary right-act cone
$\mathcal C(M,S)$.  This functor induces an isomorphism on rational
degree-two cohomology.  The opposite statement holds when the right
$M$-action is trivial.
\end{lemma}
\begin{proof}
Write $E_R$ for the right-action bar of $M$ on $S$.  The edge
double bar of the collage is literally $E_R\times BN$, and its
maps to the two vertex spaces are the action-bar projection to
$BM$ and the second projection to $BN$.
Every finite monoid has zero first rational cohomology.  K\"unneth
therefore gives
\[
 \begin{aligned}
 H^1(E_R\times BN;\mathbb Q)&=H^1(E_R;\mathbb Q),\\
 H^2(E_R\times BN;\mathbb Q)&=H^2(E_R;\mathbb Q)
 \oplus H^0(E_R;\mathbb Q)\otimes H^2(BN;\mathbb Q).
 \end{aligned}
\]
In the second direct summand, pullback from $BN$ is the diagonal
map $y\mapsto 1\otimes y$, which is injective because $E_R$ is
nonempty.  Pullback from $BM$ lands in the first summand.
Consequently the kernel of the map from the two vertex
degree-two groups to the edge group is exactly
\[
 \ker\bigl(H^2(BM;\mathbb Q)\longrightarrow
                 H^2(E_R;\mathbb Q)\bigr)\oplus0.
\]
The Mayer--Vietoris sequence for the double mapping cylinder
places its degree-two cohomology in a short exact sequence with
left term $H^1(E_R;\mathbb Q)$ and right term this kernel.
The ordinary act-cone sequence has the same two end terms.
Naturality under the stated collapse functor gives the identity
on both, and hence an isomorphism in the middle.  This argument
allows disconnected $E_R$ and does not assume that either vertex
has zero second cohomology.
\end{proof}

\begin{theorem}[Two cross arrows force zero degree-two products]
\label{C:two-cross-arrows-universal}
Let $M,N$ be arbitrary finite monoids and let $S$ be a
two-element $(N,M)$-biset.  All products of rational degree-two
classes on its two-object collage vanish.  In particular, no
such collage realizes $\mathbb CP^2$.
\end{theorem}
\begin{proof}
Consider the image of the left $N$-action in the four-element
transformation monoid of $S$.

If both constant transformations occur, any commuting right
transformation fixes both points of $S$, so the right $M$-action
is trivial.  Lemma~\ref{C:trivial-biset-side-H2}, applied with
the sides reversed, identifies the collage's rational $H^2$
with that of the ordinary two-state left-act cone over $N$.
Every positive-degree product of the latter is zero by
Lemma~\ref{C:finite-line-act-cups}, applied to the opposite
monoid.  The comparison is an actual functor, so it preserves
cup products and proves the assertion in this case.

If exactly one constant transformation occurs, no swap occurs:
composing the swap with that constant would produce the other
constant.  The constant value is therefore fixed by the left
action.  Commutation with this constant also makes the same
point fixed by the right action.  Its fixed-state bar gives a
section $BM\times BN\to E_S$ of the edge projection.
Thus the direct sum of the two vertex degree-two groups
injects into $H^2(E_S;\mathbb Q)$.

If no constant transformation occurs and the left action is
trivial, Lemma~\ref{C:trivial-biset-side-H2} reduces the assertion
to the ordinary two-state right-act cone over $M$, and again
Lemma~\ref{C:finite-line-act-cups} applies.

Finally suppose the left image contains the swap and no constant.
Every commuting right transformation is then a permutation too.
The edge double bar is the two-sheeted covering of $BM\times BN$
associated to the commuting permutation actions.  Rational
transfer makes its pullback injective.  In particular the direct
sum of the vertex degree-two groups injects into the edge group.

In the two cases with an injective vertex-to-edge map,
Mayer--Vietoris says that every collage degree-two class restricts
to zero on both vertex spaces.  It therefore lifts to relative
cohomology.  Collapsing the two vertex spaces in the double
mapping cylinder gives the suspension $\Sigma(E_{S,+})$.
All positive-degree products in its reduced cohomology vanish;
their images in absolute cohomology give the required vanishing.
These cases exhaust all possible left transformation images.
\end{proof}

\begin{remark}\label{C:two-cross-strengthens-earlier-gates}
Theorem~\ref{C:two-cross-arrows-universal} removes the hypothesis
on a vertex degree-two group from
Lemma~\ref{C:two-state-biset-H2zero}.  In particular it excludes
the three two-cross-arrow profiles $(9,5,2)$, $(8,6,2)$ and
$(7,7,2)$ in the sixteen-arrow scope without any further census.
Earlier enumerations of two-state bisets remain independent
checks of the stronger statement.
\end{remark}

\subsection{Proper biset retracts and minimal models}

\begin{lemma}[A proper biset retract preserves a hypothetical projective plane]
\label{C:proper-biset-retract-CP2}
Let $S$ be a finite $(N,M)$-biset and let $f:S\to S$ be an
equivariant idempotent with nonempty image $U$.  If the collage
on $S$ has type $\mathbb CP^2$, then the collage on $U$ has
type $\mathbb CP^2$ as well.
\end{lemma}
\begin{proof}
The inclusion and $f$ give actual functors between the two
collages, identity on the vertices, whose composite on the
smaller collage is the identity.  Write their classifying
spaces as $X_U\subseteq X_S$ and their edge bars as
$E_U\subseteq E_S$.  These are inclusions of simplicial
subspaces.  The double mapping cylinder pair has quotient
\[
 X_S/X_U\simeq\Sigma(E_S/E_U).
\]
If restriction $H^2(X_S;\mathbb Q)\to H^2(X_U;\mathbb Q)$
were zero, its generator would lift to relative cohomology.
The suspension quotient has zero positive-degree products,
contradicting the projective-plane square.

Restriction is therefore nonzero.  The retraction splits
integral homology and the fundamental group.  Since those of
$X_S\simeq\mathbb CP^2$ are $\mathbb Z$ in degrees zero,
two and four, zero in the other positive degrees, and trivial
in degree one homotopy, the same holds for $X_U$ in degree two
and for its fundamental group.  Pullback from $X_U$ is an
isomorphism in degree two.  Cup naturality forces its fourth
homology to be nonzero, hence the full $\mathbb Z$ summand.
Thus the actual inclusion $X_U\to X_S$ is an integral
homology equivalence between simply connected CW spaces.
The homological Whitehead theorem makes it a homotopy
equivalence.
\end{proof}

\begin{corollary}\label{C:biset-core-prune}
A least-arrow two-object triangular model of $\mathbb CP^2$
has no proper equivariant biset retraction.  At sixteen arrows,
every proper equivariant idempotent on the cross-biset is
excluded by the already complete fifteen-arrow lower bound.
\end{corollary}

\begin{corollary}[A trivial action side gives an actual act-cone model]
\label{C:trivial-biset-side-CP2-retract}
If one side of a finite biset action is trivial and its collage
has type $\mathbb CP^2$, the corresponding ordinary act cone
also has type $\mathbb CP^2$.  A nontrivial trivially acting
corner can consequently be removed in a minimal model.
\end{corollary}
\begin{proof}
The collapse functor of Lemma~\ref{C:trivial-biset-side-H2}
has an actual section, obtained by inserting the identity of
the collapsed monoid.  Its pullback is an isomorphism on
rational $H^2$.  The split integral-homology and cup-product
argument in Lemma~\ref{C:proper-biset-retract-CP2} therefore
makes this section an integral homology equivalence between
simply connected CW spaces.
\end{proof}

\begin{lemma}[Commuting permutation actions]
\label{C:permutation-biset-cups}
If both actions on a nonempty finite biset are by permutations,
every rational degree-two product of its collage vanishes.
\end{lemma}
\begin{proof}
The edge action bar is a finite-sheeted covering of
$BM\times BN$.  Rational transfer makes pullback injective,
and K\"unneth embeds the direct sum of the two vertex
degree-two groups in the product's degree-two group.
Thus every collage degree-two class restricts to zero on
both vertices.  The relative suspension quotient kills all
their products, as in Theorem~\ref{C:two-cross-arrows-universal}.
\end{proof}

\subsection{The unique residual pattern on three cross arrows}

\begin{lemma}[Commuting images on three points]
\label{C:three-point-commuting-images}
Let two finite transformation monoids act on a three-element set,
and suppose their actions commute.  At least one of the following
holds: they have a common fixed point; both act by permutations;
one acts trivially; there is a proper idempotent transformation
commuting with both; or, after relabeling the points and interchanging
the sides, their images are exactly
\[
 \{\id,t_0,t_1\},\qquad \{\id,c_0,c_1\},
\]
where $t_i$ fixes $0,1$ and sends $2$ to $i$, and $c_i$ is
constant with value $i$.
\end{lemma}
\begin{proof}
Here is a complete finite verification with two independent
enumerations.  Encode each of the $27$ transformations by its
three values.  The first enumeration starts with the identity
submonoid, adjoins each omitted transformation, takes product
closure, and repeats until no new closed subset occurs.  It gives
exactly $699$ submonoids.  Every submonoid occurs, since adjoining
its elements one at a time gives a path to it in this search.

The independent enumeration assigns each of the $27$ maps to
an included or excluded set.  It propagates the three implications
\[
 a,b\in P\Rightarrow ab\in P,\qquad
 a\in P,\ ab\notin P\Rightarrow b\notin P,\qquad
 b\in P,\ ab\notin P\Rightarrow a\notin P.
\]
It branches on every still undecided map, beginning with the
identity included.  Its $1{,}463$ branch nodes give exactly the
same $699$ literal subsets.  Their $1{,}854$ commuting ordered
pairs are tested directly.  In the displayed order of the
alternatives, their disjoint counts are
\[
\begin{array}{c|rrrrr}
\text{alternative}&\text{fixed point}&\text{permutations}&
 \text{trivial side}&\text{proper idempotent}&\text{residual}\\\hline
\text{ordered pairs}&595&5&1158&90&6.
\end{array}
\]
For the fourth test, all $27$ maps are checked for idempotence,
proper image, and literal commutation with both actions.  Its
ninety witnesses are retained.  The six remaining ordered pairs
are exactly the simultaneous state conjugates and side interchanges
of the pair in the statement.
The complete lists and verification are
\path{code/outputs/b_three_state_commuting_images.json} and
\path{code/outputs/b_three_state_image_classification_verified.json};
the two algorithms are
\path{code/b_three_state_commuting_images.py} and
\path{code/b_verify_three_state_image_classification.cpp}.
This classifies the action images only; it does not replace either
original corner monoid by its image.
\end{proof}

\begin{lemma}[The residual three-state collage has zero second homology]
\label{C:three-state-residual-H2zero}
Let $M,N$ be arbitrary finite monoids.  Suppose their commuting
actions on $S=\{0,1,2\}$ have right $M$-image
$\{\id,c_0,c_1\}$ and left $N$-image
$\{\id,t_0,t_1\}$, with notation as in
Lemma~\ref{C:three-point-commuting-images}.
Then the collage $\mathcal C_S$ has
$H_2(B\mathcal C_S;\mathbb Q)=0$.
\end{lemma}
\begin{proof}
Put $U=\{0,1\}$, $u=[1]-[0]$, and $v=[2]-[0]$.
The two nontrivial zero-one characters are
\[
 \chi(m)=\begin{cases}1& m\text{ acts identically},\\0&
 m\text{ acts constantly},\end{cases}
 \qquad
 \psi(n)=\begin{cases}1& n\text{ acts identically},\\0&
 n\text{ acts as }t_0\text{ or }t_1.\end{cases}
\]
As a right $M$-module, both $u$ and $v$ transform by $\chi$.
As a left $N$-module, $u$ is fixed and
\[
 n v=\psi(n)v+\gamma(n)u,
 \qquad \gamma(t_0)=0,\quad\gamma(t_1)=1.
\]
Here the last equalities mean any preimages of the indicated
transformations; $\gamma$ is zero on the identity image.
Let $F_N$ denote this two-dimensional left module.  Its extension
\begin{equation}\label{C:three-state-gamma-extension}
 0\longrightarrow\mathbb Q\longrightarrow F_N
 \longrightarrow\mathbb Q_\psi\longrightarrow0
\end{equation}
is not split.  Indeed a splitting would send the quotient generator
to $v+a u$.  A preimage of $t_0$ requires $a=0$, whereas a preimage
of $t_1$ requires $a=-1$.

Write $X_U\subseteq X_S$ for the classifying spaces of the
collages with cross-bisets $U\subseteq S$.
Their two vertex spaces are identical.  In the double action-bar
complex, the coefficient quotient is
\[
 \mathbb QS/\mathbb QU
   =\mathbb Q_\psi\otimes_{\mathbb Q}\mathbb Q_\chi.
\]
Consequently the reduced chains of the pair $(X_S,X_U)$ are
the suspension of the tensor product of the two bar complexes
representing
\[
 D_N=\mathbb Q\otimes^{\mathbf L}_{\mathbb QN}\mathbb Q_\psi,
 \qquad
 D_M=\mathbb Q_\chi\otimes^{\mathbf L}_{\mathbb QM}\mathbb Q.
\]
Both degree-zero homology groups vanish: each character takes
the value zero on some actor.  K\"unneth therefore gives
\[
 H_2(X_S,X_U;\mathbb Q)=0,
 \qquad
 H_3(X_S,X_U;\mathbb Q)
 =\operatorname{Tor}^{\mathbb QN}_1(\mathbb Q,\mathbb Q_\psi)
   \otimes
   \operatorname{Tor}^{\mathbb QM}_1(\mathbb Q_\chi,\mathbb Q).
\]

We identify the boundary of this pair explicitly.
The coefficient extension for $U\subseteq S$ is the pushout
of extensions of left-$N$, right-$M$ bimodules
\[
\begin{array}{ccccccccc}
0&\to&\mathbb Q\otimes\mathbb Q_\chi&\to&
 F_N\otimes\mathbb Q_\chi&\to&
 \mathbb Q_\psi\otimes\mathbb Q_\chi&\to&0\\
 &&\downarrow [1]-[0]&&\downarrow&&\Vert\\
0&\to&\mathbb QU&\to&\mathbb QS&\to&
 \mathbb Q_\psi\otimes\mathbb Q_\chi&\to&0.
\end{array}
\]
Thus the relative connecting map is the extension class
$\gamma$ of \eqref{C:three-state-gamma-extension}, tensored
with the identity on $\mathbb Q_\chi$, followed by the
inclusion $u:\mathbb Q_\chi\to\mathbb QU$.
This can also be read directly in normalized chains: in the
splitting $\mathbb QS=\mathbb QU\oplus\mathbb Qv$, the
right-$M$ faces have no off-diagonal term, and the only
off-diagonal left-$N$ face is $\gamma(n)u$.
Both vertex projections annihilate $v$.

On the core $U$, the left $N$-action is trivial.
The actual collapse $X_U\to B\mathcal C(M,U)$ is a rational
$H_2$ isomorphism by Lemma~\ref{C:trivial-biset-side-H2}.
The augmentation extension
\[
0\longrightarrow\mathbb Q_\chi
 \xrightarrow{u}\mathbb QU\longrightarrow\mathbb Q
 \longrightarrow0
\]
identifies the reduced act-cone complex with $D_M[1]$.
Its canonical quasi-isomorphism into the cone is induced by
$u$, not by a choice of homology splitting.  After the core
collapse, the boundary of the pair on degree three is therefore
\[
 \partial=\pm\gamma\otimes\id:
 \operatorname{Tor}^{\mathbb QN}_1(\mathbb Q,\mathbb Q_\psi)
 \otimes H_1(D_M)\longrightarrow H_1(D_M).
\]
The nonsplit extension gives a nonzero functional
$\gamma$ on the first Tor group, by the exact duality
\[
 \operatorname{Ext}^1_{\mathbb QN}(\mathbb Q_\psi,\mathbb Q)
 \cong
 \operatorname{Hom}_{\mathbb Q}\!\left(
   \operatorname{Tor}^{\mathbb QN}_1(\mathbb Q,\mathbb Q_\psi),
   \mathbb Q\right).
\]
That functional is onto $\mathbb Q$.  Thus $\partial$ is
onto $H_2(X_U;\mathbb Q)$, including when the latter is zero.
Together with $H_2(X_S,X_U;\mathbb Q)=0$, the homology
sequence of the pair proves $H_2(X_S;\mathbb Q)=0$.
\end{proof}

\begin{theorem}[Three cross arrows reduce to an ordinary act cone]
\label{C:three-cross-arrows-act-reduction}
If a two-object finite collage with three cross arrows has
type $\mathbb CP^2$, one side of the cross action is trivial,
and its corresponding ordinary three-state act cone also has
type $\mathbb CP^2$.
\end{theorem}
\begin{proof}
Apply Lemma~\ref{C:three-point-commuting-images} to the two
literal transformation images.  A common fixed state makes
the vertex degree-two groups inject into the edge; every
collage class is then relative and has zero square.
Permutation images are excluded by
Lemma~\ref{C:permutation-biset-cups}.
A proper equivariant idempotent retracts the collage to one
with at most two cross arrows; Lemma~\ref{C:proper-biset-retract-CP2}
would make that retract a projective plane, contrary to
Theorem~\ref{C:two-cross-arrows-universal} (the one-arrow
case is a join and has zero degree-two products).
The residual pattern is excluded by
Lemma~\ref{C:three-state-residual-H2zero}.
Only a trivial action side remains, and
Corollary~\ref{C:trivial-biset-side-CP2-retract} gives the
stated ordinary act cone.
\end{proof}

\begin{corollary}\label{C:three-cross-arrows-lower22}
Every finite projective-plane collage with three cross arrows
has at least twenty-two arrows.  If both corner monoids are
nontrivial, it has at least twenty-three arrows.
\end{corollary}
\begin{proof}
The ordinary three-state cone supplied by the theorem has at
least twenty-two arrows by Theorem~\ref{C:three-state-cone-minimum22}.
Removing a nontrivial trivially acting corner decreases the
number of arrows by at least one.
\end{proof}

\subsection{A two-point constant image at arbitrary state count}

\begin{theorem}\label{C:two-constant-image-bisets}
Let $S$ be a finite nonempty $(N,M)$-biset.  Suppose the right
$M$-action has exactly the three transformations
\[
 \{\id_S,c_0,c_1\},
\]
where $c_i$ is constant at a distinct point $i\in S$.
Then its two-object collage cannot have type $\mathbb CP^2$.
The same holds with the sides interchanged.  There is no bound
on the numbers of states or corner elements.
\end{theorem}
\begin{proof}
Commutation with $c_0,c_1$ forces every $N$-actor to fix both
points.  Put $U=\{0,1\}$ and write $X_U\subseteq X_S$ for
the two collage spaces.  Let $\chi:M\to\{0,1\}$ distinguish
identity and constant transformations.  As a right module,
the augmentation module $L=\ker(\mathbb QS\to\mathbb Q)$
is a direct sum of copies of $\mathbb Q_\chi$.
The vector $u=[1]-[0]$ is fixed by $N$.  With $V=L/\mathbb Qu$
as a left $\mathbb QN$-module, there is an extension
\begin{equation}\label{C:two-constant-image-extension}
 0\longrightarrow\mathbb Q\xrightarrow{u}L
 \longrightarrow V\longrightarrow0
\end{equation}
with class $\gamma\in\operatorname{Ext}^1_{\mathbb QN}(V,\mathbb Q)$.
Here $L$ in this display denotes its left-module structure;
its commuting right-module structure is scalar $\chi$.

The coefficient extension
$0\to\mathbb QU\to\mathbb QS\to V\otimes\mathbb Q_\chi\to0$
is the pushout of \eqref{C:two-constant-image-extension}
tensored with $\mathbb Q_\chi$ along the actual inclusion
$u:\mathbb Q_\chi\to\mathbb QU$.
Thus the reduced relative chains of $(X_S,X_U)$ are
\[
 \left(
  \mathbb Q\otimes^{\mathbf L}_{\mathbb QN}V
  \otimes_{\mathbb Q}
  \mathbb Q_\chi\otimes^{\mathbf L}_{\mathbb QM}\mathbb Q
 \right)[1].
\]
The degree-zero homology of the right factor is zero because
$\chi$ takes the value zero.  The K\"unneth summand
\[
 \operatorname{Tor}^{\mathbb QN}_1(\mathbb Q,V)
 \otimes
 \operatorname{Tor}^{\mathbb QM}_1(\mathbb Q_\chi,\mathbb Q)
 \ \subseteq\ H_3(X_S,X_U;\mathbb Q)
\]
therefore has the same boundary computation as in
Lemma~\ref{C:three-state-residual-H2zero}.
More explicitly, the augmentation extension of $\mathbb QU$
identifies the ordinary core act-cone complex with
$(\mathbb Q_\chi\otimes^{\mathbf L}_{\mathbb QM}\mathbb Q)[1]$.
Collapsing the trivially acting $N$-corner induces an actual
$H_2$ isomorphism on $X_U$ by
Lemma~\ref{C:trivial-biset-side-H2}.  Under this identification,
the boundary of the displayed summand is
\[
 \pm\gamma\otimes\id:
 \operatorname{Tor}^{\mathbb QN}_1(\mathbb Q,V)
 \otimes H_2(X_U;\mathbb Q)
 \longrightarrow H_2(X_U;\mathbb Q).
\]
This is a chain-level connecting map: in the splitting
$\mathbb QS=\mathbb QU\oplus V$ the only off-diagonal action
faces come from the extension \eqref{C:two-constant-image-extension};
the right faces are scalar $\chi$, and both vertex augmentations
vanish on the chosen augmentation-module lifts of $V$.

If $\gamma\ne0$, the duality
\[
 \operatorname{Ext}^1_{\mathbb QN}(V,\mathbb Q)
  \cong
 \operatorname{Hom}_{\mathbb Q}
 \bigl(\operatorname{Tor}^{\mathbb QN}_1(\mathbb Q,V),\mathbb Q\bigr)
\]
makes this boundary onto.  Hence $H_2(X_U;\mathbb Q)\to
H_2(X_S;\mathbb Q)$ is zero.  Every degree-two cohomology
class of $X_S$ is relative to $X_U$; the quotient is the
suspension of the corresponding edge-bar pair.  Its positive
products vanish, excluding a projective plane.

If $\gamma=0$, choose an $N$-linear retraction
$q:L\to\mathbb Q$ with $q(u)=1$, and set
$q(s)=q([s]-[0])$.  Then $q(0)=0$, $q(1)=1$, and $q(ns)=q(s)$.
Choose the single rounding function
$\rho:\mathbb Q\to\{0,1\}$ with $\rho(1)=1$ and
$\rho(x)=0$ for $x\ne1$.  The map $r(s)=\rho(q(s))$ fixes
$U$, is $N$-equivariant, and commutes with every right
$M$-actor, since those actors are identity or constant at
$0,1$.  It is an actual biset retraction.
If $|S|>2$, Lemma~\ref{C:proper-biset-retract-CP2} would make
$X_U$ a projective plane whenever $X_S$ were one, contrary
to Theorem~\ref{C:two-cross-arrows-universal}.  The case
$|S|=2$ follows from that theorem directly.
\end{proof}

\begin{remark}\label{C:two-constant-image-scope}
The nonsplit case in the proof gives zero rational
degree-two products without a rank or finiteness assumption.
The split case gives the stated projective-plane exclusion
by an actual retract; no general zero-product assertion is
needed there.
\end{remark}

\subsection{Three constants and one transient state}

\begin{proposition}\label{C:constant-core-one-transient-H2}
Let $S=U\sqcup\{t\}$ be a finite $(N,M)$-biset.  Suppose the
right $M$-image consists of the identity and all constant maps
with values in $U$.  Suppose the left $N$-image consists of the
identity and the maps $r_j$, for $j$ in a nonempty subset
$J\subseteq U$, where $r_j$ fixes $U$ pointwise and sends $t$ to
$j$.  Let $T$ be the actual $N$-orbit set: it identifies $t$ and
all of $J$, and leaves the other points of $U$ distinct.
The actual functor from the biset collage to the ordinary
$M$-act cone on $T$ induces an isomorphism on rational second
homology, and hence on rational second cohomology.
\end{proposition}
\begin{proof}
Put $A=\ker(\mathbb QU\to\mathbb Q)$ and
$L=\ker(\mathbb QS\to\mathbb Q)$.
The identity-versus-constant character is
$\chi:M\to\{0,1\}$, and the identity-versus-retraction
character is $\psi:N\to\{0,1\}$.
The whole right action on $L$ is scalar $\chi$.
As left $N$-modules there is an exact sequence
\begin{equation}\label{C:constant-core-vector-extension}
 0\longrightarrow A\longrightarrow L
 \longrightarrow\mathbb Q_\psi\longrightarrow0,
\end{equation}
where $A$ has trivial $N$-action.

Write $X_U\subseteq X_S$ for the two collage spaces.
The coefficient extension for this pair is the pushout of
\eqref{C:constant-core-vector-extension}, tensored with
$\mathbb Q_\chi$, along the inclusion
$A\otimes\mathbb Q_\chi\to\mathbb QU$.
Consequently its relative chains are the suspended tensor of
the two action bars for $\mathbb Q_\psi$ and $\mathbb Q_\chi$.
Both bars have zero degree-zero homology, because each character
takes the value zero.  It follows that
\[
 H_2(X_S,X_U;\mathbb Q)=0,\qquad
 H_3(X_S,X_U;\mathbb Q)=
 \operatorname{Tor}^{\mathbb QN}_1(\mathbb Q,\mathbb Q_\psi)
 \otimes
 \operatorname{Tor}^{\mathbb QM}_1(\mathbb Q_\chi,\mathbb Q).
\]
The left action on $U$ is trivial.  By
Lemma~\ref{C:trivial-biset-side-H2}, collapsing the $N$-corner
identifies its second homology with that of the ordinary core
act cone.  The augmentation sequence of $\mathbb QU$ therefore
identifies
\[
 H_2(X_U;\mathbb Q)=
 A\otimes
 \operatorname{Tor}^{\mathbb QM}_1(\mathbb Q_\chi,\mathbb Q).
\]
Under these identifications the boundary of the relative pair
is $\pm\delta\otimes\id$, where
\[
 \delta:
 \operatorname{Tor}^{\mathbb QN}_1(\mathbb Q,\mathbb Q_\psi)
 \longrightarrow A
\]
is the connecting map of
\eqref{C:constant-core-vector-extension}.
This is the same literal off-diagonal-face calculation as in
Lemma~\ref{C:three-state-residual-H2zero}, now with values in
$A$: choose $u_0\in J$ and split
$\mathbb QS=\mathbb QU\oplus\mathbb Q([t]-[u_0])$;
the only off-diagonal action faces are the left extension faces;
the right faces are scalar $\chi$, and both vertex augmentations
annihilate the chosen augmentation lifts.

Taking $N$-coinvariants in
\eqref{C:constant-core-vector-extension} shows
\[
 \operatorname{im}\delta
  =\ker(A\longrightarrow L_N)
  =\ker(\mathbb QJ\longrightarrow\mathbb Q).
\]
Indeed the coinvariant relations identify $[t]$ with each $[j]$,
$j\in J$, and impose no other relations between the basis points.
Thus $L_N$ is exactly $\ker(\mathbb QT\to\mathbb Q)$.
The pair sequence now gives
\[
 H_2(X_S;\mathbb Q)=
 \bigl(A/\ker(\mathbb QJ\to\mathbb Q)\bigr)
 \otimes
 \operatorname{Tor}^{\mathbb QM}_1(\mathbb Q_\chi,\mathbb Q).
\]
The actual orbit-collapse functor induces the displayed quotient
map on the first factor.  The same augmentation calculation for
the ordinary $M$-act $T$ gives the identical expression, so this
functor is the asserted second-homology isomorphism.
\end{proof}

\begin{corollary}\label{C:three-constant-four-state-cups}
Under the hypotheses of
Proposition~\ref{C:constant-core-one-transient-H2}, if
$|U|=3$ and $|J|\geq2$, every rational degree-two cup product
on the biset collage is zero.  In particular it cannot have
type $\mathbb CP^2$.
\end{corollary}
\begin{proof}
Here $|T|=|U|-|J|+1\leq2$.
Every degree-two class pulls back from the ordinary cone on
$T$.  Theorem~\ref{C:two-cross-arrows-universal} gives zero
products when $|T|=2$; when $|T|=1$ that cone is contractible.
Naturality of cup products proves the claim.
\end{proof}

\subsection{Five states over a three-by-two minimal band}

\begin{lemma}[Invariant small column sets in the remaining three-by-five hull]
\label{C:3x5-small-invariant-columns}
Let $H$ be a local-character-allowed actor submonoid of orbit $164$
in Theorem~\ref{B:three-by-five-targets}, containing its full
$3\times5$ band and identity.  Suppose that a nonempty subset of
at most two columns is invariant under every actor outside the
band.  Then every product of rational degree-two classes of the
ambient orbit-$164$ hull restricts to zero on $BH$.
More precisely, $H$ lies in one of five
explicit actor monoids with zero fourth rational homology.
\end{lemma}
\begin{proof}
The allowed maximal actor monoid has thirty elements, consisting
of the fifteen band elements, its identity, and fourteen optional
actors.  Test all $2^{14}$ optional subsets for multiplication
closure and for invariance of each of the five singleton and ten
two-element column subsets.  The complete counts are
\[
 \begin{array}{l|r}
 \text{closed actor monoids}&2186\\
 \text{closed monoids admitting such an invariant subset}&408\\
 \text{maximal monoids in the latter family}&5.
 \end{array}
\]
In the indexing of the complete safe-submonoid list, the five
maxima and their independently certified bar ranks are
\[
 \begin{array}{c|r|r|r|r}
 \text{index}&\text{order}&\operatorname{rank}d_2&
 \operatorname{rank}d_3&\dim C_2\\\hline
 1031&22&39&249&288\\
 1732&24&57&455&512\\
 1794&24&57&455&512\\
 1804&24&57&455&512\\
 1860&24&57&455&512.
 \end{array}
\]
The ranks are over $\mathbb F_5$ in the normalized two-sided
complex of Lemma~\ref{lem:B-band-Tor}; its degree two computes
fourth monoid homology.  In every case the two ranks sum to
the middle dimension, proving $H_4(-;\mathbb F_5)=0$, hence
$H^4(-;\mathbb Q)=0$.  Inclusions into these five monoids
therefore kill every pulled-back degree-two product.

The direct subset enumeration and every maximality comparison are
in \path{code/a_five_state_band_obstruction.py}.  It compares its
complete list with the independently certified orbit-$164$ list,
but does not use that list to generate subsets.  The output
\path{code/outputs/a_five_state_band_obstruction.json} includes
all closed sets, all $408$ retained images, and independent-column
witnesses for the five pairs of differentials.  A separate dense
table-based reconstruction, using the verifier of the earlier
three-by-three theorem with shape parameters $(3,5)$, checks
those rank witnesses and all associativity equations.
\end{proof}

\begin{theorem}[A five-state three-by-two obstruction]
\label{C:five-state-three-by-two-cones}
Let $M$ be a finite monoid with rectangular-band minimal ideal
of shape $3\times2$, and let $S$ be a right $M$-set with five
elements.  Then
\[
 B\mathcal C(M,S)\not\simeq\mathbb CP^2.
\]
There is no bound on the number of elements of $M$ in this statement.
\end{theorem}
\begin{proof}
Write $K=I\times J$, with $|I|=3$ and $|J|=2$, and suppose
that $X=B\mathcal C(M,S)$ has the indicated finite homotopy type.
The act is connected, since the act-cone homology sequence and
$H_1(BM;\mathbb Q)=0$ identify $H_1(X;\mathbb Q)$ with the
reduced degree-zero homology of the action category.
If a state is fixed by all of $K$, the first paragraph of
Lemma~\ref{C:four-state-band-dichotomy} applies without any
state-count restriction: the band grid maps onto $H_2(BM)$
through the edge space, and every degree-two product on $X$
is zero.  We may exclude this case.

First suppose that some element of $K$ acts nonconstantly.
The set $SK$ is a disjoint union of right $K$-orbits, each
invariant under $M$.  Each orbit has at least two elements and
at most two elements: its elements are obtained by varying the
column of a fixed band row.  If there were a single orbit,
idempotence would make every $K$-actor constant on all of $S$,
as in Lemma~\ref{C:four-state-band-dichotomy}.  Thus $SK$
consists of two two-point orbits, with the remaining state
outside $SK$.  Label the two orbits by their two column values.

Let $\bar S$ be the quotient in
Lemma~\ref{C:constant-action-quotient-H2}.  For either column,
its state in the first orbit is identified with its state in
the second orbit.  Therefore $|\bar S|\leq3$.
All elements of $K$ act constantly on $\bar S$, and that lemma
gives a surjection
\[
 H^2(B\mathcal C(M,\bar S);\mathbb Q)
       \longrightarrow H^2(X;\mathbb Q).
\]
Compress the quotient act cone by
Theorem~\ref{C:constant-minimal-act-compression}, obtaining
a monoid $N$ with a $3\times|\bar S|$ minimal band.
For at most two columns every balanced matrix has rank at most
one.  Lemma~\ref{B:rank-one-square} and polarization therefore
make the rational degree-two products zero, without requiring
finite homotopy type of $BN$.

For three columns, the proof of
Theorem~\ref{B:three-by-three-cups} applies to every class
pulled back to $X$.  Here is the finite-target point that is
needed for that assertion.  All cases other than its exceptional
skew-form orbit have universal cup primitives or zero fourth
cohomology witnesses, which pull back directly.  In the skew
case, an eventual local group of $N$ that maps onto a forbidden
cyclic actor group is nontrivial and therefore lies outside the
new band.  The outside elements and all products within such
a group are literally unchanged from $M$.  This supplies a
local-group map into the endomorphism monoid of
$\mathcal C(M,S)$.  The actual functor
\[
 \mathcal C(M,S)\longrightarrow\mathcal C(M,\bar S)
       \longrightarrow\operatorname{Kar}(N)
\]
pulls the integral hull cocycle back along this same group map.
Its forbidden cyclic character remains nonzero under inflation,
whereas Miller's theorem makes the map from that finite group's
classifying space to the finite CW model of $X$ null-homotopic.
This is the same contradiction as in the original skew-orbit
proof.  Thus its bad actors are excluded here as well, and the
remaining nine targets have zero fourth rational cohomology.
No finite-type assertion about $BN$ is used.
All pulled-back degree-two products consequently vanish.
Their degree-two surjection onto $X$ contradicts the square in
$H^*(\mathbb CP^2;\mathbb Q)$.

It remains that every element of $K$ acts constantly on $S$.
Compression now gives an actual homotopy equivalence
\[
 X\simeq BN,
\]
where $N$ has minimal band $I\times S$ of shape $3\times5$.
For $j\in J$, let $c_j\in S$ be the constant value of
any band element with column $j$.  The nonempty set
\[
 T=\{c_j:j\in J\},\qquad |T|\leq2,
\]
is invariant under every outside actor.  Indeed
$c_jm=c_{km}$ for any $k$ in column $j$, and $km\in K$.
Theorem~\ref{B:three-by-five-targets} reduces a possible
projective-plane source of this shape to its allowed orbit-$164$
action image.  That image fixes the entire band, and every
outside actor in the image comes from an outside element of
$N$; it therefore preserves the corresponding subset $T$.
Lemma~\ref{C:3x5-small-invariant-columns} makes the restricted
ambient degree-two square zero.  All rational degree-two classes
of $BN$ pull back from the ambient orbit-$164$ hull: the source,
its action image, and that maximal balanced-form hull have the
same rational balancing relation space.  Thus the
balanced-tensor presentation identifies their degree-two groups,
so this is again a contradiction.
\end{proof}

\begin{remark}[Scope and computational controls]
\label{C:five-state-band-cones-scope}
The preceding theorem assumes a minimal band, not a nontrivial
minimal maximal group.  It uses complete finite actor calculations
only in the already classified three-by-three and three-by-five
hulls and in the five new zero-homology witnesses above; it does
not enumerate monoids of unbounded order.
As a separate control, all $1082$ two-core-orbit five-state normal
forms over the $272$ nonsplit-row order-ten residual targets
were checked in \path{code/a_five_state_two_orbits.py}.
Of these, $1050$ fail exact rational homology through degree
three and the other $32$ have explicit mod-two square primitives.
Three complete independent action enumerations agree with the
normal-form generator.  These controls are not used to establish
the all-order theorem.
\end{remark}

\subsection{Projective orbit identifications for a two-column band}

\begin{lemma}[The kernel of the constant-action quotient]
\label{C:two-column-projective-orbit-kernel}
Let $M$ have rectangular-band minimal ideal $K=I\times\{0,1\}$,
and let $S$ be a nonempty finite right $M$-set with no state fixed
by all of $K$.  Write $o$ for the number of right $K$-orbits in
$SK$, and let $\bar S$ be the constant-action quotient of
Lemma~\ref{C:constant-action-quotient-H2}.  Put
$R=\mathbb ZM$, $L=\widetilde{\mathbb ZS}$, and
$\bar L=\widetilde{\mathbb Z\bar S}$.  For any minimal idempotent
$e\in K$ there is an exact sequence of right $R$-modules
\[
 0\longrightarrow(eR)^{\oplus(o-1)}
   \longrightarrow L\longrightarrow\bar L\longrightarrow0.
\]
In particular, if $X=B\mathcal C(M,S)$ and
$Y=B\mathcal C(M,\bar S)$, their actual quotient map induces
isomorphisms $H_n(X;\mathbb Z)\cong H_n(Y;\mathbb Z)$ for
$n\geq3$, and an exact sequence
\[
 0\longrightarrow H_2(X;\mathbb Z)
 \longrightarrow H_2(Y;\mathbb Z)
 \longrightarrow\mathbb Z^{o-1}
 \longrightarrow H_1(X;\mathbb Z)
 \longrightarrow H_1(Y;\mathbb Z)\longrightarrow0.
\]
\end{lemma}
\begin{proof}
Every $K$-orbit in $SK$ has two elements: its column map has
at most two values, and a one-point orbit would be a $K$-fixed
state.  Label these orbits
\[
 O_a=\{s_{a,0},s_{a,1}\},\qquad 0\leq a<o,
\]
by their column values.  Each orbit is $M$-invariant.  Indeed,
if $x=y(i,j)$ and $m\in M$, then
$xm=y((i,j)m)=y(i,\tau_m(j))$ stays in the same orbit.
More precisely every actor has the same column map on every orbit:
\[
 s_{a,j}m=s_{a,\tau_m(j)}.
\]
The kernel of $L\to\bar L$ is $LK$, as proved in
Lemma~\ref{C:constant-action-quotient-H2}.  Its exact integral
basis is
\[
 s_{a,j}-s_{0,j}\qquad(1\leq a<o,\ j\in\{0,1\}).
\]
These differences belong to $LK$ by applying a band element in
column $j$ to points in the two indicated core orbits.
Conversely, for fixed $k=(i,j)$, every point of $Sk$ is the
column-$j$ point of one of these orbits.  Thus every generating
difference $sk-s'k$ is a difference of the displayed kind.
The points of $S\setminus SK$ introduce no additional relations.
The formula for the action identifies each pair indexed by $a$
with the right ideal $eR$, whose basis is the two elements in
the row of $e$.  This proves the module statement, including
projectivity of its kernel.

The proof of Lemma~\ref{C:act-cone-Tor} works over $\mathbb Z$
without change: use a projective left $\mathbb ZM$-resolution
of the trivial module and the exact augmentation sequence of
the free abelian permutation module $\mathbb ZS$.  Thus
$\widetilde H_n(X;\mathbb Z)=\operatorname{Tor}_{n-1}^R(L,\mathbb Z)$
for every $n\geq1$, naturally in $S$.  Apply the Tor exact
sequence to the displayed projective-kernel sequence and note
that $eR\otimes_R\mathbb Z\cong\mathbb Z$.  This gives exactly
the asserted homology sequence and all higher isomorphisms.
\end{proof}

\begin{lemma}[A free second-homology complement]
\label{C:free-second-homology-wedge}
Let $f:X\to Y$ be a map of simply connected CW complexes.
Suppose $f_*$ is an isomorphism on integral homology in every
degree at least three, and fits into an exact sequence
\[
 0\longrightarrow H_2(X;\mathbb Z)
 \xrightarrow{f_*}H_2(Y;\mathbb Z)
 \longrightarrow\mathbb Z^r\longrightarrow0.
\]
There is a homotopy equivalence
$X\vee\bigvee_rS^2\to Y$ whose restriction to $X$ is $f$,
up to homotopy. No finiteness assumption on $X$ or $Y$ is needed.
\end{lemma}
\begin{proof}
The displayed sequence splits because its quotient is free.
Choose a basis of a complementary copy of $\mathbb Z^r$ in
$H_2(Y;\mathbb Z)$. Since $Y$ is simply connected, the Hurewicz
isomorphism $\pi_2(Y)\cong H_2(Y;\mathbb Z)$ represents this
basis by based maps from $r$ two-spheres. Wedge these maps with
$f$, after the harmless choice of compatible basepoints.
The resulting map is an isomorphism on second homology by the
chosen splitting, and on every higher homology group by the
hypothesis on $f$. It is also an isomorphism in degrees zero
and one. Both spaces are simply connected CW complexes, so
homological Whitehead makes it a homotopy equivalence.
\end{proof}

\begin{corollary}[The orbit quotient splits off two-spheres]
\label{C:two-column-orbit-wedge}
Under the hypotheses of
Lemma~\ref{C:two-column-projective-orbit-kernel}, suppose the
original act cone $X=B\mathcal C(M,S)$ is simply connected.
The actual quotient map to
$Y=B\mathcal C(M,\bar S)$ extends to a homotopy equivalence
\[
 X\vee\bigvee_{o-1}S^2\ \simeq\ Y\ \simeq\ BN,
\]
where $N$ is the constant-action compression monoid.
\end{corollary}
\begin{proof}
The compression theorem gives the displayed equivalence with
$BN$. Its minimal ideal is a rectangular band, so its universal
group and hence its fundamental group are trivial.
The projective-kernel lemma gives the required isomorphisms in
all degrees at least three. Since $H_1(X)=H_1(Y)=0$, its exact
sequence in low degrees has free cokernel $\mathbb Z^{o-1}$.
Apply Lemma~\ref{C:free-second-homology-wedge} to the actual
quotient map. The equivalence thus retains that map on the
original summand; it is not merely an abstract homology comparison.
\end{proof}

\begin{lemma}[A finite model for the quotient of a hypothetical projective plane]
\label{C:two-column-finite-quotient-model}
Under the hypotheses of
Lemma~\ref{C:two-column-projective-orbit-kernel}, suppose
$B\mathcal C(M,S)\simeq\mathbb CP^2$.
The compressed monoid $N$ associated to $\bar S$ has the same
number of outside elements as $M$, minimal band
$I\times\bar S$, and a finite four-dimensional CW model with
one zero-cell, $o$ two-cells, and one four-cell. More precisely,
\[
 BN\simeq\mathbb CP^2\vee\bigvee_{o-1}S^2.
\]
Its integral
homology is
\[
 H_2(BN)=\mathbb Z^o,\qquad H_4(BN)=\mathbb Z,
 \qquad \widetilde H_n(BN)=0\quad(n\notin\{2,4\}).
\]
The map from the original act cone to $BN$ is surjective on
integral degree-two cohomology in the pullback direction.
\end{lemma}
\begin{proof}
Compression is Theorem~\ref{C:constant-minimal-act-compression};
it leaves the outside elements unchanged and replaces the band
by $I\times\bar S$.  The preceding exact sequence gives the
homology assertion, since an extension of $\mathbb Z^{o-1}$
by $\mathbb Z$ splits as an abelian group.
Corollary~\ref{C:two-column-orbit-wedge} proves the stronger
wedge identification under the original projective-plane summand.
The usual cell structure of $\mathbb CP^2$, together with the
additional two-spheres, is the stated finite four-dimensional model.
Finally the asserted cohomology surjection is the integral
statement of Lemma~\ref{C:constant-action-quotient-H2} followed
by the compression equivalence.
\end{proof}

\begin{theorem}[Four outside elements never suffice for a two-column act cone]
\label{C:two-column-four-outside-act-cones}
Let $M$ have a rectangular-band minimal ideal with two columns
and at most four elements outside it.  For every nonempty finite
right $M$-set $S$,
\[
 B\mathcal C(M,S)\not\simeq\mathbb CP^2.
\]
There is no bound on the number of rows or states.
More generally, if the act cone has bounded rational homology, all
of its rational degree-two products vanish.  Thus the obstruction
also applies to the rational homology and cohomology ring of a
projective plane, without requiring an integral finite CW model.
\end{theorem}
\begin{proof}
We prove the stronger rational assertion.
If $S$ has a state fixed by $K$, the band-grid argument in
Lemma~\ref{C:four-state-band-dichotomy} makes every degree-two
class relative to $BM$, and their products are zero.  That argument
does not use the number of states.
Otherwise Lemma~\ref{C:two-column-projective-orbit-kernel} and
compression give a monoid $N$ with the same rational homology
as the act cone in every degree at least three.  Its rational
homology is therefore bounded.  It has at most four elements
outside its minimal band.  Theorem~\ref{B:contracted-four-Tor},
applied to its row and column augmentation modules in the band
Tor formula, implies $H_4(BN;\mathbb Q)=0$: a nonzero group in
that degree would force nonzero groups in arbitrarily high degrees.
Thus the original act cone also has zero rational fourth homology,
and its rational degree-two products vanish.  This argument uses
neither an integral finiteness inference nor a degree-two rank
restriction on the compressed monoid.
\end{proof}

\begin{corollary}[Six and seven nonconstant states over three rows]
\label{C:three-row-six-seven-nonconstant}
Let $M$ have minimal band $3\times2$, and let $S$ have six or
seven elements.  If some element of the band acts nonconstantly
on $S$, then $B\mathcal C(M,S)\not\simeq\mathbb CP^2$.
There is no bound on the order of $M$.
\end{corollary}
\begin{proof}
Suppose the cone is a projective plane.  As above, it has no
$K$-fixed state.  The act is connected: the homotopy-pushout
homology sequence identifies its first rational homology with
the reduced degree-zero homology of the action category, because
$H_1(BM;\mathbb Q)=0$.
If $SK$ had a single orbit, every band actor would be constant
on all states, since $xk=(xk)k$ and $k$ is constant on that orbit.
Thus $o\geq2$.  The $K$-orbits are individually $M$-invariant,
so $S=SK$ with $o>1$ would disconnect the act.

For six states these facts force $o=2$, with two states outside
$SK$.  The constant-action quotient identifies the two core
orbits column by column and leaves the other states distinct,
so $|\bar S|=4$.  The finite compressed model has minimal band
$3\times4$ and degree-two rank two.  The rank-two portion of
the proof of Theorem~\ref{B:three-by-four-targets} makes every
rational degree-two product zero for any finite CW source of
that rank: its only nonzero-square exceptional orbit has rank
one.  Its actual zero-product certificates and local-group
exclusions therefore apply to this compressed model.

For seven states, either $o=2$ and $|\bar S|=5$, or $o=3$ and
$|\bar S|=3$.  In the former case the finite compressed model
has degree-two rank two and minimal band $3\times5$; every
rank-two orbit in the proof of
Theorem~\ref{B:three-by-five-targets} has zero rational products.
In the latter case Theorem~\ref{B:three-by-three-cups} applies
without a degree-two rank restriction.
In all cases every rational degree-two class of the original
cone pulls back from the compressed model.  Its nonzero square
is therefore impossible.
\end{proof}

\begin{remark}[Scope of the orbit argument]
\label{C:two-column-orbit-quotient-scope}
The projective-kernel assertion uses a minimal band with exactly
two columns and the absence of one-point core orbits.  For a
larger column set, different core orbits can identify different
columns, and the displayed kernel need not be projective.
The proof does not assert projectivity in that situation.
The integral orbit calculation and the finite four-cell model
were independently audited by agents B and D; the six-state
rank-two application was checked against the full earlier
classification, not inferred merely from its named product target.
\end{remark}

\subsection{The local five-dimensional quotient on three rows}

\begin{lemma}[An affine pencil cannot give a lone positive even Tor group]
\label{C:drifting-pencil-Tor}
Let $A$ be the rational algebra with basis $1,x,y,z,c$, with
\[
 x^2=z,\qquad xz=zx=y^2=c,\qquad
 xy=\epsilon c,\qquad yx=(1-\epsilon)c,
 \qquad \epsilon\in\{0,1\},
\]
and all other products between $x,y,z,c$ zero.
For $a\in\mathbb Q$, let $V_a$ be the two-dimensional module
on which $y$ acts as a nonzero square-zero Jordan block, $x$
acts as $a$ times that block, and $z,c$ act as zero.
There is no finite-dimensional right $A$-module $W$ with
\[
 \operatorname{Tor}_1^A(W,V_a)=0,\qquad
 \dim\operatorname{Tor}_2^A(W,V_a)=1,\qquad
 \operatorname{Tor}_n^A(W,V_a)=0\quad(n\geq3).
\]
\end{lemma}
\begin{proof}
Put $\eta=1-\epsilon$, $s=\epsilon-\eta$, and
$t_n=a+ns$.  A free resolution of $V_a$ has one copy of $A$
in every degree and differentials, specified by the images of $1$,
\[
 d_{2n+1}(1)=x-t_ny,\qquad
 d_{2n+2}(1)=y+(t_n-\eta)z\qquad(n\geq0).
\]
These are maps of left modules, hence right multiplications by
the displayed elements.  Direct multiplication gives
\begin{align*}
 (y+(t-\eta)z)(x-ty)&=0,\\
 (x-(t+s)y)(y+(t-\eta)z)&=0.
\end{align*}
The left ideal generated by $x-ty$ has basis
$x-ty,z,c$ for every rational $t$, and the left ideal generated
by $y+(t-\eta)z$ has basis $y+(t-\eta)z,c$.
Thus successive images have dimensions three and two, their
composites are zero, and they are exact in the five-dimensional
free terms.  The initial cokernel is $V_a$, so this is the
claimed complete resolution.

After tensoring with $W$, write $d=\dim W$ and $r_i$ for the
rank of the $i$th differential.  The asserted Tor dimensions
give
\[
 r_1+r_2=d,\qquad r_2+r_3=d-1,
 \qquad r_n+r_{n+1}=d\quad(n\geq3).
\]
In particular $r_1-r_3=1$, while
$r_3=r_5=r_7=\cdots$.
The odd differentials are the matrix pencil
$R_x-t_nR_y$, where $R_x,R_y$ are the right-action matrices
on $W$.  Since $s=\pm1$, these parameters are distinct.
A rational matrix pencil has its maximal rank at all but
finitely many rational parameters: choose a nonzero minor of
maximal size and exclude its finitely many roots.
Consequently $r_3$ equals that maximal rank.  Its specialization
at $t_0=a$ has rank $r_1\leq r_3$, contradicting $r_1-r_3=1$.
The argument uses the stated vanishing in every degree at least
three, not just an unspecified eventual tail with additional
intervening homology.
\end{proof}

\begin{proposition}[Every local five-element contracted basis on three rows]
\label{C:local-five-three-row-Tor}
Let $A$ be a local contracted rational algebra of a finite
monoid, with $\dim A=5$.  Suppose $V$ is a two-dimensional
$A$-module with a fixed augmentation basis of a three-point
set, in which every distinguished nonzero monoid basis element
acts by the augmentation matrix of a transformation of that
set.  The transformations themselves need not descend through
the contracted zero; only their augmentation operators are
required to give the stated $A$-module.
There is no finite-dimensional right module $W$
with
\[
 \operatorname{Tor}_1^A(W,V)=0,\qquad
 \dim\operatorname{Tor}_2^A(W,V)=1,\qquad
 \operatorname{Tor}_n^A(W,V)=0\quad(n\geq3).
\]
\end{proposition}
\begin{proof}
We give the exact finite basis calculation and its independent
verification, including the only infinite-resolution issue.
In a local contracted monoid algebra there is no nonidentity
idempotent basis element and no nontrivial finite subgroup:
either would give a nontrivial idempotent in the algebra.
Every nonidentity nonzero monoid element therefore has zero
idempotent power.  The span of these four elements is a
nilpotent ideal and is the radical $J$, with quotient
$A/J=\mathbb Q$.  Its nilpotence is also checked literally
in the finite enumeration below.

If $V$ is semisimple, it is the sum of two copies of the simple
module and its degree-two Tor has even dimension.  Otherwise
$JV$ is one-dimensional and $J^2V=0$.  All radical operators
are scalar multiples of one nonzero square-zero matrix $N$.
A three-point transformation with nonzero nilpotent augmentation
operator is a length-two chain; relabeling makes it $(0,0,1)$.
In its augmentation basis $N=\left(\begin{smallmatrix}0&1\\0&0\end{smallmatrix}\right)$.
The actual three-point transformations whose augmentation is
a multiple of $N$ are the three constants, $(0,0,1)$, and
$(1,1,0)$, with coefficients $0,1,-1$ respectively.
Thus every nonidentity distinguished basis element acts by
$0,N$, or $-N$, and its coefficient vanishes on $J^2$.

Start from all $1373$ isomorphism-and-opposite classes in the
complete order-six census
\path{code/outputs/c_monoids_6.tsv}, add literal opposites,
retain tables with a zero, and contract it.  Computing powers
of the four-dimensional nonunit ideal leaves $156$ distinct
literal local bases.  On each basis enumerate all nonzero
assignments of coefficients in $\{-1,0,1\}$ to its four
nonidentity elements and impose the multiplication equations
on their two-by-two matrices.  There are $3600$ assignments.
Relabeling the four nonidentity basis elements and allowing the
common change $N\mapsto-N$ leaves $1286$ pairs $(A,V)$.
This deliberately retains assignments that might not admit
an affine lift to a literal three-point action; overenumeration
cannot discard a possible source.

Exact rational projective covers and explicit kernel bases give
the following complete partition:
\[
\begin{array}{l|r}
 \text{verified obstruction}&\text{pairs}\\\hline
 \text{a simple direct summand in a syzygy}&414\\
 \text{a syzygy with top dimension greater than radical dimension}&859\\
 \text{a split recurrence of the first syzygy}&7\\
 \text{the two affine-pencil algebras of Lemma~\ref{C:drifting-pencil-Tor}}&6.
\end{array}
\]
The first line supplies actual inclusion and retraction matrices.
If the proposed Tor profile held, $W$ would be nonprojective;
its Tor with the local simple module is then nonzero in every
degree, by Lemma~\ref{B:local-minimal-Tor}.  Dimension shifting
therefore excludes that line.  The second line is excluded by
Lemma~\ref{B:top-heavy-Tor}, again using dimension shifting.
The seven recurrences give actual split maps from the first
syzygy into a later syzygy.  They are excluded by
Lemma~\ref{B:split-syzygy-persistence}.

For the last six records, a literal permutation of the radical
basis gives exactly the multiplication in
Lemma~\ref{C:drifting-pencil-Tor}; the row module has
$a\in\{-1,0,1\}$.  That lemma supplies the complete resolution
and the rank obstruction.  Equality of successive module
dimensions is not treated as module periodicity: in these six
records the pencil parameter actually changes.

The constructor is \path{code/a_five_local_all_rows.py}.
The separate checker is
\begin{quote}\small
\path{code/a_verify_five_local_all_rows.py}.
\end{quote}
It imports none of the search, cover, or syzygy implementations
and rebuilds the
entire $3600$-assignment source coverage directly from the
census, using nilpotence of the full ideal in place of the
constructor's individual power test, and independently imposes
all matrix multiplication equations.  It checks every one of
the $1372$ stored projective-prefix steps by surjectivity,
equivariance, zero composite, and complementary exact ranks;
it checks every inclusion and retraction and every top dimension.
For the six exceptional records it verifies the two zero
composites as polynomial identities and exhibits constant unit
minors proving ranks three and two for every parameter.
The result
\path{code/outputs/a_five_local_all_rows_independently_verified.json}
is \texttt{PASS} with precisely the displayed counts and no
unresolved pair.  The source census and certificate hashes
are recorded there, so the finite scope is explicit.
\end{proof}

\begin{corollary}[Local quotients with five outside elements]
\label{C:local-five-two-column-act-cones}
Let $M$ have rectangular-band minimal ideal of shape $3\times2$,
with exactly five outside elements, and suppose
$\mathbb QM/\mathbb QK$ is local.
For every nonempty finite right $M$-set $S$,
$B\mathcal C(M,S)$ cannot have the rational homology and
cohomology ring of $\mathbb CP^2$.
In particular it cannot be homotopy equivalent to $\mathbb CP^2$.
\end{corollary}
\begin{proof}
A $K$-fixed state makes every degree-two square zero, as before.
Otherwise Lemma~\ref{C:two-column-projective-orbit-kernel} and
constant-action compression give a monoid $N$ with the same
homology as the act cone in every degree at least three.
Its contracted algebra and three-point row augmentation module
are the original $A=\mathbb QM/\mathbb QK$ and $V$.
The band--Tor formula would give
$\operatorname{Tor}_1^A(W,V)=0$,
$\dim\operatorname{Tor}_2^A(W,V)=1$, and zero Tor in every
higher degree, for the new column augmentation module $W$.
Proposition~\ref{C:local-five-three-row-Tor} excludes this.
Only the rational homology profile was used in this branch.
\end{proof}

\subsection{The last five-outside three-row act cones through nine states}

\begin{lemma}[The nine-state Nakayama calculation]
\label{C:Nakayama-eleven-nine-state}
Let $M$ have a $3\times2$ rectangular-band minimal ideal $K$ and
exactly five elements outside $K$.  Suppose its contracted rational
algebra is the Nakayama algebra
\[
 A=\mathbb QM/\mathbb QK,\qquad
 a\in e_1Ae_2,\quad b\in e_2Ae_1,\quad ab=c,\quad ba=0,
\]
and its two-dimensional row augmentation module is
$V=L=Ae_1/\mathbb Qc$.
There is no nine-element right $M$-set $S$ for which
$B\mathcal C(M,S)\simeq\mathbb CP^2$, with $SK$ consisting of
two two-element orbits.
\end{lemma}
\begin{proof}
We describe the complete finite calculation, its necessary algebraic
filters, and the actual cup certificates.  This is a calculation on
literal monoids and acts, not just on their linear representations.

First discard a split row augmentation module by
Proposition~\ref{C:split-row-act-cups}.  If $H^2(BM;\mathbb Q)=0$,
every degree-two class of the act cone is relative to $BM$; the
quotient by $BM$ is a suspension, so all its squares vanish.
Thus a possible source has a nonzero balanced form between its
row and column augmentation modules.  Since there are two columns,
such a form is a nonzero scalar-valued row functional.  Give the
three rows its affine values.  If all three values are distinct,
the finite affine self-maps of these values are the identity,
a possible common reflection, and constants.  Their augmentation
operators commute and diagonalize over $\mathbb Q$, so the row
module is split.  Otherwise relabeling and rescaling makes the
three values $(0,0,1)$.

Write an actor's row and column transformations as $\sigma$ and
$\tau$.  Compatibility with this functional is exactly
\[
 [\sigma(i)=2]-[\sigma(0)=2]
   =(\tau(1)-\tau(0))[i=2]\qquad(i=1,2).
\]
There are precisely $24$ pairs of transformations satisfying these
equations.  The generator
\path{code/a_eleven_three_by_two_targets.cpp} chooses every
multiset of four such pairs for the four nonidentity outside
elements: $\binom{27}{4}=17550$ profiles.  It discards $3550$
profiles whose row augmentation matrices commute and satisfy
$X^3=X$, which proves simultaneous splitting into rational
character lines.  For every remaining profile it tries every
outside product with the required row and column transformations,
and imposes all associativity equations.  Finally it takes orbits
under row, column, and outside-label permutations.  This gives
$120208$ associative sorted tables and $16399$ ordinary
isomorphism classes.  No passage to opposite monoids is made.
The source tables are retained in
\path{code/outputs/a_eleven_three_by_two_targets.tsv}.

The fourteen possible distinguished Nakayama bases have already
been independently classified in
Lemma~\ref{B:five-Nakayama-common-annihilator}.
The exact extraction
\path{code/a_nakayama11_sources.py} compares each contracted
table to all labelings of these bases.  It finds $93$ targets:
in $50$ the row module is the projective $Ae_2$, and in the
remaining $43$ it is $L$.  The projective case has zero
degree-two Tor and cannot give the fourth homology of a projective
plane.  Each of the $43$ records stores the full table, the
two rational characters, the row action, and exact distinguished-
basis coordinates of $e_1,b,c$.

Now suppose that a nine-state candidate is a projective plane.
Lemma~\ref{C:two-column-projective-orbit-kernel} identifies its
constant-action quotient with a compressed monoid whose column
augmentation module $W$ has dimension six and Tor dimensions
\[
 (\dim\operatorname{Tor}_0^A(W,L),
   \dim\operatorname{Tor}_1^A(W,L),
   \dim\operatorname{Tor}_2^A(W,L))=(2,0,1).
\]
The exact indecomposable table in
Lemma~\ref{B:five-exceptional-targets} forces
\[
 W=T_2\oplus e_1A\oplus e_2A.
\]
In particular $\dim We_1=\dim We_2=3$,
$\operatorname{rank}R_b=2$, and
$\operatorname{rank}R_c=1$.  If $\chi_1,\chi_2$ are the two
characters, the action on the seven-element quotient set must
therefore have
\[
 \#\operatorname{Fix}(m)=1+3\chi_1(m)+3\chi_2(m)
\]
fixed points for every actor $m$.  These are necessary conditions;
they are not used as substitutes for the final cup computation.

Here is the complete state normal form.  Label the core orbits
$\{0,1\}$ and $\{2,3\}$, with the second coordinate giving
the two columns, and label the other five states $4,\ldots,8$.
For each state $s$ let
$\gamma_s:\{0,1,2\}\to\{0,1\}$ record the core orbit reached
by the three rows.  A band element $(i,j)$ acts by
\[
 s(i,j)=2\gamma_s(i)+j.
\]
The two core orbits have constant $\gamma$ values zero and one.
Each outside actor preserves both core orbits and acts there by
its column transformation.  On the other states it must satisfy
\[
 \gamma_{sm}(i)=\gamma_s(\sigma_m(i)).
\]
Sort the five three-bit $\gamma$ values, retaining every map
within each equal-value fiber.  This gives $\binom{12}{5}=792$
profiles per source.  If all five values are constant, the act
splits into two invariant components and cannot have the required
first homology; discarding this case is valid.

The finite constraint search
\path{code/a_nakayama11_nine_states.cpp} assigns all twenty
unknown values of the four outside actors on the five remaining
states.  Its domains are exactly the displayed $\gamma$ fibers.
It propagates the fixed-point counts and every equation
$(sm)n=s(mn)$, branches on every value left in a domain, and
checks all $11^2\cdot9$ action equations on every completed act.
The $43\cdot792=34056$ profiles give $2119812$ completed acts
after the trace filter.  The ranks of $R_b,R_c$ leave $131144$
records.  Their six-by-six integer matrices have entries of
absolute value at most five.  Their ranks computed modulo
$1000000007$ are exact rational ranks: every minor is bounded
in absolute value by
$6^3\cdot5^6=3375000$, strictly less than that prime.

Permuting the two core orbits and the five other states gives
$2\cdot5!=240$ possible normal-form conjugacies.
The action of the three column-zero band elements and the four
outside actors determines the entire act; column-one band images
are obtained by adding one to the column-zero images.
The complete state-orbit calculation leaves $9139$ representatives
on $40$ of the $43$ targets.  All raw acts and representatives
are retained, so this reduction has an explicit coverage check.
The independent checker
\path{code/d_nakayama21_orbit_coverage.cpp} instead uses all
eleven actor maps as its key and expands all $240$ conjugates
of every retained representative.  Its $2193360$ lookups cover
all $131144$ raw acts, with disjoint retained orbits.
The separate \path{code/d_nakayama21_linkage.py} verifies all
normal-form equations and the literal source-to-category tables;
both saved result files report \texttt{PASS}.

For each representative the converter
\path{code/a_nakayama21_categories.py} builds the literal
$21$-arrow composition table, with identities first, and checks
the monoid and action equations again.  The normalized nerve has
\[
 (\dim C_1,\dim C_2,\dim C_3,\dim C_4)
          =(19,190,1900,19000)
\]
over $\mathbb F_2$.
The exact checker \path{code/a_nakayama21_cups.cpp} verifies
all typed associativity equations, computes a nonboundary
degree-two cocycle $u$, and solves
\[
 \delta w=u\smile u
\]
using the Alexander--Whitney product.  Every one of the $9139$
cases has $H^1=0$, $\dim H^2=1$, $H^3=0$, and such a primitive.
It then checks the primitive on every four-simplex.  In total
this checks $23112531$ typed associativity equations,
$17364100$ cocycle equations, and $173641000$ primitive equations.

The full tables and bit-vector cochains are in
\path{code/outputs/a_nakayama21_cups.jsonl}.
The separate checker \path{code/a_verify_nakayama21_cups.cpp}
imports none of the search or linear-system code: it reconstructs
the normalized chains, computes ranks using the opposite pivot
order, verifies that $u$ is not a coboundary, and evaluates every
stored cocycle and primitive equation directly.
Its full result in
\path{code/outputs/a_nakayama21_cups_independently_verified.json}
is \texttt{PASS} on all $9139$ categories and all the displayed
equations.
Thus the nonzero mod-two square required for $\mathbb CP^2$
is absent in every possible candidate.  This proves the lemma.
\end{proof}

\begin{theorem}[Five outside elements through nine states]
\label{C:eleven-three-row-nine-state}
Let $M$ have a $3\times2$ rectangular-band minimal ideal and
exactly five outside elements.  For every nonempty finite right
$M$-set $S$ with $|S|\leq9$,
\[
 B\mathcal C(M,S)\not\simeq\mathbb CP^2.
\]
\end{theorem}
\begin{proof}
A state fixed by the minimal band gives zero degree-two squares.
Otherwise let $o$ be the number of two-point core orbits and
apply Lemmas~\ref{C:two-column-projective-orbit-kernel} and
\ref{C:two-column-finite-quotient-model}.  The compressed monoid
has the same contracted algebra $A$ and row module $V$ as $M$,
and its column augmentation module $W$ satisfies
\[
 \dim V=2,\qquad \dim W=|S|-2o+1,
 \qquad(\dim\operatorname{Tor}_0,\dim\operatorname{Tor}_1,
        \dim\operatorname{Tor}_2)=(o,0,1),
\]
with zero higher Tor.  The local case is excluded for all $S$
by Corollary~\ref{C:local-five-two-column-act-cones}.
Proposition~\ref{B:five-nonlocal-exceptions} leaves, in the
nonlocal case, only the zero-composite line algebra and the
Nakayama algebra.

For the line algebra the degree-one and degree-two conditions
force $V=S_1\oplus S_3$, by the indecomposable calculation in
Lemma~\ref{B:five-exceptional-targets}, independently of
$\dim W$.  This is a split row module, so
Proposition~\ref{C:split-row-act-cups} applies.
For the Nakayama algebra the same table forces
\[
 V=L,\qquad W=T_2\oplus(e_1A)^{\oplus p}
                    \oplus(e_2A)^{\oplus q},\qquad
 o=p+q.
\]
Comparing dimensions gives
\[
 |S|=4o+p.
\]
If $o=1$, there is one core orbit, hence every band actor is
constant on all of $S$.  The formula permits only $|S|=4$ or
$5$; compression would give a projective-plane monoid of order
$5+3|S|\leq20$, contradicting its exact minimum $23$.

If $o\geq2$ and $|S|\leq9$, the only possibilities are
$(o,p,q,|S|)=(2,0,2,8)$ and $(2,1,1,9)$.
In the first case the compressed monoid has minimal band
$3\times6$ and the Nakayama module pair
$W=T_2\oplus(e_2A)^2$, $V=L$.
The actual-basis quotient argument in the proof of
Proposition~\ref{B:three-by-six-five-outside} makes every
rational degree-two product zero for this pair: its faithful
action quotient has hereditary contracted algebra, and the
degree-two map is an isomorphism.  This part of that proof
does not require the nondegenerate intersection form of its
named product target.  All degree-two classes of the original
act cone pull back from the compressed monoid, so they too
have zero products.  The nine-state case is precisely
Lemma~\ref{C:Nakayama-eleven-nine-state}.
\end{proof}

\begin{corollary}[Small endomorphism monoids through twenty-one arrows]
\label{C:eleven-act-cones21}
An act-cone category $\mathcal C(M,S)$ with $|M|\leq11$ and
$|M|+|S|+1\leq21$ cannot have classifying space
$\mathbb CP^2$.
\end{corollary}
\begin{proof}
Theorem~\ref{C:nine-act-cones} deals with $|M|\leq9$.
A minimal ideal with at most two rows is excluded by
Theorem~\ref{C:two-row-act-cups}.  A minimal ideal with one
column makes the monoid group type; its rational degree-two
cohomology is zero, so relative products in the act cone vanish.
Thus at least three rows and two columns are necessary.
At orders ten and eleven a nontrivial minimal maximal group
would require at least $3\cdot2\cdot2$ ideal elements, already
too many.  The minimal ideal is therefore a band.

The possible shapes are $3\times2$, $4\times2$, $5\times2$,
and $3\times3$, with the identity counted outside the band.
All two-column cases with at most four outside elements are
excluded by Theorem~\ref{C:two-column-four-outside-act-cones}.
The remaining two-column case is $3\times2$ with five outside
elements; here $|S|\leq9$, and the preceding theorem applies.

For the $3\times3$ band the contracted rational algebra has
dimension at most two.  It is either semisimple or is the dual
number algebra $\mathbb Q[\epsilon]/(\epsilon^2)$, as follows
directly from the multiplication of its at most one nonidentity
nonzero basis element.  The two-dimensional row augmentation
module is then either projective over the contracted algebra or
a sum of one-dimensional modules.  The latter case is covered
by Proposition~\ref{C:split-row-act-cups}.  In the former case,
the projective minimal-band ideal gives projective dimension at
most one for that row module over $\mathbb QM$.  The exact
augmentation sequence with projective middle term gives
$\operatorname{pd}_{\mathbb QM}\mathbb Q\leq2$.
The act-cone Tor formula therefore makes its fourth homology
zero.  This exhausts the shapes at these orders.
\end{proof}

\subsection{An affine primitive for act-cone products}

\begin{theorem}[A common fixed row for nonconstant column actors]
\label{C:act-common-fixed-row-cups}
Let $M$ have rectangular-band minimal ideal $I\times\{0,1\}$.
Write $\sigma_m:I\to I$ and
$\tau_m:\{0,1\}\to\{0,1\}$ for its row and column actions.
Suppose there is a row $i_0$ fixed by every $\sigma_m$ for which
$\tau_m$ is nonconstant.  Then for every nonempty finite right
$M$-set $S$, all products
\[
 H^2(B\mathcal C(M,S);\mathbb Z)\otimes
 H^2(B\mathcal C(M,S);\mathbb Z)
       \longrightarrow H^4(B\mathcal C(M,S);\mathbb Z)
\]
are zero.  In particular no such act cone is a projective plane.
There is no bound on the number of rows, states, or outside elements,
and no finiteness hypothesis on the classifying space.
\end{theorem}
\begin{proof}
By Lemma~\ref{C:constant-action-quotient-H2}, every integral
degree-two class pulls back from the constant-action quotient
$\bar S$.  Compress that quotient to a monoid $N$ with minimal
band $I\times\bar S$, as in
Theorem~\ref{C:constant-minimal-act-compression}.
Let $s_j\in\bar S$ be the constant value of an original band
element in column $j$.  This is independent of its row, and
\[
 s_jm=s_{\tau_m(j)}.
\]
The two values are allowed to coincide.  Use the base row $i_0$
and base column $s_0$ throughout.

Put $V=\widetilde{\mathbb ZI}$ and
$W=\widetilde{\mathbb Z\bar S}$.
By Lemma~\ref{lem:B-band-Tor}, integral degree-two classes on
$BN$ are precisely integral balanced bilinear forms
$L:W\times V\to\mathbb Z$.
Here $H_1(BN)=0$, so the universal-coefficient identification
has no additional degree-two extension term.
For an actor $f$ of $N$ define
\[
 a_f=[\sigma_f(i_0)]-[i_0]\in V,\qquad
 b^L_f(v)=L([\tau_f(s_0)]-[s_0],v).
\]
Let $A_f$ be its row augmentation operator.  The affine identities
and balance give
\[
 a_{fg}=a_f+A_fa_g,\qquad
 b^L_{fg}=b^L_g+b^L_fA_g.
\]
Consequently $u_L(f,g)=b^L_fa_g$ is the usual integral
two-cocycle representing $L$.

We make its transfer to the act category literal.
Put $e=(i_0,s_0)$ in the band of $N$, and take the full
two-object Karoubi subcategory on $1,e$.
The compression functor sends the monoid object to $1$, the
sink to $e$, an old actor $m$ to its image in $N$, and a cross
state $s$ to $(i_0,s)$.  Extend $u_L$ to this subcategory by the
same formula on the underlying $N$ labels.
Both object identities $1,e$ have $a=b^L=0$, so this is a
normalized category cocycle; its coboundary is zero by the
same affine identities.  Its restriction to the object $1$
is the original monoid cocycle.  The inclusion of this object
is a classifying-space equivalence, so the extension represents
the unique corresponding cohomology class.  Pulling back along
the honest compression functor gives the desired representative
on $\mathcal C(M,\bar S)$.
For a cross arrow its row vector $a$ is zero and its column
functional is $L([s]-[s_0],-)$.  All other labels are old actors.

Define the column character and translation by
\[
 \chi(m)=\tau_m(1)-\tau_m(0),\qquad
 \beta(m)=\tau_m(0).
\]
They obey
\[
 \beta(mn)=\beta(n)+\chi(n)\beta(m).
\]
For each balanced form $L$, put
\[
 \ell_L(v)=L([s_1]-[s_0],v),\qquad
 \alpha_L(m)=\ell_L(a_m).
\]
The equivariant map from the original column augmentation line
to $W$ gives $\ell_LA_m=\chi(m)\ell_L$.  Therefore
\[
 \alpha_L(mn)=\alpha_L(m)+\chi(m)\alpha_L(n),
 \qquad
 b^L_m=\beta(m)\ell_L,
 \qquad
 u_L(m,n)=\beta(m)\alpha_L(n)
\]
on old monoid arrows.
The common-fixed-row hypothesis is exactly
\[
 \chi(m)a_m=0\qquad(m\in M).
\]

Use algebra-ordered normalized chains, so their adjacent
composition is multiplication from left to right.  A chain in
an act category either consists entirely of old monoid arrows
or has one cross arrow in its first position; all subsequent
arrows belong to $M$.  For two balanced forms $L,D$ define an
integral three-cochain to be zero on every word containing an
identity, and on the nondegenerate typed domain define it by
\[
 P(f,g,h)=u_L(f,g)\,\beta(g)\,\alpha_D(h).
\]
This prescription is normalized; on a nondegenerate word the
arguments of $\beta$ and $\alpha_D$ are always old actors.
Expand its five-term coboundary on
$(f,g,h,k)$ and use the two-cocycle identity for $u_L$:
\begin{align*}
 (\delta P)(f,g,h,k)
 &=u_L(f,g)\beta(h)\alpha_D(k)\\
 &\quad+\chi(h)\beta(g)\alpha_D(k)
       \bigl(u_L(f,gh)-u_L(f,g)\bigr)\\
 &=u_L(f,g)u_D(h,k)
   +\chi(h)\beta(g)\alpha_D(k)b^L_f A_g a_h\\
 &=u_L(f,g)u_D(h,k).
\end{align*}
The last equality uses $\chi(h)a_h=0$.
Thus $\delta P=u_L\smile u_D$ as an equality of integral
normalized category cochains.  Every degree-two class is
represented this way after the quotient pullback, proving the
assertion about all mixed products.
\end{proof}

\begin{corollary}[All dual-number three-row act cones]
\label{C:dual-row-all-act-cups}
Suppose $M$ has a $3\times2$ rectangular-band minimal ideal and
the image of its rational row augmentation representation is
the nonsemisimple dual-number algebra.
For every finite right $M$-set $S$, all integral degree-two
products on $B\mathcal C(M,S)$ are zero.
In particular no such cone realizes $\mathbb CP^2$, regardless
of the sizes of $M$ and $S$.
\end{corollary}
\begin{proof}
Normalize the row labels as in
Proposition~\ref{B:three-row-dual-cup-radical}.
The actual row transformations are
\[
 (0,1,2),\quad(0,0,1),\quad(1,1,0),\quad
 (0,0,0),\quad(1,1,1),\quad(2,2,2).
\]
Their augmentation operators are $I,N,-N,0,0,0$, with
$N=\left(\begin{smallmatrix}0&1\\0&0\end{smallmatrix}\right)$.
Write $\psi$ for their common diagonal character; it is one
only for the identity row map and zero otherwise.
The only nonzero one-dimensional quotient of the rational row
module has character $\psi$, and every such quotient kills
$\operatorname{im}N$.

Retain the notation of the preceding proof for an arbitrary
integral balanced form $L$ after constant-action compression.
Its old-column covector $\ell_L$ satisfies
$\ell_LA_m=\chi(m)\ell_L$.
If the original column character $\chi$ differs from $\psi$,
every such covector is zero, since no quotient of that character
exists over $\mathbb Q$.  Hence every transferred cocycle
$u_L(m,n)$ is literally zero on pairs of old monoid arrows.
On every four-chain the second factor of a degree-two cup
product is such a pair; all products are therefore already
zero as cochains.

If $\chi=\psi$, each nonconstant column actor has $\chi=1$,
and its row transformation is the identity.  Every row is a
common fixed row, and
Theorem~\ref{C:act-common-fixed-row-cups} applies.
Both arguments use integral cochains and the integral
degree-two quotient surjection; the rational representation
classification introduces no denominator into the primitive.
\end{proof}

\begin{remark}[A completed finite control and a superseded search]
\label{C:dual-row-act-control-scope}
Before the preceding primitive was found, a bounded upper-model
search was completed for one explicit six-dimensional contracted
basis: a square-zero loop at the first vertex of a zero-composite
three-vertex chain.  Its $72$ literal $3\times2$ affine lifts and
$144$ trace/rank cases gave $593536$ eight-state acts after the
trace constraints.  Of these, $275040$ had the required exact
Tor profile, and $90240$ retained a possible rank-two balanced
class.  The $10236$ state-orbit representatives all had zero
mod-two square, with actual stored primitives checked on
$258858204$ four-simplices.
The code and outputs have prefix \path{a_dual_loop_chain}.
This finite control is not used in the all-order proof.

The broader restricted order-twelve dual-row target stream
\path{code/outputs/a_twelve_dual_targets.tsv} was intentionally
interrupted when the all-order theorem superseded it.  It is
retained as an incomplete research prefix, not a complete census
and not a premise of any exclusion.
\end{remark}

\subsection{Forced affine profiles over a three-by-two minimal band}

\begin{proposition}[A common rational affine origin]
\label{C:act-common-affine-origin-cups}
Let $M$ have a completely simple minimal ideal with row set $I$
and two columns.  Its minimal maximal group may be nontrivial.
Suppose that the row maps of all actors with nonconstant column
map fix a common point
\[
 p\in\mathbb QI,\qquad \sum_{i\in I}p_i=1.
\]
Then all rational degree-two products of every finite act cone
$B\mathcal C(M,S)$ vanish.  The point $p$ need not be a row or
have nonnegative coordinates.
\end{proposition}
\begin{proof}
Use the rational constant-action quotient of
Lemma~\ref{C:constant-action-quotient-H2} and the compression of
Theorem~\ref{C:constant-minimal-act-compression}.
This gives a band $I\times\bar S$ even when the original
minimal maximal group is nontrivial.  As proved in that
compression theorem, the constant value of an old minimal-ideal
element depends only on its column: left multiplication varies
the row and group coordinate arbitrarily while preserving that
value.  The two old columns therefore give the same equivariant
map $j\mapsto s_j$ as in
Theorem~\ref{C:act-common-fixed-row-cups}.
Use its honest Karoubi transfer and affine cocycles over
$\mathbb Q$.
Retain its old base row $i_0$, base column $s_0$, covectors $b_f$,
and augmentation operators $A_f$.  Put $t=p-[i_0]$.
For an old monoid arrow replace its row translation by
\[
 a'_m=a_m+(A_m-1)t=\sigma_m p-p.
\]
Put $a'_f=0$ on cross arrows and on the sink identity.  These
assignments are typed: the monoid object uses affine origin $p$
and the sink uses origin $i_0$.  A cross arrow has constant row
map, so its augmentation operator is zero.  Thus on every
composable pair the same affine identity
$a'_{fg}=a'_f+A_fa'_g$ holds.

The cochain $u'_L(f,g)=b^L_fa'_g$ represents the original
degree-two class.  Indeed, on normalized composable pairs its
difference from $u_L$ is
\[
 b^L_f(A_g-1)t=-\delta h(f,g),\qquad h(f)=b^L_ft,
\]
where $h$ is zero on identities.  This equation includes the
cross-arrow/old-arrow pairs; it follows from
$b^L_{fg}=b^L_g+b^L_fA_g$.
The column character $\chi$ now satisfies
$\chi(m)a'_m=0$: for $\chi(m)\ne0$ this is the fixed-point
hypothesis, and otherwise it is automatic.
Define $\alpha'_L(m)=\ell_L(a'_m)$ and use the same normalized
primitive
\[
 P(f,g,h)=u'_L(f,g)\beta(g)\alpha'_D(h)
\]
on nondegenerate chains, extended by zero on degenerate chains.
The calculation in Theorem~\ref{C:act-common-fixed-row-cups}
has defect
$\chi(h)\beta(g)\alpha'_D(k)b^L_fA_ga'_h=0$.
It therefore proves $\delta P=u'_L\smile u'_D$ over
$\mathbb Q$.  The degree-two quotient surjection proves the
assertion for the original act.
\end{proof}

\begin{lemma}[Four necessary outside profiles]
\label{C:three-row-forced-profiles}
Suppose $M$ has a minimal ideal with three rows and two columns and
some finite right act cone has a nonzero rational degree-two
product.  After relabelling rows, $M$ has actors with each of
the following four distinct row/column profiles:
\[
\begin{array}{c|c|c}
 &\sigma&\tau\\ \hline
 p_0&(0,0,2)&(0,1)\\
 p_1&(1,1,2)&(0,1)\\
 r_0&(2,2,0)&(1,0)\\
 r_1&(2,2,1)&(1,0).
\end{array}
\]
Every actor satisfies the balance equation for row values
$(0,0,1)$.  There are exactly $24$ possible row/column
transformation pairs satisfying that equation.  In particular,
the four displayed profiles and the identity require at least
five elements outside the minimal band.
\end{lemma}
\begin{proof}
First $H^2(BM;\mathbb Q)\ne0$.  Otherwise every degree-two
class of the act cone lifts to cohomology relative to $BM$.
The relative quotient is a suspension of the action space,
whose positive-degree products vanish.  Relative cup naturality
would make every degree-two product of the cone zero.

Collapse its minimal maximal group using
Lemma~\ref{lem:B-collapse-minimal-group}; this preserves rational
cohomology and all row/column actions.  The balanced-tensor
description of $H^2$ of the resulting band monoid supplies
a nonzero covector on the row augmentation module with
\[
 \ell A_m=\chi(m)\ell,\qquad
 \chi(m)=\tau_m(1)-\tau_m(0)\in\{0,1,-1\}.
\]
Write the three row values as $v_i$, up to common translation.
If they are all distinct, an actor with $\chi=1$ fixes all three
rows: an affine translation preserving a finite set is the
identity.  An actor with $\chi=-1$, if present, is the unique
reflection of that three-point set and fixes its middle row.
The preceding proposition applies, a contradiction.
Thus there are two levels.  Relabel and rescale them as
$(v_0,v_1,v_2)=(0,0,1)$.

An actor with $\chi=1$ fixes row $2$ and preserves $\{0,1\}$.
An actor with $\chi=-1$ has row map either $(2,2,0)$ or
$(2,2,1)$, and its column map is the swap.  If neither occurs,
row $2$ is a common fixed row.  If only the first occurs,
every $\chi=1$ actor also fixes row $0$: composing its row map
on the left with the first displayed flip would otherwise
produce the second flip.  Then
$( [0]+[2])/2$ is a common rational affine fixed point of every
actor with nonconstant columns.  The other single-flip case
has fixed point $( [1]+[2])/2$.  Each possibility contradicts
Proposition~\ref{C:act-common-affine-origin-cups}.
Both flips are therefore necessary.  Their squares and mixed
products force $p_0,p_1$.

For completeness, the profile count is literal.  There are four
row maps for the identity column map, two for the swap, and
nine for each constant column map.  In the latter case all row
images must lie in one level, giving $2^3+1=9$ choices.
Hence the count is $4+2+9+9=24$.
Minimal-ideal actors have constant column maps, whereas the four forced
profiles and the global identity do not; the five are distinct.
\end{proof}

\begin{lemma}[Projective row modules obstruct every act]
\label{C:row-pd-one-act-cups}
Let $M$ be a finite monoid with minimal ideal $K$, row set $I$,
contracted rational algebra $A=\mathbb QM/\mathbb QK$, and row
augmentation module $V=\widetilde{\mathbb QI}$.
If $\operatorname{pd}_A V\leq1$, then every finite right
$M$-act cone has zero rational degree-two products.
No bound on the act or its homology is required.
The assertion also permits a nontrivial minimal maximal group.
\end{lemma}
\begin{proof}
Pass to the constant-action quotient $\bar S$.
Lemma~\ref{C:constant-action-quotient-H2} gives a surjection on
rational degree-two cohomology in the pullback direction.
Compress by Theorem~\ref{C:constant-minimal-act-compression}.
The resulting monoid $N$ has minimal band $I\times\bar S$.
Its outside elements and their products modulo the minimal
ideal are literally those of $M$, so its contracted algebra
is $A$ and its row module is $V$.
Writing $\bar W=\widetilde{\mathbb Q\bar S}$, the band
dimension-shift formula gives
\[
 H_4(BN;\mathbb Q)
   =\operatorname{Tor}^{A}_2(\bar W,V)=0.
\]
All degree-two products on $BN$ are zero, and cup naturality
and the degree-two surjection give the result.  If $K=M$,
then $M$ is a group, $V=0$, and the rational group-algebra
averaging argument gives the same conclusion directly.
\end{proof}

\begin{proposition}[Five outside elements, with no state bound]
\label{C:three-row-five-outside-all-acts}
If a finite monoid has a minimal ideal with three rows and two
columns and at most five elements outside it, every finite act cone has zero
rational degree-two products.
\end{proposition}
\begin{proof}
If the necessary profiles of
Lemma~\ref{C:three-row-forced-profiles} are absent, that lemma
already proves the assertion.  Otherwise there are exactly
five outside elements: the identity and the four profiles
$p_0,p_1,r_0,r_1$.  Their products cannot land in the band,
because their column maps are nonconstant, and there is only
one element of each profile.  Their literal multiplication is
\[
 (i,\epsilon)(j,\delta)=(i,\epsilon+\delta),
 \qquad i,j\in\{0,1\},\quad\epsilon,\delta\in C_2,
\]
with $(i,0)=p_i$ and $(i,1)=r_i$.
Thus $A$ is the unitalization of $\mathbb Q[L_2\times C_2]$.
Put
\[
 e_\pm=\frac{p_0\pm r_0}{2},\qquad f=1-p_0,\qquad
 n_\pm=\frac{p_1-p_0\pm(r_1-r_0)}2.
\]
The three idempotents $e_+,e_-,f$ are orthogonal and sum to
one.  The radical has basis $n_+,n_-$, with
$n_\pm\in fAe_\pm$ and radical square zero.  This is the
path algebra of a three-vertex fork, hence hereditary.
More directly the row module itself is $Ae_-$: in the row
basis $[1]-[0],[2]-[0]$, the element $e_-$ projects onto
$[2]-[0]$ and $n_-$ sends that vector to $-[1]+[0]$.
The resulting two-dimensional equivariant map $Ae_-\to V$
is an isomorphism.  Lemma~\ref{C:row-pd-one-act-cups}
finishes the proof.
\end{proof}

\begin{lemma}[The complete forced-profile row certificates]
\label{C:forced-profile-row-certificates}
Among ordinary isomorphism classes of monoids with minimal
band $3\times2$ and all four forced profiles, the following
table is complete.  ``Simple periodic kernel'' means an exact
sequence $0\to T\to P\to V\to0$ with $P$ projective and a
one-dimensional module $T$ admitting
$0\to T\to Q\to T\to0$ with $Q$ projective.
\[
\begin{array}{c|r|r|r|r|r}
 |M|&\text{raw tables}&\text{classes}&
  V\text{ projective}&\operatorname{pd}V\leq1\text{, extra}&
  \text{simple periodic kernel}\\ \hline
 11&1&1&1&0&0\\
 12&27&10&10&0&0\\
 13&786&152&148&4&0\\
 14&32188&2980&2908&41&31
\end{array}
\]
Here raw tables use sorted outside profiles before taking
ordinary row/column/element isomorphisms; opposition is not
used to identify different orientations.
\end{lemma}
\begin{proof}
We give the finite classification and the exact module
certificates, so the scope does not depend on a published census.
Fix identity label $0$ and the six row-major band labels.
For order $n$ select the four required profiles and a multiset
of $n-11$ further profiles from the $24$ possibilities above.
The numbers of multisets for $n=11,12,13,14$ are
$1,24,300,2600$.
The identity and all products involving the band are forced.
For each ordered pair of other elements, try every label whose
profile is their composite, including a band label whenever
allowed.  Associativity propagates known products and rejects
only an incompatible equality.  A finished table is checked
on every triple.  Every actual monoid under consideration is
obtained after relabelling its outside elements, proving
coverage of this search.

The code \path{a_forced_swap_targets.cpp} implements this
enumeration.  Its fast version first sorts the transformed
outside profiles under each row/column permutation and then
tries all permutations within equal-profile fibers.  This is
a complete ordinary-isomorphism invariant; it need not be the
unrestricted lexicographic minimum.
The separate program \path{b_forced_profile_raw.cpp} repeats
the raw enumeration with reverse-cell branching and no producer
canonicalizer.  The separate
\path{b_verify_forced_profile_coverage.py} expands every retained
class under all row, column, and equal-profile permutations.
These literal orbits are disjoint and cover the independently
generated raw tables exactly.  It requires final metadata,
exact counts, and complete line lengths before comparison.
The resulting certificate is
\path{b_forced_profile_coverage_verified.json}.

For the module columns, work over the literal contracted basis
$A=\mathbb QM/\mathbb QK$.  The stored row action is the exact
two-dimensional augmentation action.  Each cover is a left
ideal $A\epsilon$, with
$\epsilon=(e-g)/2$ for a literal local pair
$e^2=e$, $g^2=e$, $eg=ge=g$.
The certificate stores the full basis of this ideal, the
equivariant surjective cover to $V$, and its embedded kernel.
For a one-dimensional kernel with character $\theta$, a
projective section is exhibited by an element $z\in A$ with
\[
 az=\theta(a)z\quad\text{on every basis element},\qquad
 \theta(z)=1.
\]
These equations imply $z^2=z$ and identify the kernel with
the projective left ideal $Az$.
The files \path{a_forced_swap_rational_rows12.json},
\path{a_forced_swap_rational_rows13.json}, and
\path{a_forced_swap_rational_rows14.json} store every cover.
The independent Fraction-elimination verifier
\path{d_verify_forced_swap_row_projects.py} imports no producer
or Flint routines.  It checks the literal monoids, ideals,
contracted products, idempotents, full left ideals, surjectivity,
equivariance, kernels, and all displayed projective sections.

There are exactly $31$ remaining one-dimensional kernels at
order $14$.  For each, the file
\path{a_forced_swap_simple_tails14.json} supplies an idempotent
$\eta$, a basis of the two-dimensional ideal $A\eta$, its
surjection to the kernel character, and a one-dimensional
kernel with exactly that same character.  Thus these are
literal exact sequences $0\to T\to A\eta\to T\to0$,
not a periodicity inference from ranks.  The separate Fraction
verifier \path{code/d_verify_forced_swap_simple_tails.py}
reconstructs every full ideal, cover, and kernel and checks the
character equality without importing producer routines or
trusting a recurrence status.  Its final output is
\path{code/outputs/d_forced_swap_simple_tails_verified.json}.
\end{proof}

\begin{theorem}[Every three-by-two act cone through order fourteen]
\label{C:three-row-fourteen-act-cones}
Let $M$ have minimal band $3\times2$.
If $|M|\leq13$, all rational degree-two products of every
finite right act cone vanish, without any homology bound.
If $|M|\leq14$ and the act cone has bounded rational homology,
its rational degree-two products also vanish.
In particular no act cone over such a monoid realizes
$\mathbb CP^2$, at any number of states.
\end{theorem}
\begin{proof}
If the four forced profiles are absent, use
Lemma~\ref{C:three-row-forced-profiles}.
For at most five outside elements use
Proposition~\ref{C:three-row-five-outside-all-acts}.
The row certificates at orders $12$ and $13$, and all but
$31$ cases at order $14$, have projective dimension at most
one, so Lemma~\ref{C:row-pd-one-act-cups} applies.

In one of the remaining $31$ cases, the two certified exact
sequences imply, for every finite-dimensional right module $W$,
\[
 \operatorname{Tor}^{A}_n(W,V)
 \cong\operatorname{Tor}^{A}_{n-1}(W,T)
 \cong\operatorname{Tor}^{A}_1(W,T)
 \qquad(n\geq2).
\]
If the act has a state fixed by the minimal band, its
degree-two products already vanish by the band-grid relative
argument of Lemma~\ref{C:four-state-band-dichotomy}.
Otherwise Lemma~\ref{C:two-column-projective-orbit-kernel}
identifies the original and constant-action-quotient homology
in every degree at least three.  Compress the latter to $N$.
It has the same contracted algebra $A$ and row module $V$,
so the band dimension shift and the displayed recurrence give
\[
 H_q(BN;\mathbb Q)\cong H_4(BN;\mathbb Q)
 \qquad(q\geq4).
\]
Bounded rational homology forces this common group to be zero.
It follows that the original cone has zero fourth rational
homology and hence zero degree-two products, as required.
\end{proof}

\begin{corollary}[The whole three-by-two branch below twenty-two arrows]
\label{C:three-row-act-cones21}
No act cone $\mathcal C(M,S)$ with minimal band $3\times2$
and $|M|+|S|+1\leq21$ realizes $\mathbb CP^2$.
\end{corollary}
\begin{proof}
The preceding theorem handles $|M|\leq14$ at every state
count.  Otherwise $|S|\leq5$.  For at most two states use
the finite-line act argument; three states require at least
$22$ arrows by Theorem~\ref{C:three-state-cone-minimum22}.
For four states, the band dichotomy gives either a fixed state
and zero products, or constant action and a compressed
$3\times4$ monoid; the latter is excluded by
Theorem~\ref{B:three-by-four-targets}.
For five states apply
Theorem~\ref{C:five-state-three-by-two-cones}.
\end{proof}

\begin{corollary}[Nontrivial minimal groups below twenty-two arrows]
\label{C:nontrivial-minimal-group-act-cones21}
No finite act cone with a nontrivial minimal maximal group and
at most $21$ arrows realizes $\mathbb CP^2$.
\end{corollary}
\begin{proof}
At most two minimal rows are excluded by
Theorem~\ref{C:two-row-act-cups}.  One minimal column makes
$BM$ rationally acyclic, so the relative suspension argument
excludes that case.  At most three states are already excluded
below $22$ arrows.  Thus a remaining example has at least three
rows, two columns, four states, and at most $16$ monoid elements.
If its minimal maximal group has order $h\geq2$, its minimal
ideal has $rch$ elements and is proper.  The only possibility
under this bound is $r=3,c=2,h=2$: every other choice costs at
least $16$ minimal-ideal elements and one identity outside it.
The monoid consequently has at most four outside elements,
contrary to the necessary five of
Lemma~\ref{C:three-row-forced-profiles}.
\end{proof}

\subsection{Three small affine obstructions for arbitrary acts}

\begin{lemma}[A local contracted algebra with two columns]
\label{C:local-two-column-act-cups}
Let $M$ have a completely simple minimal ideal with two
columns, and suppose its contracted rational algebra
$A=\mathbb QM/\mathbb QK$ is local.  Every rational degree-two
product of every finite act cone of $M$ is zero.
There is no bound on the number of rows or states, and the
minimal maximal group may be nontrivial.
\end{lemma}
\begin{proof}
Every distinguished nonzero basis element of $A$ has a high
power that is either zero or a distinguished idempotent basis
element.  A local algebra has no idempotents other than zero
and one.  If that power is one, the original basis element is
an actual finite unit in the pointed quotient monoid $M/K$.
A nontrivial cyclic subgroup would give the idempotent average
of its distinct basis elements, which is neither zero nor one.
This is impossible.  Consequently every outside basis element
other than the identity is nilpotent.

The old column augmentation defines the character
$\chi(m)=\tau_m(1)-\tau_m(0)$ of $A$.
It is zero on every such nilpotent basis element.
All nonidentity actors therefore have constant column maps.
The identity is the only actor with nonconstant columns, so
any row is a common fixed row for those actors.
Proposition~\ref{C:act-common-affine-origin-cups} applies.
\end{proof}

\begin{lemma}[Two quotient covectors on at most five rows]
\label{C:two-covector-affine-origin}
Let $M$ have a minimal ideal with at most five rows and two
columns.  Put $V=\widetilde{\mathbb QI}$ and let $\chi$ be its
column augmentation character.  If
\[
 \dim_{\mathbb Q}\operatorname{Hom}_A(V,\mathbb Q_\chi)\geq2,
 \qquad A=\mathbb QM/\mathbb QK,
\]
then every rational degree-two product of every finite act
cone is zero.  No bounded-homology assumption is needed.
\end{lemma}
\begin{proof}
Choose two independent quotient covectors and evaluate them
on $[i]-[i_0]$.  This gives a map $f:I\to\mathbb Q^2$
whose image $F$ affinely spans the plane.  Each row actor has
\[
 f(\sigma_m i)=\chi(m)f(i)+b_m.
\]
For $\chi(m)=1$, the translation $b_m$ preserves the finite
set $F$ and hence is zero.  Such an actor preserves every
fiber of $f$.  Every actor with $\chi(m)=-1$ induces a central
reflection of $F$.  All these reflections are the same:
the composite of two is a translation of $F$, necessarily
the identity.

If there is no negative actor, then $|F|\geq3$ and $|I|\leq5$
give a singleton fiber.  Its row is fixed by every actor with
nonconstant columns.  If there is a negative actor, the
centrally symmetric set $F$ has at least four points, because
it spans a plane.  With at most five rows, at most one fiber
can contain more than one point.  Thus there is an opposite
pair of singleton fibers, or a singleton central fiber.
In the first case the midpoint of the two corresponding row
basis vectors is fixed by every nonconstant-column actor;
in the second case the corresponding row is fixed.
Either gives the common rational affine origin required by
Proposition~\ref{C:act-common-affine-origin-cups}.
\end{proof}

\begin{lemma}[Three outside elements, at every row and state count]
\label{C:three-outside-all-two-column-act-cups}
If a finite monoid has a minimal ideal with two columns and
at most three outside elements, all rational degree-two
products of all its finite act cones vanish.
\end{lemma}
\begin{proof}
Let $P$ be the actors with nonconstant column map.  This is an
actual submonoid containing the identity: its column slopes
multiply in $\{1,-1\}$, so a product remains outside $K$.
Its order is at most three.

Every row action of such a monoid has a common rational affine
fixed point except possibly when $P$ consists of the identity
adjoined to a two-element left-zero semigroup.  Here is a direct
verification.  If $P$ is a group, average a row orbit.  If its
minimal ideal has one element, that zero's row image is fixed
pointwise by $P$.  A two-element minimal ideal is either a
group of order two, a left-zero band, or a right-zero band.
For the group use its local group average; the only extra
element is the global identity.  For the right-zero band,
$\sigma_x\sigma_e=\sigma_e$ makes every point of
$\operatorname{im}\sigma_e$ fixed.  These cases also cover
orders one and two.

Only the two-left-zero case remains.  It uses all three
outside elements, and all their products stay outside $K$.
Consequently $A=\mathbb QP$, the unitalization of the algebra
of two left zeros.  If these are $p_0,p_1$, the orthogonal
idempotents $p_0,1-p_0$ and radical element $p_1-p_0$ identify
$A$ with the three-dimensional path algebra of one arrow.
It is hereditary, so every row module has projective dimension
at most one.  Lemma~\ref{C:row-pd-one-act-cups} applies in
this last case, and the affine-origin proposition applies in
all other cases.
\end{proof}

\subsection{Full column orbits and the three-by-three branch}

\begin{lemma}[Projective identifications of full column orbits]
\label{C:full-column-orbit-kernel}
Let $M$ have minimal band $I\times J$, and let $S$ be a finite
right act.  Suppose every core orbit in $SK$ has exactly $|J|$
elements.  If there are $o$ such orbits, the constant-action
quotient satisfies
\[
 0\longrightarrow(e\mathbb ZM)^{o-1}
 \longrightarrow\widetilde{\mathbb ZS}
 \longrightarrow\widetilde{\mathbb Z\bar S}
 \longrightarrow0,
\]
for a minimal idempotent $e$.  In particular the quotient
act-cone map is an integral homology isomorphism in every
degree at least three.
\end{lemma}
\begin{proof}
Label each full core orbit $O_a$ by its columns,
$s_{a,j}$, $j\in J$.  Each orbit is $M$-invariant and the
action on every orbit is the same column map:
$s_{a,j}m=s_{a,\tau_m(j)}$.
The relations of the constant-action quotient identify
$s_{a,j}$ with $s_{0,j}$ for fixed $j$, and make no other
identifications.  Its kernel therefore has integral basis
\[
 s_{a,j}-s_{0,j}\qquad(1\leq a<o,\ j\in J).
\]
For each $a$ these vectors form the right ideal module
$e\mathbb ZM$, proving the exact sequence and projectivity.
The act-cone Tor formula and its long exact sequence give the
claimed homology comparison, exactly as in
Lemma~\ref{C:two-column-projective-orbit-kernel}.
\end{proof}

\begin{theorem}[Three-by-three act cones below twenty-two arrows]
\label{C:three-by-three-act-cones21}
No act cone $\mathcal C(M,S)$ with minimal band $3\times3$
and $|M|+|S|+1\leq21$ realizes $\mathbb CP^2$.
\end{theorem}
\begin{proof}
Suppose otherwise and write $q=|M|-9$ for the number of
outside elements.  At most three states are excluded by the
earlier small-state results, so $q\leq7$ and $|S|\leq11-q$.
If $q\leq2$, Corollary~\ref{C:eleven-act-cones21} already
applies.  We may therefore assume $3\leq q\leq7$.

Let $c_j\in\bar S$ be the constant value of an old minimal
element in column $j$, and put $T=\{c_j:j\in J\}$.
If $|T|=1$, every old actor fixes this one state $s_0$.
For a balanced form $L$ on the compressed band, its general
affine representative is $u_L(f,g)=b^L_fa_g$, with
$b^L_m=L([s_0m]-[s_0],-)=0$ on every old actor.
This formula and its honest Karoubi transfer are the initial
construction in Theorem~\ref{C:act-common-fixed-row-cups};
that construction does not restrict the old column count.
The transferred two-cocycles thus
vanish on all pairs of old actors, making every degree-two
cup product zero on the act category.  This is impossible.

If $|T|=2$, identify old band columns with the same $c_j$,
while retaining all outside labels.  Equivariance
$c_jm=c_{\tau_m(j)}$ makes this a monoid quotient $M'$ with
minimal band $3\times2$ and the same $q$ outside elements.
The constant act $\bar S$ factors through $M'$.
The compression monoids for $(M,\bar S)$ and $(M',\bar S)$
are literally identical: their outside labels and products
agree, and their new band is $I\times\bar S$ with the same
actions and the same images of old band products.
Consequently the two act cones are homotopy equivalent through
that common monoid.  Since $|M'|=6+q\leq13$,
Theorem~\ref{C:three-row-fourteen-act-cones} makes every
rational degree-two product of this compressed cone zero,
without any finiteness assumption.  The degree-two quotient
surjection contradicts the original projective-plane square.

It remains that $|T|=3$.  Every core orbit has all three
columns: an equality between two column points in any core
orbit would identify the corresponding $c_j$ in $\bar S$.
Apply Lemma~\ref{C:full-column-orbit-kernel} and compress to
$N$.  Its rational homology is bounded and
$H_4(BN;\mathbb Q)=\mathbb Q$, because both are true of the
original cone and the homology comparison holds in degrees
at least three.  The contracted algebra still has dimension
$q$.  If $q\leq4$, the band Tor formula and
Theorem~\ref{B:contracted-four-Tor} give an immediate
contradiction.

Thus $5\leq q\leq7$ and $|S|\leq6$.  The act is connected:
the first homology sequence of the act-cone pushout identifies
$H_1$ with the reduced degree-zero homology of its action
space, since $H_1(BM;\mathbb Q)=0$.
Two full core orbits would already use six states, and each
is individually $M$-invariant.  At this state bound they would
exhaust $S$ and disconnect the act.  There is therefore only
one core orbit.  Every band actor is then constant on all of
$S$: its image lies in that orbit, and idempotence shows it
equals the constant value of its restriction to the orbit.
Compression is now an actual homotopy equivalence from the
original cone to a monoid with
\[
 |N|=q+3|S|\leq33-2q.
\]
For $q=6,7$ this is at most $21$, contrary to the exact
monoid minimum $23$.  For $q=5$, it is at most $23$;
the only case not already excluded by the monoid minimum is
$|S|=6$.  Its compressed monoid has minimal band $3\times6$
and five outside elements, excluded by
Lemma~\ref{C:three-by-six-five-outside-CP2}.
This exhausts the cases.
\end{proof}

\begin{proposition}[Four states and four minimal rows through twenty-one arrows]
\label{C:four-state-four-row-cones21}
Let $M$ have a $4\times2$ rectangular-band minimal ideal and
at most sixteen elements.  A right $M$-set $S$ with four states
does not give $B\mathcal C(M,S)\simeq\mathbb CP^2$.
\end{proposition}
\begin{proof}
Suppose the cone has that homotopy type.
Lemma~\ref{C:four-state-band-dichotomy} forces every minimal-band
element to act constantly.  Compression gives a monoid $N$ of
the same homotopy type with a $4\times4$ minimal band and
\[
 |N|=|M|-8+16=|M|+8\leq24.
\]
The two original columns give a nonempty subset of at most two
new columns invariant under every outside element.
Proposition~\ref{B:four-by-four-image24} forces the action image
of $N$ to have at least twenty-four elements.  Consequently
$|N|=24$ and its action map is faithful.

The complete allowed-image enumeration at order twenty-four in
the three hulls has the following counts.  The third column
uses $(b_2,b_3,b_4)=(1,0,1)$ over $\mathbb F_2$.
\[
 \begin{array}{c|r|r|r|r}
 \text{hull}&\text{rank-one images}&\text{orbits}&
 \text{additive candidates}&\text{small invariant subset}\\\hline
 217&121656&7659&74&53\\
 216&7934&3976&7&7\\
 221&22170&5582&32&20.
 \end{array}
\]
The last column allows a one- or two-element invariant subset on
either side; allowing both sides preserves coverage under opposite
isomorphisms.  Of these eighty cases, nine have an invariant
singleton.  Each already has zero degree-two products for the
following structural reason.  After opposition if needed, retain
the eight outside elements and the four band elements in that
single invariant column.  This is a twelve-element submonoid
$P$: products of outside elements landing in the band must
have that column, because a composition of maps fixing the
chosen point still fixes it.  The minimal band of $P$ has one
column, so $BP$ has zero positive rational cohomology by the
minimal-corner reduction.  Its action on the four original
columns reconstructs $N$ by constant-minimal-act compression.
Every degree-two class of that cone is relative to $BP$ and
has zero products, since the relative quotient is a suspension.

The remaining $71$ cases split as $48,3,20$ in the three hulls.
For every one, an exact integral projective prefix through $P_4$
and a literal comparison to the normalized bar give a primitive
degree-two cocycle and a square which is identically zero on
every degree-four projective basis vector.  The independent
verifier \path{code/a_verify_cp24_compression_cups.py} reconstructs
all projective actions, checks Smith exactness and saturation
through $P_3$, and checks every comparison-chain and equivariance
equation.  Thus its conclusion requires no inference about an
uncomputed tail.  All $71$ checks pass in
\path{code/outputs/a_cp24_compression_cups_independently_verified.json}.

Finally \path{code/a_cp24_small_subband_coverage.py} independently
recomputes every invariant subset among the $113$ additive
candidates, verifies the nine singleton submonoids directly,
and matches the remaining cases bijectively to those $71$
checked cup certificates.  The complete closure/orbit lists are
the files with prefix \path{code/outputs/d_cp24_hull}; the
independent linkage is
\path{code/outputs/a_cp24_small_subband_coverage.json}.
Every remaining square is zero, giving the contradiction.
\end{proof}

\begin{lemma}\label{C:four-by-five-five-outside-CP2}
A finite monoid with a $4\times5$ or $5\times4$ rectangular-band
minimal ideal and exactly five elements outside that ideal cannot
have classifying space homotopy equivalent to $\mathbb CP^2$.
\end{lemma}
\begin{proof}
After opposition the band--Tor modules have dimensions
$\dim V=3$, $\dim W=4$ over a five-dimensional contracted monoid
algebra $A$.  Their required Tor dimensions are $1,0,1,0,\ldots$.
A local $A$ is excluded by Proposition~\ref{B:local-quotient-CP2}.
The nonlocal classification of
Proposition~\ref{B:five-nonlocal-exceptions} leaves the two algebras
whose indecomposables and exact Tor table are given in
Lemma~\ref{B:five-exceptional-targets}.

For the zero-composite line, nonzero degree-two Tor of dimension
one requires exactly one $T_1$ summand in $W$ and exactly one
$S_3$ summand in $V$.  Degree-one vanishing forbids $T_2$ in $W$
and $S_2$ in $V$.  The only three-dimensional possibilities are
\[
 V=S_3\oplus S_1^{\oplus2},\quad
 S_3\oplus Ae_2,\quad S_3\oplus Ae_3,
\]
where the latter two projectives are the single-arrow modules.
Write
$W=T_1\oplus T_3^{\oplus s}\oplus(e_1A)^{\oplus p}
                         \oplus(e_2A)^{\oplus q}$.
Then $s+2p+2q=3$.  The respective tensor dimensions are
$2+s+2p$, $s+p+q$, and $2s+q$.
The first is at least two.  In the second, tensor dimension one
would give $s+p+q=1$, incompatible with the dimension equation.
In the third, tensor dimension one would give $s=0,q=1$, again
incompatible with $s+2p+2q=3$.  Thus the line gives no pair.

For the Nakayama algebra, the complete table leaves only
$V=S_2\oplus Ae_2$ as a three-dimensional left module admitting
degree-one vanishing and degree-two Tor of dimension one.
Indeed the other left decompositions either give zero degree-two
Tor on every degree-one-vanishing right summand, or give only
even or only multiples-of-three degree-two dimensions.
For this remaining $V$, the allowed right summands have Tor triples
\[
\begin{array}{c|rrrr}
 &T_2&R&e_1A&e_2A\\\hline
(\dim\operatorname{Tor}_0,\dim\operatorname{Tor}_1,
 \dim\operatorname{Tor}_2)&(2,0,1)&(1,0,1)&(1,0,0)&(2,0,0).
\end{array}
\]
Tensor dimension one excludes $T_2$.  The remaining degree-two
class requires one $R$ summand, already consuming the entire
tensor dimension and two of the four module dimensions.
The remaining two-dimensional projective summand would be
$e_2A$, which contributes two more tensor dimensions.
There is no solution.  These are all cases.

For a separate exact finite check,
\path{code/b_five_three_by_four_modules.py} enumerates every
dimension-three left and dimension-four right decomposition in
the two displayed indecomposable lists, using the accepted exact
Tor table.  The line has sixteen left and thirty right modules;
the Nakayama algebra has nine left and sixteen right modules.
Neither has a pair with Tor triple $(1,0,1)$.
Its output is \path{code/outputs/b_five_three_by_four_modules.json}.
\end{proof}

\subsection{Four and five rows with five outside elements}

\begin{lemma}[A missing old column character]
\label{C:missing-column-character-act-cups}
Let $M$ have a minimal ideal with two columns and row
augmentation module $V$.  If
$\operatorname{Hom}_A(V,\mathbb Q_\chi)=0$ for its old
column character $\chi$, all rational degree-two products
of every finite act cone are zero.
\end{lemma}
\begin{proof}
After the constant-action quotient and compression, each
balanced form $L$ restricts on the old column augmentation
line to a covector $\ell_L$ satisfying
$\ell_L A_m=\chi(m)\ell_L$.
The hypothesis makes it zero.  Consequently the honest
transferred affine cocycle has
$b^L_m=\beta(m)\ell_L=0$ on every old actor, and its value
on every pair of old actors is zero.  The second pair in
any nondegenerate four-chain of the act category consists
of old actors.  Thus every degree-two product is already
zero as a cochain on the quotient cone.  Pullback along
the degree-two surjection proves the result.
\end{proof}

\begin{theorem}[Four rows and five outside elements]
\label{C:four-row-five-outside-seven-states}
Let $M$ have minimal band $4\times2$ and exactly five outside
elements.  No finite right act with at most seven states
gives a cone of type $\mathbb CP^2$.
Equivalently, this entire minimal-band/outside-size branch
is excluded through twenty-one category arrows.
\end{theorem}
\begin{proof}
Assume such a cone exists.  A state fixed by the minimal band
would make its degree-two products zero.  Otherwise use
Lemmas~\ref{C:two-column-projective-orbit-kernel} and
\ref{C:two-column-finite-quotient-model}.
If $o$ is the number of two-point core orbits, the compressed
monoid $N$ has a finite four-dimensional model and band--Tor
modules with
\[
 \dim V=3,\quad \dim W=|S|-2o+1,\qquad
 \bigl(\dim\operatorname{Tor}_0,
       \dim\operatorname{Tor}_1,
       \dim\operatorname{Tor}_2\bigr)=(o,0,1),
\]
and zero higher Tor, over the same five-dimensional contracted
algebra $A$ as $M$.
If $A$ is local, Lemma~\ref{C:local-two-column-act-cups}
already gives a contradiction.  Otherwise
Proposition~\ref{B:five-nonlocal-exceptions} leaves the line
and Nakayama algebras of Lemma~\ref{B:five-exceptional-targets}.

For the Nakayama algebra, its exact indecomposable Tor table
shows that the only three-dimensional left module compatible
with $\operatorname{Tor}_1=0$ and one-dimensional
$\operatorname{Tor}_2$ is
\[
 V=S_2\oplus Ae_2.
\]
Indeed a degree-two contributor requires exactly one copy of
$S_2$ or $L$; any added $S_1$ forces degree-one Tor, and
$L\oplus S_2$ contributes twice.  The remaining projective
and semisimple cases contribute zero or an even multiplicity.
The only one-dimensional quotient character of the displayed
module is $S_2$, and it occurs twice.
If the old column character differs, use
Lemma~\ref{C:missing-column-character-act-cups}; if it agrees,
use Lemma~\ref{C:two-covector-affine-origin}.
Both contradict the projective-plane square.

For the line algebra write
$L_a=Ae_2$, $L_b=Ae_3$, $R_a=e_1A$, $R_b=e_2A$.
The degree-two condition requires one $S_3$ in $V$ and one
$T_1$ in $W$; the degree-one condition excludes $S_2$ from
$V$ and $T_2$ from $W$.
Thus the only possible row modules are
\[
 2S_1\oplus S_3,\qquad S_3\oplus L_a,
 \qquad S_3\oplus L_b.
\]
The first is split into characters and is excluded by
Proposition~\ref{C:split-row-act-cups}.
The last has only the quotient character $S_3$, with
multiplicity two, so the preceding affine argument excludes it.

For $V=S_3\oplus L_a$, write
\[
 W=T_1\oplus T_3^{\oplus s}
       \oplus R_a^{\oplus p}\oplus R_b^{\oplus q}.
\]
The exact tensor dimensions give
\[
 o=s+p+q,\qquad
 1+s+2p+2q=|S|-2o+1,
 \qquad |S|=4o-s.
\]
Since $s\leq o$ and $|S|\leq7$, one has $o\leq2$.
For $o=1$, $|S|\leq4$, already excluded by the small-state
theorems.  For $o=2$ the only cases are
\[
 \begin{array}{c|c|c}
 |S|&s&(p,q)\\ \hline
 6&2&(0,0)\\
 7&1&(1,0)\text{ or }(0,1).
 \end{array}
\]
For six states, both degree-two forms factor through the
single $S_3$ row summand, so they have matrix rank one.
For seven states with $(p,q)=(0,1)$, the contributions from
$T_3\otimes S_3$ and $R_b\otimes L_a$ give a basis of two
rank-one balanced forms.  Every basis class has square zero
by Lemma~\ref{B:rank-one-square}.
Their pullbacks span the one-dimensional $H^2$ of the original
cone.  At least one has a nonzero pullback, contradicting the
projective-plane square.  This uses a spanning set of zero
squares, not an unsupported claim about their mixed product.

For the remaining pair
\[
 V=S_3\oplus L_a,
 \qquad W=T_1\oplus T_3\oplus R_a,
\]
both representation kernels in $A$ are exactly $\mathbb Qb$.
The complete actual-basis result
Lemma~\ref{B:five-line-bases}, applied to the opposite algebra
if necessary, writes $b$ as a nonzero multiple of either one
distinguished basis element or a difference of two such elements.
In the first case both augmentation operators are zero, so
both finite set maps are constant and the actor profile lies
in the band.  In the difference case the two augmentation
operators agree on each side and are nonzero: otherwise two
independent basis elements would lie in the one-dimensional
kernel $\mathbb Qb$.  Equality of nonzero augmentation
operators determines equality of the actual finite set maps.
Thus the faithful profile quotient of $N$ kills $b$ modulo
the band.  Its contracted algebra is a quotient of
\[
 A/(b)\cong\mathbb Q\times\mathbb Q(1\longrightarrow2).
\]
Every quotient of this particular one-arrow algebra is
hereditary.  Its fourth rational homology is zero, whereas
the natural map preserves the balanced-tensor description
of $H^2$.  All degree-two products on $N$ pull back from
zero products in that quotient, a final contradiction.
\end{proof}

\begin{lemma}[The four-by-four module dimensions at five outside elements]
\label{C:five-by-five-five-outside-CP2}
A monoid with minimal band $5\times5$ and exactly five
outside elements cannot have type $\mathbb CP^2$.
\end{lemma}
\begin{proof}
The row and column augmentation modules both have dimension
four and require Tor dimensions $(1,0,1,0,\ldots)$ over a
five-dimensional algebra.  The local case is excluded by
Proposition~\ref{B:local-quotient-CP2}.  The nonlocal
classification leaves the same two algebras as above.

For the line algebra the degree-one and degree-two conditions
give exactly the following row and column possibilities.
Their degree-zero Tor dimensions are
\[
\begin{array}{c|ccc}
 &T_1\oplus3T_3&T_1\oplus T_3\oplus R_a&
                  T_1\oplus T_3\oplus R_b\\ \hline
 3S_1\oplus S_3&6&7&4\\
 S_1\oplus S_3\oplus L_a&4&4&3\\
 S_1\oplus S_3\oplus L_b&7&4&4
\end{array}
\]
None is one.
For the Nakayama algebra the only remaining left modules are
$S_2\oplus Ae_1$ and $L\oplus Ae_2$.
For the former, the admissible four-dimensional right modules
with one degree-two contribution are
$T_2\oplus e_1A$ and $R\oplus e_2A$;
for the latter, only $T_2\oplus e_1A$ remains.
Their tensor dimensions are all three, not one.

These lists can also be checked directly from the full
indecomposable table: \path{code/a_five_four_by_four_modules.py}
enumerates all thirty dimension-four modules on each side
for the line and all sixteen on each side for the Nakayama
algebra, checking all $900+256$ pairs.
The exact output is
\path{code/outputs/a_five_four_by_four_modules.json}; it has
no pair with Tor triple $(1,0,1)$.
\end{proof}

\begin{corollary}[Five rows and five outside elements]
\label{C:five-row-five-outside-five-states}
Let $M$ have minimal band $5\times2$ and five outside elements.
No act with at most five states gives a projective-plane cone.
Thus this entire branch is excluded through twenty-one arrows.
\end{corollary}
\begin{proof}
At most three states are already excluded below twenty-two
arrows.  For four states the band dichotomy gives zero products
or constant action.  In the latter case compression has band
$5\times4$ with five outside elements, excluded by
Lemma~\ref{C:four-by-five-five-outside-CP2}.
For five states, a fixed state again gives zero products.
Otherwise there are one or two two-point core orbits.
With one, the action is constant and compression has band
$5\times5$, excluded by the preceding lemma.
With two, the projective-kernel and finite-model lemmas give
a compressed monoid with band $5\times3$, degree-two rank
two and fourth homology $\mathbb Z$.
The rank-two part of
Theorem~\ref{B:three-by-five-targets}, after opposition,
makes every rational degree-two product zero on this finite
source.  Its degree-two classes pull back onto those of the
original cone, contradicting its nonzero square.
\end{proof}

\subsection{Small outside quotients and at least three columns}

\begin{lemma}[A small column image of the constant-action quotient]
\label{C:old-column-image-two-transfer}
Let $M$ have minimal band $I\times J$, and let $\bar S$ be
the constant-action quotient of a finite right act $S$.
Write $c_j\in\bar S$ for the constant value of an old band
element in column $j$, and put $T=\{c_j:j\in J\}$.
If $|T|=1$, every rational degree-two product of the original
act cone is zero.  If $|T|=2$, those products pull back from
an act cone over a monoid with minimal band $I\times T$ and
the same number of outside elements as $M$.
The latter comparison is an actual homotopy equivalence
after replacing $S$ by $\bar S$.
\end{lemma}
\begin{proof}
For $|T|=1$, the unique old column value $s_0$ is fixed by
every old actor.  Each transferred balanced cocycle has
\[
 u_L(f,g)=b^L_fa_g,\qquad
 b^L_m=L([s_0m]-[s_0],-)=0
\]
on old actors.  Its degree-two products therefore vanish
on all normalized act-category four-chains, as in
Theorem~\ref{C:three-by-three-act-cones21}.

For $|T|=2$, identify old columns with the same $c_j$ and
retain every outside element.  Equivariance of $j\mapsto c_j$
makes this a monoid quotient $M'$.  The act $\bar S$ factors
through it.  The compression monoids of
$(M,\bar S)$ and $(M',\bar S)$ are literally the same: both
have the old outside elements and the band $I\times\bar S$,
and every product and action is unchanged.
The compression theorem proves the asserted homotopy
comparison.  Finally
Lemma~\ref{C:constant-action-quotient-H2} is surjective on
degree-two cohomology in the pullback direction, which gives
the stated transfer to the original act cone.
\end{proof}

\begin{proposition}[Small outside quotients and small acts]
\label{C:three-outside-six-state-band-cones}
Let $M$ have rectangular-band minimal ideal $I\times J$ and
$q$ outside elements.  Its act cone cannot have type
$\mathbb CP^2$ in any of the following cases:
\[
 q=1;\qquad q\leq3\ \hbox{and }|S|\leq6;
 \qquad q\leq4\ \hbox{and }|S|\leq4.
\]
There is no restriction on the band dimensions.
\end{proposition}
\begin{proof}
For $q=1$, the contracted algebra is $\mathbb Q$, so the
row module is projective and
Lemma~\ref{C:row-pd-one-act-cups} applies at every state count.
Assume a projective-plane cone in one of the other cases.
It has no state fixed by the band, by the band-grid relative
argument.  Its act is connected, by the degree-one homology
sequence of the act-cone pushout.

If the band action is constant, compression gives an actual
projective-plane monoid with the same $q\leq4$ outside
elements.  The band Tor formula and
Theorem~\ref{B:contracted-four-Tor} exclude this.
If the action is not constant, there are at least two core
orbits.  Each has at least two states, each is individually
$M$-invariant, and their union cannot exhaust a connected
act.  Thus at most four states cannot occur at all.

For at most six states, at least one core orbit has exactly
two points: two orbits of three or more points would already
exhaust all six states.  The column values of that orbit
identify the old column image $T$ in the constant-action
quotient into at most two points.  Apply
Lemma~\ref{C:old-column-image-two-transfer}.
If $|T|=1$, the products are zero directly.  If $|T|=2$, the
comparison monoid has two columns and $q\leq3$ outside
elements, so
Lemma~\ref{C:three-outside-all-two-column-act-cups} makes
all its act-cone products zero without requiring bounded
homology for that intermediate cone.  Pullback gives the
same contradiction on the original projective plane.
\end{proof}

\begin{corollary}[At least three minimal columns below twenty-two arrows]
\label{C:three-or-more-column-act-cones21}
An act cone with rectangular-band minimal ideal having at
least three columns and at most twenty-one total arrows
cannot realize $\mathbb CP^2$.
\end{corollary}
\begin{proof}
At most two rows are excluded by the all-order two-row act
theorem, and at most three states are excluded below
twenty-two arrows.  Thus the band has at least three rows
and the act has at least four states.  The arrow budget and
the identity outside the band leave only
\[
 3\times3,\quad3\times4,\quad4\times3,
 \quad3\times5,\quad5\times3.
\]
The first shape is Theorem~\ref{C:three-by-three-act-cones21}.
For the last two, the budget forces $q=1$.
For the twelve-element bands, $1\leq q\leq4$ and
$|S|\leq8-q$.  The case $q=1$ is unrestricted; the cases
$q=2,3$ have at most six states; and $q=4$ has at most four.
The preceding proposition covers each case.
\end{proof}

\subsection{The nonconstant six-state four-row branch}

\begin{lemma}[Six bounded rank-two action hulls]
\label{C:six-rank-two-bounded-hulls}
Consider the six rank-two $4\times4$ action hulls numbered
\[
 137,169,172,188,222,225
\]
in Proposition~\ref{B:four-by-four-product25}.
Delete every actor on whose eventual cyclic group one of the
two stored integral degree-two classes restricts nontrivially.
Among closed remaining submonoids containing the band and identity,
of order at most twenty-two and rational degree-two rank two,
the following is a complete independently checked count.
\[
\begin{array}{c|r|r|r}
\text{hull}&\text{closed subsets}&\text{orbits}&
 (b_2,b_3,b_4)_{\mathbb F_2}=(2,0,\text{positive})\\\hline
137&332&168&4\\
169&570&157&0\\
172&708&367&0\\
188&332&168&4\\
222&7&4&0\\
225&7&4&0
\end{array}
\]
The third column uses automorphisms and anti-isomorphisms.
All eight representatives in the last column have order twenty.
In either orientation, if their outside actors preserve a two-element
subset of the four column labels, then the identity profile is the
only outside profile whose restriction to that subset is nonconstant.
\end{lemma}
\begin{proof}
The independent fixed-point closure program
\path{code/a_rank_two_act_closure.cpp} starts with the sixteen constant
profiles and the identity, adjoins one allowed actor at a time, and
repeatedly adjoins every product until closure.  It rejects a closure
only when it contains a forbidden actor or more than five optional
actors.  Thus every allowed closed subset in the stated size range is
reached.  Rational elimination on the nine balanced-tensor generators
retains exactly relation rank seven.  It uses primitive integer rows;
the original row squared norms are at most eight, so the nine-column
minor bound also bounds the intermediate integer arithmetic.

The separate checker \path{code/a_rank_two_act_audit.py} reconstructs
every literal actor product, every associativity equation, and every
balancing relation.  It independently verifies the cyclic-character
exclusions, the input allowed-actor bits, the complete closure sets,
and the disjoint orbit coverage under every supplied automorphism or
anti-isomorphism.  It then computes the three contracted two-sided
bar differentials over $\mathbb F_2$ in an independent tuple ordering,
checking all $827{,}910$ relevant $d^2=0$ equations.  The band double
shift identifies their homology with degrees two, three and four.
All $868$ representative results agree with the original producer.
The exact record is
\path{code/outputs/a_rank_two_act_audit.json}, whose completed flag is
true; the six individual closure logs are retained alongside it.

Finally the checker tests all six two-column subsets in both
orientations of every surviving representative.  With the original
zero-based case labels, the only possibilities are
\[
\begin{array}{c|c|c|c}
\text{hull}&\text{case}&\text{orientation}&\text{old columns}\\\hline
137&152,169&\text{ordinary}&\{2,3\}\\
188&168,204&\text{opposite}&\{0,2\}.
\end{array}
\]
The other cases, $279,291$ in hull $137$ and $208,210$ in hull $188$,
have no such subset in either orientation.  In each displayed case
the sole nonconstant outside profile on the subset is the identity,
whose row map is also the identity.  Testing both orientations is
necessary because the orbit reduction permits anti-isomorphisms.
\end{proof}

\begin{theorem}[Nonconstant six-state four-row acts]
\label{C:four-row-nonconstant-six-states}
Let $M$ have minimal band $4\times2$ and at most six outside elements.
If a finite right act has at most six states and some minimal-band
actor is nonconstant, its act cone does not have type
$\mathbb CP^2$.
In particular, the nonconstant case of the twenty-one-arrow profile
$(|M|,|S|)=(14,6)$ is excluded.
\end{theorem}
\begin{proof}
Suppose the cone $X$ is a projective plane.  A state fixed by the
minimal band would make all its degree-two products zero, so there
is no such state.  The act is connected.  Every core orbit therefore
has two points and is individually $M$-invariant.  Nonconstant band
action requires at least two core orbits, and they cannot exhaust
the connected act.  At most four states are impossible.

For five states there are two core orbits.  The projective-orbit
kernel and finite-quotient model lemmas give a compressed monoid
with minimal band $4\times3$, a finite four-dimensional model,
$H_2=\mathbb Z^2$, $H_4=\mathbb Z$, and no other positive homology.
The rank-two conclusion in the proof of
Theorem~\ref{B:three-by-four-targets}, after opposition, kills all
its rational degree-two products.  Its degree-two classes pull back
onto those of $X$, a contradiction.

For six states there are again exactly two core orbits: three would
exhaust the act and disconnect it.  The constant-action quotient
has four states.  Its compression $N$ has a $4\times4$ minimal band,
the same $q\leq6$ outside elements, and therefore order at most
twenty-two.  Lemmas~\ref{C:two-column-projective-orbit-kernel} and
\ref{C:two-column-finite-quotient-model} give
\[
 H_2(BN;\mathbb Z)=\mathbb Z^2,\qquad
 H_4(BN;\mathbb Z)=\mathbb Z,
\]
zero other positive homology, and an actual finite four-dimensional
CW model.  Moreover $H^2(BN;\mathbb Q)\to H^2(X;\mathbb Q)$ is
surjective.  Thus a square on $BN$ must have nonzero pullback.

Use the complete rank-two hull classification in the proof of
Proposition~\ref{B:four-by-four-product25}.  The rank-one-form,
repeated-point, zero-fourth-homology and exact-zero-product cases
have all rational degree-two products zero, and so cannot occur.
The three common-radical hulls numbered $161,237,400$ also have all
products zero.  The separately verified orbit-$346$ subset argument
gives zero fourth rational homology for every possible image of
order at most twenty-four.  Consequently only the six hulls in
Lemma~\ref{C:six-rank-two-bounded-hulls} remain.  We do not exclude
them merely from the existence of a radical line: the map to $X$
could kill that line.

Let $H$ be the faithful row-and-column image of $N$ in one of those
hulls.  This quotient fixes its band pointwise.  It preserves
integral $H_2$, by the identical balanced-tensor presentation, and
is onto on $H_3(-;\mathbb F_2)$ by
Lemma~\ref{B:band-quotient-H3-surjection}.
Thus $b_2(H;\mathbb F_2)=2$ and $b_3(H;\mathbb F_2)=0$.
If $b_4(H;\mathbb F_2)=0$, its rational fourth homology is zero,
and the isomorphism on rational $H^2$ again kills the source products.
The forbidden cyclic-character actors cannot occur: a local cyclic
group in the image lifts surjectively from an eventual local cyclic
group of $N$, and a nonzero integral character remains nonzero on
that lift.  Miller's theorem makes every such restriction zero for
the finite-dimensional source $BN$.
Lemma~\ref{C:six-rank-two-bounded-hulls} therefore applies to $H$.

Let $T$ be the image in the four quotient states of the original two
minimal columns.  If $|T|=1$, every old column term $b_m$ in the
transferred affine cocycles is zero, which already kills all products.
Otherwise $T$ is a two-element subset preserved by every outside
actor.  This invariance passes literally to the profile image $H$.
The lemma says that an image profile nonconstant on $T$ must be the
identity profile.  Any original outside actor mapping to that
profile has the identity row map; an actor mapping to a band profile
has constant columns.  Hence every actor of the original $M$ with
nonconstant old column map fixes every row.
The common-fixed-row cup primitive now makes all products on $X$
zero, the final contradiction.

This last argument explicitly permits duplicate outside profiles.
No conclusion about the higher homology or finite type of a proper
image is substituted for cup-product naturality.
\end{proof}

\subsection{Four rows and six outside elements}

\begin{lemma}[Necessary four-row source profiles]
\label{C:four-row-six-outside-source-census}
Suppose a monoid with minimal band $4\times2$ and six outside
elements has an act cone with a nonzero rational degree-two square.
Up to ordinary monoid isomorphism it belongs to one of two complete
source lists: seventeen monoids for row-level partition $3+1$, or
$6534$ monoids for row-level partition $2+2$.
\end{lemma}
\begin{proof}
If the old column character has no row quotient covector, all products
vanish by Lemma~\ref{C:missing-column-character-act-cups}.
If it has two independent quotient covectors, use
Lemma~\ref{C:two-covector-affine-origin}.
Thus its quotient-covector space is a line.  A nonzero covector gives
a function $f$ on the four rows, considered up to affine rescaling.
The nonconstant-column actors act on the finite image of $f$ by
translations or central reflections.  Translations are zero, and all
reflections have the same center.

If there are four values, all row fibers are singletons and their
uniform affine average is fixed.  For three values, either there is
no reflection and a singleton fiber gives a fixed row, or the values
are symmetric.  In the latter case a singleton central fiber is
fixed; if the central fiber is doubled, both extreme fibers are
singletons and their midpoint is fixed.  These are all possibilities
on four rows.  The common-affine-origin lemma excludes them.
Consequently $f$ has two values, which can be normalized to zero and
one, with multiplicities $3+1$ or $2+2$.

For each partition, enumerate all transformation pairs satisfying
the exact balancing equation
\[
 f(\sigma(i))-f(\sigma(0))
   =\bigl(\tau(1)-\tau(0)\bigr)\bigl(f(i)-f(0)\bigr).
\]
There are respectively $194$ and $96$ profiles.
The profiles with nonconstant columns form a monoid of respectively
$30$ and $32$ elements.  Enumerate every closed submonoid containing
the identity of order at most six, and discard it only if the
equations for a common rational affine row fixed point are consistent.
For each remaining profile image $P$, choose every multiset of the
other $6-|P|$ outside profiles from $P$ and the constant-column
profiles.  Duplicate identity profiles, duplicate nonidentity profiles,
and outside actors whose profiles coincide with band elements are
all retained.

For each multiset, try every product value with the required
transformation profile and impose all associativity equations.
The unique quotient-covector line makes the permitted row relabelings
exactly those preserving or interchanging its two level sets; include
both column permutations and all permutations within equal-profile
outside fibers.  This yields the stated ordinary-isomorphism lists.

The producer files are \path{code/a_four_row_profile_inputs.py} and
\path{code/a_four_row_sources.cpp}.  Independent verification uses
\path{code/b_verify_four_row_source_inputs.py} and
\path{code/b_four_row14_raw_audit.cpp}.
It reconstructs all profile inputs, tests affine consistency with
\texttt{Fraction} arithmetic, regenerates the $528$ and $119948$
profile multisets, and uses a different product-assignment order with
no propagation or producer canonicalizer.  Full literal relabeling
expansions of the seventeen and $6534$ retained sources are disjoint
and equal the independently generated sets of twenty and $16884$
tables.  All $46{,}384{,}576$ associativity equations are checked.
The complete independent record is
\path{code/outputs/b_four_row14_sources_independently_verified.json}.
\end{proof}

\begin{proposition}[The six-dimensional row-module reduction]
\label{C:four-row-six-outside-row-reduction}
For the source lists in the preceding lemma, put
$A=\mathbb QM/\mathbb QK$ and $V=\widetilde{\mathbb QI}$.
The following exact partition holds.
\[
\begin{array}{c|r|r|r|r}
\text{row levels}&\operatorname{pd}_A V\leq1&
 \substack{\text{stable Tor}\\\text{recurrence}}&\substack{\text{local-corner}\\\text{obstruction}}&
 \substack{\text{remaining}\\\operatorname{pd}_A V\leq2}\\\hline
3+1&10&7&0&0\\
2+2&4374&1977&10&173
\end{array}
\]
In every entry of the middle two columns, bounded
$\operatorname{Tor}_*^A(W,V)$ forces
$\operatorname{Tor}_2^A(W,V)=0$ for every finite-dimensional right
$A$-module $W$.
\end{proposition}
\begin{proof}
The full actual contracted multiplication and row action are retained
for every source.  The radical is recovered from the regular trace
form and checked to be a nilpotent two-sided ideal.  When the entire
semisimple quotient is split basic, exact lifted orthogonal
idempotents give projective covers.  Every projective term is stored
as a direct sum of literal ideals $A\varepsilon$, with a full basis
inside $A$.  The cover maps, all action equations, and the complete
embedded kernels are checked over $\mathbb Q$.

In the $2+2$ list this directly gives $2268$ projective rows and
$2092$ rows with projective first kernel.  The $3+1$ list gives ten
projective rows.  For $1975$ further $2+2$ rows and seven $3+1$ rows,
put $K=\Omega V$.  The certificates supply the literal identity
\[
 1_K=\sum_i g_i f_i+\sum_j u_jv_j,
\]
where each $g_if_i$ factors through a specified
$\Omega^sK$, with $s=1$ or $2$, and each $u_jv_j$ factors through a
finite free $A$-module.  Every displayed factor is an actual module
map.  Applying any positive Tor functor makes the projective terms
zero.  Therefore nonzero $\operatorname{Tor}_1(W,K)$ forces nonzero
Tor in arbitrarily high degrees by dimension shifting.  This is the
claimed obstruction to nonzero $\operatorname{Tor}_2(W,V)$ in a
bounded complex.

Sixteen source quotients are not split basic.  Fourteen have a
projective radical, certified by an identity factoring through free
modules.  Then $A/\operatorname{rad}A$ has projective dimension at
most one; each simple module is its direct summand, and finite
composition series give $\operatorname{pd}_A V\leq1$.
In the other two cases, an actual central cyclic group of order three
gives the central idempotent
$(1+g+g^2)/3$.  It acts as the identity on $V$; its complementary
cyclotomic summand acts as zero.  Literal projective covers in the
four-dimensional central summand give the same stable identity as
above.  These account for the additional fourteen entries in the
first column and two in the second.

For the ten local-corner entries there is an idempotent $e$ with
\[
 Ae=eAe=B,\qquad \dim B=4,
 \qquad K\cong\operatorname{rad}B
\]
as left $A$-modules, with the right side induced through $Ae$.
The certificate gives the full bases and an invertible intertwiner;
$B$ has radical dimension three and radical powers of dimensions
$3,1,0$, so it is local with residue field $\mathbb Q$.
Inducing a $B$-projective resolution through $Ae$, which is free of
rank one as a right $B$-module and projective on the left, gives
\[
 \operatorname{Tor}_n^A(W,V)
   \cong\operatorname{Tor}_n^B(We,\mathbb Q),\qquad n\geq2.
\]
In a minimal right free resolution of $We$, tensoring with the residue
field makes every differential zero.  Eventual Tor vanishing therefore
makes that resolution finite.  Thus $We$ has finite projective dimension
over the finite-dimensional local algebra $B$, hence is projective.
For completeness, a nonzero last differential in a minimal finite
free resolution would have all entries in the radical; multiplying
a basis vector by a nonzero element annihilated on the left by that
radical gives a kernel element, contradicting injectivity.
Thus the degree-two Tor group is zero as claimed.

Only the $173$ final rows remain.  For each, the stored exact
projective resolution has length two, proving the required bound
$\operatorname{pd}_A V\leq2$.  No strict lower bound on the projective
dimension is needed for the subsequent act calculation.
The producer records are the two pairs
\path{code/outputs/a_four_row_covers14_*.json.gz} and
\path{code/outputs/a_four_row_stable14_*.json.gz}, together with
\path{code/outputs/a_four_row_local_corner_tails.json},
\path{code/outputs/a_four_row_nonsplit_radicals.json}, and
\path{code/outputs/a_four_row_cyclotomic_tails.json}.
The independent Fraction-only checker
\path{code/d_verify_four_row_certificates.py} imports no producer
or Flint routines.  It reconstructs all $6551$ actual sources,
checks $17{,}975{,}944$ monoid associativity equations,
$11{,}066$ exact idempotent-projective presentations,
$132{,}792$ cover/kernel equivariance matrix equations, and
$1998$ identity factorizations.  It also checks all local-corner and
central-cyclotomic identifications literally.  The complete record is
\path{code/outputs/d_four_row_certificates_verified.json}.
\end{proof}

\begin{theorem}[Four rows and six outside elements]
\label{C:four-row-six-outside-six-state-cones}
Let $M$ have minimal band $4\times2$ and exactly six outside elements.
No finite right act with at most six states gives a cone of type
$\mathbb CP^2$.
\end{theorem}
\begin{proof}
The fixed-state and nonconstant-band cases are excluded by the
relative grid argument and
Theorem~\ref{C:four-row-nonconstant-six-states}.
At most four states are already excluded by the earlier small-state
theorems.  We may therefore suppose that the band acts constantly,
with five or six states.  If the two old column values coincide,
all old affine column terms vanish and the products are zero.
Otherwise label those two states zero and one.
Constant-action compression gives an actual homotopy-equivalent
monoid with the same contracted algebra $A$ and row module $V$.
The preceding two statements therefore leave exactly the $173$
sources with the remaining length-two projective resolutions.

For each source, enumerate the right acts in the indicated normal
form.  The identity action and all eight band actions are fixed.
Generator images range over every transformation with the required
two core-state images, subject only to their actual power relations.
Multiplication closure determines further actor images; the final
table is tested on every action equation.  Canonical augmentation
uses the exact stabilizer among permutations of the remaining states.
Already determined generators require no extra stabilizer change,
since their images are words in the preceding images.
Thus the complete counts are
\[
\begin{array}{c|r|r}
\text{states}&\text{constant acts}&
 (\operatorname{Tor}_0,\operatorname{Tor}_1,
 \operatorname{Tor}_2)_{\mathbb Q}=(1,0,1)\\\hline
5&110418&0\\
6&1060108&14544.
\end{array}
\]

The Tor calculation tensors the exact three-term projective
resolution.  For each summand $A\varepsilon$ it uses
$W\otimes_AA\varepsilon=W\varepsilon$; a differential is
reconstructed from its value on the literal idempotent generator.
All coordinate inclusions and both composite differentials are
checked exactly over $\mathbb Q$.
Hence a projective-plane cone must lie among the $14544$ six-state
survivors.  The complete literal twenty-one-arrow categories for
these survivors have $H^1(-;\mathbb F_2)=0$,
$H^2(-;\mathbb F_2)=\mathbb F_2$, and square zero.
For each, the normalized category bar contains an explicit
degree-three primitive for the square of a nonboundary degree-two
cocycle.  In total $607{,}110{,}192$ primitive equations are checked.
This contradicts the projective-plane square.

The source, act, tensor and cup routines are respectively
\path{code/a_four_row_sources.cpp},
\path{code/a_constant_two_column_acts.cpp},
\path{code/a_four_row_constant_tensor.py}, and
\path{code/a_four_row_constant_cups.cpp}.
Their terminal-count records are
\path{code/outputs/a_four_row_constant_acts14_s5.out},
\path{code/outputs/a_four_row_constant_acts14_s6.out},
the two \path{code/outputs/a_four_row_constant_tensor14_s*.json}
files, and \path{code/outputs/a_four_row_constant_cups_all.out}.
The complete literal primitives and source linkage are retained in
\path{code/outputs/a_four_row_constant_cups_all.jsonl} and its
\path{code/outputs/a_four_row_constant_cups_all_index.json} index.
The independent coordinate-free tensor checker
\path{code/b_prepare_four_row_mod_tensor.py} first verifies all
$173$ projective resolutions over $\mathbb F_{65521}$, including
all literal actions, idempotents, chain equations and exact ranks.
Then \path{code/b_verify_four_row_mod_tensor.cpp} computes the images
$W\varepsilon$ and the right-multiplication differential blocks in
ambient coordinates for every emitted act.  It independently checks
$1{,}354{,}896{,}648$ action equations and reproduces the candidate
counts $0$ and $14544$.  This is an exact finite-field necessary test
for a projective-plane cone, not an unproved inference from modular
to rational rank.
A separate checker reconstructs every survivor's literal category
from its source action, verifies the full index and table linkage,
and checks all cocycles and cup primitives without producer imports.
The combined independent record is
\path{code/outputs/b_four_row_constant_acts_independently_verified.json}.
The act-enumeration completeness uses the audited generator-image
and stabilizer proof above; no second numerical act enumerator is
asserted by these verification records.
\end{proof}

\begin{corollary}[Reduction to the final act-cone profile]
\label{C:act-cone-twenty-one-frontier}
Every finite act-cone category of type $\mathbb CP^2$ has at least
twenty-one arrows.  If it has at most twenty-one arrows, its minimal
ideal is a rectangular band, the band acts constantly, and its
parameters must be
\[
 (\text{rows},\text{columns},\text{outside elements},\text{states})
   =(4,2,7,5).
\]
Its unique old-column quotient covector has row-level partition $2+2$.
This final profile is excluded below in
Theorem~\ref{C:four-row-seven-outside-five-state-cones}.
\end{corollary}
\begin{proof}
Combine the preceding theorem with the all-order two-row theorem,
the nontrivial-minimal-group exclusion, the three-row forced-profile
theorem, the at-least-three-column theorem, the universal two-column
small-outside obstructions, and the four/five-row five-outside
theorems.  The earlier exact budget reduction leaves only
$(4,2,6,5)$, $(4,2,6,6)$, $(4,2,7,5)$, and $(5,2,6,4)$.
The first two have now been excluded, and
Theorem~\ref{C:five-row-six-outside-four-state-cones} excludes the last.
Nonconstant action in the
five-state case compresses to a finite rank-two $4\times3$ source
whose products vanish.  Thus the sole remaining profile has constant
band action and exactly twenty-one total arrows.
Corollary~\ref{C:four-row-seven-outside-split-act-cones} excludes
the row-level partition $3+1$, leaving only $2+2$.
\end{proof}

\subsection{The three-plus-one row family with seven outside elements}

\begin{proposition}[The complete three-plus-one source family]
\label{C:four-row-seven-outside-split-source}
Let $M$ have minimal band $K=I\times J$ with $|I|=4$, $|J|=2$,
and $|M\setminus K|=7$.  Let $\chi$ be its old column character.
Suppose the space of $\chi$-quotient covectors on
$V=\widetilde{\mathbb QI}$ is one-dimensional, its nonzero covector
has row-level multiplicities $3+1$, and the actors with $\chi\ne0$
have no common rational affine row fixed point.
There are exactly $414$ ordinary-isomorphism classes of such sources.
For every finite-dimensional right $A=\mathbb QM/\mathbb QK$-module
$W$, bounded $\operatorname{Tor}_*^A(W,V)$ implies
$\operatorname{Tor}_2^A(W,V)=0$.
\end{proposition}
\begin{proof}
The profile construction of
Lemma~\ref{C:four-row-six-outside-source-census} applies with seven
outside elements.  Normalize the row levels to $(0,0,0,1)$.
The same exact affine balancing equations give $194$ possible
transformation profiles, of which $30$ have nonconstant columns.
An independent closure tree gives $567$ closed nonconstant-column
profile monoids of order at most seven containing the identity.
Exactly $70$ have no common rational affine row fixed point.
For every such image $P$, all multisets of $7-|P|$ further outside
profiles are included.  In particular duplicate identity, band,
and nonidentity profiles are retained.
There are $46711$ multisets, $3989$ after the permitted profile
relabelings, and $127$ with closed profile support.  Trying every
compatible product and all associativity equations gives $669$
raw tables and $414$ ordinary-isomorphism classes.

The source producers are \path{code/a_four_row_profile_inputs.py}
and \path{code/a_four_row_sources.cpp} at order fifteen.
The independent files
\path{code/b_verify_four_row_source_inputs.py},
\path{code/b_four_row_source_raw_audit.cpp}, and
\path{code/b_verify_four_row15_split31.py}
use a separate closure tree, a different product-assignment order,
and no producer canonicalizer.  They regenerate all $669$ raw
tables, check $2257875$ associativity equations, and verify that the
full literal relabeling expansions of the $414$ retained classes
are disjoint and cover precisely those tables.  The terminal record is
\path{code/outputs/b_four_row15_split31_sources_independently_verified.json}.

For each source retain its actual contracted multiplication table
and its actual three-dimensional row module.  Identical pairs
$(A,V)$ may share a certificate, but equality is checked entry by
entry; neither ranks nor module dimensions are used as a substitute.
The $414$ sources have $363$ distinct such pairs.  Exact projective
presentations give the following source partition:
\[
\begin{array}{c|r}
\text{verified algebraic alternative}&\text{sources}\\\hline
V\text{ projective}&197\\
\operatorname{pd}_A V\leq1&63\\
\text{stable positive-Tor recurrence}&153\\
\operatorname{rad}A\text{ projective}&1.
\end{array}
\]
Every projective is a literal idempotent left ideal, and every cover
has its full embedded kernel.  In the recurrence cases,
$L=\Omega V$ has an exact identity factoring through a finite sum
of $\Omega^sL$, $s=1$ or $2$, and free modules.  Thus a nonzero
$\operatorname{Tor}_1(W,L)$ forces nonzero Tor in arbitrarily high
degrees.  This is exactly the dimension-shifting argument in
Proposition~\ref{C:four-row-six-outside-row-reduction}.

The remaining algebra has dimension seven, a two-dimensional
square-zero radical, and a five-dimensional semisimple quotient.
Its radical identity factors through finite free modules, so the
radical is projective.  Consequently $A/\operatorname{rad}A$ has
projective dimension at most one.  Each simple is a direct summand
of this semisimple regular module, and finite composition series
give projective dimension at most one for every finite-dimensional
module.  This also supplies the desired degree-two Tor vanishing.

The producer certificates are
\path{code/outputs/a_four_row_q7_modules_31.jsonl.gz} and
\path{code/outputs/a_q7_split31_radical.json}, with the complete
source-to-certificate index
\path{code/outputs/a_four_row_q7_modules_31.json}.
The checker \path{code/a_verify_q7_split31_certificates.py} uses
only the independent Fraction matrix routines of
\path{code/d_verify_four_row_certificates.py}; it imports no
producer, Flint, cover finder, or algebra classifier.
It checks all $414$ literal sources, $1397250$ monoid associativity
equations, $552$ exact projective presentations, $7728$ cover and
kernel equivariance equations, and $127$ exact identity
factorizations.  The radical and its semisimple quotient are
reconstructed independently.  Its complete record is
\path{code/outputs/a_q7_split31_certificates_verified.json}.
\end{proof}

\begin{corollary}[Three-plus-one acts]
\label{C:four-row-seven-outside-split-act-cones}
For any source in Proposition~\ref{C:four-row-seven-outside-split-source}
and any finite right act $S$, bounded rational homology of the act
cone implies that all products of rational degree-two classes vanish.
In particular no such cone has type $\mathbb CP^2$.
Among the remaining profile $(4,2,7,5)$ in
Corollary~\ref{C:act-cone-twenty-one-frontier}, only the row-level
partition $2+2$ can therefore occur.
\end{corollary}
\begin{proof}
A state fixed by the minimal band makes all degree-two cone classes
relative to the monoid vertex; the relative suspension product is
zero.  Otherwise every minimal-band orbit has two points, and the
projective orbit-kernel sequence of
Lemma~\ref{C:two-column-projective-orbit-kernel} preserves rational homology
in every degree at least three when the act is replaced by its
constant-action quotient.  Constant-action compression is an actual
homotopy equivalence to a monoid with the same $A$ and $V$.
Its degree-four homology is $\operatorname{Tor}_2^A(W,V)$ for the
new column augmentation module $W$, and its higher homology is the
corresponding shifted Tor.  The proposition makes this degree-four
group zero under the boundedness hypothesis.  The orbit-kernel
comparison therefore makes the original cone's rational
degree-four cohomology zero as well.  This proves the product claim.
The source-profile necessity and the affine-origin exclusions used
in Lemma~\ref{C:four-row-six-outside-source-census} do not depend
on having six rather than seven outside elements, so the final
statement follows.
\end{proof}

% General comparison lemma; numerical q7 conclusion is a separate gate.
\subsection{Small projective certificates for act-cone squares}

\begin{lemma}[The literal act-cone kernel]
\label{C:act-cone-literal-bar-kernel}
Let $M$ be a monoid, $S$ a nonempty right $M$-set, $k$ a field,
$R=kM$, and $L=\ker(kS\to k)$.
Write $C(M,S)$ for the two-object act category.
The kernel of the normalized-chain map
\[
 C_*(BC(M,S);k)\longrightarrow C_*(BC(M,*);k)
\]
is naturally the suspension of the normalized two-sided bar
$B_*(L,R,k)$.  In degree $n+1$ the comparison sends
\[
 w[a_1|\cdots|a_n]\longmapsto
 (-1)^{n(n-1)/2}\sum_{s\in S}w_s
       [a_n|\cdots|a_1|s],\qquad w=\sum_s w_s s.
\]
Consequently it induces
$\operatorname{Tor}_n^R(L,k)\cong\widetilde H_{n+1}(BC(M,S);k)$.
\end{lemma}
\begin{proof}
The two categories have the same objects and endomorphisms.
A nondegenerate simplex which changes objects consists of monoid
arrows followed by exactly one arrow labelled by a state.
The kernel on its state coefficients is precisely $L$.
All other normalized chains map identically.  This identifies the
kernel degreewise with the displayed two-sided bar.

Let $T_n$ denote the displayed map without its sign.
The last category face cancels because $\sum_s w_s=0$.
The face next to the state is its right $M$-action; the remaining
faces multiply consecutive monoid letters or remove the other end.
Reversing the letters gives
$\partial T_n=(-1)^nT_{n-1}d$.
The sign in the statement changes this to the usual suspended
differential $-d$.  The chain map is therefore a literal isomorphism
onto the kernel.  The quotient category $C(M,*)$ has a terminal
object and contractible classifying space.  The exact sequence of
normalized chains gives the asserted homology isomorphism.
\end{proof}

\begin{proposition}[A projective comparison for the affine square]
\label{C:act-projective-affine-square}
Suppose $M$ has minimal band $I\times J$ acting constantly on $S$.
Choose base elements $i_0\in I$ and $s_0\in S$ as in the
constant-action compression, and put
$V=\widetilde{kI}$ and $W=\widetilde{kS}=L$.
Let $F:W\times V\to k$ be a balanced bilinear form and let $u_F$
be its affine degree-two class on $C(M,S)$.
Set
\[
 \alpha_a=e_{\sigma_a(i_0)}-e_{i_0},\qquad
 \beta_b=e_{s_0b}-e_{s_0}.
\]
On the two-sided bar of Lemma~\ref{C:act-cone-literal-bar-kernel},
the pullback of $u_F^2$ in degree three is, up to the one common
minus sign supplied by that lemma,
\[
 w[a|b|c]\longmapsto F(w,\alpha_a)F(\beta_b,\alpha_c).
\]

Let $P_\bullet\to k$ be a projective resolution of the trivial
left $R$-module, with each displayed summand a literal ideal
$R\varepsilon$, $\varepsilon^2=\varepsilon$.
Projective terms $P_0,\ldots,P_3$ with exactness through $P_2$
(so $P_3$ surjects onto $\ker d_2$), together with their comparison
to the normalized free left bar resolution, suffice to certify that
this square is a coboundary by exhibiting a primitive on
$L\otimes_RP_\bullet$.
No termination of the undisplayed resolution is required.
\end{proposition}
\begin{proof}
The affine representative on an old pair $g$ then $f$ is
$F(\beta_f,\alpha_g)$; on an old arrow $g$ followed by the state
$s$ it is $F(e_s-e_{s_0},\alpha_g)$.
These are the normalized cocycles supplied by the actual
constant-action compression of
Theorem~\ref{C:constant-minimal-act-compression}.
The two Alexander--Whitney halves of
$[c|b|a|s]$ give
\[
 F(\beta_b,\alpha_c)F(e_s-e_{s_0},\alpha_a).
\]
Summing with coefficients $w_s$ gives the formula because their
sum is zero.  In degree three the sign of the preceding lemma is
$-1$; multiplying a cochain and its primitive by $-1$ does not
change whether the class is zero.

Here is a completely explicit comparison on a projective prefix.
The free left bar resolution has basis $(r;a_1,\ldots,a_n)$,
where each $a_i\ne1$, and its underlying augmented complex has
contraction
\[
 s(r;a_1,\ldots,a_n)=
 \begin{cases}(1;r,a_1,\ldots,a_n),&r\ne1,\\0,&r=1.
 \end{cases}
\]
The degree-minus-one contraction sends $1$ to $(1;\varnothing)$.
For each idempotent generator define recursively
\[
 \phi_n(\varepsilon)
   =\varepsilon\,s\bigl(\phi_{n-1}(d\varepsilon)\bigr),
\]
using the augmentation in degree zero.
The argument of $s$ is a cycle.  It is fixed by $\varepsilon$,
since $d\varepsilon$ is fixed by that idempotent and the previous
comparison is $R$-linear.  Thus
$d\phi_n(\varepsilon)=\phi_{n-1}(d\varepsilon)$.
The identity
$\varepsilon\phi_n(\varepsilon)=\phi_n(\varepsilon)$ makes the
extension $r\varepsilon\mapsto r\phi_n(\varepsilon)$ well defined
and $R$-linear.  This proves all comparison and module equations
on a generating set, and hence everywhere.

After tensoring, a term $(r;a,b,c)$ in
$\phi_3(\varepsilon)$ contributes
\[
 F(wr,\alpha_a)F(\beta_b,\alpha_c).
\]
This gives the square cochain on the small tensor complex.
If its row vector is the product of a row vector in degree two
with the tensor differential $d_3$, that vector is an explicit
degree-three primitive for the square on the act cone.
The projective prefix extends to a full projective resolution by
first covering $\ker d_3$ and then continuing to cover kernels;
the same contraction extends
$\phi$.  The comparison of these resolutions computes the derived
tensor product in every degree.  Therefore the certified primitive
has the asserted meaning in ordinary cohomology, independently of
whether the remaining resolution is finite.

For clarity, tensoring a map
$R\varepsilon\to R\eta$ whose generator has value $z$ uses
$w\mapsto wz$.  Its literal support is
$z=\varepsilon z=z\eta$.
On column-coordinate matrices for a right module this says
$R_z=E_\eta R_z E_\varepsilon$.
The genuine tensor terms are the images $W\varepsilon$.
Extending the matrix to ambient copies of $W$ preserves its rank,
and a supported cochain is a coboundary exactly when it belongs
to the ambient differential's row space.  This is the finite
linear-algebra test used by the certificates.
\end{proof}

\begin{remark}[What a finite-field square certificate proves]
\label{C:act-projective-square-certificate-scope}
In the application over $\mathbb F_{65521}$, the exact contracted
projective prefix gives
$\dim\operatorname{Tor}_0^{kM/kK}(W,V)=1$.
A nonzero balanced form is therefore the nonzero degree-two
cohomology class, by the affine representative and compression
comparison.  A primitive from the preceding proposition proves
that this actual field class has square zero and excludes
$\mathbb CP^2$.  It does not infer rational rank from a modular
rank, nor infer vanishing of higher homology from a finite prefix.
The routines are
\path{code/a_monoid_bar_prefix_modp.py} and
\path{code/a_act_bar_square_modp.py}; all projective ideals, full
embedded kernels, idempotent supports and sparse comparison chain
equations are retained explicitly.
\end{remark}

\subsection{Five rows, six outside elements, and four states}

\begin{lemma}[The necessary five-row profile lists]
\label{C:five-row-six-outside-source-census}
Suppose a monoid with minimal band $5\times2$ and six outside
elements has a finite act cone with a nonzero rational degree-two
square. Up to ordinary monoid isomorphism it belongs to the following
three necessary source lists:
\[
\begin{array}{c|r|r}
\text{row levels}&\substack{\text{nonconstant image}\\\text{orbits}}&
 \text{source monoids}\\\hline
4+1&14&33\\
3+2&183&60223\\
(-1,-1,0,0,1)&2&12.
\end{array}
\]
\end{lemma}
\begin{proof}
The missing-column-character lemma and the two-covector affine-origin
lemma leave a unique nonzero old-column quotient covector on the row
augmentation module. Let $f$ be its affine function on the five rows.
Every nonconstant-column actor induces a translation or a central
reflection on the finite set of values of $f$. Translations are zero,
and every reflection has the same center.

If there is no reflection and $f$ has at least three values, a
singleton value fiber is fixed by every such actor. If there are
reflections and four or five values, at least one opposite pair of
fibers consists of two singletons; its midpoint is a common affine
fixed point. For three values, a singleton central fiber gives a
fixed row. A central fiber of size three leaves two singleton extreme
fibers, whose midpoint is fixed. The only remaining three-value
case has fiber sizes $2,2,1$, with the doubled middle fiber; after
affine rescaling and reversing the values it is
$(-1,-1,0,0,1)$. With two values the multiplicities are $4+1$ or
$3+2$. The common-affine-origin primitive excludes all discarded
cases.

For each retained level vector enumerate every row transformation
$\sigma$ and two-column transformation $\tau$ satisfying
\[
 f(\sigma(i))-f(\sigma(0))
 =\bigl(\tau(1)-\tau(0)\bigr)\bigl(f(i)-f(0)\bigr).
\]
The full profile universes have respectively $2310,730,154$ elements;
the nonconstant-column submonoids have $260,180,24$ elements.
Enumerate their submonoids containing the identity through order six,
up to the actual level-preserving row-conjugacy group, and solve the
common affine fixed-point equations over $\mathbb Q$.
This leaves the image-orbit counts in the statement.

There is a shorter exact enumeration for $4+1$. Write $*$ for the
singleton upper row. Every reflection is
\[
 r_i:\quad\{0,1,2,3\}\longmapsto *,\qquad *\longmapsto i.
\]
With no reflection, $*$ is fixed. With only $r_i$, every
identity-column actor fixes $i$ because its product with $r_i$ is
again $r_i$; hence the midpoint of $i$ and $*$ is fixed. Thus two
different reflections $r_i,r_j$ are necessary. Their products are
the two distinct maps which fix $*$ and collapse the four lower
rows to $i$ or $j$. Together with the identity these are five
distinct profiles. Three different reflections would require seven
profiles, which is impossible here. Normalize $i=0,j=1$ and adjoin
at most one further nonconstant profile. This yields exactly the
fourteen image orbits used above.

For each remaining nonconstant image $P$, select one outside actor
for every nonidentity element of $P$, then enumerate every multiset
of the remaining $6-|P|$ profiles. These may be constant-column
profiles or further copies of profiles in $P$. Duplicate identity
profiles, duplicate nonidentity profiles, and outside actors with
the same profile as a band element are all retained.
Since $P$ has already been normalized up to conjugacy, multiset
pruning uses only its \emph{actual stabilizer}. This prevents an
independent normalization from discarding a relative position.

For each multiset, every product value with the required profile
is tried. Associativity propagation is followed by all $16^3$
literal associativity equations at each leaf. The row quotient
covector must be unique; equivalently its relation matrix has rank
three. The implementation tests this rank modulo $32003$, which
is exact here: the given covector makes the rank at most three,
and every $3\times3$ minor has absolute value at most
$12^{3/2}<42$.
Finally quotient by all level-preserving row permutations, both
column permutations, and every permutation within equal-profile
outside fibers. The unique covector line ensures that these are
exactly the ordinary source isomorphisms in this normal form.

The profile-multiset counts are $2068,29279353,136$.
After the necessary profile tests, the numbers of associative
labelled source tables are $36,94744,15$, giving the ordinary
source counts in the statement. The producer files are
\path{code/d_five_row_profile_universes.py},
\path{code/d_five_row_profile_submonoids.cpp},
\path{code/d_five_row_source_inputs.py}, and
\path{code/d_five_row_sources.cpp}.
The complete tables and terminal counts are retained as
\path{code/outputs/d_five_row_sources16_41.tsv},
\path{code/outputs/d_five_row_sources16_32.tsv}, and
\path{code/outputs/d_five_row_sources16_221.tsv}, with their
corresponding output logs.
The independent source gate reconstructs the profile algebras and
actual relabelings.  In the $3+2$ case,
\path{code/d_verify_five_row32_nonconstant.py} generates all $10159$
labelled closed nonconstant-profile monoids without orbit pruning,
then recovers exactly the $1075$ final image orbits.  The separate
multiset enumeration reproduces all $29279353$ multisets and $8734$
closed profile lists.  A different associative-extension search assigns
products in reverse column order, without the producer's propagation.
Its complete raw output agrees exactly with the disjoint literal
isomorphism-orbit expansion of all retained source tables.
Together with the independent $4+1$ and three-value checks, this gives
$60268$ ordinary classes and $94795$ raw associative tables, checking
all $388280320$ associativity equations.  The source algorithms are
\path{code/b_verify_five_row_source_inputs.py},
\path{code/b_rectangular_source_raw_audit.cpp}, and the separate
$3+2$ multiset and expansion checkers.  Their exact complete linkage is
\path{code/outputs/b_five_row_sources_acts_complete_verified.json}.
\end{proof}

\begin{proposition}[The five-row algebra partition]
\label{C:five-row-six-outside-row-reduction}
Put $A=\mathbb QM/\mathbb QK$ and
$V=\widetilde{\mathbb QI}$ for the preceding source lists.
Their exact certificate partition is
\[
\begin{array}{c|r|r|r|r|r}
\text{levels}&\operatorname{pd}V\leq1&\substack{\text{stable}\\\text{Tor}}&
 \substack{\text{local simple}\\\text{tail}}&\substack{\text{even degree-two}\\\text{Tor}}&
 \substack{\text{remaining}\\\operatorname{pd}V\leq2}\\\hline
4+1&24&9&0&0&0\\
3+2&35683&22951&22&20&1547\\
(-1,-1,0,0,1)&12&0&0&0&0.
\end{array}
\]
For a stable or local simple tail, eventual vanishing of the Tor
groups implies
\[
 \operatorname{Tor}_2^A(W,V)=0
\]
for every finite-dimensional right $A$-module $W$.
In the next column this degree-two group has even dimension for every $W$.
\end{proposition}
\begin{proof}
As in Proposition~\ref{C:four-row-six-outside-row-reduction}, every
projective term is a literal direct sum of ideals $A\varepsilon$.
Full bases in $A$, idempotents, cover maps, embedded kernels, and
all equivariance equations are retained. For the $3+2$ list the
first two covers give $16266$ projective rows and $19375$ rows of
projective dimension at most one. They leave $24536$ higher kernels
and $46$ nonsplit quotients.

For $22932$ of those kernels, an exact identity factors through
a specified later syzygy and a finite free module. The syzygy
distance is one or two. Applying positive Tor and dimension shifting
propagates any nonzero degree-two Tor to arbitrarily high degrees.
Fifteen further cases have the same explicit stable-retract identity
at distance three. The nine nonprojective $4+1$ rows have the
same certificates. All twelve three-value rows are projective.

After the fifteen distance-three cases, the other $42$ split-basic
exceptions lie at one local corner. For $22$ there is an idempotent $e$
such that
\[
 Ae=eAe=B,\qquad\dim B=3,\qquad
 \dim\operatorname{rad}B=2,\qquad(\operatorname{rad}B)^2=0.
\]
The actual first kernel $K=\Omega_A V$ satisfies $eK=K$, and
the displayed later syzygy is
$\Omega_A^2K\cong\mathbb Q^{\oplus2}$, with the residue action
through $B$. Since $Ae=B$, induction is exact, takes projectives
to projectives, and gives
\[
 \operatorname{Tor}_i^A(W,K)
  =\operatorname{Tor}_i^B(We,K).
\]
The same formula applies to the displayed simple tail. Eventual
vanishing therefore implies eventual vanishing of
$\operatorname{Tor}_i^B(We,\mathbb Q)$. A minimal right free
resolution of $We$ has zero differentials after tensoring with
the residue field, so it must be finite. A finite-dimensional
local algebra has no nonprojective finite-dimensional module of
finite projective dimension: a last minimal free differential
has radical entries and kills a nonzero vector from the left
socle. Thus $We$ is projective and
$\operatorname{Tor}_2^A(W,V)=\operatorname{Tor}_1^B(We,K)=0$.

For the other twenty cases, $Ae=eAe=B$ is four-dimensional local,
with radical powers of dimensions $3,1,0$, and an explicit
invertible module intertwiner gives
\[
 K\cong(\operatorname{rad}B)^{\oplus2}.
\]
The two dimension shifts give, for $n\geq2$,
\[
 \operatorname{Tor}_n^A(W,V)
  \cong\operatorname{Tor}_n^B(We,\mathbb Q)^{\oplus2}.
\]
In particular its degree-two dimension cannot be one.
All uses of induction here require and certify the equality
$Ae=eAe$, not merely the existence of the corner algebra $eAe$.

Of the $46$ nonsplit quotients, $37$ have an exact identity on
their radical factoring through finite free modules. Their radical
is projective, so all row modules have projective dimension at most
one. In the other nine, a literal local cyclic group
$\{e,g,g^2\}$ of order three supplies the central cyclotomic
idempotent
\[
 c=(2e-g-g^2)/3.
\]
Its two-dimensional ideal is retained, even when it acts nontrivially
on the row module. Orthogonal idempotents in the complementary
four-dimensional algebra, together with $c$, give literal
projective presentations. Five row modules are projective and four
have exact stable-retract identities. No assumption that averaging
acts as the identity on $V$ is used.

Only the $1547$ final rows remain. Each has an exact projective
resolution of length at most two. Its last term is certified
projective by the zero kernel of the next displayed cover; no
claim of exact projective dimension is needed.

The producer records are the three
\path{code/outputs/d_five_row_covers16_*.json.gz} files, the
\path{code/outputs/d_five_row_stable16_41.json.gz} and
\path{code/outputs/d_five_row_stable16_32.json.gz} files, and
\path{code/outputs/d_five_row_exceptional_local_periodic.json}
and \path{code/outputs/d_five_row_exceptional_nonsplit.json}.
The separate Fraction checker is
\path{code/d_verify_five_row_certificates.py}; it imports no
producer, Flint routine, or representation classifier.
Its complete terminal record is
\path{code/outputs/d_five_row_certificates_verified.json}.
It checks all $60268$ literal source tables, including $246857728$
associativity equations, $106480$ exact projective presentations,
$1277760$ cover and kernel equivariance equations, and $22997$
exact identity factorizations. It separately verifies the local
corner and cyclotomic certificates. Source-census completeness
and the finite-resolution action gates are distinct checks.
\end{proof}

\begin{theorem}[Five rows and four states]
\label{C:five-row-six-outside-four-state-cones}
Let $M$ have minimal band $5\times2$ and six outside elements.
No four-state right act has cone of type $\mathbb CP^2$.
\end{theorem}
\begin{proof}
A disconnected action or a state fixed by the minimal band is
excluded by the preceding homology and relative-grid arguments.
The four-state band dichotomy
(Lemma~\ref{C:four-state-band-dichotomy}) therefore makes the
band action constant. If the old two column states coincide,
the old affine column terms vanish and every relevant square is
zero. Otherwise label these two states zero and one.

Constant-action compression is an actual homotopy equivalence to
a monoid with band $5\times4$, with the same contracted algebra
$A$ and row module $V$. Its column augmentation module $W$ has
dimension three. A projective plane requires
\[
 \bigl(\dim\operatorname{Tor}_0^A(W,V),
       \dim\operatorname{Tor}_1^A(W,V),
       \dim\operatorname{Tor}_2^A(W,V)\bigr)=(1,0,1)
\]
and zero higher Tor. The preceding algebra partition leaves exactly
the $1547$ finite-resolution sources.

The generator-image and stabilizer enumeration tries every
constant-band four-state act of these sources. All band actions and
the two labelled core states are fixed; only the two other states
are permuted. Power constraints are necessary conditions only;
multiplication closure and all literal action equations are checked
at each leaf. It yields $115229$ acts. Tensoring the exact
projective resolutions with their actual column modules gives
$4094$ candidates with the required rational Tor triple.

The literal twenty-one-arrow categories for all $4094$ candidates
have zero mod-two square. Each retained certificate contains a
nonboundary degree-two cocycle and a degree-three primitive of its
square. The producer checks $21243766$ typed associativity equations,
$17501850$ cocycle equations, and $262527750$ primitive equations.
Thus none has the projective-plane cup product.

The complete act and tensor records are
\path{code/outputs/d_five_row_constant_acts16_s4.tsv},
\path{code/outputs/d_five_row_constant_acts16_s4.out}, and
\path{code/outputs/d_five_row_constant_tensor16_s4.json}.
Every candidate is linked to its literal category in
\path{code/outputs/d_five_row_constant_categories21_index.json}.
All cup primitives are retained in
\path{code/outputs/d_five_row_constant_cups_all.jsonl}.
The cup producer has capacity for $8192$ normalized three-chains;
these categories have $4275$, so no four-row capacity assumption
is silently reused.
A separate pointwise-domain action enumerator, using neither the
producer's generator choices nor its closure order, recovers all
$216087$ labelled acts and exactly the same $115229$ orbits under
the permutation of the two remaining states.  It checks each actual
source table and every action equation.  A separate static projective
checker and tensor calculation over $\mathbb F_{65521}$ recover the
same $4094$ candidates.  The projective presentations and their exact
kernels are checked in that field.  This is a finite-field necessary
test for a projective-plane cone; no equality of rational and modular
ranks is assumed.

The independent literal-category verifier reconstructs every category
from its monoid and act, checks that its degree-two cocycle is
nonboundary, and checks every displayed square primitive.  It
verifies all $21243766$ typed associativity equations, $17501850$
cocycle and differential-square equations, and $262527750$ primitive
equations.  The complete action, tensor and cup record is
\path{code/outputs/b_five_row_constant_act_tensor_cups_verified.json}.
The final gate
\path{code/outputs/b_five_row_sources_acts_complete_verified.json}
links every source index and literal finite-resolution table to
these action records and all $4094$ checked categories.  All three
source lists, all algebra certificates, and every action branch in
the stated four-state scope are therefore covered.
\end{proof}

% Drafting record; all independent gates closed on 2026-11-16 at 20:13 UTC.
\subsection{The final four-row act family}

\begin{lemma}[The complete two-plus-two source family]
\label{C:four-row-seven-outside-two-plus-two-sources}
Let $M$ have minimal band $4\times2$, seven outside elements, and
a one-dimensional old-column quotient-covector space with row-level
partition $2+2$.  Suppose that its nonconstant-column actors have
no common rational affine row fixed point.  There are exactly
$194497$ ordinary-isomorphism classes of such monoids.
For $A=\mathbb QM/\mathbb QK$ and its actual row module
$V=\widetilde{\mathbb QI}$, the following source partition holds:
\[
\begin{array}{c|r}
\text{verified alternative}&\text{sources}\\\hline
\operatorname{pd}_A V\leq1&117061\\
\text{stable positive-Tor recurrence}&63656\\
\text{remaining exact length-two resolutions}&11537\\
\text{remaining exact prefixes through degree two}&2243.
\end{array}
\]
The last prefixes have projective terms $P_0,\ldots,P_3$, with
$P_3$ surjecting onto $\ker d_2$.  No termination claim is made
about their undisplayed tails.
\end{lemma}
\begin{proof}
Normalize the row levels to $(0,0,1,1)$.  The exact profile
equations of Lemma~\ref{C:four-row-six-outside-source-census}
give $96$ transformation pairs, with $32$ having nonconstant
columns.  There are $303$ closed nonconstant-column profile monoids
of order at most seven containing the identity; $138$ have no
common rational affine fixed point.  Keeping all multiplicities
of outside profiles gives $2201951$ profile multisets, $139863$
after the permitted relabelings, and $3311$ with closed support.
Distinct outside actors with identical profiles, including
identity and band profiles, remain distinct actors.
Every compatible product value is tried, and all associativity
equations are imposed.  This yields $937564$ raw tables and
$194497$ ordinary-isomorphism classes.

The independent source check uses a different closure tree and
reverse-column product assignment, without the producer's
canonicalizer.  Full literal relabeling expansions of the retained
classes are disjoint and equal the newly generated raw set.
The raw search has $71388326$ nodes and checks
$3164278500$ associativity equations.  Its complete record is
\path{code/outputs/b_four_row15_split22_sources_independently_verified.json}.
The literal source stream is
\path{code/outputs/a_four_row_sources15_22.tsv}.

Sharing a module certificate requires entrywise equality of the
contracted multiplication table and all row-action matrices.
There are $101449$ distinct such actual pairs $(A,V)$.
The initial exact covers give $55453$ projective rows, $61238$
rows with projective first kernel, and $63604$ stable recurrences.
In each recurrence, $L=\Omega V$ has a literal identity factoring
through a finite sum of $\Omega^sL$, $s=1$ or $2$, and free modules.
The positive-Tor dimension-shifting argument of
Proposition~\ref{C:four-row-six-outside-row-reduction} shows that
bounded $\operatorname{Tor}_*^A(W,V)$ forces
$\operatorname{Tor}_2^A(W,V)=0$.

The remaining certificates are also literal idempotent-projective
presentations.  A projective radical gives $303$ further sources
with projective dimension at most one, since every simple is a
summand of $A/\operatorname{rad}A$.  Explicit alternative projective
covers give another $22$ such sources.  In the remaining nonsplit
cases a central finite-cyclic averaging idempotent acts as the
identity on $V$.  Its full central summand, its complementary
zero action on $V$, and every cover are checked literally.
This gives $33$ projective rows, $12$ more rows of projective
dimension at most one, and $52$ additional stable recurrences.
The first two columns of the stated partition therefore total
$117061$ and $63656$.

The other $13780$ sources share $7301$ actual projective prefixes.
For $11537$ sources the row resolution ends at degree two;
the other $2243$ have an explicitly verified next projective
surjection.  These prefixes determine $\operatorname{Tor}_0$,
$\operatorname{Tor}_1$, and $\operatorname{Tor}_2$ for every right
module, which is all the subsequent necessary-homology test uses.
The complete source-to-module index is
\path{code/outputs/a_four_row_q7_modules_22_resumed.json};
the closed original certificate chunks have prefix
\path{code/outputs/a_q7_22_closed_modules_}.
The supplementary certificates and exact status partition are
\path{code/outputs/a_q7_residual_row_prefixes.jsonl.gz} and
\path{code/outputs/a_q7_residual_row_prefixes.json}.
The interrupted historical stream is explicitly not used as a
complete index.

The separate finite-field verifier
\path{code/d_verify_q7_modular_prefixes.py} imports no producer
or Flint routines.  It checks all $7301$ prefixes, including
$45682$ full idempotent ideals, $23298$ exact covers,
$326172$ cover/kernel equivariance equations, $111979$ chain
equations and $11270639$ denominator checks.  The exported
idempotents and differential blocks agree literally with the
certificates, as do all $13780$ source-to-module links.  Its
complete record is
\path{code/outputs/d_q7_modular_prefixes_verified.json}.
The full rational checker
\path{code/d_verify_q7_row_certificates.py} uses independent
\texttt{Fraction} elimination and imports neither producer code
nor Flint.  It checks all $194497$ literal sources and all
$101449$ actual module pairs: $656427375$ source associativity
equations, $183156$ exact projective presentations,
$2564184$ cover/kernel equivariance matrix equations, and
$32165$ exact identity factorizations.  It also reconstructs the
nonsplit radicals, semisimple quotients and central cyclic blocks.
Its complete partition is
\path{code/outputs/d_q7_row_certificates_verified.json}.
\end{proof}

\begin{theorem}[Seven outside elements and five states]
\label{C:four-row-seven-outside-five-state-cones}
Let $M$ have minimal band $4\times2$ and seven outside elements.
No right act with five states gives a cone of type $\mathbb CP^2$.
\end{theorem}
\begin{proof}
Suppose such a cone $X$ were a projective plane.  A band-fixed
state makes every degree-two class relative and its square zero.
Without such a state every core orbit has two points.  If the
band action is nonconstant, connectedness of the five-state act
forces exactly two core orbits and one other state.  The
projective-orbit kernel of
Lemma~\ref{C:two-column-projective-orbit-kernel} and its finite-model
consequence give a compressed $4\times3$ monoid with a finite
four-dimensional model, $H_2=\mathbb Z^2$ and $H_4=\mathbb Z$.
The rank-two conclusion in the proof of
Theorem~\ref{B:three-by-four-targets}, after opposition, kills all
its rational degree-two products.  Its degree-two classes pull
back onto those of $X$.  This excludes the nonconstant case
without a bound on the number of outside actors.

We may therefore suppose the band acts constantly.
A common affine row origin kills the
degree-two products.  The quotient-covector arguments leave only
the two row-level partitions, and
Corollary~\ref{C:four-row-seven-outside-split-act-cones} excludes
$3+1$.  Thus the source belongs to the preceding lemma.
If the two old columns have the same constant state value, the
affine cochain vanishes on every old/old pair and its square is
zero.  Otherwise label the two core states zero and one.

Constant-action compression is an actual homotopy equivalence
to a monoid with the same contracted algebra $A$ and row module
$V$.  The projective-dimension and stable-Tor alternatives in the
lemma exclude its having bounded homology and nonzero $H_4$.
It remains to examine the $13780$ sources with the recorded
projective prefixes.

For each of these, fix the identity transformation and the eight
constant band transformations.  A finite generating set among
the six other outside actors is found by literal multiplication
closure.  Its images range over every five-state transformation
with the required two core values and the necessary power
relations.  Multiplication closure propagates each assignment,
and every full action satisfies all $15^2$ action equations.
At each generator, canonical augmentation retains exactly one
orbit under the stabilizer of the preceding assigned images in
the permutation group of the three remaining states.  An image
already determined by previous generators needs no new stabilizer
restriction, since it is a word in those images.  This is the
same complete generator-image argument used for the six-outside
family.

There are exactly $19338315$ such canonical actions.  The exact
projective prefixes remain exact over $k=\mathbb F_{65521}$,
as checked on their full ideals, embedded kernels and cover maps.
For each action put $W=\widetilde{kS}$.  Tensoring a summand
$A\varepsilon$ gives the image $W\varepsilon$, and a differential
is right multiplication by its literal generator value.
All ranks are calculated over this field, without inferring a
rational rank from a modular calculation.  Exactly $870480$
actions satisfy the necessary condition
\[
 (\dim\operatorname{Tor}_0^A(W,V),
   \dim\operatorname{Tor}_1^A(W,V),
   \dim\operatorname{Tor}_2^A(W,V))=(1,0,1).
\]
The completed interval records, all per-source counts and hashes,
and every survivor are indexed by
\path{code/outputs/a_q7_constant_act_progress.json}.
The generating and tensor routines are
\path{code/a_q7_constant_act_tensor.cpp} and
\path{code/a_q7_constant_tensor_static.py}.
Coverage here is proved by the audited generator-image and exact
stabilizer enumeration; no second independent numerical act
enumerator is claimed.  The recorded FNV64 values are diagnostic
stream-integrity checks only.  They are never used to identify
orbits, assert table equality, or establish completeness.

First, $467648$ survivors admit short exact rational exclusions.
For each, a nonzero integral $4\times3$ balanced form $F$ is
retained and all balancing equations are verified over $\mathbb Z$.
The same relation matrix has nullity one over $k$; thus the
rational balanced space is exactly the line spanned by $F$.
In $337486$ cases its old-column row is zero.  The actual affine
cocycle then vanishes on every old/old pair; each normalized
four-simplex has at most one cross arrow, so one of its two cup
factors is zero.  In the other $130162$ cases $F$ has rank one,
as certified by all integral two-by-two minors.  The universal
rank-one affine-square primitive excludes these cases.
The independent checker
\path{code/a_verify_q7_affine_partition.cpp} imports no producer
code and checks $84176640$ integral balancing equations,
$467648$ independent nullity-one tests, and $526104000$ literal
action equations.

For the remaining $402832$ actions, use
Proposition~\ref{C:act-projective-affine-square}.
There are $3273$ sources.  For each source a separate full
monoid-algebra projective prefix resolves the trivial module
through degree two; its degree-three map surjects onto the preceding
kernel.  Every term is a literal idempotent ideal and every
comparison generator is an explicit normalized free-bar chain.
The recursion and all support and boundary equations are retained.
The nonzero balanced field form is the actual nonboundary
degree-two class.  Its square pulls back to the zero cochain in
$396159$ actions and to an explicit coboundary with a nonzero
primitive in the other $6673$.  Thus every remaining square is
zero over $k$, which excludes $\mathbb CP^2$.

The prefix and comparison index is
\path{code/outputs/a_q7_Rbar_prefix_progress.json};
the literal cup and primitive records are
\path{code/outputs/a_q7_Rbar_square_certificates.tsv},
with complete counts in
\path{code/outputs/a_q7_Rbar_square_summary.json}.
The fast evaluator is \path{code/a_q7_Rbar_squares.cpp}.
The separate ordinary-category control
\path{code/a_q7_chronological_control.py}
constructs one of the nonzero pulled-back square cochains from
its literal twenty-one-arrow table.  Its normalized dimensions
are $19,266,3724$, its boundary ranks are $18,247$, and its
affine class is a nonboundary cocycle.  All $704$ literal
Alexander--Whitney evaluations agree with the stored projective
cochain.  This pins down the chronological convention on an
example requiring a primitive, rather than only a zero cochain.

Finally \path{code/a_verify_q7_affine_partition.py} checks the
literal disjoint partition of all $870480$ survivors, exact
remaining-stream equality, and each source/action/cup record's
linkage.  Its complete record is
\path{code/outputs/a_q7_affine_partition_independently_verified.json}.
The independent routines
\path{code/b_verify_q7_Rbar_prefixes.py} and
\path{code/b_verify_q7_Rbar_squares.cpp} import no producer code.
They check all $3273$ full-monoid prefixes, $25200$ literal
idempotent ideals and supported comparison generators, and every
generator boundary equation.  Their separate action/cup pass
checks all $402832$ actions, $453186000$ action equations,
$72509760$ balancing equations, $44370992$ chronological
Alexander--Whitney evaluations and $53512$ primitive equations.
Every one of the stated $396159$ zero cochains and $6673$
nontrivial primitives passes.  The complete independent records are
\path{code/outputs/b_q7_Rbar_prefixes_independently_verified.json} and
\path{code/outputs/b_q7_Rbar_squares_independently_verified.json}.
The aggregate record
\path{code/outputs/b_q7_bar_act_complete_verified.json}
checks all twenty-nine contiguous completed source intervals,
all per-source counts, the exact static export, and every source,
action, partition and cup-certificate linkage by literal comparison
and SHA256.  It explicitly retains the distinction between the
audited action-enumeration proof and the independent numerical
algebra and cup checks.
\end{proof}

\begin{corollary}[The exact act-cone minimum]
\label{C:act-cone-minimum-twenty-two}
The least number of arrows in a finite act-cone category with
classifying space of type $\mathbb CP^2$ is twenty-two.
\end{corollary}
\begin{proof}
All smaller act cones are excluded by the preceding theorem and
Corollary~\ref{C:act-cone-twenty-one-frontier}.
The construction in Theorem~\ref{C:CP2-upper22} attains the bound
through an actual homotopy equivalence to the certified
twenty-three-element monoid.
\end{proof}


% Complete proof independently audited by root and agents B and D.
\subsection{Small act cones under finite-target nullity}

\begin{lemma}[The precise finite-target uses in the small band hulls]
\label{C:small-band-finite-target-cups}
Let $P$ be a finite monoid with rectangular-band minimal ideal.
Assume that every integral class in $H^2(BP;\mathbb Z)$ restricts
to zero along every map from the classifying space of a finite group.
Then the following consequences of the earlier hull calculations hold.
\begin{enumerate}
\item For a $3\times3$ band, all rational degree-two products vanish.
\item For a $3\times4$ band with rational second-cohomology rank two,
all rational degree-two products vanish.
\item For a $3\times4$ band with rational second-cohomology rank one,
the square of an integral generator cannot be nonzero and primitive modulo torsion
in fourth cohomology.
\end{enumerate}
The same statements hold after opposition. There is no finite-CW-type
assumption on $BP$ in this lemma.
\end{lemma}
\begin{proof}
Read the proofs of Theorems~\ref{B:three-by-three-cups} and
\ref{B:three-by-four-targets} with the stated nullity hypothesis.
Every use there of finite CW type is to rule out an actor with a
nonzero integral character on its eventual cyclic group. A preimage
actor has an eventual finite group surjecting onto that cyclic group;
inflation of the character is injective. The corresponding integral
hull class pulls back to a class on $BP$, contradicting the present
hypothesis. All other cases use actual cup primitives, rank-one
square primitives, or zero fourth rational homology in a containing
hull. Those conclusions pull back without any finiteness premise.

The only nonzero-square case for a $3\times4$ band is the rank-one
even-Hopf hull. Its integral second homology is $\mathbb Z$, and the
square of its primitive degree-two class evaluates evenly on every
integral four-cycle. The balanced-tensor presentation makes the map
from $P$ onto this hull surjective on integral second homology.
Even if $H_2(BP;\mathbb Z)$ has torsion, its free rank is one,
and a surjection to $\mathbb Z$ is an isomorphism on the free part.
The hull generator therefore pulls back to a generator on $BP$.
Its square still evaluates evenly, which prevents primitivity modulo
torsion. In rank two this exceptional hull cannot occur at all.
This proves the three assertions, including the integral precision.
\end{proof}

\begin{lemma}[Nullity survives the projective orbit quotient]
\label{C:orbit-quotient-local-nullity}
Let $X=B\mathcal C(M,S)$ be simply connected, with minimal band
$I\times\{0,1\}$ and no band-fixed state. Suppose every integral
second-cohomology class of $X$ restricts to zero on every finite $BG$.
The constant-action quotient and its compression $Y\simeq BP$ have
the same nullity property in degree two.
\end{lemma}
\begin{proof}
Corollary~\ref{C:two-column-orbit-wedge} gives
$Y\simeq X\vee\bigvee_{o-1}S^2$, under the actual quotient map.
Every degree-two class on this wedge is a sum of pullbacks from the
two projections to $X$ and to $\bigvee_{o-1}S^2$.
The first restrictions vanish by hypothesis. The second vanish by
Miller's theorem with the finite target $\bigvee_{o-1}S^2$.
Thus every such restriction vanishes on $Y$ as asserted.
\end{proof}

\begin{lemma}[The rational five-outside row argument at six states]
\label{C:rational-four-row-five-outside-six-states}
Let $M$ have minimal band $4\times2$ and exactly five outside
elements. An act cone on five or six states cannot have the rational
homology and cohomology ring of $\mathbb CP^2$.
\end{lemma}
\begin{proof}
The fixed-state case has zero degree-two products. Otherwise the
projective orbit kernel gives a constant-action compression with the
same rational homology in degrees at least three. If $o$ is the
number of core orbits, its row and column modules satisfy
\[
 \dim V=3,\qquad \dim W=|S|-2o+1,
 \qquad \dim\operatorname{Tor}_{0,1,2}^A(W,V)=(o,0,1),
\]
and higher Tor vanishes. Here $A$ is the five-dimensional contracted
algebra of $M$. This uses only the rational homology assumptions.

The local case is the universal affine obstruction
Lemma~\ref{C:local-two-column-act-cups}. The nonlocal classification
and the complete module tables used in
Theorem~\ref{C:four-row-five-outside-seven-states} leave the line
and Nakayama algebras. For Nakayama, the only possible row module
is $S_2\oplus Ae_2$, and its sole quotient character occurs twice;
the missing-character or two-covector affine-origin lemma gives zero
products. For the line, the split-character row and the row with a
twice-occurring quotient character are excluded by the same universal
arguments. The remaining row is $V=S_3\oplus L_a$, and the module
table gives
\[
 W=T_1\oplus T_3^{\oplus t}\oplus R_a^{\oplus p}
                         \oplus R_b^{\oplus q},\qquad
 o=t+p+q,\qquad |S|=4o-t,
 \qquad 0\leq t\leq o.
\]
Five states admit no solution: $o=1$ gives at most four states,
whereas $o\geq2$ gives at least six. For six states the only
possibility is $o=2,t=2,p=q=0$. Both balanced degree-two forms
then factor through the single $S_3$ row summand and have matrix
rank one. Their squares are zero by Lemma~\ref{B:rank-one-square}.
Their pullbacks span the one-dimensional second cohomology of the
original cone; one pulls back nontrivially and contradicts its
projective-plane square. No finite model for the compression is used.
\end{proof}

\begin{theorem}[A finite-target obstruction through twenty act arrows]
\label{C:finite-target-small-act-twenty}
Let $X=B\mathcal C(M,S)$ for a nonempty finite right act, with
$|M|+|S|+1\leq20$. The following conditions cannot hold simultaneously:
\begin{enumerate}
\item $X$ is simply connected and its rational homology is that of
$\mathbb CP^2$;
\item an integral generator $v\in H^2(X;\mathbb Z)\cong\mathbb Z$
has square nonzero and primitive modulo torsion in $H^4(X;\mathbb Z)$;
\item for every finite group $G$ and every map $BG\to X$, the
restriction of $v$ is zero;
\item over $\mathbb F_{65521}$, the Betti numbers in degrees
two, three and four are $(1,0,1)$.
\end{enumerate}
In particular no finite-CW-type assumption on $X$ is required.
\end{theorem}
\begin{proof}
Write the minimal Rees ideal as $I\times H\times J$, put
$r=|I|$, $c=|J|$, $q=|M|-rc|H|$, and $s=|S|$.
The nonzero square requires $H^2(BM;\mathbb Q)\ne0$; otherwise
every degree-two class is relative to $BM$ and its products vanish
in the suspension quotient. The rational balanced presentation
therefore gives $r,c\geq2$. The all-order two-row act theorem gives
$r\geq3$. The finite-line argument gives $s\geq3$. The identity
lies outside the minimal ideal in these cases, so $q\geq1$.

If $H$ is nontrivial, the arrow budget leaves only
$r=3,c=2,|H|=2,q\leq4$. Lemma~\ref{C:three-row-forced-profiles},
which permits a nontrivial minimal group, requires at least five
outside elements for a nonzero product. Thus $H$ is trivial below.
The original square is nonzero rationally by its primitivity.

For every compression used below, an actual homotopy equivalence
preserves the generator and primitive square. If an orbit quotient
with two-column projective kernel is used instead, its higher
rational homology is unchanged, its second-cohomology rank is the
number $o$ of core orbits, and
Lemma~\ref{C:orbit-quotient-local-nullity} supplies nullity for all
its integral degree-two classes. We explicitly indicate these uses.

\emph{Three states.}
If a band actor is nonconstant, the constant-action quotient has at
most two states. Its degree-two cohomology surjects onto that of the
original cone, and its finite-line products are zero. Otherwise
compression gives an actual equivalent monoid with band $r\times3$
and $q$ outside elements. Here $|M|\leq16$.
For $r=3$ use the first part of
Lemma~\ref{C:small-band-finite-target-cups}; for $r=4$ use its
primitive-square part. For $r=5$, its order is
$15+q\leq15+(16-2r)=21$, so
Proposition~\ref{B:three-by-five-finite-target-transfer} applies after
opposition. If $r\geq6$, the same inequality gives $q\leq4$,
and Theorem~\ref{B:contracted-four-Tor} contradicts bounded
rational homology with nonzero fourth homology.

\emph{Four states.}
Lemma~\ref{C:four-state-band-dichotomy} uses only first rational
homology vanishing. It gives zero products or constant band action.
In the latter case the equivalent compression has band $r\times4$,
with $|M|\leq15$. Three rows are excluded by the primitive-square
part of Lemma~\ref{C:small-band-finite-target-cups}.
For four rows and two old columns, the compressed order is
$16+q\leq23$; apply
Proposition~\ref{B:four-by-four-finite-target-transfer}.
For four rows and at least three old columns, $q\leq3$, so the
bounded Tor obstruction already applies. For five rows, necessarily
$c=2$ and $q\leq5$. Values $q\leq4$ again give that obstruction;
at $q=5$ the compressed $5\times4$ band has augmentation dimensions
four and three. The proof and complete module table of
Lemma~\ref{C:four-by-five-five-outside-CP2} exclude the rational
Tor sequence $(1,0,1,0,\ldots)$ without any finiteness premise.
For at least six rows, $q\leq3$ and
Lemma~\ref{C:three-outside-all-two-column-act-cups} applies.

\emph{At least five states and two columns.}
For three rows the budget gives $q\leq8$, hence $|M|\leq14$.
The bounded-rational assertion of
Theorem~\ref{C:three-row-fourteen-act-cones} excludes the square.
For at least five rows, $q\leq4$ and the stronger rational
assertion of Theorem~\ref{C:two-column-four-outside-act-cones}
applies at every state count.
It remains that $r=4$, where $q+s\leq11$.
The cases $q\leq4$ have the same Tor obstruction; $q=5$ gives
$s=5$ or $6$ and the preceding rational module lemma applies.
The sole remaining possibility is $q=6,s=5$.

In this last case a fixed state makes the products zero.
If the band action is nonconstant, the connected act has exactly
two two-point core orbits and one additional state. Its quotient
has three states. The projective orbit kernel and actual wedge
equivalence give a compression with band $4\times3$, rational
second-cohomology rank two, and nullity of all integral degree-two
group restrictions. The rank-two clause of
Lemma~\ref{C:small-band-finite-target-cups} makes its products
zero, and pullback gives the contradiction.

If instead the band acts constantly, use the exact source census
and row-module partition in
Lemma~\ref{C:four-row-six-outside-source-census} and
Proposition~\ref{C:four-row-six-outside-row-reduction}.
Projective-dimension-one rows give zero products universally;
the stable and local-corner rows give zero fourth rational homology
under boundedness. For the remaining $173$ rows, the independently
verified exact projective resolutions over $\mathbb F_{65521}$ and
the exhaustive five-state act calculation in
Theorem~\ref{C:four-row-six-outside-six-state-cones} apply.
All $110418$ constant five-state acts fail the necessary tensor
dimensions $(1,0,1)$ in Tor degrees zero, one and two.
These are the actual Betti numbers in degrees two, three and four
of the compressed cone. Condition (4) is therefore impossible.
This uses exact field resolutions and a field homology obstruction,
not an inference from a rank measured at an uncontrolled prime.

\emph{At least five states and at least three columns.}
The budget leaves only bands $3\times3$, $3\times4$, and $4\times3$.
For either twelve-element band, $q+s\leq7$. Thus $q=1$ gives a
projective row module; otherwise $q=2$ and $s=5$. If the action is
constant, bounded rational Tor with $q\leq4$ is impossible.
If it is not, connectedness and the absence of a fixed state give
at least two core orbits, at least one of which has two points.
Its column map makes the old-column image in the constant-action
quotient have at most two points. The universal transfer in
Lemma~\ref{C:old-column-image-two-transfer}, followed by
Lemma~\ref{C:three-outside-all-two-column-act-cups}, kills products.

For the $3\times3$ band, $q+s\leq10$.
An old-column image of size one gives zero old/old affine cochains.
Size two gives a literal common compression over a $3\times2$
monoid with the same $q\leq5$ outside elements, whose act products
vanish universally by Theorem~\ref{C:three-row-fourteen-act-cones}.
If all three old-column values remain distinct, every core orbit
has three points. Lemma~\ref{C:full-column-orbit-kernel} then
preserves bounded higher rational homology under the quotient.
For $q\leq4$ use Theorem~\ref{B:contracted-four-Tor}.
The only other possibility is $q=5,s=5$. Two full core orbits
cannot fit, so the band already acts constantly. Compression is an
actual equivalence to a $3\times5$ band monoid of order twenty,
excluded by Proposition~\ref{B:three-by-five-finite-target-transfer}.
This exhausts the arrow budget and proves the theorem.
\end{proof}

\begin{corollary}[Nontrivial projective opposite corners through twenty-one arrows]
\label{C:projective-opposite-corner-twenty-one}
Let $C$ be a triangular two-object category with endomorphism monoids
$M,N$, cross-biset ${}_NS_M$, and projective trivial right
$\mathbb QN$-module. If $|N|\geq2$ and $|C|\leq21$, then
$BC\not\simeq\mathbb CP^2$.
The same conclusion holds with $|C|\leq22$ and $|N|\geq3$.
Both statements also hold after opposition.
\end{corollary}
\begin{proof}
Assume otherwise. The actual orbit quotient of
Theorem~\ref{C:finite-target-opposite-corner} has at most
$|C|-|N|+1\leq20$ arrows.
Corollary~\ref{C:finite-target-corner-CP2} supplies simple
connectivity, rational projective-plane homology, the primitive
integral square, and nullity of every integral degree-two class
on every finite group. The group order of the opposite corner
is at most $22$, so $65521$ is a good prime and the same comparison
supplies condition (4). The preceding theorem gives a contradiction.
\end{proof}

\begin{remark}[The remaining two-object scope]
\label{C:finite-target-small-act-scope}
The argument treats the precise finite-target hypotheses above,
rather than declaring every rational projective-plane act cone to
have finite CW type. At twenty-one arrows, an unexcluded triangular
category must fail both projective-corner tests: its source corner
must have at least two minimal rows and its target corner at least
two minimal columns. This follows by applying the corollary also
to the opposite category. It does not dispose of these nonprojective
opposite corners. Strongly connected two-object categories are resolved
separately by Theorem~\ref{C:twentytwo-strongly-connected}; the
present theorem itself concerns only triangular categories. The
original unrestricted category minimum is not asserted here.
\end{remark}

\subsection{A fixed incoming arrow and the extreme eight-element corner}

\begin{lemma}[A fixed incoming arrow]
\label{C:fixed-incoming-cups}
Let $C$ be a finite category whose objects consist of a full subcategory
$D$ and one additional object $v$, with no arrows from $v$ to $D$.
Write $M=\operatorname{End}_C(v)$.  Suppose
$H^2(BD;\mathbb Q)=0$ and there is an arrow $f:d\to v$, for some
$d\in D$, such that $mf=f$ for every $m\in M$.
Then
\[
 H^2(BC;\mathbb Q)\otimes H^2(BC;\mathbb Q)
       \longrightarrow H^4(BC;\mathbb Q)
\]
is the zero map.  The opposite assertion holds for a source object
and a fixed outgoing arrow.
\end{lemma}
\begin{proof}
The correspondence $d\mapsto\operatorname{Hom}_C(d,v)$ has commuting
right $D$- and left $M$-actions.  Its two-sided bar space $E$ gives the
natural homotopy-pushout description
\[
 BC\simeq BD\mathop{\cup}\limits_E^h BM.
\]
This description follows directly by subdividing each mixed nerve simplex
at its unique crossing from $D$ to $v$; the pieces are the double bar,
with its two endpoint projections.

The chosen arrow gives an actual section $BM\to E$ of the projection
$E\to BM$.  In the double bar, keep the $D$-coordinate at the object $d$
and the correspondence-coordinate at $f$, and retain all the $M$-arrows.
The equalities $mf=f$ verify every endpoint face; the other faces and
degeneracies are the ordinary bar identities.  Thus this is a simplicial
section, rather than a merely linear splitting on homology.  Consequently
$H^2(BM;\mathbb Q)\to H^2(E;\mathbb Q)$ is injective.

Let $X$ be the double-mapping-cylinder model of the displayed pushout,
and let $A=BD\sqcup BM\subset X$ be its two endpoint spaces.
The Mayer--Vietoris sequence and $H^2(BD;\mathbb Q)=0$ show that
$H^2(X;\mathbb Q)\to H^2(A;\mathbb Q)$ is zero.  Every degree-two
class therefore lifts to $H^2(X,A;\mathbb Q)$.  The quotient is
\[
 X/A\cong (E\times[0,1])/(E\times\{0,1\})
       \cong \Sigma(E_+),
\]
where the entire endpoint subspace is collapsed to one basepoint.
Positive-degree products in its reduced cohomology vanish because it is
a suspension.  Relative products followed by the map to absolute
cohomology now prove the claim.
\end{proof}

\begin{corollary}[The extreme eight-element corner at fifteen arrows]
\label{C:fifteen-extreme-eight}
None of the following linear three-object cardinality profiles, nor
any of their opposites, can have type $\mathbb{CP}^2$:
\[
 (m,n,p,x,y,z)=(1,1,8,2,1,2),\quad
 (1,1,8,2,2,1),\quad (2,1,8,1,1,2).
\]
Here $(m,n,p)$ are the endomorphism orders at $0,1,2$ and
$(x,y,z)$ are the cardinalities of the hom-sets from $0$ to $1$,
from $1$ to $2$, and from $0$ to $2$, respectively.
No classification of the eight-element corner is required.
\end{corollary}
\begin{proof}
Take $D$ to be the full subcategory on $0,1$ and the sink to be $2$.
For the first two profiles, $BD$ is the circle formed by two parallel
arrows, so $H^2(BD;\mathbb Q)=0$.  In the first profile the singleton
hom-set from $1$ to $2$ supplies a fixed incoming arrow.  In the second,
the singleton hom-set from $0$ to $2$ supplies it.  For the third profile,
object $1$ is terminal in $D$, because its endomorphism monoid and the
hom-set from $0$ to $1$ are singletons.  Hence $BD$ is contractible, and
the singleton hom-set from $1$ to $2$ again supplies the fixed arrow.
Lemma~\ref{C:fixed-incoming-cups} applies in all three cases, whereas
the rational degree-two generator of $\mathbb{CP}^2$ has nonzero square.
\end{proof}


\subsection{The complete twelve-arrow triangular exclusion}

\begin{lemma}[A singleton return hom-set]\label{C:singleton-return}
Let $C$ be a category with two objects $a,b$, with
$\Hom(b,a)=\{t\}$ and $\Hom(a,b)\ne\varnothing$.
Then $BC$ is contractible. More generally, the same conclusion holds
for any category with an object $a$ such that every
$\Hom(b,a)$, $b\ne a$, is a singleton and at least one
$\Hom(a,b)$, $b\ne a$, is nonempty. The opposite criterion also holds.
\end{lemma}
\begin{proof}
Put $M=\End(a)$ and choose $s:a\to b$, where $b\ne a$.
If $t:b\to a$ is the unique return arrow, every $m\in M$
satisfies $mt=t$. Thus $e=ts$ satisfies $me=e$ for all $m$.
It is the component of a natural transformation from the constant
functor to the identity functor on $M$, so $BM$ is contractible.
For the full inclusion $i:\{a\}\to C$, the comma category
$(a\downarrow i)$ has initial object $1_a$. For $b\ne a$,
$(b\downarrow i)$ has its single object $t$ and endomorphism
monoid $M$, since $mt=t$. Every comma category is therefore
contractible. Quillen's Theorem A gives $BM\simeq BC$.
Applying the same proof to $C^{\mathrm{op}}$ gives the dual assertion.
No finiteness assumption is needed.
\end{proof}

\begin{remark}[Retract pruning in a twelve-arrow search]\label{C:twelve-retract-pruning}
If $u:x\to y$ and $v:y\to x$, $x\ne y$, satisfy $vu=1_x$,
the full subcategory obtained by deleting $x$ is Cauchy dense.
Adding or removing a retract preserves the classifying space, as in
Lemma~\ref{C:two-object-cup-pruning}. At twelve arrows this smaller
category has at most eleven arrows, so the completed eleven-arrow
exclusion rules out $\mathbb CP^2$. In particular a trivial
endomorphism vertex with arrows both ways can be pruned before any
composition is chosen. There is one further deletion.  Suppose a vertex $x$ has exactly
one arrow from $x$ to any of the other vertices combined, say
$x\to y$, and $\End(y)=\{1_y\}$.  Its own endomorphisms and
incoming arrows are unrestricted.  For the full inclusion deleting
$x$, the comma category $(x\downarrow i)$ is the single arrow
$x\to y$ with endomorphism monoid $\End(y)=\{1_y\}$.
Retained vertices have identity initial objects in their comma
categories.  Quillen's Theorem A therefore permits deletion of $x$.
The criterion with incoming arrows is dual.  The file
\texttt{c\_categories12\_cofinal.cpp} uses precisely this stronger
deletion; the preceding \texttt{reduced} version used only its
special case for a leaf with trivial endomorphisms.
The other reductions are used in the new
\texttt{c\_categories12\_minimal.cpp} search; they do not alter the
recorded scopes or counts of the earlier enumerations.
\end{remark}

\begin{lemma}\label{C:twelve-triangular}
No triangular two-object category with at most twelve arrows has
classifying space homotopy equivalent to $\mathbb CP^2$.
\end{lemma}
\begin{proof}
Let $(m,n,s)$ denote the two endomorphism orders and the nonzero
cross-hom cardinality, with $m\geq n$. The preceding eleven-arrow
proof and Lemma~\ref{C:two-object-cup-pruning} leave, at twelve arrows,
only
\[
 (8,1,3),\ (7,2,3),\ (6,3,3),\ (5,4,3),\
 (5,3,4),\ (5,5,2),
\]
or act-cone cases over at most seven elements. To see that this list
is exhaustive, singleton cross-hom sets are excluded by the join
description. For two cross arrows a vertex of order at most four
is rationally acyclic and the two-point argument applies; the only
remaining endomorphism split is $(5,5)$. For $s\geq3$, two vertex
orders at most four again give zero rational products. One may
therefore assume $m\geq5$, and enumerate $m+n+s=12$ with $m\geq n$.
All cases with $n=1,m\leq7$ are covered for arbitrary $s$ by
Proposition~\ref{C:seven-act-cones}. That proposition also covers
$(7,2,3)$, because each two-element monoid algebra is semisimple
over $\mathbb Q$.

Here is an independent complete enumeration for $(8,1,3)$.
If $H^2(BM;\mathbb Q)=0$, the relative-cup argument already
excludes the act cone. Otherwise the minimal ideal has at least
two rows and two columns. Its local group must be trivial: a
nontrivial group would require at least eight ideal elements and
a further global identity. Its band is consequently $2\times2$,
$2\times3$, or $3\times2$.

For a $2\times2$ band, label its elements $x_{ij}$ and write the
other elements as $1,a,b,c$. Each extra element specifies a row
and column transformation on two points. The balanced presentation
for rational $H_2$ forces equality of their augmentation characters.
There are exactly six such translation pairs. Choose their types
in nondecreasing order; this removes only permutations of $a,b,c$.
There are $\binom{8}{3}=56$ type profiles. The nine products among
$a,b,c$ may be any of the eight elements with the prescribed
composite translations. Exhaustive associativity propagation on
these choices leaves 4032 associative tables and 492 ordinary
isomorphism classes. Isomorphism recognition uses exactly the
permutations of the two rows, two columns, and three extra elements;
the minimal ideal is characteristic, so this is the complete
isomorphism test in the chosen coordinates.

For a $2\times3$ band there is one extra element. There are
$4\cdot27=108$ translation pairs, and its square is tried at
every compatible table entry. This gives 56 associative tables
and twelve ordinary isomorphism classes. The opposite orientation
gives another twelve classes. Thus the complete list has 516
tables. The program \path{code/c_eight_H2_targets.cpp} constructs
this list without using a published census. The independent Python
check verifies every identity and all $8^3$ associativity equations
of each table. Its exact rational balanced-module ranks discard
six tables with $H_2(BM;\mathbb Q)=0$, leaving 510 targets.

For each of the 510 targets, enumerate all right actions on three
points. Choose a generating set, assign every transformation to
each generator subject only to its exact power relations, extend
along generator words, and verify all multiplication equations.
State permutations remove precisely isomorphic actions. Because
ordinary opposite tables are included, both arrow orientations
are covered. There are 8828 action classes. The exact normalized
category computation excludes 2266 by $b_1(-;\mathbb F_2)$ and
3321 by $b_2(-;\mathbb F_2)$. Of the others, 3226 have an explicit
degree-three primitive for the degree-two generator's square.
The remaining fifteen have third mod-two Betti number one, and
therefore are not $\mathbb CP^2$. The tables, actions, cocycles,
primitives, and complete residual ranks are retained in
\path{code/outputs/c_eight_all_act_cones.json}.

It remains to treat the other four profiles in the displayed list.
For $(6,3,3)$, the relative-cup argument again restricts the larger
monoid to those with nonzero rational $H^2$. Its minimal band is
$2\times2$ with one extra element. The same translation-pair and
square enumeration gives exactly five ordinary isomorphism classes.
For $(5,4,3)$ and $(5,3,4)$, the larger monoid must be the
five-element sphere monoid. In $(5,5,2)$, both monoids must be that
monoid: every other five-element monoid has rationally acyclic
classifying space and is covered by the two-point reduction.

The own three- and four-element monoid censuses, with opposites
restored, supply respectively seven and 35 ordinary isomorphism
classes for the smaller vertex. The program
\path{code/c_twelve_triangular_bisets.py} enumerates every right
action of the larger monoid up to state permutation and \emph{all}
conjugates of every left action of the smaller monoid. It retains
exactly the commuting pairs, checking commutation first on
generators and then on every pair of elements. Thus no relative
position of the two actions is lost. Its complete counts are
\[
\begin{array}{c|r|r|r|r|r}
(m,n,s)&\text{bisets}&b_1\ne0&b_2\ne1&x^2=0&x^2\ne0\\\hline
(6,3,3)&1740&460&890&390&0\\
(5,4,3)&2264&572&1244&448&0\\
(5,3,4)&2102&851&759&491&1\\
(5,5,2)&8&1&1&6&0.
\end{array}
\]
The tests are applied in the indicated order, so the columns are
disjoint and sum to the biset count. The sole nonzero-square
residual has third mod-two Betti number one. Every asserted zero
square has a full cochain primitive checked on all normalized
four-simplices. All four profiles are therefore excluded, which
completes the proof.
\end{proof}

\subsection{A structural twelve-arrow exclusion with arrows in both directions}

\begin{lemma}[The minimal ideal of a finite strongly connected category]
\label{C:finite-category-Rees}
Let $C$ be a nonempty finite category in which every hom-set is
nonempty. There are a finite group $G$, positive integers $r_i,c_i$
for each object $i$, and a least nonempty two-sided ideal $K$ of
arrows, with
\[
 K(i,j)\cong R_j\times G\times C_i,
 \qquad |R_j|=r_j,\quad |C_i|=c_i.
\]
After choosing these coordinates, composition has the form
\[
 (r_k,g,c_j)(r_j,h,c_i)
       =(r_k,g\,p_{j,c_j,r_j}\,h,c_i),\qquad p_{j,c_j,r_j}\in G.
\]
In particular,
\[
 |K|=|G|\Bigl(\sum_i r_i\Bigr)\Bigl(\sum_i c_i\Bigr).
\]
If every $r_i=1$, or every $c_i=1$, then $BC\simeq BG$.
If $1_i\in K(i,i)$, object $i$ is a retract of every object of $C$.
\end{lemma}
\begin{proof}
Fix an object $a$, choose the least nonempty ideal of the finite
monoid $\operatorname{End}(a)$, and choose an idempotent $e$ in
that ideal. Such an idempotent is a sufficiently large power of
any element of the ideal. The monoid
$G=e\operatorname{End}(a)e$ is a group. Indeed, for every
$z\in G$, minimality gives $e=u z v$ for some endomorphisms
$u,v$; after sandwiching, this is
$e=(eue)z(eve)$ in the finite monoid with identity $e$.
In a finite monoid, every factor of a product equal to the
identity is a unit, so $z$ is a unit.

Put
\[
 U_j=C(a,j)e,\qquad V_i=eC(i,a),\qquad K(i,j)=U_jV_i.
\]
These sets are nonempty and the $K(i,j)$ form an ideal.
Every nonempty ideal of $C$ meets $\operatorname{End}(a)$:
compose any one of its arrows with arrows from and to $a$.
Its intersection with that endomorphism monoid consequently
contains $e$, and it contains every $U_jV_i$. Thus $K$ is the
least nonempty ideal.

The right $G$-action on $U_j$ and the left $G$-action on $V_i$
are free. For example, if $u\in U_j$ and $v\in V_j$, then
$vu\in G$, so $ug=u$ implies $(vu)g=vu$ and hence $g=e$.
The other action is analogous. Moreover, multiplication gives
a bijection
\[
 U_j\times_G V_i\longrightarrow K(i,j).
\]
Surjectivity is the definition. If $uv=u'v'$, choose $w\in U_i$
and adjust it by an element of $G$ so that $vw=e$.
Then $u=u'(v'w)$, with $v'w\in G$. Applying a left inverse
of $u'$ to the original equality gives the corresponding
relation between $v$ and $v'$. This is precisely the balanced
$G$-equivalence relation. Choosing orbit representatives now
gives the stated coordinates; the intervening products
$v_{j,c}u_{j,r}\in G$ give the displayed sandwich entries.

Adjoin the retract object represented by $e$, inside the
idempotent completion, and call the resulting category $C^+$.
This does not change the classifying space. To recall the
elementary Quillen-A argument, write
$u:e\to a$, $v:a\to e$ with $vu=1_e$.
In $(e\downarrow C)$, the object $u$ admits a natural
transformation from its constant functor to the identity,
whose component at $f:e\to i$ is $fv:a\to i$.
This contracts the comma category; commas for objects already
in $C$ have initial identity objects.

The one-object full subcategory at $e$ has endomorphism group
$G$. If every $r_i=1$, then every $U_i=C^+(e,i)$ is a free
transitive right $G$-set. The comma category $(G\downarrow i)$
is its translation groupoid and is contractible. Quillen's
Theorem A gives $BG\simeq BC^+\simeq BC$.
If every $c_i=1$, apply the opposite argument.

Finally, if $1_i=uv\in K(i,i)$, then $vu\in G$ is idempotent
and therefore equals $e$. For any object $j$, choose
$u_j\in U_j$, $v_j\in V_j$ and adjust so that $v_ju_j=e$.
The arrows $u_jv:i\to j$ and $uv_j:j\to i$ have composite
$uv_ju_jv=uv=1_i$, proving the retract assertion.
\end{proof}

\begin{lemma}[The rectangular minimal ideal Tor comparison for a category]
\label{C:category-band-Tor}
Suppose the group in Lemma~\ref{C:finite-category-Rees} is
trivial. Put $R=\mathbb QC$ and $A=R/\mathbb QK$, and choose
the minimal-corner idempotent $e$ used in that proof.
Let $\mathbf 1_\ell$ and $\mathbf 1_r$ be the constant left
and right $R$-modules, with one copy of $\mathbb Q$ at every
object and every arrow acting as the identity between its
source and target copies. Define
\[
 V=\ker(Re\longrightarrow\mathbf1_\ell),\qquad
 W=\ker(eR\longrightarrow\mathbf1_r),
\]
where the maps send each arrow basis element to one in its
object component. Then $V,W$ are $A$-modules and, for $q\geq2$,
\[
 H_q(BC;\mathbb Q)
       \cong\operatorname{Tor}_{q-2}^{A}(W,V).
\]
Here $\dim V=\sum_i(r_i-1)$ and
$\dim W=\sum_i(c_i-1)$.
\end{lemma}
\begin{proof}
The module $Re$ has $r_i$ basis vectors at object $i$, and
$eR$ has $c_i$ basis vectors there. Their augmentations are
therefore surjective and have the stated kernel dimensions.
Every arrow in $K$ acts constantly on the corresponding row
or column set, so $K$ annihilates both augmentation kernels.
The coordinate description also gives an $R$-bimodule
isomorphism
\[
 Re\otimes_{\mathbb Q}eR\xrightarrow{\ \mathrm{multiplication}\ }
 \mathbb QK.
\]
In fact its arrow-basis tensors map bijectively to the row-column
basis of $K$. It follows that $\mathbb QK$ is projective on
both sides.

The normalized category bar computes
$H_q(BC;\mathbb Q)=\operatorname{Tor}_q^R(\mathbf1_r,\mathbf1_\ell)$.
Dimension shifting first through $Re$ and then through $eR$
gives
\[
 H_q(BC;\mathbb Q)\cong\operatorname{Tor}_{q-2}^R(W,V)
 \qquad(q\geq2).
\]
At $q=2$, the second shift uses the exact sequence
\[
 0\longrightarrow\operatorname{Tor}_1^R(\mathbf1_r,V)
 \longrightarrow W\otimes_R V\longrightarrow eR\otimes_R V.
\]
The last term is $eV=0$, so this is indeed the same formula
with ordinary tensor product in degree zero.

There is no extra derived change-of-rings term here.
The projective left resolution
$0\to\mathbb QK\to R\to A\to0$ gives
\[
 W\otimes_R\mathbb QK
 \cong We\otimes_{\mathbb Q}eR=0,
\]
and therefore $\operatorname{Tor}_j^R(W,A)=0$ for $j>0$,
while $W\otimes_R A=W$. The change-of-rings spectral sequence
\[
 \operatorname{Tor}_i^A\bigl(\operatorname{Tor}_j^R(W,A),V\bigr)
 \Longrightarrow\operatorname{Tor}_{i+j}^R(W,V)
\]
has only its $j=0$ row. This proves the asserted comparison.
\end{proof}

\begin{proposition}[The nontriangular two-object case through twelve arrows]
\label{C:twelve-strongly-connected}
No category with two objects, nonempty hom-sets in both
directions, and at most twelve arrows has classifying space
homotopy equivalent to $\mathbb CP^2$.
\end{proposition}
\begin{proof}
Suppose otherwise, and use the notation of
Lemma~\ref{C:finite-category-Rees}. If an object identity
belongs to $K$, that object is a retract of the other object.
Deleting it preserves the classifying space by the same comma
contraction used in that lemma, and leaves at most eleven
arrows. This contradicts Proposition~\ref{C:CP2-lower12}.
Thus both identities lie outside $K$.

If all row numbers, or all column numbers, are one, the same
lemma gives $BC\simeq BG$. A finite group has no positive
rational homology, so this cannot be $\mathbb CP^2$.
Otherwise both total row and column numbers are at least
three. The arrow count now gives
\[
 |G|\,(r_a+r_b)(c_a+c_b)+2\leq12.
\]
It forces $G=1$ and
\[
 r_a+r_b=c_a+c_b=3.
\]
The ideal therefore has exactly nine arrows, and there is at
most one additional arrow beyond its arrows and the two
identities. Up to exchanging the objects and taking opposites,
there are just two configurations.

First suppose the doubled row and column occur at the same
object $a$, so $(r_a,r_b)=(c_a,c_b)=(2,1)$.
The minimal hom-cardinalities are
\[
 \begin{pmatrix}4&2\\2&1\end{pmatrix}.
\]
At least one cross hom-set has exactly its two minimal arrows.
Put $M=\operatorname{End}(a)$; its cardinality is at most six.
If $C(b,a)=K(b,a)$, this left $M$-set is isomorphic to $Mf$
for any idempotent $f\in K(a,a)$: both consist of the two
row coordinates with a fixed column. The comma category
$(b\downarrow M)$ is its action category and is contractible.
For clarity, the action category of $M$ on $Mf$ contracts
by the natural transformation from the constant object $f$
to the identity, whose component at $z\in Mf$ is the arrow
labelled $z$, since $zf=z$.
The other comma category has an initial identity object.
Quillen's Theorem A gives $BM\simeq BC$.
If instead $C(a,b)=K(a,b)$, use the right $M$-set $fM$ and
the opposite comma argument. In either case a monoid of at
most six elements would realize $\mathbb CP^2$, again
contradicting Proposition~\ref{C:CP2-lower12}.

Second suppose the doubled row and column occur at different
objects. Choose names so that $(r_a,r_b)=(2,1)$ and
$(c_a,c_b)=(1,2)$. The minimal endomorphism cardinalities
are two at each object, and the cross cardinalities are
$|K(a,b)|=1$, $|K(b,a)|=4$.
If no additional arrow is placed in $C(a,b)$, the singleton
return criterion, Lemma~\ref{C:singleton-return}, makes $BC$
contractible. The only remaining possibility thus has exactly
twelve arrows: the nine arrows of $K$, two identities, and
one extra arrow $t:a\to b$.

In this last case $A=\mathbb QC/\mathbb QK$ has basis
$1_a,1_b,t$ and is precisely the path algebra of one arrow
$a\to b$. It is hereditary. Explicitly, for a representation
$f:X_a\to X_b$, its projective resolution is
\[
 0\longrightarrow P_b\otimes X_a
 \longrightarrow(P_a\otimes X_a)\oplus(P_b\otimes X_b)
 \longrightarrow X\longrightarrow0,
\]
with the maps at vertex $b$ given by
$x\mapsto(x,-f(x))$ and $(x,y)\mapsto f(x)+y$.
Lemma~\ref{C:category-band-Tor} consequently gives
\[
 H_4(BC;\mathbb Q)=\operatorname{Tor}_2^A(W,V)=0.
\]
This is the final contradiction. The proof covers every
composition law satisfying the indicated minimal-ideal
structure; no twelve-arrow composition-table census is used.
\end{proof}

\subsection{Rank-one augmentation modules in a finite category}

\begin{lemma}[Affine representatives for balanced degree-two classes]
\label{C:category-affine-balanced}
Let $C$ be a finite strongly connected category with rectangular-band
minimal ideal.  Use the row and column sets $R_i,C_i$ and the
augmentation modules $V,W$ of Lemma~\ref{C:category-band-Tor}.
Suppose $\dim V=1$.  Then every class in $H^2(BC;\mathbb Q)$ has
an affine cocycle representative of the following form.  There is a
unique object $a$ with $|R_a|=2$, while every other $R_i$ is a
singleton.  Choose functions $a_i:R_i\to\mathbb Q$ and
$b_i:C_i\to\mathbb Q$, both zero away from object $a$, and with
chosen distinguished values zero.  For each arrow $f:i\to j$ there
are scalars $\chi_f,a_f,b_f$ such that
\begin{align*}
 a_j\sigma_f&=\chi_f a_i+a_f,&
 b_i\tau_f&=\chi_f b_j+b_f.
\end{align*}
Here $\sigma_f:R_i\to R_j$ and $\tau_f:C_j\to C_i$ are the
row and column actions, and $\chi_f=0$ unless $i=j=a$.
For a composable pair written in multiplication order,
\[
 u(f,g)=b_fa_g
\]
represents the prescribed degree-two class.
\end{lemma}
\begin{proof}
The comparison of Lemma~\ref{C:category-band-Tor} identifies
degree-two cohomology with the dual of $W\otimes_A V$.
Object idempotents show that a balanced functional is supported on
$W_a\otimes V_a$.  Choose $a_a$ to take the values zero and one
on the two row points.  Any balanced functional is then
$\lambda(w,v)=b_a(w)a_a(v)$ for a linear functional $b_a$ on
the column augmentation space; extend $b_a$ to a function on
$C_a$ with value zero at its distinguished point.

On $V_a$, an endomorphism acts by a scalar $\chi_f$.
Balancing says that its column action has the same scalar on
$b_a$.  For a cross arrow leaving $a$, balancing says that
$b_a\tau_f$ is constant.  For a cross arrow entering $a$, its
row source is a singleton, so $a_a\sigma_f$ is constant.
These observations give exactly the displayed affine equations,
with $\chi_f=0$ on all other arrows.  The maps on the one-dimensional
module $V$ show that $\chi_{fg}=\chi_f\chi_g$ whenever $fg$ is
defined.  Composition of the affine maps gives
\begin{equation}\label{C:category-affine-composition}
 a_{fg}=a_f+\chi_f a_g,\qquad
 b_{fg}=b_g+\chi_g b_f.
\end{equation}
Thus $u$ is a normalized category two-cocycle, by substitution into
$\delta u(f,g,h)=u(g,h)-u(fg,h)+u(f,gh)-u(f,g)$.

We verify the asserted identification with the balanced functional,
rather than only the cocycle equation.  Let $o$ be the object of a
chosen minimal idempotent $e$.  Choose row arrows
$u_r:o\to i$ and column arrows $v_c:i\to o$ as in
Lemma~\ref{C:finite-category-Rees}.  Since the corner group is
trivial, $v_cu_r=e$ for every such pair.  For augmentation vectors
$w=\sum_c w_cv_c$ and $v=\sum_r v_ru_r$ at object $i$, set
\[
 z(w,v)=\sum_{c,r}w_cv_r[\,v_c\mid u_r\,].
\]
This is a two-cycle: in
$\partial[v_c\mid u_r]=[u_r]-[e]+[v_c]$, every summand vanishes
after one of the two zero coefficient sums is taken.  For
$f:i\to j$, the boundary of
$\sum_{c,r}w_cv_r[v_c\mid f\mid u_r]$ is
$z(w,fv)-z(wf,v)$, so these cycles respect the balancing relation.
They are the inverse of the two augmentation connecting maps in
Lemma~\ref{C:category-band-Tor}: lift their final coefficient at
$o$ to $e\in Re$; the first boundary is the bar one-cycle with
coefficients $v$ and column entries $v_c$; lifting its other
coefficient to $e\in eR$ gives boundary $w\otimes v$.
In particular they represent every degree-two homology class.
Finally the minimal-ideal actions are constant, and hence
$a_{u_r}=a_i(r)$ and $b_{v_c}=b_i(c)$.  Therefore
\[
 \langle u,z(w,v)\rangle
 =\sum_{c,r}w_cv_r b_i(c)a_i(r)=\lambda(w,v).
\]
This proves the claimed cohomology representative.
\end{proof}

\begin{proposition}[The category rank-one cup obstruction]
\label{C:category-rank-one-cups}
Under the hypotheses of Lemma~\ref{C:category-band-Tor}, if
$\dim V\leq1$ or $\dim W\leq1$, then every product of two
rational degree-two classes is zero.  Thus a nonzero such product
requires
\[
 \sum_i r_i\geq |\operatorname{Ob}C|+2,
 \qquad
 \sum_i c_i\geq |\operatorname{Ob}C|+2.
\]
No finite-homotopy hypothesis is needed.
\end{proposition}
\begin{proof}
Dimension zero gives $H^2=0$.  Taking the opposite category
exchanges $V,W$ and preserves the assertion, so assume $\dim V=1$
and use Lemma~\ref{C:category-affine-balanced}.  Let $e_i$ be one
at object $a$ and zero elsewhere.  The offsets $a_f,b_f$ are
supported, respectively, on arrows with target $a$ and source $a$.
The zero class needs no argument, so suppose that the balanced
functional is nonzero; both scalar value sets are then nonconstant.
Whenever $\chi_f\ne0$, the induced affine maps permute the finite
nonconstant value sets of $a_a$ and $b_a$.  Their slopes are
$1$ or $-1$.  A slope-one permutation has zero translation, and
all slope-minus-one permutations of one finite value set have
its same midpoint.  Choose rational midpoints $A,B$ for the two
sets, using zero if no negative-slope map occurs, and put
\[
 a'_f=a_f+A\chi_f-Ae_{t(f)},\qquad
 b'_f=b_f+B\chi_f-Be_{s(f)}.
\]
Then $\chi_fa'_f=\chi_fb'_f=0$ for every arrow.  The modified
offsets still satisfy \eqref{C:category-affine-composition},
because $\chi_f$ is supported on endomorphisms of $a$.

Set $u'(f,g)=b'_fa'_g$ and
$h(f)=(e_{s(f)}+e_{t(f)})/2$.  For a composable pair,
$\delta h(f,g)=e_{s(f)}=e_{t(g)}$.  The support of the original
offsets and $\chi_f$ gives the exact identity
\[
 u'-u=-\delta\bigl(A b+B a+AB(\chi-h)\bigr).
\]
The cochain in parentheses vanishes on every identity arrow.
Finally set
\[
 w(f,g,h)=b'_fa'_gb'_ga'_h.
\]
Expanding its coboundary with the affine composition equations
gives, on every composable quadruple,
\[
 \delta w(f,g,h,k)
 =u'(f,g)u'(h,k)
   +b'_fa'_k\,\chi_g b'_g\,\chi_h a'_h
 =u'(f,g)u'(h,k).
\]
Every degree-two square is therefore zero.  Over $\mathbb Q$,
graded commutativity and the square of a sum show that every
product of two degree-two classes is zero as well.
\end{proof}

\begin{lemma}[Rational Tor comparison with a finite minimal group]
\label{C:category-group-Tor}
The Tor formula of Lemma~\ref{C:category-band-Tor} remains valid
when the minimal group $G$ is any finite group, if $V,W$ are
defined by the orbit augmentations
\[
 P=Re\otimes_{\mathbb QG}\mathbb Q\longrightarrow\mathbf1_\ell,
 \qquad
 Q=\mathbb Q\otimes_{\mathbb QG}eR\longrightarrow\mathbf1_r.
\]
Their dimensions are still $\sum_i(r_i-1)$ and
$\sum_i(c_i-1)$, and $A=R/\mathbb QK$.
\end{lemma}
\begin{proof}
The trivial $\mathbb QG$-module is projective, by the averaging
idempotent $|G|^{-1}\sum_{g\in G}g$.  Thus $P,Q$ are projective
left and right $R$-modules, respectively.  Their bases are the
row and column orbits, and their augmentation kernels are
annihilated by $K$.  The Rees multiplication bijection in
Lemma~\ref{C:finite-category-Rees} gives
\[
 Re\otimes_{\mathbb QG}eR\cong\mathbb QK.
\]
The row and column $G$-actions are free, so this ideal is
projective on each side.  Also $eV=We=0$ and
$Q\otimes_RV=\mathbb Q\otimes_{\mathbb QG}eV=0$.
The two augmentation sequences therefore give
$H_q(BC;\mathbb Q)\cong\operatorname{Tor}_{q-2}^{R}(W,V)$,
including $q=2$.  Finally
$W\otimes_R\mathbb QK=We\otimes_{\mathbb QG}eR=0$;
the projective resolution $0\to\mathbb QK\to R\to A\to0$
and the same change-of-rings argument as before replace $R$
by $A$.
\end{proof}

\begin{theorem}[The strongly connected two-object exclusion through twenty]
\label{C:twenty-strongly-connected}
A finite category with two objects, nonempty hom-sets in both
directions, and at most twenty arrows cannot have classifying
space homotopy equivalent to $\mathbb{CP}^2$.
\end{theorem}
\begin{proof}
If either object identity is in the minimal ideal, the retract
deletion in Lemma~\ref{C:finite-category-Rees} leaves an
equivalent one-object category with at most nineteen arrows.
This contradicts Corollary~\ref{B:CP2-lower21}.  Thus both
identities are outside $K$.

If all row numbers or all column numbers are one, the same
Rees lemma gives $BC\simeq BG$, which has no positive rational
homology.  Otherwise both total row and column numbers are at
least three.  If $G\ne1$, the arrow bound
\[
 |G|\Bigl(\sum_i r_i\Bigr)\Bigl(\sum_i c_i\Bigr)+2\leq20
\]
forces $|G|=2$, both sums equal to three, and exactly two arrows
outside $K$, namely the identities.  Then $A\cong\mathbb Q^2$
is semisimple, so Lemma~\ref{C:category-group-Tor} gives
$H_4(BC;\mathbb Q)=0$.

It remains to treat $G=1$.  The nonzero degree-two square of
$\mathbb{CP}^2$ and Proposition~\ref{C:category-rank-one-cups}
force both row and column sums to be at least four.  Thus
$|K|\geq16$ and $\dim A\leq4$.  The two object identities
survive as nonzero orthogonal idempotents in $A$, so $A$ is
nonlocal.  By Lemma~\ref{B:nonlocal-four-Tor}, a nonzero
$\operatorname{Tor}_2^A(W,V)$ forces nonzero Tor in arbitrarily
high degrees.  Lemma~\ref{C:category-band-Tor} would therefore
make nonzero $H_4$ incompatible with eventual vanishing of
$H_q(BC;\mathbb Q)$.  This is the required contradiction.
\end{proof}

\subsection{Two minimal rows at each object and a stronger two-object bound}

\begin{proposition}\label{C:objectwise-two-row-cups}
Let $C$ be a finite strongly connected category with a
rectangular-band minimal ideal.  If each object has at most two
minimal rows, every product of rational degree-two classes on
$BC$ is zero.  The same statement holds if each object has at
most two minimal columns.  No finite-homotopy hypothesis is needed.
\end{proposition}
\begin{proof}
Use the row and column augmentation modules $V,W$ of
Lemma~\ref{C:category-band-Tor}.  A balanced functional decomposes
objectwise as $\lambda_i:W_i\otimes V_i\to\mathbb Q$.
Since $\dim V_i\leq1$, choose a row coordinate $a_i$ with values
zero and one when the row set has two elements, and value zero
when it has one.  Write $\lambda_i=b_i\otimes a_i$, normalizing
$b_i$ to be zero at a distinguished column point.  At a one-row
object choose $b_i=0$.

For $f:i\to j$, let $\chi_f$ be its row augmentation scalar,
zero if either row module is zero.  The row action and balancing
give affine equations
\[
 a_j\sigma_f=\chi_f a_i+a_f,
 \qquad b_i\tau_f=\chi_f b_j+b_f.
\]
These also cover the zero-dimensional endpoint cases.  If the
source has two rows and the target has one, balancing says that
$b_i\tau_f$ is constant; if the source has one row, $b_i=0$
makes the second equation automatic.  Composition gives
\[
 \chi_{fg}=\chi_f\chi_g,\qquad
 a_{fg}=a_f+\chi_f a_g,\qquad
 b_{fg}=b_g+\chi_g b_f.
\]
The band-grid cycles and their balancing three-chains in the proof
of Lemma~\ref{C:category-affine-balanced} apply separately at
each object.  They show that $u(f,g)=b_fa_g$ represents the
prescribed functional, and hence that every rational degree-two
class has this form.  That comparison does not require the
sum of the dimensions of the $V_i$ to be one.

Put $A_i=1/2$ at each two-row object and $A_i=0$ at each
one-row object.  Replace the offsets by
\[
 a'_f=a_f+\chi_f A_{s(f)}-A_{t(f)}.
\]
If $\chi_f\ne0$, the row map is a bijection between two-point
sets, so it preserves their midpoints and $a'_f=0$.
Consequently $\chi_fa'_f=0$ for every arrow.  The new offsets
still satisfy the typed affine composition equation.
For the normalized one-cochain $v(f)=A_{s(f)}b_f$ one has
\[
 \delta v(f,g)=b_f\bigl(A_{t(g)}-\chi_g A_{s(g)}\bigr),
 \qquad u'-u=-\delta v,
 \qquad u'(f,g)=b_fa'_g.
\]
The normalized three-cochain
$w(f,g,h)=b_fa'_g b_ga'_h$ satisfies
\[
 \delta w(f,g,h,k)
 =u'(f,g)u'(h,k)
  +b_fa'_k\,\chi_gb_g\,\chi_ha'_h
 =u'(f,g)u'(h,k).
\]
All equations are on composable strings, with
$s(f)=t(g)$ making the midpoint indices agree.
Every square therefore vanishes, and polarization over
$\mathbb Q$ gives every mixed product.  Passing to the opposite
category proves the assertion for columns.
\end{proof}

\begin{proposition}\label{C:one-active-three-by-three}
Let $C$ be a finite category with two objects $A,B$ and a
rectangular-band minimal ideal.  Suppose its row counts are
$(3,1)$ and its column counts are $(3,1)$, with both triples at
$A$.  If $BC$ has finite CW homotopy type, every product of two
rational degree-two classes on $BC$ is zero.
\end{proposition}
\begin{proof}
Let $H$ be the finite transformation-pair image of
$\operatorname{End}_C(A)$ on its three row and three column
points.  It contains the complete constant band
$\{x_{ij}:0\leq i,j\leq2\}$.  Put $e=x_{00}$ and let
$D$ be the full subcategory of $\operatorname{Kar}(H)$ on
objects $1,e$.
There is an explicit functor $q:C\to D$.
Send $A$ to $1$, $B$ to $e$, endomorphisms of $A$ to their
transformation pairs, and every endomorphism of $B$ to $e$.
An arrow $A\to B$ has the unique singleton row map and selects
a column $j$ at $A$; send it to $x_{0j}\in eH$.
An arrow $B\to A$ selects a row $i$ at $A$ and has the unique
singleton column map; send it to $x_{i0}\in He$.
The action equations and $x_{ij}x_{kl}=x_{il}$ check every
composable pair.  In particular, a composite through $A$ from
$B$ to itself maps to $x_{00}=1_e$, so identities are preserved.

The object $e$ is a retract of $1$, so $BH\to BD$ is a
classifying-space equivalence by the retract-generator lemma.
Moreover $q$ induces an isomorphism on $H_2(-;\mathbb Q)$.
Indeed $V_B=W_B=0$.  In the natural balanced tensor presentation
only $W_A\otimes V_A$ remains, and its only nonzero balancing
relations come from $\operatorname{End}_C(A)$; every cross-arrow
relation has a zero factor.  The presentation is therefore the
same as the one for $H$.  The band-grid cycles identify this
isomorphism with the actual functor-induced map.

Apply the complete three-by-three action-hull classification and
its exact certificates from the proof of
Theorem~\ref{B:three-by-three-cups}.  The maximal hull has exactly
the same balanced relation space as $H$, so every degree-two
class under consideration pulls back from that hull through
its two-object Karoubi category.  The rank-one, zero-product,
and fourth-homology-zero cases of that classification immediately
give the required vanishing after pullback.

We spell out the remaining skew-hull case, because finite type
is required only for the original $BC$ here.
The integral affine cocycle of that thirty-element hull extends
to the cohomology of its two-object Karoubi category through the
classifying-space equivalence.  Pull it back along $q$ to an
integral class on $BC$.
If one of its eight forbidden actors were the image of an element
of $\operatorname{End}_C(A)$, the eventual cyclic group of that
element would surject onto the actor's eventual cyclic group.
The nonzero integral character computed in
\eqref{B:3x3-cyclic-character} remains nonzero under this
surjection.  It is the restriction of the just-constructed class
on $BC$.  On the other hand, the local finite group maps to $C$
by adjoining a fresh identity if necessary, and Miller's theorem
makes the induced map to a finite CW model of $BC$ null-homotopic.
This is a contradiction.

Thus the actual image avoids all eight forbidden actors.  As
proved and independently checked in
Theorem~\ref{B:three-by-three-cups}, it lies in one of the nine
safe submonoids, each with zero fourth rational homology.  The
functor $q$ then factors through the corresponding Karoubi
category, and the pulled-back square vanishes.  Polarization
finishes the proof.  No finite-homotopy assertion about
$B\operatorname{End}_C(A)$, $BH$, or any quotient is used.
\end{proof}

\begin{theorem}[The strongly connected two-object exclusion through twenty-two]
\label{C:twentytwo-strongly-connected}
A finite category with two objects, nonempty hom-sets in both
directions, and at most twenty-two arrows cannot have classifying
space homotopy equivalent to $\mathbb{CP}^2$.
\end{theorem}
\begin{proof}
If either object identity lies in the minimal ideal, the retract
deletion in Lemma~\ref{C:finite-category-Rees} gives an equivalent
one-object category of smaller order.  This contradicts the monoid
minimum twenty-three in Theorem~\ref{B:CP2-minimum23}.  Thus both
identities survive outside the minimal ideal $K$.
Write $R=\sum_i r_i$ and $S=\sum_i c_i$ for its total row and
column counts.  If either augmentation module has dimension zero,
the band--Tor comparison gives $H_2=0$.

Suppose first that the minimal group $G$ is nontrivial.  Nonzero
$H_2$ requires $R,S\geq3$, and the arrow bound forces
$|G|=2$, $R=S=3$, and $|K|=18$.
The contracted category algebra has dimension at most four and
contains the two nonzero orthogonal object idempotents; hence it
is nonlocal.  Lemma~\ref{B:nonlocal-four-Tor} and the rational
comparison of Lemma~\ref{C:category-group-Tor} make a nonzero
fourth homology group incompatible with eventual homology
vanishing.  This excludes the group case.

Now $K$ is a band.  The existing rank-one obstruction in
Proposition~\ref{C:category-rank-one-cups} requires $R,S\geq4$.
If $RS\geq20$, at most two arrows lie outside $K$; they must be
exactly the two identities.  The contracted algebra is then
$\mathbb Q^2$, and the band--Tor formula gives $H_4=0$.
It follows that $R=S=4$.
Proposition~\ref{C:objectwise-two-row-cups} excludes a distribution
$(2,2)$ on either side.  Thus the row and column distributions
are both $(3,1)$, up to exchanging objects.
If the triples belong to different objects, each tensor
$W_i\otimes V_i$ is zero and $H_2=0$.
If they belong to the same object,
Proposition~\ref{C:one-active-three-by-three} gives zero
degree-two products.  Both conclusions contradict
$\mathbb{CP}^2$, completing the proof.
\end{proof}

\subsection{The minimum within strongly connected categories}

\begin{proposition}\label{C:strong-three-objects-small-cups}
Let $C$ be a finite strongly connected category with at least
three objects and at most twenty-four arrows. Then every product
of two rational degree-two classes on $BC$ is zero.
No finite-homotopy hypothesis is imposed.
\end{proposition}
\begin{proof}
Let $k$ be the number of objects. In the minimal Rees ideal write
$r_i,c_i$ for the positive row and column counts at object $i$,
and $G$ for its finite maximal group. Put
$R=\sum_i r_i$ and $S=\sum_i c_i$. The minimal ideal has exactly
$|G|RS$ arrows, by Lemma~\ref{C:finite-category-Rees}.

If $G$ is nontrivial, nonzero rational second homology requires
both augmentation modules in
Lemma~\ref{C:category-group-Tor} to be nonzero. Their dimensions
are $R-k$ and $S-k$, respectively. Thus $R,S\geq k+1$, and
\[
 |\operatorname{Mor}C|\geq |G|RS
 \geq2(k+1)^2\geq32,
\]
which contradicts the arrow bound. In this case $H^2(BC;\Q)=0$.

If $G$ is trivial, its minimal ideal is a rectangular band.
By Proposition~\ref{C:objectwise-two-row-cups}, a nonzero
degree-two product requires at least one object with three or
more rows, and at least one object with three or more columns.
The two objects need not be the same. Since every other row
and column count is positive, $R,S\geq k+2$. Consequently
\[
 |\operatorname{Mor}C|\geq RS\geq(k+2)^2\geq25,
\]
again a contradiction. These two alternatives exhaust the
minimal group, proving the proposition.
\end{proof}

\begin{corollary}\label{C:strong-category-CP2-minimum23}
Among finite strongly connected categories whose classifying
spaces have homotopy type $\CP^2$, the least number of arrows,
identities counted, is twenty-three.
\end{corollary}
\begin{proof}
For one object this is Theorem~\ref{B:CP2-minimum23}.
Two objects and at most twenty-two arrows are excluded by
Theorem~\ref{C:twentytwo-strongly-connected}. Three or more
objects are excluded by
Proposition~\ref{C:strong-three-objects-small-cups}, since a
projective plane has a rational degree-two class with nonzero
square. The twenty-three-element monoid of
Theorem~\ref{B:CP2-minimum23}, regarded as a one-object category,
attains the bound. This is a minimum within strongly connected
categories; general finite categories remain a separate question.
\end{proof}

% Root proof; independent literal audit by D: PASS. Preserved draft in notes.
\subsection{Removing large strong components from minimal candidates}

\begin{lemma}\label{C:extreme-trivial-Rees-block}
Let $C$ be a finite category partitioned into a full strongly
connected source block $S$ and a full complementary category $D$,
with no arrows from $D$ to $S$. Suppose that the minimal ideal
of $S$ has trivial maximal group and exactly one row and one
column at every object. There is an actual functor from $C$ to
the category obtained by replacing $S$ by one trivial object and
each incident right $S$-act by its orbit set. This functor induces
a homotopy equivalence of classifying spaces. The dual statement
holds for a sink block.
\end{lemma}
\begin{proof}
Write $k_{ij}:i\to j$ for the unique minimal-ideal arrow with
these endpoints. The ideal property and uniqueness give
$f k_{hi}=k_{hj}$ and $k_{j\ell}f=k_{i\ell}$ for every
$f:i\to j$ in $S$. In particular the arrows $k_{ij}$ form a
system of compatible zero arrows, and $BS$ is contractible:
the arrows $k_{ia}:i\to a$ give a natural transformation from
the identity functor to the constant functor at a chosen object $a$.

For a finite right $S$-act $F:S^{\mathrm{op}}\to\mathrm{Set}$,
use the category of elements with arrows
$(i,x)\to(j,y)$ represented by $f:i\to j$ with $x=yf$.
Put $e=k_{aa}$ and $T=F(a)e$. Define
\[
 p(i,x)=xk_{ai}\in T,
 \qquad s(t)=(a,t).
\]
Every arrow has the same value under $p$, since
$f k_{ai}=k_{aj}$. Also $ps=\mathrm{id}_T$, and
$k_{ai}$ gives an arrow $s p(i,x)\to(i,x)$ in the element
category. The displayed zero-arrow equations make these arrows
a natural transformation $sp\Rightarrow\mathrm{id}$.
Thus $p$ is a homotopy equivalence onto the discrete set $T$,
which is canonically the orbit set of $F$. The construction is
natural in $F$, including every action of arrows of $D$.

The two-block correspondence description of $BC$ is a homotopy
pushout of the source and target block spaces over their two-sided
category of cross arrows. Applying the preceding natural
equivalence objectwise in $D$ gives a homotopy equivalence of
these cross-arrow spaces to the correspondence for the orbit
category. The source block map $BS\to *$ is an equivalence,
and the target block is unchanged. Homotopy invariance of the
pushout gives the assertion. Composition in the quotient is the
actual induced $D$-action on orbits; there is no path leaving
and then reentering $S$, so no additional composition needs
definition. Opposition proves the dual statement.
\end{proof}

\begin{proposition}\label{C:large-SCC-minimal-candidate21}
Let $C$ be a connected finite category with at most twenty-one
arrows. Assume no object is a retract of another object.
If $C$ has a strongly connected component containing at least
three objects, then $BC$ is not homotopy equivalent to $\CP^2$.
Consequently, a category of type $\CP^2$ with the fewest arrows
and at most twenty-one arrows can have only one- or two-object
strongly connected components.
\end{proposition}
\begin{proof}
A component equal to the whole category is excluded by
Proposition~\ref{C:strong-three-objects-small-cups}.
Let $S$ be a proper component with $t\geq3$ objects.
Because no object is a retract of another, every object identity
lies outside the minimal ideal of $S$. The same adjustment by a
minimal-group inverse as in
Theorem~\ref{C:sixteen-three-SCC} proves this assertion.
Its minimal ideal has at least $t^2$ arrows, so $|S|\geq t^2+t$.
An incident arrow forces an arrow to or from every object of $S$.
The proper complementary category therefore contributes at least
one identity and $t$ cross arrows. For $t\geq4$ this already
requires $t^2+2t+1\geq25$ arrows. Thus $t=3$ and $|S|\geq12$.

If $k$ is the total number of objects, contract $S$ in the
undirected support. Connectivity requires at least $k-3$ support
edges outside $S$, one of which expands to three arrows. With
the remaining identities, this gives
\[
 |C|\geq12+(k-3)+(k-1)=2k+8.
\]
Hence there are at most three objects outside $S$.

If $S$ has both an incoming and an outgoing arrow in the
condensation order, take outside objects $a,b$ with
$a\to S\to b$. They are distinct, since $S$ is a full strongly
connected component. Transitivity supplies three arrows in each
interface and one arrow $a\to b$. Together with $|S|\geq12$
and the two outside identities this costs twenty-one arrows.
Equality is forced. There are no further objects, the outside
corners are trivial and every displayed cross-hom is a singleton.
The object $a$ is initial, so this case is contractible.
We may therefore use opposition to assume that $S$ is a source
block. Write $D=C\setminus S$ and let $T$ be the total number
of cross arrows from $S$ to $D$.

One has $T\geq3$ and $|S|\leq17$. The rational second homology
of $S$ is zero: the categorical Rees Tor formula would otherwise
require total row and column counts at least four, hence at least
sixteen minimal arrows and the three outside identities, giving
$|S|\geq19$. If $H^2(BD;\Q)=0$ as well,
Lemma~\ref{C:two-block-zero-H2-products} makes all degree-two
products of $BC$ zero. Suppose henceforth that $H^2(BD;\Q)\ne0$.

Each connected component of $D$ receives at least one interface
from $S$, hence at least three cross arrows. If $D$ is disconnected,
the component with nonzero second cohomology needs at least five
arrows, by the complete lower-arrow census; another component
needs at least its identity. Thus
$|S|+|D|+T\geq12+6+6=24$, impossible. Therefore $D$ is connected.

If $D$ has three objects, nonzero second cohomology requires at
least eight arrows by Lemma~\ref{C:three-small-H2-zero}; then
$12+8+3>21$. If it has one object with monoid $M$, nonzero
second cohomology requires $|M|\geq5$, whence $|S|\leq13$.
The minimal ideal of $S$ must consequently have trivial group
and total row and column counts both three: the next possible
band size twelve, together with three outside identities, would
already exceed thirteen. It has one row and column per object.
Lemma~\ref{C:extreme-trivial-Rees-block} replaces $S$ by a point
and gives an actual homotopy-equivalent ordinary act cone over
$M$. Its arrow count is at most
\[
 |M|+T+1\leq |C|-|S|+1\leq10.
\]
This contradicts the exact act-cone minimum twenty-two.

It remains that $D$ has two objects. Nonzero second cohomology
requires at least six arrows by
Lemma~\ref{C:two-object-five-H2-zero}. Thus the budget forces
$|S|=12$, $|D|=6$ and $T=3$. As before $S$ has one minimal
row and column per object and trivial minimal group. The same
actual collapse gives a category with three objects and at most
$1+6+3=10$ arrows. Its classifying space is still that of $C$,
contrary to Theorem~\ref{C:CP2-lower16}.
Every case has now been excluded.

Finally a minimal-arrow model cannot contain a retract object:
deleting it preserves the classifying-space type and strictly
decreases the arrow count. The last assertion follows.
\end{proof}

\subsection{Structural reductions for three objects and twelve arrows}

\begin{lemma}[A singleton cross row or column is a join]
\label{C:singleton-cross-join}
Let a category have objects $0<1<2$, with cross arrows only
in the increasing directions. If both $C(0,1)$ and $C(0,2)$
are singletons, then
\[
 C=\operatorname{End}(0)\star C|_{\{1,2\}},
\]
where $\star$ is the categorical join. Dually, singleton
$C(0,2)$ and $C(1,2)$ give
$C=C|_{\{0,1\}}\star\operatorname{End}(2)$.
In either case every positive-degree integral cohomology
product of $BC$ vanishes.
\end{lemma}
\begin{proof}
The categorical join retains the two full subcategories and
has exactly one arrow from every object of the first to every
object of the second. Its cross compositions are forced.
These are exactly the given hom-sets and compositions, including
the actions of an arbitrary extreme endomorphism monoid.
The nerve of a categorical join realizes to the topological
join of the two classifying spaces. Both spaces are nonempty
and connected here, so their join is homotopy equivalent to
the suspension of their smash product. Positive-degree products
on a suspension vanish. The opposite case is identical.
\end{proof}

\begin{lemma}[A rational consequence of the all-act obstruction]
\label{C:eight-bounded-cups}
Consider a triangular two-object category with one endomorphism
monoid $M$ of order at most eight and another monoid $N$ whose
trivial module is projective over the rational algebra on the
side needed for its coinvariants. If its rational homology
vanishes in all sufficiently high degrees, every product of
two degree-two rational cohomology classes is zero.
\end{lemma}
\begin{proof}
Projectivity of the trivial module makes $BN$ rationally
acyclic and its coinvariant functor exact. If
$H^2(BM;\mathbb Q)=0$, the relative-suspension argument of
Lemma~\ref{C:small-triangular-vertices} proves the assertion.
Otherwise the complete proof of
Proposition~\ref{C:eight-act-cones}, applied to the exact
coinvariant augmentation kernel in
Lemma~\ref{C:act-cone-Tor}, says that nonzero fourth rational
homology would persist in arbitrarily high degrees. Boundedness
therefore forces $H_4=0$, which also forces every displayed
degree-two product to be zero. This argument uses the proved
Tor criterion, rather than an inference from the homotopy-type
statement alone.
\end{proof}

\begin{lemma}[A singleton leg in a star]
\label{C:star-singleton-join}
Let a three-object category have cross arrows only on two
legs of a star, both directed toward the center or both away
from it. If one cross-hom set is a singleton and the opposite
leaf has rationally acyclic classifying space, every positive
degree rational cohomology product of the category vanishes.
The endomorphism monoid at the singleton leaf may be arbitrary.
\end{lemma}
\begin{proof}
Write $X$ for the center, $Y$ for the singleton leaf, and $Z$
for the other leaf, using the same letters for their
classifying spaces. The singleton collage is $J=X*Y$,
and the inclusion $X\to J$ is null-homotopic. Write $D$ for
the other collage and $E$ for its two-sided edge bar. Since
there are no composable pairs of cross arrows, the classifying
space of the star is exactly the homotopy pushout of the two
collages along $X$. Nullity of the first inclusion gives
\[
 BC\simeq J\vee\operatorname{hocofib}(X\to D)
       \simeq J\vee\operatorname{hocofib}(E\to Z).
\]
The join of the connected based spaces $X,Y$ is homotopy
equivalent to a suspension. Since $Z$ is rationally acyclic,
its map to a point induces a rational homology equivalence
from the second summand to the unreduced suspension of $E$.
These maps induce cohomology ring isomorphisms over $\mathbb Q$.
Both summands consequently have zero positive-degree products,
as does their wedge.
\end{proof}

\begin{lemma}[All three-object stars through twelve arrows]
\label{C:twelve-all-stars}
No three-object star category with at most twelve arrows has
classifying space homotopy equivalent to $\mathbb CP^2$.
\end{lemma}
\begin{proof}
Let $m,n,p$ be the three endomorphism orders and let $s,t\geq1$
be the leg cardinalities. Thus $m+n+p+s+t\leq12$.
Every monoid of order at most four has rationally acyclic
classifying space, by Theorem~\ref{E:all-six-types}.
In particular, a vertex with nonzero rational $H^2$ has
endomorphism order at least five. There can be at most one
such vertex, since two orders at least five, a third identity,
and two legs would already require thirteen arrows.

If all three vertex spaces have zero rational $H^2$, every
degree-two class lifts to cohomology relative to their union.
After collapsing that union, the graph-of-spaces description
is a wedge of suspensions of the two edge bars with disjoint
basepoints. Relative degree-two products are therefore zero,
and so are the absolute products. This excludes $\mathbb CP^2$.

Suppose first that the unique vertex with nonzero rational
$H^2$ is the center, of order $m\geq5$.
The leaves have order at most four and are rationally acyclic.
A singleton leg is handled by
Lemma~\ref{C:star-singleton-join}. Otherwise $s,t\geq2$ and
the budget gives $n+p\leq3$. Both leaf algebras are therefore
semisimple over $\mathbb Q$, since their monoids have order
one or two. Take their exact coinvariants in the two edge
bars. If the resulting right $\mathbb QM$-modules are $T_1,T_2$,
the high-degree part of the total star chain complex is the
mapping cone of
\[
 (T_1\oplus T_2)\otimes_{\mathbb QM}P_\bullet
 \longrightarrow\mathbb Q\otimes_{\mathbb QM}P_\bullet,
\]
where $P_\bullet$ is a projective resolution of the trivial
module and the map is the sum of the two augmentations.
The two leaf copies of $\mathbb Q$ affect only degrees zero
and one. Thus, for $q\geq2$,
\[
 H_q(BC;\mathbb Q)
 \cong\operatorname{Tor}_{q-1}^{\mathbb QM}(L,\mathbb Q),
 \qquad L=\ker(T_1\oplus T_2\to\mathbb Q).
\]
Here $m\leq8$ and $L$ is finite-dimensional. The same
all-module Tor criterion used in
Proposition~\ref{C:eight-act-cones} says that nonzero fourth
homology persists arbitrarily high. This contradicts the
homology of $\mathbb CP^2$.

Finally suppose the distinguished vertex is a leaf, of order
$m\geq5$. The center and the other leaf are rationally acyclic.
Collapsing the center therefore induces a rational cohomology
ring isomorphism and splits the star into a wedge of its two
leg cofibers. Each leg cofiber is also rationally equivalent,
by a cohomology ring isomorphism, to its original collage.
The collage between the two acyclic vertices has zero
positive-degree products. A nonzero degree-two square must
therefore occur in the collage containing the distinguished
leaf, and that collage has bounded rational homology if the
whole star does.

Let $s$ be this distinguished leg's cardinality and $n$ the
center's endomorphism order. If $s\leq2$,
Lemma~\ref{C:two-object-cup-pruning} gives zero products.
If $s\geq3$, the other leaf and other leg contribute at least
two arrows, so
\[
 m+n\leq12-s-2\leq7.
\]
Hence $n\leq2$, its rational algebra is semisimple, and the
distinguished leaf has order at most eight.
Lemma~\ref{C:eight-bounded-cups} again gives zero degree-two
products. This is the final contradiction.
\end{proof}

\begin{proposition}[Only the linear three-object condensation remains]
\label{C:twelve-three-condensation}
Let $C$ be a three-object category with at most twelve arrows
whose classifying space is homotopy equivalent to $\mathbb CP^2$.
After labelling its objects, it must have cross arrows exactly
in the three directions
\[
 0\longrightarrow1,\qquad1\longrightarrow2,\qquad0\longrightarrow2,
\]
with all three displayed hom-sets nonempty and no reverse arrows.
\end{proposition}
\begin{proof}
Any object that is a retract of another may be deleted without
changing the classifying space. That would leave at most eleven
arrows, contradicting Proposition~\ref{C:CP2-lower12}.
We may therefore assume there are no such retracts.

If all three objects form one strongly connected component,
apply Lemma~\ref{C:finite-category-Rees}. None of their
identities belongs to its minimal ideal, by the retract
assertion. If neither all row numbers nor all column numbers
are one, both total numbers are at least four, so the ideal
alone contains at least sixteen arrows. Otherwise the
classifying space has finite-group type and is rationally
acyclic. Both alternatives exclude $\mathbb CP^2$.

Suppose the strongly connected components have sizes two
and one. The two-object block has at least six arrows:
both endomorphism monoids have order at least two, since a
trivial one would make that object a retract, and there are
arrows in both directions. Connectedness supplies a cross
arrow between this block and the singleton. Composing with
arrows inside the block supplies cross arrows for both block
objects, so there are at least two such arrows.
The singleton's endomorphism order is therefore at most four,
and its classifying space is rationally acyclic.

Conversely, the two-object block has at most nine arrows.
By Lemma~\ref{C:finite-category-Rees}, a block without a
retract that is not of finite-group type needs total row and
column numbers at least three, nine minimal-ideal arrows,
and its two identities outside that ideal. It would therefore
need at least eleven arrows. Our block is consequently of
finite-group type and is rationally acyclic as well.

The mixed-nerve bar decomposition now writes $BC$ as a
homotopy pushout of the two diagonal block spaces over their
cross-bimodule bar. Replacing the two rationally acyclic
vertex spaces by points gives a rational homology equivalence
to a suspension, by the homotopy-pushout exact sequence.
It induces a rational cohomology ring isomorphism, so all
positive products vanish, again a contradiction.

The three strongly connected components are therefore
singletons. Their connected condensation poset is either a
two-leg star or a linear order. The star is excluded by
Lemma~\ref{C:twelve-all-stars}. In a linear order, composition
of the two adjacent nonempty hom-sets makes the third
displayed hom-set nonempty. This is exactly the asserted
remaining configuration.
\end{proof}

\subsection{The complete exclusion through twelve arrows}

\begin{theorem}\label{C:CP2-lower13}
Every finite category whose classifying space has the homotopy
type of $\mathbb CP^2$ has at least thirteen arrows.
Consequently the category minimum for $\mathbb CP^2$ lies in
the interval $[13,40]$ with the explicit upper model below.
\end{theorem}
\begin{proof}
Proposition~\ref{C:CP2-lower12} already excludes at most eleven
arrows. We therefore consider exactly twelve. A connected
category with $k$ objects has its $k$ identities and at least
$k-1$ further arrows connecting its underlying undirected
object graph. Thus $12\geq2k-1$ and $k\leq6$; there are no
additional cases with seven or more objects.

For one object the stronger monoid lower bound in
Corollary~\ref{B:fourdim-lower20} applies.
For two objects, categories with arrows in just one direction
are excluded by Lemma~\ref{C:twelve-triangular}, and those
with arrows in both directions by
Proposition~\ref{C:twelve-strongly-connected}.
Disconnected categories cannot have the connected homotopy
type under consideration.

For three objects,
Proposition~\ref{C:twelve-three-condensation} reduces the
problem to three singleton strongly connected components in
a linear order. The program
\path{code/c_categories12_linear.cpp} enumerates precisely
this remaining scope. Label the objects uniquely in increasing
order $0<1<2$; each of the three upper cross hom-sets is
nonempty and every lower cross hom-set is empty. Identity
arrows are labelled $0,1,2$, and the other arrows are sorted
by their source and target. The program enumerates every
such hom-cardinality profile and every type-correct product,
propagating all associativity equations. It discards only the
proved singleton cross-row or cross-column joins of
Lemma~\ref{C:singleton-cross-join}, and the previously proved
cofinal-object deletions. There is no comparison with an
incomplete earlier enumeration.

For four, five, and six objects, the complete
\path{code/c_categories12_cp2.cpp} runs enumerate all
connected hom-cardinality profiles and all associative
type-correct compositions after discarding initial- or
terminal-object profiles. All four completed enumeration
scopes have the following exact counts. The table count is
the number of associative labelled tables in the retained
scope; candidates have $H^1(-;\mathbb F_2)=0$ and
$\dim H^2(-;\mathbb F_2)=1$. The class count identifies
isomorphic or opposite candidate categories.
\[
\begin{array}{c|r|r|r|r|r}
 k&\text{tables}&\text{candidates}&\text{classes}
      &x^2=0&x^2\ne0,\ H^3\ne0\\\hline
 3\text{, linear}&1\,205\,250&165\,629&3\,624&3\,500&124\\
 4&36\,993\,658&3\,094\,704&4\,554&4\,545&9\\
 5&718\,340&77\,260&220&220&0\\
 6&51\,984&6\,000&6&6&0
\end{array}                                                    \tag{\ref{C:CP2-lower13}.1}
\]

Here is the independently checked obstruction for every
candidate class. The Python routine
\path{code/a_category_cups.py} verifies the identity,
source-target and associativity equations anew, constructs
the complete normalized nerve in the needed degrees, and
uses exact binary elimination to find the nonzero degree-two
cohomology class $x$. It computes its Alexander--Whitney
square. Every asserted zero square has an explicit degree-three
cochain primitive, checked on every normalized four-simplex.
For the 133 classes with nonzero square, a separate full
boundary-rank calculation gives $H^3(-;\mathbb F_2)\ne0$.
The generator's early $H^1,H^2$ tests and all these later
tests are necessary conditions for $\mathbb CP^2$, so each
case is excluded without inferring an all-degree homotopy
type from a homology cutoff.

The three-object stream is checked by
\path{code/c_twelve_stream_cups.py}, which marks completion
only after the generator's final summary exists and the
candidate file has been fully read. Its completed output is
\path{code/outputs/c_categories_12_3_linear_CP2_cups.json}:
all 165629 candidates have been read, all 3624 classes checked,
and the unexcluded list is empty.
The other completed outputs are
\path{code/outputs/c_categories_12_4_CP2_cups.json},
\path{code/outputs/c_categories_12_5_CP2_cups.json}, and
\path{code/outputs/c_categories_12_6_CP2_cups.json}.
The four-object file also contains its nine independent
residual rank calculations. Exact isomorphism caches use full
relabelled table strings, so a hash collision cannot suppress
a class. The known nonzero-square six-element class 190 is
rechecked as a positive control.

The earlier twelve-arrow streams labelled minimal, reduced,
or cofinal remain explicitly incomplete prefixes. None is
used for the completeness assertion in this theorem.
The completed scopes in (1), together with the structural
one-, two-, and three-object exclusions, exhaust every
possible object count and finish the proof.
\end{proof}

\subsection{The thirteen-arrow structural reduction}

\begin{lemma}\label{C:thirteen-three-SCC}
A three-object category with at most thirteen arrows and a
nontrivial strongly connected component cannot have classifying
space of type $\mathbb{CP}^2$.
\end{lemma}
\begin{proof}
An object that is a retract of another can be deleted without
changing the classifying space; the result has at most twelve
arrows and is excluded by Theorem~\ref{C:CP2-lower13}.
Assume there is no such retract.
For a strongly connected three-object category, the finite
categorical Rees description gives finite-group type if all row
numbers or all column numbers are one. Otherwise the two total
numbers are at least four, so its minimal ideal has at least
sixteen arrows, with all three identities outside that ideal.
Both alternatives exclude the proposed type.

The remaining possibility is a strongly connected block on two
objects and one singleton block. The first block has at least six
arrows: each endomorphism monoid has at least two elements and
there is an arrow in each cross direction. A cross arrow between
the blocks yields, by composition inside the first block, at least
one cross arrow for each of its two objects. Thus the first block
has at most ten arrows. A two-object block without a retract that
is not of finite-group type requires at least nine arrows in its
minimal ideal and two identities outside it, hence at least eleven.
Our block is therefore of finite-group type and rationally acyclic.

The singleton endomorphism monoid has order at most five. If its
order is at most four, it too is rationally acyclic. The mixed-nerve
bar decomposition then gives a rational cohomology ring equivalence
from the whole category to a suspension, so every positive-degree
product is zero. If its order is five, every count is forced:
the first block has six arrows and there are exactly two cross
arrows, one for each of its objects. These singleton cross hom-sets
make the category the categorical join of its two blocks, in one
order or the other. Its realization is the join of two connected
spaces, homotopy equivalent to a suspension of their smash product.
Its positive-degree products vanish as well.
\end{proof}

\begin{lemma}\label{C:thirteen-many-objects}
A connected category with thirteen arrows and at least seven
objects has contractible classifying space.
\end{lemma}
\begin{proof}
With $k$ objects connectedness requires at least $k-1$ nonidentity
cross arrows, so there are at least $2k-1$ arrows in all. Thus
$k=7$, there are exactly six cross arrows, and there are no
nonidentity endomorphisms. The underlying unoriented graph must
be a tree, with no parallel or reverse pair: either would leave
too few distinct edges for connectedness. There cannot be a
composable pair of cross arrows, since its composite would give
an additional edge joining the endpoints of a length-two tree
path. The nerve is consequently just that tree, with degeneracies,
and its realization is contractible.
\end{proof}

\begin{remark}\label{C:thirteen-search-scope}
The remaining thirteen-arrow computation for three through six
objects uses \texttt{c\_categories13\_minimal.cpp}. It enumerates
all connected hom-set cardinalities, retains one orbit under
object permutations and opposition, and solves every typed
associativity equation by propagation and exhaustive branching.
At each completed table it keeps the mod-two conditions
$H_1=0$, $\dim H_2=1$, and writes one representative under
identity-fixing permutations within each hom-set.
The latter reduction uses actual relabelings of the entire table;
no partial-table isomorphism heuristic is used.

The only extra deletions are a retract, an initial or terminal
object, the proved singleton-comma deletion, the singleton-join
cases, and Lemma~\ref{C:thirteen-three-SCC} for three objects.
Each remaining representative is checked independently for its
cup square, followed by third and fourth mod-two homology when
the square is nonzero. Unfinished output prefixes are explicitly
marked incomplete and support no complete-census claim. The
executable rejects two-object input; that case uses the separate
triangular and strongly connected arguments.
\end{remark}

\begin{proposition}\label{C:thirteen-four-plus}
A category with thirteen arrows and at least four objects cannot
have classifying space homotopy equivalent to $\mathbb{CP}^2$.
\end{proposition}
\begin{proof}
Disconnected categories cannot have the connected target type.
Lemma~\ref{C:thirteen-many-objects} handles seven or more objects.
For four, five and six objects, the complete algorithm and proved
pruning rules of Remark~\ref{C:thirteen-search-scope} give the
following finished counts. Object hom-cardinality profiles are
already taken modulo relabeling and opposition; the table counts
retain all labels within each such profile.
\[
\begin{array}{c|rrrrr}
\text{objects}&\text{tables}&\text{Hom orbits}
 &\text{classes}&x^2=0&b_3\ne0\\\hline
4&23649952&20692&20123&20103&20\\
5&12902&434&414&414&0\\
6&52&4&3&3&0.
\end{array}
\]
All retained categories have $H_1(-;\mathbb F_2)=0$ and
$\dim H_2(-;\mathbb F_2)=1$. The last two columns exhaust them.
The zero squares have explicit normalized cochain primitives;
the other twenty have nonzero third mod-two homology, computed
using both adjacent boundary ranks. The files
\path{code/outputs/c_categories_13_4_minimal_CP2_cups.json},
\path{code/outputs/c_categories_13_5_minimal_CP2_cups.json}, and
\path{code/outputs/c_categories_13_6_minimal_CP2_cups.json}
are all marked complete, with row counts matching the final
nonempty generator summaries. Their corresponding programs,
input tables and output summaries are supplied. Thus this proof
uses complete enumerations, not a prefix of the still-running
three-object search.
\end{proof}

\begin{lemma}\label{C:trivial-central-fork}
Let $C$ be a three-object category with a fork as its condensation
diagram, no arrows between the two leaves, and only the identity
endomorphism at the central object.  If $BC\simeq\CP^2$, then
one of its two full two-object subcategories also has classifying
space homotopy equivalent to $\CP^2$.  Consequently such a fork
with thirteen arrows cannot have type $\CP^2$.
\end{lemma}
\begin{proof}
Write $C_1,C_2$ for the two full subcategories containing the central
object and one leaf.  Every nerve simplex lies in one of these
subcategories, and their nerves intersect in precisely the central
vertex.  Realization is therefore literally the CW wedge
$BC_1\vee BC_2$.

If this wedge is $\CP^2$, van Kampen makes both factors simply
connected.  Integral reduced homology splits as the direct sum of
the homology of the two factors.  Since $H_2(\CP^2)=\mathbb Z$,
exactly one factor contains that group.  Products on a wedge are
componentwise, and the square of its primitive degree-two class
generates $H^4(\CP^2;\mathbb Z)$.  Thus the same factor contains
$H_4=\mathbb Z$.  All other positive homology groups of both
factors vanish.  The other simply connected factor is integrally
acyclic, hence contractible by the homological Whitehead theorem.
The first factor is therefore homotopy equivalent to $\CP^2$.
Each proper two-object block has at most twelve arrows, contradicting
Theorem~\ref{C:CP2-lower13}.
\end{proof}

\begin{theorem}\label{C:thirteen-all-stars}
No three-object category whose cross arrows form a star and
whose total number of arrows is at most thirteen has classifying
space homotopy equivalent to $\CP^2$.
\end{theorem}
\begin{proof}
Theorem~\ref{C:CP2-lower13} handles smaller totals.  Write
$m,n,p$ for the endomorphism orders and $s,t$ for the two nonempty
leg sizes, so $m+n+p+s+t=13$.  The relative-product argument of
Lemma~\ref{C:twelve-all-stars} excludes the case in which all
three vertex spaces have zero rational $H^2$.

A vertex with nonzero rational $H^2$ has order at least five.
If two vertices have that property, every size is forced:
the endomorphism orders are $5,5,1$ and both legs are singletons.
If the trivial vertex is central, apply
Lemma~\ref{C:trivial-central-fork}.  Otherwise one singleton leg
has a rationally acyclic opposite leaf, and
Lemma~\ref{C:star-singleton-join} gives zero products.

Assume there is exactly one vertex with nonzero rational $H^2$.
Each other vertex has order at most five; by the complete
classification in Theorem~\ref{E:all-six-types}, its classifying
space is rationally acyclic.  Suppose first that the distinguished
vertex is the center.  A singleton leg is again excluded by
Lemma~\ref{C:star-singleton-join}.  Otherwise $s,t\geq2$, and
the leaf orders have sum at most four.  If both are at most two,
their rational algebras are semisimple.  The exact-coinvariant
Tor calculation in Lemma~\ref{C:twelve-all-stars} applies verbatim:
the center has order at most seven, and its nonzero fourth homology
would persist in arbitrarily high degrees, a contradiction.

The remaining central case has leaf orders $3,1$, central order
five and leg sizes $2,2$.  The central monoid is necessarily the
five-element rectangular-band sphere monoid $P_2$, so its
classifying space is $S^2$.  Each two-object collage has zero
rational positive-degree products by the two-point assertion of
Lemma~\ref{C:two-object-cup-pruning}.  Mayer--Vietoris for their
union along $BP_2$ gives an injection
\[
 H^4(BC;\mathbb Q)\longrightarrow
 H^4(BC_1;\mathbb Q)\oplus H^4(BC_2;\mathbb Q),
\]
because $H^3(BP_2;\mathbb Q)=0$.  Every degree-two product
restricts to zero in both terms, so it is zero in $BC$ as well.

Finally suppose the distinguished vertex is a leaf, of order
$m\geq5$, and let $n$ be the central order and $s$ the size of
this distinguished leg.  Collapsing the rationally acyclic
center splits the star into its two leg cofibers by a rational
cohomology ring isomorphism.  Each cofiber is rationally equivalent,
through such a ring isomorphism, to its original collage.  The
other collage has two acyclic vertices and zero products.
Thus a nonzero square must occur in the distinguished collage,
whose rational homology must be bounded.

For $s\leq2$, Lemma~\ref{C:two-object-cup-pruning} excludes that
collage.  For $s\geq3$, the other leaf and leg contribute at least
two arrows, hence $m+n\leq8$.  If $n\leq2$, the all-act bounded
criterion of Lemma~\ref{C:eight-bounded-cups} applies.  Otherwise
the only profile is
\[
 \text{center }3,\qquad \text{distinguished leaf }5,
 \qquad \text{other leaf }1,
 \qquad (s,t)=(3,1).
\]
The five-element monoid is again exactly $P_2$.

The last profile is completely enumerated by
\path{code/c_thirteen_star_residual.py}.  It uses all seven ordinary
isomorphism classes of order-three monoids from the independent
census.  Every right action on the three-point leg is enumerated
up to state conjugacy, every left $P_2$-action in every relative
state position is retained, and every commuting pair is tested.
The singleton leg has forced actions.  Taking the opposite covers
inward stars.  The result is exactly $357$ commuting bisets:
$107$ have nonzero mod-two $H_1$, $196$ have the wrong mod-two
$H_2$ dimension, and $54$ have a zero degree-two square with an
explicit normalized cochain primitive.  The program verifies all
$111741$ typed associativity equations of these complete tables.
Every table, action and certificate is retained in
\path{code/outputs/c_thirteen_star_residual.json}.
There is no unexcluded case.  This finishes all star profiles.
\end{proof}

\begin{proposition}\label{C:thirteen-triangular}
No triangular two-object category with at most thirteen arrows
has classifying space homotopy equivalent to $\mathbb{CP}^2$.
\end{proposition}
\begin{proof}
Write $(m,n,s)$ for the two endomorphism orders and the cross-biset
size, with $m\geq n$. The twelve-arrow result handles smaller totals.
A singleton cross-biset gives a join. For two cross arrows, the
rationally acyclic vertex criterion leaves only $(6,5,2)$, with
both vertices having nonzero rational $H^2$. For $s\geq3$, the
relative-cup argument requires a vertex with nonzero rational
$H^2$, hence order at least five. The all-module obstruction through
order eight covers $n=1,m\leq8$ and $n=2,m\leq8$, since the
rational algebras of the two-element monoids are semisimple.
The remaining new profiles are
\[
\begin{gathered}
(9,1,3),\quad(5,3,5),\quad(5,4,4),\quad(6,3,4),\\
(5,5,3),\quad(6,4,3),\quad(7,3,3),\quad(6,5,2).
\end{gathered}
\]

For $(9,1,3)$, a monoid with nonzero rational $H^2$ has at least
two rows and two columns in its minimal ideal. The minimal-band
shapes $2\times2$, $2\times3$, and $3\times2$ are enumerated
directly by their translation pairs and all compatible outside
products, as in the eight-element construction. The complete counts
are
\[
\begin{array}{c|r|r|r}
\text{band}&\text{type profiles}&\text{associative sorted tables}
 &\text{ordinary isomorphism classes}\\\hline
2\times2&126&158631&7784\\
2\times3&5886&1672&191\\
3\times2&5886&1672&191.
\end{array}
\]
For the first shape there are four nonidentity outside elements
and the same six matching augmentation-character types, giving
$\binom{9}{4}=126$ profiles. For the other two shapes there are
two such elements and $108$ types, giving
$\binom{109}{2}=5886$ profiles. Associativity is propagated and
then checked on every triple. Isomorphism uses all row, column
and outside-element permutations; the minimal ideal is
characteristic. Both orientations are retained.
The independent program \texttt{c\_nine\_H2\_targets.cpp}
therefore produces exactly $8166$ tables. Rational balancing
discards $124$ zero-rank cases, leaving $8042$ targets.

The other possible nine-element minimal ideals are the
$2\times4$ and $4\times2$ bands, and a $2\times2$ Rees ideal
with maximal group of order two. In each case the only outside
element is the global identity. The contracted rational algebra
is $\mathbb Q$, while the minimal group algebra is semisimple.
The two augmentation sequences of the group-averaged band Tor
comparison give a projective resolution of the trivial rational
module of length at most two. Hence the act-cone formula has
zero fourth rational homology for every finite act. These cases
cannot give the target either.

For each of the $8042$ remaining tables, the program
\texttt{c\_nine\_all\_act\_cones.py} checks all right actions
on three points. It assigns transformations to a generating set,
uses only necessary power relations, and propagates every known
multiplication. At each choice the stabilizer of earlier images
removes only conjugate choices. Every completed assignment is
checked on all multiplication pairs. Thus this incremental
construction loses no action; the opposite tables cover left acts.
There are $175106$ action classes. The exact normalized-nerve
calculation excludes $39249$ by $H_1(-;\mathbb F_2)$ and
$69892$ by $\dim H_2(-;\mathbb F_2)\ne1$. Another $65755$
have explicit cochain primitives for the generator's square.
The remaining $210$ have $(b_3,b_4)=(1,1),(2,1),(2,3)$ with
multiplicities $204,4,2$, respectively, over $\mathbb F_2$.
All therefore have nonzero third homology. The full tables,
actions, primitives and residual ranks are in
\path{code/outputs/c_nine_all_act_cones.json}.

For the other seven profiles, the larger vertex ranges over
one, five, or forty-three independently enumerated monoids of
orders five, six, or seven with nonzero rational $H^2$.
The smaller vertex ranges over all ordinary classes of its order:
seven at order three, thirty-five at order four, and $228$ at
order five. In $(6,5,2)$ the latter may be restricted to the
five-element sphere monoid, because all other five-element
monoids are rationally acyclic. The order-seven list is generated
directly by all minimal-band extensions; the filtered universal-
group census is not used as an exhaustive source here.

The biset enumeration retains all right actions up to state
permutation and every relative conjugate of every left action,
then checks all commutation equations. The resulting complete
counts are as follows. The rejection columns are disjoint,
with $H_1$ tested first, then $H_2$, then the square.
\[
\begin{array}{c|rrrrr}
(m,n,s)&\text{bisets}&b_1\ne0&b_2\ne1&x^2=0&x^2\ne0\\\hline
(5,3,5)&17619&8097&6308&3209&5\\
(5,4,4)&23075&7756&10898&4413&8\\
(6,3,4)&14312&4913&5915&3452&32\\
(5,5,3)&22559&4656&13802&4095&6\\
(6,4,3)&13992&3042&8032&2918&0\\
(7,3,3)&20215&4342&11307&4538&28\\
(6,5,2)&44&5&9&29&1.
\end{array}
\]
Every one of the eighty nonzero-square cases has nonzero third
mod-two homology. Every zero square has a full normalized
three-cochain primitive. The completed certificate is
\path{code/outputs/c_thirteen_triangular_bisets.json}.
This excludes every profile in the displayed exhaustive list.
\end{proof}

\subsection{The triangular two-object exclusion through fourteen arrows}

\begin{lemma}[At most two elements outside the minimal ideal]
\label{C:two-outside-act-obstruction}
Let $M$ be a finite monoid with minimal ideal $K$, and suppose
$|M\setminus K|\leq2$.  For every finite-dimensional right
$R=\mathbb QM$-module $L$, eventual vanishing of
$\operatorname{Tor}_n^R(L,\mathbb Q)$ implies its vanishing for
every $n\geq3$.  In particular no finite right-act cone over $M$
has type $\mathbb{CP}^2$.
\end{lemma}
\begin{proof}
If $K=M$, the monoid is a finite group and its rational algebra
is semisimple.  Otherwise put $A=R/\mathbb QK$.
The Rees factorization
$\mathbb QK\cong Re\otimes_{\mathbb QH}eR$, with $eR$ free on
the left over the maximal group algebra, shows that $\mathbb QK$
is projective on the left over $R$.  Hence
$\operatorname{pd}_R A\leq1$.
Averaging $H$ gives a projective left module $P$ with its row-set
basis and an exact sequence
\[
0\longrightarrow V\longrightarrow P\longrightarrow\mathbb Q
\longrightarrow0.
\]
Every element of $K$ acts constantly on the row set, so it
annihilates the augmentation module $V$.  Thus $V$ is an $A$-module.

A rational algebra of dimension at most two is semisimple or is
$\mathbb Q[\varepsilon]/(\varepsilon^2)$.  In the semisimple case
$V$ is projective over $A$, so its projective dimension over $R$
is at most one.  This gives the asserted vanishing directly.
In the remaining case, Jordan normal form for the square-zero
operator gives
$V\cong A^a\oplus\mathbb Q^b$ as an $A$-module.
For $n\geq3$, dimension shifting therefore identifies
\[
\operatorname{Tor}_n^R(L,\mathbb Q)
 \cong\operatorname{Tor}_{n-1}^R(L,\mathbb Q_A)^b,
\]
where $\mathbb Q_A$ is the simple dual-number module restricted
to $R$.
If $b=0$ the conclusion follows.  If $b>0$, eventual vanishing
of these groups and duality give eventual vanishing of
$\operatorname{Ext}_R^j(\mathbb Q_A,DL)$.
The derived adjunction in Lemma~\ref{C:derived-Frobenius-Tor}
identifies these groups with
\[
\operatorname{Ext}_A^j\bigl(\mathbb Q,
                     \operatorname{RHom}_R(A,DL)\bigr).
\]
The second argument has cohomology only in degrees zero and one.
The minimal-injective-complex argument proved in that lemma
therefore makes these groups zero for every $j\geq2$.
Dimension shifting proves the claim.  Finally apply
Lemma~\ref{C:act-cone-Tor}; a projective-plane cone would have
nonzero $\operatorname{Tor}_3$ and an eventually zero tail.
\end{proof}

\begin{lemma}[A two-arrow biset opposite the sphere monoid]
\label{C:two-arrow-sphere-opposite}
Let $P=\{1\}\sqcup(\{0,1\}\times\{0,1\})$ be the
five-element sphere monoid.  If $H^2(BM;\mathbb Q)=0$, no
triangular biset category with vertex monoids $M,P$ and two
cross arrows has the rational cohomology ring of
$\mathbb{CP}^2$.  This requires no rational acyclicity of $BM$.
\end{lemma}
\begin{proof}
Pass to the opposite category if needed and write the biset as
a left $P$-set and right $M$-set.  Every element of the minimal
band of $P$ is idempotent.  Its image in the transformations of
a two-point set is therefore either the identity or a constant.
The rectangular-band equations imply, up to permuting the two
states, exactly three types: every band element acts identically;
every band element is the same constant; or the two rows act by
the two distinct constants.  Mixing an identity and a constant
would give a nontrivial semilattice image, which a rectangular
band cannot have.

In the first type the edge bar is $E_S\simeq E_R\times BP$,
and its map to $BP$ is projection.  Pullback on $H^2(BP;\mathbb Q)$
is injective.  Since $H^2(BM;\mathbb Q)=0$, the vertex restriction
of every degree-two class on the collage is zero.  Such classes
come from the relative cohomology of the two vertex spaces;
their products vanish because the corresponding relative quotient
is a suspension of the edge bar with a disjoint basepoint.

In the common-constant type, commutation forces the constant state
to be fixed by $M$ as well.  This gives a section
$BM\times BP\to E_S$ of the edge projection.  The map from
the direct sum of the two vertex $H^2$ groups into the edge group
is injective, so the same relative argument applies.

In the last type, a transformation commuting with both constants
fixes both states, so the $M$-action is trivial.  The left $P$-set
is its two-element minimal left ideal.  Its rational permutation
module is projective, and its action bar is rationally acyclic.
Thus $E_S\to BM$ is a rational homology equivalence.  The
homotopy-pushout description makes $BP\to BC$ a rational homology
equivalence, and in particular $H^4(BC;\mathbb Q)=0$.
All three cases exclude the stated ring.
\end{proof}

\begin{proposition}[The ten-element three-row act calculation]
\label{C:ten-three-row-three-point-acts}
No act-cone category with a ten-element vertex monoid and three
cross arrows has type $\mathbb{CP}^2$.
\end{proposition}
\begin{proof}
A vertex with zero rational $H^2$ is excluded by the relative-cup
argument.  A vertex with at most two minimal rows is excluded by
Theorem~\ref{C:two-row-act-cups}.  Of the remaining minimal-ideal
shapes at order ten, $4\times2$ and $3\times3$ bands have at
most two elements outside the ideal, and
Lemma~\ref{C:two-outside-act-obstruction} applies.  A nontrivial
minimal group with at least two rows and columns must have order
two and shape $2\times2$, already covered by the two-row theorem.
Thus only a $3\times2$ rectangular-band ideal remains.

The complete translation-table enumeration
\path{code/b_ten_non_square_targets.cpp} places the identity at
zero, the six band elements at $1,\ldots,6$, and the three other
elements at $7,8,9$.  There are $3^3 2^2=108$ possible linked
row/column transformation pairs.  Every multiset of three pairs
is tried, giving $\binom{110}{3}=215820$ profiles.  For every
ordered pair of outside elements, the possible product labels
are exactly those with the composed transformation pair.
All these labels are branched over, with associativity propagation;
every completed table is checked on every triple.  Isomorphism
uses every row, column, and outside permutation.  Since the
minimal ideal is characteristic, this is a complete enumeration.
It gives $58187$ associative sorted tables and $4260$ ordinary
isomorphism classes.  The opposite $2\times3$ enumeration gives
the same class count and is retained as a separate control.

Exact balanced ranks discard $1676$ of the $3\times2$ tables.
Another $1276$ have independently checked integral projective
resolutions of length at most two, so every act has zero fourth
rational homology.  For each of the remaining $1308$ tables,
all right actions on three states are enumerated by incremental
multiplication propagation, with exact stabilizers removing only
state conjugates.  There are $22741$ action classes.  The
disjoint rejection counts are
\[
\begin{array}{l|r}
\text{disconnected act, hence nonzero cone }H_1&5457\\
\text{wrong rational homology in degrees one through three}&14251\\
\text{wrong mod-two second Betti number}&88\\
\text{explicit mod-two primitive for the degree-two square}&2936\\
\text{nonzero square but nonzero mod-two }H_3&9.
\end{array}
\]
All nine last cases have $(b_3,b_4)=(1,1)$ over $\mathbb F_2$.
The rational calculations tensor an exact projective prefix through
$P_3$ with the act augmentation module, which computes
$\operatorname{Tor}_0,\operatorname{Tor}_1,\operatorname{Tor}_2$;
no omitted tail is used.  Every prefix is checked for literal
equivariance, zero composites, integral saturation, and exact rank
sums.  No resolution search failed on this list.
The normalized-bar primitives are checked on every degree-four
simplex.  Tables, actions, prefixes, primitives, and residual ranks
are in \path{code/outputs/b_ten_non_square_acts.json} and its
referenced compressed certificate chunks.
\end{proof}

\begin{proposition}[The triangular two-object exclusion through fourteen]
\label{C:fourteen-triangular}
No triangular finite category with two objects and at most fourteen
arrows has classifying space homotopy equivalent to
$\mathbb{CP}^2$.
\end{proposition}
\begin{proof}
The thirteen-arrow theorem handles smaller totals.  A singleton
cross-biset is a join and has zero positive-degree products.
For two cross arrows, Lemma~\ref{C:two-object-cup-pruning} leaves
the vertex-order pairs $(7,5)$ and $(6,6)$.  In the first pair,
Lemma~\ref{C:two-arrow-sphere-opposite} handles a seven-element
vertex with zero rational $H^2$; thus a possible $S^3$ vertex is
not silently called acyclic.  In the second pair, every six-element
vertex with zero rational $H^2$ is rationally acyclic by the
complete small-monoid classification.  It remains to enumerate
the displayed two-arrow profiles with both vertices having
nonzero rational $H^2$.

For at least three cross arrows, some vertex must have nonzero
rational $H^2$.  Choose it as the first vertex, passing to an
opposite category when appropriate.  If the other vertex has
order one or two, its rational coinvariants are exact.
Theorem~\ref{C:nine-act-cones}, the actual orbit-act comparison,
and Proposition~\ref{C:ten-three-row-three-point-acts} handle
all these cases.  The remaining profiles have both vertex orders
at least three and appear in the table below.  The additional
$(5,6,3)$ row covers the case where only the smaller vertex has
nonzero $H^2$.

The first vertex ranges over the complete lists of one, five,
forty-three, or $516$ possible $H^2$-positive tables at orders
five through eight, retaining the harmless zero-rank entries
in the order-eight list.  The other vertex ranges over all
ordinary classes of its order, except in the two-arrow rows,
where the preceding reductions permit both $H^2$-positive lists.
Every other-vertex monoid with projective trivial right rational
module is discarded using the all-act theorem through order nine;
each such discard is certified by an actual vector $v$ with
$vm=v$ for all $m$ and augmentation one.

For each remaining right act $S$, the program first enumerates
its complete endomorphism monoid by assigning the image of a
state and propagating equivariance on every monoid element.
It then enumerates every homomorphism from the opposite of the
other vertex monoid into this endomorphism monoid.  All power
and multiplication relations are checked.  No state permutation
is removed at this second stage, so every relative position of
the left action is retained.  This proves completeness of the
commuting-biset construction.  The disjoint checks give
\[
\begin{array}{c|r|r|r|r|r}
(m,n,s)&\text{bisets}&b_1\ne0&b_2\ne1&x^2=0&x^2\ne0\\\hline
(8,3,3)&43270&6636&24619&11985&30\\
(7,3,4)&30986&6400&13217&11343&26\\
(6,3,5)&27041&7181&13303&6557&0\\
(5,3,6)&55095&16503&19929&18663&0\\
(7,4,3)&25357&3782&13533&8014&28\\
(6,4,4)&27222&5743&13724&7735&20\\
(5,4,5)&67385&15627&36291&15467&0\\
(6,5,3)&23888&3773&12342&7743&30\\
(5,5,4)&55739&12078&26440&17212&9\\
(5,6,3)&53201&8623&26472&18076&30\\
(7,5,2)&430&43&129&257&1\\
(6,6,2)&254&25&79&135&15.
\end{array}
\]
All Betti numbers in this table are over $\mathbb F_2$.
Every one of the $189$ nonzero-square cases has nonzero third
mod-two homology.  Each zero square has an explicit normalized
degree-three primitive, checked on all degree-four simplices.
The complete run is \path{code/b_fourteen_small_bisets.py}, with
summary \path{code/outputs/b_fourteen_small_bisets.json} and
twelve referenced compressed files.  These contain every vertex
table, every projectivity splitting used to skip a vertex, every
primitive, and every residual homology calculation.  No numerical
tail vanishing is inferred.  All possible profiles are therefore
excluded.
\end{proof}

\begin{remark}\label{C:fourteen-triangular-verification-scope}
The above table and the ten-element calculation are complete exact
runs of their stated enumeration and cochain algorithms.  The additional
independent biset gates, which were still running at the ninth checkpoint,
are now complete.  The fresh pointwise-domain constraint solver
\path{code/b_verify_fourteen_biset_coverage.cpp} uses neither the
producer's monoid generators nor its endomorphism-monoid enumeration.
It independently enumerates $185847$ labelled right actions, passes
through $1176994$ constraint-search nodes, and reconstructs all
$409868$ commuting bisets.  Every per-vertex first- and second-homology
count agrees, and there is an exact bijection between its $123376$
degree-two survivors and the saved certificate records.  Its complete
output is \path{code/outputs/b_fourteen_biset_coverage_verified.jsonl}.

The separate verifier
\path{code/b_verify_fourteen_biset_primitives.cpp} reconstructs every
category from its two monoids and two actions.  It checks $73075962$
typed associativity equations, $43924444$ degree-two cocycle equations,
and every degree-two dimension and nonboundary generator independently.
All $123187$ zero-square certificates satisfy their literal coboundary
identity on $237022004$ normalized quadruples in total.  All $189$
remaining records have independently recomputed positive third homology.
The output is \path{code/outputs/b_fourteen_biset_primitives_verified.jsonl}.
The rational splittings used to skip projective-trivial opposite corners
are separately verified with exact fraction arithmetic in
\path{code/b_prepare_fourteen_biset_audit.py}.

The broader
order-nine resolution, stable-retract, rational-cup, and three/four-state
action calculations have already passed the separate checker
\path{code/b_verify_nine_small_acts.py}.  Its fresh all-element
action CSP independently matches all $18144$ residual action orbits
without using the original generator or stabilizer construction.
The global category lower bound is not advanced here: categories
with at least three objects require their own complete coverage.
\end{remark}

\subsection{The triangular two-object exclusion through fifteen arrows}

\begin{lemma}[The rational all-act obstruction at order ten]
\label{C:ten-rational-act-cones}
Let $|M|=10$ and let $S$ be any nonempty finite right $M$-set.
If the act cone has bounded rational homology, then all of its
rational degree-two products vanish.  In particular it cannot
have the rational cohomology ring of $\mathbb{CP}^2$.
\end{lemma}
\begin{proof}
A vertex with zero rational $H^2$ gives only relative degree-two
classes, whose products vanish.  Otherwise its minimal ideal has
at least two rows and columns.  At most two rows are covered by
Theorem~\ref{C:two-row-act-cups}, including nontrivial minimal
groups.  The order bound leaves only band shapes $3\times2$,
$4\times2$, and $3\times3$ when there are at least three rows.
For $3\times2$ there are four outside elements, and the rational
bounded-homology form of
Theorem~\ref{C:two-column-four-outside-act-cones} applies.
The other two shapes have at most two outside elements; by
Lemma~\ref{C:two-outside-act-obstruction}, bounded rational
homology forces their act cones to have zero fourth rational
homology.  This proves the assertion in every case.
\end{proof}

\begin{lemma}[Three states opposite a two-element right-zero band]
\label{C:three-state-right-zero-biset}
Let $N=\{1,p,q\}$, with $p^2=p$, $q^2=q$, $pq=q$, and $qp=p$.
Let $S$ be a three-element $(N,M)$-biset.  If its collage has
type $\mathbb{CP}^2$, then an ordinary right $M$-act cone on at
most three states has that type.  Consequently no such collage
with $|M|\leq17$ has type $\mathbb{CP}^2$.
\end{lemma}
\begin{proof}
Write $P,Q:S\to S$ for the two left actions.  The identities
$PQ=Q$ and $QP=P$ imply that both maps are retractions onto
the same nonempty subset $U$, on which both are the identity.
The right $M$-action commutes with both retractions.
If $|U|=1$, its point is fixed by both monoids.  The resulting
section $BM\times BN\to E_S$ of the edge double bar makes every
degree-two class on the collage relative to the vertex spaces.
Its square vanishes because the relative quotient is a suspension.
If $|U|=3$, the left action is trivial.  Since $BN$ is
contractible, collapsing this vertex gives an actual homotopy
equivalence with the ordinary three-state right-act cone over $M$.

Suppose $|U|=2$, and denote the remaining point by $z$.
If $Pz\ne Qz$, the left-action category is the union of the two
full subcategories on $\{z,Pz\}$ and $\{z,Qz\}$.  Each has
initial object $z$: its unique arrow from $z$ to its other object
is fixed under postcomposition by every $N$-endomorphism.  Their
intersection is the singleton $z$, and no simplex crosses between
the two arms.  Its nerve is therefore contractible.  Realizing
the right $M$-action bar now makes $E_S\to BM$ a homotopy
equivalence, by objectwise homotopy invariance of the bar construction.
The collage is contractible as well, since $BN$ is contractible.

If $Pz=Qz$, the common retraction $S\to U$ is an actual biset
retraction.  It gives a retraction from the collage $X_S$ to
$X_U$.  The latter is equivalent to the ordinary two-state
$M$-act cone and has zero positive-degree rational products by
Lemma~\ref{C:finite-line-act-cups}.  The inclusion of double bars
$E_U\subset E_S$ is a simplicial cofibration, and the quotient
of the two collage models is $\Sigma(E_S/E_U)$.  If
$\dim H^2(X_S;\mathbb Q)=1$, the retraction leaves two possibilities.
When $H^2(X_U;\mathbb Q)\ne0$, its pullback spans $H^2(X_S;\mathbb Q)$,
so the square vanishes.  When $H^2(X_U;\mathbb Q)=0$, every
degree-two class of $X_S$ is relative to $X_U$; the displayed
suspension quotient again makes its square zero.
These cases leave only the ordinary three-state act cone.
The last assertion follows from
Theorem~\ref{C:three-state-cone-minimum22}; the one- and two-state
cases have zero positive products directly.
\end{proof}

\begin{proposition}[The complete fifteen-arrow triangular calculation]
\label{C:fifteen-triangular}
No triangular finite category with two objects and at most fifteen
arrows has classifying space homotopy equivalent to $\mathbb{CP}^2$.
\end{proposition}
\begin{proof}
Proposition~\ref{C:fourteen-triangular} handles smaller totals.
Write $(m,n,s)$ for the two vertex orders and the number of cross
arrows.  A singleton biset gives a join, with zero positive
products.  For two cross arrows,
Lemma~\ref{C:two-state-biset-H2zero} permits restriction to
two vertices with nonzero rational $H^2$, leaving $(8,5,2)$
and $(7,6,2)$ up to opposition.
For $s\geq3$, some vertex has nonzero rational $H^2$; choose
it as the first vertex, taking the opposite category if necessary.

If the other vertex has order at most two, its rational
coinvariants are exact, and its actual orbit-act comparison
reduces the question to an ordinary act cone.  Vertices of order
at most nine are covered by the bounded-rational statement of
Lemma~\ref{C:nine-bounded-rank-one}.
The remaining profiles have a ten-element vertex and at most
four states, or an eleven-element vertex and at most three states.
The former are excluded by Lemma~\ref{C:ten-rational-act-cones},
which applies directly to the rational equivalence produced by
orbit coinvariants.  For the latter profile the opposite vertex
has order one, so this is an actual three-state act cone and
Theorem~\ref{C:three-state-cone-minimum22} applies.
One or two states were already excluded by the finite-line lemma.
Thus all residual profiles are precisely the seventeen rows below.
The larger order may occur in the second position when it is
the smaller vertex that supplies the nonzero degree-two class.

The first-vertex catalogs have $1,5,43,516$ tables at orders
five through eight, as in the preceding complete calculation.
At order nine, retain all $8166$ translation tables of the
$2\times2$, $2\times3$, and $3\times2$ minimal-band enumeration,
including its $124$ harmless rational-rank-zero entries.  Add
the identity-adjoined $2\times4$ and $4\times2$ bands and both
normalized $2\times2$ Rees ideals over $C_2$.  This gives a
complete $8170$-table cover of the possible positive-degree-two
vertices.  The last four tables are included literally: a
projective-dimension bound for an ordinary act cone alone would
not justify discarding them opposite a nonprojective corner.

For the second vertex take every ordinary class of orders
three through six.  At order seven, the independent complete
$17730$-class isomorphism-and-opposition census and its transposes
give $32823$ distinct literal tables.  This is an exhaustive
oriented cover; it is not asserted to be an ordinary isomorphism
class count.  In the two-state rows, use the positive-degree-two
catalogs for both vertices.  Any second vertex with projective
trivial right rational module is skipped only after supplying
and checking an actual vector $v$ satisfying
$vn=v$ for every $n$ and $\epsilon(v)=1$.
The orbit-act comparison and the bounded-rational statement of
Lemma~\ref{C:nine-bounded-rank-one} then apply.

The pointwise constraint solver assigns every value of every
right action and enforces all multiplication equations.  It
retains one representative under the full state-permutation
group.  For each such right action it assigns all left-action
values, imposes the opposite multiplication equations and every
commuting equation, and retains every solution.  No quotient by
a second, independently chosen state permutation is taken.
Thus every relative position of the two actions occurs.
All completed tables and action equations are checked exactly.
The disjoint normalized-bar outcomes are
\[
\begin{array}{c|r|r|r|r|r}
(m,n,s)&\text{bisets}&b_1\ne0&b_2\ne1&x^2=0&x^2\ne0\\\hline
(9,3,3)&817136&112228&479153&225333&422\\
(8,3,4)&480725&88987&222031&169460&247\\
(8,4,3)&366630&48500&207170&110750&210\\
(7,3,5)&358971&81252&182780&94735&204\\
(7,4,4)&333826&59897&180226&93477&226\\
(7,5,3)&271272&35609&154775&80384&504\\
(6,3,6)&368606&99538&149247&119786&35\\
(6,4,5)&433128&91248&240070&101810&0\\
(6,5,4)&355692&68244&181157&106003&288\\
(6,6,3)&324811&46065&174566&104030&150\\
(5,3,7)&921591&272192&482520&166879&0\\
(5,4,6)&1684292&318335&1071600&294357&0\\
(5,5,5)&1121618&233893&625012&262684&29\\
(5,6,4)&863223&174107&416858&271958&300\\
(5,7,3)&902825&131912&468763&301828&322\\
(8,5,2)&5604&516&1996&3084&8\\
(7,6,2)&2492&215&987&1199&91.
\end{array}
\]
Here every Betti number is over $\mathbb F_2$.
All $3036$ cases in the last column have nonzero third homology.
The other $2507757$ retained cases have explicit normalized
degree-three primitives for the square of a certified nonboundary
degree-two cocycle.  Consequently all $9612442$ bisets are excluded.
The first row also follows from
Lemma~\ref{C:three-state-right-zero-biset}: among the seven
three-element monoids, only the displayed right-zero extension
fails the checked rational right-trivial splitting test.

The catalogs and all projectivity witnesses are saved in
\path{code/outputs/b_fifteen_biset_input_manifest.json} and its
seventeen compressed input files.  The complete producer is
\path{code/b_fifteen_bisets.cpp}, with disjoint completion logs
and all retained certificates under
\path{code/outputs/b_fifteen_bisets_00.jsonl.gz} and
\path{code/outputs/b_fifteen_bisets_rest.jsonl.gz}.
The aggregate counts are in
\path{code/outputs/b_fifteen_bisets_summary.json}.
\end{proof}

\begin{remark}[Independent verification gate]
\label{C:fifteen-triangular-verification}
The independent pointwise action construction used here was
checked against every outcome of the fourteen-arrow enumeration.
For the new certificates, a separate normalized-bar implementation
rebuilds every literal category, checks all typed associativity
equations, verifies its second Betti number and the nonboundary
generator, and checks every claimed primitive on every degree-four
simplex.  The alternative certificates are checked by independent
third-homology ranks.  All seventeen profiles pass this independent
verification: $2510793$ retained certificates, $1859788417$ typed
associativity equations, $1167434995$ cocycle equations, and
$7023259440$ primitive equations.  Every one of the $3036$
alternative cases has independently recomputed positive third
homology.  The input preparation also checks all $31232$ saved
rational projectivity splittings with exact fractions.  The
separate final records are
\path{code/outputs/b_fifteen_biset_primitives_verified00.jsonl}
and
\path{code/outputs/b_fifteen_biset_primitives_verified_rest.jsonl}.
The verifier imports none of the producer's action enumeration,
cup solver, or cohomology routines.
\end{remark}

\begin{corollary}\label{C:fifteen-two-objects}
A finite category with at most two objects and classifying space
homotopy equivalent to $\mathbb{CP}^2$ has at least sixteen arrows.
\end{corollary}
\begin{proof}
For one object use Theorem~\ref{B:CP2-minimum23}.
A disconnected category cannot realize a connected space.
For two objects in one strongly connected component use
Theorem~\ref{C:twentytwo-strongly-connected}; the remaining
connected case is Proposition~\ref{C:fifteen-triangular}.
\end{proof}

\begin{theorem}\label{C:CP2-lower14}
Every finite category whose classifying space is homotopy
equivalent to $\CP^2$ has at least fourteen arrows, identities
counted.  The current minimum-arrow interval is $14$--$23$.
\end{theorem}
\begin{proof}
Theorem~\ref{C:CP2-lower13} excludes smaller totals.  At thirteen
arrows the monoid case follows from Theorem~\ref{B:CP2-minimum23}.
Two-object categories are excluded by the strongly connected
two-object theorem and the complete triangular theorem above.
Proposition~\ref{C:thirteen-four-plus} handles at least four objects.
For three objects, Lemma~\ref{C:thirteen-three-SCC} excludes every
nontrivial strongly connected block, and
Theorem~\ref{C:thirteen-all-stars} excludes the star condensations.
It remains to check linear condensation with all three cross-hom
sets nonempty.

Write $(m,n,p)$ for the endomorphism orders and $(x,y,z)$ for
the cardinalities in directions $0\to1$, $1\to2$, $0\to2$.
There are exactly $792$ positive six-tuples with sum thirteen.
The independent profile checker
\path{code/c_thirteen_linear_coverage.py} enumerates these
six-tuples directly.  It excludes $204$ by the singleton-extreme
join and $161$ more by the singleton cofinal deletion to a trivial
corner.  The other $427$ profiles form $228$ orbits under object
relabeling and opposition.  Every one occurs in the certified
completed portion of the associativity search.

The complete searched portion has $301$ canonical hom-cardinality
profiles and $87519771$ associative labelled tables in those
profiles.  Its exact mod-two $H_1=0$, $\dim H_2=1$ filter retains
$11945202$ labelled tables, represented by $114896$ literal
hom-set-permutation orbits.  Every representative is independently
checked on all typed associativity triples and on its normalized
cup square.  The final partition is
\[
 \begin{array}{c|r}
 \text{explicit zero-square primitive}&113695\\
 \text{nonzero square and }H_3(-;\mathbb F_2)\ne0&1201.
 \end{array}
\]
There are $15670686$ typed associativity checks in this verification.
The two complete certificate scopes are
\path{c_categories_13_3_completed_profiles_CP2_hom_cups.json}
and \path{c_categories_13_3_completed_extra_CP2_hom_cups.json}
in \path{code/outputs/}; their source tables, row hashes and
cochain certificates are supplied.  The separate file
\path{code/outputs/c_thirteen_linear_coverage.json} certifies
the full cardinality-profile coverage.

For precision, the original search was stopped in a subsequently
excluded star branch.  A completion line in its progress log was
not treated as sufficient evidence of a flushed output file:
the last completed linear profile had $27$ missing buffered rows.
That entire profile was rerun and flushed, then independently
checked in the second complete scope.  The first scope contains
only profiles whose exact recorded row counts match their literal
complete tables.  Thus no incomplete stream is a premise of this
exclusion.  The final complementary profile run has no further
associative tables.  All thirteen-arrow cases are now excluded.
The $23$-arrow upper bound is the explicit monoid of
Theorem~\ref{B:CP23}.
\end{proof}

\subsection{The fourteen-arrow structural reduction}

\begin{lemma}\label{C:fourteen-three-SCC}
A three-object category with at most fourteen arrows and a
nontrivial strongly connected component cannot have classifying
space homotopy equivalent to $\mathbb{CP}^2$.
\end{lemma}
\begin{proof}
Retract deletion reduces to at most thirteen arrows, already
excluded, so suppose no object is a retract of another.
If the category is strongly connected, the categorical Rees
description gives finite-group type when all row numbers or all
column numbers are one. Otherwise its minimal ideal has at least
$4\cdot4=16$ arrows, which already exceeds the budget.

The remaining condensation has a strongly connected block $S$
on two objects and a singleton monoid $M$. Write $s,m,t$ for
their respective arrow counts and the number of arrows between
the blocks. We have $s\geq6$, $t\geq2$, and $s+m+t\leq14$.
If $BS$ is not of finite-group type, its minimal ideal has at
least nine arrows and both identities lie outside it; hence
$s\geq11$. Equality is forced, with $m=1$ and $t=2$.
There is then exactly one cross arrow for each object of $S$,
so the entire category is a categorical join with a point and
has contractible classifying space.

Thus $BS$ has finite-group type and is rationally acyclic.
If $m\leq4$, then $BM$ is also rationally acyclic. The actual
mixed-nerve homotopy pushout, followed by collapse of the two
vertex spaces, induces a rational cohomology ring isomorphism
to a suspension. All positive products vanish.

If $m\geq5$, the budget gives $s\leq7$ and $m\leq6$.
Let $r_1,r_2,c_1,c_2$ and $G$ be the Rees data of $S$.
Its two identities lie outside the minimal ideal, and
\[
 |G|(r_1+r_2)(c_1+c_2)+2\leq7.
\]
Consequently $G=1$ and every row and column number is one.
For a minimal idempotent $e$, the representable rational modules
$\mathbb QS e$ and $e\mathbb QS$ have one basis vector at each
object, and their augmentations to the constant modules are
isomorphisms. Both constant modules of $S$ are therefore projective.

Collapse the whole block $S$ to a single trivial object and
replace its cross-arrow diagram by its set of $S$-coinvariants.
This defines an actual functor to a triangular two-object
category $C'$ with monoid $M$ at its other vertex. In the
mixed two-sided bar, exactness of the $S$-coinvariant functor
identifies the edge bar rationally with the corresponding
finite $M$-act bar. On vertex spaces the collapse $BS\to *$
is a rational homology equivalence. The homotopy-pushout
comparison therefore proves that $BC\to BC'$ is a rational
homology equivalence and hence induces a cohomology ring
isomorphism. Its rational homology is bounded if $BC$ has
the proposed type. Lemma~\ref{C:eight-bounded-cups}, with
$|M|\leq6$ and the other vertex trivial, now makes every
degree-two product zero. This excludes the last case.
\end{proof}

\begin{lemma}[Generic star deletions]\label{C:generic-star-deletions}
Consider an outward three-object star with central endomorphism
order $m$, leaf orders $a,b$, and corresponding leg sizes $x,y$.
Within a total budget of fourteen arrows, the following conditions
exclude $\mathbb{CP}^2$; the same statements hold for inward stars.
\begin{enumerate}
\item $m=1$;
\item $x=1$ and $b\leq4$, or $y=1$ and $a\leq4$, or $x=y=1$;
\item $m\leq8$ and $a,b\leq2$;
\item $m,b\leq4$ and $x\leq2$, or $m,a\leq4$ and $y\leq2$;
\item $m\leq2$, $b\leq4$, $a\leq8$, or the same condition
with the two leaves exchanged.
\end{enumerate}
If none of the three vertex monoids has nonzero rational $H^2$,
the target is also excluded.
\end{lemma}
\begin{proof}
Opposition reverses the star without changing its classifying
space. The last assertion is the relative-product argument in
Lemma~\ref{C:twelve-all-stars}. For (1), the nerve is a literal
wedge of its two collages. Lemma~\ref{C:trivial-central-fork}
places a hypothetical projective-plane type in a proper collage,
whose at most thirteen arrows contradict
Theorem~\ref{C:CP2-lower14}.

The first two alternatives of (2) are
Lemma~\ref{C:star-singleton-join}. If both legs are singletons,
the two collages are joins, and the inclusion of the central
space into each is null-homotopic. Their homotopy pushout is a
wedge of the two joins and a suspension of the central space.
All positive products vanish.

For (3), the two leaf algebras are rationally semisimple. Their
exact coinvariants give a map from the sum of the two edge
modules to the central constant module. The augmentation kernel
of this \emph{sum}, rather than either individual kernel, gives
the same all-module Tor comparison as in
Lemma~\ref{C:twelve-all-stars}. The all-module criterion through
order eight in Proposition~\ref{C:eight-act-cones} says that
nonzero fourth homology then persists in arbitrarily high degrees,
unless the central $H^2$ is zero; in that case the relative-product
argument applies. Both alternatives exclude the target.

In (4), the center and the nondistinguished leaf are rationally
acyclic. Collapsing the center induces a rational cohomology-ring
isomorphism from the star to the wedge of its two collage cofibers.
The nondistinguished collage has two rationally acyclic vertices
and zero products. The distinguished collage has a leg of at most
two elements, so Lemma~\ref{C:two-object-cup-pruning} applies.
For (5), the same wedge comparison transfers bounded rational
homology to the distinguished collage. Exact central coinvariants
and Lemma~\ref{C:eight-bounded-cups} exclude its nonzero square.
\end{proof}

\begin{theorem}\label{C:fourteen-all-stars}
A three-object star category with at most fourteen arrows cannot
have classifying space homotopy equivalent to $\mathbb{CP}^2$.
\end{theorem}
\begin{proof}
Smaller totals are already excluded. Enumerate the positive
five-tuples $m+a+b+x+y=14$, modulo exchange of the two leaves,
and use Lemma~\ref{C:generic-star-deletions}. A vertex with
nonzero rational $H^2$ has at least five elements. The eleven
remaining cardinality profiles and complete finite checks are
as follows. The homology columns use $\mathbb F_2$.
\[
\begin{array}{c|r|r|r|r|r}
(m,a,b,x,y)&\text{cases}&b_1\ne0&b_2\ne1&u^2=0&b_3\ne0\\\hline
(3,1,5,1,4)&3549&1231&1440&878&0\\
(3,1,5,2,3)&863&502&279&82&0\\
(3,1,6,1,3)&2555&633&1193&729&0\\
(3,2,5,1,3)&714&214&392&108&0\\
(4,1,5,1,3)&2264&574&1297&393&0\\
(5,1,3,2,3)&831&449&195&187&0\\
(5,1,3,3,2)&270&159&81&30&0\\
(5,1,4,2,2)&849&371&206&272&0\\
(5,1,5,1,2)&8&1&2&5&0\\
(5,2,3,2,2)&450&191&144&115&0\\
(6,1,3,2,2)&877&349&262&259&7
\end{array}
\]
At an order-five vertex with positive rational $H^2$, the monoid
is $P_2$; at order six the complete minimal-band list is the
function \texttt{six\_H2} used in the independently checked
thirteen-arrow triangular computation. All ordinary isomorphism
classes at orders at most four are restored from the independent
monoid census, including opposites. In the unique profile with
two potentially positive vertices, $(5,1,5,1,2)$, only $P_2,P_2$
need be retained. If the other order-five monoid has zero rational
$H^2$, its rational acyclicity from the complete small-order
classification, followed by the singleton or two-point collage
argument, gives zero products.

For each leg, enumerate every right central action up to state
conjugacy and every left leaf action in every relative state
position; retain precisely the commuting pairs. The two legs
have independent state relabelings and no composites with one
another, so taking every pair of these bisets is exhaustive.
The complete tables and certificates are in
\path{code/outputs/c_fourteen_star_residuals.jsonl}, generated by
\path{code/c_fourteen_star_residuals.py}.

The standalone verifier \path{code/d_verify_fourteen_stars.py}
independently recovers the eleven budget profiles, checks every
one of the $13230$ literal tables and all $5096182$ typed
associativity equations, recomputes the low homology, and checks
all $2861687$ stored cup-primitive equations. In the seven
nonzero-square cases it independently verifies nonzero third
homology. Its complete result is
\path{code/outputs/d_verify_fourteen_stars.json}.
The displayed reasons exhaust every case and prove the theorem.
\end{proof}

\subsection{An eleven-arrow strong pair with a three-arrow exit}

\begin{lemma}[Deleting a retract-supported vertex]
\label{C:retract-supported-vertex}
Let $C$ contain distinct objects $a,b$ such that
$\End_C(b)=\{1_b,e\}$ with $e^2=e\ne1_b$.
Suppose $e$ fixes every arrow incident to $b$ and another object:
$ep=p$ for $p:x\to b$, $pe=p$ for $p:b\to x$, $x\ne b$.
If there are $u:b\to a$, $v:a\to b$ with $vu=e$, then deleting
$b$ preserves the homotopy type of the classifying space.
\end{lemma}
\begin{proof}
Write $C'$ for the full subcategory omitting $b$, and $J:C'\to C$
for its inclusion.  Define $G:C\to C'$ by $G(b)=a$, by fixing
all retained objects and arrows, and by
\[
 G(1_b)=1_a,\quad G(e)=uv,\quad
 G(p:x\to b)=up,\quad G(q:b\to x)=qv.
\]
The element $uv$ is idempotent.  A composite through $b$ is
preserved because $qvu p=qep=qp$.  A composite through a retained
object and beginning and ending at $b$ cannot equal $1_b$:
it is fixed by multiplication by $e$, whereas $1_b$ is not.
It therefore equals $e$, and its image is $uv$ as required.
The other functor equations follow from the stated support
identities.  Thus $GJ=\id_{C'}$.  The components $1_x$ at retained
objects and $v$ at $b$ give a natural transformation
$JG\Rightarrow\id_C$.  Hence $BJ$ and $BG$ are homotopy inverses.
\end{proof}

\begin{proposition}\label{C:eleven-strong-three-exit}
Let $S$ be the eleven-arrow two-object category formed by a
rectangular minimal ideal with three total rows and three total
columns, and one additional identity at each object.  Suppose
each object has at least one row and one column.  Form $C$ by
adjoining a sink $c$ with $\End(c)=\{1_c\}$ and three total
arrows from $S$ to $c$, with arbitrary compatible actions.
Then $BC$ does not have the homotopy type of $\mathbb CP^2$.
\end{proposition}
\begin{proof}
Write the two objects as $a,b$.  The row counts are $2,1$, and
the column counts are $2,1$.  The quotient of the rational
category algebra by its minimal ideal is
$\mathbb Q1_a\oplus\mathbb Q1_b$.  The natural categorical
balanced-tensor presentation therefore identifies $H_2(BS;\mathbb Q)$
with the tensor product of the row and column augmentation modules
over this semisimple algebra.  If the two-dimensional row and
column sets occur at different objects, the augmentation modules
have disjoint supports, so $H^2(BS;\mathbb Q)=0$.
Every degree-two class of the act cone is then relative to $BS$;
its products vanish because the relative quotient is the suspension
of the action category.  This excludes the indicated case.

Suppose instead that both two-dimensional sets occur at $a$.
Then $\End_S(a)=P_5$, the five-element sphere monoid, while
$\End_S(b)=\{1_b,e\}$ and both cross-hom sets have two elements.
The same balanced-tensor formula shows that the full corner
inclusion induces an isomorphism
\[
 H_2(BP_5;\mathbb Q)\longrightarrow H_2(BS;\mathbb Q).
                                                               \tag{\ref{C:eleven-strong-three-exit}.1}
\]
Both source objects have at least one arrow to $c$, since they
are in the same strongly connected component.  Their arrow
counts are therefore either $(2,1)$ or $(1,2)$.

In the first case, the unique arrow $b\to c$ is fixed by $e$,
as are all the original arrows between $a$ and $b$.
Any $u:b\to a$, $v:a\to b$ satisfy $vu=e$ by the rectangular
composition law.  Lemma~\ref{C:retract-supported-vertex} deletes
$b$.  The resulting category is an act cone over $P_5$ with two
states and has eight arrows.  Its degree-two products vanish
by Theorem~\ref{C:two-row-act-cups}.

In the second case there is a unique arrow $a\to c$, fixed by
every element of $P_5$.  It gives a functor from $P_5$ to the
action category $E$ of the exit, whose composite with
$E\to S$ is the full-corner inclusion.  By
(\ref{C:eleven-strong-three-exit}.1), the map
$H_2(BE;\mathbb Q)\to H_2(BS;\mathbb Q)$ is surjective.
Thus $H^2(BS;\mathbb Q)\to H^2(BE;\mathbb Q)$ is injective.
The cofiber exact sequence again makes every degree-two class
of $BC$ relative to $BS$, and the relative-suspension argument
makes every degree-two product zero.  All cases contradict the
nonzero square of the rational generator on $\mathbb CP^2$.
\end{proof}

\subsection{Deleting a corner fixed on one side}

\begin{lemma}[One-sided corner deletion]
\label{C:one-sided-corner-deletion}
Let $C$ be a finite category, let $a\ne b$, and suppose there are
arrows $u:b\to a$ and $v:a\to b$ such that
\[
 vu=e,\qquad e^2=e\ne1_b.
\]
Assume that $eg=g$ for every arrow $g:d\to b$ with $d\ne b$,
and that every unit of $\End_C(b)$ fixes every such incoming arrow.
Then the inclusion of the full subcategory $D=C\setminus\{b\}$
induces a homotopy equivalence $BD\simeq BC$.
The unit hypothesis holds in particular if the identity is the
only unit, or if $e$ is a zero of $\End_C(b)$.
There is a dual assertion with outgoing arrows in place of incoming
arrows.  No condition on outgoing arrows is needed in the stated
incoming version.
\end{lemma}
\begin{proof}
Put $A=\End_C(b)$ and write $A^\times$ for its unit group.
In a finite monoid, the nonunits form a two-sided ideal: a
one-sided inverse is an inverse, as follows by applying the
corresponding injective self-map of the finite monoid.
Define a functor $G:C\to D$ on objects by $G(b)=a$, fixing
every other object.  It fixes arrows of $D$ and is given on
the remaining arrows by
\begin{align*}
 G(h)&=1_a &&(h\in A^\times),\\
 G(h)&=uhv &&(h\in A\setminus A^\times),\\
 G(g)&=ug &&(g:d\to b,\ d\ne b),\\
 G(f)&=fv &&(f:b\to d,\ d\ne b).
\end{align*}
These formulas preserve identities.  We check every other
composition involving $b$.

For two nonunit endomorphisms $h,k$, their product is a nonunit
and
\[
 G(h)G(k)=uhvukv=uhekv=uhkv=G(hk),
\]
because $kv:a\to b$ is an incoming cross arrow fixed by $e$.
If one factor is a unit, its action on every incoming arrow is
the identity.  Thus $hkv=kv$ when $h$ is a unit, and $kv=v$
when $k$ is a unit.  These identities verify both mixed
endomorphism products; a product of two units is immediate.

Composing a nonunit $h$ after an incoming arrow $g$ gives
\[
 G(h)G(g)=uhvug=uheg=uhg=G(hg).
\]
The unit case uses $hg=g$.  Composing a nonunit $h$ before an
outgoing arrow $f$ gives
\[
 G(f)G(h)=fvuhv=fehv=fhv=G(fh),
\]
because $hv$ is incoming; the unit case uses $hv=v$.
For a composite through $b$ between retained objects,
\[
 G(f)G(g)=fvug=feg=fg.
\]
For a composite $gf:b\to b$ through a retained object,
$e(gf)=gf$.  It cannot be a unit, since multiplying that
identity by its inverse would imply $e=1_b$.  Consequently
\[
 G(g)G(f)=ugfv=G(gf).
\]
Compositions with other retained arrows follow directly by
associativity.  These cases exhaust all typed pairs, so $G$
is indeed a functor.

Let $J:D\to C$ be the inclusion.  Then $GJ=\id_D$.
The components $1_d$ for $d\ne b$ and $v$ at $b$ define a
natural transformation $JG\Rightarrow\id_C$.  For an incoming
$g$, its naturality equation is $vug=eg=g$; for an outgoing
$f$, it is the definition $G(f)=fv$.  For a nonunit $h$ it is
$vuhv=ehv=hv$, while for a unit it is $v=hv$.
Thus $BJ$ and $BG$ are homotopy inverses.
Finally, if $e$ is a zero of $A$ and $h$ is a unit, then
$hg=heg=eg=g$, which verifies the stated sufficient condition.
Opposition proves the dual assertion.
\end{proof}

\begin{corollary}[Arbitrary exits from the eleven-arrow strong pair]
\label{C:eleven-strong-arbitrary-exit}
Let $S$ be the eleven-arrow strong two-object category in
Proposition~\ref{C:eleven-strong-three-exit}.  Adjoin a sink
with trivial endomorphism monoid and any nonempty finite
compatible family of outgoing arrows from $S$.
Every rational degree-two square of the resulting classifying
space is zero.  In particular it cannot be a projective plane.
The dual assertion holds for an adjoined source.
\end{corollary}
\begin{proof}
If the two-dimensional row and column sets occur at different
objects, the balanced-tensor computation in the cited proposition
gives $H^2(BS;\mathbb Q)=0$.  Every degree-two class of the
whole act cone is then relative to $BS$ and pulls back from
the suspension of the action category.  Its square is zero.

Otherwise both two-dimensional sets occur at one object $a$,
with corner $P_5$, while the other object $b$ has corner
$\{1_b,e\}$.  Every incoming cross arrow at $b$ comes from
$a$ and belongs to the rectangular ideal, so is fixed by $e$.
The rectangular multiplication supplies $u,v$ with $vu=e$.
Lemma~\ref{C:one-sided-corner-deletion} deletes $b$ without
imposing any condition on the outgoing arrows to the new sink.
The resulting category is a finite act cone over $P_5$.
Theorem~\ref{C:two-row-act-cups} makes all its degree-two
squares zero, with no bound on the number of states.
\end{proof}

\subsection{Structural reductions at fifteen arrows}

\begin{lemma}\label{C:connected-tree-budget}
A connected finite category on $k$ objects has at least $2k-1$
arrows.  If equality holds, its classifying space is contractible.
Consequently a fifteen-arrow category realizing $\mathbb{CP}^2$
has at most seven objects.
\end{lemma}
\begin{proof}
There are $k$ identity arrows, and the connected underlying
undirected graph needs $k-1$ cross arrows on a spanning tree.
In the equality case these are all arrows: every tree edge has
one orientation, there are no other endomorphisms, and there
are no parallel or reverse cross arrows.  A composable pair
of nonidentity arrows would force a third edge on a cycle of
this tree.  Thus no such pair exists.  The ordinary nerve is
the tree itself and is contractible.  For $k\geq8$ the stated
fifteen-arrow consequence follows from the inequality and its
equality case.
\end{proof}

\begin{lemma}\label{C:nine-bounded-rank-one}
Let $|M|\leq9$, let $S$ be a nonempty finite right $M$-set,
and suppose the act cone $X=B\mathcal C(M,S)$ has bounded rational
homology and $\dim H^2(X;\mathbb Q)\leq1$.  Every degree-two square
is zero.  No assumption on $H^1(X;\mathbb Q)$ is needed.
\end{lemma}
\begin{proof}
This records exactly the obstructions in the proof of
Theorem~\ref{C:nine-act-cones}, without its unnecessary target
specialization.  Through order eight use
Lemma~\ref{C:eight-bounded-cups}.  At order nine, zero vertex
$H^2$ makes every degree-two class relative and its square zero.
Two minimal rows give Theorem~\ref{C:two-row-act-cups}.
The projective-dimension-two cases have $H_4(X;\mathbb Q)=0$.
The exact periodic and stable-retract cases make nonzero fourth
homology recur in arbitrarily high degrees, contradicting the
stated bound.  The seven remaining row modules split into
characters, so Proposition~\ref{C:split-row-act-cups} excludes a
nonzero square when $H^2$ has dimension at most one.
These are all rows of the independently verified order-nine
partition.  None of these arguments uses vanishing first homology.
\end{proof}

\begin{proposition}\label{C:fifteen-three-SCC}
A three-object category with at most fifteen arrows and a
nontrivial strongly connected component does not realize
$\mathbb{CP}^2$.
\end{proposition}
\begin{proof}
Deleting an object which is a retract of another removes its
identity and at least two opposite cross arrows, hence leaves
at most twelve arrows.  The proved smaller exclusion applies.
Assume therefore that no such retract exists.

If all three objects form a strongly connected component,
the categorical Rees description gives finite-group homotopy
type when all row counts or all column counts are one.
Otherwise both total counts are at least four, so the minimal
ideal alone has at least sixteen arrows.  Both alternatives
exclude the proposed type.

It remains to consider a strong two-object block $S$, a singleton
monoid $M$, and cross arrows in one direction.  Write $s,m,t$
for their sizes.  Then $s\geq6$, $t\geq2$, and $s+m+t\leq15$.
If $t=2$, there is one cross arrow at each object of $S$.
The category is the categorical join of $S$ with $M$ (in one
of the two orders).  Its classifying space is a join and has
zero positive-degree rational products.
We may suppose $t\geq3$.

If $BS$ is not of finite-group type, its minimal ideal has at
least nine arrows and the two identities lie outside it.
Thus $s\geq11$, forcing $(s,m,t)=(11,1,3)$.
The minimal group is trivial, the total row and column counts
are three, and only the two identities lie outside the ideal.
Proposition~\ref{C:eleven-strong-three-exit}, or its opposite,
excludes this case.

Otherwise $BS$ is rationally acyclic.  If $m\leq4$, both
vertex-block spaces in the correspondence-bar pushout are
rationally acyclic.  Collapsing them to points gives a rational
cohomology-ring equivalence to a suspension, whose products vanish.
If $m\geq5$, the inequalities $t\geq3$ and $s+m+t\leq15$
give $s\leq7$ and $m\leq6$.  As in
Lemma~\ref{C:fourteen-three-SCC}, the minimal Rees data of $S$
have trivial group and one row and one column at each object.
Both constant rational $S$-modules are projective.
The actual orbit-collapse functor therefore identifies the
whole category rationally, including its cohomology ring,
with a finite act cone over $M$.  Its rational homology is
bounded and its second cohomology is one-dimensional if the
original category has the proposed type.  Apply
Lemma~\ref{C:nine-bounded-rank-one}.
\end{proof}

\begin{lemma}[Star reductions at fifteen arrows]
\label{C:fifteen-star-deletions}
For an outward three-object star of total size fifteen, let
$(m,a,b,x,y)$ denote the central and leaf endomorphism orders
and the two leg sizes.  The generic star exclusions in
Lemma~\ref{C:generic-star-deletions} still hold, with each
bound eight there replaced by nine.  There is also an exclusion
when $m\leq4$ and $x,y\leq2$.  The same assertions hold for
inward stars.  If every vertex has zero rational $H^2$, the
star is excluded as before.
\end{lemma}
\begin{proof}
The singleton-leg, join, and relative-product arguments in the
earlier lemma do not depend on its total-size bound.  For a trivial
center the nerve is a wedge of the two proper collages; each has
at most thirteen arrows, because the other leaf and leg each
contribute at least one arrow.  A wedge of connected spaces of
type $\mathbb{CP}^2$ must have a projective-plane summand: the
free-product formula for $\pi_1$, the direct sum in reduced
homology, and the primitive nonzero degree-two square show
that the other summand is simply connected and acyclic, hence
contractible.  The smaller-category exclusion rules this out.

For $m\leq9$ and $a,b\leq2$, exact rational coinvariants at
the two leaves replace them by trivial monoids and replace the
legs by two actual nonempty finite central acts $S_1,S_2$.
The quotient functor is a rational homology equivalence on each
vertex and edge bar, hence an isomorphism of rational cohomology
rings on the star.  Identify the two sink vertices of this new
star.  On nerves this gives precisely the act cone for
$S_1\sqcup S_2$, since no nonidentity arrow exits a sink.
On realizations it identifies two points in a connected CW space;
the resulting space is homotopy equivalent to the star wedged
with $S^1$.  Thus a hypothetical projective-plane star would
give an act cone with bounded rational homology, one-dimensional
$H^2$, and a nonzero square.  This contradicts
Lemma~\ref{C:nine-bounded-rank-one}.

For the earlier distinguished-leaf case with $m\leq2$, other
leaf order at most four, and distinguished leaf order at most
nine, collapse the rationally acyclic central vertex.
This is a rational cohomology-ring equivalence, and its quotient
is the wedge of the two collage cofibers.  The other collage
has two rationally acyclic vertices and zero products.
Exact central coinvariants turn the distinguished collage
rationally into a finite act cone over its leaf monoid.
Bounded rational homology and dimension at most one for its
$H^2$ follow from the wedge decomposition of a hypothetical
projective-plane star.  The same lemma excludes its square.

Finally, if $m\leq4$ and both legs have at most two arrows,
collapse the rationally acyclic center in the same way.
Each of the two collage summands has a rationally acyclic
vertex and at most two cross arrows, so
Lemma~\ref{C:two-object-cup-pruning} gives zero positive products
on each summand.  Their wedge has zero products as well.
Opposition proves the inward assertions.
\end{proof}

\subsection{The completed fifteen-arrow star calculation}

\begin{theorem}\label{C:fifteen-all-stars}
A three-object category with star support and at most fifteen
arrows cannot have classifying space homotopy equivalent to
$\mathbb {CP}^2$.
\end{theorem}
\begin{proof}
Theorem~\ref{C:CP2-lower15} handles totals below fifteen.
Pass to an opposite category if necessary, so that the star is
outward.  Write $(m,a,b,x,y)$ for its central endomorphism order,
the two leaf endomorphism orders, and the two leg sizes.  All
five integers are positive and have sum fifteen.  Relabel the
leaves so that $(a,x)\leq(b,y)$ lexicographically.
Lemma~\ref{C:fifteen-star-deletions}, together with its preceding
generic singleton and relative-product reductions, leaves exactly
the following profiles:
\[
\begin{array}{c|l}
m&(a,b,x,y)\\\hline
3&(1,5,1,5),(1,5,2,4),(1,5,3,3),(1,6,1,4),\\
 &(1,6,2,3),(1,7,1,3),(2,5,1,4),(2,5,2,3),\\
 &(2,6,1,3),(3,5,1,3)\\[2pt]
4&(1,5,1,4),(1,5,2,3),(1,6,1,3),(2,5,1,3)\\[2pt]
5&(1,3,2,4),(1,3,3,3),(1,3,4,2),(1,4,2,3),\\
 &(1,4,3,2),(1,5,1,3),(1,5,2,2),(1,6,1,2),\\
 &(2,3,2,3),(2,3,3,2),(2,4,2,2),(2,5,1,2),\\
 &(3,3,2,2)\\[2pt]
6&(1,3,2,3),(1,3,3,2),(1,4,2,2),(1,5,1,2),\\
 &(2,3,2,2)\\[2pt]
7&(1,3,2,2).
\end{array}
\]
This list is regenerated independently from the five positive
integers by the verification program below.  In particular, a
profile is not omitted merely because a previously searched
profile looked similar.

At least one vertex must have nonzero rational second cohomology:
otherwise every degree-two class is relative to the three vertex
spaces, and the star quotient by those spaces is a wedge of
suspensions.  Its positive products vanish.  At orders at most
four all vertex spaces are rationally acyclic.  At orders five,
six, and seven, use the complete positive-$H^2$ catalogs already
proved in the small-monoid classification.

There is a necessary distinction in using these catalogs.  In
twenty-eight of the displayed profiles only one corner can have
positive $H^2$; that corner ranges over its complete positive
catalog and the other corners range over every ordinary monoid
of their stated orders.  In the other five profiles,
\[
 (5,1,5,1,3),\quad(5,1,5,2,2),\quad(5,1,6,1,2),\quad
 (5,2,5,1,2),\quad(6,1,5,1,2),
\]
both potentially positive corners range over \emph{every} ordinary
monoid of their orders.  Only a choice with zero rational $H^2$
at every vertex is discarded.  Thus a positive corner paired with
a rationally acyclic corner of the same cardinality is retained.
The positivity test is the exact minimal-band augmentation
pairing; for the two-by-two band it tests equality of the two
one-dimensional action characters.  The complete small-monoid
classification justifies this test in the stated size range.

For fixed vertex tables, each leg is an arbitrary commuting
biset.  The central right action is enumerated up to exact state
conjugacy.  Every state conjugate of every action of the opposite
leaf monoid is then restored, and every commutation equation is
checked on all pairs of monoid elements.  These independent
state changes may be made separately on the two legs.  There
is no composition between different legs of an outward star,
so taking every pair of the resulting leg bisets imposes all
and only the categorical associativity conditions.  This proves
the completeness of the action enumeration.  It is implemented
in \path{code/c_fifteen_star_residuals.py}; no identification of
relative left-action positions is made after fixing a central
right action.

The complete calculation has $381897$ literal categories.
The mutually exclusive outcomes, with the homology tests taken
in the producer's stated order, are
\[
\begin{array}{l|r}
\dim H_2(-;\mathbb F_2)\ne1&182559\\
H_1(-;\mathbb F_2)\ne0&109749\\
\text{an explicit primitive for the degree-two square}&89004\\
\text{nonzero square and nonzero }H_3(-;\mathbb F_2)&585.
\end{array}
\]
All categories in the final line have independently checked
positive third homology.  A projective plane satisfies none of
these alternatives.

The complete primary summary and literal certificate stream are
\path{code/outputs/c_fifteen_star_residuals.json} and
\path{code/outputs/c_fifteen_star_residuals.jsonl}.
The independent verifier
\path{code/d_verify_fifteen_stars_stream.py} imports neither the
table producer nor its cohomology routines.  It reconstructs
the typed bar chains and matrices from each literal category,
checks all $189231356$ typed associativity equations, recomputes
the first and second mod-two homology tests, and checks the
degree-two cocycles and nonboundary generators.  Each of the
$89004$ primitives is checked on every normalized quadruple,
for $129159493$ exact coboundary equations in total.  It separately
recomputes the positive third homology in all $585$ remaining
cases.  Its final output is
\path{code/outputs/d_verify_fifteen_stars.json}, with status
\texttt{PASS}; its exact profile list, all row counts, and final
stream hash agree with the completed primary stream.  The
independent gate certifies the complete cardinality scope and
every saved table and cochain.  Completeness of the finite
action search is supplied by the explicit enumeration argument
above; it is not inferred from the number of saved certificates.
This exhausts all fifteen-arrow star categories, and opposition
gives the inward case.
\end{proof}

\subsection{Large corners with four objects and fifteen arrows}

\begin{proposition}\label{C:fifteen-four-large-corner}
A four-object category with at most fifteen arrows, one of whose
endomorphism monoids has at least eight elements, cannot realize
$\mathbb{CP}^2$.
\end{proposition}
\begin{proof}
Suppose the category is connected, since the proposed type is.
Write $v$ for the large corner, of order $q\geq8$.
Deleting a retract removes at least its identity and two opposite
cross arrows, leaving at most twelve arrows and a previously
excluded category.  Assume no retract remains.

The underlying connected graph on four objects needs at least
three cross arrows.  A nontrivial strongly connected component
needs at least one further arrow.  The resulting budget
$q+3+4\leq15$ forces $q=8$ and makes all other endomorphism
monoids trivial.  Any two objects in that component then include
one with trivial endomorphism monoid; a pair of opposite arrows
exhibits it as a retract.  This is impossible.  Thus the support
is acyclic.  There are at most four cross arrows, and the other
endomorphism orders sum to at most four.

If $v$ has both an incoming and an outgoing cross arrow, those
arrows and their composite occupy a three-object chain.  Connecting
the fourth object uses the fourth cross arrow.  Therefore all
four hom-sets are singletons, $q=8$, and all other corners are
trivial.  The fourth object can attach only by an arrow from the
bottom of that chain or by an arrow to its top: attaching to its
middle, or orienting either extreme attachment the other way,
would force an additional composite hom-set.  In the first case
the bottom is initial, and in the second the top is terminal.
The category is contractible.

By opposition we may therefore suppose $v$ is a sink.  Let $D$
be the full subcategory on the other three objects.  Its corners
have orders at most two, hence rationally semisimple algebras.
If some incoming hom-set to $v$ is a singleton, then $D$ has at
most six arrows.  We claim $H^2(BD;\mathbb Q)=0$.
Unless its acyclic support contains a three-object chain, the
relative bar over the semisimple diagonal has length at most one.
If it does contain that chain, the three required cross arrows
and three identities exhaust its budget; it is the ordinary
three-element chain.  This proves the claim in both cases,
also allowing disconnected $D$.
The singleton incoming arrow is fixed by the whole large monoid,
so Lemma~\ref{C:fixed-incoming-cups} makes every degree-two
product zero.

Finally suppose every nonempty hom-set into $v$ has at least
two elements.  There can be only one neighboring object $w$:
two neighbors would use four cross arrows and still leave the
fourth object disconnected.  The remaining subcategory $D$ is
connected and needs two cross arrows.  Thus $q=8$, every corner
of $D$ is trivial, the hom-set from $w$ to $v$ has two elements,
and $D$ has exactly two singleton cross arrows.  Its support is
a three-object star.  No arrow can enter $w$, since composing it
with $w\to v$ would create another incoming hom-set to $v$.
If $w$ is the central source of the star, delete its two leaves
by the singleton cofinal deletion.  If $w$ is a source leaf of
an inward star, first delete the other source leaf and then the
central sink.  Each deletion has its singleton adjacent trivial
corner, and neither deleted vertex is adjacent to $v$.
The result is the two-object act cone with two cross arrows over
the large monoid.  These are actual homotopy-equivalence deletions.
Lemma~\ref{C:two-object-cup-pruning} gives zero products, which
again contradicts the projective-plane square.
\end{proof}

\subsection{Three reductions for four-object stars}

\begin{proposition}\label{C:fifteen-four-star-prunes}
Let $C$ be a four-object outward or inward star category with
at most fifteen arrows.  Write $M$ for its central endomorphism
monoid, $A_1,A_2,A_3$ for the leaf monoids, and $s_1,s_2,s_3$
for the nonempty leg cardinalities.  Each of the following
conditions excludes $BC\simeq\mathbb {CP}^2$:
\begin{enumerate}
\item Every $|A_i|\leq4$, and some $s_i=1$.
\item $|M|\leq9$ and every $|A_i|\leq2$.
\item At least two of $s_1,s_2,s_3$ equal one.
\end{enumerate}
\end{proposition}
\begin{proof}
Opposition interchanges inward and outward stars without changing
the classifying-space homotopy type.  For (1), take the inward
orientation.  The full subcategory $D$ on the leaves is a
disjoint union of three monoids of orders at most four, so
$H^2(BD;\mathbb Q)=0$.  A singleton incoming hom-set supplies
an arrow fixed by every element of $M$.
Lemma~\ref{C:fixed-incoming-cups}, which allows disconnected
$D$, makes all degree-two products zero.

For (2), take the outward orientation.  Every rational leaf
coinvariant functor is exact.  The actual orbit-collapse
functors replace the three leaves by trivial monoids and the
three legs by nonempty finite right $M$-sets $S_1,S_2,S_3$.
On each edge correspondence bar and each vertex they are
rational homology equivalences.  Hence they give a rational
cohomology-ring equivalence from $BC$ to the resulting star
$Z$.  Identify its three sink points.  Since no nonidentity
arrow exits a sink, the resulting nerve is exactly the act
cone for $S_1\sqcup S_2\sqcup S_3$.  Identifying three points
of a connected CW space gives the homotopy type
\[
 Z\vee S^1\vee S^1.
\]
If $BC\simeq\mathbb {CP}^2$, this act cone has bounded rational
homology, one-dimensional rational $H^2$, and a nonzero square.
Lemma~\ref{C:nine-bounded-rank-one} excludes this possibility;
the additional first homology is explicitly allowed there.

For (3), take the outward orientation and suppose $s_1=s_2=1$.
Put $X=BM$, and let $Y_i$ be the classifying space of the full
two-object subcategory on the center and leaf $i$.  The star
is the homotopy union of these three spaces along $X$; the
nerve inclusions are subcomplex inclusions.  A singleton leg
makes $Y_i$ the join $X*BA_i$.  Thus the maps $X\to Y_1$ and
$X\to Y_2$ are null-homotopic.  The ordinary homotopy-pushout
formula, first for these two maps and then for the third,
gives an actual homotopy equivalence
\begin{equation}\label{C:four-star-two-singleton-splitting}
 BC\simeq Y_1\vee Y_2\vee\Sigma X
       \vee\operatorname{hocofib}(X\longrightarrow Y_3).
\end{equation}
For the second pushout the common map from $X$ to
$Y_1\cup_X^hY_2$ is still null, since it factors through
the null map to $Y_1$.  This justifies the displayed splitting
without a choice of compatible inverse equivalences.

Suppose the left side were a projective plane.  The reduced
integral homology of the wedge on the right is the direct sum
of that of its four summands.  Its free-product fundamental
group is trivial, so each summand is simply connected.  The
unique degree-two integral generator and its primitive
nonzero square must belong to the same wedge summand; that
summand contains both nonzero homology groups, in degrees
two and four.  A suspension has zero positive cup products,
so $\Sigma X$ cannot be that summand.  It is therefore integrally
acyclic.  Since $X$ is connected, suspension homology shows
that $X$ is integrally acyclic as well.

Theorem~\ref{A:bounded} now makes the finite monoid space
$X=BM$ simply connected.  The homological Whitehead theorem
gives $X\simeq *$.  Consequently the two joins $Y_1,Y_2$ are
contractible and
$\operatorname{hocofib}(X\to Y_3)\simeq Y_3$.
Equation~\ref{C:four-star-two-singleton-splitting} reduces the
whole category to its center--third-leaf subcategory.  Its
number of arrows is at most
\[
 15-(|A_1|+1)-(|A_2|+1)\leq11.
\]
The already proved smaller-category exclusion is a contradiction.
This proves (3) and completes all three reductions.
\end{proof}

\subsection{Five-object acyclic categories through twenty-one arrows}

\begin{proposition}\label{C:five-object-acyclic21}
Let $C$ have five objects, trivial endomorphism monoids, and
acyclic support.  If $C$ has at most twenty-one arrows, then
$BC\not\simeq\mathbb {CP}^2$.
\end{proposition}
\begin{proof}
The normalized nerve has no simplices above dimension four.
If it has no four-simplex, its fourth homology is zero and the
claim follows.  Otherwise a four-simplex visits all five objects,
which can therefore be labeled $0<1<2<3<4$.  Transitivity
forces every forward hom-set to be nonempty.  These are all
the cross-hom sets, since the support is acyclic.  Put
$h_{ij}=|C(i,j)|$ for $i<j$.

There are three immediate necessary conditions for a nonzero
four-cycle.  If $h_{01}=1$, the face deleting object $0$ maps
the basis of $C_4(NC)$ injectively into the summand of
$C_3(NC)$ on objects $1,2,3,4$.  None of the other faces
lands in that summand, so $\partial_4$ is injective.  The
same argument with the last face applies when $h_{34}=1$.
If both $h_{12}=h_{23}=1$, the face deleting object $2$
replaces their unique arrows by their fixed composite in
$C(1,3)$.  It again sends the four-simplex basis injectively
into its own three-simplex summand.  Thus nonzero fourth
homology requires
\[
 h_{01}\geq2,\qquad h_{34}\geq2,\qquad
 h_{12}>1\ \text{or}\ h_{23}>1.
\]
In particular, there must be at least eighteen arrows:
five identities, ten nonempty forward hom-sets, and these
three additional arrows.  This argument works over every
coefficient ring and does not assume that the category is
a poset.

For totals eighteen through twenty-one, enumerate every positive
ten-tuple $(h_{ij})_{i<j}$ of the required sum satisfying the
displayed conditions.  Prescribe the five identities, let every
typed product range over its entire possible hom-set, and
propagate every associativity triple.  There are no other
structural profile discards.  Permutations within hom-sets
are used only to remove repeated completed table representatives.
The exact generator is
\path{code/d_five_object_line_upper.cpp}.  Its full completed
counts, followed by the independent bar check, are
\[
\begin{array}{c|r|r|r|r}
\text{arrows}&\text{profiles}&\text{associative tables}
 &H_1=0,\ \dim H_2=1\text{ orbits}&\text{zero squares}\\\hline
18&2&2&0&0\\
19&19&67&4&4\\
20&100&1426&56&56\\
21&385&27546&472&472.
\end{array}
\]
The first and second homology tests in this table are over
$\mathbb F_2$.  For each of the $532$ surviving literal
hom-set-permutation orbits,
\path{code/c_generic_hom_cups.py} independently verifies all
typed associativity equations, reconstructs the normalized
bar complex, checks a nonboundary degree-two cocycle, and
checks an explicit primitive for its square.  There are
$13382$ typed associativity checks in this gate.  Hence no
survivor has the mod-two cohomology ring of a projective plane.

The output stems are \path{code/outputs/d_line5_18} through
\path{code/outputs/d_line5_21}; the table streams, complete
profile logs, and literal cochain certificates are retained.
The separate coverage program
\path{code/d_five_acyclic_scope.py} regenerates every required
ten-tuple directly and checks equality with the completed
profile lists, not just their cardinalities.  It also checks
all generator and checker completion sentinels, exact row
counts, and final byte positions.  Its record
\path{code/outputs/d_five_acyclic_scope.json} has status
\texttt{PASS}.  Combined with the top-cycle argument, these
finite checks cover every category in the stated scope.
\end{proof}

\subsection{All acyclic categories through twenty-one arrows}

\begin{lemma}[All degree-two products on the two small five-chain levels]
\label{C:five-chain-nineteen-products}
Let $D$ have five linearly ordered objects, trivial endomorphisms,
and a nonempty hom-set in every forward direction.  If $D$ has
at most nineteen arrows, then every product of two classes in
$H^2(BD;\mathbb F_2)$ is zero.  No assumption on its first homology
or the dimension of its second homology is required.
\end{lemma}
\begin{proof}
The top-face argument in Proposition~\ref{C:five-object-acyclic21}
shows that $H_4(BD;\mathbb F_2)=0$ unless
\[
 h_{01},h_{34}\geq2,\qquad h_{12}>1\text{ or }h_{23}>1.
\]
These conditions already require eighteen arrows.  At totals
eighteen and nineteen there are respectively two and nineteen
positive cardinality profiles satisfying them.

The generator \path{code/d_five_object_line_all.cpp} enumerates
all completed tables on these profiles, with no homology filter.
It gives two and sixty-seven labelled associative tables, or two
and thirty-one orbits under permutations within hom-sets.
The standalone checker \path{code/d_five_chain_all_cups.py}
checks every typed composition law and every bar differential,
constructs a full basis of second mod-two cohomology, and finds
an explicit degree-three primitive for every ordered product
of basis classes.  Thus it verifies all products, not just the
square in a previously selected one-dimensional case.

There is also an independent enumeration with no associativity
propagation: \path{code/d_five_chain_independent.py} forms every
Cartesian product of the allowed values of the unprescribed
typed products, then checks all associativity equations directly.
The two levels have respectively two and $167$ raw assignments,
of which two and sixty-seven are associative.  Their entire
hom-set-permutation orbit sets agree with the completed streams.
Every one of these sixty-nine labelled tables is then subjected
to the full-basis cup checker.  This fresh check verifies $9280$
typed associativity equations and all eight ordered basis
products occurring before removal of repeated table orbits.
Every product has the stored explicit primitive.  The final
record \path{code/outputs/d_five_chain_independent.json} has
status \texttt{PASS}; the separate orbit-level record is
\path{code/outputs/d_five_chain_all_cups.json}.
Linearity of the cup product now proves the assertion for all
classes.  The argument uses only mod-two cohomology and makes
no inference about divisibility of integral cup products.
\end{proof}

\begin{theorem}[Acyclic support with trivial endomorphisms]
\label{C:all-acyclic21}
Let $C$ be a finite category with trivial endomorphism monoids
and acyclic support on distinct objects.  If $C$ has at most
twenty-one arrows, then $BC$ is not homotopy equivalent to
$\mathbb {CP}^2$.
\end{theorem}
\begin{proof}
We may assume that $C$ is connected.  Write $k$ for its number
of objects.  If its nerve has no four-simplex, fourth homology
vanishes.  Otherwise fix a four-simplex, and let $D$ be the
full subcategory on its five objects.  Those objects form a
linear order, and all ten forward hom-sets are nonempty.
If $k=5$, Proposition~\ref{C:five-object-acyclic21} applies;
if $k<5$, no four-simplex exists.  We treat the remaining
object counts by their cross-arrow budget.

Suppose a second four-simplex uses $r$ objects outside $D$.
For $1\leq r\leq4$, its five-object chain forces at least
\[
 r(5-r)+\binom r2
\]
new cross-hom sets involving those $r$ objects.  Its union
with the original chain is connected because they share an
object.  Connecting the other $k-5-r$ objects therefore
requires at least $k-5-r$ further cross arrows.  Including
the ten original cross arrows and the $k$ identities, the
total is at least
\begin{equation}\label{C:outside-four-simplex-cost}
 k+10+r(5-r)+\binom r2+(k-5-r)
 =2k+5+\frac{7r-r^2}{2}.
\end{equation}
For $r=1,2,3,4$ these bounds are respectively
$2k+8,2k+10,2k+11,2k+11$.  If $r=5$, the two chains are
disjoint.  The second chain contributes ten new cross arrows;
one further arrow is needed to connect the two components,
and the remaining $k-10$ objects need at least that many
additional arrows.  The total is then at least $2k+11$.
These estimates count only required nonempty hom-sets, so
parallel arrows can only increase them.

For $k\geq7$, all these bounds exceed twenty-one.  Hence
every four-simplex lies in $D$.  Every five-simplex would
have a four-dimensional face using an outside object, so
there are no five-simplices.  Thus the normalized chain
groups and boundary maps give an actual isomorphism
\[
 H_4(BD;\mathbb F_2)\longrightarrow H_4(BC;\mathbb F_2).
\]
Connecting the $k-5$ outside objects costs at least $k-5$
cross arrows as well as their identities.  Consequently
\[
 |D|\leq21-2(k-5)\leq17.
\]
The top-face argument in
Proposition~\ref{C:five-object-acyclic21} makes $H_4(BD)$
zero at this size, a contradiction.

It remains to consider $k=6$, with outside object $v$.
First suppose that $v$ lies in a four-simplex.  It is
comparable in the support order with at least four of the
five objects of $D$.  If it is comparable with all five,
the six objects form a chain.  The fifteen required cross
arrows and six identities exhaust the budget, so $C$ is
the ordinary six-element chain poset and is contractible.

Otherwise $v$ is incomparable with exactly one object $a$
of $D$.  Transitivity forces $v$ and $a$ to have the same
order position relative to the other four objects.  The
support is therefore a five-level order with one doubled
level.  It has fourteen required cross arrows and six
identities, leaving room for at most one additional
parallel arrow.  Its only five-object chains are $D$ and
the chain obtained by replacing $a$ with $v$.

The normalized four-chains consequently split into two
summands, one for each choice at the doubled level.  In
either summand, choose a face which deletes neither its
exclusive vertex nor either endpoint of the possible
parallel adjacent arrow.  At most three of the five
vertices are forbidden, so such a face exists.  It retains
the exclusive vertex, hence its target three-chain summand
receives no boundary contribution from the other branch.
If the parallel arrow occurs among the four adjacent arrows,
the selected face also retains its label; all other arrows
are unique.  This face map is therefore injective on the
four-simplex basis.  If the extra arrow is not adjacent in
this chain, that basis has one element and the same conclusion
holds.  Projecting the boundary to the two selected face
summands proves that $\partial_4$ is injective.  Again
$H_4(BC;\mathbb F_2)=0$.

Finally suppose that $v$ lies in no four-simplex.  All
four-simplices and all their faces are literally in $D$,
and there are no five-simplices.  Inclusion again induces
an isomorphism on fourth mod-two homology, hence restriction
induces an isomorphism
\[
 H^4(BC;\mathbb F_2)\longrightarrow H^4(BD;\mathbb F_2).
\]
The outside object requires its identity and at least one
cross arrow, so $|D|\leq19$.  By
Lemma~\ref{C:five-chain-nineteen-products}, every product
of degree-two classes on $D$ is zero.  Naturality and this
fourth-cohomology isomorphism imply that every degree-two
product on $C$ is zero as well.  This contradicts the
nonzero square of the generator of
$H^2(\mathbb {CP}^2;\mathbb F_2)$ and completes all cases.
\end{proof}

\begin{lemma}[Semisimple diagonal corners at fifteen arrows]
\label{C:semisimple-fifteen}
Let $C$ be a connected finite category with at most fifteen
arrows, at least four objects, and acyclic object support.
If every rational endomorphism algebra is semisimple, then
$H_4(BC;\mathbb Q)=0$.
\end{lemma}
\begin{proof}
By Lemma~\ref{C:semisimple-diagonal-height}, it suffices to
bound the length of a chain of distinct objects by three.
Suppose instead that five objects occur on a four-step chain.
Transitivity supplies all ten forward cross-hom sets among them.
If $k$ is the total number of objects, contracting these five
vertices in the connected underlying undirected support graph
leaves $k-5$ other vertices.  At least $k-5$ additional cross
arrows are necessary to connect them.  With the $k$ identities,
the total arrow count is therefore at least
\[
 10+(k-5)+k=2k+5.
\]
The assumed budget forces $k=5$ and equality throughout.
Then every endomorphism is an identity, every forward hom-set
is a singleton, and there are no other arrows.  This is the
five-element chain poset, whose nerve is contractible.
In every other case the relative bar has length at most three,
and the cited semisimple-diagonal argument gives the claimed
vanishing.  No assertion is made for cyclic object support.
\end{proof}

\subsection{A universal obstruction for the two-circle chart family}

\begin{lemma}[A right-zero action has a natural graph model]
\label{C:right-zero-action-graph}
Let $R=\{1,p,q\}$, with $ab=b$ for $a,b\in\{p,q\}$,
act on the left on a finite set $T$.  Its two nonidentity
transformations $P,Q$ are retractions onto the same subset
$U\subseteq T$.  The action category has a natural classifying
space equivalence to the graph with vertices $U$ and, for each
$t\in T\setminus U$, one oriented edge from $P(t)$ to $Q(t)$.
An equivariant set map sends this edge to the edge indexed by
the image of $t$, or collapses it to a vertex if that image
belongs to the new common retraction image.
\end{lemma}
\begin{proof}
The equations $PQ=Q$ and $QP=P$ show that
$\operatorname{im}P=\operatorname{im}Q=U$, and both maps
are the identity on $U$.  In the action category every vertex
of $U$ has endomorphism monoid $R$.  A vertex outside $U$ has
only its identity and two outgoing arrows, labeled $p$ and $q$,
to $P(t)$ and $Q(t)$.  There are no other arrows between
different objects.  Postcomposition fixes each of these two
arrow labels, since $ap=p$ and $aq=q$ for $a\in R$.

Consequently its classifying space is the homotopy colimit of
a graph of spaces: the vertices in $U$ carry $BR$, outside
vertices carry points, and each of the two incident edges at
an outside vertex carries $BR$, mapping identically to its
$U$-endpoint and constantly to the outside point.  The space
$BR$ is contractible, since $R$ has a right zero.  Collapsing
each copy of $BR$ gives an equivalence to the indicated
bipartite graph.  Removing its bivalent outside vertices gives
the graph in the statement.

This equivalence can be chosen naturally, rather than merely
separately for each action.  A nerve simplex either stays at a
vertex of $U$, or has one initial arrow $p$ or $q$ from an
outside vertex, followed by endomorphisms of its endpoint.
Collapse the first kind to that vertex; send the second along
the corresponding half-edge, with position equal to the sum
of its endpoint barycentric coordinates.  The right-zero
identities make this compatible with every face and degeneracy.
An equivariant map preserves that half-edge label, unless it
sends its outside vertex into $U$, in which case both sides
collapse it.  This proves the asserted naturality.
\end{proof}

\begin{lemma}[Two commuting right-zero actions]
\label{C:two-retraction-cubical-model}
Suppose that $P_0,P_1$ and $Q_0,Q_1$ are two commuting pairs
of retractions on a finite set $T$, each pair having a common
image.  Put $U=\operatorname{im}P_0=\operatorname{im}P_1$
and $V=\operatorname{im}Q_0=\operatorname{im}Q_1$.
The double action correspondence has a natural finite CW model
$K(T)$ with cells
\[
\begin{array}{c|l}
0&U\cap V\\
1&U\setminus V\quad\text{and}\quad V\setminus U\\
2&T\setminus(U\cup V).
\end{array}
\]
Its two-cell indexed by $t$ has the four faces indexed by
$P_0t,P_1t,Q_0t,Q_1t$, with a face collapsed when its index
lies in $U\cap V$.  Equivariant maps preserve these cells
or collapse a coordinate.
\end{lemma}
\begin{proof}
First apply Lemma~\ref{C:right-zero-action-graph} to the $P$
action.  The $Q$ action commutes with it and therefore acts
on its graph by the natural cellular maps in that lemma.
Regard the graph as a one-dimensional simplicial set, including
its degenerate edges.  Its set of vertices is $U$, and its set
of edges together with degenerate edges is literally $T$.
Apply the same graph construction to the $Q$ action in each
simplicial degree.  Naturality makes this a simplicial graph.
Iterated realization and the two action bars commute, and
the levelwise graph equivalences give the asserted equivalence.
The degeneracies are subcomplex inclusions, so this use of
levelwise equivalences is valid.

In degree zero the graph has vertices $U\cap V$ and edges
$U\setminus V$.  In degree one, the new vertices are
$V\setminus U$ and the new edges are
$T\setminus(U\cup V)$.  Higher simplicial degrees are
degenerate copies of these.  Realization therefore gives
exactly the displayed finite two-dimensional cell structure.
The four boundary maps are the original four retractions;
their corner equations are precisely their commutation
relations.  This also proves the stated cellwise naturality.
\end{proof}

\begin{theorem}[The two-circle chart family has zero degree-two products]
\label{C:affine-chart-family-zero}
Let $C$ have objects $0,1,2$ and only forward cross arrows.
Let $\operatorname{End}(0)$ be the identity-adjoined two-element
left-zero band,
$\operatorname{End}(1)=1$, and let
$\operatorname{End}(2)$ be the identity-adjoined two-element
right-zero band.  Thus both these corner monoids have order three.
Suppose
\[
 C(0,1)=\{a_0,a_1\},\qquad C(1,2)=\{b_0,b_1\},
\]
with every nonidentity element of the first corner sending
both $a_i$ to $a_0$ by precomposition, and every nonidentity
element of the last corner sending both $b_i$ to $b_0$ by
postcomposition.  The set $C(0,2)$ may be any finite nonempty
biset, and composition may be any equivariant map.
Then every rational product of degree-two classes on $BC$ is
zero.  In particular this family never realizes
$\mathbb {CP}^2$, regardless of its total number of arrows.
\end{theorem}
\begin{proof}
Put $T=C(0,2)$ and $t=b_1a_1$.  Precomposition gives two
retractions $P_0,P_1$ with common image $U$, and postcomposition
gives two retractions $Q_0,Q_1$ with common image $V$.
They commute.  Equivariance says
\[
 P_0t=P_1t=b_1a_0,\qquad
 Q_0t=Q_1t=b_0a_1.
\]
These equations determine the entire composition map from $t$.

The bar decomposition by chains of distinct objects expresses
$BC$ as a triangle of spaces.  Its vertex spaces are the three
corner classifying spaces; all are contractible.  The two
adjacent edge correspondences have the graph models
$A=B=S^1$: each has one fixed state and one transient state
whose two retraction images agree.  The long edge has the
model $K=K(T)$ of Lemma~\ref{C:two-retraction-cubical-model}.
Because the middle corner is trivial, the face correspondence
is the product $P=A\times B=T^2$.  Its three face maps are
the two projections and the cellular map
\[
 h:A\times B\longrightarrow K
\]
induced by composition.  These comparisons are compatible
with the face maps by the naturality in the two lemmas.
Thus replacing the original bar diagram by this diagram is
an objectwise equivalence, and gives the same homotopy colimit.
Equivalently, the triangle of spaces is obtained from the
three edge mapping cylinders by attaching
$P\times\Delta^2$ along its three faces.

The product of the two one-edge circle models has exactly
one two-cell.  Its image under $h$ is the two-cell indexed
by $t$, if $t\notin U\cup V$, and is degenerate otherwise.
In the first case its opposite faces agree by the two
displayed equations.  Therefore this cell is a cellular
two-cycle.  It represents a nonzero primitive class of
$H_2(K;\mathbb Z)$: its coefficient on its own basis cell
is one, and $K$ has no three-cells.

For clarity, the homology spectral sequence of the triangle
filtration has, in total degree four, only
$H_2(P)$ in triangle degree two.  Its first differential is
the map to the sum of the second homology of the three edges,
namely
\[
 h_*:H_2(T^2;\mathbb Q)\longrightarrow H_2(K;\mathbb Q),
\]
since the other two edges are circles.  No higher differential
can hit a nonzero vertex group in this degree.  The preceding
cell calculation makes $h_*$ injective.  Thus
$H_4(BC;\mathbb Q)=0$, and all degree-two products vanish.

In the remaining case, suppose for example that $t\in U$.
Both $P_i$ fix $t$, so the composition map is independent of
the $a$-coordinate.  On the natural graph models,
$h$ factors through the projection to $B$:
$h=\alpha\operatorname{pr}_B$ for a cellular map
$\alpha:B\to K$.  There is then an actual map from the
triangle of spaces to the unreduced suspension $\Sigma K$.
On the long edge use the identity suspension, on the $B$ edge
use $\Sigma\alpha$, and collapse the $A$ edge.  On the face
use the barycentric coordinate of the common endpoint of the
long and $B$ edges as the suspension coordinate, with space
coordinate $\alpha(b)$.  On the third edge that coordinate
is an endpoint of the suspension, so the definitions agree.
This gives a continuous map
\[
 BC\longrightarrow\Sigma K.
\]
It induces an isomorphism on rational degree-two cohomology.
Indeed the triangle spectral sequence identifies that group
with
\[
 \ker\bigl(H^1(A)\oplus H^1(B)\oplus H^1(K)
        \longrightarrow H^1(A\times B)\bigr).
\]
The first two coordinates map isomorphically to the two
summands of $H^1(A\times B)$, so projection of this kernel
onto $H^1(K)$ is an isomorphism.  The suspension map just
constructed induces exactly its inverse, up to the harmless
edge orientation signs.  There is no additional degree-two
term: the map from the edge degree-zero groups to the face
degree-zero group is surjective.  This remains true if $K$
is disconnected.

A suspension has zero positive cohomology products, so all
degree-two products of $BC$ vanish in this case as well.
If $t\in V$, interchange the two adjacent edges.  These
possibilities exhaust $T$ and prove the theorem.
\end{proof}

\begin{remark}[The completed sixteen-arrow control]
\label{C:affine-chart16-control}
For $|T|=5$, \path{code/d_affine_chart16.py} enumerates all
$1536$ ordered pairs of retractions onto a common subset.
They give $24$ right-action orbits under state permutations.
Restoring every relative left-action position gives $3374$
commuting bisets and $12437$ equivariant composition choices.
Exactly $2311$ literal tables have $H_1(-;\mathbb F_2)=0$
and one-dimensional $H_2$.  The independent bar checker
verifies all $249588$ typed associativity equations and an
explicit zero-square primitive for every one of these tables.
The complete records have prefix
\path{code/outputs/d_affine_chart16_CP2}.
This calculation is a control of the universal theorem, not
its proof.
\end{remark}

\subsection{The chart family with two-edge circles}

\begin{lemma}[Boundary values distinguish degenerate cells]
\label{C:retraction-boundary-injective}
For the commuting retractions in
Lemma~\ref{C:two-retraction-cubical-model}, the map
\[
 U\cup V\longrightarrow T^4,\qquad
 t\longmapsto(P_0t,P_1t,Q_0t,Q_1t)
\]
is injective.
\end{lemma}
\begin{proof}
Two elements of $U$ are distinguished by $P_0$, and two
elements of $V$ by $Q_0$.  If $x\in U$, $y\in V$ have the
same boundary tuple, then $P_0y=x$ and $Q_0x=y$.
Commutation gives
\[
 x=P_0Q_0x=Q_0P_0x=Q_0x=y.
\]
These cases exhaust the union.
\end{proof}

\begin{theorem}[The two-edge-circle chart family has zero products]
\label{C:affine-four-circle-family-zero}
Keep the corner monoids and support of
Theorem~\ref{C:affine-chart-family-zero}, but take the two
adjacent hom-sets to be
\[
 A=\{u_0,u_1,a_0,a_1\},\qquad
 B=\{v_0,v_1,b_0,b_1\}.
\]
The two nonidentity precomposition maps fix both $u_\eta$
and send each $a_j$ to $u_0$ or $u_1$, respectively.
The two nonidentity postcomposition maps fix both $v_\eta$
and send each $b_i$ to $v_0$ or $v_1$, respectively.
Allow any finite nonempty long-edge biset $T=C(0,2)$ and
any equivariant composition.  Every rational product of
two degree-two classes on $BC$ is zero.  Thus this entire
family fails to realize $\mathbb {CP}^2$ at every size.
\end{theorem}
\begin{proof}
The two adjacent action graphs are circles, each with two
vertices and two equally oriented parallel edges.  Write
$\Gamma_A,\Gamma_B$ for them.  As in the proof of
Theorem~\ref{C:affine-chart-family-zero}, the natural graph
and cubical models give a triangle of spaces with contractible
vertices, edges $\Gamma_A,\Gamma_B,K(T)$, and face
$\Gamma_A\times\Gamma_B$.  Its face maps are the two
projections and the cellular composition map
\[
 h:\Gamma_A\times\Gamma_B\longrightarrow K(T).
\]
Put $t_{ij}=b_i a_j$.  Equivariance gives, for every
$\varepsilon,\delta\in\{0,1\}$,
\begin{equation}\label{C:four-circle-boundary-equalities}
 P_\varepsilon t_{i0}=P_\varepsilon t_{i1},\qquad
 Q_\delta t_{0j}=Q_\delta t_{1j}.
\end{equation}
The remaining composition values are their boundary values.

Let $[t]$ denote the two-cell of $K(T)$ indexed by $t$ when
$t\notin U\cup V$, and put $[t]=0$ for a degenerate cell.
The fundamental cellular two-cycle of the product of the
two circles maps to
\begin{equation}\label{C:four-circle-top-chain}
 z=[t_{00}]-[t_{01}]-[t_{10}]+[t_{11}].
\end{equation}
This is a cycle because $h$ is a cellular map.  If $z\ne0$
as a cellular chain, it is nonzero in $H_2(K(T);\mathbb Q)$,
since the target has no three-cells.  The only total-degree-four
term in the triangle spectral sequence is
$H_2(\Gamma_A\times\Gamma_B;\mathbb Q)=\mathbb Q$.
Its first differential sends its generator to $z$; there
is no other nonzero term in total degree four.  Thus
$H_4(BC;\mathbb Q)=0$ in this case.

Suppose now that $z=0$.  First assume that one of its four
states is nondegenerate.  Cancellation in the free cellular
two-chain group forces a positive occurrence to equal a
negative occurrence as an actual state.  Such positions
are adjacent in the two-by-two array.  Interchanging edge
names, and exchanging the two coordinates in this argument
if necessary, reduces to
\[
 t_{00}=t_{01}\notin U\cup V.
\]
Equation~\eqref{C:four-circle-top-chain} gives
$[t_{10}]=[t_{11}]$.  If this common cell is nonzero, its
indices are equal.  Otherwise both states belong to $U\cup V$.
Their $P$-boundaries agree by
\eqref{C:four-circle-boundary-equalities}; their $Q$-boundaries
agree by the same equations and $t_{00}=t_{01}$.
Lemma~\ref{C:retraction-boundary-injective} again gives
$t_{10}=t_{11}$.  Thus the entire two-by-two array has equal
columns.  All its boundary values agree accordingly, so
$h$ factors through the quotient of $\Gamma_A$ identifying
its two equally oriented edges.  That quotient is an interval.
Contracting the interval to its initial vertex gives a based
homotopy of $h$ to a map factoring through
$\operatorname{pr}_{\Gamma_B}$.  The equal-row case gives
the other projection.

It remains to consider the case where all four $t_{ij}$
are degenerate.  The image of $h$ then lies in the one-skeleton
of $K(T)$, a graph.  The two circle generators of the source
map to closed edge paths in its two distinct edge colors:
one uses only the cells of $U\setminus V$, and the other
uses only the cells of $V\setminus U$.  These based loops
commute in the graph fundamental group, because they come
from a map of a torus.  They cannot both be nontrivial.
Indeed, reduce each edge path by cancellation of an edge
with its inverse.  If both reduced paths were nonempty,
the word formed by their commutator would have no cancellation
at any of its four junctions: adjacent factors have disjoint
edge colors.  It would therefore be a nontrivial reduced
loop in a graph, a contradiction.

At least one coordinate loop is consequently null-homotopic.
The connected graph component containing the image is a
$K(F,1)$.  The map $h$ is based homotopic to the map obtained
by restricting to the other coordinate circle and then
projecting onto it, since these maps induce the same
homomorphism on fundamental groups.  This conclusion also
covers the case where both loops are trivial.

We have proved that whenever $z=0$, the face map $h$ is
homotopic to a map through one of the adjacent-edge
projections.  Replacing this face map by that homotopic map
preserves the triangle homotopy colimit; the maps to its
vertex spaces are constant, so no additional compatibility
is required.  The suspension construction in the proof of
Theorem~\ref{C:affine-chart-family-zero} now gives a map
from this homotopy colimit to $\Sigma K(T)$ inducing an
isomorphism on rational degree-two cohomology.  Explicitly,
if $h=\alpha\operatorname{pr}_{\Gamma_B}$, use the identity
suspension on the long edge, $\Sigma\alpha$ on the $B$ edge,
and collapse the $A$ edge; the barycentric coordinate of
their common extreme vertex supplies the face formula.
The degree-two group is the kernel of
\[
 H^1(\Gamma_A)\oplus H^1(\Gamma_B)\oplus H^1(K(T))
 \longrightarrow H^1(\Gamma_A\times\Gamma_B),
\]
and projection of this kernel onto $H^1(K(T))$ is an
isomorphism.  The edge-to-face degree-zero map is onto,
so there is no additional degree-two term, even when
$K(T)$ is disconnected.  Since positive products on a
suspension vanish, every claimed product on $BC$ vanishes.
\end{proof}

\begin{remark}[Completed controls through twenty-one arrows]
\label{C:affine-four-circle-controls}
The exact generator \path{code/d_affine_four_circles.py}
enumerates all commuting retraction bisets and all equivariant
composition maps for $1\leq|T|\leq6$, using only actual
target-state and transient-edge automorphisms to remove
repetitions.  The respective arrow totals and numbers of
literal $H_1=0$, $\dim H_2=1$ mod-two candidates are
\[
\begin{array}{c|rrrrrr}
\text{arrows}&16&17&18&19&20&21\\\hline
\text{candidates}&0&6&44&340&2414&19652.
\end{array}
\]
The independent bar checker verifies a zero-square primitive
for all $22456$ candidates, including $4092904$ typed
associativity equations.  All six generator and checker
streams are complete; their prefixes are
\path{code/outputs/d_affine_four_circles_1} through
\path{code/outputs/d_affine_four_circles_6}.
These computations corroborate the theorem and are not used
in its all-size proof.
\end{remark}

\subsection{A one-edge circle against an arbitrary circle graph}

\begin{theorem}[The one-edge-circle chart obstruction]
\label{C:one-edge-circle-chart-zero}
Use the corner monoids and three-object support of
Theorem~\ref{C:affine-chart-family-zero}.  Suppose one adjacent
act is the two-state act whose two retractions both send its
transient state to its fixed state.  Suppose the natural graph
of the other adjacent act is connected and has first Betti
number one.  The long-edge biset and composition may be arbitrary.
Then every rational product of degree-two classes on $BC$ is
zero.  The same assertion holds with the two adjacent acts
interchanged.
\end{theorem}
\begin{proof}
Let $A$ be the one-edge circle and $B$ the other action graph.
Its unique essential circuit is a graph deformation retract.
Restricting to that circuit in the triangle-of-spaces model
is an objectwise homotopy equivalence: the face becomes
$A\times B_0\to A\times B$, where $B_0$ is that circuit,
and the long edge stays unchanged.  We may therefore assume
that $B$ is a circle with any finite number of oriented edges.
Write $\varepsilon_e\in\{1,-1\}$ for their coefficients in
an oriented fundamental cycle.

Use the natural cubical model $K=K(T)$ from
Lemma~\ref{C:two-retraction-cubical-model}, and write
\[
 h:A\times B\longrightarrow K
\]
for composition.  For an edge $e$ of $B$, let $t_e$ be the
long-edge state assigned to the product of $e$ with the unique
edge of $A$.  The two endpoints of that $A$-edge are the same
fixed state.  Consequently
\[
 P_0t_e=P_1t_e.
\]
As before, the cellular image of the torus fundamental class is
\[
 z=\sum_e\varepsilon_e[t_e]\in C_2(K;\mathbb Z),
\]
where $[t_e]$ is zero for a degenerate cell.  If $z\ne0$,
then $h_*:H_2(A\times B;\mathbb Q)\to H_2(K;\mathbb Q)$
is injective, since $K$ has no three-cells.  The triangle
spectral sequence gives $H_4(BC;\mathbb Q)=0$, proving the
claim in this case.

Suppose $z=0$.  If some $t_e$ belongs to $U$, the entire
$A$-coordinate collapses over that edge.  In particular the
$A$-loop at an endpoint of $e$ maps to a point.  Transporting
the basepoint along $B$ shows that its image is null-homotopic
at every choice of $B$-basepoint.  Thus
\begin{equation}\label{C:loop-chart-one-coordinate}
 h_*H_1(A;\mathbb Z)=0.
\end{equation}

Otherwise no $t_e$ on the essential circuit belongs to $U$.
At the fixed vertex of $A$, the image of the $B$-edge $e$
is the one-cell indexed by $P_0t_e$, or a vertex when that
state lies in $U\cap V$.  If $t_e\in V$, commutation makes
$P_0t_e\in U\cap V$, so its contribution is zero.  Every
remaining $t_e$ is a nondegenerate two-cell.  Apply to the
literally zero chain $z$ the linear map
\[
 C_2(K;\mathbb Z)\longrightarrow C_1(K;\mathbb Z),
 \qquad [t]\longmapsto
 \begin{cases}
 [P_0t],&P_0t\in U\setminus V,\\
 0,&P_0t\in U\cap V.
 \end{cases}
\]
Its value is exactly the cellular image of the $B$-cycle.
Hence in this second case
\begin{equation}\label{C:loop-chart-other-coordinate}
 h_*H_1(B;\mathbb Z)=0.
\end{equation}
We do not claim that the full map $h$ factors through a
projection: only one of these two homology statements is
needed.

Here is the cohomological consequence, which also explains
why homology in the preceding step suffices.  The triangle
spectral sequence identifies
\[
 H^2(BC;\mathbb Q)\cong
 \ker\bigl(H^1(A;\mathbb Q)\oplus H^1(B;\mathbb Q)
             \oplus H^1(K;\mathbb Q)
           \longrightarrow H^1(A\times B;\mathbb Q)\bigr).
\]
Projection of this kernel onto $H^1(K;\mathbb Q)$ is an
isomorphism, because the first two edge coordinates map
isomorphically onto the two summands of the target.  The
edge-to-face map in degree zero is onto, so there is no
additional degree-two term.  This remains valid when $K$
is disconnected.

For any integral class $u\in H^1(K;\mathbb Z)$, choose a
representing map $f_u:K\to S^1$.  Replacing the long edge
by $S^1$ and its face map by $f_u h$ gives a triangle of
spaces $X_u$ and an actual map
\[
 BC\longrightarrow X_u.
\]
The generator of $H^2(X_u;\mathbb Q)$ pulls back to the
class corresponding to $u$ in the preceding kernel
identification.  By \eqref{C:loop-chart-one-coordinate} or
\eqref{C:loop-chart-other-coordinate}, the map
$f_u h:A\times B\to S^1$ has degree zero on one coordinate.
Since $S^1$ is a $K(\mathbb Z,1)$, it is homotopic to a map
through the other projection.  The suspension construction
of Theorem~\ref{C:affine-chart-family-zero} then supplies a
map $X_u\to\Sigma S^1$ inducing an isomorphism on degree-two
cohomology.
The generator on $X_u$ therefore has zero square, and so
does its pullback on $BC$.

The finite CW complex $K$ has an integral first-cohomology
lattice spanning $H^1(K;\mathbb Q)$.  Multiplying any rational
class by an integer reduces it to the preceding integral
case.  Thus every rational degree-two square on $BC$ is
zero.  Polarization gives the vanishing of every product
of two such classes.
\end{proof}

\begin{corollary}[A one-edge circle against any connected action graph]
\label{C:one-edge-arbitrary-graph-cups}
In Theorem~\ref{C:one-edge-circle-chart-zero}, the other adjacent
action graph may be any connected finite graph, with no restriction
on its first Betti number.  All rational degree-two products still
vanish.
\end{corollary}
\begin{proof}
Write $X$ for the same triangle-of-spaces model, with adjacent
edges $A=S^1$ and $B$, long edge $K=K(T)$, and composition
$h:A\times B\to K$.  All edge spaces have dimension at most two,
the vertex spaces are points, and the face has dimension two.
The integral triangle spectral sequence therefore identifies
\[
 H_4(X;\mathbb Z)=
 \ker\bigl(H_2(A\times B;\mathbb Z)
                 \xrightarrow{h_*}H_2(K;\mathbb Z)\bigr).
\]
There are no other terms in total degree four and no higher
differential involving this kernel.  Since
$H_2(A\times B;\mathbb Z)=H_1(A;\mathbb Z)\otimes H_1(B;\mathbb Z)$,
each element of this kernel has the form $[A]\otimes\gamma$.

Choose a based closed edge word $\lambda:S^1\to B$ representing
$\gamma$.  This is always possible by choosing a maximal tree
in the connected graph and writing a word in its remaining edges.
Pulling back the triangle along $\lambda$ gives a triangle
$X_\gamma$ and an actual map $X_\gamma\to X$.  Its face map is
$h_\gamma=h(1\times\lambda):A\times S^1\to K$.
The cellular image of its fundamental two-cycle is literally
zero: it represents $h_*([A]\otimes\gamma)=0$, and $K$ has
no three-cells.  Consequently $H_4(X_\gamma;\mathbb Z)$ has
a generator mapping to the specified element of $H_4(X;\mathbb Z)$.

The zero-chain part of the preceding proof applies verbatim
to this closed edge word.  More explicitly, each traversal
copies the corresponding product square, with its orientation
reversed when the edge is traversed backwards.  If any used
top state lies in $U$, the $A$-loop has constant image at a
point of the other circle and hence zero image in first homology.
Otherwise every used degenerate state lies in $V$, and applying
the displayed $P_0$-face map to the literally zero signed
two-chain shows that the other circle has zero image in first
homology.  Repeated edges merely repeat their signed summands;
no embedded-circuit assumption is used in either argument.
For every map $K\to S^1$, the resulting face map therefore has
one coordinate degree zero.  The same suspension comparison
proves that every degree-two square on $X_\gamma$ is zero.
It is not necessary for $X_\gamma$ itself to be the classifying
space of a category in the stated family.

By naturality, the square of any rational degree-two class on
$X$ evaluates to zero on every integral element of $H_4(X)$.
These integral kernel elements span $H_4(X;\mathbb Q)$, and
rational cohomology is detected by evaluation on rational
homology.  Every rational square is therefore zero; polarization
proves the assertion for all mixed products.
\end{proof}

\begin{remark}[The complete longer-circle controls]
\label{C:longer-circle-chart-controls}
The exact program \path{code/d_circle_charts.py}, with the
typed filter \path{code/d_circle_charts_filter.cpp}, also
checks the following two bounded families.  Circle orientations
are retained explicitly: a two-edge circuit is either two
equally directed parallel edges or a directed two-cycle;
a three-edge circuit is either cyclically directed or
transitively directed.
\[
\begin{array}{c|c|c|r}
\text{adjacent state counts}&|T|&\text{orientation pairs}
 &\text{checked candidates}\\\hline
(4,6)&1,\ldots,4&4&5844\\
(2,6)&1,\ldots,6&2&44304.
\end{array}
\]
Every commuting long-biset retraction pair and every equivariant
composition is included.  The only orbit reductions are actual
state automorphisms of the two adjacent acts and the long biset.
The completed independent bar checks give explicit zero-square
primitives for all $50148$ candidates.  The output prefixes
\path{code/outputs/d_circle_charts_} record the chosen graph
names and long-biset size; \texttt{parallel2}, \texttt{cycle2},
\texttt{cycle3}, and \texttt{acyclic3} specify these four
orientations.  The additional suffix \texttt{minimal} denotes
a separately restricted fixed four-state target, not the full
biset search.  The theorem above makes the second row an
all-size structural exclusion.  The first row is stated only
in its completed bounded scope.
\end{remark}

% B proof; root and D full independent literal audits PASS.
\subsection{Bouquet graphs in the right-zero chart construction}

\begin{proposition}[A bouquet on one side of a chart]
\label{C:bouquet-chart-rank-one-square}
Use the three corner monoids of
Theorem~\ref{C:affine-chart-family-zero}: the initial corner is
the identity-adjoined left-zero two-band, the middle corner is
trivial, and the final corner is the identity-adjoined right-zero
two-band.  Let $A,B$ be the natural graphs of the two adjacent
acts, and let $K=K(T)$ be the cubical model of the long-edge
biset.  Suppose every component of $A$ is a bouquet of loops;
isolated vertices are allowed.  The graph $B$ and the composition
may be arbitrary.  If the whole category has rational second
cohomology of dimension at most one, its degree-two square is
zero.  The same conclusion holds with $A,B$ interchanged.
\end{proposition}
\begin{proof}
Write $X$ for the actual triangle-of-spaces model with vertex
spaces points, edge spaces $A,B,K$, face space $A\times B$,
and long face map $h:A\times B\to K$.  The other two face
maps are the projections.  This is the natural model supplied
by Lemmas~\ref{C:right-zero-action-graph} and
\ref{C:two-retraction-cubical-model}.

Replace each edge space by its component set and the face by
$\pi_0A\times\pi_0B$, retaining the component map induced by
$h$.  Call the resulting two-dimensional triangle space $X_0$.
The actual component maps give $q:X\to X_0$.  The strict-face
chain filtration gives the exact cohomology sequence
\[
 H^2(X_0;\Q)\xrightarrow{q^*}H^2(X;\Q)
 \longrightarrow
 \left\{(\alpha,\beta,\lambda):
 h^*\lambda=\operatorname{pr}_A^*\alpha+
                 \operatorname{pr}_B^*\beta\right\}
 \longrightarrow0,
\]
where $\alpha\in H^1(A;\Q)$, $\beta\in H^1(B;\Q)$,
and $\lambda\in H^1(K;\Q)$.  Indeed the bottom row of the
filtration is the cellular complex of $X_0$, and the next
row is the restriction map from the three edge first-cohomology
groups to $H^1(A\times B)$.  There is no vertex cohomology in positive
degree and hence no further differential in total degree two.
If $q^*$ has nonzero image, that image spans the assumed
one-dimensional group.  Its square vanishes because $X_0$
is two-dimensional.

We may therefore assume $q^*=0$.  Every rational degree-two
class is then uniquely described by a displayed triple.
After multiplying by a nonzero integer, its three entries
are integral.  Represent them by actual maps
\[
 \alpha:A\to S^1,\qquad \beta:B\to S^1,
 \qquad \lambda:K\to S^1.
\]
Their equality in first cohomology says, on each component,
that $\lambda h$ is homotopic to the addition map
$\alpha+\beta:A\times B\to S^1$.  Thus they give a map
of triangle homotopy colimits to the universal triangle
whose three edge spaces are $S^1$, whose face is
$S^1\times S^1$, and whose long map is addition.
Write $Y$ for that universal triangle.  The same filtration
gives $H^2(Y;\Q)=\Q$ and $H_4(Y;\Q)=\Q$.
The pullback of its degree-two generator is our class:
the triples agree, and the possible difference is in
the zero image of $q^*$.

Since $K$ is two-dimensional and $A,B$ are graphs, the
degree-four filtration identifies
\[
 H_4(X;\Q)=\ker\bigl(H_2(A\times B;\Q)
                            \xrightarrow{h_*}H_2(K;\Q)\bigr).
\]
Under the map to $Y$, a class $z$ in this kernel maps to
the evaluation of $\alpha\times\beta$ on $z$.  We show
that this evaluation is zero.  This suffices by cup
naturality, without computing the square in $Y$.

Choose a cellular integral one-cocycle representing
$\lambda$ on $K$.  Write $a_i$ for the literal oriented
loop edges of $A$, based at vertices $v_i$.  Every element
of $H_2(A\times B;\Q)$ has the form
\[
 z=\sum_i a_i\otimes\gamma_i,
 \qquad \gamma_i\in Z_1(B;\Q).
\]
Use $P_0,P_1$ for the horizontal retractions and $Q_0,Q_1$
for the vertical retractions on $T$.  For an edge $b$ of
$B$, let $t_{ib}$ be the actual composition state of
$a_i$ and $b$.  Since $a_i$ is a loop,
$P_0t_{ib}=P_1t_{ib}$.  Its horizontal lower face is
$Q_0t_{ib}$ and its vertical left face is $P_0t_{ib}$;
a face in the common core is a collapsed edge and is
assigned cocycle value zero.

If this top state is horizontally degenerate, the image
of $a_i$ is a point and $\alpha(a_i)=0$.  If it is
vertically degenerate, its vertical edge has value zero.
Otherwise the cohomology identity on $A\times B$ gives
\[
 \alpha(a_i)=\lambda(Q_0t_{ib}).
\]
For the vertical cycles it gives, even when $B$ is
disconnected,
\[
 \beta(\gamma_i)
   =\sum_b\gamma_i(b)\lambda(P_0t_{ib}).
\]
Consequently evaluation of $\alpha\times\beta$ on $z$
factors through its literal cellular image in $C_2(K;\Q)$:
assign to the two-cell indexed by $t$ the number
\[
 c(t)=\lambda(Q_0t)\lambda(P_0t).
\]
Degenerate top states contribute zero as just checked.
For $z\in\ker h_*$, its cellular image is the zero
two-chain, because $K$ has no cellular three-boundaries.
Thus $\langle\alpha\times\beta,z\rangle=0$.
The induced map on fourth homology is zero, so the
pullback square vanishes over $\Q$.  Rational scaling
proves the proposition for the original class.
\end{proof}

\begin{corollary}\label{C:three-state-adjacent-chart-square}
With the same three corner monoids, if either adjacent
act has at most three states, the category cannot have
the rational cohomology ring of $\CP^2$, regardless of
the other two hom-set sizes.
\end{corollary}
\begin{proof}
A right-zero graph on at most three states is either
a forest or a disjoint union of bouquets of loops.
In the forest case its first homology is zero, so the
displayed fourth-homology kernel is zero.  In the other
case apply the proposition.
\end{proof}

% Root and D independently completed literal audits; promoted after checkpoint 18.
\subsection{A singleton long hom-set with rationally acyclic corners}

\begin{lemma}\label{C:singleton-long-trivial-middle-suspension}
Let $C$ be a three-object linear category, with all three
forward hom-sets nonempty.  Suppose its middle
corner is trivial and its long hom-set has one element.
If the two endpoint monoids have rationally acyclic
classifying spaces, every positive-degree rational cohomology
product of $BC$ is zero.
\end{lemma}
\begin{proof}
Write $M,N$ for the two endpoint monoids.  Let $A$ be the
action-bar space of the right $M$-set $C(0,1)$ and $B$ the
action-bar space of the left $N$-set $C(1,2)$.  The strict-object
bar model has three vertex spaces $BM,*,BN$, three edge
spaces $A,B,BM\times BN$, and face space $A\times B$.
The face maps to the first two edges are the projections;
the third is induced by composition into the singleton long
hom-set.  The product form of the face uses the literally
trivial middle monoid.

Replace the three vertex spaces and the long-edge space by
points.  Keep $A,B$ and $A\times B$ unchanged.  These maps
give an actual map of the strict-object homotopy-colimit
diagrams.  They are rational homology equivalences on every
diagram entry: for the long edge this is the rational
K\"unneth theorem and the assumed acyclicity of $BM,BN$.
The finite homotopy-colimit filtration therefore makes the
realized map a rational homology equivalence and hence a
rational cohomology-ring isomorphism.

In the resulting model, collapse the long-edge interval
and its two endpoint vertices.  This is a contractible
subcomplex, so the collapse is a homotopy equivalence.
There remain two suspension vertices, one arc for every
point of $A$ and of $B$, and, for every $(a,b)\in A\times B$,
the suspended interval joining those two arcs.  This is
precisely the unreduced suspension of the join $A*B$.
Equivalently, the original triangular face modulo its
long side is a bigon, the suspension of the interval
from $a$ to $b$ in that join.  This description respects
all internal cells of $A$ and $B$ and their identifications.
The suspension has zero positive-degree cohomology
products.  Pullback through the preceding rational
equivalence proves the lemma.
\end{proof}

\begin{theorem}[A singleton long edge with acyclic corner bars]
\label{C:singleton-long-acyclic-corners}
Let $C$ be a three-object linear category with all three
forward hom-sets nonempty and with $|C(0,2)|=1$.
Suppose all three corner monoids have rationally acyclic
classifying spaces.  Then every positive-degree rational
cohomology product on $BC$ is zero.  The middle corner
need not be trivial or have projective trivial modules.
\end{theorem}
\begin{proof}
Write $M,N,P$ for the corners at $0,1,2$, respectively,
and write ${}_NS_M=C(0,1)$, ${}_PT_N=C(1,2)$,
and ${}_PU_M=C(0,2)$.  Use the standard strict-object
bar model, retaining every internal corner-bar direction.
Its vertex spaces are $BM,BN,BP$.  The edge spaces are
the mixed action bars
\[
 A=B(N,S,M),\quad B=B(P,T,N),\quad K=B(P,U,M).
\]
Its face space is the realized iterated two-sided bar
\[
 F=B(P,T,N,S,M).
\]
Before realization, its multi-simplices are strings
\[
 (p_1,\ldots,p_r;\ t;\ n_1,\ldots,n_q;\ s;\
                                      m_1,\ldots,m_l).
\]
The maps to $A$ and $B$ forget the opposite outer string
and the opposite cross arrow.  The map to $K$ retains
the $P,M$ strings and replaces the middle portion by
$t n_1\cdots n_q s\in U$.  Every middle-bar face and
degeneracy preserves this product, and outer faces
commute with it by the biset identities.  These are
therefore actual maps of the bar models.  The further
maps to vertices retain just the appropriate corner
string; the two possible composites agree literally.

This gives a strict diagram over the faces of the
triangle.  Its homotopy colimit is $BC$: decomposing a
nerve chain at its strict object changes gives precisely
these one-, two-, and three-corner bars, and subdividing
their products with the corresponding simplex gives
the usual nerve realization.  In particular $F$ is
not replaced by $A\times B$ when $N$ is nontrivial.

Since $U$ is a singleton, $K=BP\times BM$.  Replace
the three vertex spaces and $K$ by points while keeping
$A,B,F$ and their two maps $F\to A,B$ unchanged.
This is an actual map of diagrams and an entrywise
rational homology equivalence.  The finite homotopy-
colimit filtration makes its realization a rational
homology equivalence, preserving the rational ring.

Collapse the resulting long-edge interval.  Its endpoints
become one suspension vertex, and the middle vertex is
the other.  The two remaining edge spaces give the
suspended copies of $A$ and $B$.  Each triangular
prism $F\times\Delta^2$, with its long side collapsed,
is a suspended interval parametrized by $F$; its two
ends attach by the retained maps $F\to A,B$.
Thus the quotient is precisely
\[
 \Sigma\operatorname{hocolim}(A\longleftarrow F
                                      \longrightarrow B),
\]
with unreduced suspension and the ordinary double
mapping cylinder as the inner homotopy colimit.
All three spaces are nonempty.  The collapsed interval
is a contractible subcomplex, so this collapse preserves
homotopy type.  Suspension products vanish, proving
the theorem through the actual rational comparison.
\end{proof}

\begin{corollary}\label{C:small-corners-singleton-long}
Every three-object linear category with singleton long
hom-set and all three corner orders at most four has
zero positive-degree rational cohomology products.
\end{corollary}
\begin{proof}
Apply the theorem and the rationally acyclic classification
of the corner monoids of order at most four.
\end{proof}

\begin{corollary}\label{C:small-ends-singleton-long}
The lemma applies whenever both endpoint monoids have
order at most four.  There is no bound on the adjacent
hom-set sizes.
\end{corollary}
\begin{proof}
Every finite monoid of order at most four has rationally
acyclic classifying space, by the accepted small-monoid
classification.
\end{proof}

% Root and D completed independent full literal audits.
\subsection{A small right-zero endpoint behind a singleton long arrow}

\begin{theorem}\label{C:right-zero-small-last-stalk}
Let $C$ be a three-object linear category with all three forward
hom-sets nonempty. Suppose that its last corner is
$R_3=\{1,p,q\}$, with $p^2=p$, $q^2=q$, $pq=q$, and $qp=p$, that
$|C(0,2)|=1$, and that $|C(1,2)|\leq3$.
If $C$ has at most eighteen arrows, then $BC$ is not homotopy
equivalent to $\mathbb CP^2$. There is no projectivity
assumption on either of the other corners. The opposite
statement also holds.
\end{theorem}
\begin{proof}
Let $D$ be the full category on $0,1$ and put $S=C(1,2)$.
The incoming right $D$-act at the last object has a singleton
stalk at $0$ and the left $R_3$-set $S$ at $1$. Write $P,Q$
for the actions of $p,q$ on $S$ and $U=PS=QS$ for their
common image. The two retractions fix $U$ pointwise.

The natural action-category graph has vertex set $U$ and,
for each $s\in S\setminus U$, one edge from $Ps$ to $Qs$.
It can be obtained from the natural labeled graph model by
collapsing all edges labeled $p$, which form disjoint stars.
This is an equivariant weak equivalence: commuting right
maps preserve those stars and send a remaining edge at unit
speed to an edge or collapse it. At object $0$ the graph is
a point. Every map across the unique direction in $D$ takes
the graph at $1$ to that point. Thus these are actual natural
graph diagrams, not merely objectwise homology identifications.
Since $BR_3$ is contractible, the collage has the homotopy model
\[
 X=\operatorname{hocofib}\bigl(\Gamma_{hD}\longrightarrow BD\bigr).
\]

If $|S|\leq3$, the graph at $1$ is either a forest or a disjoint
union of bouquets, with isolated vertices allowed. Indeed
$|U|=1$ gives only loops, $|U|=2$ permits at most one edge,
and $|U|=3$ gives no edges.

In the forest case, replace both graph stalks by their component
sets. The comparison is an objectwise weak equivalence and is
natural for every arrow of $D$. Equivalently, the actual functor
from each $R_3$-action category to its component set has this
property. Collapse the last corner $R_3$ to the trivial monoid,
and replace $S$ by its set of action components. The singleton
long hom-set is unchanged; all composition maps descend.
This defines an actual finite category $C_0$ and a functor
$C\to C_0$ inducing a homotopy equivalence, by the action-bar
homotopy-pushout comparison. Its arrow count is
\[
 |C_0|=|C|-2-\bigl(|S|-|\pi_0\Gamma|\bigr)\leq16.
\]
Theorem~\ref{C:CP2-lower17} excludes this case.

In the remaining case every edge is a loop, so $Ps=Qs$ for
all $s$. The common retraction $P=Q:S\to U$ commutes with
the right action of the middle monoid. It also commutes with
the cross maps to the singleton stalk at $0$. The component
collapse therefore gives an actual functor $q:C\to C_0$,
where the last corner is trivial and its incoming hom-sets
are the singleton and $U$. There is an actual section
$j:C_0\to C$: send the last identity to the identity of
$R_3$, and include the core states $U$ among the arrows.
The singleton long hom-set guarantees that this preserves
all two-cross-arrow compositions. Hence $qj=1$ and
$|C_0|\leq|C|-2\leq16$.

Double each graph loop, fixing its vertex and every point
stalk. This is a strict natural graph map, since all graph
actions send an oriented edge to one edge or to a vertex.
It induces a self-map of $X$, acting as the identity on $BD$.
In the cellular action-bar chain model, graph vertices have
eigenvalue $1$ and graph edges have eigenvalue $2$; the graph
boundary is zero because every edge is a loop. Consequently
the entire cofiber chain complex splits into its component
complex and its loop complex. In every cohomological degree
the only eigenvalues of this self-map are $1,2$, and the
eigenvalue-$1$ subspace is exactly the image of
$q^*:H^*(BC_0;\mathbb Q)\to H^*(BC;\mathbb Q)$.

Suppose now that $BC\simeq\mathbb CP^2$. Its one-dimensional
second rational cohomology is an eigenspace. If its eigenvalue
is $2$, the square of its generator has eigenvalue $4$, which
does not occur in degree four, so that square is zero. This
contradicts the projective-plane ring.

If the eigenvalue is $1$, the generator pulls back from $BC_0$.
Its nonzero square forces both $H^2(BC_0;\mathbb Q)$ and
$H^4(BC_0;\mathbb Q)$ to be nonzero. Since $BC_0$ is an actual
homotopy retract of $BC$, its fundamental group is trivial and
its integral homology is a direct summand of that of
$\mathbb CP^2$. In degrees two and four the nonzero summands
are therefore the whole copy of $\mathbb Z$; there is no
homology in the other positive degrees. The section $Bj$ is
an integral homology equivalence between simply connected
CW spaces and hence a homotopy equivalence. Thus
$BC_0\simeq\mathbb CP^2$, again contradicting its at most
sixteen arrows by Theorem~\ref{C:CP2-lower17}. This proves the theorem.
\end{proof}

\begin{corollary}\label{C:small-last-stalk-seventeen-profile}
There is no seventeen-arrow projective-plane category with
linear profile $(m,n,3;x,y,1)$ satisfying $m\leq2$ and
$y\leq3$. The opposite profiles are also excluded.
\end{corollary}
\begin{proof}
For an ordinary monoid of order three, either its trivial right
rational module is projective or it is the right-zero monoid
$R_3$ above. In the latter case apply the theorem. In the former,
both extreme trivial modules are projective, since $m\leq2$;
Theorem~\ref{C:semisimple-endpoints-seventeen} excludes the
profile. This argument uses the actual projectivity alternatives
and does not infer finite CW type from a rational quotient.
\end{proof}

% Root and D completed independent full literal audits.
\subsection{A semisimple two-object block opposite a small corner}

\begin{lemma}[A natural graph for a right act on a semisimple pair]
\label{C:semisimple-pair-element-graph}
Let $D$ be a two-object linear category with corner monoids
$N,P$, nonempty hom-set ${}_PX_N$, and semisimple rational
corner algebras $\mathbb QN,\mathbb QP$. Let $F$ be a finite
right $D$-act, with stalks $U=F(0)$ and $T=F(1)$.
There is an actual natural rational homology equivalence from
the element category of $F$ to a finite graph. Its vertices are
\[
 (U/N)\ \sqcup\ (T/P),
\]
and its edges are the balanced orbit set
\[
 (T\otimes_P X)/N.
\]
The edge represented by $(t,x)$ goes from $[tx]$ to $[t]$.
Every commuting left action of a finite monoid induces strict
unit-speed graph maps and makes this comparison equivariant.
Moreover $H_j(BD;\mathbb Q)=0$ for all $j\geq2$.
\end{lemma}
\begin{proof}
Use the element-category convention in which an arrow
$(i,s)\to(j,t)$ is an arrow $a:i\to j$ with $ta=s$.
Thus a cross arrow is the pair $(t,x)$, with endpoints
$(0,tx)$ and $(1,t)$. Send the objects to their corner
orbits, all corner arrows to identities, and this cross
arrow to the displayed balanced orbit. The relations
\[
 (tp,x)\sim(t,px),\qquad (t,xn)\sim(t,x)
\]
are exactly the relations imposed by composing with the two
corner actions. The endpoint formulas respect these relations.
There are no composable pairs of cross arrows. This defines an
actual functor to the category whose only nonidentity arrows
are the indicated graph edges.

Decompose a nerve chain by whether it contains a cross arrow.
The two vertex bars compute
\[
 \mathbb QU\otimes^{\mathbb L}_{\mathbb QN}\mathbb Q,
 \qquad
 \mathbb QT\otimes^{\mathbb L}_{\mathbb QP}\mathbb Q.
\]
The cross-arrow bar computes the iterated tensor
\[
 \mathbb QT\otimes^{\mathbb L}_{\mathbb QP}
 \mathbb QX\otimes^{\mathbb L}_{\mathbb QN}\mathbb Q.
\]
Semisimplicity makes all three complexes concentrated in
degree zero. Their degree-zero groups are precisely the
linearizations of the two orbit sets and the displayed
balanced orbit set, with the actual endpoint maps. The
strict-object filtration, or its graph-of-bars realization,
therefore proves that the functor is a rational homology
equivalence. No component or orbit representatives are chosen.
A commuting left action preserves all the displayed relations,
so its comparison is strictly equivariant. Finally apply the
same argument to the constant singleton right act: its element
category is $D$, and its quotient is a finite graph.
\end{proof}

\begin{theorem}\label{C:semisimple-pair-six-endpoint-cups}
Let $C$ be a three-object linear category with all forward
hom-sets nonempty. Suppose its first two corner algebras are
semisimple over $\mathbb Q$, and its final corner monoid has
order at most six. If $BC$ has bounded rational homology,
then every rational degree-two product on $BC$ is zero.
The opposite assertion also holds. There is no bound on the
three cross hom-set sizes.
\end{theorem}
\begin{proof}
Write $D$ for the full subcategory on the first two objects,
$M$ for the last corner, and ${}_MF_D=C(-,2)$ for the incoming
biset. The element category of $F$ has its natural commuting
left $M$-action. Let $E$ denote its homotopy-orbit space.
The standard two-block collage model is the actual homotopy
pushout
\[
 BC\simeq\operatorname{hocolim}(BD\longleftarrow E
                                      \longrightarrow BM).
\]
Collapse the $BD$ vertex space, forming the homotopy cofiber
\[
 X=\operatorname{hocofib}(E\longrightarrow BM).
\]
This gives an actual map $BC\to X$. Since $H^2(BD;\mathbb Q)=0$,
the cofiber exact sequence gives a surjection
\[
 H^2(X;\mathbb Q)\longrightarrow H^2(BC;\mathbb Q).
\]
The same sequence and $H_j(BD;\mathbb Q)=0$ for $j\geq2$
show that $X$ has bounded rational homology whenever $BC$ does.
These are rational cofiber statements; no nullhomotopy of
$BD\to BC$ or integral finite-type conclusion is needed.

Apply Lemma~\ref{C:semisimple-pair-element-graph} to $F$.
Its actual, $M$-equivariant rational equivalence to a finite
graph $\Gamma$ remains a rational equivalence after taking
homotopy orbits: the free monoid bar is nonnegative and
preserves a quasi-isomorphism of coefficient complexes.
It consequently gives an actual rational equivalence of
cofibers
\[
 X\longrightarrow
 X_\Gamma=\operatorname{hocofib}(\Gamma_{hM}\longrightarrow BM).
\]
The finite graph has the unit-speed edge action required by
Theorem~\ref{C:two-column-graph-localization}; use its opposite
version for the current left action.

If $H^2(BM;\mathbb Q)=0$, all second-cohomology classes of
$X_\Gamma$ pull back from the suspension obtained by collapsing
its $BM$ subspace. Their products therefore vanish.
Otherwise, at order at most six, the minimal ideal of $M$ is
a $2\times2$ rectangular band. Its contracted algebra has
dimension one or two. In dimension two it has basis $1,a$
with $a^2$ equal to $0,a$, or $1$. It is respectively local
with rational residue, or semisimple. The row augmentation
for $M$ or $M^{\mathrm{op}}$ is one-dimensional. Thus either
it is projective or the local clause of the graph-localization
theorem applies. Bounded rational homology of $X_\Gamma$
has already been proved. In every case all its degree-two
products vanish.

Pull back through the actual rational equivalence and then
through the surjection in second cohomology. This proves that
all mixed degree-two products on $BC$ vanish, as claimed.
\end{proof}

\begin{corollary}\label{C:two-small-corners-six-endpoint}
There is no finite projective-plane category with a linear
profile $(m,n,p;x,y,u)$ for which $m,n\leq2$ and $p\leq6$.
The opposite profiles are excluded as well, with no bound
on $x,y,u$.
\end{corollary}
\begin{proof}
Both rational monoid algebras of orders at most two are
semisimple, and a projective-plane classifying space has
bounded rational homology. Apply the theorem.
\end{proof}

% Root and agents B and D have passed the full literal proof; independent finite gates cited below are complete.
\subsection{A singleton adjacent orbit at a projective endpoint}

\begin{proposition}\label{C:projective-initial-singleton-middle}
Let $C$ be a three-object linear category with corners $N,P,M$
and nonempty forward hom-sets $X=C(0,1)$, $T=C(1,2)$,
$U=C(0,2)$. Suppose the trivial left $\mathbb QN$-module
is projective and the right $N$-set $X$ has exactly one orbit.
If $|M|\leq11$ and $BC$ has bounded rational homology, all
rational degree-two products on $BC$ vanish. The opposite
assertion holds with right projectivity at the final endpoint.
\end{proposition}
\begin{proof}
The actual initial-corner orbit quotient of
Proposition~\ref{C:finite-target-extreme-corner} replaces $N$
by the trivial monoid and $X,U$ by $X/N,U/N$.
It is a rational homology and cohomology-ring equivalence,
because the asserted projectivity makes the initial
coinvariant functor exact. Denote the resulting category by
$C'$. Its first adjacent hom-set is a singleton, say $\{x\}$.

The full subcategory $C''$ on objects $0,2$ is an actual
homotopy retract of $C'$. Explicitly, define $G:C'\to C''$
by sending object $1$ to $0$, its whole endomorphism monoid
$P$ to the identity at $0$, the arrow $x$ to that identity,
and each $t\in T$ to $tx\in U/N$; leave all other arrows
unchanged. Since every element of $P$ fixes the singleton
$\{x\}$, these formulas preserve composition. The inclusion
$J:C''\to C'$ satisfies $GJ=1$. There is a natural
transformation $JG\Rightarrow1_{C'}$ whose component at
$1$ is $x$, and whose other components are identities.
Thus $BC'\simeq BC''$.

The category $C''$ is the ordinary left-act cone on the
finite $M$-set $U/N$. It has bounded rational homology by
the preceding equivalence and the initial rational quotient.
Lemma~\ref{C:eleven-bounded-all-products}, applied in the
opposite orientation if needed, kills all its degree-two
products. Pullback proves the proposition.
\end{proof}

\begin{corollary}\label{C:small-projective-initial-singleton}
There is no projective-plane category with linear profile
$(m,n,p;1,y,u)$ where $m\leq2$ and $p\leq11$.
The opposite profiles are also excluded. There is no
restriction on $n,y,u$ beyond finiteness and nonemptiness.
\end{corollary}
\begin{proof}
The initial trivial rational module is projective at order
at most two, and its singleton adjacent set has one orbit.
Apply the proposition.
\end{proof}

% Root, A, and D full literal audits passed. The frozen154-profile computation is preserved.
\subsection{Two equal incoming hom-set sizes at a monoid endpoint}

\begin{proposition}\label{C:two-small-incoming-homs}
Let $\mathcal C$ have objects $0,1,2$, with endomorphism monoids
$1,1,M$, respectively, and no backward arrows.  Suppose
\[
 X=\mathcal C(0,1)\ne\varnothing,
 \qquad |\mathcal C(1,2)|=|\mathcal C(0,2)|=2.
\]
Then $B\mathcal C$ is not homotopy equivalent to $\mathbb CP^2$.
No bound on the order of the finite monoid $M$ is required.
\end{proposition}

\begin{proof}
Put $T=\mathcal C(1,2)$ and $U=\mathcal C(0,2)$, regarded as left
$M$-sets.  Each $x\in X$ gives an equivariant map
$\beta_x:T\longrightarrow U$, $t\longmapsto tx$.
If one of these maps is constant, its value is fixed by every element
of $M$.  The lower full subcategory on $0,1$ is a graph, and in
particular has zero rational second cohomology.  The fixed incoming
arrow lemma therefore makes every rational degree-two product of
$B\mathcal C$ zero.

We may consequently assume that every $\beta_x$ is a bijection.
Choose $x_0\in X$ and identify $U$ with $T$ by $\beta_{x_0}$.
Every $\beta_x$ is now either the identity or the transposition
$\tau$ of this two-element set.  If both occur, $\tau$ commutes with
the $M$-action, so every actor is the identity or $\tau$.
Label the two states by $\mathbb Z/2$.  There is an actual functor
\[
 \mathcal C\longrightarrow B(\mathbb Z/2)
\]
which sends an endomorphism in $M$ to its permutation, an arrow in
$T$ or $U$ to its state label, and $x\in X$ to the parity of
$\beta_x$.  The equivariance and composition formulas verify every
functor equation.  If $\beta_x=\tau$, the loop represented by $x$
and the reverse of $x_0$ maps nontrivially.  Thus the fundamental
group of $B\mathcal C$ surjects onto $\mathbb Z/2$.

It remains that all $\beta_x$ are the same bijection.  Collapse $X$
to its selected arrow $x_0$, obtaining an actual functor
$q:\mathcal C\longrightarrow\mathcal D$ with a literal section.
Let $E$ be the action category of the left $M$-set $T$.
If $E$ is disconnected, both elements of $T$ are fixed by every
actor, and the fixed incoming arrow argument applies again.
Assume therefore that $E$ is connected.

Here the normalized rational chain kernel of $q$ can be read off
without a spectral sequence.  In degree one it has basis
$[x]-[x_0]$, for $x\ne x_0$.  In degree $n\ge2$ its basis consists
of the differences of chronological simplices
\[
 [x,t,a_1,\ldots,a_{n-2}]
 -[x_0,t,a_1,\ldots,a_{n-2}],
\]
where $t\in T$ and the $a_i$ are composable nonidentity arrows of
$M$.  The face forgetting $x$ cancels, as does the face composing
$t$ with $x$, since all $\beta_x$ coincide.  The remaining faces
are precisely the augmented action-category boundary, up to the
uniform shift sign.  Writing $\widetilde{\mathbb Q X}$ for the
augmentation kernel, this identifies the chain kernel with
\[
 \widetilde{\mathbb Q X}\otimes
 \operatorname{Cone}\bigl(C_*(BE;\mathbb Q)
                    \xrightarrow{\epsilon}\mathbb Q\bigr)[1].
\]
Since $E$ is nonempty and connected, the displayed chain kernel
has zero homology in degrees one and two.  Its long exact sequence consequently gives
an isomorphism
$q^*:H^2(B\mathcal D;\mathbb Q)\longrightarrow
H^2(B\mathcal C;\mathbb Q)$.

Finally, the single arrow $0\to1$ in $\mathcal D$ deletes object
$1$ by an actual deformation-retract functor: send $1$ to $0$,
send $t:1\to2$ to $tx_0$, and retain the monoid $M$.
The classifying space is therefore the ordinary left-act cone on
the two-element set $T$.  The universal two-cross-arrow theorem
makes all its rational degree-two products zero.  Naturality and
the preceding $H^2$ isomorphism finish the proof.
\end{proof}

\begin{proposition}\label{C:two-by-three-incoming-small-endpoint}
Let $\mathcal C$ have the same three-object support and trivial
first two corners, with $|M|\le8$, $X\ne\varnothing$,
$|\mathcal C(1,2)|=2$, and $|\mathcal C(0,2)|=3$.
Then every rational degree-two square on $B\mathcal C$ is zero.
\end{proposition}

\begin{proof}
If $H^2(BM;\mathbb Q)=0$, both blocks in the collage of the
graph $B\mathcal C|_{\{0,1\}}$ with $BM$ have zero second
cohomology.  The relative suspension quotient then kills all
degree-two products.  Assume $H^2(BM;\mathbb Q)\ne0$.

The minimal-ideal augmentation description forces at least two
minimal rows and at least two minimal columns.  Its maximal group
is trivial: a nontrivial group would make the ideal contain at
least eight elements, in addition to the identity outside that
ideal.  (If the identity were in the minimal ideal, $M$ would
itself be a group and have zero positive rational cohomology.)
Consequently every homomorphism from $M$ to a group is trivial.
Indeed its restriction to the rectangular minimal band is the
identity element of the group, and multiplication by any band
element then forces the same value on every element of $M$.

Again write $\beta_x:T\to U$ for the equivariant composition
maps.  A constant $\beta_x$ gives an $M$-fixed incoming arrow
and hence zero degree-two products.  Otherwise every $\beta_x$
is injective, with two-element invariant image in the three-element
set $U$.  Two distinct images have a singleton invariant
intersection, giving the same fixed-arrow conclusion.
We may thus suppose all images are one common set $W$.
If two of the maps differ, their ratio is the transposition of
$W$, so the $M$-action on $W$ commutes with that transposition.
Every such actor is a permutation; the preceding triviality of
group quotients makes the action trivial.  This again supplies
a fixed incoming arrow.

The only remaining case has all $\beta_x$ equal.  Collapse $X$
to $x_0$.  The explicit normalized chain-kernel calculation in
Proposition~\ref{C:two-small-incoming-homs} does not require
surjectivity of $\beta_{x_0}$: the two faces still cancel because
all composition maps coincide.  If the action category of $T$
is disconnected, its two states and therefore their images in
$U$ are fixed.  Otherwise the collapse induces an isomorphism
on rational second cohomology.  Its target deletes the middle
object using $t\mapsto tx_0$, leaving the ordinary $U$-act cone.
Every ordinary act cone on a monoid of order at most nine has
all rational degree-two products zero, by
Lemma~\ref{C:nine-unbounded-act-cups}.  Pullback proves the
asserted square vanishing; polarization also gives all mixed
degree-two products.
\end{proof}

\begin{remark}\label{C:two-small-incoming-homs-scope}
The propositions and their opposites remove, in particular, the
seventeen-arrow linear profiles $(1,1,8;3,2,2)$,
$(1,1,8;2,2,3)$, and $(1,1,7;4,2,2)$.
The first two are the only remaining order-eight corner profiles.
This observation is separate from the frozen
154-profile computation: no running source stream is altered by
this note.
\end{remark}

\subsection{Four cross states opposite a right-zero corner}

\begin{lemma}[The cofiber of a connected graph correspondence]
\label{C:connected-graph-cone-degree}
Let a finite monoid $M$ act on the right on a connected finite
graph $\Gamma$ by cellular maps, and suppose $H_1(\Gamma;\mathbb Q)$
is one-dimensional.  Write $\rho_m$ for the induced degree of
$m$ on that vector space.  Put
\[
 X=\operatorname{hocofib}(\Gamma_{hM}\longrightarrow BM).
\]
Then
\[
 \widetilde H_n(X;\mathbb Q)
   \cong\operatorname{Tor}^{\mathbb QM}_{n-2}
       (H_1(\Gamma;\mathbb Q),\mathbb Q)\qquad(n\geq2).
\]
In particular $H_2(X;\mathbb Q)\ne0$ forces $\rho_m=1$ for
every $m$.  If an $M$-equivariant self-map of $\Gamma$ has
degree $d\notin\{0,1\}$, then every positive-degree rational
cohomology product of $X$ is zero.
\end{lemma}
\begin{proof}
The augmented cellular chain complex
$C_1(\Gamma)\to C_0(\Gamma)\to\mathbb Q$ is a complex of
right $\mathbb QM$-modules.  Its only nonzero homology is
$H_1(\Gamma)$ in the first position.  Inserting this kernel is
an equivariant quasi-isomorphism.  Tensor with a free left
bar resolution and take the mapping cone of augmentation.
This computes the indicated homotopy cofiber and shifts the
kernel by two, proving the displayed formula.  It also gives
$\widetilde H_1(X;\mathbb Q)=0$.
In degree two the result is the coinvariant space
$H_1(\Gamma)\otimes_{\mathbb QM}\mathbb Q$; a single degree
$\rho_m\ne1$ makes it zero.

An equivariant degree-$d$ self-map acts as multiplication by
$d$ on the kernel in this chain comparison.  It therefore acts
as multiplication by $d$ on every positive homology and
cohomology group of the cofiber.  For positive-degree classes
$u,v$, naturality and multiplicativity give
$d(u v)=d^2(u v)$.  Over $\mathbb Q$ this forces $uv=0$.
\end{proof}

\begin{theorem}[Four-state right-zero reduction]
\label{C:right-zero-four-state-reduction}
Let $R=\{1,p,q\}$ with $ab=b$ for $a,b\in\{p,q\}$, and
let $S$ be a nonempty $(R,M)$-biset with $|S|\leq4$.
If its two-object collage has type $\mathbb{CP}^2$, an ordinary
right $M$-act cone on at most four states has that type.
\end{theorem}
\begin{proof}
The cases with at most three states follow from
Lemma~\ref{C:three-state-right-zero-biset}, the common-fixed-state
argument, and the finite-line lemma.  Suppose $|S|=4$ and let
$P,Q$ be the two retractions onto their common image $U$.
Use the natural graph model of
Lemma~\ref{C:right-zero-action-graph}: its vertices are $U$ and
its edges are the states outside $U$, with edge $s$ directed
from $Ps$ to $Qs$.  The graph carries the commuting right
$M$-action.  Since $BR$ is contractible, the collage is the
homotopy cofiber of $\Gamma_{hM}\to BM$.

If $|U|=1$, its point is fixed by both monoids, and every
degree-two product on the collage vanishes.  If $|U|=4$, the
left action is trivial, and the collage is equivalent to the
ordinary four-state $M$-act cone.
If $|U|=3$ and its single graph edge is not a loop, each graph
component is contractible.  The equivariant map to its two-point
component set is a homotopy equivalence, also after taking the
action bar.  This gives an ordinary two-state act cone.

Whenever $P=Q$, including the remaining $|U|=3$ case, the map
$P:S\to U$ is a biset retraction.  Let $Y$ be the collage on
$U$.  It is equivalent to the ordinary $M$-act cone on $U$,
and is a homotopy retract of the original collage $X$.
The relative quotient is $\Sigma(E_S/E_U)$, as in
Lemma~\ref{C:three-state-right-zero-biset}.
If $X\simeq\mathbb{CP}^2$ and $H_2(Y;\mathbb Z)=0$, every
rational degree-two class of $X$ is relative, so its square
would vanish.  Otherwise split integral homology makes
$H_2(Y)\to H_2(X)$ an isomorphism.  The nonzero primitive square
then makes $H_4(Y)\to H_4(X)$ an isomorphism as well.
All other reduced homology groups vanish, and $\pi_1(Y)=1$
because it is a retract of $\pi_1(X)=1$.  Homological Whitehead
makes $Y\to X$ a homotopy equivalence.  Thus this case gives
the required ordinary act cone.

It remains that $|U|=2$ and the two-edge graph is connected.
It has one-dimensional first homology.  A hypothetical
projective plane forces every actor to have degree $+1$ by
Lemma~\ref{C:connected-graph-cone-degree}.
There are three possible graph types.

If there is one connecting edge and one loop, a degree-one
actor must carry the unique loop to itself.  It therefore fixes
its vertex, which is also fixed by $R$.  This common fixed state
makes all degree-two classes relative to the two vertex spaces,
and their products vanish.  The connecting edge need not be
invariant: it may map to the loop.  Only the fixed loop vertex
is used.

If the two edges connect the two vertices in the same direction,
a degree-one actor must fix both vertices and both edges.
Swapping the edges has degree $-1$, while folding or collapsing
them has degree zero.  Hence there is again a common fixed
state and zero degree-two products.

Finally suppose the two connecting edges have opposite directions.
A degree-one actor is either the identity or the half-rotation
interchanging the two vertices and edges.  Identify this graph
with the circle so that the half-rotation is $z\mapsto-z$.
The degree-three map $z\mapsto z^3$ commutes with every actor.
Lemma~\ref{C:connected-graph-cone-degree} makes every positive
rational product zero, again excluding a projective plane.
These alternatives exhaust the four-state graph and finish
the reduction.
\end{proof}

\begin{corollary}\label{C:right-zero-four-state-bounds}
If $|M|\leq12$, a biset collage with the right-zero corner $R$
and at most four cross arrows cannot have type $\mathbb{CP}^2$.
For an arbitrary three-element opposite corner, the same
conclusion holds when $|M|\leq10$ and there are four cross arrows.
In particular the sixteen-arrow profile $(9,3,4)$ is excluded
without a new enumeration.
\end{corollary}
\begin{proof}
For $R$, use the reduction and
Corollary~\ref{C:four-state-cones-seventeen},
Theorem~\ref{C:three-state-cone-minimum22}, and the one- and
two-state finite-line cases.
Every other three-element monoid has projective trivial right
rational module, as checked by the splitting vectors in the
seven-monoid catalog of the preceding biset calculation.  The actual
orbit-act functor is then a rational homology and ring
equivalence.  For $|M|\leq9$ apply
Lemma~\ref{C:nine-bounded-rank-one}; for $|M|=10$ apply
Lemma~\ref{C:ten-rational-act-cones}.  Both use precisely the
bounded rational homology and one-dimensional degree-two
cohomology of a hypothetical projective plane.
\end{proof}

% Root and agent B passed the complete literal proof; the control certificate has an independent gate.
\subsection{Connected five-state graphs opposite a right-zero corner}

\begin{lemma}[The natural connected-graph Tor comparison]
\label{C:connected-graph-cone-natural-Tor}
Let a finite monoid $M$ act cellularly on the right on a connected
finite graph $\Gamma$. Put
\[
 X=\operatorname{hocofib}(\Gamma_{hM}\longrightarrow BM),
 \qquad W=H_1(\Gamma;\mathbb Z).
\]
Then $H_1(X;\mathbb Z)=0$ and, naturally in the graph and action,
\[
 \widetilde H_n(X;\mathbb Z)
   \cong\operatorname{Tor}_{n-2}^{\mathbb ZM}(W,\mathbb Z),
       \qquad n\geq2.
\]
The analogous formula holds over every field.
If the action factors through a surjective monoid map $M\to M'$, the
actual cofiber comparison is an isomorphism on $H_2$ and a surjection
on $H_3$, integrally and over every field.
If an $M$-equivariant self-map of $\Gamma$ acts on $W$ as
$d\,1_W$ for an integer $d\notin\{0,1\}$, every positive rational
cohomology product of $X$ vanishes.
\end{lemma}
\begin{proof}
The augmented cellular complex
$C_1(\Gamma)\to C_0(\Gamma)\to\mathbb Z$ is a complex of
right $\mathbb ZM$-modules. Its only nonzero homology is $W$
in degree one, and the literal inclusion of cycles is an
equivariant quasi-isomorphism from $W$ in that degree.
Tensor with a free left bar resolution of the trivial module.
The mapping cone of the augmentation computes the homotopy
cofiber and supplies the further degree shift. This gives the
displayed formula and its naturality, including the field version.

For a quotient of the acting monoid, the degree-zero two-sided bar
is the identical module $W$, and its degree-one map is surjective.
The degree-zero differentials agree, so every target one-cycle lifts
to a source one-cycle. This proves the $H_2$ isomorphism and $H_3$
surjection in the stated directions.
For a scalar self-map, the literal map on graph cycles is $d$;
the same natural chain comparison makes its induced map on every
positive-degree rational cohomology group multiplication by $d$.
Multiplicativity gives $d(xy)=d^2xy$, hence $xy=0$ over $\mathbb Q$.
\end{proof}

\begin{theorem}[The connected graph case at five cross states]
\label{C:right-zero-five-connected}
Let $R=\{1,p,q\}$ with $ab=b$ for $a,b\in\{p,q\}$, and let $S$
be an $(R,M)$-biset with at most five elements. Write $P,Q:S\to U$
for the two retractions onto their common image, and form the
natural action graph $\Gamma$ of
Lemma~\ref{C:right-zero-action-graph}: its vertices are $U$, and
each $s\in S\setminus U$ is an oriented edge from $Ps$ to $Qs$.
If this graph is connected, the two-object collage cannot have
classifying space $\mathbb CP^2$. There is no bound on $|M|$.
\end{theorem}
\begin{proof}
The graph comparison is natural for the right $M$-action, and
$BR$ is contractible. Thus the collage is the homotopy cofiber
in the preceding lemma. A vertex fixed by all of $M$ gives zero
degree-two products by the opposite form of
Lemma~\ref{C:fixed-incoming-cups}, since $H^2(BR;\mathbb Q)=0$.
An actor inducing zero on $W$ makes its coinvariants zero and
therefore makes $H_2(X;\mathbb Q)=0$.
We may exclude these possibilities throughout.

A connected graph on at most five total vertices and edges has
at most three vertices. With one vertex it has the fixed vertex
just discussed; with three vertices it is a tree, so its map to
a point is an equivariant underlying equivalence and the cofiber
is contractible. With two vertices and at most two edges, the
connected cases are precisely those treated in the proof of
Theorem~\ref{C:right-zero-four-state-reduction}: a tree, a loop
and a connecting edge, or two parallel connecting edges. Their
fixed-vertex, zero-degree and circle-power arguments exclude a
projective plane without any bound on the source monoid.

It remains that $|U|=2$ and there are three edges. Label the
vertices $0,1$ and let $k$ be the number of edges joining distinct
vertices. Since the graph is connected, $k\in\{1,2,3\}$.
Every actor maps an edge to an edge of the corresponding ordered
endpoint type or collapses it to a vertex. In particular it cannot
reverse an edge while fixing its endpoints. We enumerate the
following possibilities up to reversing all orientations and
interchanging the vertices.

\emph{Three connecting edges.}
If all have the same orientation, no actor can interchange the
vertices. A vertex-constant actor kills all graph homology;
otherwise both vertices are fixed. Thus the only remaining shape
has two edges $a,b:0\to1$ and one edge $c:1\to0$.
A vertex-constant actor again kills all graph homology, while in
the absence of a vertex interchange both vertices are fixed.
Hence every actor fixes or interchanges the vertices, with at least
one interchange. The quotient identifying $a$ with $b$ is strictly
$M$-equivariant and is the two-edge circle. Write
$z_a=a+c$, $z_b=b+c$ for the integral cycle basis.
An interchanging actor sends both basis cycles to the same one,
so their difference is zero in coinvariants. Every actor acts
as the identity on the quotient cycle. The induced map on cycle
coinvariants is therefore an isomorphism to $\mathbb Z$.
The natural Tor lemma makes the cofiber comparison an isomorphism
on second rational cohomology in the pullback direction.
On the quotient circle every actor is the identity or the
half-rotation. The degree-three map commutes with both, so its
cofiber has zero rational products by the scalar-map clause.
Pullback excludes the original nonzero square.

\emph{Two connecting edges and one loop.}
Put the loop at vertex zero. If the connecting edges have the
same orientation, a vertex interchange is impossible. A constant
map to vertex one kills all graph homology; in its absence vertex
zero is fixed. Thus the connecting edges must have opposite
orientations. If there is no interchanging actor, the same
alternative applies. An interchanging actor collapses the loop,
so its class is zero in cycle coinvariants. A vertex-constant
actor sends the connecting cycle to a multiple of this loop
class or to zero; it would then kill the remaining coinvariant
class as well. Consequently all actors fix or interchange the
vertices. Collapsing the loop is now a strictly equivariant map
to the two-edge circle. It induces an isomorphism on the
coinvariants of graph homology, since the loop class was already
killed and the remaining circle class has trivial action.
Use the same natural comparison and degree-three circle map.

\emph{One connecting edge and two loops.}
No actor can interchange the vertices, because no reverse
connecting edge exists. If both loops are based at one vertex,
an actor constant at the other vertex kills all graph homology;
in its absence the loop vertex is fixed. It remains that there
is one loop $\ell_i$ at each vertex $i$, and a connecting edge
$a:0\to1$.

To avoid a common fixed vertex, the action image must contain
actors constant at each of the two vertices. A nonzero invariant
functional on the cycle coinvariants forces both loop classes
to have the same nonzero value. Indeed a constant-at-$i$ actor
sends both loops either to $\ell_i$ or to its vertex; if any
such loop image is collapsed, the relations from constant actors
at both vertices force both functional values to vanish.
Likewise a vertex-fixing actor must fix both loops. Thus every
vertex-fixing actor acts as the identity on the entire graph,
and every other actor has exactly the form
\[
 c_{i,\delta}:\quad
 0,1\longmapsto i,\qquad
 \ell_0,\ell_1\longmapsto\ell_i,\qquad
 a\longmapsto
 \begin{cases}i,&\delta=0,\\\ell_i,&\delta=1.
 \end{cases}
\]
Here $i,\delta\in\{0,1\}$.

If all such actors have $\delta=0$, fix the connecting edge and
map each loop to its square. This is a strictly equivariant
graph self-map acting as twice the identity on graph homology.
If all have $\delta=1$, map each loop to its square and map the
connecting edge to the path $\ell_0a$, using two equal parameter
subintervals. A constant-at-$i$ actor sends this path to
$\ell_i^2$, exactly the image required by equivariance; the
vertex-fixing actors are identities. This again gives a scalar
degree-two self-map. Both uniform cases have zero rational products.

If both values of $\delta$ occur, closure and the presence of both
vertex values force all four maps $c_{i,\delta}$.
For the right-action convention their multiplication is
\[
 c_{i,\delta}c_{j,\epsilon}=c_{j,\delta}.
\]
The faithful image of $M$ is therefore exactly the identity
adjoined to this four-element rectangular band. The original
collage maps by an actual functor to the thirteen-arrow collage
with this five-element source image, the same corner $R$, and
the same five states. The graph-Tor comparison makes this
map an isomorphism on $H_2(-;\mathbb F_2)$.

The literal target table and a normalized two-cocycle and square
primitive are stored in
\path{code/outputs/a_rightzero_five_connected_control.json}.
Its normalized chain dimensions are $11,50,212,872$ in degrees
one through four. Exact elimination gives $H_1(-;\mathbb F_2)=0$
and $\dim H^2(-;\mathbb F_2)=1$; the supplied generator is not
a coboundary, and all $872$ equations verify its square primitive.
The standalone checker \path{code/b_verify_rightzero_five_control.py}
reconstructs the two retractions, all right maps, and the entire
category table without producer imports. It checks every typed
associativity equation, all consecutive-boundary equations, the
nonboundary generator and all literal square identities. Its complete
record is
\path{code/outputs/b_rightzero_five_connected_control_verified.json}.
Thus its degree-two square is zero. Naturality and the actual
second-homology isomorphism make the square zero on the original
collage, contradicting the mod-two ring of $\mathbb CP^2$.
These graph types exhaust the connected case.
\end{proof}

\begin{remark}[Remaining five-state graphs]
\label{C:right-zero-five-connected-scope}
This theorem concerns the connected graph of the left $R$-action,
not merely a connected action category after taking the right
$M$-bar. A disconnected graph can have its components joined by
the right action. Those cases are handled separately by
Theorem~\ref{C:right-zero-five-state-reduction} below.
The thirteen-arrow certificate treats one faithful action image;
the natural graph-Tor comparison is what extends that calculation
to arbitrary finite source monoids with that image.
\end{remark}

% Root and B full literal audits passed; D independently verified both H3 controls.
\subsection{Connected six-state graphs opposite the right-zero two-band}

\begin{lemma}[Polynomial graph actions and a degree-two square]
\label{C:connected-graph-polynomial-square}
Let a finite monoid act on a connected graph, and let $X$ be
the cofiber of its action-bar projection. Suppose an actual
equivariant self-map induces an operator $F$ on
$W=H_1(\Gamma;\mathbb Q)$ with $(F-1)(F-2)=0$.
If its pullback multiplies a degree-two class $v$ by two,
then $v^2=0$.
\end{lemma}
\begin{proof}
The natural connected graph comparison identifies every
positive homology group of $X$ with a shifted derived tensor
of $W$. Functoriality and additivity make
$(f^*-1)(f^*-2)=0$ on every positive rational cohomology
group. The square is multiplied by four, so that polynomial
applied to it gives $(4-1)(4-2)v^2=0$.
\end{proof}

\begin{theorem}[The connected six-state right-zero obstruction]
\label{C:right-zero-six-connected}
Let $R=\{1,p,q\}$ be the identity-adjoined right-zero
two-band, and let $|S|\leq6$ in an $(R,M)$-biset.
If its natural left-action graph is connected, the collage
cannot have type $\mathbb CP^2$. There is no bound on $|M|$.
\end{theorem}
\begin{proof}
A singleton core gives a common fixed state. Connected
graphs with at most five states were treated in
Theorem~\ref{C:right-zero-five-connected}.
The remaining graph has either three vertices and three
edges, which is unicyclic and excluded by
Lemma~\ref{C:unicyclic-graph-correspondence}, or two vertices
and four edges. A three-vertex tree is contractible and
cannot occur either. Label the two core vertices $0,1$.
We assume the second rational cohomology is one-dimensional,
and write $\lambda$ for its nonzero invariant functional
on $W=H_1(\Gamma;\mathbb Q)$.
All graph quotients used below are actual equivariant maps.
When their first-homology coinvariants have the same
nonzero one-dimensional functional, the natural graph-Tor
comparison makes their cofiber pullback an isomorphism on
second rational cohomology.

\emph{No loops.}
A constant vertex map kills all first homology, so none
can occur. Without a vertex swap a core vertex is fixed.
Thus both cross-edge orientations occur and a swap occurs.
For three edges in one direction and one in the other,
any swap kills all differences among the three parallel
edges. Identifying them gives the two-edge circle and an
isomorphism on first-homology coinvariants. Its degree-three
map kills the square.

For two edges in each direction let $P$ be the two-dimensional
space of differences within the two parallel pairs. If an
actor folds one pair, then either a swap transfers that
relation to the other pair, or the square of a nonbijective
swap kills both differences. Thus $P$ has zero coinvariants.
If every actor is a permutation, but a nonidentity
vertex-fixing permutation occurs, it reverses one or both
parallel differences; a vertex swap again forces both to
vanish rationally. These cases reduce to the same two-edge
circle quotient. Otherwise the entire graph action image
is just an involution interchanging the vertices. It acts
freely on vertices and edges, and its quotient has one
vertex and two edges. Its first-homology coinvariants have
dimension two, contrary to the assumed second homology.

\emph{Loops at only one vertex.}
There must again be a vertex swap, or a core vertex is fixed
unless a constant actor at the other vertex occurs; that
actor would kill all first homology. The swap kills the
whole loop subspace, because the other vertex has no loops.
Any constant actor then has image in a subspace already
zero in coinvariants, so makes all coinvariants zero.
Collapse the loops. This is an equivariant graph quotient,
with at most three cross edges, and is an isomorphism on
first-homology coinvariants. The connected five-state
rational square obstruction, proved in the connected part
of Lemma~\ref{C:nine-five-state-graph-square}, applies.
That connected argument places no bound on $M$.

There are only three remaining loop distributions.

\emph{One loop at each vertex and two parallel cross edges.}
Write the loops as $x_0,x_1$ and the parallel edges as
$a,b:0\to1$. There is no vertex swap, so both constant
vertex maps must occur. Invariance forces
$\lambda(x_0)=\lambda(x_1)=c\ne0$ and every constant
actor sends both loops to the loop at its target.
Every core-fixing actor fixes the loops.
On $y=a-b$, invariance gives
$\lambda(y)/c\in\{-1,0,1\}$.

If that value is zero, identifying $a=b$ is equivariant
and preserves the nonzero coinvariant functional, hence
reduces to a connected five-state graph. If it is one,
every core-fixing actor is the identity and every constant
actor sends $a$ to the target loop and $b$ to the target
vertex. Define $f(b)=b$ and $f(a)=ab^{-1}a$.
Map each loop twice around itself with a constant pause
between the two traversals. Parameterize these as three
equal intervals, so the middle collapsed $b^{-1}$ segment
matches the pause exactly under every constant actor.
This is a strictly equivariant map acting as twice the
identity on $W$. The scalar graph lemma kills the square.
The value minus one is the same argument with $a,b$
interchanged.

\emph{One loop at each vertex and opposite cross edges.}
Write $a:0\to1$, $b:1\to0$, and $z=a+b$.
If there is no constant vertex actor, a nonzero class on
the loop subspace either vanishes under a collapsed loop
and its swapped mate, or all actors act through the single
simultaneous swap of loops and cross edges. The latter
case has two-dimensional first-homology coinvariants.
Without a swap there is a fixed core vertex. Thus a rank-one
square in the no-constant case pulls back from the two-edge
circle quotient and is zero.

With a constant actor, avoiding a fixed core vertex forces
both constant vertex types, possibly by composing with a
swap. Invariance now gives
$\lambda(x_0)=\lambda(x_1)=c\ne0$, and the loops are
never collapsed. Core-fixing actors are identities; a swap,
if present, is the unique simultaneous swap.
For a constant actor write $\delta_a,\delta_b\in\{0,1\}$
for whether the corresponding cross edge enters the target
loop rather than the target vertex. The sum
\[
 d=\delta_a+\delta_b=\lambda(z)/c
\]
is the same for every constant actor.

For $d=0$, square both loops and fix $a,b$. The resulting
map is strictly equivariant, its operator on $W$ has
eigenvalues one and two, and $\lambda$ is multiplied by
two. For $d=2$, square the loops and put
$f(a)=x_0a$, $f(b)=x_1b$, with equal halves on each
concatenation. Constant actors send both halves to the
same loop, and the swap preserves the construction.
Here $F(x_i)=2x_i$ and $F(z)=z+x_0+x_1$, so again
$(F-1)(F-2)=0$ and $\lambda F=2\lambda$.
Lemma~\ref{C:connected-graph-polynomial-square} handles
both cases.

If $d=1$ and only the pattern $(1,0)$ occurs, there is
no swap. Square the loops, put $f(a)=x_0a$, and fix $b$.
Then $F(z)=z+x_0$, giving the same polynomial and eigenvalue.
The pattern $(0,1)$ is symmetric.
Finally, if both patterns occur, closure forces all four
constant actors. The actual right-action image is the
five-element monoid consisting of these actors and the
identity, or the six-element monoid also containing the
simultaneous swap.

The corresponding fourteen- and fifteen-arrow literal
collages are stored in
\path{code/outputs/a_rightzero_six_opposite_edges.json}.
Each has a nonzero third homology class over $\mathbb F_2$,
certified by an explicit normalized three-cycle and a
three-cocycle pairing to one. Their normalized four-simplex
counts are respectively $992$ and $1859$.
The standalone checker
\path{code/d_verify_rightzero_six_H3.py} reconstructs both
literal categories from their corner tables and actions,
checks all typed associativity triples, and independently
enumerates every normalized chain. It verifies both
three-cycles, all $992+1859$ three-cocycle equations, and
the two pairings equal to one, without producer imports
or a matrix-rank calculation. Its complete record is
\path{code/outputs/d_rightzero_six_H3_independently_verified.json}.
The actual quotient from $M$ to either action image is
onto on third homology of the connected graph cofiber:
the natural comparison is $\operatorname{Tor}_1(W,\mathbb F_2)$,
whose two-sided bar has identical degree zero and a
surjection in degree one. Every target one-cycle therefore
lifts to a source one-cycle. This contradicts the third
homology of $\mathbb CP^2$ and excludes both images,
including arbitrary nonfaithful source inflations.

\emph{Three loops and one cross edge.}
Suppose there are loops $a,b$ at vertex zero, one loop
$c$ at vertex one, and the unique cross edge.
There is no swap, and both constant vertex types are
necessary. The constant-one actor forces
$\lambda(a),\lambda(b)\in\{0,\lambda(c)\}$,
with $\lambda(c)\ne0$. The constant-zero actor forces
at least one of the two values to equal $\lambda(c)$.
If one is zero, its loop and its vertex form a right-invariant
collapse: no actor can send the zero-valued loop to a
nonzero-valued loop. Collapsing it gives the connected
five-state graph and an isomorphism on the rank-one
coinvariant space. If both values equal $\lambda(c)$,
identify $a=b$ instead. This is again equivariant; a
constant-one actor kills their difference, so the map
preserves the same coinvariant space. The rational
connected five-state obstruction finishes both cases.

These alternatives exhaust every connected graph with
two vertices and four edges, proving the theorem.
\end{proof}

\begin{remark}[Scope of the image controls]
\label{C:right-zero-six-connected-control-scope}
The two small third-homology certificates have passed their
standalone independent verification. They assert nonzero
third homology of precisely the two displayed literal
categories; the comparison in the proof supplies the
separate nonfaithful-source consequence.
The separate ninety-three-case faithful-image search on a
five-state component plus an isolated core vertex is not used
here and supplies no nonfaithful-source completeness claim.
\end{remark}

% Root and B literal audits passed; independent equal-loop and isolated-image/lift gates passed.
\subsection{Disconnected six-state graphs with an eight-element source}

\begin{lemma}[Low homology under an action-image quotient]
\label{C:graph-action-image-low-homology}
Let a graph $\Gamma$ carry a cellular right action of a monoid $M$
which factors through a surjection $M\to H$.  The actual map of
action-bar cofibers
\[
 X_M=\operatorname{hocofib}(\Gamma_{hM}\to BM)
       \longrightarrow
 X_H=\operatorname{hocofib}(\Gamma_{hH}\to BH)
\]
is an isomorphism on $H_1$ and a surjection on $H_2$, integrally
and over every field.  The graph need not be connected.
\end{lemma}
\begin{proof}
Use the augmented graph complex
$A_\Gamma=[C_1(\Gamma)\to\widetilde C_0(\Gamma)]$,
in degrees one and zero.  A natural reduced chain model for the
cofiber is the suspension of the total two-sided monoid bar on
$A_\Gamma$.  Before that suspension, the quotient map is literally
the identity in total degree zero and is onto in total degree one.
The images of the two differentials into degree zero coincide.
Consequently it is an isomorphism on degree-zero homology.
Every degree-one target cycle lifts to a source chain; its boundary
is zero since degree zero is unchanged.  This proves surjectivity
on degree-one homology, hence the asserted shifted statements.
\end{proof}

\begin{lemma}[A graph quotient with no cycle coinvariants in its kernel]
\label{C:graph-cycle-kernel-quotient}
Suppose an equivariant graph quotient $\Gamma\to\Gamma'$ is
bijective on components and onto on first homology.  Write $L$
for the kernel of its first-homology map over a field $k$.
If $L\otimes_{kM}k=0$, its action-bar cofiber map induces an
isomorphism on $H_2(-;k)$.
\end{lemma}
\begin{proof}
The actual cellular chain maps used here are degreewise
surjective.  Their first-homology surjectivity and component
isomorphism imply that the chain kernel has homology only $L$
in degree one.  Graphs have no cellular two-boundaries, so the
literal inclusion of kernel cycles is a natural equivariant
quasi-isomorphism from $L$ in that degree.  For an arbitrary
cellular comparison with the stated homology properties, use its
derived fiber and the same mapping-cone homology sequence instead.
After taking the free monoid bar and suspending, the comparison
fiber has homology $L\otimes_{kM}k$ in degree two and zero in
degree one.  The homology exact sequence gives the assertion.
\end{proof}

\begin{proposition}[Two bounded disconnected image controls]
\label{C:right-zero-six-disconnected-controls}
Consider the following two six-state right-zero action graphs,
up to relabeling states and interchanging $p,q$ (which reverses
all edge orientations).
\begin{enumerate}
\item There is a loop at $0$, a directed edge $0\to1$, and a
separate loop at $2$.  Every right actor preserves the functional
which is one on both loops.
\item There is a loop at each of $0,1$, an edge $0\to1$, and an
isolated vertex $5$.  The restriction of the right-action image
to the five-state connected component is the complete five-element
mixed image from Theorem~\ref{C:right-zero-five-connected}.
\end{enumerate}
If $|M|\leq8$, neither collage has type $\mathbb CP^2$.
\end{proposition}
\begin{proof}
For the first graph, the two left retractions on the state list
$(0,1,2,\ell_0,\ell_2,a)$ are
\[
 P=(0,1,2,0,2,0),\qquad Q=(0,1,2,0,2,1).
\]
There are exactly $55$ commuting state maps.  Exactly ten preserve
the displayed loop functional, and these form a monoid.  Among its
$1024$ subsets, exactly $56$ contain the identity, are closed, and
have cardinality at most eight.  Their literal collages split as
follows: seven have nonzero $H_1(\mathbb F_2)$; three have
$\dim H_2(\mathbb F_2)=2$; and forty-six have one-dimensional
$H^2(\mathbb F_2)$ with an explicit zero-square primitive.
No image has zero second homology.

The producer record is
\path{code/outputs/a_rightzero_six_loop_tail_images.json}; only
its first, equal-loop hull is used here.  The standalone verifier
\path{code/d_verify_rightzero_six_equal_loop.py} independently
enumerates all $6^6$ state maps and all $1024$ subsets, rebuilds
every literal category, and verifies its binary ranks and actual
cup equations.  Its complete record
\path{code/outputs/d_rightzero_six_equal_loop_independently_verified.json}
checks $23169$ typed associativity equations, $11659$ cocycle
equations, and $58771$ primitive equations.
Lemma~\ref{C:graph-action-image-low-homology} transfers these
obstructions to every source with that image.  In the square case,
the target's nonzero one-dimensional second homology and the
surjection from the hypothetical projective plane make pullback
an isomorphism on second cohomology; this nonzero-target condition
is essential.

For the second graph put
\[
 P=(0,1,0,1,0,5),\qquad Q=(0,1,0,1,1,5).
\]
A right actor restricts to one of the five prescribed maps on
$0,\ldots,4$ and sends $5$ to one of the three core vertices
$0,1,5$.  These are precisely fifteen maps, closed under
composition.  Exhausting their subsets of cardinality at most
eight whose restriction contains all five component maps gives
the following $93$ literal image monoids:
\[
\begin{array}{c|rrrr}
\text{image order}&5&6&7&8\\\hline
H_1(\mathbb F_2)\ne0&1&0&0&0\\
\dim H_2(\mathbb F_2)=2&1&0&0&0\\
\dim H_2(\mathbb F_2)=1,\ \text{square zero}&0&10&12&6\\
H_2(\mathbb F_2)=0&0&0&19&44
\end{array}
\]
The first three rows are excluded by the same actual comparison.
An order-eight image must be the source itself, so its zero
second homology also excludes it.  An order-seven image leaves
only the source itself or an eight-element source with one
two-element quotient fiber.

If $H^2(BM;\mathbb Q)=0$, the collage's degree-two classes are
relative to its two vertex spaces, and the suspension quotient
kills their squares.  Otherwise every order-eight source occurs
in the complete $516$-table list proved in
Theorem~\ref{C:eight-act-cones}.  No published-census completeness
is needed for that list.  For each such table, try every pair
of elements as the sole nontrivial congruence class.  There are
$989$ actual one-pair congruences.  Construct the seven-element
quotient and try every literal isomorphism to each of the nineteen
remaining image monoids.  This yields $660$ distinct right acts
on the fixed six states.  Every one of their actual seventeen-arrow
collages has $H_2(\mathbb F_2)=0$.

The complete producers and records are
\path{code/a_rightzero_six_isolated_images.py},
\path{code/a_rightzero_six_isolated_inflations.py}, and their
same-prefix files in \path{code/outputs/}.
The latter record stores the source index, unique quotient pair,
literal quotient map, full category table, and homology result
for every act.  It links the original $516$-table source and the
$93$-image record by their exact hashes.
The standalone verifier
\path{code/b_verify_rightzero_six_isolated.py} imports no producer
or matrix-library routines.  It reconstructs the fifteen maps
from all $6^6$ state maps and exhausts all $16384$ subsets containing
the identity.  It independently checks all $516$ source tables,
all $989$ pair congruences, and all $19\cdot720$ target relabelings;
their literal quotient lifts reproduce exactly the $660$ acts.
All $753$ image and lifted categories have their typed laws and
normalized bar equations checked, including $89048$ supplied
cup-primitive equations.  Its complete record is
\path{code/outputs/b_rightzero_six_isolated_independently_verified.json}.
Thus the nineteen nonfaithful-source cases are excluded by a
complete small lift calculation, not by transferring zero second
homology from an action image.
\end{proof}

\begin{theorem}[Six right-zero states with a source of order at most eight]
\label{C:right-zero-six-small-sources}
Let $R=\{1,p,q\}$ be the identity-adjoined right-zero two-band.
For $|M|\leq8$ and $|S|\leq6$, an $(R,M)$-biset collage cannot
have type $\mathbb CP^2$.  In particular the seventeen-arrow
profile $(|M|,|R|,|S|)=(8,3,6)$ is excluded.
\end{theorem}
\begin{proof}
The connected case is Theorem~\ref{C:right-zero-six-connected}.
Let $\Gamma$ be disconnected and put $T=\pi_0\Gamma$.
The ordinary $T$-act cone has zero rational degree-two products
by Lemma~\ref{C:nine-unbounded-act-cups}.  If a hypothetical
projective-plane generator pulled back from that cone, its square
would therefore be zero.  By the actual augmented-graph exact
sequence of Lemma~\ref{C:graph-component-eigenvalues}, its image
is consequently a nonzero invariant functional
\[
 \lambda:H_1(\Gamma;\mathbb Q)\longrightarrow\mathbb Q.
\]
This assertion does not identify the entire second cohomology
with cycle coinvariants when the graph is disconnected.

A forest is an actual ordinary-act reduction.  If all edges are
loops, the vertex retraction is an ordinary-act homotopy retract
with suspension relative quotient.  Its second cohomology either
spans that of the collage, or every degree-two class is relative;
both possibilities give zero square.  We may assume there is a
nonloop edge and a cycle.

With four vertices and at most two edges, a unique loop forces
its vertex to be fixed by every actor unless its cycle
coinvariants vanish.  The only other cyclic possibility is a
two-edge circle and isolated vertices.  Nonzero $\lambda$ forces
degree $+1$ on the circle and preservation of its component.
The degree-three circle map fixes its two vertices and commutes
with identity or half-rotation; extend it by the identity on all
isolated vertices.  It remains equivariant even when an isolated
vertex enters the circle.  The component eigenvalue lemma kills
the square.  More vertices leave no further nonforest case.

Thus there are three vertices, three edges, and two components.
If all cycles lie in a single component, either its vertex is
fixed, or it has two vertices and three edges, with the remaining
vertex isolated.  Call that connected component $C$.  Every actor
must preserve $C$: sending it to the isolated vertex annihilates
the nonzero functional $\lambda$.

Apply the literal connected five-state analysis to the action on
$C$.  Its reductions use a nonzero invariant cycle functional and
absence of a common fixed core vertex.  With three connecting
edges, a swap kills the parallel-edge difference.  With two
connecting edges and one loop, a swap kills the loop difference.
The corresponding equivariant quotients preserve the isolated
vertex and have zero coinvariants in their cycle kernel, so
Lemma~\ref{C:graph-cycle-kernel-quotient} gives a second-homology
isomorphism.  Their circle power maps fix every graph vertex and
therefore extend equivariantly over the isolated vertex.

For one connecting edge and two loops, both loops must be based
at different endpoints; otherwise their common loop vertex is
fixed.  The nonzero functional and absence of a fixed core vertex
force both constant-vertex types and equal nonzero loop values.
If all constants have the same collapse choice, the strict
degree-two maps in Theorem~\ref{C:right-zero-five-connected}
fix every vertex and extend over the isolated vertex.  The
component eigenvalue lemma again applies.  The only remaining
case is the complete mixed five-element component image, which
is precisely the second control of the preceding proposition.
No coinvariant-rank assumption about the disconnected collage
has been inserted in this reduction.

Suppose instead that both components contain a cycle.  One
component has two vertices and the other is a single vertex.
If the two-vertex component is a two-edge circle, the other is
a loop.  A nonzero value on that unique loop forces its vertex
to be fixed by every actor.  Otherwise $\lambda$ is nonzero on
the circle and zero on the loop.  No actor can send the circle
to the loop component, and its degree on the circle must be
$+1$.  Apply the degree-three map to both the circle and the
loop.  It fixes all vertices and commutes also with an actor
which collapses the loop to a circle vertex.  Thus it is a strict
scalar graph map and the component eigenvalue lemma applies.

The remaining two-vertex component is a loop with a connecting
tail, while the singleton component also carries one loop.
If the two loop values of $\lambda$ are unequal, a loop with
nonzero value must be sent to itself by every actor, forcing
its core vertex to be fixed.  If they are equal, they are both
nonzero, and every actor belongs to the equal-loop hull in the
first control above.  This excludes the last case.
Finally, a tree component and a singleton component carrying
both loops also force the latter vertex to be fixed: sending
that component into the tree annihilates $\lambda$.

This exhausts the disconnected graphs and proves the theorem.
\end{proof}

% Root, B, and D literal audits passed; no new numerical gate.
\subsection{All five-state actions opposite a right-zero corner}

\begin{lemma}[The two eigenvalues of a component-preserving graph map]
\label{C:graph-component-eigenvalues}
Let a finite monoid $M$ act cellularly on a finite graph $\Gamma$,
put $T=\pi_0\Gamma$, and form the actual comparison
\[
 q:X=\operatorname{hocofib}(\Gamma_{hM}\to BM)
       \longrightarrow
 X_0=\operatorname{hocofib}(T_{hM}\to BM).
\]
Write $W=H_1(\Gamma;\mathbb Q)$ as a right $\mathbb QM$-module.
If $W\otimes_{\mathbb QM}\mathbb Q=0$, the pullback
$H^2(X_0;\mathbb Q)\to H^2(X;\mathbb Q)$ is surjective.

More generally, suppose an actual $M$-equivariant graph self-map
induces the identity on $T$ and multiplication by
$d\in\mathbb Z\setminus\{-1,0,1\}$ on $W$.
On every rational cohomology group of $X$, its induced map $f^*$
satisfies
\[
 (f^*-1)(f^*-d)=0,
 \qquad
 \ker(f^*-1)=\operatorname{im}\bigl(q^*:H^*(X_0;\mathbb Q)
                                   \to H^*(X;\mathbb Q)\bigr).
\]
If $\dim H^2(X;\mathbb Q)=1$ and all degree-two products on $X_0$
vanish, then the degree-two square on $X$ vanishes as well.
\end{lemma}
\begin{proof}
Let $A_\Gamma$ be the augmented graph complex with
$C_1(\Gamma;\mathbb Q)$ in degree one and
$\widetilde C_0(\Gamma;\mathbb Q)$ in degree zero.
Projection to the component set gives a degreewise surjection
of right-module complexes
\[
 A_\Gamma\longrightarrow\widetilde{\mathbb QT}.
\]
Its kernel has $C_1(\Gamma;\mathbb Q)$ in degree one and
$\ker(C_0(\Gamma;\mathbb Q)\to\mathbb QT)$ in degree zero.
This kernel has only the homology $W$ in degree one; the literal
inclusion of graph cycles is a natural equivariant quasi-isomorphism.

Tensor this short exact sequence with the free left monoid bar
resolution and suspend once. These are natural reduced chain models
for the two cofibers and their actual map $q$. The kernel is
quasi-isomorphic to the twice-suspended two-sided bar of $W$.
Consequently the cohomology exact sequence has the segment
\[
 H^2(X_0;\mathbb Q)\longrightarrow H^2(X;\mathbb Q)
 \longrightarrow
 \operatorname{Hom}_{\mathbb Q}
       (W\otimes_{\mathbb QM}\mathbb Q,\mathbb Q),
\]
which proves the first assertion.

For the self-map, the action on the component complex is the
identity and its action on the kernel's bar homology is scalar $d$.
This is a natural statement on the actual chain models: graph
cycles have no two-dimensional boundaries, so their induced map
is literally multiplication by $d$. In the cohomology exact
sequence, $f^*-d$ therefore carries every class on $X$ into the
image of $q^*$, where $f^*$ is the identity. This proves the
polynomial identity. A class fixed by $f^*$ maps to a class both
fixed and multiplied by $d$ in the kernel term; since $d\ne1$,
that image is zero. Exactness gives the stated eigenvalue-one
pullback assertion.

On the one-dimensional space $H^2(X;\mathbb Q)$, $f^*$ is scalar
either $1$ or $d$. In the first case its generator pulls back from
$X_0$, so its square is zero by hypothesis. In the second case its
square is multiplied by $d^2$. The displayed polynomial in degree
four annihilates this square, and
$(d^2-1)(d^2-d)\ne0$. The square is again zero.
This argument makes no finite-type inference about $X_0$.
\end{proof}

\begin{theorem}[Five-state right-zero reduction]
\label{C:right-zero-five-state-reduction}
Let $R=\{1,p,q\}$, with $ab=b$ for $a,b\in\{p,q\}$, and let
$S$ be an $(R,M)$-biset with $|S|\leq5$.
If its two-object collage has classifying space $\mathbb CP^2$,
then an ordinary right $M$-act cone on at most $|S|$ states has
that homotopy type.
\end{theorem}
\begin{proof}
Use the two retractions $P,Q:S\to U$ and their natural right
$M$-action graph $\Gamma$. Its collage is
$X=\operatorname{hocofib}(\Gamma_{hM}\to BM)$.
The connected case is excluded by
Theorem~\ref{C:right-zero-five-connected}.

If $\Gamma$ is a forest, its projection to the component set
$T$ is an equivariant underlying homotopy equivalence. The bar
construction preserves it, so $X$ is actually equivalent to the
ordinary $T$-act cone. Since $|T|\leq|S|$, this is the required
reduction.

If $P=Q$, their common retraction is an actual biset retraction
$S\to U$. The collage $Y$ on $U$ is a homotopy retract of $X$
and is equivalent to the ordinary $U$-act cone, because the left
$R$-action there is trivial and $BR$ is contractible.
The relative quotient of the two double-mapping-cylinder models
is $\Sigma(E_S/E_U)$, precisely as in the $P=Q$ part of
Theorem~\ref{C:right-zero-four-state-reduction}.
If $H^2(Y;\mathbb Q)=0$, all degree-two classes of $X$ are relative
and their squares vanish. Otherwise the split integral homology
inclusion into $H_2(X;\mathbb Z)=\mathbb Z$ is an isomorphism.
The nonzero square forces the same conclusion in degree four.
Every other reduced homology group of $Y$ is zero because $Y$
is a retract of $X\simeq\mathbb CP^2$, and its fundamental group
is trivial for the same reason. The inclusion $Y\to X$ is an
integral homology equivalence between simply connected CW spaces,
so it is a homotopy equivalence. This obtains the actual ordinary
act of the assertion; it does not replace finite type by rational
homology alone.

It remains that $\Gamma$ is disconnected, is not a forest, and
has at least one nonloop edge, since $P\ne Q$.
The bound $|U|+|S\setminus U|\leq5$ forces exactly three
vertices and two edges. There is one cyclic component and one
acyclic component, so $T$ has two points and $W=H_1(\Gamma;\mathbb Q)$
has dimension one. The cyclic component is either a two-edge circle,
a loop with a connecting tail, or an isolated loop with a separate
interval component. These exhaust two edges in a disconnected graph
with a nonloop edge and a cycle.

If some actor has degree different from $+1$ on $W$, its scalar
relation kills $W\otimes_{\mathbb QM}\mathbb Q$.
Lemma~\ref{C:graph-component-eigenvalues} then makes every degree-two
class pull back from the ordinary two-state $T$-act cone, whose
products vanish by Lemma~\ref{C:finite-line-act-cups}.
We may therefore assume that every actor induces degree $+1$.
Every actor must preserve the unique cyclic component; mapping it
to the other component would induce degree zero.

When that component contains a loop edge, degree $+1$ forces
the unique loop edge to map to itself, hence forces its vertex
to be fixed by all of $M$. That vertex is also fixed by $R$.
The opposite form of Lemma~\ref{C:fixed-incoming-cups} makes
every degree-two product zero, a contradiction.

The only remaining case is a two-edge circle and an isolated
vertex. If the circle edges have the same orientation, a degree-one
actor fixes both vertices and both edges; an edge interchange has
degree minus one and a fold has degree zero. If they have opposite
orientations, a degree-one actor is either the identity or the
half-rotation interchanging both vertices and edges.
Identify the circle with the unit circle so its two vertices are
$1$ and $-1$. The map $z\mapsto z^3$ fixes these vertices and
commutes with every permitted actor on the cyclic component.
Extend it by the identity on the isolated vertex. An actor can
send that isolated vertex only to a vertex of the graph, all of
which the self-map fixes. Thus this is a strictly $M$-equivariant
graph map, inducing the identity on $T$ and multiplication by
three on $W$.

The ordinary two-state $T$-act cone has zero degree-two products.
Lemma~\ref{C:graph-component-eigenvalues}, with $d=3$ and
$\dim H^2(X;\mathbb Q)=1$, contradicts the projective-plane square.
Thus the remaining disconnected cases cannot occur, leaving only
the actual ordinary-act reductions above.
\end{proof}

\begin{corollary}[The exact small-state right-zero-corner minimum]
\label{C:right-zero-five-state-minimum24}
Among triangular two-object categories with the fixed corner
$R=\{1,p,q\}$ above and at most five cross arrows, the least
arrow count realizing $\mathbb CP^2$ is twenty-four.
\end{corollary}
\begin{proof}
The preceding theorem gives an ordinary act cone with the same
source monoid and at most as many states. Its arrow count is at
least twenty-two by Theorem~\ref{C:act-cone-minimum-twenty-two}.
Restoring the three-element corner in place of the singleton adds
two arrows, so the original category has at least twenty-four.

For equality take either twenty-two-arrow act-cone model of
Theorem~\ref{C:CP2-upper22}, with an eighteen-element source
monoid and three states. Replace its trivial sink corner by $R$
and give all elements of $R$ the identity action on those states.
Its edge action category is the disjoint union of copies of $R$;
the projection to the state set is equivariantly a homotopy
equivalence because $BR$ is contractible. The collage is therefore
equivalent to the original act cone, and has $18+3+3=24$ arrows.
\end{proof}

% Root and B full literal audits passed; numerical records are controls only.
\subsection{Three cross states opposite a corner of order at most four}

\begin{lemma}[A literal unicyclic graph correspondence]
\label{C:unicyclic-graph-correspondence}
Let a finite monoid $M$ act on a connected finite graph $\Gamma$ by
maps that send every oriented edge at unit speed to one oriented edge,
or collapse it to a vertex. Suppose $\dim H_1(\Gamma;\mathbb Q)=1$.
Then the cofiber of $\Gamma_{hM}\to BM$ cannot have type
$\mathbb CP^2$.
\end{lemma}
\begin{proof}
If the second rational homology of the cofiber is nonzero,
Lemma~\ref{C:connected-graph-cone-degree} forces every actor to
induce degree $+1$ on $H_1(\Gamma;\mathbb Q)$.
The graph has a unique circuit $C$, with, say, $l$ edges.
The image of its circuit word is a closed walk of length at most
$l$. Degree $+1$ requires a traversal of that circuit exactly
once; there is no length available for a tree detour or a
collapsed circuit edge. Thus $C$ is strictly invariant, and
every actor restricts to an orientation-preserving cellular
automorphism of $C$.

The inclusion $C\to\Gamma$ is an actual equivariant underlying
homotopy equivalence. It remains an equivalence on action bars
and their cofibers. Identify the equally subdivided circuit with
the unit circle. Its actor image is a group of rotations by
multiples of $2\pi/l$, so $z\mapsto z^{l+1}$ commutes strictly
with every actor. The same scalar-degree lemma makes all positive
rational products on the cofiber zero. This contradicts the
projective-plane square. The unit-speed hypothesis is used in
the closed-walk length argument.
\end{proof}

\begin{lemma}[Three states over a two-element right-zero minimal ideal]
\label{C:right-zero-minimal-two-three-states}
Suppose that a finite monoid $N$ has minimal ideal
$K=\{p,q\}$ with $ab=b$ for $a,b\in K$.
For every $(N,M)$-biset $S$ with $|S|\leq3$, a collage of type
$\mathbb CP^2$ yields an ordinary right $M$-act cone on three
states of that actual homotopy type.
There is no bound on the orders of $N$ and $M$.
\end{lemma}
\begin{proof}
Every $n\in N$ fixes each element of $K$ by left multiplication:
for $k,h\in K$, $nk=n(hk)=(nh)k=k$.
The left actions $P,Q$ of $p,q$ are retractions onto a common
subset $U$, and every $N$-actor fixes $U$ pointwise.
The right $M$-action preserves $U$ and commutes with both
retractions. The monoid $N$ has a right zero, so $BN$ is
contractible.

If $|U|=1$, its unique state is fixed by both monoids and the
common-fixed-state argument kills the degree-two products.
If $U=S$, the left action is trivial and the collage is actually
equivalent to the ordinary $S$-act cone.
If $P=Q$, the common map is an $N$-equivariant retraction:
$nP=P$ and $Pn=P$, the latter since $pn\in K$ and both elements
of $K$ have the same action. It is also $M$-equivariant.
The core collage is the ordinary $U$-act cone, and the relative
quotient is a suspension. Here $|U|\leq2$, whose ordinary act
has zero degree-two products. A one-dimensional degree-two
space on the original collage either pulls back from this
retract or lifts to its suspension quotient; its square is
zero in either case.

The sole remaining case has $U=\{u_0,u_1\}$ and
$S=U\sqcup\{z\}$, with $Pz=u_0$, $Qz=u_1$.
Let $Z=B(N\ltimes S)$ be the actual left action category.
It is connected and simply connected. To see the latter
directly, pass to its fundamental groupoid. At each $u_i$
all endomorphism labels become identities, since their
endomorphism monoid is $N$ and $BN$ is contractible.
For any arrow $n:z\to nz$, compose with the arrow labeled
$p$ from $nz$. The resulting label $pn$ lies in $K$.
If $nz=z$, distinctness of $Pz,Qz$ forces $pn=p$, so
$p_z n_z=p_z$ and $n_z$ is an identity in the groupoid.
If $nz=u_i$, the same equation identifies $n_z$ with
the unique distinguished arrow $p_z$ or $q_z$ to $u_i$.
Thus the groupoid has precisely the two tree arrows from
$z$ to $u_0,u_1$, their inverses, and their composites;
there is no nontrivial loop.

In particular $\widetilde H_0(Z;\mathbb Z)=H_1(Z;\mathbb Z)=0$.
The augmented cellular or simplicial chains of $Z$, tensored
with the free left $M$-bar resolution, compute the projection
$Z_{hM}\to BM$. Its cofiber has zero second homology:
the reduced coefficient complex starts in homology degree two,
and the cofiber suspends it once. Since $BN$ is contractible,
this is the original collage. It therefore cannot be a
projective plane. Only the actual ordinary-act case remains,
and one or two states in it are excluded by the finite-line
lemma.
\end{proof}

\begin{theorem}[All opposite corners of order at most four on three states]
\label{C:small-opposite-corner-three-states}
Let $|N|\leq4$ and $1\leq|S|\leq3$ in a triangular biset
category with corners $M,N$. If its classifying space has type
$\mathbb CP^2$, then an ordinary three-state right $M$-act cone
has that actual homotopy type. In particular $|M|\geq18$.
For each fixed $n\in\{1,2,3,4\}$, the least total arrow count
among such categories with $|N|=n$ is exactly $21+n$.
\end{theorem}
\begin{proof}
When the trivial right rational $N$-module is projective, apply
Lemma~\ref{C:three-state-exact-coinvariants}.
Otherwise the minimal Rees ideal has at least two columns.
The order bound forces its maximal group and row set to be
singletons: a two-row or nontrivial-group ideal with two
columns already has at least four elements, and its identity
would have to be outside it. Hence the minimal ideal is a
right-zero band of size two or three.
The size-two case is covered by the preceding lemma.
In the size-three case, $N$ is exactly that right-zero band
with identity adjoined.

For this last monoid, the common image $U$ of its three
minimal retractions is fixed pointwise by $N$.
The cases $|U|=1$, $U=S$, or all three retractions equal
are dealt with exactly as in the preceding proof.
Otherwise $S$ has two core states and one outside state.
The left action category has a natural graph model with
these three states as vertices and one edge from the outside
state to its image under each of the three right-zero labels.
Indeed the endomorphism categories at core states are copies
of the contractible $BN$; postcomposition fixes every
right-zero edge label. Collapsing these categories yields
the graph naturally for every commuting right action, by
the same nerve-coordinate construction as
Lemma~\ref{C:right-zero-action-graph}.
Both core states occur among the three edge endpoints, so
this is a connected three-vertex, three-edge graph.
Every right actor sends an edge at unit speed to the edge
with the same label, or collapses it to a core vertex.
Lemma~\ref{C:unicyclic-graph-correspondence} excludes this case.

The resulting ordinary act has three states, so the exact
three-state act-cone minimum gives $|M|\geq18$.
It follows that the original total is at least $18+3+n=21+n$.
For equality use the accepted twenty-two-arrow ordinary act
model, whose source has order eighteen and whose act has three
states. For $n\geq2$, replace its trivial sink by the
identity-adjoined right-zero band of size $n-1$, acting
trivially on the three states. Its classifying space is
contractible, so the natural edge-bar projection gives an
actual equivalence with the original act cone. For $n=1$
use that cone itself. These examples have exactly $21+n$
arrows.
\end{proof}

\begin{remark}[Exact small-corner inventory and a bounded control]
\label{C:small-opposite-corner-inventory}
The elementary Rees argument above is sufficient for the theorem.
As an independent bounded control,
\path{code/a_small_nonprojective_corners.py} expands the complete
own order-two through order-four monoid catalogs under opposition
and ordinary isomorphism. There are respectively $2,7,35$
ordinary classes, and $0,1,6$ nonprojective right-trivial cases.
The six order-four cases consist of the identity-adjoined
right-zero three-band and five extensions of a right-zero
two-band $K$ by an extra actor $x$ and the identity.
Their $(x^2,\tau_x)$ possibilities are
\[
 (1,\mathrm{id}),\quad(1,\mathrm{swap}),\quad
 (x,\mathrm{id}),\quad(x,\mathrm{constant}),\quad
 (k,\mathrm{constant}_k).
\]
The exact nonprojective tables and their Rees data are recorded in
\path{code/outputs/a_small_nonprojective_corners.json}.
A separate full three-state left-action control is
\path{code/a_four_corner_three_states.py}, with output
\path{code/outputs/a_four_corner_three_states.json}.
It exhausts all labeled actions and retains the counts and literal normalized
$d_1,d_2$ matrices in the two-core, distinct-retraction branch;
all twenty-four such instances for the two-band extensions
have $H_0=\mathbb Q$ and $H_1=0$. These calculations are
controls of the structural proof, not premises for it.
\end{remark}

% Root and B full literal audits passed; no numerical premise.
\subsection{Four cross states opposite the small nonprojective corners}

\begin{lemma}[The minimal right-zero graph surjects on fundamental groupoids]
\label{C:right-zero-graph-fundamental-quotient}
Let a finite monoid $N$ have minimal ideal $K=\{p,q\}$, a
right-zero band, and let $S$ be an $(N,M)$-biset.
The actual inclusion
\[
 (K^1\ltimes S)\longrightarrow(N\ltimes S)
\]
induces a bijection on components and a surjection on every
fundamental-groupoid hom set. These assertions and the canonical
maps from action categories to their fundamental groupoids are
natural for the commuting right $M$-action.
\end{lemma}
\begin{proof}
For any arrow $n_s:s\to ns$ of the larger action category,
composition with $p_{ns}:ns\to pns$ has label $pn\in K$.
After inverting arrows,
\[
 n_s=p_{ns}^{-1}(pn)_s.
\]
Both arrows on the right belong to the smaller action category.
Thus each new arrow joins objects already joined there, and
every path in the larger fundamental groupoid is the image
of a path in the smaller one. The identity-on-objects inclusion
and the defining equations commute with every right actor.
Groupoid completion itself is functorial, giving all the stated
naturality.
\end{proof}

\begin{lemma}[The low-degree component and groupoid comparisons]
\label{C:action-space-low-degree-comparison}
Let a finite monoid $M$ act on a CW space $Z$, with the usual
natural action-bar projection to $BM$.
If every component of $Z$ has zero first rational homology,
projection to $T=\pi_0Z$ gives a surjection
\[
 H^2\bigl(\operatorname{hocofib}(T_{hM}\to BM);\mathbb Q\bigr)
 \longrightarrow
 H^2\bigl(\operatorname{hocofib}(Z_{hM}\to BM);\mathbb Q\bigr).
\]
If $Z$ is connected, then
\[
 H_2\bigl(\operatorname{hocofib}(Z_{hM}\to BM);\mathbb Q\bigr)
   =H_1(Z;\mathbb Q)\otimes_{\mathbb QM}\mathbb Q.
\]
Consequently a natural map of connected coefficient spaces
inducing an isomorphism on $H_1$ induces an isomorphism on the
second homology of these actual cofibers. In particular this
applies to the classifying map from an action category to its
fundamental groupoid.
\end{lemma}
\begin{proof}
Use the augmented singular or simplicial chain complex and
tensor it with the free left monoid bar resolution. Its
suspension is the reduced cofiber chain model. The kernel of
the projection onto $\widetilde{\mathbb QT}$ has zero
degree-zero homology, with first homology $H_1(Z;\mathbb Q)$.
The first-quadrant total complex therefore gives the displayed
surjection when that first homology vanishes.
For connected $Z$, the only term contributing to total degree
two after suspension is its first homology in bar degree zero;
there is no lower nonzero coefficient row from which a
differential could enter. This gives the tensor formula and
its naturality. These statements require neither bounded
homology nor finite CW type for $Z$.
\end{proof}

\begin{theorem}[Four-state reduction for an arbitrary two-element right-zero ideal]
\label{C:right-zero-minimal-two-four-states}
Let a finite monoid $N$ have minimal ideal a two-element
right-zero band. For every $(N,M)$-biset $S$ with $|S|\leq4$,
a collage of type $\mathbb CP^2$ yields an ordinary right
$M$-act cone on at most $|S|$ states of that actual type.
For each such fixed $N$, the least arrow count among these
collages realizing $\mathbb CP^2$ is exactly $21+|N|$.
\end{theorem}
\begin{proof}
Write $P,Q$ for the two minimal retractions, with common
image $U$. All $N$-actors fix $U$ pointwise, and $BN$ is
contractible. The cases with at most three states were proved
in Lemma~\ref{C:right-zero-minimal-two-three-states}.
A singleton $U$ gives a common fixed state; $U=S$ gives an
actual ordinary act cone.

If $P=Q$, this common map is a biset retraction onto $U$:
$nP=P$ and $Pn=P$ follow respectively from $nK=K$
pointwise and $Kn\subseteq K$. The core collage $Y$ is an
ordinary $U$-act cone and a homotopy retract of the original
collage $X$. The relative quotient is a suspension.
If $H^2(Y;\mathbb Q)=0$, the projective-plane square is
therefore impossible. Otherwise the split integral homology
argument used in Theorem~\ref{C:right-zero-five-state-reduction}
makes $Y\to X$ an integral homology equivalence: degree two
is an isomorphism by splitting, degree four is nonzero by the
square and hence an isomorphism by splitting, and all other
reduced homology vanishes. Both spaces are simply connected
because $Y$ is a retract of $X$. Homological Whitehead gives
the required actual equivalence.

Suppose now that $P\ne Q$, and use the natural graph $\Gamma$
of the $K^1$ action. If $|U|=3$, its single edge is nonloop,
so its two components are trees. By
Lemma~\ref{C:right-zero-graph-fundamental-quotient}, the space
$Z=B(N\ltimes S)$ has the same two components and each is
simply connected. Its component set $T$ is a two-state right
act. Lemma~\ref{C:action-space-low-degree-comparison} makes
the second cohomology of its ordinary act cone surject onto
that of the collage. The finite-line lemma kills every square.

The remaining graph has two vertices and two edges and is
connected. Its fundamental group is infinite cyclic. The
same groupoid-surjection lemma says that $Z$ is connected
and its fundamental group is a quotient of $\mathbb Z$.
If this quotient is finite, $H_1(Z;\mathbb Q)=0$ and the
collage has zero second rational homology. If it is infinite,
the actual composite
\[
 B(K^1\ltimes S)\longrightarrow B(N\ltimes S)
       \longrightarrow B\Pi(N\ltimes S)
\]
is a weak equivalence: its domain is the unicyclic graph and
its target is a connected $K(\mathbb Z,1)$, and the induced
surjection $\mathbb Z\to\mathbb Z$ is an isomorphism.
This entire comparison is $M$-equivariant. Taking bars gives
an equivalence on cofibers between the graph cofiber and the
fundamental-groupoid cofiber. Meanwhile the actual map from
the original cofiber to the latter induces an isomorphism on
$H_2(-;\mathbb Q)$ by
Lemma~\ref{C:action-space-low-degree-comparison}.

A nonzero second homology forces every right actor to have
degree $+1$ on the graph cycle. The proof of
Lemma~\ref{C:unicyclic-graph-correspondence} makes every
positive rational product on its cofiber zero. Pullback
along the actual second-homology isomorphism therefore kills
the projective-plane square on the original collage. This
excludes the last case.

The ordinary-act reduction and its exact minimum twenty-two
give the lower bound $21+|N|$. For equality let $N$ act
trivially on the three states of the accepted eighteen-element
source model. Since $BN$ is contractible, its collage is
actually equivalent to the twenty-two-arrow act cone and
has $18+3+|N|$ arrows.
\end{proof}

\begin{theorem}[Four states opposite the three-element right-zero band]
\label{C:right-zero-three-four-states}
Let $N$ be the identity-adjoined three-element right-zero band,
and let $|S|\leq4$. A collage of type $\mathbb CP^2$ yields
an ordinary right $M$-act cone on at most $|S|$ states of that
actual type. In particular this fixed opposite corner has
exact least arrow count twenty-five in the stated state range.
\end{theorem}
\begin{proof}
Use the natural graph with core vertices $U$, the outside
states as additional vertices, and one edge from each outside
state for each of the three right-zero labels. Its core is
fixed pointwise by $N$. A singleton core gives a common fixed
state; a trivial action gives an ordinary act. If the three
retractions coincide, their common biset retraction reduces
a projective plane to an actual ordinary act by the same
split-homology argument as above.

With three core states and one outside state $z$, three
distinct endpoints give a tree, hence a contractible collage.
Two distinct endpoints give a unicyclic component and one
isolated core vertex. If any right actor sends $z$ into the
core, it kills the cycle class, so its coinvariants vanish.
The component comparison of
Lemma~\ref{C:graph-component-eigenvalues} makes every degree-two
class pull back from a two-state act cone, and its square is
zero. Otherwise every right actor fixes $z$ and both of its
core endpoints; either endpoint is a common fixed state.

It remains that $U=\{u_0,u_1\}$ and there are two outside
states $z,w$. Record the three retraction values at an outside
state as a word $v_s\in\{0,1\}^3$. At least one of the two
words is nonconstant, or all minimal retractions coincide.
The graph is therefore connected and has first Betti number
three. If a right actor induces the vertex map
$\sigma:\{0,1\}\to\{0,1\}$, commutation says that the word
of the image of $s$ is $\sigma\circ v_s$; core states have
words $000$ and $111$.

Suppose first that both outside words are nonconstant.
A constant vertex map then sends every state to that constant
core vertex, and kills the entire first homology. It would
make the collage's second homology zero, so such actors may
be excluded in the putative projective-plane case.
If every vertex map is the identity, there is a common fixed
core state. Otherwise a vertex swap occurs, which forces
$v_w$ to be the complement of $v_z$. The words are distinct,
so an actor fixing the core fixes all states, while an actor
swapping it is the unique simultaneous interchange of
$u_0,u_1$ and $z,w$. The graph action image is exactly $C_2$.
This involution acts freely on its four vertices and six
edges. Its quotient graph has two vertices and three edges,
so the rational first-homology coinvariants have dimension
two. The connected graph-Tor comparison gives
$\dim H_2(X;\mathbb Q)=2$, again excluding a projective plane.

Finally suppose, after relabeling, that $v_z$ is nonconstant
and $v_w=000$. A vertex swap is impossible, since the
complement of $v_z$ is not the word of either outside state.
A constant-one vertex map collapses all states to $u_1$ and
kills first homology. Unless the second homology is zero,
all vertex maps are therefore the identity or constant zero;
they all fix $u_0$. The common-fixed-state argument kills
the degree-two products. These cases are exhaustive.
The lower and upper twenty-five bounds follow from the
ordinary-act reduction and the trivial-action construction.
\end{proof}

\begin{corollary}[The nonprojective order-four opposite corners]
\label{C:nonprojective-four-corners-four-states}
For every four-element monoid $N$ whose trivial right rational
module is nonprojective, a collage with at most four cross
states and type $\mathbb CP^2$ has at least twenty-five arrows.
The bound is attained for each of the six such monoids.
In particular the seventeen-arrow profile $(9,4,4)$ is
excluded, including all nonprojective opposite corners.
\end{corollary}
\begin{proof}
The elementary inventory in
Theorem~\ref{C:small-opposite-corner-three-states} leaves
precisely a two-element or a three-element right-zero minimal
ideal. Apply the two preceding theorems. The projective
opposite corners in the stated seventeen-arrow profile were
already excluded by
Corollary~\ref{C:projective-opposite-corner-twenty-one}.
\end{proof}

% Root and D full literal audits passed; independent rational control passed.
\subsection{Five states for the two-band four-element opposite corners}

\begin{lemma}[Unbounded rational cup vanishing for source order at most nine]
\label{C:nine-unbounded-act-cups}
For every finite monoid $M$ of order at most nine and every
finite right $M$-act, all rational degree-two products on its
ordinary act cone vanish. No bounded-homology assumption is
needed.
\end{lemma}
\begin{proof}
If $H^2(BM;\mathbb Q)=0$, the degree-two classes lift relative
to $BM$, and the suspension quotient kills their products.
Otherwise the minimal ideal has at least two rows and two
columns. At most two rows are covered by
Theorem~\ref{C:two-row-act-cups}. With at least three rows,
the order bound forces a trivial-group band of shape
$3\times2$ or $4\times2$, with at most three outside actors.
Lemma~\ref{C:three-outside-all-two-column-act-cups} applies
without a boundedness hypothesis and proves the assertion.
\end{proof}

\begin{lemma}[The rational five-state graph obstruction for small sources]
\label{C:nine-five-state-graph-square}
Let $M$ have order at most nine and act on a graph with vertex
set $U$ and oriented edge set $E$, where $|U|+|E|\leq5$.
Every actor sends each oriented edge at unit speed to another
oriented edge or collapses it to a vertex. If the cofiber of
its action-bar projection to $BM$ has one-dimensional second
rational cohomology, the square of that class is zero.
\end{lemma}
\begin{proof}
Regard the graph as the natural action graph for the right-zero
two-band on the state set $U\sqcup E$. The connected analysis
of Theorem~\ref{C:right-zero-five-connected} uses only
one-dimensional second rational cohomology, except for its
last literal control. That same thirteen-arrow control has
an exact rational square primitive as well.
The record
\path{code/outputs/a_rightzero_five_rational_control.json}
contains a nonboundary rational two-cocycle and a three-cochain
with only zero and one as coefficients. Its $872$ signed
four-simplex equations verify the square identity.
The standalone checker
\path{code/b_verify_rightzero_five_rational_control.py}
imports no producer or matrix-library routines. It rebuilds
all signed boundaries, checks their consecutive composites,
computes rational ranks $10,39$ for the first two
coboundaries, and verifies that the supplied cocycle is
nonboundary and satisfies all $872$ primitive equations.
Its complete record is
\path{code/outputs/b_rightzero_five_rational_control_verified.json}.
Thus the rational second cohomology is one-dimensional,
and the actual graph-image second-homology
isomorphism used in that proof also gives a rational
zero-square pullback. All connected cases are excluded
under the present rational hypothesis.

For a forest, the component projection is an equivariant
underlying equivalence, reducing to an ordinary act cone;
apply Lemma~\ref{C:nine-unbounded-act-cups}.
When every edge is a loop, the common vertex retraction gives
an ordinary-act homotopy retract and a suspension quotient.
If the retract has nonzero second cohomology its pullback
spans the assumed rank-one space and squares to zero. If it
has zero second cohomology, the class is relative and again
squares to zero. The remaining disconnected graphs are
exactly the three-vertex/two-edge cases in
Theorem~\ref{C:right-zero-five-state-reduction}. Its cycle
coinvariant, fixed-vertex, and two-eigenvalue arguments use
only the current rational rank-one assumption and the
two-state finite-line lemma. They therefore apply unchanged.
No intermediate space is asserted to have finite CW type.
\end{proof}

\begin{lemma}[An explicit normalized action-groupoid presentation]
\label{C:four-corner-normalized-action-groupoid}
Let $N=\{1,p,q,x\}$ have minimal right-zero ideal
$\{p,q\}$, and let $S$ be a left $N$-set. Write $P,Q,E$
for the actions of $p,q,x$, and let $U$ be the common image
of $P,Q$. In the fundamental groupoid of $N\ltimes S$,
identify each object $s$ with $Ps$ using the arrow labeled
$p$. The resulting groupoid on $U$ has one oriented
generator $g_s:Ps\to Qs$ for each $s\notin U$, with
$g_u=1$ for $u\in U$, and precisely these relations:
\[
\begin{array}{c|c}
\text{right action of $x$ on $\{p,q\}$}&\text{relations}\\\hline
\mathrm{id}&g_{Es}=g_s\\
\mathrm{constant}_p&g_{Es}=1\\
\mathrm{swap}&g_{Es}=g_s^{-1}.
\end{array}
\]
The constant case may be arranged to be $p$ by relabeling.
The presentation and its functor from the actual action
category are natural for every commuting right action.
\end{lemma}
\begin{proof}
All $N$-actors fix $U$ pointwise. In the fundamental groupoid,
all core endomorphisms vanish because $BN$ is contractible.
Normalize an arrow $n_s:s\to ns$ by the paths labeled $p$.
The equation $p_{ns}n_s=(pn)_s$ shows that its normalized
value is the identity if $pn=p$, and is $g_s$ if $pn=q$.
Thus every arrow is represented by the displayed generators.
The same equation with $q$ gives exactly the listed relation
for $E$; the relations for $P,Q$ only say $g_u=1$.

Conversely, define a functor from the action category to the
displayed groupoid by $s\mapsto Ps$, and send the arrow
$n_s$ to $1$ or $g_s$ according as $pn=p$ or $q$.
For $n=1$, use the identity. Each equation for a composable
pair follows from the displayed relation and the right
translation identity $\tau_{ab}=\tau_b\tau_a$.
In particular the $x^2$ equation adds nothing: in the first
two cases $x_s$ is sent to the identity; in the swap case
its product is $g_{Es}g_s=1$ and $x^2=1$.
This functor and the generator assignment into the completed
action category are inverse after the stated identifications
of objects. A commuting right map preserves $P,Q,E$ and
the equations, proving naturality.
\end{proof}

\begin{theorem}[Five states for the five two-band order-four corners]
\label{C:four-corner-two-band-five-states}
Suppose $|N|=4$, its minimal ideal is a right-zero band of
size two, $|M|\leq9$, and $|S|\leq5$ in an $(N,M)$-biset.
If the collage has one-dimensional second rational cohomology,
its generator has square zero. In particular none of these
collages has type $\mathbb CP^2$.
\end{theorem}
\begin{proof}
A singleton core $U$ is a common fixed state, so assume
$|U|\geq2$. Use the preceding presentation of the actual
left action category $Z=N\ltimes S$.

If $\tau_x$ is the identity, $E$ preserves both endpoints of
each graph edge. It is either an involution or a retraction.
The relations identify parallel edges in its orbits or with
their retraction images; an edge whose image is a core vertex
is a loop and is collapsed. The resulting groupoid is the
fundamental groupoid of an oriented graph $Q$ with the same
core vertices and at most $|S\setminus U|$ edges.
The right action sends an edge to an edge with the same
orientation or to a vertex, because it commutes with $E$.

If $\tau_x$ is constant, $PE=P$ and $QE=P$ show that each
edge indexed by an element of $ES\setminus U$ is a loop.
The relations simply collapse those loops, leaving another
oriented graph $Q$ with the same vertex and edge bounds.
The collapsed set is right-invariant, again giving a strict
right graph action. In these two cases the normalized
functor $Z\to\Pi Q$ induces isomorphisms on components and
first integral homology. The natural graph map
$Q\to B\Pi Q$ is an equivariant weak equivalence.
Consequently the actual collage maps to the $Q$ cofiber
through a natural equivalence zigzag, and its second
rational cohomology is the pullback of that cofiber's second
cohomology. The low-degree augmented bar comparison gives
an isomorphism here: the cone of the augmented chain map has
zero homology through degree two, since the target graph
has no second homology and the map is an isomorphism in
degrees zero and one. Its free-bar tensor has the same
connectivity. Lemma~\ref{C:nine-five-state-graph-square}
kills the square.

Finally suppose $\tau_x$ swaps $p,q$. Then $x^2=1$, so
$E$ is an involution fixing $U$ pointwise and reversing
the endpoints of each indexed edge. Each two-element
outside orbit contributes one unoriented edge. An outside
fixed point indexes a loop with relation $g_s^2=1$.
Kill these order-two loop generators. The resulting
groupoid is the fundamental groupoid of a graph $Q$ with
vertex set $U$ and at most
$\lfloor(|S|-|U|)/2\rfloor$ edges. This quotient is natural
for right actors: they preserve the set of $E$-fixed states,
and send a two-element edge orbit to another such orbit,
a fixed loop, or a vertex. Reversing an edge is allowed in
this last graph model.

The killed loops contribute only two-torsion to the
abelianization. Hence $Z\to\Pi Q$ still induces a
component bijection and a first rational homology
isomorphism, and the actual cofiber comparison induces
an isomorphism on second rational homology.
Since $|U|\geq2$ and $|S|\leq5$, this graph has at most
one edge. If that edge is a loop, it is the sole cycle.
An actor of degree other than $+1$ kills its rational
coinvariants, giving component-cone pullback and zero
products by Lemma~\ref{C:nine-unbounded-act-cups}.
Otherwise every actor fixes the unique loop vertex, giving
a common fixed state in the original collage. Thus the
loop case also has zero square.
For the forest case, the component projection reduces its
cofiber to an ordinary act cone. Apply
Lemma~\ref{C:nine-unbounded-act-cups}. Pullback along the
actual second-homology isomorphism finishes the proof.
\end{proof}

\begin{remark}[The remaining order-four shape]
\label{C:four-corner-five-state-partial-scope}
The theorem treats five of the six nonprojective right-trivial
four-element corners, and excludes their seventeen-arrow
profile $(8,4,5)$ as well as $(9,4,5)$ at eighteen arrows.
The identity-adjoined three-element right-zero band has a
different minimal ideal and is handled separately by
Theorem~\ref{C:right-zero-three-five-small-sources} below.
No actual finite CW type is inferred for any graph-quotient cofiber.
\end{remark}

% Root and D full literal audits passed; independent rational control passed.
\subsection{Five states opposite the three-element right-zero band}

\begin{lemma}[A vertex-fixing scalar map of a double-cover graph]
\label{C:double-cover-graph-scalar}
Let a connected finite graph $A$ have a free cellular $C_2$-action
on vertices and edges. For every odd positive integer $d$, there
is a $C_2$-equivariant self-map fixing every vertex and inducing
$d$ times the identity on $H_1(A;\mathbb Z)$.
\end{lemma}
\begin{proof}
The quotient $B=A/C_2$ is a connected graph, and its connected
double cover is classified by a surjection
$\eta:F_r=\pi_1(B)\to C_2$.
Choose a free basis $a,b_1,\ldots,b_{r-1}$ with
$\eta(a)=1$ and $\eta(b_j)=0$; replacing another odd basis
element by its product with $a^{-1}$ gives such a basis.
The endomorphism $a\mapsto a^d$, $b_j\mapsto b_j^d$
preserves $\eta$. Represent it by a cellular map of $B$
that fixes a chosen maximal tree pointwise. Its lift fixing
a chosen lifted base vertex is equivariant for the deck
involution, and fixes the two lifts of every tree vertex,
hence every vertex of $A$.

A Schreier basis for the kernel of $\eta$ is
\[
 a^2,\qquad b_j,\qquad ab_ja^{-1}.
\]
The first two types are multiplied by $d$ after abelianization.
Writing $d=2k+1$, the image of the last type is
$a^{2k}(ab_j^da^{-1})a^{-2k}$, so its abelianization is
also multiplied by $d$. This is exactly the claimed action
on the cover's first homology.
\end{proof}

\begin{lemma}[Removing constant outside stars]
\label{C:right-zero-constant-star-removal}
Let $R$ be the identity-adjoined right-zero three-band acting
on the left of a finite biset $S$, with core $U$.
Suppose that every outside state on which the three minimal
actors agree has its common image at the same $u\in U$.
Collapse all those states to $u$. This gives an actual right
$M$-equivariant quotient of the natural action graph.
Either $u$ is fixed by all of $M$, or the induced map on
collage cofibers is an isomorphism on second rational
cohomology by pullback.
\end{lemma}
\begin{proof}
Let $A$ be the vector space on these outside states, with
the partial right action that sends a basis state to zero
when it enters the core. Such a state cannot map to a
nonconstant outside state, by commutation with the three
retractions. If it maps to another constant outside state,
the image of $u$ is $u$; if it maps into the core, that
image is also the image of $u$. Thus the quotient relation
is right-invariant.

The graph-chain kernel of the collapse has homology only
in degree one, where it is $A\otimes\widetilde{\mathbb Q^3}$.
The second factor records the three edge colors and is
fixed by every right actor. If $A\otimes_{\mathbb QM}\mathbb Q$
is nonzero, its state-transition relations have a nonempty
component with no transition into the core. Each actor
therefore sends a state of that component to another
constant outside state attached to $u$; commutation forces
that actor to fix $u$. Hence $u$ is a common fixed state.
Otherwise the kernel's first-homology coinvariants are zero.
After the free bar tensor and the cofiber suspension, the
kernel complex has zero homology in both degrees one and
two. Its homology exact sequence therefore makes the actual
cofiber comparison an isomorphism in degree two. Duality
over $\mathbb Q$ proves the stated cohomology isomorphism.
\end{proof}

\begin{theorem}[The remaining four-element corner on five states]
\label{C:right-zero-three-five-small-sources}
Let $R$ be the identity-adjoined three-element right-zero band,
let $|M|\leq9$, and let $|S|\leq5$ in an $(R,M)$-biset.
If its collage has one-dimensional second rational cohomology,
that class has square zero. In particular the collage cannot
have type $\mathbb CP^2$.
\end{theorem}
\begin{proof}
Use the natural graph with core $U$, one additional vertex
for every outside state, and three colored edges from that
state to its three retraction images. Right actors preserve
colors and send an edge at unit speed to another edge or
collapse it. A singleton core gives a common fixed state.
If all three retractions coincide, their common biset
retraction gives an ordinary-act retract and a suspension
quotient. Lemma~\ref{C:nine-unbounded-act-cups}, together
with the assumed rank one in degree two, kills the square.
This also includes the trivial left action.

The proof of Theorem~\ref{C:right-zero-three-four-states}
works with this same rational rank-one hypothesis when
$|S|\leq4$: replace its ordinary-act conclusion in the
common-retraction branch by that universal zero-product
lemma. All its other branches already give zero products
or a contradiction in second rational homology.
It therefore remains to treat exactly five states.

If constant outside states occur only at one core vertex,
Lemma~\ref{C:right-zero-constant-star-removal} either gives
a common fixed state or identifies the second cohomology
with that of a graph on at most four states. Its rank remains
one, so the preceding
paragraph finishes that case. If constant outside states
occur at two different core vertices and a nonconstant
outside state also exists, the state budget forces exactly
two core vertices, one constant outside state at each, and
one nonconstant outside state. We treat this configuration
last. Hence for the moment all outside states may be
assumed nonconstant.

\emph{Two core vertices.}
There are three outside states. Each has a nonconstant
binary word $v_s\in\{0,1\}^3$ recording its colored
endpoints. A vertex-constant actor sends every state to
that constant core vertex, so kills all first homology
of the connected graph and makes the collage's second
homology zero. Without such actors, the core actions are
identity or swap. If no swap occurs, a core vertex is a
common fixed state.
If a swap occurs, the distinct outside words occur in
complementary pairs. There are at most three words, so
there is exactly one complementary pair. Identify all
outside vertices having the same word, preserving each
colored edge. This is an equivariant graph quotient and
is onto on first homology: choosing one representative
for each word gives an underlying graph section.
The quotient graph has four vertices and six edges, and
its right-action image is the free $C_2$ action interchanging
both pairs of vertices. Its orbit graph has two vertices
and three edges, so its first rational homology coinvariants
have dimension two. The original connected graph's
coinvariants surject onto them. The graph-Tor comparison
contradicts one-dimensional second homology.

\emph{Three core vertices.}
There are two nonconstant outside words $v_z,v_w$.
If one word uses all three core vertices, relabel it as
the identity word. If the other also uses all three, a
nonidentity core actor must carry the first word to the
second and vice versa. It is a transposition on the three
core vertices; otherwise its square could not be a word
present in the action. A core vertex is therefore fixed
by every actor. If the two full words agree, all core
actors are already the identity.
If the other word uses only two vertices, the only
possible nonconstant core actors are the identity and
that one rank-two map $\sigma$. A constant core actor
collapses every state, killing second homology.
Closure then forces $\sigma^2=\sigma$; a constant square
would give the same contradiction, and a rank-three square
is impossible. Thus a vertex in its nonempty image is fixed
by every actor. This disposes of every full word.

Suppose both words use two core vertices. If they use the
same pair, the graph has a cyclic component on those two
core vertices and the two outside states, and one isolated
core vertex. An actor collapsing that component's first
homology makes its coinvariants zero, so degree-two classes
pull back from the two-state component act and square to
zero. Otherwise the four-state word analysis shows either
a common fixed core vertex, or complementary words with
the free $C_2$ graph action just described.
In the latter case Lemma~\ref{C:double-cover-graph-scalar}
gives a degree-three scalar map on the component fixing
all its vertices. Extend it by the identity on the isolated
vertex. This is strictly equivariant: actors restrict to
the identity or deck involution on the component, and
send the isolated core vertex only to core vertices.
The component-eigenvalue lemma and the two-state finite-line
lemma kill the rank-one square.

If the endpoint pairs differ, their union is all three
core vertices, and the graph is connected with two
independent cycles, one for each outside state. If the
two words repeat different pairs of color positions, an
actor cannot send one outside state to the other: applying
a vertex map preserves equal color positions. Each cycle
therefore carries its own zero-or-one character, recording
whether that outside state is fixed or enters the core.
Nonzero coinvariants require at least one outside state
to be fixed by all actors. Its core endpoints are then
common fixed states, again killing products.

If the two words repeat the same pair of color positions,
their cycle module is the partial permutation module on
$\{z,w\}$. A nonzero coinvariant class determines a
nonempty transition component that is invariant and never
enters the core. A singleton such component gives a fixed
outside state and hence a fixed core endpoint.
Otherwise both outside states stay outside under every
actor. Regard each as an oriented edge from its repeated
endpoint to its remaining endpoint. The two edges form
an oriented path on the three core vertices. Every actor
sends each edge to an edge with its orientation, without
collapse. For the two possible star orientations the
middle vertex is fixed, and for a directed path only the
identity edge assignment is possible. Thus the middle
vertex is a common fixed state in every orientation.

\emph{Four core vertices.}
There is one nonconstant outside state. Three distinct
endpoints give a forest, whose component quotient is an
ordinary act cone and has zero degree-two products.
Two distinct endpoints give one unicyclic component.
If an actor sends the outside state to the core it kills
the cycle coinvariants; the class then pulls back from
the component act, and Lemma~\ref{C:nine-unbounded-act-cups}
applies. Otherwise that outside state is fixed by every
actor, and either of its endpoints is a common fixed state.
This finishes the case of no constant outside states.

\emph{Two constant stars and one nonconstant state.}
Write the core as $u_0,u_1$, the constant outside states
as $w_0,w_1$, and the remaining state as $z$.
The word at $z$ has two equal entries and one other entry.
A core swap is impossible because its complementary word
is absent. Without both constant core maps, a core vertex
is fixed, so assume both occur.

Put $V=\widetilde{\mathbb Q^3}$, the color-difference space.
The first homology of the connected graph decomposes as
$V_{w_0}\oplus V_{w_1}\oplus L_z$, where $L_z\subset V$
is the one-dimensional line comparing the equal color
positions of $z$. A nonzero invariant functional has
restrictions $\lambda_0,\lambda_1,\lambda_z$ to these
summands. For an actor constant at $u_i$, each $w_j$ or
$z$ is sent either to $u_i$ or $w_i$.
Invariance under the constant actors at both vertices
forces $\lambda_0=\lambda_1=\lambda\ne0$, and forces
both $w_0,w_1$ to be sent to $w_i$ by every such actor.
It also forces every core-fixing actor to fix both $w_j$
and $z$, hence to be the identity on all five states.

Every remaining actor is therefore $c_{i,\delta}$:
it sends the two core states to $u_i$, both constant
outside states to $w_i$, and $z$ to $u_i$ if $\delta=0$
or $w_i$ if $\delta=1$. Invariance says
$\lambda_z=\delta\lambda|_{L_z}$.
If all actors have the same $\delta$, the invariant
functional space has dimension two, since $\lambda\in V^*$
is arbitrary. This contradicts the rank-one assumption.
Otherwise closure under
$c_{i,\delta}c_{j,\epsilon}=c_{j,\delta}$ gives all
four actors and the identity. It forces
$\lambda|_{L_z}=0$ and $\lambda_z=0$; their invariant
functional space is exactly one-dimensional.

Now identify in $R$ the two right-zero labels occupying
the repeated positions of the word at $z$. The left action
on every state factors through this actual quotient
$R\to\{1,p,q\}$. Simultaneously replace $M$ by its
five-element action image. The resulting actual collage
is precisely the thirteen-arrow control of
Lemma~\ref{C:nine-five-state-graph-square}.
On first graph-homology coinvariants its pullback identifies
the target generator with the line
$\lambda_0=\lambda_1=\lambda$, $\lambda|_{L_z}=0$,
$\lambda_z=0$ just computed. The actual functor therefore
induces an isomorphism on second rational cohomology.
The signed rational square primitive of that literal
control pulls back and finishes the last case.
\end{proof}

\begin{corollary}[All order-four opposite corners at the first five-state profiles]
\label{C:four-corners-five-state-eighteen}
Every triangular two-object category with $|N|=4$,
$|S|\leq5$, and $|M|\leq9$ is excluded from type
$\mathbb CP^2$. In particular the seventeen-arrow profile
$(8,4,5)$ is excluded completely.
\end{corollary}
\begin{proof}
Projective trivial right modules are covered by the accepted
finite-target opposite-corner comparison and
Theorem~\ref{C:finite-target-small-act-twenty}.
The five nonprojective corners with a two-element minimal
band are covered by
Theorem~\ref{C:four-corner-two-band-five-states}.
The remaining corner is the one treated above.
\end{proof}

\subsection{Three cross states and small opposite corners}

\begin{lemma}[Exact coinvariants with at most three states]
\label{C:three-state-exact-coinvariants}
Let $N$ have projective trivial right rational module, and let
$S$ be an $(N,M)$-biset with $1\leq|S|\leq3$.
If its collage has type $\mathbb{CP}^2$, an ordinary right
$M$-act cone on three states has that type.
In particular $|M|\geq18$.
\end{lemma}
\begin{proof}
The actual orbit-biset functor replaces $N$ by the trivial
monoid and $S$ by the right $M$-set $N\backslash S$.
Exact rational coinvariants give a rational homology and
cohomology-ring equivalence on the edge bars and vertex spaces,
hence on the collages.  If the $N$-action is nontrivial, its
orbit set has at most two states; the finite-line lemma makes
every positive rational product zero.

If the $N$-action is trivial, this collapse has an actual
section: include the trivial monoid into $N$ by its identity
and keep the $M$ vertex and cross arrows unchanged.  Thus the
ordinary act cone is a rationally equivalent homotopy retract
of a hypothetical projective plane.  Its integral homology is
a direct summand in each degree of that of the projective plane.
The rational equivalence forces each such split map to be an
integral isomorphism, since those groups are either zero or
$\mathbb Z$.  Its fundamental group is also trivial, being a
retract of the trivial group.  Homological Whitehead gives an
actual homotopy equivalence with the projective plane.
One or two states are excluded by the finite-line lemma;
Theorem~\ref{C:three-state-cone-minimum22} gives $|M|\geq18$
in the remaining case.
\end{proof}

\begin{proposition}[The five-element sphere corner on three states]
\label{C:P5-three-state-bisets}
Let $P=\{1\}\sqcup(\{0,1\}\times\{0,1\})$ be the
five-element sphere monoid, let $N$ be any finite monoid, and
let $S$ be a nonempty $(P,N)$-biset with $|S|\leq3$.
If $H^2(BN;\mathbb Q)=0$, its collage cannot have type
$\mathbb{CP}^2$.  When $|S|=3$ and the left $P$-action has
no fixed state, the same conclusion holds without the
hypothesis on $H^2(BN;\mathbb Q)$.
\end{proposition}
\begin{proof}
One state gives a join; two states are covered by
Lemma~\ref{C:two-state-biset-H2zero}.  Suppose there are three
states.  If one is fixed by $P$, it gives an actual section
$BP\to E_S$ of the edge projection to $BP$.  If
$H^2(BN;\mathbb Q)=0$, the complete vertex degree-two group
therefore injects into the edge group.  All degree-two classes
of the collage are relative to its vertices, and their products
vanish because that relative quotient is a suspension.

It remains to prove the stronger assertion with no $P$-fixed
state.  Write the band elements as $k_{ij}$.  Each left orbit
of the band in its image has one or two states; a one-point
orbit would be a fixed state.  There is therefore one two-point
core $U=\{u_0,u_1\}$ and a third state $z$.  The idempotent
identity $k_{ij}(k_{ij}s)=k_{ij}s$ forces each band actor to
be constant on all three states, with
\[
 k_{ij}s=u_i\qquad(s\in S).
\]
Commutation makes every right $N$-actor fix both $u_0,u_1$.
It can only fix $z$ or send it to one of these core states.
Define $\chi_n=1$ in the first case and $\chi_n=0$ otherwise;
put $\beta_n=0$ when $\chi_n=1$, and put $\beta_n=i$ when
$zn=u_i$.  The action equations give
\begin{equation}\label{C:P5-three-state-affine}
 \chi_{mn}=\chi_m\chi_n,\qquad
 \beta_{mn}=\beta_m+\chi_m\beta_n.
\end{equation}

We need the actual equivariant model of the left action bar
$Z=B(P\ltimes S)$.  The subcategory on $U$ has contractible
classifying space.  Indeed the arrows $k_{i0}:u_0\to u_i$
give a natural transformation from the constant functor at
$u_0$ to its identity functor.  Every arrow leaving $z$ has
label $k_{ij}$, and postcomposition preserves $j$.  Thus its
nerve is the union of two cones on the core action category,
one for each column $j$, with their cone vertices both equal
to $z$.  Their intersection consists of the core category
and the singleton $z$.

Collapse the contractible core category to one vertex.  This
gives an equivalence to the graph with two vertices and two
parallel edges, hence a circle.  It is $N$-equivariant: if
$zn=z$, the actor is the identity on the whole action category;
otherwise it sends the whole category into its core, so it
acts constantly on the quotient graph.  In particular its
action on the circle's first homology is $\chi$.
The quotient circle has a fixed core vertex, and its cellular
chain complex splits over $\mathbb QN$ as
\[
 \mathbb Q\ \oplus\ \mathbb Q_\chi[1]
\]
plus a contractible two-term summand.  The edge double bar
$E_S=Z_{hN}$ therefore satisfies, compatibly with projection
and its core-state section,
\begin{equation}\label{C:P5-three-state-edge-homology}
 H_n(E_S;\mathbb Q)\cong H_n(BN;\mathbb Q)
       \oplus\operatorname{Tor}^{\mathbb QN}_{n-1}
                     (\mathbb Q_\chi,\mathbb Q)
 \qquad(n\geq1).
\end{equation}

If $\chi$ is identically one, the $N$-action is trivial and
$E_S=BN\times Z$.  Its map to $BP\simeq S^2$ factors through
$Z\simeq S^1$ and is null-homotopic.  Since a finite monoid
has zero first rational homology, Mayer--Vietoris gives
$\dim H_2(BC;\mathbb Q)=2$, excluding a projective plane.

Suppose $\chi$ is nontrivial.  Its coinvariants are zero, so
$H_1(E_S;\mathbb Q)=0$.  The section in
\eqref{C:P5-three-state-edge-homology} maps constantly to $BP$.
After cancelling its $BN$ summands, the relevant
Mayer--Vietoris segment is therefore
\[
 0\longrightarrow H_3(BC;\mathbb Q)
 \longrightarrow T_1\xrightarrow{\alpha}\mathbb Q
 \longrightarrow H_2(BC;\mathbb Q)\longrightarrow0,
 \qquad
 T_1=\operatorname{Tor}^{\mathbb QN}_1
                 (\mathbb Q_\chi,\mathbb Q).
\]
A projective plane would force $\alpha=0$ to retain its
one-dimensional second homology, and then force $T_1=0$
by its vanishing third homology.

Let $K$ be the minimal ideal of $N$, with row set $I$.
The character $\chi$ is zero on $K$: the ideal generated by
any zero-character element contains $K$.  Averaging the finite
minimal maximal group gives a projective left row module
$P_I=\mathbb QI$ and an exact sequence
\[
 0\longrightarrow V\longrightarrow P_I
   \longrightarrow\mathbb Q\longrightarrow0,
 \qquad V=\widetilde{\mathbb QI}.
\]
Since $P_I=KP_I$, tensoring with $\mathbb Q_\chi$ kills
$P_I$.  Projectivity and dimension shifting give
\[
 T_1\cong\mathbb Q_\chi\otimes_{\mathbb QN}V.
\]
Equation~\eqref{C:P5-three-state-affine} makes $\beta$ constant
on each minimal right ideal, so write its value on row $i$
as $b_i$.  Left multiplication by $n$ induces the row map
$\sigma_n$, and the same equation gives
\[
 b_{\sigma_n(i)}=\beta_n+\chi_n b_i.
\]
If the $b_i$ are not all equal, their differences define a
nonzero module homomorphism
$V\to\mathbb Q_\chi$.  Tensoring it gives a surjection
$T_1\to\mathbb Q$, contradicting $T_1=0$.

Consequently $b_i=c$ for a single $c\in\{0,1\}$, and the
last equation gives $\beta_n=c(1-\chi_n)$ for every $n$.
Every nonidentity action on $z$ therefore collapses it to the
same core state $u_c$.  The map $S\to U$ fixing $U$ and
sending $z$ to $u_c$ is an actual biset retraction.
The collage on $U$ is equivalent to $BP$: its edge is the
product of the contractible core action bar with $BN$, and
its projection to $BN$ is an equivalence.
The original collage thus retracts onto $S^2$.  If its second
rational cohomology had dimension one, the pulled-back sphere
class would span it and have square zero.  This again excludes
the projective plane, completing every case.
\end{proof}

\begin{corollary}\label{C:sixteen-P5-opposite-eight}
In the sixteen-arrow triangular two-object calculation, the
profile $(5,8,3)$ with rationally degree-two acyclic
eight-element corner is excluded structurally.  If that corner
has nonzero rational $H^2$, opposition places the case in
the complete $(8,5,3)$ catalog.
No exhaustive catalog of degree-two acyclic eight-element
monoids is required for this reduction.
\end{corollary}
\begin{proof}
The only five-element monoid with nonzero rational $H^2$
is the sphere monoid $P$, which is isomorphic to its opposite.
Apply the preceding proposition and then take the opposite
category when appropriate.
\end{proof}

\subsection{A sphere corner with one core orbit or at most four states}

\begin{proposition}[One core orbit for the five-element sphere corner]
\label{C:P5-one-core-bisets}
Let $P=\{1\}\sqcup(\{0,1\}\times\{0,1\})$ be the
five-element sphere monoid, with minimal band $K$, and let
$S$ be a finite nonempty $(P,N)$-biset.  Suppose that $KS$
is one two-point left $K$-orbit and that $P$ has no fixed
state in $S$.  If the collage $X$ has
\[
 \dim H^2(X;\mathbb Q)=1,\qquad H_3(X;\mathbb Q)=0,
\]
its degree-two generator has square zero.
In particular $X$ cannot have type $\mathbb{CP}^2$.
There is no bound on $|S|$ or $|N|$.
\end{proposition}
\begin{proof}
Label the core $U=\{u_0,u_1\}$ so that every band actor
$k_{ij}$ sends every state to $u_i$.  This follows from
idempotence exactly as in Proposition~\ref{C:P5-three-state-bisets}.
Commutation forces $N$ to fix both core points.  Every other
state is sent by a right actor either to another noncore
state or to one of these two points.

Put $A=\mathbb Q(S\setminus U)$, with the partial right
$N$-action that sends a basis vector to zero when its state
enters $U$.  The core of the left action category
$P\ltimes S$ is contractible.  For each noncore state there
are two cones on that core, one for each band column, with
their apex equal to that state.  Collapsing the core therefore
gives a graph with two parallel edges for each noncore state,
all sharing the core vertex.  This graph model is natural
for the right $N$-action: an actor carries a column-labeled
edge to the corresponding edge of the image state, or
collapses it to the core vertex.

For clarity, there is an explicit map from the quotient
of the actual nerve by its core to this graph.  On a
column-labeled cone send a point to the same cone-height
coordinate on its corresponding edge; equivalently this
is the barycentric weight of its apex.  The core is
contractible, so this map is a weak equivalence.  It is
$N$-equivariant: an actor preserves the height and column
label when its apex stays outside the core, and otherwise
collapses the whole cone into the core.  The graph chain complex is
$A[1]$ plus the contractible pair that sends a chosen first
edge to its noncore vertex minus the core vertex.  The
chosen edge and vertex pairs transform by the same partial
action, so this is a splitting of right $\mathbb QN$ chain
complexes.

Consequently the edge double bar has, compatibly with its
projection and core-state section to $BN$,
\[
 H_n(E_S;\mathbb Q)=H_n(BN;\mathbb Q)
       \oplus\operatorname{Tor}_{n-1}^{\mathbb QN}(A,\mathbb Q)
 \qquad(n\geq1).
\]
Its section maps constantly to $BP$.  Thus the inclusion
$BN\to X$ is null-homotopic, and cancellation of the $BN$
summands in Mayer--Vietoris gives
\[
 0\longrightarrow H_3(X;\mathbb Q)
 \longrightarrow T_1\longrightarrow\mathbb Q
 \longrightarrow H_2(X;\mathbb Q)
 \longrightarrow T_0\longrightarrow0,
 \qquad T_j=\operatorname{Tor}^{\mathbb QN}_j(A,\mathbb Q).
\]

If $H^2(X;\mathbb Q)\to H^2(BP;\mathbb Q)$ is zero,
the degree-two classes restrict to zero on both vertices.
They are relative classes, so their products vanish on the
suspension quotient by the vertex spaces.  In the remaining
case the dual map $H_2(BP;\mathbb Q)\to H_2(X;\mathbb Q)$
is an isomorphism.  The displayed sequence and the assumed
vanishing of $H_3(X;\mathbb Q)$ imply $T_0=T_1=0$.

There is an exact sequence of right $\mathbb QN$-modules
\[
 0\longrightarrow\mathbb Q(u_1-u_0)
 \longrightarrow\widetilde{\mathbb QS}
 \longrightarrow A\longrightarrow0.
\]
Its first term is the trivial module.  Rational duality
of a projective resolution of the right module $A$ gives
\[
 \operatorname{Ext}^1_{(\mathbb QN)^{\mathrm{op}}}(A,\mathbb Q)
   \cong D\operatorname{Tor}^{\mathbb QN}_1(A,\mathbb Q)=0.
\]
The sequence therefore splits.  A retraction onto its first
term defines an $N$-invariant function $q:S\to\mathbb Q$
with $q(u_i)=i$, by applying it to $s-u_0$.
Compose $q$ with the fixed function
$\rho:\mathbb Q\to\{0,1\}$ that is one at $1$ and zero
at every other rational number.  Sending $s$ to
$u_{\rho(q(s))}$ is an actual $N$-equivariant retraction
onto $U$.  It is also $P$-equivariant, since every band
actor is constant onto its prescribed $u_i$ on both sets.

The core collage is equivalent to $BP\simeq S^2$: its edge
bar is the product of a contractible core action category
with $BN$, and projection to $BN$ is an equivalence.
Thus $X$ retracts onto $S^2$.  The pullback of its sphere
class spans the one-dimensional $H^2(X;\mathbb Q)$ and
has square zero, as required.
\end{proof}

\begin{proposition}[Four states opposite the sphere corner]
\label{C:P5-four-state-bisets}
Let $S$ be a nonempty $(P,N)$-biset with $|S|\leq4$.
If $H^2(BN;\mathbb Q)=0$, its collage cannot have type
$\mathbb{CP}^2$.  When the left $P$-action has no fixed
state, the same conclusion holds without that cohomological
hypothesis on $N$.
\end{proposition}
\begin{proof}
The cases of at most three states were proved in
Proposition~\ref{C:P5-three-state-bisets}.  A $P$-fixed
state gives a section $BP\to E_S$; with the stated
vanishing at the other vertex, every degree-two collage
class is relative and its square is zero.

If there is no fixed state, every left $K$-orbit in $KS$
has two points.  One orbit is covered by
Proposition~\ref{C:P5-one-core-bisets}.  The only other
possibility at four states is $S=KS$ with two complete
orbits.  Write $S=O\times\{u_0,u_1\}$, $|O|=2$.
The right action preserves the row coordinate and acts
only on $O$.

The projection $P\ltimes S\to O$ has an $N$-equivariant
section choosing row zero.  The arrows $k_{i0}$ give an
$N$-equivariant natural transformation from that section
followed by projection to the identity.  The map of the
edge double bar to $BP$ is therefore coherently
$N$-equivariantly null-homotopic, while its map to $BN$
is identified with the action-bar projection for $O$.
The collage consequently has type
\[
 BP\ \vee\ B\mathcal C(N,O).
\]
Every positive rational product on the second summand is
zero by the two-state finite-line lemma, and the first
summand is a sphere.  The wedge has zero degree-two
products as well, completing the proof.
\end{proof}

\begin{corollary}\label{C:sixteen-P5-opposite-seven}
The sixteen-arrow profile $(5,7,4)$ needs no separate
enumeration.  If its seven-element vertex has zero
rational degree-two cohomology, apply
Proposition~\ref{C:P5-four-state-bisets}; otherwise its
opposite is in the complete $(7,5,4)$ catalog.
Any already computed $(5,7,4)$ records are independent
controls and need not be a premise of the exclusion.
\end{corollary}
\begin{proof}
Use the unique positive-degree-two five-element monoid
$P\cong P^{\mathrm{op}}$, and take the opposite category
in the second case.
\end{proof}

\subsection{Graph cofibers over the five-element band monoid}

\begin{lemma}[The literal two-step trivial-module resolution]
\label{C:P5-trivial-two-step}
Let $P=\{1\}\sqcup\{x_{ij}:i,j\in\{0,1\}\}$, where
$x_{ij}x_{kl}=x_{il}$, and let $A=\Q P$.
Put
\[
 e=x_{00},\qquad f=1-e,\qquad
 a=x_{10}-x_{00},\qquad b=x_{01}-x_{00}.
\]
The trivial left $A$-module has the exact projective resolution
\[
 0\longrightarrow Ae\xrightarrow{\ {-}\,b\ }Af
   \xrightarrow{\ {-}\,a\ }Ae
   \xrightarrow{\ \varepsilon\ }\Q\longrightarrow0.
\]
For every right $A$-module $W$ this identifies
\[
 \operatorname{Tor}_2^A(W,\Q)
   =\ker\bigl(We\xrightarrow{\ {-}\,b\ }Wf\bigr).
\]
\end{lemma}
\begin{proof}
The idempotents $e,f$ make both displayed left ideals projective.
The first has basis $x_{00},x_{10}$, and the second has basis
\[
 f,\qquad b=x_{01}-x_{00},\qquad c=x_{11}-x_{10}.
\]
Right multiplication by $b$ sends the two basis vectors of $Ae$
to $b,c$.  Right multiplication by $a$ sends $f$ to $a$ and
annihilates $b,c$.  The augmentation kernel in $Ae$ is exactly
$\Q a$.  These literal formulas prove exactness, including
injectivity at the left end.  Tensoring gives the asserted kernel.
\end{proof}

\begin{theorem}[A universal graph obstruction over $P_5$]
\label{C:P5-graph-cofiber-cups}
Let $P$ be the five-element monoid of the preceding lemma, and let
it act on the right of a nonempty finite graph $\Gamma$.
Each actor is required to send vertices to vertices and each
oriented edge at unit speed to an oriented edge or to collapse
it to a vertex.  Write
\[
 X=\operatorname{hocofib}(\Gamma_{hP}\longrightarrow BP).
\]
Then either the action has a common fixed vertex, or
$H_4(X;\Q)=0$.  In either case all products of two rational
degree-two classes on $X$ vanish.  No connectedness or
finite-CW-type hypothesis is imposed on $X$ or $\Gamma_{hP}$.
\end{theorem}
\begin{proof}
Let $r_{ij}$ denote the graph map for right multiplication by
$x_{ij}$.  Thus $r_{kl}r_{ij}=r_{il}$.  The image subgraphs
\[
 \Gamma_0=r_{00}\Gamma,\qquad \Gamma_1=r_{01}\Gamma
\]
are retracts, and the image of $r_{ij}$ depends only on $j$.
Indeed, for a fixed $j$ the two idempotent maps restrict to
the identity on each other's image, by the band identities.

If a vertex $v$ belongs to both image graphs, it is fixed by
$x_{00}$ and $x_{01}$.  For every $i,j$,
\[
 v x_{ij}=(v x_{00})x_{ij}=v x_{0j}=v.
\]
It is therefore fixed by all of $P$.  The fixed-vertex inclusion
gives a simplicial section $BP\to\Gamma_{hP}$ of the projection.
The section makes the projection injective on cohomology.
The cofiber exact sequence therefore shows that every
positive-degree cohomology class of the cofiber is a pullback
from the suspension of the action bar.  Its positive-degree
products vanish by naturality.  This is the fixed-state
relative-cup argument and does not require a further claim
about the homotopy type of the cofiber.

Assume, therefore, that $\Gamma_0$ and $\Gamma_1$ have disjoint
vertex sets.  They also have disjoint edge sets.  The maps
$r_{01}:\Gamma_0\to\Gamma_1$ and
$r_{00}:\Gamma_1\to\Gamma_0$ are inverse graph isomorphisms:
their two composites are the corresponding idempotent
retractions, hence are identities on these image graphs.

Set $W=H_1(\Gamma;\Q)$.  The retraction identifies
$We$ with the actual embedded subspace $H_1(\Gamma_0;\Q)$.
For $w$ in that subspace the map in the preceding resolution is
\[
 wb=w x_{01}-w x_{00}=r_{01*}w-w.
\]
A graph has no cellular two-boundaries.  Its first homology is
the subspace of cycles in its edge-chain vector space.
The two terms in this difference are supported on the disjoint
edge sets of $\Gamma_1$ and $\Gamma_0$.  Their difference can
therefore vanish only when $w=0$.  This proves
\[
 \operatorname{Tor}_2^A(H_1(\Gamma;\Q),\Q)=0.
\]
The injectivity uses actual graph realizability; an arbitrary
module of the same dimension need not have this property.

For completeness, this Tor group is the required fourth
homology even when the graph is disconnected.  Let
$L=\widetilde H_0(\Gamma;\Q)$ and let $K$ be the augmented
cellular chain complex
\[
 K=[C_1(\Gamma;\Q)\longrightarrow
       \ker(C_0(\Gamma;\Q)\to\Q)],
\]
in degrees one and zero.  The action-bar cofiber is represented
in positive homology by $(K\otimes_A^{\mathbb L}\Q)[1]$.
The triangle
$W[1]\to K\to L\to W[2]$ gives the exact segment
\[
 \operatorname{Tor}_4^A(L,\Q)\longrightarrow
 \operatorname{Tor}_2^A(W,\Q)\longrightarrow
 H_4(X;\Q)\longrightarrow
 \operatorname{Tor}_3^A(L,\Q).
\]
Both outer groups vanish by the length-two projective
resolution.  Hence $H_4(X;\Q)=0$ in the disjoint-image case,
as asserted.  Over a field this also makes $H^4(X;\Q)=0$.
\end{proof}

\begin{corollary}[The $P_5$--right-zero-three collage]
\label{C:P5-right-zero-three-all-states}
Let $R=\{1,p,q\}$ be the identity-adjoined two-element right-zero
band.  For every finite $(R,P)$-biset $S$, the two-object collage
has zero rational degree-two products.  In particular it cannot
have type $\CP^2$, at any number of states.  The opposite
collage is excluded as well.
\end{corollary}
\begin{proof}
If $S$ is empty, the classifying space is the disjoint union
$BP\sqcup BR$.  Its fourth rational cohomology is zero by
the displayed resolution and the contractibility of $BR$.
Assume henceforth that $S$ is nonempty.
The natural graph model of Lemma~\ref{C:right-zero-action-graph}
is equivariant for the commuting right $P$-action.
Since $BR$ is contractible, the collage has the homotopy type
of the graph cofiber in the theorem.  Its graph vertices are
the common image of the two retractions, so a common fixed
graph vertex is also a literal common fixed biset state.
The theorem applies without a bound on $|S|$.
\end{proof}

% Root and D full literal audits passed, including all derived shifts and module sides.
\subsection{Two-column graph actions and their contracted algebra}

\begin{lemma}[A two-term complex over a local algebra]
\label{C:local-two-term-derived-tensor}
Let $A$ be a finite-dimensional local rational algebra with
radical $J$ and residue field $A/J=\mathbb Q$.  Let $D$ be a
complex of finite-dimensional right $A$-modules concentrated in
homological degrees zero and one.  If
$H_i(D\otimes_A^{\mathbb L}\mathbb Q)$ vanishes for all
sufficiently large $i$, then it vanishes for every $i\geq2$.
\end{lemma}
\begin{proof}
Choose a bounded-below minimal resolution of $D$ by finitely
generated free right $A$-modules, in nonnegative degrees.
One construction starts with a bounded-below free resolution
and successively splits off each invertible matrix entry as a
contractible two-term free summand.  Over a local algebra the
remaining differentials have all entries in $J$; each degree
requires only finitely many cancellations.  This is the usual
minimal projective-complex construction, applied to the given
two-term complex.

Tensoring that minimal complex with the left residue module
$\mathbb Q$ makes all differentials zero.  The stated eventual
vanishing and Nakayama's lemma therefore make the free resolution
bounded.  Let $N$ be its largest nonzero degree.  If $N>1$, the
vanishing of $H_N(D)$ forces the top differential to be injective.
This is impossible.  A matrix of a right-module map has entries
acting on the left of the input coefficients.  Since $J$ is
nilpotent, there is a nonzero $a\in A$ with $Ja=0$; take a
nonzero element of the last nonzero radical power, or $a=1$ if
$J=0$.  The vector which is $a$ in one top free coordinate and
zero elsewhere is killed by every matrix entry of that minimal
differential.  It is a nonzero kernel vector.  Thus $N\leq1$,
which proves the asserted tensor-homology vanishing.
\end{proof}

\begin{theorem}[The two-column graph localization]
\label{C:two-column-graph-localization}
Let $M$ be a finite monoid with minimal rectangular band
$K=I\times\{0,1\}$, put $R=\mathbb QM$ and
$A=R/\mathbb QK$, and write
$V=\widetilde{\mathbb QI}$ for the left row augmentation module.
Let $M$ act on the right of a nonempty finite graph $\Gamma$,
mapping vertices to vertices and every edge at unit speed to an
edge, in either orientation, or to a vertex.  Set
\[
 X=\operatorname{hocofib}(\Gamma_{hM}\to BM).
\]
Either the action has a common fixed vertex, or there is an actual
two-term complex $D$ of finite-dimensional right $A$-modules such
that, naturally for every $n\geq4$,
\[
 H_n(X;\mathbb Q)
   \cong H_{n-2}(D\otimes_A^{\mathbb L}V).
\]
In particular:
\begin{enumerate}
\item If $V$ is projective over $A$, then all rational degree-two
products on $X$ vanish, without a bounded-homology hypothesis.
\item If $|I|=2$ and $A$ is local, then bounded rational homology
of $X$ implies that all its rational degree-two products vanish.
The same conclusion holds when the one-dimensional $V$ is
supported on a local direct-factor algebra of $A$.
\end{enumerate}
\end{theorem}
\begin{proof}
Choose $i_0\in I$ and put $e=x_{i_0,0}$ and $e_1=x_{i_0,1}$.
The image graphs $\Gamma_0=\Gamma e$ and $\Gamma_1=\Gamma e_1$
are retract subgraphs.  The image of any $x_{ij}$ depends only
on its column, since the same-column idempotents fix each other's
images.  Right multiplication by $e_1$ and $e$ gives mutually
inverse graph isomorphisms between $\Gamma_0$ and $\Gamma_1$.

If a vertex $u$ belongs to both image graphs, it is fixed by
both $e,e_1$ and hence by every band element:
$ux_{ij}=(ue)x_{ij}=ux_{i_0,j}=u$.
For arbitrary $m\in M$, the equality
$um=(ue)m=u(em)=u$ then makes it a common fixed vertex.
Its action-bar section makes the projection injective on
cohomology.  The cofiber exact sequence therefore makes every
positive-degree cohomology class of the cofiber a pullback
from the suspension of the action bar.  Products of such
classes vanish, as in Theorem~\ref{C:P5-graph-cofiber-cups}.

Otherwise $\Gamma_0$ and $\Gamma_1$ have disjoint vertex sets
and therefore disjoint edge sets.  Let
\[
 C=[C_1(\Gamma;\mathbb Q)\longrightarrow
       \widetilde C_0(\Gamma;\mathbb Q)]
\]
be its augmented graph complex of right $R$-modules.  Multiplication
gives a chain map
\[
 C e\otimes_{\mathbb Q}eR\longrightarrow C.
\]
This is degreewise injective.  In degree one its two column
summands have disjoint edge supports.  In degree zero,
$C_0e=\widetilde{\mathbb Q(\operatorname{Vert}\Gamma_0)}$;
the two summands have disjoint vertex supports and each has
augmentation zero.  These are precisely the image subspaces of
the two column copies.  The same argument allows reversed edge
orientations, since their only effect on chain coordinates is a
sign.

The image is exactly $C\mathbb QK$, by the literal band
factorization $\mathbb QK=Re\otimes_{\mathbb Q}eR$.
Consequently the quotient
\[
 D=C/(Ce\otimes_{\mathbb Q}eR)
\]
is literally a two-term right $A$-complex.  The subcomplex has
projective right $R$-modules in its two degrees.  After derived
tensoring with the trivial left $R$-module, it becomes the finite
complex $Ce$ in those same degrees, because $eR\otimes_R\mathbb Q
\cong\mathbb Q$.  Thus the short exact sequence of complexes
induces
\[
 H_j(C\otimes_R^{\mathbb L}\mathbb Q)
 \cong H_j(D\otimes_R^{\mathbb L}\mathbb Q),
 \qquad j\geq3.
\]

The row augmentation sequence is
\[
 0\longrightarrow V\longrightarrow Re\longrightarrow\mathbb Q
       \longrightarrow0.
\]
The ideal $\mathbb QK$ annihilates $V$, and $De=0$.
Since $Re$ is projective, tensoring this sequence gives
\[
 D\otimes_R^{\mathbb L}\mathbb Q
       \simeq (D\otimes_R^{\mathbb L}V)[1].
\]
The comparison to $A$ has no omitted derived correction.
Indeed $0\to\mathbb QK\to R\to A\to0$ is a left-projective
resolution of $A$, and a right $A$-module $U$ satisfies
$U\otimes_R\mathbb QK=Ue\otimes eR=0$.  Hence
$U\otimes_R^{\mathbb L}A\simeq U$ and
$U\otimes_R^{\mathbb L}V\simeq U\otimes_A^{\mathbb L}V$.
The same natural comparison applies to the finite complex $D$.
Finally the graph-cofiber chain model is
$(C\otimes_R^{\mathbb L}\mathbb Q)[1]$ in reduced homology.
Combining the two shifts proves the displayed formula exactly
in the asserted range $n\geq4$.

If $V$ is $A$-projective, its tensor with $D$ is concentrated
in degrees zero and one, so $H_4(X;\mathbb Q)=0$.
If $|I|=2$, the module $V$ is one-dimensional.  When $A$ is
local, its action on $V$ identifies the residue division algebra
with $\mathbb Q$, and $V$ is the residue module.  Bounded rational
homology of $X$ makes $D\otimes_A^{\mathbb L}V$ homologically
bounded.  Lemma~\ref{C:local-two-term-derived-tensor} gives
$H_2(D\otimes_A^{\mathbb L}V)=0$, so again $H_4(X;\mathbb Q)=0$.
For a local direct factor $B=Az$, with $z$ central and $zV=V$,
replace $D$ by $Dz$ and $A$ by $B$ in the same argument.
This proves all assertions, including the fixed-vertex alternative.
\end{proof}

\begin{corollary}[All sources of order at most six opposite $R_3$]
\label{C:six-source-right-zero-all-states}
Let $R_3$ be the identity-adjoined two-element right-zero band.
For every finite $(R_3,M)$-biset with $|M|\leq6$, bounded rational
homology of its collage implies that all rational degree-two
products vanish.  In particular no such collage is a projective
plane, at any number of cross states.
\end{corollary}
\begin{proof}
The left action category has the natural equivariant finite graph
model, and $BR_3$ is contractible.  Thus the collage is precisely
the graph cofiber in the theorem.
If $H^2(BM;\mathbb Q)=0$, every degree-two class is relative to
the two vertex spaces and the suspension quotient kills products.
Otherwise the minimal ideal must have at least two rows and two
columns.  At order at most six it is a $2\times2$ band with at
most two outside elements; a nontrivial local group would already
make the ideal have at least eight elements.

The contracted algebra is either $\mathbb Q$, or has basis
$1,a$ with $a^2$ equal to $0,a$, or $1$.  In the last two cases
it is semisimple, so the row module is projective.  In the first
it is the local dual-number algebra.  The preceding theorem
therefore applies in every case.  This argument uses bounded
rational homology only in the local-algebra alternative, and
does not infer finite CW type for an intermediate graph space.
\end{proof}

% Root and B full literal audits passed; D's finite counterexample gate passed.
\subsection{Natural graph comparisons for every opposite corner of order at most four}

\begin{lemma}[A natural graph for an identity-adjoined right-zero band]
\label{C:right-zero-band-natural-graph}
Let $N=\{1\}\sqcup K$, where $K$ is a nonempty finite right-zero
band, and let $S$ be a finite left $N$-set.  Put $U=KS$.
There is a natural finite graph $G_N(S)$ and an actual weak
equivalence
\[
 B(N\ltimes S)\longrightarrow G_N(S).
\]
Its vertex set is $S$.  For each $s\notin U$ and $k\in K$,
it has a labeled edge from $s$ to $ks$.  All maps of left
$N$-sets induce graph maps taking an edge at unit speed to an
edge with the same label or collapsing it to a vertex.
In particular the comparison is equivariant for every commuting
right monoid action.
\end{lemma}
\begin{proof}
Every $k\in K$ fixes $U$ pointwise.  Regard the displayed graph
as a category with the indicated edges as its only nonidentity
arrows; there are no composable nonidentity pairs.  Define the
functor from the action category by keeping its objects and
its arrows from $S\setminus U$ to $U$, and sending every
endomorphism at a core object to its identity.  The right-zero
multiplication makes this a functor.

For a core vertex $u$, the comma category $(u\downarrow q)$
has one object and endomorphism monoid $N$, whose classifying
space is contractible.  For a noncore vertex $s$, the comma
category $(s\downarrow q)$ has the initial object $(s,1_s)$.
Indeed the other objects are indexed by the individual labeled
edges from $s$, and the corresponding labeled action arrow is
the unique morphism from $(s,1_s)$ to each of them.  Quillen's
Theorem A proves the weak equivalence.  A map of acts preserves
the core, and sends a noncore vertex either to another noncore
vertex or to the core.  The stated edge rule therefore makes
both the graph construction and the functor strictly natural.
\end{proof}

\begin{lemma}[The higher homology of a two-band four-element action category]
\label{C:four-corner-action-tail}
Let $N=\{1,p,q,x\}$ have minimal ideal the right-zero band
$K=\{p,q\}$.  Let $S$ be a finite left $N$-set, $U=KS$,
$R=\mathbb QN$, and $A=R/\mathbb QK$.  Write
$W=\mathbb Q(q-p)$ as a right $A$-module and
$T=\mathbb QS/\mathbb QU$ as a left $A$-module.  Then
\[
 H_n(B(N\ltimes S);\mathbb Q)
   \cong \operatorname{Tor}_{n-1}^{A}(W,T),
 \qquad n\geq2.
\]
The comparison is natural for every commuting right action
on $S$.  In four of the five such four-element monoids, $A$
is semisimple, so these homology groups vanish.  The sole
exception, up to relabeling $p,q$, has
\[
 px=qx=p,\qquad xp=p,\qquad xq=q,\qquad x^2=p.
\]
For this exception $A=\mathbb Q[\epsilon]/(\epsilon^2)$
and $W$ is its residue module.
\end{lemma}
\begin{proof}
Put $e=p$.  Then $Re=\mathbb Qp$ and
$eR=\mathbb QK=\mathbb Qp\oplus\mathbb Qq$.
The multiplication isomorphism
$\mathbb QK=Re\otimes_{\mathbb Q}eR$ gives
\[
 \mathbb QK\otimes_R\mathbb QS
   =Re\otimes_{\mathbb Q}e\mathbb QS
   \cong\mathbb QU.
\]
Its map to $\mathbb QS$ is the literal inclusion of the
core-state span.  Since
$0\to\mathbb QK\to R\to A\to0$ is a projective resolution
of the right $R$-module $A$, this proves the natural derived
identification
\[
 A\otimes_R^{\mathbb L}\mathbb QS\simeq T.
\]
The right augmentation sequence
$0\to W\to eR\to\mathbb Q\to0$ has projective middle term.
Dimension shifting and derived associativity consequently give
\[
 H_n(B(N\ltimes S);\mathbb Q)
 =\operatorname{Tor}_{n}^{R}(\mathbb Q,\mathbb QS)
 =\operatorname{Tor}_{n-1}^{R}(W,\mathbb QS)
 =\operatorname{Tor}_{n-1}^{A}(W,T)
\]
in exactly the stated range.

For completeness the five possibilities are as follows.
If the right translation of $x$ on $K$ is the identity, then
$x^2=1$ or $x$.  If it swaps $p,q$, then $x^2=1$.
If it is constant, relabeling makes that constant $p$,
and $x^2=x$ or $p$.  These possibilities follow by applying
the right translation to $x^2$, and excluding a larger
minimal ideal.  The first four contracted algebras have
basis $1,\bar x$ with $\bar x^2=1$ or $\bar x$, hence
are semisimple over $\mathbb Q$.  The final possibility
has $\bar x^2=0$.  There the column-difference character
annihilates $\bar x$, as asserted.
\end{proof}

\begin{proposition}[The full natural small-corner comparison]
\label{C:small-opposite-natural-graph-comparison}
Let $|N|\leq4$, let $M$ be any finite monoid, and let $S$
be a nonempty finite $(N,M)$-biset.  There is a finite graph
$\Gamma$, with a right $M$-action by unit-speed edge maps
(allowing reversal) and edge collapses, and a natural
comparison from the collage to
\[
 Y=\operatorname{hocofib}(\Gamma_{hM}\longrightarrow BM),
\]
through actual maps and equivariant weak-equivalence zigzags,
which induces
\[
 H^2(Y;\mathbb Q)\xrightarrow{\;\cong\;}H^2(BC;\mathbb Q).
\]
Thus products on the collage are pullbacks of products on
the graph cofiber.  If $N$ is not the exceptional nilpotent
four-element monoid of Lemma~\ref{C:four-corner-action-tail},
this comparison is a rational homology equivalence in every
degree.  For that exception only the displayed low-degree
assertion is made.
\end{proposition}
\begin{proof}
If the trivial right $\mathbb QN$-module is projective,
the actual action-category projection
$N\ltimes S\to N\backslash S$ is a rational homology
equivalence: its chain homology is computed by the exact
coinvariant functor.  Also $BN$ is rationally acyclic.
Take $\Gamma$ to be the discrete right act $N\backslash S$.
The natural collage quotient, or its double-bar homotopy
pushout, gives the asserted full rational comparison.

The elementary small-monoid inventory leaves, in the
nonprojective case, $1+$ a right-zero two-band, the five
two-band four-element monoids, and $1+$ a right-zero
three-band.  The identity-adjoined bands have the natural
actual graph model of
Lemma~\ref{C:right-zero-band-natural-graph}.  Their classifying
spaces are contractible, so their collages are the indicated
graph cofibers up to homotopy.

For the five remaining monoids use the normalized
groupoid presentation of
Lemma~\ref{C:four-corner-normalized-action-groupoid}, for
arbitrary $S$, with no state bound.  Its objects are the
core $U$, and the generators are $g_s:Ps\to Qs$.
When the right translation of $x$ is the identity, quotient
the parallel edges by $g_{Es}=g_s$, collapsing an edge when
its image is a core vertex.  When the translation is constant,
collapse the loops $g_{Es}$; this is an invariant set.
Both constructions give finite oriented graphs with strict
right actions by edge maps or collapses.

In the swap case $E^2=1$.  Each two-element outside orbit
gives one unoriented edge.  An $E$-fixed outside point gives
a loop with relation $g_s^2=1$; kill these order-two loop
generators.  The result is the fundamental groupoid of a
finite graph.  Commuting right maps preserve the fixed
points and the pairs, so they give strict graph maps,
possibly reversing an edge.  Killing the order-two
generators preserves first rational homology.

In all three cases the actual normalized functor from
$N\ltimes S$ to the resulting graph groupoid is natural
and induces a component bijection and an isomorphism on
first rational homology.  The natural one-truncation
equivalence of a graph with the classifying space of its
fundamental groupoid is an equivalence of $M$-diagrams.
One may realize this equivalence using singular simplicial
sets and their natural one-truncations; it does not require
choosing equivariant spanning trees or edge representatives.

The target graph has no second homology.  Consequently the
cone of the augmented chain comparison has zero homology
in degrees zero, one and two.  Tensoring with the free left
monoid bar preserves that connectivity.  Suspending the
augmented bar comparison gives an isomorphism on second
homology of the actual cofibers, and hence the displayed
second-cohomology isomorphism.  The cofiber maps are actual
maps of spaces, so cup products pull back.

In the four semisimple cases,
Lemma~\ref{C:four-corner-action-tail} also shows that the
source action category has no rational homology above one.
Its graph comparison is therefore a rational homology
equivalence in all degrees.  The free bar and the homotopy
pushout preserve that equivalence.  No higher-degree
equivalence is asserted in the remaining nilpotent case.
\end{proof}

\begin{remark}[A necessary restriction on the graph comparison]
\label{C:small-opposite-graph-higher-homology-warning}
The exceptional qualification is essential.  In the
nilpotent four-element monoid take $S=\{u,t\}$, with every
nonidentity actor sending both states to $u$.  Then
$T=\mathbb Q\{t\}$ and $W$ are both the residue module of
$A=\mathbb Q[\epsilon]/(\epsilon^2)$.  The periodic free
resolution with every differential multiplication by
$\epsilon$ gives
\[
 H_n(B(N\ltimes S);\mathbb Q)=\mathbb Q
 \quad\hbox{for every }n\geq2.
\]
The normalized graph is one circle; the same augmentation
sequence, or its fundamental-groupoid presentation, gives
$H_1=\mathbb Q$ as well.  Thus this finite action category
has nonzero rational homology in every positive degree.
Its natural graph comparison still induces the stated
second-cohomology isomorphism on collages, but cannot
transfer bounded higher homology from a collage to its
graph replacement.  No spectral-sequence degeneration
or arbitrary finite-type inference is being used.

The standalone program
\path{code/d_verify_nilpotent_graph_counterexample.py}
checks all $64$ multiplication triples and $32$ action
equations, the literal projective right ideals with bases
$(p,q)$ and $(1-p,x-p)$, and every kernel and image in the
periodic resolution.  It also rebuilds the eight-arrow
action category and verifies the actual functor to the
one-circle groupoid, together with a first-homology
cycle--cocycle pairing equal to one.  Its complete record,
\path{code/outputs/d_nilpotent_graph_counterexample_verified.json},
certifies all the stated finite data and the all-degree
periodic conclusion without producer imports.
\end{remark}

\begin{corollary}[Unconditional small-corner graph cup consequences]
\label{C:small-opposite-P5-all-states}
For every $|N|\leq4$ and every finite $(N,P_5)$-biset,
all rational degree-two products on the collage vanish,
with no state bound or bounded-homology assumption.
The same conclusion holds for a source $M$ whose minimal
ideal is a two-column rectangular band and whose row
augmentation is projective over $\mathbb QM/\mathbb QK$.
In particular the seventeen-arrow profile $(5,4,8)$ is
excluded in full.
\end{corollary}
\begin{proof}
Apply the second-cohomology isomorphism in the proposition.
The unconditional projective-row clause of
Theorem~\ref{C:two-column-graph-localization} kills all
degree-two products on the graph cofiber, including graphs
with reversed edges.  For $P_5$ its contracted algebra is
$\mathbb Q$, so its row augmentation is projective.
The empty-biset case is a disjoint union of the two vertex
spaces, whose fourth rational cohomology is zero.
\end{proof}

\begin{corollary}[The bounded small-source consequence away from the exception]
\label{C:small-opposite-six-source-regular}
Suppose $|M|\leq6$, $|N|\leq4$, and $N$ is not the
exceptional nilpotent four-element corner.  If the collage
of a finite $(N,M)$-biset has bounded rational homology,
all its rational degree-two products vanish, at every
state size.
\end{corollary}
\begin{proof}
The natural graph comparison is now a full rational
homology equivalence, so its graph cofiber has bounded
rational homology.  If $H^2(BM;\mathbb Q)=0$, the vertex
restrictions vanish and the suspension quotient kills
products.  Otherwise the minimal ideal of $M$ is a
$2\times2$ band and the contracted algebra has dimension
at most two.  The projective or local alternatives in
Theorem~\ref{C:two-column-graph-localization} apply exactly
as in Corollary~\ref{C:six-source-right-zero-all-states}.
The exceptional $N$ is deliberately excluded here.
\end{proof}

\begin{corollary}[The corresponding star comparisons]
\label{C:small-opposite-graph-stars}
Let a finite star category have central endomorphism
monoid $M$, sink endomorphism monoids $N_i$ of order at
most four, and nonempty cross bisets $S_i$.  If $M=P_5$,
or if $M$ has a two-column minimal band with projective
contracted row augmentation, all rational degree-two
products on its classifying space vanish.  If $|M|\leq6$,
the same conclusion holds under bounded rational
homology provided none of the $N_i$ is the exceptional
nilpotent corner.  All assertions have opposite versions.
\end{corollary}
\begin{proof}
Apply the natural comparison separately to each edge
bar in the star's homotopy-pushout description.  The
comparisons fix the central $BM$, and each $BN_i$ is
rationally acyclic.  They therefore give the same
second-cohomology isomorphism, and in the nonexceptional
case a full rational homology equivalence, with the
star of graph cofibers.

Identify its finitely many cone vertices.  For a connected
space, identifying $r$ distinct points adds a wedge of
$r-1$ circles up to homotopy.  The resulting space is
the single cofiber associated to the disjoint union
of the coefficient graphs.  This operation preserves
the cohomology ring in degrees at least two and preserves
bounded rational homology.  The graph-localization
theorem applies to that disjoint union and proves
the asserted conclusions.  The only boundedness transfer
used is the full rational comparison in the explicitly
nonexceptional case.
\end{proof}

% Root, B and D full literal audits PASS; independent96 signed rational gate PASS.
\subsection{Rational six-state graph obstructions and four-element two-band corners}

\begin{lemma}[The connected six-state obstruction over the rationals]
\label{C:connected-six-rational-square}
Let a finite monoid $M$ act on a connected graph with
$|V|+|E|\leq6$ by orientation-preserving unit-speed
edge maps and edge collapses.  If the graph cofiber
$X$ has one-dimensional $H^2(X;\mathbb Q)$, its
degree-two square is zero.  There is no bound on
$|M|$ and no bounded-homology hypothesis.
\end{lemma}
\begin{proof}
The complete connected classification in
Theorem~\ref{C:right-zero-six-connected} is rational
under this rank-one assumption, except for its last
two literal image controls.  Its connected five-state
dependency has the exact rational square primitive
of Lemma~\ref{C:nine-five-state-graph-square}; the
connected part of that proof places no bound on $M$.
All other reductions in the six-state proof use
equivariant graph maps, invariant functionals, and
the stated scalar or polynomial self-maps.

The two remaining action-image monoids are the
literal five- and six-element tables in
\path{code/outputs/a_rightzero_six_opposite_edges.json}.
The first is $P_5$.  The second has a $2\times2$
minimal band and two outside elements $1,u$ with
$u^2=1$; its contracted algebra is
$\mathbb Q[C_2]$.  Thus the row augmentation is
projective in both cases.  These identifications
are literal in the stored tables: in the six-element
table the minimal band consists of labels $2,3,4,5$,
in row-major order, and the units are $0,1$.
All of these multiplication tables and their
graph actions were independently checked in
\path{code/outputs/d_rightzero_six_H3_independently_verified.json}.
The unconditional clause of
Theorem~\ref{C:two-column-graph-localization}
therefore makes the fourth rational homology of
each image graph cofiber zero.

For a connected graph, its cofiber second homology
is $H_1(\Gamma;\mathbb Q)\otimes_{\mathbb QM}\mathbb Q$.
Passing to the actual action image leaves this
coinvariant space unchanged, because the module
action factors through that image.  The actual
cofiber comparison is an isomorphism on second
rational homology.  The zero image square therefore
pulls back to the original square, including for
arbitrary source inflations.  No mod-two homology
test is transferred in this argument.
\end{proof}

\begin{lemma}[The complete rational image controls through order seven]
\label{C:six-state-rational-image-controls}
For the two disconnected graph families remaining
in Theorem~\ref{C:right-zero-six-small-sources},
restrict the independently complete action-image
lists to order at most seven.  The equal-loop family
has $53$ images and the isolated-vertex family has
$43$.  Their exact rational classifications are
\[
\begin{array}{c|r}
H_1\ne0&5\\
\dim H_2\geq2&4\\
H_1=0,\ H_2=\mathbb Q,\ \text{square zero}&68\\
H_1=H_2=0&19.
\end{array}
\]
All nineteen zero-second-homology images have order
seven and belong to the isolated-vertex family.
\end{lemma}
\begin{proof}
The underlying finite sets of action images are
the ones already independently regenerated in
\path{code/outputs/d_rightzero_six_equal_loop_independently_verified.json}
and
\path{code/outputs/b_rightzero_six_isolated_independently_verified.json}.
The selection retains every order-at-most-seven
image in the equal-loop $(1,1)$ hull and every
order-at-most-seven isolated-vertex image.
No group-order or binary-homology filter is used
to discard a rational case.

The producer
\path{code/a_rightzero_six_rational_images.py}
rebuilds the literal categories and all signed
normalized cochain matrices.  Each positive
first-homology record has an exact rational
cycle--cocycle pairing equal to one.
Each rank-at-least-two second-homology record has
two exact cocycles and dual cycles.
For the zero-dimensional second homology records,
modulo-$65521$ ranks give an upper bound zero
on the rational second Betti number.
For each of the other $68$ records the same
rank bounds give upper bound one, and its exact
cycle--cocycle pairing gives the matching lower
bound.  A rational three-cochain is then checked
against every signed four-simplex equation for
the square of that actual nonboundary generator.
The modular solve proposes coefficients only:
every final identity is checked over $\mathbb Q$.

The complete source-linked rational certificates
are in
\path{code/outputs/a_rightzero_six_rational_images.json}.
The standalone checker
\path{code/b_verify_rightzero_six_rational_images.py}
imports no producer or matrix-library routines.
It rebuilds every category from the literal actor inventory,
checks the exact $53+43$ selection and source hashes,
and independently verifies the signed faces, all boundary
squares, the finite-field rank upper bounds, and every
exact rational cycle, dual pairing and primitive.
Its complete record
\path{code/outputs/b_rightzero_six_rational_images_independently_verified.json}
checks $53072$ typed associativity equations,
$204344$ signed boundary-square equations,
$21380$ cocycle equations, $81$ rational cycles,
$89$ exact dual pairings and $101748$ signed
square-primitive equations.  It independently recovers
all four displayed counts and verifies that all nineteen
zero-second-homology images have order seven.
These certificates concern the displayed image
categories only; the source comparison is supplied
in the following theorem.
\end{proof}

\begin{theorem}[The rational six-state graph obstruction for sources through seven]
\label{C:seven-source-six-state-rational-graph}
Let $|M|\leq7$ and let $M$ act on a graph with
$|V|+|E|\leq6$ by orientation-preserving unit-speed
edge maps and edge collapses.  If its graph cofiber
$X$ has $H_1(X;\mathbb Q)=0$ and
$H^2(X;\mathbb Q)=\mathbb Q$, the square of its
degree-two class is zero.  No finite-CW-type,
third-homology or bounded-homology assumption is needed.
\end{theorem}
\begin{proof}
The connected case is
Lemma~\ref{C:connected-six-rational-square}.
In the disconnected case follow the complete
graph-shape analysis of
Theorem~\ref{C:right-zero-six-small-sources}.
Its noncontrol arguments use only the present
rational hypotheses: the component-act pullback
has zero products by
Lemma~\ref{C:nine-unbounded-act-cups}; the remaining
edge functional is nonzero; scalar maps and the
two-eigenvalue argument kill its square; and each
equivariant graph quotient has a cycle kernel
with zero coinvariants, giving the actual
second-cohomology isomorphism.  In particular the
proof does not identify the entire disconnected
second homology with the cycle coinvariants.

Only the equal-loop and isolated-vertex image
families remain.  The actual graph-action
quotient is an isomorphism on first homology
and a surjection on second homology over
$\mathbb Q$, by
Lemma~\ref{C:graph-action-image-low-homology}.
Therefore the five positive-first-homology and
four rank-at-least-two records are impossible.
In each of the $68$ rank-one records the
surjection is an isomorphism, so its exact
rational zero square pulls back to the square
on $X$.

Finally every zero-second-homology image has
seven elements.  A surjective action map from
$M$, with $|M|\leq7$, is therefore an isomorphism
in these cases.  Its graph cofiber is the
literal image cofiber and has zero second
rational homology, again impossible.
Thus none of these nineteen records needs a
nonfaithful-source inflation argument.
This finishes every disconnected graph case.
\end{proof}

\begin{lemma}[The small swap graph at six states]
\label{C:four-corner-swap-six-state-graph}
For an order-four two-band opposite corner whose
outside actor swaps the two minimal columns,
a collage with at most six cross states and
source order at most seven cannot have type
$\mathbb CP^2$.
\end{lemma}
\begin{proof}
The actual normalized graph comparison of
Proposition~\ref{C:small-opposite-natural-graph-comparison}
has core vertices $U$ and at most
$\lfloor(6-|U|)/2\rfloor$ edges.
A singleton core is a common fixed state.
For $|U|\geq2$ there are at most two edges;
right actors may reverse them.
The comparison preserves both first and second
rational homology: its augmented-chain cone
has zero homology in degrees zero, one and two,
and the cofiber construction adds one suspension.

For a forest, the component projection is an
equivariant underlying equivalence and reduces
to an ordinary act cone, whose degree-two
products vanish by
Lemma~\ref{C:nine-unbounded-act-cups}.
If every edge is a loop,
simultaneous loop doubling, fixing every vertex,
commutes with all actors, including reversals.
The graph-component two-eigenvalue argument
makes the eigenvalue-one classes pull back
from the ordinary component cone, and makes
an eigenvalue-two square vanish.

Any remaining cyclic graph has two vertices
and two edges and is connected and unicyclic:
it is a loop with a bridge or a two-edge circle.
Indeed two quotient edges force $|U|=2$ by
the stated bound, so there are no additional
isolated vertices.  A nonzero rational cycle
coinvariant forces every actor to have degree
$+1$ on its unique circuit.  The unit-speed
length argument makes that circuit strictly
invariant and each restricted actor an
orientation-preserving circle automorphism.
Its inclusion is an equivariant underlying
equivalence.  For a circuit of $L=1$ or $2$
edges, the equal-edge degree-$(L+1)$ map
commutes with these automorphisms and kills
the square.
These alternatives exhaust the small graph.
Pullback through the actual comparison proves
the assertion.
\end{proof}

\begin{theorem}[The five two-band order-four corners at the remaining profile]
\label{C:seven-four-six-two-band-excluded}
Let $|M|\leq7$, let $|S|\leq6$, and let the
four-element opposite corner have minimal ideal
a two-element right-zero band.  Its collage
cannot have type $\mathbb CP^2$.
\end{theorem}
\begin{proof}
For the identity or constant right translation
of the outside actor, the normalized graph is
oriented and has at most $|S|$ vertex-plus-edge
states.  The actual small-corner comparison
preserves first and second rational homology
and pulls back cup products.  A projective
plane would therefore give the hypotheses
of Theorem~\ref{C:seven-source-six-state-rational-graph},
contradicting its zero square.
This uses only the low-degree comparison even
for the exceptional nilpotent corner, and does
not transfer bounded higher homology there.
The swap case is the preceding lemma.
\end{proof}

% Root, B and D full literal audits PASS; independent149 signed/hash gate PASS.
\subsection{The rational six-state graph obstruction through eight source elements}

\begin{lemma}[All image controls through order eight]
\label{C:six-state-rational-images-eight}
The complete equal-loop and isolated-vertex action-image inventories
through order eight contain $56+93=149$ literal monoids.  Their
rational graph-cofiber classifications are
\[
\begin{array}{c|r}
H_1\ne0&5\\
\dim H_2\geq2&4\\
H_1=0,\ H_2=\mathbb Q,\ \text{square zero}&77\\
H_1=H_2=0&63.
\end{array}
\]
The last row consists precisely of nineteen order-seven images
and forty-four order-eight images, all in the isolated-vertex family.
The nineteen order-seven records are exactly the ones in the
previous complete $660$-lift calculation.
\end{lemma}
\begin{proof}
The complete finite inventories are the independently verified ones
in Proposition~\ref{C:right-zero-six-disconnected-controls}.
The new producer
\path{code/a_rightzero_six_rational_images_eight.py}
retains every order-at-most-eight image in the equal-loop $(1,1)$
hull and every order-at-most-eight isolated-vertex image.
It writes a separate atomic file,
\path{code/outputs/a_rightzero_six_rational_images_eight.json},
preserving the earlier order-seven output.

Every positive-rank assertion has exact rational cocycles and
dual cycles.  Rational upper bounds are independently testable
from ranks modulo $65521$.  Each rank-one image has a signed
rational three-cochain whose coboundary is the square of an
actual nonboundary generator, checked on every normalized
four-simplex.  Thus the added fifty-three records contribute
nine exact rational square primitives and forty-four additional
zero-second-homology records; there is no nonzero-square survivor.

The standalone checker
\path{code/b_verify_rightzero_six_rational_images_eight.py}
imports no producer or matrix-library routines.  It independently
reconstructs all $149$ categories, verifies the exact $56+93$
selection, all rational pairings and primitives, and the identity
of the nineteen order-seven records with the accepted lift scope.
Its complete record
\path{code/outputs/b_rightzero_six_rational_images_eight_independently_verified.json}
checks $112485$ typed associativity equations, $529711$ signed
boundary-square equations, $28157$ cocycle equations, $90$
rational cycles, $98$ exact dual pairings and $149259$ signed
square-primitive equations.  It checks the current SHA256
identities of all three unchanged lift inputs and the old
independent $660$-lift gate itself.  That previous gate's
arithmetic is reused with exact linkage, not described as a
new second enumeration.  No producer completion flag alone
is used as an independent premise.
\end{proof}

\begin{theorem}[Rational six-state graph cups through source order eight]
\label{C:eight-source-six-state-rational-graph}
Let $|M|\leq8$, and let $M$ act on a finite oriented graph with
$|V|+|E|\leq6$ by orientation-preserving unit-speed edge maps
and edge collapses.  Suppose its graph cofiber $X$ satisfies
\[
 H_1(X;\mathbb Q)=0,\qquad H^2(X;\mathbb Q)=\mathbb Q.
\]
Then its degree-two square vanishes.  No finite-CW-type or bounded
higher-homology hypothesis is needed.
\end{theorem}
\begin{proof}
The connected case is the all-source-order
Lemma~\ref{C:connected-six-rational-square}.  The noncontrol
disconnected reductions in
Theorem~\ref{C:seven-source-six-state-rational-graph} use the
source order only through the unconditional ordinary-act cup
theorem for monoids of order at most nine.  They therefore still
apply here, with the actual component contribution retained.

For either remaining image family the action-image map is an
isomorphism on $H_1$ and a surjection on $H_2$ over $\mathbb Q$.
The first two rows of the preceding lemma are impossible.
In each of the $77$ rank-one rows the actual second-homology
map is an isomorphism, so the rational zero square pulls back.
An order-eight image of a source with at most eight elements is
the source itself.  Its zero second homology excludes all
forty-four such records.  The same reasoning excludes an
order-seven source mapping onto an order-seven zero-$H_2$ image.

Only an eight-element source mapping onto one of the nineteen
order-seven zero-$H_2$ images remains.  If
$H^2(BM;\mathbb Q)=0$, all degree-two classes of the graph cofiber
are relative to $BM$ and pull back from its suspension quotient,
so their products are already zero.  Otherwise $M$ is covered by
the complete $516$-table source list of
Theorem~\ref{C:eight-act-cones}.  A surjection from eight to seven
elements has exactly one two-element fiber and all other fibers
singletons.  Its unique pair is a monoid congruence.  The accepted
independent enumeration tries every such pair, every quotient
isomorphism to the nineteen target monoids, and therefore every
possible right action in this residual scope.

Specifically,
\path{code/b_verify_rightzero_six_isolated.py}
checks all $516$ source tables, all $989$ one-pair congruences,
and all $19\cdot720$ identity-preserving target relabelings.
Its $660$ distinct lifted six-state acts have literal graph
cofibers with $H_2(-;\mathbb F_2)=0$.  The complete independent
record is
\path{code/outputs/b_rightzero_six_isolated_independently_verified.json}.
The exact source-linked files are
\path{code/outputs/c_eight_H2_targets.tsv},
\path{code/outputs/a_rightzero_six_isolated_images.json}, and
\path{code/outputs/a_rightzero_six_isolated_inflations.json}.
Their recorded SHA256 identities, and their equality to the
current input identities, are part of the new linkage gate.
The source-list completeness is the proved own-census statement,
not an unstated external-census assumption.

Each literal finite category has finitely generated integral
homology in each fixed degree.  The universal-coefficient
injection
\[
 H_2(-;\mathbb Z)\otimes\mathbb F_2
       \longrightarrow H_2(-;\mathbb F_2)
\]
therefore shows that $H_2(-;\mathbb F_2)=0$ forces the free
rank of $H_2(-;\mathbb Z)$ to be zero, and hence forces
$H_2(-;\mathbb Q)=0$.  This is the only change of coefficient
field used for these $660$ lifts.  It contradicts the assumed
rational rank one.  No mod-two square primitive is substituted
for a rational square primitive, and no zero second homology
is pulled back from a proper image without the complete lift
calculation.
\end{proof}

\begin{corollary}[Five four-element corners at the first new six-state profile]
\label{C:eight-four-six-two-band-excluded}
Let $|M|\leq8$, $|S|\leq6$, and suppose the four-element
opposite corner has minimal ideal a two-element right-zero band.
Its collage cannot have type $\mathbb CP^2$.
In particular all five such corners are excluded in the
eighteen-arrow profile $(8,4,6)$.
\end{corollary}
\begin{proof}
For identity or constant right translation of the extra opposite
actor, the actual normalized graph comparison is oriented and
has at most six vertex-plus-edge states.  It induces the required
rational $H_1,H_2$ isomorphisms and preserves cup products.
The preceding theorem applies, including for the nilpotent
four-element corner.  No comparison of higher homology is used.

For the swap corner, the proof of
Lemma~\ref{C:four-corner-swap-six-state-graph} applies at source
order eight as well: its only source-order input is the ordinary
component-act cup theorem through order nine.  The same
at-most-two-edge graph alternatives, strict bouquet maps and
unique-circuit maps exclude the square.  This completes the
five two-band corners.  The identity-adjoined right-zero
three-band is a separate sixth corner and is not claimed
excluded at source order eight by this corollary.
\end{proof}

% Root and agents B and D have passed the full literal proof; independent finite gates cited below are complete.
\subsection{The exceptional six-by-four opposite-corner profile}

\begin{lemma}[The remaining six-element source is unique]
\label{C:six-local-source-normal-form}
Suppose $|M|=6$, $H^2(BM;\mathbb Q)\ne0$, and its row
augmentation is not projective over the contracted
minimal-band algebra.  Then, up to isomorphism, $M$ is
\[
 M=\{1\}\sqcup\{x_{ij}:i,j=0,1\}\sqcup\{h\},
\]
with the usual rectangular-band multiplication and
\[
 hx_{ij}=x_{0j},\qquad x_{ij}h=x_{i0},\qquad h^2=x_{00}.
\]
In particular the contracted algebra is the dual-number
algebra and the row module is its residue module.
\end{lemma}
\begin{proof}
The minimal ideal must have at least two rows and columns
and is therefore the $2\times2$ band.  The contracted
algebra has basis $1,\bar h$ with $\bar h^2$ equal to
$0,\bar h$, or $1$.  In the latter two cases it is
semisimple and the row module is projective.
In the first case the row character of $h$ is zero,
so its row map is constant; the nonzero balanced
degree-two form forces the same zero column character,
so its column map is constant as well.
Relabel the row and column constants to be zero.
All mixed products are then the displayed ones.
The profile of $h^2$ forces its band value to be
$x_{00}$, proving uniqueness.
\end{proof}

\begin{lemma}[The exceptional corner has only one remaining core pattern]
\label{C:six-four-seven-core-pattern}
Let $M$ be the preceding source and let
$N=\{1,p,q,x\}$ be the exceptional nilpotent four-element
opposite corner.  If a seven-state $(N,M)$-collage has
type $\mathbb CP^2$, its minimal-band state core consists
of exactly two full two-state orbits.  On these orbits
the three nonidentity left maps $P,Q,E$ all induce the
same graph: one vertex orbit and one loop-edge orbit,
with each sending the loop orbit to the vertex orbit.
\end{lemma}
\begin{proof}
Use the actual second-cohomology comparison of
Proposition~\ref{C:small-opposite-natural-graph-comparison}.
For the nilpotent corner its normalized graph is obtained
from the $P,Q$ graph by deleting precisely those loop
edges indexed by noncore states in $ES$.
The comparison is only a second-cohomology isomorphism;
no boundedness of its graph cofiber is inferred.

The projective-plane generator pulls back from a
rank-one nonzero square on this graph cofiber.
If its invariant edge functional were zero, the class
would pull back from an ordinary component-act cone.
Lemma~\ref{C:nine-unbounded-act-cups} would kill its square.
Thus that edge functional is nonzero.  For a band
idempotent $e$, it is determined by its restriction
to $We=H_1(\Gamma e;\mathbb Q)$, which is nonzero.

A singleton band-state orbit gives a band-fixed vertex
after applying a left retraction if necessary.  This
vertex is fixed by all $M$ and all $N$, so the fixed-state
argument gives zero products.  Consequently every band
orbit has two states and there are at most three.
The image graph $\Gamma e$ has at most one state for
each such orbit.  One orbit cannot support a cycle.

With three orbits there are six core states and one
possible transient.  Before the loop deletions, the
three-state core graph must have positive first homology,
so all its edges are loops and $P=Q$ on the band core.
If $P=Q$ everywhere, the normalized graph is a disjoint
union of bouquets.  Strict loop doubling fixes its
vertices and commutes with every actor.  The graph-component
argument of Lemma~\ref{C:graph-component-eigenvalues}
gives eigenvalues one and two on its positive cohomology;
the eigenvalue-one second classes pull back from the
ordinary component cone.  They have zero squares by
Lemma~\ref{C:nine-unbounded-act-cups}, while an
eigenvalue-two class has square of eigenvalue four
and therefore zero.  Its rank-one square is impossible.
Otherwise the sole transient $z$ has two distinct
core endpoints in the same band orbit, hence in
opposite columns.  It is a nonloop edge.

Every state in $ES$ satisfies $P=Q$, because
$PE=QE=P$.  Thus this nonloop edge is not deleted by
the normalized graph quotient.  In the graph-localization
complex of Theorem~\ref{C:two-column-graph-localization}
it is the sole surviving noncore edge, and its boundary
is the nonzero difference of the two column classes.
Again the complex is
$[\mathbb Q\xrightarrow{\pm1}\mathbb Q]$.
The graph cofiber has zero fourth rational homology
without any boundedness assumption, a contradiction.

There must therefore be two band orbits.  Positive
$H_1(\Gamma e)$ forces one to be a vertex orbit and the
other a loop-edge orbit.  On the orbit set $P,Q$ send
the second to the first.  The equality $E^2=P$ and
the fact that $E$ fixes the common left core force
$E$ to do the same: it cannot fix the loop orbit.
All commuting maps between full two-column orbits
preserve the column labels.  This proves the stated
literal normal form.
\end{proof}

\begin{lemma}[The complete exceptional seven-state biset control]
\label{C:six-four-seven-nilpotent-control}
In the normal form of the preceding lemma there are
$244$ labeled-profile right actions and $4340$ raw
left actions.  Exact state conjugacies leave $2237$
literal categories; every one either fails the required
second mod-two Betti number or has a zero square.
\end{lemma}
\begin{proof}
Label the two core orbits by $(0,1)$ and $(2,3)$.
The first is the vertex orbit, and the second is the
loop orbit.  Each of the three transient states has
a row-orbit profile in $\{0,1\}^2$; sort their multiset.
The extra right actor $h$ restricts to $(0,0,2,2)$.
For a transient $s$, its image must have constant profile
$(\gamma_s(0),\gamma_s(0))$.  Every such image is tried,
then $h^2=x_{00}$ and every source product are checked.

The left map $E$ restricts to $(0,1,0,1)$ on the
four core states.  Its transient values must be
$0,1$ or transients of profile $(0,0)$.
Try every such assignment commuting with the complete
right action and retain $E^3=E^2$.  Set $P=E^2$.
Then try every commuting idempotent $Q$ with the same
image as $P$, imposing $QE=P$.  These conditions are
exactly the left-corner equations:
\[
\begin{gathered}
 P^2=P,\quad Q^2=Q,\quad PQ=Q,\quad QP=P,\\
 E^2=P,\quad EP=P,\quad EQ=Q,\quad PE=QE=P.
\end{gathered}
\]
For example $E$ fixes the image of $P$ because
$EP=E^3=P$, so $EQ=Q$ follows for a common-image
idempotent $Q$.  All sixteen left products and all
mixed equations are also rechecked literally.
The case $P=Q$ was already excluded.

Only permutations of equal-profile transient labels
are used to identify these actions.  In particular,
$p,q$ are not interchanged: the fixed nilpotent corner
distinguishes $p=x^2$.  The producer
\path{code/a_six_nilpotent_seven_acts.py}
stores every retained right action and each $P,Q,E$,
as well as the actual mod-two category certificate.
Its complete output is
\path{code/outputs/a_six_nilpotent_seven_acts.json}.
Exactly $1198$ categories have the wrong second
Betti number; all $1039$ others have explicit
nonboundary generators and square primitives.
The standalone checker
\path{code/b_verify_six_nilpotent_seven_acts.py}
uses fresh exhaustive $7^3$ pointwise assignments for
the extra right actor and the transient left maps.
It reproduces all $244$ right actions and $4340$
raw left actions, and expands the $2237$ retained
orbits disjointly, without interchanging $p,q$.
Its independent literal bars check $1816444$ typed
associativity, $1308645$ boundary-square, $514305$
cocycle and $2711790$ square-primitive equations.
The complete output is
\path{code/outputs/b_six_nilpotent_seven_acts_independently_verified.json}.
\end{proof}

\begin{theorem}[The full seventeen-arrow six-by-four profile]
\label{C:six-four-seven-excluded}
No two-object triangular category with corner orders
six and four and seven cross arrows has type
$\mathbb CP^2$.
\end{theorem}
\begin{proof}
Use opposition to write the cross arrows from the
order-six corner to the order-four corner; the other
orientation is covered by applying the same statement
to the opposite category.
If the source has zero second rational cohomology,
both vertex restrictions vanish and the relative
suspension kills the square.
Every nonexceptional four-element opposite corner
is covered by
Corollary~\ref{C:small-opposite-six-source-regular}.
For the nilpotent exception, every projective contracted
source row is covered by
Corollary~\ref{C:small-opposite-P5-all-states}.
The remaining source is exactly
Lemma~\ref{C:six-local-source-normal-form}.
Lemmas~\ref{C:six-four-seven-core-pattern} and
\ref{C:six-four-seven-nilpotent-control} exhaust its
actual seven-state bisets and exclude them.
\end{proof}

% Root, B and D full literal audits passed; both independent numerical gates passed.
\subsection{Seven source elements and seven states opposite the right-zero two-band}

\begin{lemma}[The exact square-band source partition at order seven]
\label{C:seven-graph-source-algebras}
In the accepted complete list of forty-three order-seven
monoids with positive rational second cohomology, forty-one
have minimal band $2\times2$.  Their contracted algebra
$A=\mathbb QM/\mathbb QK$ and one-dimensional left row
module $V$ have the following exact partition:
\[
\begin{array}{c|r}
V\text{ projective over }A&23\\
A\text{ commutative, hence }V\text{ on a local factor}&12\\
A\cong T_2(\mathbb Q),\quad V\text{ nonprojective}&6.
\end{array}
\]
The six final source indices are $3,10,14,16,19,36$ in
\path{code/outputs/a_seven_cone_targets.json}.
\end{lemma}
\begin{proof}
The complete source construction is the one used in
Lemma~\ref{C:triangular-eleven}: it gives $229$ normalized
tables and $41$ ordinary square-band classes, followed by
the identity-adjoined $2\times3$ and $3\times2$ bands.
The new program
\path{code/a_seven_graph_source_algebras.py}
checks the identity and every associativity triple of all
forty-three literal sources, then reconstructs each
three-dimensional contracted algebra on the outside
labels $0,5,6$ and its actual row character $\chi$.

For each projective case it supplies $z\in A$ with
$az=\chi(a)z$ on every basis element and $\chi(z)=1$.
These equations imply $z^2=z$ and split the row-character
quotient by the projective left ideal $Az$.
The next twelve algebras are checked literally to be
commutative.  A finite-dimensional commutative rational
algebra is a product of Artinian local algebras, and the
one-dimensional character is supported on one local factor
with residue $\mathbb Q$.

For each remaining source the certificate supplies a
basis $e,f,a$ with $e+f=1$, $e^2=e$, $f^2=f$, and
\[
 ef=fe=a^2=fa=ae=0,\qquad ea=af=a,
 \qquad \chi(e)=\chi(a)=0,\quad\chi(f)=1.
\]
Thus these are literal upper-triangular algebras with the
stated nonprojective row character.  All rational entries,
basis ranks and products are checked exactly using
\texttt{Fraction}; no modular-rank inference is used.
The complete certificates are in
\path{code/outputs/a_seven_graph_source_algebras.json}.
The independent checker
\path{code/d_verify_seven_graph_source_algebras.py}
imports no producer or matrix library.  It checks all
$14749$ multiplication triples in the forty-three
sources, derives every row character from its literal
band action, and reproduces all $23+12+6$ cases.
Exact inconsistency checks confirm the nonprojective
cases, and every Peirce product and basis rank is
verified over $\mathbb Q$.  Its complete output is
\path{code/outputs/d_seven_graph_source_algebras_independently_verified.json}.
\end{proof}

\begin{lemma}[The necessary core graph and its three-orbit case]
\label{C:seven-state-two-column-core}
Let $M$ have a $2\times2$ minimal band and $|M|\leq9$.
Suppose that the right-zero-two-band collage on at most
seven states has type $\mathbb CP^2$.  After removing the
already impossible common-fixed-state and common-retraction
cases, the minimal band has exactly two state orbits, each
of size two.  Their core graph has one vertex orbit and
one loop-edge orbit.
\end{lemma}
\begin{proof}
Write $P,Q$ for the two left retractions and $\Gamma$ for
their natural graph.  Let $W=H_1(\Gamma;\mathbb Q)$.
The graph-component exact sequence makes the kernel of
the edge map
\[
 H^2(X;\mathbb Q)\longrightarrow
 \operatorname{Hom}_{\mathbb QM}(W,\mathbb Q)
\]
the image of the ordinary component-act cofiber.
By Lemma~\ref{C:nine-unbounded-act-cups} that image has
zero squares.  A projective-plane generator must therefore
have a nonzero edge functional $\lambda$.
For every band idempotent $e$, invariance gives
$\lambda(w)=\lambda(we)$, so $We\ne0$.

A singleton minimal-band orbit would give a band-fixed
state.  If it is an edge state, either endpoint is a
band-fixed vertex.  A band-fixed vertex is $M$-fixed,
since $vm=(ve)m=v(em)=v$.  This gives zero products.
Thus every orbit in $SK$ has two states, one for each
column.  There are at most three such orbits.

The image graph $\Gamma e$ has one state for each of
these orbits; the left retractions give its vertex and
edge structure.  Its first homology is $We\ne0$.
One orbit is insufficient.  With two orbits the only
possibility is one vertex and one loop edge, as asserted.

Suppose there are three orbits.  A graph with three
states and positive first homology has either one vertex
and two loops, or a loop together with an isolated vertex.
Hence $P=Q$ on all six states of $SK$.
If they agree on every state, the common-retraction
argument of Theorem~\ref{C:right-zero-five-state-reduction}
applies without its state bound: an ordinary-act retract
and a suspension quotient exclude a rank-one nonzero
square.

Otherwise there is exactly one further state $z$, and
$Pz\ne Qz$.  Both are core vertices.  Indeed, the
common-image condition rules out either endpoint being
$z$: membership of $z$ in that image would force
$Pz=Qz=z$.  Commutation with
the right band and equality of $P,Q$ on $SK$ give
$(Pz)k=(Qz)k$ for every $k\in K$.  In a disjoint union
of full two-column orbits this says that the two endpoints
are the two different states of one orbit.

Use the literal quotient complex $D$ from
Theorem~\ref{C:two-column-graph-localization}.
Its only remaining edge is $z$.  The two whole column
core graphs have been killed, so its degree-zero space
is the one-dimensional difference of the two column
vertex classes.  The endpoints of $z$ are in different
columns, giving
\[
 D=[\,\mathbb Q\xrightarrow{\ \pm1\ }\mathbb Q\,].
\]
It is acyclic.  The localization theorem makes
$H_4(X;\mathbb Q)=0$, a contradiction.  This excludes
the three-orbit case.
\end{proof}

\begin{lemma}[The complete two-orbit seven-state control]
\label{C:seven-state-triangular-graph-control}
For the six upper-triangular sources of
Lemma~\ref{C:seven-graph-source-algebras}, no seven-state
right-zero-two-band collage whose minimal band has the
two-orbit graph in Lemma~\ref{C:seven-state-two-column-core}
has type $\mathbb CP^2$.
\end{lemma}
\begin{proof}
Label the two minimal-band orbits by $(0,1)$ and $(2,3)$,
with the first as the vertex orbit and the second as the
loop-edge orbit.  The two left retractions both restrict
to $(0,1,0,1)$ on these four states.
Each of the three remaining states has a row-orbit
profile $\gamma_s\in\{0,1\}^2$.  State relabeling sorts
their multiset.  The actual action of each of the two
outside source elements must send a transient state to
a state with profile $\gamma_s\circ\sigma_m$, and its
action on the two core orbits is the given column map.
Every possible such image is tried.  All four outside
products and then all $49$ source multiplication laws
are checked literally.

For each right action, the program tries every possible
commuting idempotent extending $(0,1,0,1)$, and pairs
exactly those with the same image.  Their transient
values can only be $0,1$ or a transient of profile
$(0,0)$.  This is precisely the necessary and sufficient
condition imposed before the idempotence and full
commutation checks.  Common retractions were already
excluded.  The only identifications used are permutations
of equal-profile transient states and interchanging
$P,Q$.

The complete literal counts are
\[
\begin{array}{c|r|r|r}
\text{source index}&\text{right actions}&
   \text{raw pairs}&\text{retained categories}\\\hline
3&446&2450&1270\\
10&198&2185&1152\\
14&446&2450&1270\\
16&198&1396&792\\
19&1426&3666&1769\\
36&1360&3714&1781.
\end{array}
\]
Thus $4074$ right actions and $15861$ raw pairs give
$8034$ literal seventeen-arrow categories.  Of these,
$4161$ fail the required mod-two second Betti number,
and every one of the other $3873$ has a nonboundary
two-cocycle with an explicitly checked square primitive.
The source and complete certificates are
\path{code/a_seven_graph_T2_acts.py} and
\path{code/outputs/a_seven_graph_T2_acts.json}.

The standalone checker
\path{code/b_verify_seven_graph_T2_acts.py}
uses no producer or matrix-library imports.  It freshly
tries all $7^3$ outside transient assignments, then every
$7^3$ retraction assignment, and expands every retained
orbit under the full stabilizer.  It reproduces all
six counts above and checks $7415382$ typed associativity,
$5495256$ boundary-square and $13756896$ square-primitive
equations.  Its complete source-linked output is
\path{code/outputs/b_seven_graph_T2_acts_independently_verified.json}.
This is a full independent numerical enumeration of the
stated normal form, not merely a sample check.
\end{proof}

\begin{lemma}[The nonsquare seven-element band source]
\label{C:two-by-three-seven-state-graph}
Let $M=1+(I_2\times J_3)$.  No right-zero-two-band
collage on at most seven states has type $\mathbb CP^2$.
\end{lemma}
\begin{proof}
Retain the nonzero invariant edge functional $\lambda$
and the exclusion of band-fixed vertices and common
retractions from the preceding arguments.  Every
minimal-band state orbit is a quotient of its three
columns, so has two or three states.  Its column
partition records exactly when two columns give the
same state.  A commuting map between core orbits can
only coarsen this partition.

For $e=x_{00}$, the graph $\Gamma e$ has one state for
each core orbit and has nonzero first homology.
Thus there are two or three core orbits.
With two, there is one vertex orbit and one loop-edge
orbit; the edge partition refines the vertex partition.

If that edge orbit has three states, use the literal
left trivial resolution
\[
 0\longrightarrow(\mathbb QM e)^2
 \xrightarrow{(b_1,b_2)}
 \mathbb QM(1-e)\xrightarrow{\ {-}\,a\ }
 \mathbb QM e\longrightarrow\mathbb Q\longrightarrow0,
\]
where $a=x_{10}-e$ and $b_j=x_{0j}-e$.
It is exact by the same basis calculation as
Lemma~\ref{C:P5-trivial-two-step}: the four vectors
$x_{ij}-x_{i0}$, $i=0,1$, $j=1,2$, are the kernel
of the middle map.  In particular the trivial module
has projective dimension two.
The space $We$ is spanned by the one core loop $c_0$,
whose two other column images are distinct loop edges
$c_1,c_2$.  The second tensor differential is
\[
 (\alpha c_0,\beta c_0)\longmapsto
       \alpha(c_1-c_0)+\beta(c_2-c_0).
\]
Distinct edge supports make it injective.  Therefore
$\operatorname{Tor}_2^{\mathbb QM}(W,\mathbb Q)=0$.
The component terms in the graph-complex exact sequence
are $\operatorname{Tor}_3$ and $\operatorname{Tor}_4$,
which vanish by this resolution.  Hence $H_4(X;\mathbb Q)=0$.

If both core orbits have two states, they have the same
column partition.  The whole right action then factors
through the column-collapse quotient $q:M\to P_5$,
since every image of a band element lies in these
core orbits.  Choosing a representative of each column
block gives an actual monoid section of $q$.
The corresponding graph cofiber over $P_5$ is an actual
homotopy retract of the cofiber over $M$.
The edge-functional map under this section is the
identity on
$\operatorname{Hom}_{M}(W,\mathbb Q)
=\operatorname{Hom}_{P_5}(W,\mathbb Q)$.
Thus the nonzero $\lambda$ makes the projective-plane
generator restrict nontrivially to that retract.
Its pullback spans $H^2(X;\mathbb Q)$, contradicting
the universal $P_5$ graph cup theorem.

Finally suppose that there are three core orbits.
The core graph has only loops, so $P=Q$ on $SK$.
If any orbit has three states, $|SK|\geq7$ and there
are no remaining states; this is the common-retraction
case.  Otherwise all three orbits have two states and
there is at most one transient.

If the core graph has one vertex orbit and two loop-edge
orbits, all three column partitions coincide.  The same
split $P_5$ quotient applies.  If it has a loop-edge
orbit $B$ at a vertex orbit $A$, and an isolated vertex
orbit $C$, the partitions of $A,B$ agree.  Equal
partitions for $C$ again give the split quotient.
Assume instead that its partition is different.

For $P\ne Q$ there must be one transient $z$.
As before, $Pz,Qz$ are the two different points of
one core vertex orbit.  If that orbit is $A$,
commutation with each row of the band forces
$zK\subseteq A\cup B$.  If it is $C$, commutation
forces $zK\subseteq C$.  In either case the graph's
components supported on $A\cup B$ and those supported
on $C$ remain in two disjoint nonempty invariant
families under $M=1+K$.  The component right act
therefore has at least two orbits, so
\[
 H_1(X;\mathbb Q)
   =\widetilde{\mathbb Q\pi_0\Gamma}
      \otimes_{\mathbb QM}\mathbb Q\ne0.
\]
This excludes the final case.
\end{proof}

\begin{theorem}[Seven source elements and seven states opposite $R_3$]
\label{C:seven-source-seven-state-right-zero}
If $R_3$ is the identity-adjoined two-element right-zero
band, no $(R_3,M)$-collage with $|M|\leq7$ and at most
seven cross states has type $\mathbb CP^2$.
In particular the entire seventeen-arrow profile
$(7,3,7)$ is excluded.
\end{theorem}
\begin{proof}
Sources of order at most six are covered for every
state size by
Corollary~\ref{C:six-source-right-zero-all-states}.
For a seven-element source with zero rational second
cohomology, the two vertex restrictions vanish and
the relative suspension kills the square.
Otherwise use the complete forty-three-source list.

For the forty-one square-band sources, the projective
and local-factor alternatives of
Theorem~\ref{C:two-column-graph-localization} cover the
first $23+12$ cases.  The remaining six reduce by
Lemma~\ref{C:seven-state-two-column-core} to the complete
control in Lemma~\ref{C:seven-state-triangular-graph-control};
fewer than seven states were already excluded by
Theorem~\ref{C:right-zero-six-small-sources}.
The source $1+(I_3\times J_2)$ has contracted algebra
$\mathbb Q$ and projective row module, so is covered
by the same localization theorem.  The other nonsquare
source is Lemma~\ref{C:two-by-three-seven-state-graph}.
These cases exhaust the source list.
\end{proof}

% Root and A full literal audits passed.  The110-case independent
% raw-scope/orbit/literal-bar gate is complete and source-hash linked.
\subsection{Eight states for the nonsquare seven-element source}

Write $M=1+(I_2\times J_3)$ and
$R_3=\{1,p,q\}$, where $pq=q$ and $qp=p$.
For an $(R_3,M)$-biset $S$, the left retractions $P,Q$ have
one common image $U$.  Its natural right $M$-graph $\Gamma$
has vertices $U$ and one edge $Ps\to Qs$ for each
$s\in S\setminus U$.  The collage has the actual homotopy model
\[
 X=\operatorname{hocofib}(\Gamma_{hM}\longrightarrow BM).
\]
Throughout this subsection maps between such models are the
actual graph and monoid maps, not hypothetical finite replacements.

\begin{lemma}[Split source quotients and full three-column loops]
\label{C:two-by-three-eight-state-reductions}
In the preceding model suppose $X\simeq\mathbb CP^2$.
Put $K=I_2\times J_3$ and $W=H_1(\Gamma;\mathbb Q)$.
Then no orbit of $SK$ is a singleton, and the projective-plane
class has a nonzero invariant edge functional
$\lambda:W\to\mathbb Q$.  The following cases are impossible:
\begin{enumerate}
\item the right action factors through either split quotient
$M\to1+J_3$ or $M\to P_5$;
\item $H_1(\Gamma e;\mathbb Q)$ is one-dimensional, generated by
a loop whose corresponding core outside-state orbit has all three
columns.
\end{enumerate}
\end{lemma}
\begin{proof}
The graph-component exact sequence identifies the kernel of
the edge-functional map with the image of the ordinary
component-act cone.  Its degree-two products vanish by
Lemma~\ref{C:nine-unbounded-act-cups}.  Thus a class with nonzero
square has $\lambda\ne0$.  A singleton core orbit supplies a
$K$-fixed graph vertex, after applying $P$ if necessary; that
vertex is also $M$-fixed and gives the usual section obstruction.
For every $e\in K$, invariance gives $\lambda=\lambda e$.
Consequently $We=H_1(\Gamma e;\mathbb Q)$ is nonzero.

For either split quotient, its monoid section gives an actual
collage retract, using the same graph with the factored action.
Restriction of the chosen second class retains the same nonzero
edge functional.  The retract's second-cohomology pullback
therefore spans the projective-plane line.  Over $1+J_3$, both
corner spaces have zero rational second cohomology, so the
relative suspension kills all degree-two products.  Over $P_5$
use Theorem~\ref{C:P5-graph-cofiber-cups}.  In either case
the pulled-back nonzero class has square zero.  This argument
does not infer finite homotopy type for the retract.

For the second assertion write $e=x_{00}$,
$a=x_{10}-e$, and $b_j=x_{0j}-e$ for $j=1,2$.
The exact left projective resolution already proved in
Lemma~\ref{C:two-by-three-seven-state-graph} is
\[
0\longrightarrow (Re)^2\xrightarrow{(b_1,b_2)}R(1-e)
 \xrightarrow{{-}\,a}Re\longrightarrow\mathbb Q
 \longrightarrow0,\qquad R=\mathbb QM.
\]
Here $We$ is one-dimensional, generated by the loop $c_0$.
The degree-two tensor map has basis images
$c_1-c_0,c_2-c_0$.  The three full-column core loops have
distinct literal edge supports, so these two cycles are
independent.  Hence $\operatorname{Tor}_2^R(W,\mathbb Q)=0$.
The component-module terms in the augmented-graph exact
sequence have Tor degrees three and four; both vanish by
the displayed length-two resolution.  Therefore
$H_4(X;\mathbb Q)=0$, a contradiction.
\end{proof}

\begin{lemma}[The unequal-partition eight-state control]
\label{C:two-by-three-eight-state-control}
Suppose $SK$ consists of three two-state orbits $A,B,C$.
Both left retractions send $B$ to $A$ and fix $A,C$;
their column partitions are respectively $001,001,010$.
There are two remaining transient states.  This normal form
contains no projective-plane collage.
\end{lemma}
\begin{proof}
If $P=Q$, the graph is a union of bouquets and strict loop
doubling excludes a rank-one nonzero square, using the
component-act zero-product theorem.  Otherwise label the
core by $A=\{0,1\}$, $B=\{2,3\}$, $C=\{4,5\}$.
The common restrictions of $P,Q$ are
$(0,1,0,1,4,5)$.  Each transient has an arbitrary ordered
row-orbit profile in $\{A,B,C\}^2$.  Their unordered
multisets give exactly $45$ right-action normal forms.
All transient retraction values are tried.  Idempotence,
commutation with every right actor, and equality of the
two images are checked literally.  There are $126$ unordered
pairs on these sorted profiles.  Quotienting only by the
actual transient-state stabilizer and interchange of $p,q$
retains $110$ cases.

The producer and complete records are
\path{code/a_eight_state_nonsquare_graphs.py} and
\path{code/outputs/a_eight_state_nonsquare_graphs.json}.
Exactly $100$ have nonzero first mod-two homology,
$4$ have the wrong second Betti number, and the other
$6$ have explicit square primitives for their nonboundary
second cocycles.

The independent checker
\path{code/d_verify_eight_state_nonsquare_graphs.py}
does not import the producer.  For each transient it
tries all $6^2$ first-column images in the two rows,
forces the remaining images by multiplication, and
checks all band equations.  It retains both transient
labels throughout, obtaining $81$ ordered right actions.
It then tries all $8^2$ retraction assignments and every
ordered distinct common-image pair, yielding $420$ pairs.
The complete expansions of the $110$ stored cases under
both transient permutations and both left-label permutations
are disjoint and equal exactly to that fresh raw set.
The sorted-profile counts independently recover $45$ and $126$.

The checker rebuilds each identity-first eighteen-arrow
category and uses the already separate literal bar verifier
from \path{code/b_verify_rightzero_six_isolated.py}.  It checks
$110220$ typed associativity equations, $81840$ boundary-square
equations, $3840$ cocycle equations and $23232$ square-primitive
equations.  No producer or matrix-library routine is imported.
The complete atomic record is
\path{code/outputs/d_eight_state_nonsquare_graphs_independently_verified.json}.
Thus every case in precisely this normal form is excluded.
\end{proof}

\begin{theorem}[The nonsquare source with eight right-zero states]
\label{C:two-by-three-eight-state-right-zero}
Let $M=1+(I_2\times J_3)$.  An $(R_3,M)$-collage with
at most eight states cannot have type $\mathbb CP^2$.
\end{theorem}
\begin{proof}
At most seven states are already excluded by
Theorem~\ref{C:seven-source-seven-state-right-zero}.
Assume therefore $|S|=8$ and use the nonzero $\lambda$
of Lemma~\ref{C:two-by-three-eight-state-reductions}.
Each orbit in $SK$ is a two- or three-state quotient of
$J_3$.  Let $o$ be the number of these orbits.

If $S=SK$, every band action is row-independent, so it
factors through the split row-collapse $M\to1+J_3$.
The same holds with one core orbit, even in the presence
of transients: both row images necessarily enter that
same orbit with its prescribed column.  These cases are
excluded by the first clause of the reduction lemma.
There are at most four core orbits, and four would use
all eight states.  Thus only $o=2,3$ with transients remain.

For $o=2$, the graph $\Gamma e$ has positive first
homology only when one orbit $A$ is a vertex orbit and
the other $B$ supplies a loop.  Both retractions send
$B$ to $A$, so the column partition of $B$ refines that
of $A$.  If $|B|=3$, the full-column loop clause applies.
If $|B|=2$, both orbits have the same two-part partition.
Every transient band image enters one of these two
orbits, so the entire right action factors through the
split column-collapse $M\to P_5$.  Again it is excluded.

Now let $o=3$.  The idempotent graph has three states.
Its nonzero first homology leaves only one vertex and
two loops, or two vertices and one loop.  In particular
the sole edge in the two-vertex case cannot be nonloop:
that would make $We=0$.  If every edge of $\Gamma$ is
a loop, simultaneous loop doubling and the component
argument finish at once.  We retain only graphs with
at least one nonloop edge.

Suppose there is one vertex orbit $A$ and two loop
orbits $B,C$.  If the core has six states, all three
two-state partitions agree, because the maps to $A$
are column-equivariant.  The split $P_5$ quotient applies.
If the core has seven states, $A$ has size two and one
of $B,C$ has size three.  There is exactly one transient.
It must be the nonloop edge joining the two vertices
of $A$; a transient vertex or loop would leave an
all-loop graph.  The resulting graph is connected,
and its first-homology module is literally
$\mathbb Q[B]\oplus\mathbb Q[C]$, with the transient
edge a bridge.  Its cycle coinvariants have dimension
two, one for each core orbit.  The connected graph
comparison then gives $\dim H_2(X;\mathbb Q)=2$,
contrary to the projective-plane rank.

Finally there are two vertex orbits $A,C$ and one loop
orbit $B$, mapped to $A$.  If the core has seven states
and $|B|=3$, the full-column loop clause applies.  Otherwise
$A,B$ both have size two and $C$ has size three.
There is one transient nonloop edge.  Its endpoints cannot
lie in different core vertex orbits: every band image of
the edge is a core state, on which $P=Q$, while distinct
core orbits remain distinct under the band.  If its
endpoints lie in $A$, all its band images lie in $A$ or
$B$; if they lie in $C$, all lie in $C$.  The graph
therefore has two nonempty component families preserved
separately by every actor.  The component act has at
least two orbits, giving $H_1(X;\mathbb Q)\ne0$.

If instead the core has six states, it has two transients.
The partitions of $A,B$ agree.  If that of $C$ also agrees,
the split $P_5$ quotient applies.  Otherwise relabeling the
three columns and the two states in each orbit gives
exactly $001,001,010$.  This is the complete normal form
of Lemma~\ref{C:two-by-three-eight-state-control}.
All core sizes and all surviving vertex/loop distributions
have now been excluded, proving the theorem.
\end{proof}

% Root and A full literal audits PASS; independent1276 raw-orbit/literal-bar gate PASS.
\subsection{Six states opposite the three-element right-zero band}

Write $R_4=\{1,p_0,p_1,p_2\}$, where $p_ip_j=p_j$.
For a left $R_4$-set $S$, its three retractions $P_0,P_1,P_2$
have a common image $U$ and fix that image pointwise.  The natural
action-category graph has vertex set $U$ and, for every
$s\in S\setminus U$, two labeled edges
\[
 (s,1):P_0s\longrightarrow P_1s,\qquad
 (s,2):P_0s\longrightarrow P_2s.
\]
This is the natural graph model used in
Theorem~\ref{C:right-zero-three-five-small-sources}: collapse the
$p_0$ edges in the labeled action graph, which form disjoint
stars.  Commuting right maps preserve the stars and send each
remaining labeled edge at unit speed to an edge of the same label
or to a vertex.  Thus for an $(R_4,M)$-biset this is an actual
right $M$-graph $\Gamma$, and the collage is naturally equivalent
to
\[
 X=\operatorname{hocofib}(\Gamma_{hM}\longrightarrow BM).
\]
In particular this is a full homotopy model, not merely a
second-cohomology comparison.

\begin{lemma}[The square-band six-state reduction]
\label{C:rightzero-three-six-square-core}
Let $|M|=7$ and suppose its minimal ideal is a $2\times2$
rectangular band.  If an $(R_4,M)$-collage on six states has
type $\mathbb CP^2$, then $M$ is one of the six upper-triangular
sources in Lemma~\ref{C:seven-graph-source-algebras}.  The
minimal-band state core consists of two full two-state orbits,
one a vertex orbit and the other an outside-state orbit.  All
three left retractions send the latter to the former, preserving
column labels.  There are exactly two remaining transient states.
\end{lemma}
\begin{proof}
The other $23+12$ square-band sources have a projective row
augmentation or a one-dimensional row augmentation supported on
a local factor.  Theorem~\ref{C:two-column-graph-localization}
excludes their graph cofibers: in the local case bounded rational
homology follows from the hypothesized projective plane and the
actual graph equivalence.  The exact source-algebra partition,
including all rational sections and Peirce bases, is independently
verified in
\path{code/outputs/d_seven_graph_source_algebras_independently_verified.json}.

Put $W=H_1(\Gamma;\mathbb Q)$.  By the graph-component exact
sequence, the kernel of the edge-functional map
\[
 H^2(X;\mathbb Q)\longrightarrow
 \operatorname{Hom}_{\mathbb QM}(W,\mathbb Q)
\]
is the image of the ordinary component-act cone.  All its
degree-two products vanish by Lemma~\ref{C:nine-unbounded-act-cups}.
Consequently a projective-plane generator has a nonzero invariant
edge functional $\lambda$.  For every band idempotent $e$,
$\lambda(w)=\lambda(we)$, so $We\ne0$.

A singleton orbit in $SK$ supplies a band-fixed vertex, after
applying $P_0$ if necessary.  A band-fixed vertex is fixed by
every $m\in M$, since $vm=(ve)m=v(em)=v$.  Its section makes
all cofiber products zero.  Thus each core orbit has two states,
and there are at most three.  The graph $\Gamma e$ has one state
per core orbit.  One orbit cannot support a nonzero cycle.

If there are three orbits, all six states belong to $SK$.
Each right $M$ actor preserves each core orbit individually
and has the same column transformation on every orbit.
Every commuting left retraction preserves the column label.
It follows literally that $\Gamma$ is the disjoint union of
two identical column copies of a fixed finite graph $G$;
$M$ changes only the column-copy index.

Choose a maximal tree in each component of $G$.  There is a
vertex-fixing graph self-map inducing multiplication by two
on $H_1(G;\mathbb Q)$: keep tree edges fixed and replace each
non-tree edge by the path which traverses its associated based
fundamental loop twice, with the tree paths needed to retain its
original endpoints.  Apply the same map to both column copies.
It commutes strictly with $M$, fixes the component set, and is
scalar two on $W$.  Lemma~\ref{C:graph-component-eigenvalues},
together with the zero-product component-act cone, excludes a
rank-one nonzero square.  This rules out three core orbits.

There are therefore two.  Their graph $\Gamma e$ has nonzero
first homology only when it has one vertex state and one outside
state; the outside state supplies its two loop edges.  Each
$P_i$ sends the outside orbit to the vertex orbit.  A map between
full two-column orbits commuting with the band preserves the
column labels, proving the stated normal form.

We shall also omit the case $P_0=P_1=P_2$ on the whole state set.
Its graph is a disjoint union of bouquets.  Strict simultaneous
loop doubling commutes with every actor, so the same component
and eigenvalue argument excludes its rank-one nonzero square.
This omission does not exclude triples having exactly two equal
left retractions.
\end{proof}

\begin{lemma}[The nonsquare source cases]
\label{C:rightzero-three-six-nonsquare}
Neither identity-adjoined rectangular band of order seven
supports an $(R_4,M)$-collage on at most six states with type
$\mathbb CP^2$.
\end{lemma}
\begin{proof}
For $M=1+(I_3\times J_2)$ the contracted algebra is $\mathbb Q$
and the row augmentation is projective.  Apply
Theorem~\ref{C:two-column-graph-localization}.

Let $M=1+(I_2\times J_3)$.  Retain the nonzero invariant edge
functional $\lambda$ and the exclusion of fixed vertices from
the preceding argument.  Each orbit in $SK$ is a two- or
three-state quotient of the column set $J_3$.  The graph
$\Gamma e$ has one state per orbit and positive first homology,
so there are two or three such orbits.

First suppose $S=SK$.  Right multiplication by a band element
depends only on its column on every core orbit.  Thus the entire
right action factors through the row-collapse homomorphism
\[
 1+(I_2\times J_3)\longrightarrow 1+J_3.
\]
Choosing one fixed row gives an actual monoid section.  The
collage over $1+J_3$ is therefore an actual homotopy retract of
the original collage.  Restriction preserves $\lambda$ under
the identical edge-functional spaces, so a projective-plane
generator restricts nontrivially to this retract.  Both corner
spaces of the retract have zero rational second cohomology;
the relative suspension kills all degree-two products there.
Its pulled-back nonzero second class spans the original
rank-one space, a contradiction.  This deals in particular
with three two-state core orbits, and with two three-state
core orbits.

We may assume there are two core orbits and at least one
transient state.  As before, one orbit $A$ consists of vertices
and one orbit $B$ consists of outside states, and all three
$P_i$ map $B$ to $A$.  The column partition of $B$ refines
that of $A$.

If $B$ has three states, put $e=x_{00}$,
$a=x_{10}-e$ and $b_j=x_{0j}-e$ for $j=1,2$.
The literal exact left projective resolution from
Lemma~\ref{C:two-by-three-seven-state-graph} is
\[
 0\longrightarrow(Re)^2\xrightarrow{(b_1,b_2)}
 R(1-e)\xrightarrow{{-}\,a}Re\longrightarrow\mathbb Q
 \longrightarrow0,
 \qquad R=\mathbb QM.
\]
Here $We$ is spanned by the two colored core loops
$c_0^{(1)},c_0^{(2)}$.  The tensor differential sends
$(w_1,w_2)$ to $w_1b_1+w_2b_2$.  Its four basis images are
\[
 c_j^{(t)}-c_0^{(t)},\qquad j,t\in\{1,2\}.
\]
The six core edges have distinct literal supports.  These four
vectors are independent as graph cycles, so
$\operatorname{Tor}_2^R(W,\mathbb Q)=0$.  The component-module
terms in the graph-complex sequence have Tor degrees three and
four, and vanish by the displayed resolution.  Hence
$H_4(X;\mathbb Q)=0$, a contradiction.

If $B$ has two states, so does $A$, with exactly the same
column partition.  Every band image of every transient lies
in these two core orbits.  The whole action therefore factors
through the column-collapse homomorphism $M\to P_5$, which
has an actual section obtained by choosing column representatives.
The $P_5$ graph cofiber is a homotopy retract.  The nonzero
edge functional restricts unchanged, so its second-cohomology
pullback spans the projective-plane line.  The universal
$P_5$ graph cup theorem~\ref{C:P5-graph-cofiber-cups} gives
the final contradiction.
\end{proof}

\begin{lemma}[The complete six-source three-retraction control]
\label{C:rightzero-three-six-T2-control}
In the two-orbit normal form of
Lemma~\ref{C:rightzero-three-six-square-core}, the six remaining
sources have the following complete finite cover:
\[
\begin{array}{c|r|r}
\text{source index}&\text{right-action orbits}&
\text{retained categories}\\\hline
3&46&208\\
10&27&200\\
14&46&208\\
16&27&172\\
19&92&244\\
36&90&244\\\hline
\text{total}&328&1276.
\end{array}
\]
Every retained category fails the projective-plane mod-two
homology or cup-square condition.
\end{lemma}
\begin{proof}
Label the vertex orbit by $0,1$ and the outside-state orbit by
$2,3$.  Each of the two transient states has a row-orbit profile
$\gamma_s\in\{0,1\}^2$.  State relabeling sorts their multiset.
For each of the two nonidentity outside source actors $m$, its
core action is the literal common column map.  Its transient
image $sm$ must have profile $\gamma_s\circ\sigma_m$.
The program tries every such image, then checks all $49$ source
products on all six states.

The three left retractions restrict to $(0,1,0,1)$ on the core.
Their transient values can only be $0,1$ or a transient with
profile $(0,0)$.  Every such map is tried and tested for
idempotence and commutation with the complete right action.
Three idempotents represent the nonidentity actors of $R_4$
exactly when they have the same image: then $P_iP_j=P_j$.
All triples are retained except the all-equal triple already
excluded structurally.  In particular exactly-two-equal triples
are retained.

Right actions are normalized only by permutations of the two
transient labels preserving their profiles.  Left triples are
then normalized by the actual stabilizer of that right action
and by all six permutations of the three nonidentity left labels.
The latter are automorphisms of $R_4$.  Full orbit expansions are
checked disjoint and equal to the pre-normalization sets.  No
independent state quotient of the commuting left action is used.

The producer \path{code/d_rightzero_three_six_T2.py} stores all
right actions, stabilizers and retained left triples, and builds
the identity-first seventeen-arrow multiplication tables.
The complete input, literal TSV and primary cup certificates are
\path{code/outputs/d_rightzero_three_six_T2_input.json},
\path{code/outputs/d_rightzero_three_six_T2_categories.tsv}, and
\path{code/outputs/d_rightzero_three_six_T2_cups.jsonl}.
There are $1016$ wrong second mod-two Betti numbers and $260$
nonboundary second classes with explicit square primitives.

The standalone checker
\path{code/a_verify_rightzero_three_six_T2.py} imports no producer
or matrix library.  It freshly tries all $6^2$ transient values
for each outside actor, all $6^2$ retraction values, and every
ordered common-image triple.  It recovers $430$ raw right actions
and $5490$ ordered left triples, exactly covered by the stored
state and left-label orbit expansions.  It reconstructs and
links every literal table and certificate, and independently
checks $1231340$ typed associativity equations, $918720$
boundary-square equations, $161460$ cocycle equations and
$989820$ square-primitive equations.  The complete atomic record is
\path{code/outputs/a_rightzero_three_six_T2_independently_verified.json}.
This is a full independent normal-form enumeration and literal
bar gate, not merely a verification of selected examples.
\end{proof}

\begin{theorem}[The right-zero-three-band six-state profile]
\label{C:rightzero-three-six-state-excluded}
Let $R_4$ be the identity-adjoined three-element right-zero band.
No $(R_4,M)$-collage with $|M|\leq7$ and at most six cross states
has classifying space $\mathbb CP^2$.  In particular the entire
seventeen-arrow profile $(7,4,6)$ is excluded for this corner.
\end{theorem}
\begin{proof}
If either numerical bound is strict, the category has at most
sixteen arrows and Theorem~\ref{C:CP2-lower17} applies.  Thus
$|M|=7$ and there are six states.  If $H^2(BM;\mathbb Q)=0$,
both vertex restrictions vanish and the relative suspension
kills every degree-two product.  Otherwise the accepted complete
forty-three-source list applies.  Its forty-one square-band
sources are exhausted by
Lemmas~\ref{C:rightzero-three-six-square-core} and
\ref{C:rightzero-three-six-T2-control}; its two nonsquare sources
are Lemma~\ref{C:rightzero-three-six-nonsquare}.  Every branch
has been excluded with its stated actual maps or complete
literal certificates.
\end{proof}

% Root and D full literal audits passed.  Separate independent complete
% source/orbit and literal-bar gates both passed with observed exit zero.
\subsection{The sixth four-element corner with an eight-element source}

Let $R_4=\{1,p_0,p_1,p_2\}$ with $p_ip_j=p_j$.
For an $(R_4,M)$-biset $S$, write $P_i$ for the three
left retractions and $U$ for their common image.  Its natural
right $M$-graph has vertices $U$ and two labeled edges
\[
 (s,1):P_0s\longrightarrow P_1s,
 \qquad (s,2):P_0s\longrightarrow P_2s
 \quad(s\notin U).
\]
The natural right-zero-band comparison identifies the collage
up to actual homotopy equivalence with the graph cofiber
$X=\operatorname{hocofib}(\Gamma_{hM}\to BM)$.

\begin{lemma}[The all-core six-state obstruction]
\label{C:eight-R4-six-all-core}
Suppose $|M|\leq9$, its minimal ideal is a rectangular band
with two or three columns, $|S|\leq6$, and $S=SK$.
Then the graph cofiber above cannot have type $\mathbb CP^2$.
\end{lemma}
\begin{proof}
If a core orbit is a singleton, its image under $P_0$ is
a band-fixed vertex.  Such a vertex is fixed by every
source actor: $vm=(ve)m=v(em)=v$.  The section of the
action-bar projection makes every positive cohomology class
of the cofiber pull back from its suspension quotient, so
all products vanish.

Assume there are no singleton core orbits.  Each orbit is
a two- or three-state quotient of the column set.  Every
right source actor preserves each orbit and acts on it by
the same column transformation.  Each orbit is either
entirely in $U$ or entirely outside $U$.  A commuting map
between two core orbits is the unique column-label map
when it exists; its source column partition must refine
its target partition.

There are at most three core orbits.  If there is at most
one outside orbit and only one vertex orbit, all $P_i$
coincide.  More generally, if there is only one vertex
orbit, uniqueness makes all $P_i$ coincide on every orbit.
The graph is then a union of bouquets.  Strict simultaneous
loop doubling fixes all vertices and commutes with every
right actor.  The component/eigenvalue argument of
Lemma~\ref{C:graph-component-eigenvalues}, using the ordinary
component-act cup theorem through order nine, excludes a
rank-one nonzero square.

The only other case with edges has three core orbits:
two vertex orbits $A,C$ and one outside orbit $B$.
All three have two states, since their total size is at
most six.  If every $P_i$ sends $B$ to the same vertex
orbit, all graph edges are loops and the same argument
applies.  Otherwise $B$ maps onto both $A$ and $C$.
All three two-part column partitions must therefore agree.
The graph is literally two identical column copies of
a fixed labeled graph $G$.  Every right actor changes
only the column-copy index and is the identity on $G$.

Choose a maximal forest in $G$.  A vertex-fixing graph map
which fixes that forest and doubles each associated
fundamental loop induces scalar two on $H_1(G;\mathbb Q)$.
Apply it identically to both column copies.  It is strictly
$M$-equivariant and fixes the component set.  The same
component/eigenvalue argument excludes the square.
If there are no outside orbits, the graph is discrete
and the ordinary-act theorem applies directly.  These
cases exhaust every all-core configuration in the stated
size range.
\end{proof}

\begin{lemma}[The complete two-core-orbit normal forms]
\label{C:eight-R4-six-core-normalforms}
If an $(R_4,M)$-collage with $|M|=8$, $|S|=6$ has type
$\mathbb CP^2$, it occurs among the following normal forms
over the complete $516$-table eight-element source cover:
\[
\begin{array}{c|c|c|c}
\text{minimal band}&|A|&|B|&\text{transient states}\\\hline
2\times2\text{ or }3\times2&2&2&2\\
2\times3&2&2&2\\
2\times3&2&3&1.
\end{array}
\]
Here $A$ is the core vertex orbit, $B$ the core outside-state
orbit, and every $P_i$ is the unique column-label map from
$B$ to $A$.  In the second row their two-part column
partitions agree.  In the last row the partition of $B$
is discrete.  All three two-part partitions of the original
three-column set are retained.
\end{lemma}
\begin{proof}
If $H^2(BM;\mathbb Q)=0$, both corner restrictions vanish
and the relative suspension kills the square.  Otherwise
the minimal ideal has at least two rows and two columns.
It has trivial maximal group: a nontrivial group would
require at least eight minimal-ideal elements and a
further identity.  Thus its shape is $2\times2$, $2\times3$
or $3\times2$.  The complete own-census source cover of
Theorem~\ref{C:eight-act-cones} has $492,12,12$ tables of
these shapes.  Its six rational-$H_2$-zero extras are
harmless and are retained in the finite enumeration.

Put $W=H_1(\Gamma;\mathbb Q)$.  The graph-component
sequence sends a hypothetical projective-plane generator
to a nonzero invariant functional $\lambda:W\to\mathbb Q$:
otherwise it would pull back from the ordinary component
cone, whose products vanish by
Lemma~\ref{C:nine-unbounded-act-cups}.  For every band
idempotent $e$, $\lambda=\lambda e$, so
$We=H_1(\Gamma e;\mathbb Q)\ne0$.

Singleton core orbits are excluded as in the preceding
lemma.  A single core orbit makes $\Gamma e$ a single
vertex, contrary to $We\ne0$.  There are therefore two
or three core orbits.
Three would use all six states and are already excluded.
With two orbits, positive first homology of $\Gamma e$
forces one to be a vertex orbit $A$ and the other to be
an outside-state orbit $B$.  All $P_i$ send $B$ to $A$.
Their column equivariance gives the asserted refinement
of partitions.  The all-core case is excluded, so
$|A|+|B|<6$.  Since each size is two or three and the
map $B\to A$ is surjective, precisely the three displayed
rows remain.  This also proves that no core-orbit count,
partial three-column orbit, or transient vertex is omitted.
\end{proof}

\begin{proposition}[The finite three-retraction enumeration]
\label{C:eight-R4-six-finite-control}
The normal forms in Lemma~\ref{C:eight-R4-six-core-normalforms}
have the following complete, independently verified enumeration:
\[
\begin{array}{l|r}
\text{source tables}&516\\
\text{column-partition/gamma profiles}&5856\\
\text{raw right actions on sorted profiles}&58968\\
\text{right-action state orbits}&45692\\
\text{unordered left triples over right representatives}&144487\\
\text{ordered left triples over right representatives}&503832\\
\text{retained literal categories}&126414.
\end{array}
\]
Every retained category is excluded by its mod-two homology
or its degree-two square.  The separate independent source
and literal-bar gates are recorded below.
\end{proposition}
\begin{proof}
The source file is
\path{code/outputs/c_eight_H2_targets.tsv}.
The identity and row-major band labels are checked literally
in every table.  Write $x_{ij}$ for a band element.  On each
transient state $s$, choose its arbitrary row-orbit profile
$\gamma_s:I\to\{A,B\}$.  Then the band action is exactly
\[
 s x_{ij}=\text{the state of orbit }\gamma_s(i)
                 \text{ with column class }j.
\]
The two transient labels, when present, sort the multiset
of these profiles.  With three columns the vertex partition
is one of $001,010,011$; the outside partition is either
the same partition or $012$.

For an outside actor $m$, its core maps are forced by its
literal column transformation $\tau_m$.  Every retained
partition must be invariant under that transformation.
The possible image of a transient satisfies the necessary
and exhaustive condition
\[
 \gamma_{sm}(i)=\gamma_s(\sigma_m(i)),
\]
where $\sigma_m$ is the literal row transformation.
Every state with that profile is tried.  A source-power
test removes only maps violating an actual power equality.
The search chooses the first still-undetermined source
actor, tries each remaining map, and closes all products
using the right-action law $R_{ab}=R_b\circ R_a$.
Any contradiction rejects the assignment.  Each complete
assignment satisfies all $64$ source products on all six
states.  This deterministic generator-image procedure
retains every actual right action and does not use a
hash collision as orbit equality.

The only state permutations retained in this normal form
interchange the two transient labels when their profiles
agree.  The right action is first normalized by that
group; its actual stabilizer then acts on the left maps.
Every left map has the stated forced core restriction,
and every transient image with constant-$A$ row profile
is tried.  Idempotence and commutation with every source
actor are checked literally.  Three such idempotents
represent $R_4$ exactly when they have the same image.
All multisets of three maps are tried; repeated labels
are retained.  Only all-equal triples are excluded by
strict loop doubling.  Sorting the three left maps uses
the full $S_3$ automorphism group of the nonidentity
right-zero labels.  The remaining state stabilizer gives
the retained categories in the table.

The producer is \path{code/a_eight_rightzero_three_six.cpp}.
It verifies $17545728$ source-action equations,
$6826356$ left-action equations and $18203616$ mixed
equations, and saves all right actions, retained left
triples, per-source counts and literal category tables.
Its three compressed streams are
\path{code/outputs/a_eight_rightzero_three_six_rights.jsonl.gz},
\path{code/outputs/a_eight_rightzero_three_six_cases.jsonl.gz}, and
\path{code/outputs/a_eight_rightzero_three_six_categories.tsv.gz}.

The primary literal bar runner is
\path{code/a_eight_rightzero_three_six_cups.cpp}.
It checks every typed category law and chronological
normalized bar.  Its complete $126414$ records split as
$95602$ wrong second Betti numbers, $30794$ explicit
zero-square primitives for nonboundary second generators,
and $18$ nonzero-square cases with positive third homology.
There are no remaining candidates.  The complete record
checks $157764672$ typed associativity equations,
$26005328$ cocycle equations and $183593828$ primitive
equations.  The certificate stream is
\path{code/outputs/a_eight_rightzero_three_six_cups.jsonl.gz}.

The manifest
\path{code/outputs/a_eight_rightzero_three_six_manifest.json}
records observed exit-zero completions, exact source and
output hashes, every gzip CRC/EOF, and every expected
line count.  It preserves the primary-completion status recorded before
the independent checks.  The later complete independent
records below pin those same immutable producer bytes.
The independent source/orbit gate is
\path{code/d_eight_rightzero_three_six_raw.cpp} followed by
\path{code/d_verify_eight_rightzero_three_six_coverage.py},
with final record
\path{code/outputs/d_eight_rightzero_three_six_coverage_independently_verified.json}.
It uses all $9360$ ordered transient profiles, reverse
outside-actor assignment, every literal transient state map,
and direct multiplication constraints, with no generator
closure, power pruning or orbit pruning.  It finds $85116$
fully labeled right actions and $844470$ ordered nonequal
left triples.  Full transient-$S_2$ and left-label-$S_3$
expansions of the $45692$ right representatives and $126414$
retained categories are disjoint and agree exactly with
these fresh raw sets.  Restriction to sorted transient
profiles gives precisely the producer's $58968$ raw right
actions.  Every one of the $516$ per-source counts, every
fixed manifest hash and every exhausted gzip footer agrees.
The independent engine checks $32684544$ right-action
equations and tries $3058356$ literal retraction candidates.
This is a separate complete enumeration, not just a source
read of the producer.

The separate checker
\path{code/b_verify_eight_rightzero_three_six.py}
reconstructs every literal category directly from its source
monoid, saved right action and ordered left triple, then
uses its independent normalized-bar implementation.
Its complete final record is
\path{code/outputs/b_eight_rightzero_three_six_independently_verified.json}.
It checks all $157764672$ typed associativity equations,
$106693416$ boundary-square equations, $26005328$ cocycle
equations and $183593828$ primitive equations.  The result
is exactly $95602$ wrong second Betti numbers, $30794$
zero squares, and $18$ positive-third-homology cases,
with no survivors.  Every source/case/certificate link,
compressed and decoded hash, gzip CRC and terminal record
is checked, and the independent subprocess has observed
exit code zero.  The source/orbit gate and this literal-bar
gate therefore certify different, necessary parts of the
complete exclusion.  Their joined final record is
\path{code/outputs/a_eight_rightzero_three_six_complete_verified.json}.
\end{proof}

\begin{theorem}[The right-zero-three-band corner through eight sources]
\label{C:eight-source-rightzero-three-six}
No $(R_4,M)$-collage with $|M|\leq8$ and $|S|\leq6$
has type $\mathbb CP^2$.
\end{theorem}
\begin{proof}
For $|M|\leq7$ use
Theorem~\ref{C:rightzero-three-six-state-excluded}.
For at most five states use the accepted all-order-four-corner
five-state result, Corollary~\ref{C:four-corners-five-state-eighteen}.
The remaining exact profile is $|M|=8$, $|S|=6$.
The preceding normal-form lemma and complete independently
verified finite proposition exclude every case of that profile.
\end{proof}

\begin{corollary}[All order-four corners with eight source elements and six states]
\label{C:eight-four-six-excluded}
The eighteen-arrow two-object profile $(8,4,6)$ cannot
realize $\mathbb CP^2$.
\end{corollary}
\begin{proof}
Projective opposite coinvariants are already excluded by
the finite-target small-act theorem.  The five nonprojective
two-band corners are
Corollary~\ref{C:eight-four-six-two-band-excluded}.
The remaining nonprojective corner is the $R_4$ theorem above.
This is a single profile exclusion, not a claim about every
eighteen-arrow two-object category.
\end{proof}

\subsection{Two arrows in every hom-set around a middle corner}

\begin{theorem}[The double-hom middle-corner obstruction]
\label{C:middle-double-homs}
Let $C$ have objects $0,1,2$, with endomorphism monoids
$\{1\},M,\{1\}$ respectively, no backward arrows, and
\[
 |C(0,1)|=|C(1,2)|=|C(0,2)|=2.
\]
Here $M$ is any finite monoid.  Then $BC$ cannot have the
homotopy type of $\mathbb{CP}^2$.
More precisely, after relabeling the three two-point hom-sets,
the composition map has one of the following forms:
\[
\begin{array}{c|c}
\text{composition type}&\text{obstruction}\\\hline
\text{not onto}&H_1(BC;\mathbb Q)\ne0\\
\text{row or column}&\pi_1(BC)\cong\mathbb Z\\
\text{Latin square}&\pi_1(BC)\cong C_2\\
\text{one exceptional entry}&H_2(BC;\mathbb Q)=0.
\end{array}
\]
\end{theorem}
\begin{proof}
Write $S=C(0,1)$, $T=C(1,2)$ and $W=C(0,2)$, with
elements $s_i,t_i,w_i$ indexed by $i\in\{0,1\}$.
Composition is a map
\[
 \kappa:T\times S\longrightarrow W,
 \qquad (t,s)\longmapsto ts,
\]
balanced for the right $M$-action on $T$ and left $M$-action
on $S$.  If its image is a singleton, the unused arrow in
$W$ belongs to no nondegenerate simplex of dimension at least
two.  It is an additional edge between two vertices already
joined in the rest of the nerve, so it contributes a free
summand to first homology.

Suppose that $\kappa$ is onto.  A binary $2$ by $2$ matrix,
up to row, column and value interchanges, is a row matrix, a
column matrix, the Latin square $i+j\pmod2$, or the matrix
with its unique value $1$ at $(1,1)$.  These are exhaustive:
a value occurring once gives the last case, while two
occurrences are adjacent or diagonal.

For the row case, $\kappa(t_i,s_j)=w_i$.  Balance forces
every $M$-actor on $T$ to be the identity.  Choose $s_0,t_0$
as a maximal tree in the nerve's one-skeleton.  In its
fundamental-group presentation they become $1$, and the
composition relations give
\[
 w_0=1,\qquad s_1=1,\qquad w_1=t_1.
\]
The relations $t_0m=t_0$ kill every middle monoid generator.
No remaining relation involves the free generator $t_1$.
Thus the fundamental group is $\mathbb Z$.  The column case
is its opposite.

For the Latin square, balance forces a character
$c:M\to C_2$ such that both actions are translation by
$c(m)$.  With the same tree, composition gives
\[
 w_0=1,\qquad s_1=t_1=w_1,\qquad w_1^2=1,
 \qquad m=w_1^{c(m)}.
\]
The middle multiplication relations impose no further
relation because $c$ is a character.  Conversely these
assignments define a functor from $C$ to the one-object
groupoid $C_2$ and send $w_1$ to its nonidentity element.
The resulting fundamental group is exactly $C_2$.

It remains to handle
\[
 \kappa(t_i,s_j)=w_{ij},
\]
where $ij$ is ordinary multiplication of the binary indices.
If the row and column actions of $m$ are $\tau_m,\sigma_m$,
balance says
\[
 \tau_m(i)j=i\sigma_m(j)\qquad(i,j\in\{0,1\}).
\]
Putting one index equal to zero and then both equal to one
shows that the action pair is either $(\mathrm{id},\mathrm{id})$
or $(c_0,c_0)$, where $c_0$ is the constant map to zero.
There is consequently a character
$\chi:M\to\{0,1\}$ describing the common action on both
augmentation lines.  The following explicit chain
calculation proves the last assertion in the table.
\end{proof}

\begin{lemma}[The one-exception composition chain complex]
\label{C:AND-middle-bar}
In the last case of Theorem~\ref{C:middle-double-homs}, put
$R=\mathbb QM$ and let $\mathbb Q_\chi$ denote the
one-dimensional module given by $\chi$, on either side.
Then
\[
 H_1(BC;\mathbb Q)=H_2(BC;\mathbb Q)=0,
 \qquad
 H_n(BC;\mathbb Q)\cong
 \operatorname{Tor}_{n-2}^{R}
          (\mathbb Q_\chi,\mathbb Q_\chi)
 \quad(n\geq3).
\]
\end{lemma}
\begin{proof}
Let $C_0\subset C$ contain all objects, all middle
endomorphisms, and just the arrows $s_0,t_0,w_0$ between
different objects.  It is a subcategory because these states
are fixed by every middle actor and $t_0s_0=w_0$.
Its object $0$ is initial and its object $2$ terminal.
Thus $BC_0$ is contractible, and the normalized relative
nerve chains compute the reduced homology of $BC$.

Here is the full decomposition of those relative chains.
Write
\[
 e=s_1-s_0,\qquad f=t_1-t_0,\qquad v=w_1-w_0.
\]
As actual left and right modules,
\[
 \mathbb QS=\mathbb Qs_0\oplus\mathbb Q_\chi e,
 \qquad
 \mathbb QT=\mathbb Qt_0\oplus\mathbb Q_\chi f.
\]
For a right module $U$ and left module $V$, denote by
$B_*(U,R,V)$ the normalized two-sided monoid bar.  Its
degree-$q$ basis uses $q$ nonidentity middle arrows.  A
simplex with only an $S$-crossing contributes
$B_*(\mathbb Q,R,\mathbb Q_\chi)$ shifted by one; a
simplex with only a $T$-crossing contributes
$B_*(\mathbb Q_\chi,R,\mathbb Q)$ shifted by one.
The coefficient splittings just displayed make the
two-crossing simplices the direct sum of the three bars
with endpoint coefficients
\[
 (t_0,e),\qquad (f,s_0),\qquad (f,e),
\]
each shifted by two.  The fourth endpoint summand
$(t_0,s_0)$ lies in $C_0$ and is zero in relative chains.
Finally the long arrow contributes $\mathbb Qv$ in degree
one and in no other degree.

These are direct-sum decompositions of chain complexes,
with the following two endpoint differentials.  Deleting
the final crossing on the $(t_0,e)$ bar gives, with its
usual boundary sign, the identity onto the single-$S$
bar.  Its other endpoint deletion is zero since the
augmentation of $e$ is zero.  Likewise the $(f,s_0)$ bar
maps identically, up to sign, onto the single-$T$ bar and
has zero other endpoint deletion.  Their length-two
composition faces vanish because
\[
 \kappa(t_0,e)=0,\qquad \kappa(f,s_0)=0.
\]
Each of these two pairs is therefore the mapping cone of
the identity on its normalized bar, with a degree shift;
both are contractible by the identity-cone contraction.
All middle faces preserve the displayed summands, since
the module splittings are actual $R$-module splittings.
This verifies the cancellation on every simplex, rather
than only on a homology spectral sequence.

On the remaining $(f,e)$ bar both endpoint deletions are
zero.  Its only extra face is the degree-two composition
map
\[
 f\otimes e\longmapsto
 w_1-w_0-w_0+w_0=v.
\]
Consequently the remaining relative complex has
$B_{n-2}(\mathbb Q_\chi,R,\mathbb Q_\chi)$ in degree
$n\geq2$, $\mathbb Qv$ in degree one, and differential
from degree two to one given by this augmentation.
The degree-zero homology of that bar is
\[
 \mathbb Q_\chi\otimes_R\mathbb Q_\chi\cong\mathbb Q:
\]
the balancing relation on an element $m$ equates the
same scalar $\chi_m$ on the two sides.  The displayed
composition map induces the identity on this copy of
$\mathbb Q$.  It therefore kills both degree-one and
degree-two homology of the relative complex.  All higher
degrees are precisely the claimed shifted Tor groups.
\end{proof}

\begin{lemma}[The rational Latin-square calculation]
\label{C:Latin-middle-bar}
In the Latin-square case of Theorem~\ref{C:middle-double-homs},
let $\chi:M\to\{1,-1\}$ be its sign character.  Then
\[
 H_1(BC;\mathbb Q)=H_2(BC;\mathbb Q)=0,
 \qquad
 H_n(BC;\mathbb Q)\cong
 \operatorname{Tor}_{n-2}^{\mathbb QM}
       (\mathbb Q_\chi,\mathbb Q_\chi)
 \quad(n\geq3).
\]
In particular the obstruction survives a rational homology
and cohomology-ring equivalence to this category; the
fundamental-group computation alone is not needed for that use.
\end{lemma}
\begin{proof}
Use the rational vectors
\begin{gather*}
 s_+=\tfrac12(s_0+s_1),\quad
 t_+=\tfrac12(t_0+t_1),\quad
 w_+=\tfrac12(w_0+w_1),\\
 e=s_1-s_0,\qquad f=t_1-t_0,\qquad v=w_1-w_0.
\end{gather*}
The first two averaged vectors are fixed by every $M$-actor,
and the two augmentation lines have the same sign character.
The composition map satisfies
\[
 \kappa(t_+,s_+)=w_+,
 \quad\kappa(t_+,e)=\kappa(f,s_+)=0,
 \quad\kappa(f,e)=-2v.
\]

Collapse all three hom-sets to singletons, retaining $M$,
to obtain a category $C'$ with initial object $0$ and
terminal object $2$.  The induced normalized rational
chain map has a literal section: replace each singleton
cross arrow by its averaged vector, and on a simplex with
both adjacent crossings use their tensor product.  The
three identities just displayed, invariance of the averages,
and their augmentation one show that this section commutes
with every face and degeneracy.  Its image is therefore
the normalized chain complex of the contractible $BC'$.
This is a section of rational simplicial vector spaces;
no section by actual arrows is asserted.

Quotienting by this contractible summand leaves the same
two single-crossing bars and three double-crossing bars as
in Lemma~\ref{C:AND-middle-bar}, with $s_+,t_+$ replacing
$s_0,t_0$.  The two mixed pairs are again literal identity
cones.  The remaining complex has the two-sided
$\chi,\chi$ bar shifted by two and $\mathbb Qv$ in degree
one, with degree-two augmentation multiplication by $-2$.
Since
$\mathbb Q_\chi\otimes_{\mathbb QM}\mathbb Q_\chi\cong\mathbb Q$,
this is an isomorphism on its degree-zero bar homology.
The low-degree vanishing and higher shifted Tor formula
follow exactly as before.
\end{proof}

\begin{corollary}[Rational double-hom obstruction]
\label{C:middle-double-homs-rational}
Every category in Theorem~\ref{C:middle-double-homs} has
either nonzero first rational homology or zero second
rational homology.  Thus no space rationally homology
equivalent to such a category has the rational homology
of $\mathbb{CP}^2$.
\end{corollary}
\begin{proof}
Use the first-homology or fundamental-group computations
in the non-surjective, row and column cases, and
Lemmas~\ref{C:AND-middle-bar} and~\ref{C:Latin-middle-bar}
in the two remaining cases.
\end{proof}

\subsection{Semisimple endpoints and a small middle corner}

\begin{proposition}\label{C:semisimple-ends-small-middle}
Let $C$ be a finite category on objects $0,1,2$, with no arrows
from a larger object to a smaller one.  Suppose the two endpoint
monoid algebras are semisimple over $\mathbb Q$, and the middle
endomorphism monoid has at most three elements.  If
$H_*(BC;\mathbb Q)$ is bounded, then
\[
 H_4(BC;\mathbb Q)=0.
\]
In particular, such a category cannot have type $\mathbb CP^2$.
There is no bound on the endpoint orders or cross-hom cardinalities.
\end{proposition}
\begin{proof}
Put $R_i=\mathbb Q\operatorname{End}_C(i)$, and write
$S_{ij}=\mathbb Q\operatorname{Hom}_C(i,j)$, with its usual
left-$R_j$, right-$R_i$ actions.  Work throughout over
$\mathbb Q$.

The complete list of monoids of order at most three gives the
following elementary alternatives for $R_1$: it is semisimple;
it is the algebra of upper triangular two-by-two matrices or its
opposite; or it is
\[
 \mathbb Q\times D,\qquad D=\mathbb Q[\epsilon]/(\epsilon^2).
\]
For completeness, the two hereditary cases are an identity
adjoined to a two-element left-zero or right-zero band.  Their
difference of band elements spans the one-dimensional radical,
and the two primitive idempotents identify the algebra with the
displayed triangular algebra.  The nonhereditary case is
$\{1,a,z\}$ with $z$ a zero and $a^2=z$.  Its central
idempotent $z$ splits off the first factor, and $a-z$ is the
square-zero generator in the other factor.  The remaining
order-three cases are the cyclic group of order three, the
order-two group with a zero, that group with a new identity,
and the three-element semilattice chain; their rational
algebras are semisimple.  The order-one and order-two cases
are semisimple as well.

In every case both trivial $R_1$-modules have projective
dimension at most one.  In the dual-number case the trivial
module is precisely the projective factor supported at $z$.
Also the middle monoid has zero positive rational homology:
in the band or zero cases its classifying space is contractible,
and the other cases have finite-group homotopy type.

Filter the normalized nerve chains by the number of strict
changes of object.  That number is at most two.  The associated
graded terms are the ordinary corner bars for no change,
the two-sided correspondence bars for one change, and the
three-corner bar for the path $0\to1\to2$ for two changes.
The crossings contribute the horizontal degree $p$; internal
bar degree is $q$.  Exact endpoint coinvariants therefore give
\[
 E^1_{0,q}=0\quad(q>0),\qquad
 E^1_{1,q}=0\quad(q\geq2),\qquad
 E^1_{2,q}=\operatorname{Tor}^{R_1}_q(W,V),
\]
where
\[
 W=\mathbb Q\otimes_{R_2}S_{12},\qquad
 V=S_{01}\otimes_{R_0}\mathbb Q.
\]
The first equality uses the positive-homology assertion for
all three corners.  For the second equality, the direct
$0\to2$ correspondence has semisimple endpoints.  In each
other edge correspondence one endpoint is semisimple and the
other has trivial-module projective dimension at most one.
The formula remains valid if a cross hom-set is empty, in
which case the relevant module or bar term is zero.

For $q\geq2$ no spectral-sequence differential can enter
$(2,q)$.  Its first differential lands at $(1,q)$, which is
zero, and its possible second differential lands at
$(0,q+1)$, which is zero.  There are no larger horizontal
degrees.  There are also no other terms of total degree at
least four.  Consequently the filtration yields actual
vector-space isomorphisms
\begin{equation}\label{C:small-middle-high-homology}
 H_n(BC;\mathbb Q)
 \cong \operatorname{Tor}^{R_1}_{n-2}(W,V)
 \qquad(n\geq4).
\end{equation}
This does not assert that the lower-degree crossing
differentials vanish.

If $R_1$ is semisimple or hereditary, its Tor in degree two
and above is zero.  In the remaining case, only the
dual-number summands of $W,V$ can contribute positive Tor.
Every finite-dimensional $D$-module is a direct sum of free
modules and copies of $\mathbb Q=D/(\epsilon)$, by the
Jordan form of its square-zero operator.  If the respective
simple multiplicities are $b,a$, the elementary periodic
resolution with every differential multiplication by
$\epsilon$ gives
\[
 \operatorname{Tor}^D_j(W_D,V_D)\cong\mathbb Q^{ab}
 \qquad\text{for every }j\geq1.
\]
Thus a nonzero degree-two Tor would, by
\eqref{C:small-middle-high-homology}, give nonzero homology
in every degree at least four.  Bounded homology excludes
that possibility, proving the proposition.
\end{proof}

\begin{corollary}[A small middle corner with a zero]
\label{C:semisimple-ends-small-zero-middle}
The same conclusion holds if the middle monoid has a two-sided
zero and at most five elements, while the two endpoint algebras
remain semisimple over $\mathbb Q$.
\end{corollary}
\begin{proof}
Let $z$ be the middle zero.  The central idempotent $z$ gives
\[
 R_1=\mathbb Qz\times A,\qquad
 A=\mathbb QN/\mathbb Qz,
\]
where $A$ is the contracted algebra of the pointed finite
monoid $N$ and has dimension at most four.  The trivial
left and right $R_1$-modules are the projective factor
$\mathbb Qz$.  Thus the same strict-object filtration has
zero positive internal homology in horizontal degrees
zero and one.  It gives the stronger identification
\[
 H_n(BC;\mathbb Q)
 \cong\operatorname{Tor}^{R_1}_{n-2}(W,V)
 \cong\operatorname{Tor}^{A}_{n-2}(W_A,V_A)
 \qquad(n\geq3).
\]
The second isomorphism is in positive Tor degrees; the
semisimple $\mathbb Qz$ factor contributes none.
If degree-two Tor were nonzero, the contracted
four-dimensional persistence theorem
\ref{B:contracted-four-Tor} would make Tor nonzero in
arbitrarily high degrees.  The displayed identification
would contradict bounded rational homology of $BC$.
Consequently $H_4(BC;\mathbb Q)=0$.
\end{proof}

\begin{corollary}[Semisimple endpoints and two-point cross hom-sets]
\label{C:semisimple-ends-binary-homs}
Let $C$ have three objects $0,1,2$, no backward arrows,
semisimple rational algebras at $0$ and $2$, and
\[
 1\leq|C(0,1)|,\ |C(1,2)|,\ |C(0,2)|\leq2.
\]
The middle monoid is arbitrary.  Then $BC$ cannot have type
$\mathbb CP^2$.
\end{corollary}
\begin{proof}
Pass to the actual cross-act orbits at the two extreme
objects, replacing their corner monoids by trivial ones.
The two orbit-collapse functors induce rational homology,
and hence rational cohomology-ring, equivalences.  Indeed,
endpoint coinvariants are exact; the strict-object bar
filtration used above identifies this functor on each
associated-graded correspondence bar with that exact
coinvariant comparison.  Every resulting forward hom-set
is still nonempty and has at most two elements.

If an adjacent hom-set is a singleton, the singleton-cofinal
deletion removes the middle object; the result has two
trivial corners and is a graph.  Otherwise both adjacent
hom-sets have size two.  If the long hom-set has size one,
triangle compression gives a monoid with a $2\times2$
minimal band, whose rational degree-two products vanish
by Lemma~\ref{B:rank-one-square}.  If the long hom-set
also has size two, Corollary~\ref{C:middle-double-homs-rational}
gives nonzero first rational homology or zero second
rational homology.  Each alternative contradicts the
rational homology or cup square of a projective plane.
\end{proof}

\subsection{Augmentation kernels for a linear category}

\begin{proposition}\label{C:linear-augmentation-Tor}
Let $C$ be a finite category on objects $0,1,2$, with no
backward arrows and with trivial endomorphism monoids at
$0$ and $2$.  Suppose the three forward hom-sets are nonempty.
Put
\[
 N=\operatorname{End}_C(1),\quad R=\mathbb QN,\quad
 V=\mathbb QC(0,1),\quad W=\mathbb QC(1,2),\quad
 U=\mathbb QC(0,2).
\]
Here $V$ is a left $R$-module, $W$ is a right $R$-module,
and all three permutation spaces have their usual
augmentation to $\mathbb Q$.  Write $\bar V,\bar W,\bar U$
for the respective augmentation kernels.  Then, naturally,
\begin{equation}\label{C:linear-augmentation-high}
 H_n(BC;\mathbb Q)\cong
 \operatorname{Tor}^{R}_{n-2}(\bar W,\bar V)
 \qquad(n\geq3).
\end{equation}
There is also an exact sequence
\[
 0\longrightarrow H_2(BC;\mathbb Q)
 \longrightarrow\bar W\otimes_R\bar V
 \xrightarrow{\bar\mu}\bar U
 \longrightarrow H_1(BC;\mathbb Q)
 \longrightarrow0,
\]
where $\bar\mu$ is induced by the actual composition
$C(1,2)\times C(0,1)\to C(0,2)$.
\end{proposition}
\begin{proof}
Let $G$ be the full subcategory on $0,2$.  Its nerve is
the graph with those two vertices and one edge for each
element of $C(0,2)$.  It is connected.  Let $L_*$ be its
rational normalized chain complex, regarded as a
subcomplex of the normalized chains $C_*(NC;\mathbb Q)$.
It has no terms above degree one,
$H_0(L_*)=\mathbb Q$, and $H_1(L_*)=\bar U$.

Describe the quotient complex by the number of strict
changes of object.  With zero changes, only the middle
corner bar $B_*(\mathbb Q,R,\mathbb Q)$ remains.
With one change, its two summands are
$B_*(\mathbb Q,R,V)$ and $B_*(W,R,\mathbb Q)$.
With two changes, it is $B_*(W,R,V)$.  The one-change
and two-change terms have total-degree shifts one and
two.  The horizontal faces are the two augmentation
maps, with the ordinary alternating signs.  The face
composing the two cross arrows has its value in $U$;
it is zero in the quotient.  This description follows
by retaining the ordered string of middle endomorphisms
between the cross arrows in each nondegenerate simplex.
It includes the empty string and all its boundary faces.

Consequently the quotient complex is the total two-sided
bar for the two complexes
\[
 [\,W\xrightarrow{\varepsilon_W}\mathbb Q\,]
 \quad\hbox{and}\quad
 [\,V\xrightarrow{\varepsilon_V}\mathbb Q\,],
\]
where $W,V$ occupy degree one and $\mathbb Q$ degree zero.
Their augmentations are onto, so these complexes are
quasi-isomorphic to $\bar W[1]$ and $\bar V[1]$.
The two-sided bar computes their derived tensor product
over $R$.  These quasi-isomorphisms can also be used
without an abstract derived comparison: in each fixed
internal bar degree all other tensor factors are vector
spaces over $\mathbb Q$, hence preserve exact sequences;
the first-quadrant total-complex comparison then preserves
quasi-isomorphisms.  Thus the quotient has homology
\[
 H_n\bigl(C_*(NC;\mathbb Q)/L_*\bigr)
 \cong\operatorname{Tor}_{n-2}^R(\bar W,\bar V),
\]
with the right side understood as zero for $n<2$.

The long exact sequence for $L_*$ gives
\eqref{C:linear-augmentation-high} in degrees at least
three and gives the displayed low-degree exact sequence.
Its connecting map is composition: on a tensor
$(w-w_0)\otimes(v-v_0)$, the two remaining augmentation
faces cancel, and the direct-edge face is
\[
 wv-wv_0-w_0v+w_0v_0.
\]
This lies in $\bar U$, respects the $R$-balancing
relations, and is precisely $\bar\mu$, up to the
irrelevant simultaneous sign from the chain convention.
\end{proof}

\begin{corollary}\label{C:semisimple-ends-middle-four}
Let $C$ have three linearly ordered objects and nonempty
forward hom-sets.  Suppose its two endpoint monoid
algebras are semisimple over $\mathbb Q$.  If its middle
monoid has at most four elements, or has a two-sided zero
and at most five elements, bounded rational homology of
$BC$ implies
\[
 H_4(BC;\mathbb Q)=0.
\]
In particular these categories cannot have type
$\mathbb CP^2$, at any cross-hom cardinalities.
\end{corollary}
\begin{proof}
Collapse each endpoint monoid to the trivial monoid,
replacing every incident hom-set by its actual orbit set.
The orbit-collapse functors and the exact-coinvariant
comparison of
Corollary~\ref{C:semisimple-ends-binary-homs}
give a rational homology equivalence.  The same
strict-object bar argument applies with arbitrary
nonempty hom-set sizes: semisimplicity makes endpoint
coinvariants exact on every associated-graded bar.
We may apply Proposition~\ref{C:linear-augmentation-Tor}
to this resulting category.

If $|N|\leq4$, adjoin a new absorbing zero $z$ to $N$,
even if $N$ already has another zero.  The contracted
algebra $\mathbb Q(N\sqcup\{z\})/\mathbb Qz$ is
exactly $\mathbb QN$, of dimension at most four.
Theorem~\ref{B:contracted-four-Tor} says that a nonzero
degree-two Tor group for the displayed augmentation
kernels forces nonzero Tor in arbitrarily high degrees.
Equation~\eqref{C:linear-augmentation-high} would then
contradict bounded rational homology of $BC$.

If $N$ has a two-sided zero and at most five elements,
its algebra is the product of the semisimple zero factor
$\mathbb Qz$ and the contracted algebra of dimension at
most four.  Positive Tor is contributed only by the
contracted factor, so the same theorem gives the same
contradiction.  Thus degree-two Tor is zero, and the
degree-four case of
\eqref{C:linear-augmentation-high} proves the conclusion.
\end{proof}

\begin{proposition}[Collapsing the long hom-set]
\label{C:linear-long-hom-collapse}
Under the hypotheses of
Proposition~\ref{C:linear-augmentation-Tor}, let $D$
retain every object and arrow of $C$ except that all
arrows $0\to2$ are identified to a single arrow.
Composition into that hom-set is the unique possible
map.  This defines a category and an actual functor
$C\to D$.  If $BC$ is simply connected and
$u=|C(0,2)|$, then
\[
 BD\simeq BC\vee\bigvee^{u-1}S^2.
\]
In particular, if $BC\simeq\mathbb CP^2$, the
triangle-compression monoid of order
$|N|+|C(0,1)|\,|C(1,2)|$ realizes
$\mathbb CP^2\vee\bigvee^{u-1}S^2$.
\end{proposition}
\begin{proof}
The only nondegenerate simplices containing a direct
$0\to2$ arrow are those edges themselves: the endpoint
corners are trivial.  Thus $ND$ is the simplicial pushout
of $NC$ and the single-arrow nerve along $NG$.
The inclusion $NG\hookrightarrow NC$ is a simplicial
subcomplex inclusion, so realization gives its homotopy
pushout.  The single-arrow nerve is contractible.
The graph $BG$ is a wedge of $u-1$ circles, and its
map into the simply connected space $BC$ is
null-homotopic.  The pushout is therefore
$BC\vee\Sigma BG$, which is the stated wedge.

For completeness the compression monoid is the disjoint
union $N\sqcup(S\times T)$, where
$S=C(0,1)$ and $T=C(1,2)$.  Its old multiplication is
that of $N$, and its other products are
\[
 (s,t)(s',t')=(s,t'),\qquad
 n(s,t)=(ns,t),\qquad (s,t)n=(s,tn).
\]
The identity is the identity of $N$.  The action laws
give associativity in every mixed case.  The triangle-compression comparison of
Theorem~\ref{C:triangle-band-compression} identifies its
classifying space with $BD$; here the long hom-set is
exactly a singleton, so all hypotheses of that comparison
hold.  Its arrow count is the displayed sum.
\end{proof}

% Root and agents B and D independently passed the full literal proof.
\subsection{Extreme corners and projective endpoint coinvariants}

\begin{proposition}[Finite-target removal of an extreme corner]
\label{C:finite-target-extreme-corner}
Let $C$ be a finite category with an object $b$ having no arrows
to any other object. Put $N=C(b,b)$ and let $D$ be the full
subcategory on the other objects. Suppose the trivial right
$\mathbb QN$-module is projective. Replace $N$ by the trivial
monoid and each incoming hom-set $C(d,b)$ by its left $N$-orbit
set, leaving $D$ unchanged. This defines an actual quotient
functor $Q:C\to\bar C$.
Let $G$ be the minimal local group of $N$. Then $BQ$ is a
rational cohomology-ring equivalence and an $\mathbb F_p$
cohomology-ring equivalence whenever $p\nmid|G|$. For every
connected finite CW complex $F$, precomposition gives a weak
equivalence
\[
 \operatorname{Map}(B\bar C,F)\longrightarrow
 \operatorname{Map}(BC,F).
\]
The quotient is surjective on fundamental groups. If $BC$ has
finite CW type, $BQ$ has a homotopy left inverse.
The dual statement for an object with no incoming arrows uses
projectivity of the trivial left $\mathbb QN$-module and right
orbit sets.
\end{proposition}
\begin{proof}
For each $d\in D$ put $S(d)=C(d,b)$. Composition makes $S$ a
contravariant functor from $D$ to finite left $N$-sets. The
mixed-nerve decomposition has the actual natural homotopy-pushout
model
\[
 BC\simeq\operatorname{hocolim}
       \bigl(BD\longleftarrow E\longrightarrow BN\bigr).
\]
Here $E$ is the realization of the category bar whose degree-$k$
term is the disjoint union of $B(S(d_k)\mathbin{/\mkern-6mu/}N)$
over composable strings $d_0\to\cdots\to d_k$ in $D$.
The maps to $BD$ forget the state and its $N$-action string;
the maps to $BN$ retain that action string. This is the same
mixed-nerve model as in
Theorem~\ref{C:finite-target-opposite-corner}, with an arbitrary
category in place of the one-object category $M$.

Choose the minimal idempotent $e$ of
Lemma~\ref{C:projective-coinvariants-minimal-column}.
For every $d$, Lemma~\ref{C:opposite-corner-action-groupoid}
gives an actual functor
\[
 S(d)\mathbin{/\mkern-6mu/}N
   \longrightarrow eS(d)\mathbin{/\mkern-6mu/}G.
\]
These functors commute strictly with precomposition by every arrow
of $D$; the comma contractions have exactly the same naturality.
The groupoid has component set $N\backslash S(d)$ and finite
stabilizer groups contained in $G$. Empty incident hom-sets cause
no exception: both corresponding spaces are empty.

The maps to these component sets are rational and good-prime
homology equivalences. The category bar spectral sequence and the
pushout prove the asserted cohomology equivalences; the actual
functor makes them ring equivalences.
For a finite CW target $F$, Miller's theorem makes the canonical
constant-map inclusions from the mapping spaces of the component
sets weak equivalences onto the mapping spaces of these groupoids.
They are natural for the whole $D$-diagram. Taking the derived
homotopy limits that map the category bars into $F$, and then the
homotopy pullbacks for the pushouts, proves the mapping-space
assertion. No choice of equivariant basepoints or contractions is
needed. This repeats the finite-target proof of
Theorem~\ref{C:finite-target-opposite-corner} objectwise.

The quotient fixes the objects and is onto each hom-set. Every
edge loop in its nerve therefore lifts to an edge loop with the
same sequence of objects, proving fundamental-group surjectivity.
If $BC$ has finite CW type, surjectivity on mapping-space components
for a finite model of $BC$ produces the homotopy left inverse.
Taking opposite categories proves the dual statement.
\end{proof}

\begin{corollary}[The finite-target triangle after endpoint removal]
\label{C:finite-target-semisimple-endpoints}
Let $C$ have three linearly ordered objects, no backward arrows,
and all three forward hom-sets nonempty. Suppose the trivial left
module at the initial endpoint and the trivial right module at the
final endpoint are projective over their rational monoid algebras,
and $BC\simeq\mathbb CP^2$. Collapse both endpoints to their orbit
sets as in the preceding proposition. Write $M$ for the unchanged
middle monoid and write $s,t,u$ for the three resulting forward
hom-set cardinalities, with $u$ the long hom-set.
Then the triangle-compression monoid
\[
 P=M\sqcup(S\times T),\qquad |P|=|M|+st,
\]
has the following properties:
\begin{enumerate}
\item $BP$ is simply connected and has rational homology of
$\mathbb CP^2\vee\bigvee_{u-1}S^2$;
\item every integral degree-two class on $BP$ restricts to zero
along every map from a finite group;
\item its rational degree-two cup pairing has rank one, and it
has an integral degree-two class whose square is nonzero and
primitive modulo torsion in degree four.
\end{enumerate}
No finite-CW-type assertion about $BP$ is included.
\end{corollary}
\begin{proof}
Apply the dual extreme-corner proposition at the initial object
and the extreme-corner proposition at the final object. The orbit
sets and all their remaining actions commute with the other
endpoint quotient, so both operations are actual functors.
Let $Y$ be the classifying space of the resulting category.
It is simply connected by fundamental-group surjectivity, has
the rational cohomology ring of $\mathbb CP^2$, and the composite
quotient has a homotopy left inverse $h:Y\to BC$ by the
mapping-space assertion. As in
Corollary~\ref{C:finite-target-corner-CP2}, this gives an integral
generator $v$ with nonzero primitive free square and nullity of
every integral degree-two group restriction.

The long-hom collapse of
Proposition~\ref{C:linear-long-hom-collapse} is valid for any
simply connected $Y$, without finiteness. It gives an actual
equivalence under $Y$
\[
 BP\simeq Y\vee\bigvee_{u-1}S^2.
\]
Compose an inverse equivalence with
$h\vee1$ and a chosen finite model
$BC\simeq\mathbb CP^2$. The resulting map to
$\mathbb CP^2\vee\bigvee_{u-1}S^2$ is an isomorphism on
integral degree-two cohomology. Miller's theorem therefore gives
the asserted nullity for every such class. The rational pairing
and the primitive free square follow from the actual wedge
decomposition and the corresponding assertions for $Y$.
\end{proof}

\begin{theorem}[Projective endpoint coinvariants through sixteen arrows]
\label{C:semisimple-endpoints-sixteen}
Let $C$ have three linearly ordered objects, no backward arrows,
and nonempty forward hom-sets. Suppose the trivial left module at
the initial endpoint and the trivial right module at the final
endpoint are projective over their rational monoid algebras.
If $C$ has at most sixteen arrows, then
$BC\not\simeq\mathbb CP^2$. In particular this applies when both
endpoint algebras are semisimple over $\mathbb Q$. There is no
separate bound on the middle monoid or the individual hom-set sizes.
\end{theorem}
\begin{proof}
Suppose otherwise and apply the endpoint-removal corollary. If the
middle monoid has at most four elements, apply
Corollary~\ref{C:semisimple-ends-middle-four} to the collapsed
category, whose endpoint algebras are now $\mathbb Q$. This
contradicts its bounded rational homology and nonzero fourth
homology. Thus write $m=|M|\geq5$.
Orbit sets can only decrease cardinalities, so
\[
 m+s+t+u+2\leq16,\qquad u\geq1.
\]
The triangle-compression monoid has minimal band $s\times t$.
If either $s$ or $t$ is at most two, its rational degree-two
products vanish by the balanced presentation and
Corollary~\ref{cor:B-two-row-general}. Their pullbacks to $Y$
would vanish, contradicting its projective-plane square.
Thus $s,t\geq3$.

The numerical inequality now leaves exactly the following
possibilities, up to opposition:
\[
\begin{array}{c|c|c}
 (s,t)&\text{middle and long-arrow bounds}&\text{obstruction}\\\hline
 (3,3)&m\geq5,\ u\leq8-m\leq3&
            \text{all-rank three-by-three zero products}\\
 (3,4)&m\geq5,\ u\leq7-m\leq2&
            \text{rank-two zero products or rank-one parity}\\
 (3,5)&m=5,\ u=1&\text{three-by-five source of order }20\\
 (4,4)&m=5,\ u=1&\text{four-by-four source of order }21.
\end{array}
\]
For the first two rows, apply
Lemma~\ref{C:small-band-finite-target-cups} to the actual
compression. Its rational second-cohomology rank is $u$ and
all integral group restrictions vanish by the preceding corollary.
The rank-one case has a nonzero primitive square, so its even-square
exception is excluded as well. For the last two rows, use
Propositions~\ref{B:three-by-five-finite-target-transfer} and
\ref{B:four-by-four-finite-target-transfer}, respectively.
Here $u=1$, so the entire rational homology is that of
$\mathbb CP^2$, including vanishing third homology. All stated
transfer hypotheses are therefore supplied by the actual maps,
not by an inferred finite type of the compression.
Each possibility is impossible, proving the theorem.
\end{proof}

% Root literal audit and independent algebra audit by D: PASS.
\subsection{Projective endpoint coinvariants through seventeen arrows}

\begin{lemma}[Three degree-two classes at a five-dimensional quotient]
\label{C:three-by-four-five-outside-rank-three}
Let $P$ have rectangular-band minimal ideal of shape $3\times4$
or $4\times3$, with exactly five elements outside that ideal.
Suppose its rational homology is that of
$\mathbb CP^2\vee S^2\vee S^2$.  Then every product of two
rational degree-two classes on $BP$ is zero.  No finite-CW-type
or local-group-nullity hypothesis is required in this lemma.
\end{lemma}
\begin{proof}
After opposition the band--Tor modules have dimensions
$\dim V=2$, $\dim W=3$ over the five-dimensional contracted
algebra $A$, and their Tor dimensions are
\[
 \dim\operatorname{Tor}_{0,1,2}^A(W,V)=(3,0,1),\qquad
 \operatorname{Tor}_j^A(W,V)=0\quad(j\geq3).
\]
If $A$ is local, the dimension of its residue division algebra
divides both module dimensions, so that residue algebra is
$\mathbb Q$.  If $V$ has top dimension two, then $V$ is the
sum of two copies of the simple module; its degree-zero Tor
has even dimension, a contradiction.  Otherwise $V$ is cyclic.
Write $0\to I\to A\to V\to0$, where $I$ is nonzero.
Since degree-one Tor vanishes, tensoring gives an injection
$W\otimes_A I\to W$.  The following surjection
$W\to W\otimes_A V$ is an isomorphism, because both spaces
have dimension three.  Hence $W\otimes_A I=0$.
But this tensor product surjects onto
\[
 (W/WJ)\otimes_{\mathbb Q}(I/JI),
\]
which is nonzero by Nakayama's lemma.  This is impossible.

For nonlocal $A$, the Tor-persistence classification in
Proposition~\ref{B:five-nonlocal-exceptions} leaves the
zero-composite line and Nakayama algebras.  The complete
indecomposable calculation in Lemma~\ref{B:five-exceptional-targets}
applies before its particular module-dimension specialization.
For the line, it forces
\[
 V=S_1\oplus S_3,\qquad
 W=T_1\oplus T_3^{\oplus s}\oplus(e_1A)^{\oplus p}
                         \oplus(e_2A)^{\oplus q}.
\]
The equations are $s+2p+2q=2$ and $1+s+p=3$;
thus $W=T_1\oplus T_3^{\oplus2}$.
The actual three-point row augmentation is a sum of rational
characters, so Corollary~\ref{B:three-row-semisimple-cups}
makes all degree-two products zero.
For the Nakayama algebra the same table forces $V=L$ and
$W=T_2\oplus(e_1A)^{\oplus p}\oplus(e_2A)^{\oplus q}$.
Its dimension equation $3p+2q=2$ gives $p=0,q=1$,
and then degree-zero Tor has dimension one rather than three.
This excludes the last case and proves the lemma.
\end{proof}

\begin{theorem}[Projective endpoint coinvariants through seventeen arrows]
\label{C:semisimple-endpoints-seventeen}
Let $C$ have three linearly ordered objects, no backward arrows,
and nonempty forward hom-sets.  Suppose that the trivial left
module at the initial endpoint and the trivial right module at
the final endpoint are projective over their rational monoid
algebras.  If $C$ has at most seventeen arrows, then
$BC\not\simeq\mathbb CP^2$.
\end{theorem}
\begin{proof}
Suppose otherwise.  The actual endpoint-removal functors of
Corollary~\ref{C:finite-target-semisimple-endpoints} give a
triangle with trivial endpoints and sizes $(m;s,t,u)$.
Its compression has rational homology
$\mathbb CP^2\vee\bigvee_{u-1}S^2$, rational pairing of
rank one, and an integral class with nonzero primitive free
square.  Every integral degree-two class restricts to zero
on every finite group.  The budget is
\[
 m+s+t+u+2\leq17.
\]
The existing argument for sixteen arrows excludes $m\leq4$
and $\min(s,t)\leq2$ without any bound on the other sizes.
Thus take $m\geq5$, $3\leq s\leq t$, and $u\geq1$.
The entire remaining numerical scope is
\[
\begin{array}{c|c|c}
 (s,t)&\text{bounds}&\text{remaining treatment}\\\hline
 (3,3)&m+u\leq9&\text{all-rank zero products}\\
 (3,4)&m+u\leq8&u=3\text{ is the new case}\\
 (3,5)&m+u\leq7&\text{rank one or two}\\
 (3,6)&(m,u)=(5,1)&\text{five-outside module obstruction}\\
 (4,4)&m+u\leq7&\text{rank one or two}\\
 (4,5)&(m,u)=(5,1)&\text{five-outside module obstruction}.
\end{array}
\]
The first two rows use Lemma~\ref{C:small-band-finite-target-cups},
except for the newly proved rank-three case, which uses
Lemma~\ref{C:three-by-four-five-outside-rank-three}.
The rank-one $3\times5$ and $4\times4$ cases use
Propositions~\ref{B:three-by-five-finite-target-transfer} and
\ref{B:four-by-four-finite-target-transfer}; their compression
orders are at most twenty-one and twenty-two, respectively.
The two rank-two cases have compression orders twenty and
twenty-one, so Theorem~\ref{B:rank-two-rational-through22}
applies with the stated local nullity and rational homology.

Finally the proofs of
Lemmas~\ref{C:three-by-six-five-outside-CP2} and
\ref{C:four-by-five-five-outside-CP2} use only the rational
Tor dimensions $(1,0,1,0,\ldots)$ and zero rational squares,
not finite CW type.  In each, the local case has a rank-one
balanced form, and the two nonlocal algebras are excluded by
the displayed module equations or by the rank-one square
primitive.  Those proofs therefore apply to the present
finite-target compressions as well.
Every numerical case is excluded.  At no point is a finite
CW model inferred for an endpoint quotient or compression.
\end{proof}

\subsection{The five-element middle with four and six adjacent arrows}

\begin{proposition}\label{C:middle-five-four-six}
Let $C$ have objects $0<1<2$, trivial endomorphism monoids at $0$ and
$2$, a five-element middle monoid $N$, and hom-set cardinalities
\[
 |C(0,1)|=4,\qquad |C(1,2)|=6,\qquad |C(0,2)|=1.
\]
Then $BC$ is not homotopy equivalent to $\CP^2$.
The same conclusion holds for the opposite category.
\end{proposition}

\begin{proof}
Write $A=\Q N$, $V=\ker(\Q C(0,1)\to\Q)$ and
$W=\ker(\Q C(1,2)\to\Q)$.  Their dimensions are three and five.
The exact low-degree sequence and high-degree formula of
Proposition~\ref{C:linear-augmentation-Tor}, with its single long arrow,
would give
\[
 \dim\operatorname{Tor}_{0,1,2}^{A}(W,V)=(1,0,1),
 \qquad \operatorname{Tor}_i^A(W,V)=0\quad(i\geq3).
\]
The algebra $\Q N$ is nonlocal.  Indeed, locality would force every
idempotent of the finite monoid $N$ to be its identity, since a distinct
monoid basis element is neither the zero algebra element nor the algebra
identity.  A finite monoid with only that idempotent is a group.
Maschke's theorem then makes $\Q N$ semisimple; a semisimple local algebra
is a division algebra, whose augmentation onto $\Q$ is an isomorphism.
This contradicts $|N|=5$.

Adjoining a new absorbing zero exhibits $A$ as a contracted monoid
algebra, so Proposition~\ref{B:five-nonlocal-exceptions} leaves precisely
the two exceptional algebras.  Use the indecomposable module notation and
Tor tables of Lemma~\ref{B:five-exceptional-targets}.
For the line algebra, the degree-two condition forces exactly one
$S_3$ in $V$ and one $T_1$ in $W$; the degree-one condition excludes
$S_2$ from $V$ and $T_2$ from $W$.  Hence
\[
 V\in\{2S_1\oplus S_3,\ S_3\oplus L_a,\ S_3\oplus L_b\},
 \qquad W=T_1\oplus T_3^s\oplus R_a^p\oplus R_b^q,
 \quad s+2p+2q=4.
\]
In these three cases the tensor dimensions are respectively
$2+s+2p$, $s+p+q$, and $2s+q$.
The first cannot equal one.  In the second, combining a value of one
with the displayed dimension equation would give $s=-2$.
The third gives $s=0$, $p=q=1$.  Thus the only line-algebra pair is
\begin{equation}\label{C:middle-five-four-six-pair}
 V=S_3\oplus L_b,\qquad W=T_1\oplus R_a\oplus R_b.
\end{equation}
For the Nakayama algebra, its table gives exactly two pairs with these
module dimensions and with $(\dim\operatorname{Tor}_1,
\dim\operatorname{Tor}_2)=(0,1)$:
\[
 \begin{array}{c|c|c}
 V&W&\dim(W\otimes_A V)\\\hline
 S_2\oplus Ae_2&R\oplus e_1A&2\\
 S_2\oplus Ae_2&T_2\oplus(e_2A)^2&6.
 \end{array}
\]
Neither has tensor dimension one.  The small exact enumeration of these
displayed tables is retained in
\path{code/b_five_middle_Tor_profiles.py}; it makes no finite-set
realizability assumption.

It remains to examine actual finite-set realizations of
\eqref{C:middle-five-four-six-pair}.  The complete multiplicative-basis
classification in \path{code/outputs/b_five_line_monoid_bases.json}
has 48 oriented contracted line bases.  An ordinary five-element monoid
basis is exactly a retained basis whose multiplication never uses the
contracted zero.  There are eight such records, with indices
$0,1,2,3,30,31,32,33$.  This is complete because adjoining a new zero to
any ordinary five-element middle produces a member of the complete
six-element source census.

For each of these eight literal monoids, enumerate all left actions on
four points and all right actions on six points.  The already audited
pointwise action constraint algorithm enforces every equation
$L_{mn}=L_mL_n$ or $R_{mn}=R_nR_m$, starting with the literal identity
action; it does not assume that abstract modules lift to finite sets.
Let $\chi_1,\chi_2,\chi_3$ be the stored ordered simple characters and
$a,b$ the stored radical-arrow vectors.  The exact selection conditions
on the augmentation modules are
\[
 \begin{array}{c|c|c|c}
 &\operatorname{tr}(m)&\operatorname{rank}(a)&\operatorname{rank}(b)\\\hline
 V&\chi_2(m)+2\chi_3(m)&0&1\\
 W&2\chi_1(m)+2\chi_2(m)+\chi_3(m)&1&1.
 \end{array}
\]
Here the trace is the number of fixed points minus one.  Zero and
rank-one conditions are checked integrally, using all two-by-two minors.
The simple multiplicities and the ranks of the two arrows uniquely
specify the decompositions in \eqref{C:middle-five-four-six-pair}.

The complete action search examines 5,112 labelled four-point actions
and 1,488,464 labelled six-point actions.  State conjugacy gives the
following retained action counts:
\[
 \begin{array}{c|rrrrrrrr}
 \text{basis index}&0&1&2&3&30&31&32&33\\\hline
 \text{left actions}&1&1&1&7&1&7&1&2\\
 \text{right actions}&8&4&28&2&29&2&11&4\\
 \text{pairs}&8&4&28&14&29&14&11&8.
 \end{array}
\]
The two sets may be relabelled independently because their composition
has a one-point target.  Thus these 116 pairs cover every possible
category in the remaining module case.

For all 116 literal categories the normalized bar complex over $\F_2$
has chain dimensions $(15,80,320,1280)$ in degrees one through four,
boundary ranks $(2,13,66,254)$, and therefore
$(b_1,b_2,b_3)=(0,1,0)$.  A nonboundary degree-two cocycle and an explicit
degree-three primitive for its square are retained for every category.
The separate Python verifier rebuilds the normalized complexes and
checks all 70,528 typed associativity equations, 37,120 cocycle equations,
and 148,480 square-primitive equations.  It also checks the full
Cartesian linkage from actual action pairs to the literal tables.
Thus every remaining degree-two generator has zero square over $\F_2$,
contrary to $H^*(\CP^2;\F_2)$.
\end{proof}

\begin{remark}\label{C:middle-five-four-six-certificates}
The source/action wrapper is
\path{code/b_middle5_triangle18_actions.cpp}; the literal categories and
their square primitives are
\path{code/outputs/b_middle5_triangle18_categories.json} and
\path{code/outputs/b_middle5_triangle18_cups.jsonl}.
The independent bar/action-linkage check is
\path{code/b_verify_middle5_triangle18.py}, with result
\path{code/outputs/b_middle5_triangle18_independently_verified.json}.
The latter checks the complete retained output and all exact equations;
it does not claim a second independent numerical action enumeration.
\end{remark}

\subsection{A six-element middle with three and six adjacent arrows}

\begin{proposition}\label{C:middle-six-three-six}
Let $C$ have objects $0<1<2$, trivial endomorphism monoids at its
extremes, a six-element middle monoid $N$, and
\[
 |C(0,1)|=3,\qquad |C(1,2)|=6,\qquad |C(0,2)|=1.
\]
Then $BC\not\simeq\CP^2$.  The opposite profile is excluded as well.
\end{proposition}

\begin{proof}
Suppose otherwise.  Triangle compression
(Theorem~\ref{C:triangle-band-compression}) gives an actual monoid
with the same homotopy type, a $3\times6$ minimal band, and contracted
algebra $A=\Q N$.  Write $V,W$ for its row and column augmentation
modules.  Thus $\dim V=2$, $\dim W=5$, and
\[
 (\dim\operatorname{Tor}_0^A(W,V),
   \dim\operatorname{Tor}_1^A(W,V),
   \dim\operatorname{Tor}_2^A(W,V))=(1,0,1),
 \qquad \operatorname{Tor}_i^A(W,V)=0\ (i\geq3).
\]
If $V$ is simple, Lemma~\ref{B:six-dimensional-simple-row} makes its
degree-two Tor zero or even-dimensional.  If it is the sum of two
one-dimensional modules, Corollary~\ref{B:three-row-semisimple-cups}
annihilates all products.  For a nonsplit self-extension,
Proposition~\ref{B:three-row-dual-cup-radical} explicitly annihilates
the square when second cohomology has rank one.
Consequently the actual augmentation action on $V$ has the full
triangular two-by-two algebra as its image.

We use the complete row-module calculation underlying
Proposition~\ref{B:six-row-census}, retaining its hypotheses precisely.
Although that proposition states the target tensor dimension as two,
every discard in its proof is independent of that dimension.
Projectivity, projective dimension one, the square-zero path criterion,
and the explicit split-syzygy recurrences exclude nonzero degree-two
Tor followed by eventual vanishing.  The nonsplit-factor lemma has the
same premise.  The last local-corner discard uses only $\dim W=5$:
the putative perfect complex has Euler characteristic $5-2d$, odd
rather than divisible by four.  None of these arguments uses
$\dim(W\otimes_A V)=2$.  They therefore apply to the displayed
rank-one target without modification.

Adjoin a new zero to $N$ and use the accepted complete seven-element
zero-monoid census.  An ordinary six-element middle is recognized by
the exact condition that products of the six nonzero labels never
equal the new zero.  Among the previously checked triangular row
modules, the ordinary sources have $3065$ rational module records.
Their disjoint disposition is
\[
 \begin{array}{l|r}
 \text{nonsplit or nonbasic factor obstruction}&31\\
 \text{projective or projective dimension one}&2369\\
 \text{square-zero path persistence}&507\\
 \text{explicit split-syzygy recurrence}&82\\
 \text{odd four-dimensional local-corner Euler obstruction}&4\\
 \text{remaining square-zero modules}&52\\
 \text{remaining projective-dimension-two modules}&20.
 \end{array}
\]
These are subsets of the already independently verified source,
homomorphism, module-partition, and exact projective-kernel records;
there is no new algebra classification assumption.

For the remaining 72 rational modules, retain every stored literal
homomorphism to the thirteen-element augmentation-image monoid from
Lemma~\ref{B:three-point-triangular-target}.  Each nonzero image
specifies its unique nonconstant three-point transformation.
For each zero image try all three constant transformations, and check
all literal left-action equations.  The 1160 assignments contain
84 actual left-action orbits on 45 ordinary six-element middle
sources.  An independent generator-image enumeration starts instead
with every actual three-point action on those ordinary monoids.
It finds 566 action orbits in total and, by exact rational extension
coordinates, recovers precisely the same 84 selected actions.
Thus no affine constant lift is omitted.

For each of the 45 middle tables enumerate all right actions on six
points by the previously audited pointwise constraint algorithm.
There are $24296103$ labelled right actions.  State conjugacy is taken
independently on the two adjacent hom-sets, which is legitimate
because their composition has a one-point target.
Before constructing a category, compute the exact balanced tensor
over $\F_2$.  Its ten generators are the elementary tensors of the
five column and two row augmentation basis vectors, and its relations
are exactly
\[
 (w m)\otimes v-w\otimes(mv),\qquad m\in N.
\]
Lemma~\ref{lem:B-band-Tor}, or its literal triangle comparison,
identifies this tensor dimension with $b_2(BC;\F_2)$.
Retaining only dimension one is therefore a necessary target filter,
not an assumption about rational ranks or torsion.

The filter leaves 33867 literal eighteen-arrow categories.
Every one has a nonboundary normalized two-cocycle over $\F_2$ and
an explicit normalized three-cochain whose coboundary is its square.
The separate bar verifier reconstructs all chains and checks
25467984 typed associativity equations, 14901480 cocycle equations,
and 74507400 primitive equations.  A separate linkage pass reconstructs
each literal table from its specified middle monoid and two actions,
checks every action equation, verifies every retained record exactly
once, and matches all 45 complete source counts and the final footer.
All the squares are zero, contradicting the mod-two cohomology ring
of $\CP^2$.
\end{proof}

\begin{remark}\label{C:middle-six-three-six-verification}
The exact residual extraction and literal search are
\path{code/b_middle6_three_six_inputs.py} and
\path{code/b_middle6_three_six_search.cpp}.
The independent three-point action gate is
\path{code/b_verify_middle6_three_point_lifts.py}, with result
\path{code/outputs/b_middle6_three_point_lifts_independently_verified.json}.
The independent literal-category linkage and normalized-bar checks are
\path{code/b_verify_middle6_three_six.py} and
\path{code/b_verify_middle6_three_six_cups.cpp}; their final result is
\path{code/outputs/b_middle6_three_six_independently_verified.json}.
All 33867 actual cocycles and square primitives are retained in
\path{code/outputs/b_middle6_three_six_cups.jsonl}.
The right-action completeness premise is the audited pointwise
constraint algorithm; no second independent numerical enumeration of
all 24296103 labelled right actions is claimed.
\end{remark}

\subsection{Constant composition in a linear category}

\begin{lemma}[The graph retraction for constant composition]
\label{C:linear-constant-composition-retract}
Let $C$ have objects $0,1,2$, trivial endomorphism monoids at
$0$ and $2$, middle endomorphism monoid $N$, and nonempty forward
hom sets $S=C(0,1)$, $T=C(1,2)$ and $U=C(0,2)$.
Suppose there are no other arrows between distinct objects.
If all composites $ts$, with $s\in S$ and $t\in T$, equal a
single $u_0\in U$, the two-object category with $|U|$ parallel
arrows is an actual categorical retract of $C$.
In particular $\bigvee^{|U|-1}S^1$ is a homotopy retract of $BC$.
\end{lemma}
\begin{proof}
Include the parallel-arrow category as the full subcategory on
$0,2$. Define a functor back by identifying objects $0,1$,
sending every middle endomorphism and every arrow in $S$ to
the identity, every arrow in $T$ to $u_0$, and each arrow in
$U$ to itself. Middle actions preserve these assignments;
the only composition involving two nonidentity interobject
arrows is preserved by the assumed constant-composition formula.
The composite with the inclusion is the identity functor.
The nerve of the parallel-arrow category realizes to a graph
with two vertices and $|U|$ edges, proving the assertion.
\end{proof}

\begin{lemma}[Constant profiles force constant composition]
\label{C:linear-constant-profiles}
In the setting of Lemma~\ref{C:linear-constant-composition-retract},
suppose there is $t_0\in T$ such that for every $s\in S$ some
$n_s\in N$ acts constantly on $S$ with value $s$ and constantly
on $T$ with value $t_0$. Then all composites $ts$ have the same
value. The analogous assertion holds after interchanging $S,T$
and taking the opposite category.
\end{lemma}
\begin{proof}
Writing $\beta(t,s)=ts$, associativity gives
\[
 \beta(t,s)=\beta(t,n_s s')
 =\beta(tn_s,s')=\beta(t_0,s')
 \qquad(t\in T,\ s,s'\in S).
\]
The right side can be evaluated at any one fixed $s'$, so the
left side is constant. Apply the same argument to $C^{\mathrm{op}}$
for the symmetric form.
\end{proof}

\begin{corollary}[The eight closed-actor controls]
\label{C:rank-two-closed-actor-controls}
Consider the eight specified faithful rank-two order-twenty
action models in
\path{code/outputs/d_rank_two_small_inflation_controls.json}.
Generate a middle actor submonoid from the identity and all
three nonband actors. With adjacent hom-set sizes four and
four and a two-element long hom set, none of these controls
produces a category of type $\mathbb CP^2$ through twenty-one
arrows.
\end{corollary}
\begin{proof}
For hull $137$, cases $152,169,279,291$, the generated actor
submonoid orders are respectively $8,12,16,16$; the opposite
four controls in hull $188$ have the same orders. The category
has $|N|+12$ arrows, so only the two order-eight cases fit the
stated bound. In case $152$ the generated actor indices are
\[
 \{0,4,8,12,16,17,18,19\}.
\]
The four constant actors $4,8,12,16$ have every possible row
value and the common column value $3$.
Lemma~\ref{C:linear-constant-profiles} applies, as does its
opposite in case $168$ of hull $188$. The graph retraction
therefore supplies an $S^1$ retract, excluding a projective plane.

The direct independent construction in
\path{code/c_closed_actor_triangle_controls.py} checks the eight
literal generated submonoids and every associativity equation.
It also computes the actual balanced tensor set by all relations
$(s,tn)\sim(ns,t)$; it has one component in each control.
Consequently the two maps to the two-element long hom set are
both constant. All four resulting twenty-arrow tables are
constructed literally and satisfy every category axiom; their
normalized bar complexes have first mod-two Betti number one.
The complete record is
\path{code/outputs/c_closed_actor_triangle_controls.json}.
This is an exclusion of these eight specified controls, not a
classification of all closed actor submonoids of order at most nine.
\end{proof}

\subsection{Closed middle actors with four adjacent arrows}

\begin{theorem}[Two long arrows and four arrows on each adjacent side]
\label{C:linear-four-by-four-two-long-21}
Let $C$ have three linearly ordered objects $0,1,2$, trivial
endomorphism monoids at $0$ and $2$, and no backward arrows.
Suppose
\[
 |C(0,1)|=|C(1,2)|=4,\qquad |C(0,2)|=2.
\]
If $C$ has at most twenty-one arrows, then $BC$ is not homotopy
equivalent to $\mathbb CP^2$.
\end{theorem}

\begin{proof}
Write $N=\operatorname{End}_C(1)$, $I=C(0,1)$ and
$J=C(1,2)$.  The arrow budget gives $|N|+12\le21$, hence
$|N|\le9$.  Suppose for a contradiction that $BC\simeq\CP^2$.
Collapse the long hom-set to one arrow and apply
Proposition~\ref{C:linear-long-hom-collapse}.  The resulting
triangle-compression monoid is
\[
 M=N\sqcup(I\times J),\qquad BM\simeq\CP^2\vee S^2.
\]
The old copy of $N$ is a \emph{closed} submonoid.  In particular,
an old arrow which acts by a constant transformation pair is still
an old arrow; it is not identified with its corresponding band
element in this construction.

Let $D$ be the image of $N$ in the row-and-column transformation
monoid.  It is a closed submonoid containing the identity, with
$|D|\le9$.  The faithful transformation image of $M$ is
$H=(I\times J)\cup D$.  Both maps preserve the actual sixteen
band elements and their actions.  Thus $M\to H$ preserves
integral second homology and is onto on third homology over every
field by Lemma~\ref{B:band-quotient-H3-surjection}.  In particular,
\[
 b_2(H;\mathbb F_2)=2,\qquad b_3(H;\mathbb F_2)=0.
\]
If $H_4(BH;\mathbb F_2)=0$, every rational degree-two product
on $M$ is zero by pullback, which is impossible.

Use the complete rank-two four-by-four relation-space
classification in Proposition~\ref{B:four-by-four-product25}.
The rank-one-form, repeated-point, zero-fourth-homology and
explicit zero-product cases all give zero products on $M$.
The three zero common-radical cases $161,237,400$ do so as well.
Orbit $346$ also cannot occur: on its complete rational
degree-four cycle space the two degree-two classes have a single
hyperbolic pairing
\[
 \begin{pmatrix}0&-1\\-1&0\end{pmatrix}.
\]
The natural map from $M$ to this hull is an isomorphism on
rational $H^2$.  Pullback to the one-dimensional $H^4(BM;\Q)$
therefore gives either the zero pairing or a nondegenerate
pairing.  The pairing of $\CP^2\vee S^2$ instead has rank one.
The exact full-cycle calculation cited in that proposition
therefore excludes orbit $346$ without a size assumption on its
ambient hull.

It remains to consider the six hulls
\[
 137,169,172,188,222,225.
\]
Their forbidden local-character actors are excluded on the original
finite-dimensional space $BM$.  Every finite cyclic image group
lifts surjectively from an eventual local group of $M$, and its
nonzero integral degree-two character stays nonzero on the lift.
Miller's theorem therefore gives precisely the stated forbidden
actor tests, without a finite-type assertion about $H$ or $D$.

The actual composition $\beta:J\times I\to C(0,2)$ is balanced
over $D$.  If the balanced set $J\times_D I$ has one component,
then $\beta$ is constant, and
Lemma~\ref{C:linear-constant-composition-retract} gives a circle
retract of $BC$.  Consequently this balanced set must have at
least two components.

We enumerate closed middle submonoids $D$ of size at most nine,
including every constant profile they contain, rather than just
the nonband actors of $H$.  The exact rational balancing rank is
seven, so that the degree-two rank is two.  After imposing the
at-least-two-component condition, the complete counts are below.
The columns count closed middle submonoids, distinct faithful
images, and survivors after the image-homology tests.
\[
\begin{array}{c|r|r|r}
 \text{hull}&\text{middle submonoids}&
 \text{images}&\text{survivors}\\\hline
 137&79&35&0\\
 169&1314&468&0\\
 172&1912&559&2\\
 188&79&35&0\\
 222&2&1&0\\
 225&2&1&0.
\end{array}
\]
The image filter retains $b_2=2$, $b_3=0$ and $b_4>0$ over
$\mathbb F_2$.  All excluded cases therefore have one of the
homological obstructions just proved.  No uncomputed homology
tail is used.

The producer \path{code/b_closed_actor_inventory.cpp} uses
canonical closure from the identity, with the forbidden-actor,
size and component tests.  Each component test is monotone:
adjoining actors only identifies additional pairs in the balanced
set.  The independent program
\path{code/d_verify_closed_actor_rank_two.py} instead performs
unrestricted breadth-first adjunction with global duplicate removal.
It reconstructs the literal transformation products, the local
character tests, exact integer relation ranks and every balanced
set.  It reproduces all $3388$ closed middle submonoids.  Its fresh
Python two-sided bar matrices also reproduce every image filter,
checking $945927$ differential-square equations.  The complete
scope record is
\path{code/outputs/d_closed_actor_rank_two_independently_verified.json}.

The only two remaining middle monoids lie in hull $172$, at
candidate indices $338$ and $895$.  Each has nine elements, and
its balanced set has exactly two components.  Since
$N\twoheadrightarrow D$ and $|N|\le9$, the original middle
monoid is isomorphic to this literal $D$; there is no remaining
duplicate-profile inflation.  A nonconstant composition to the
two-element long hom-set is unique up to relabeling that hom-set.
The two resulting twenty-one-arrow category tables are saved in
\path{code/outputs/b_hull172_triangle21_candidates.json}.

Both tables have $b_1=0$, $b_2=1$ over $\mathbb F_2$, and their
degree-two generator has zero square.  The normalized chain
dimensions in degrees one through four are
\[
 (18,144,1152,9216).
\]
For each table a nonboundary two-cocycle and an explicit
three-cochain primitive of its square are supplied.  The
independent checker verifies all $3282$ typed associativity
equations, $2304$ cocycle equations and $18432$ primitive
equations across the two tables.  Its record is
\path{code/outputs/b_hull172_triangle21_cups_independently_verified.json}.
A further independent Python reconstruction produces new
cocycles and primitives in
\path{code/outputs/c_hull172_triangle21_independently_verified.json}.
Either zero-square certificate contradicts the mod-two cohomology
ring of a projective plane.  These two final cases are excluded,
completing the proof.
\end{proof}

\begin{remark}\label{C:closed-middle-two-long-scope}
The closed middle submonoid is essential to the size bound.
Arbitrary monoids with the same small faithful image can have
additional elements, and products of their outside elements may
enter the minimal band.  Here the triangle-compression construction
retains the original closed middle monoid, so its entire
transformation image, including constant profiles, costs at most
nine elements.  The theorem makes no claim for nontrivial endpoint
monoids or for larger adjacent hom-sets.
\end{remark}

\begin{lemma}[Second homology under the actual middle-action quotient]
\label{C:linear-middle-profile-H2}
Let $C$ be a finite category on linearly ordered objects $0,1,2$,
with trivial endpoint endomorphism monoids, nonempty forward
hom-sets $I,J,U$, and no backward arrows.  Replace the middle
monoid $N$ by its joint action image $D$ on $I,J$, retaining
the three hom-sets and their actual composition map.  This gives
a finite category $C_D$ and an actual quotient functor
\[
 F:C\longrightarrow C_D.
\]
For every field $k$, the induced maps on $H_1(-;k)$ and
$H_2(-;k)$ are isomorphisms.  Consequently every degree-two
cohomology class on $BC$ is the pullback of one on $BC_D$,
with its cup products preserved.
\end{lemma}
\begin{proof}
The composition map $J\times I\to U$ is balanced over $D$,
since every element of $D$ has a lift in $N$.  Thus the claimed
category and functor are defined literally, with the functor
equal to the identity on $I,J,U$ and both endpoint identities.
The normalized-chain proof of
Proposition~\ref{C:linear-augmentation-Tor} works over any field.
In its low-degree exact sequence, the augmentation modules of
$kI$ and $kJ$ factor through $kD$, and
\[
 \overline{kJ}\otimes_{kN}\overline{kI}
 =\overline{kJ}\otimes_{kD}\overline{kI}.
\]
The map from this tensor product to $\overline{kU}$ is the same
actual composition map on both sides.  Its kernel and cokernel
are therefore unchanged.  Naturality of that exact sequence gives
the asserted homology isomorphisms.  Duality over $k$ and cup
naturality for the actual functor give the last assertion.
\end{proof}

\begin{theorem}[One long arrow and a small joint middle action]
\label{C:linear-four-by-four-one-long-21}
Let $C$ have three linearly ordered objects with trivial endpoint
endomorphism monoids, no backward arrows, four arrows on each
adjacent side, and exactly one long arrow.  If the joint action
image of its middle monoid on the adjacent hom-sets has at most
ten elements, then $BC\not\simeq\mathbb CP^2$.
In particular no such category with at most twenty-one total
arrows has projective-plane type.
\end{theorem}
\begin{proof}
Let $N$ be the middle monoid and $D$ its joint action image.
Suppose $BC\simeq\CP^2$.  The triangle-compression theorem
identifies $BC$ with the classifying space of
\[
 M=N\sqcup(I\times J),\qquad |I|=|J|=4.
\]
Its faithful transformation image is $H=(I\times J)\cup D$.
The actual map $M\to H$ fixes the band and preserves all its
actions.  Hence it is an isomorphism on second homology and
onto on third homology over $\mathbb F_2$.  We must therefore
have $b_2(H;\mathbb F_2)=1$ and $b_3(H;\mathbb F_2)=0$.
If $H_4(BH;\mathbb F_2)=0$, the square on $M$ vanishes
rationally by pullback, a contradiction.

Proposition~\ref{B:four-by-four-CP2-hulls} reduces the complete
rank-one relation classification to hulls $216,217,221$.
This reduction is not bounded by the order of $M$: its zero
and even-square computations and its integral local-character
restrictions are made on the original finite-CW source.
Delete the forbidden local-character actors as in that proof.
The set $D$ is a closed monoid containing the identity, with
at most ten elements and balancing relation rank eight.
All its constant profiles must be counted as elements of $D$;
they are not free additions from the separate band.

The complete closed-middle enumeration and image filter give
\[
\begin{array}{c|r|r|r}
 \text{hull}&\text{closed middle submonoids}&
 \text{distinct faithful images}&\text{remaining literal categories}\\\hline
 216&8416&3259&56\\
 217&87656&39360&680\\
 221&18906&8142&112.
\end{array}
\]
Every omitted image has wrong second mod-two homology, positive
third mod-two homology, or zero fourth mod-two homology, which
are the obstructions established above.  The independent
breadth-first closure program uses global duplicate removal
and full-round multiplication closure, with no producer
canonical-augmentation rule.  It reconstructs respectively
$24150$, $161269$, and $44269$ closed sets before the final
rank filter, and exactly the candidate sets in the table.
The independent image checker rebuilds the two-sided bars in
a separate basis and checks every source-to-image linkage.
Its complete record is
\path{code/outputs/d_closed_actor_rank_one_independently_verified.json}.

For each of the $848$ surviving closed middle monoids, construct
the literal category $C_D$, retaining every middle element,
including all constant-profile elements, as a distinct old
arrow.  Composition into the singleton long hom-set is forced.
The middle orders are distributed as
\[
\begin{array}{c|rrrr}
 |D|&7&8&9&10\\\hline
 \text{categories}&2&126&226&494.
\end{array}
\]
All $848$ categories have second mod-two cohomology of dimension
one, and a normalized two-cocycle representing its generator
has an explicit three-cochain primitive for its square.
These are certificates on the actual category nerve, rather
than only on a monoid with its constant profiles collapsed.

The source tables and index are
\path{code/outputs/b_closed_rank_one_triangles.tsv} and
\path{code/outputs/b_closed_rank_one_triangles_index.json}.
The independent linkage checker reconstructs every middle
product and both adjacent actions, checking $354388$ literal
table entries.  Its separate normalized-bar checker verifies
$1560960$ typed associativity equations, $1119648$ cocycle
equations, the same number of boundary-square equations, and
$9599436$ square-primitive equations.  No producer bar code is
imported.  The complete record is
\path{code/outputs/b_closed_rank_one_triangles_cups_independently_verified.json}.

Lemma~\ref{C:linear-middle-profile-H2}, with $k=\mathbb F_2$,
now pulls this zero generator square back to the original
category.  This contradicts the projective-plane cohomology
ring even if $N\to D$ has nontrivial fibres.  Finally the
total arrow count is $|N|+11$, so at most twenty-one arrows
implies $|D|\le|N|\le10$, proving the stated cardinality
consequence.
\end{proof}

\subsection{Three and five adjacent arrows with a closed middle action}

\begin{theorem}[One long arrow and a three-by-five joint action]
\label{C:linear-three-by-five-one-long-21}
Let $C$ have three linearly ordered objects, trivial endpoint
endomorphism monoids, no backward arrows, three and five arrows
on its adjacent sides in either order, and exactly one long arrow.
If the joint action image of the middle monoid on the two adjacent
hom-sets has at most ten elements, then $BC\not\simeq\CP^2$.
In particular no such category with at most twenty-one total
arrows has projective-plane type.
\end{theorem}

\begin{proof}
Write $N$ for the middle monoid and $I,J$ for its adjacent left
and right sets.  Taking opposites if necessary, assume
$|I|=3$ and $|J|=5$.  Suppose $BC\simeq\CP^2$.
The triangle-compression theorem gives
\[
 M=N\sqcup(I\times J),\qquad BM\simeq BC.
\]
The copy of $N$ is a closed submonoid, disjoint from the newly
adjoined minimal rectangular band.  Let $D$ be its joint
transformation image and put
\[
 K=I\times J,\qquad H=K\cup D.
\]
The actual quotient $M\to H$ fixes $K$ and preserves its
actions.  Its balanced tensor presentation therefore preserves
integral $H_2$, so $H_2(BH;\Z)=\Z$.  Every degree-two
class on $BM$ is the pullback of one on $BH$.

Apply the complete three-by-five classification in
Theorem~\ref{B:three-by-five-targets}.  Its rank-one-form,
repeated-column, zero-fourth-homology and explicit zero-product
cases exclude projective-plane type by cup naturality.
The repeated-column parity case pulls a primitive degree-two
class back to a primitive class and makes its square even.
All forbidden local-character tests are applied to the original
finite-CW space $BM$: the detected finite cyclic groups have
surjective lifts from eventual local groups of $M$, and Miller's
theorem excludes their nonzero integral restrictions.
These reductions impose no size condition on $M$ and no
finite-type condition on $H$ or $D$.  They leave only the
thirty-element allowed hull in orbit $164$.

In this hull, the fifteen-element band and the identity are
mandatory, and there are fourteen optional nonband profiles.
For any band-containing image $H$, let
\[
 S=H\setminus K,\qquad E=\langle S\rangle
\]
where $E$ is the ordinary submonoid generated by $S$ inside the
transformation hull.  In particular $E$ includes every constant
profile forced by products in $S$.  Since $S\subseteq D$ and
$D$ is closed, $E\subseteq D$ and $|E|\le10$.
Constants cannot be omitted from this count merely because
they also belong to the separate band.

Checking all $2^{14}$ optional subsets reproduces exactly the
$2186$ closed band-containing images of the complete hull
classification.  Exactly $485$ satisfy $|\langle S\rangle|\le10$.
Their orders are distributed as follows:
\[
\begin{array}{c|rrrrrrr}
 |H|&16&17&18&19&20&21&22\\\hline
 \text{number}&1&12&56&134&166&102&14.
\end{array}
\]
Thus $|H|\le22$ is a conclusion of the closed-middle
calculation, not an additional restriction imposed on the source.
Among these $485$ images, exactly $186$ have rational second
homology of rank one; the complete Smith-form gate in
Proposition~\ref{B:three-by-five-finite-target-transfer}
shows that all $186$ have $H_2=\Z$.
Every one is among the previously independently certified
images with $H_4(-;\mathbb F_5)=0$.
Consequently $H_4(BH;\Q)=0$, and the square of its
degree-two class pulls back to zero on $M$, a contradiction.

The exact calculation is saved in
\path{code/b_closed_three_by_five_small_middle.py}
and
\path{code/outputs/b_closed_three_by_five_small_middle_verified.json}.
It reconstructs the literal transformation multiplication,
checks all optional subsets, computes the full ordinary
generated submonoid for every image, and links each of the
$186$ surviving indices to the independently verified Smith
and dense fourth-homology records.  No third-homology
prefilter and no uncomputed higher-homology tail are used.
The separate program
\path{code/d_verify_closed_three_by_five_small_middle.py}
reconstructs the ordinary generated middles by queue closure,
checks all $34^3$ actor associativity products, and repeats
all $16384$ subset tests.  It independently recovers the
$2186$, $485$ and $186$ scopes and checks every index and
input hash in the existing Smith and fourth-homology records.
Its complete output is
\path{code/outputs/d_closed_three_by_five_small_middle_independently_verified.json}.
The four exceptional twenty-two-element orbit-$164$ images
each generate a middle monoid of order seventeen, so none
can satisfy the present bound.  Likewise the two explicit
twenty-three-element projective-plane monoids generate
middle monoids of order eighteen.

Finally, the stated arrow budget is
$|N|+3+5+1+2=|N|+11\le21$.  It implies
$|D|\le|N|\le10$ and gives the last assertion.
\end{proof}

\subsection{The completed fourteen-arrow exclusion}

\begin{lemma}[Semisimple diagonal and the length of the relative bar]
\label{C:semisimple-diagonal-height}
Let $C$ be a finite category whose support on distinct objects is
acyclic.  Suppose that every algebra
$\mathbb Q\operatorname{End}_C(c)$ is semisimple, and that a chain
of nonempty cross-hom sets has length at most $d$.  Then
$H^q(BC;\mathbb Q)=0$ for $q>d$.
\end{lemma}
\begin{proof}
Put $A=\mathbb QC$ and
$B=\bigoplus_c\mathbb Q\operatorname{End}_C(c)$.
The span $J$ of the arrows between distinct objects is a
$B$-bimodule, $A=B\oplus J$, and $J^{\otimes_B n}=0$ for
$n>d$: a nonzero elementary tensor would give a chain of $n$
cross arrows.  Since $B$ is semisimple, the normalized relative
bar resolution over $B$ is a projective $A$-resolution of the
constant module.  Its degree-$n$ term is induced from
$J^{\otimes_B n}$ and is zero above $d$.  Applying the constant
coefficient functor gives the asserted vanishing.
\end{proof}

\begin{theorem}[At least three objects and fourteen arrows]
\label{C:fourteen-many-objects}
A finite category with at least three objects and at most fourteen
arrows cannot have classifying space homotopy equivalent to
$\mathbb {CP}^2$.
\end{theorem}
\begin{proof}
The smaller totals are excluded by Theorem~\ref{C:CP2-lower14}.
We may therefore take the total to be exactly fourteen and the
category to be connected.  If it has $k$ objects, it has at least
$k$ identities and $k-1$ cross arrows.  Thus $k\leq7$.

For three objects, Lemma~\ref{C:fourteen-three-SCC} excludes a
nontrivial strongly connected block, and
Theorem~\ref{C:fourteen-all-stars} excludes the star supports.
The remaining support is a three-object linear order with all
three cross-hom sets nonempty.  Write its endomorphism orders
as $(m,n,p)$ and its cross-hom cardinalities as $(x,y,z)$ in
directions $0\to1$, $1\to2$, $0\to2$.  Direct enumeration of
six positive integers with sum fourteen gives $1287$ tuples.
Of these, $285$ satisfy $x=z=1$ or $y=z=1$ and are excluded by
the singleton-extreme join of
Lemma~\ref{C:singleton-cross-join}.  Of the remaining tuples,
$232$ satisfy $y=p=1$ or $x=m=1$; deleting the adjacent
trivial extreme object is the singleton cofinal deletion.
The other $770$ tuples give exactly $406$ cardinality profiles
up to object relabeling and opposition.  Every one of these
$406$ profiles, and no further linear profile, occurs as a
completed profile in the final three-object run.

For four objects, Proposition~\ref{C:fifteen-four-large-corner}
excludes every corner of order at least eight, already under the
larger budget of fifteen.  For at least five objects the elementary
bound
\[
 |\operatorname{End}(c)|\leq14-2(k-1)
\]
makes every corner have order at most six.  These observations
justify the corner bounds in the finite enumerations below.

Here is the precise table-generation procedure.  Identities are
the first $k$ arrow labels, and all other labels are grouped by
their specified source and target.  Every connected nonnegative
cross-hom cardinality distribution and every positive diagonal
distribution with total fourteen is considered, up to object
relabeling and opposition.  The identity products are prescribed.
For each remaining composable pair, its product ranges over the
entire hom-set having the required endpoints.  Associativity
propagation discards an assignment only when a typed triple has
inconsistent products.  At a completed assignment every typed
triple is checked.  Initial or terminal objects, retracts,
singleton cofinal deletions, and singleton-extreme joins are
discarded by the homotopy and cup-product reductions already
proved above.  In the three-object run the strongly connected
and star supports are discarded only by the two theorems cited
at the start of this proof.

The three- and four-object runs fill the diagonal corners from
the complete ordinary isomorphism-class catalogs through order
seven, restoring both orientations when the catalog was stored
up to isomorphism and opposition.  In particular, the order-seven
catalog consists of $17730$ isomorphism/opposition classes and
accounts for all $18507440$ labelled monoids.  The local arrow
indices in every corner product are converted to the corresponding
global category indices.  This is a finite replacement of repeated
diagonal searches, and does not restrict the cross products.
The fast output reduction identifies only permutations within
cross-hom sets; retaining different corner labelings can therefore
produce extra representatives but cannot remove a category.
The runs with five through seven objects use the full typed
product recursion directly.

One additional discard in the three- and four-object runs is
certified independently of a homology computation.  For each
corner its regular trace matrix is
\[
 \bigl(\operatorname{Tr}(L_{ab})\bigr)_{a,b}.
\]
Full rank modulo $32003$ proves full rank over $\mathbb Q$.
The radical of a finite-dimensional characteristic-zero algebra
is contained in the kernel of this form, since every product
of a radical element with another element is nilpotent and its
left multiplication has trace zero.  Thus this test proves
semisimplicity.  It is used only for acyclic support.  When it
succeeds at every corner, Lemma~\ref{C:semisimple-diagonal-height}
gives $H^4(BC;\mathbb Q)=0$, because the support has at most
four objects.  A failed modular test makes no assertion and
does not discard the table.

Every surviving table with $H_1(-;\mathbb F_2)=0$ and
$\dim H_2(-;\mathbb F_2)=1$ is written to a complete literal
table stream, modulo the stated hom-set permutations.  The
separate checker reconstructs all normalized bar differentials,
checks every typed associativity equation, and verifies a
nonboundary degree-two cocycle.  It either checks an explicit
degree-three cochain whose coboundary is the square, or verifies
that the nonzero-square table has nonzero third mod-two homology.
The complete results are
\[
\begin{array}{c|r|r|r|r}
\text{objects}&\text{completed profiles}&\text{table rows}
 &x^2=0&x^2\ne0,\ H_3\ne0\\\hline
3&406&504126&495104&9022\\
4&2367&192952&192863&89\\
5&1131&4472&4470&2\\
6&130&56&56&0\\
7&1&0&0&0\\\hline
\text{total}&&701606&692493&9113.
\end{array}
\]
The independent checks include $161927064$ typed associativity
equations.  A projective-plane class has nonzero mod-two square
and zero third mod-two homology, so every listed table is excluded.
No assertion about an uncomputed homological tail is used.

The three exact generators, all in \path{code/}, are
\begin{quote}\small
\path{c_categories14_modular_fast.cpp},\\
\path{c_categories14_modular4_fast.cpp},\\
\path{c_categories14_minimal.cpp}.
\end{quote}
The independent bar checker is
\path{code/c_generic_hom_cups.py}.
Their complete output stems in \path{code/outputs/} are
\path{c_categories_14_3_modular_fast_CP2},
\path{c_categories_14_4_modular_fast_CP2}, and
\path{c_categories_14_5_minimal_CP2} through
\path{c_categories_14_7_minimal_CP2}.
For each stem the table stream, completed-profile log, generator
summary, and all-row cup certificates are retained.
The separate coverage program
\path{code/d_fourteen_category_coverage.py} regenerates the
$1287$ three-object tuples directly, verifies the $406$ required
profile orbits, and checks every final completion sentinel,
table-row count, and exact byte position against the independent
checker.  Its complete record
\path{code/outputs/d_fourteen_category_coverage.json} has status
\texttt{PASS} and no pending scope.  Thus no unfinished stream
or unverified cardinality profile enters the conclusion.
\end{proof}

\begin{theorem}[The global lower bound of fifteen arrows]
\label{C:CP2-lower15}
Every finite category whose classifying space is homotopy
equivalent to $\mathbb {CP}^2$ has at least fifteen arrows,
identities counted.  Together with the explicit upper model,
the current minimum-arrow interval is $15$--$22$.
\end{theorem}
\begin{proof}
For one object, Theorem~\ref{B:CP2-minimum23} gives the stronger
lower bound twenty-three.  A connected two-object category is
either strongly connected or triangular.  The former is excluded
through twenty-two arrows by
Theorem~\ref{C:twentytwo-strongly-connected}; the latter is
excluded through fourteen by
Proposition~\ref{C:fourteen-triangular}, including its completed
independent commuting-biset coverage and cochain verifications.
Theorem~\ref{C:fourteen-many-objects} handles every category with
at least three objects.  These cases exhaust all categories with
at most fourteen arrows.  The twenty-two-arrow construction is
Theorem~\ref{C:CP2-upper22}.
\end{proof}

\subsection{The completed fifteen-arrow exclusion}

\begin{theorem}[At least three objects and fifteen arrows]
\label{C:fifteen-many-objects}
A finite category with at least three objects and at most fifteen
arrows cannot have classifying space homotopy equivalent to
$\mathbb {CP}^2$.
\end{theorem}
\begin{proof}
Theorem~\ref{C:CP2-lower15} handles smaller totals.  At total
fifteen the category is connected, and
Lemma~\ref{C:connected-tree-budget} reduces the object count to
three through seven.  The equality case at eight objects is
a tree and is contractible, so no eight-object search is required.

For three objects, Proposition~\ref{C:fifteen-three-SCC} excludes
nontrivial strongly connected blocks, and
Theorem~\ref{C:fifteen-all-stars} excludes the two star supports.
The remaining support is a three-object chain with every forward
hom-set nonempty.  Write its corner orders as $(m,n,p)$ and its
cross-hom cardinalities as $(x,y,z)$ in directions $0\to1$,
$1\to2$, and $0\to2$.  There are $2002$ positive six-tuples
with sum fifteen.  In the following disjoint order, the proved
geometric reductions remove:
\[
\begin{array}{c|r}
\text{reduction}&\text{six-tuples}\\\hline
\text{singleton extreme join }(x=z=1\text{ or }y=z=1)&385\\
\text{singleton cofinal deletion }(y=p=1\text{ or }x=m=1)&321\\
\text{fixed extreme arrow, Lemma~\ref{C:fixed-incoming-cups}}&148\\
\text{triangle compression, Corollary~\ref{C:triangle-band-CP2-size}}&21.
\end{array}
\]
For clarity, the fixed-arrow test is
\[
 \bigl((n=x=1\text{ or }m=n=1)\text{ and }(y=1\text{ or }z=1)\bigr)
\]
or its opposite.  The lower two-object subcategory in the first
case has zero rational second cohomology, and one of its arrows
into the extreme object is fixed by every endomorphism there.
The triangle test is $m=p=z=1$ and $n+xy<23$; the compressed
monoid would contradict Theorem~\ref{B:CP2-minimum23}.
The remaining $1127$ tuples have corner orders at most seven
and give exactly $588$ profiles up to object permutation and
opposition.  The final generator's completed profile set is
exactly this set, not merely a set of the same cardinality.

For four objects, Proposition~\ref{C:fifteen-four-large-corner}
excludes corners of order at least eight.  For $k\geq5$ the
connected support budget gives
\[
 |\operatorname{End}(c)|\leq15-2(k-1),
\]
so the corner orders are at most seven, five, and three for
$k=5,6,7$.  The three-, four-, and five-object generators use
every ordinary corner isomorphism class through order seven,
restoring both orientations of an isomorphism/opposition
catalog.  Every local product index is converted to its global
category arrow index.  The six- and seven-object generators
instead enumerate every typed product directly.

In all cases identities are the first $k$ arrow labels, and
the remaining labels are grouped by their sources and targets.
Every remaining product ranges over the entire hom-set with
the required endpoints.  Associativity propagation removes
an assignment only when a typed triple is inconsistent.
The generic initial/terminal, retract, singleton cofinal, and
singleton-join reductions are the same proved reductions used
in Theorem~\ref{C:fourteen-many-objects}; every object deletion
leaves a smaller category covered by that theorem.  No
unproved minimality assumption is used.

The modular corner generators also discard an acyclic support
when every corner regular trace form is nonsingular modulo
$32003$.  As in the preceding fourteen-arrow proof, this
certifies semisimplicity over $\mathbb Q$.  The three-object
relative bar then has length at most two; for four or five
objects use Lemma~\ref{C:semisimple-fifteen}.  In each case
$H_4(-;\mathbb Q)=0$.  This discard is not applied to cyclic
support.  A failed modular test retains the corner choice.

The original four-object stream contains $2114185$ literal
table rows.  Proposition~\ref{C:fifteen-four-star-prunes}
removes $287757$ rows with two singleton legs, $78195$ rows
with small leaves and a fixed arrow, and $8150$ rows by exact
leaf coinvariants.  The remaining $1740083$ rows are checked
individually.  An independent implementation reconstructs
the undirected star support, checks its orientation and
cardinality conditions, and compares every retained row
literally with the original stream, in its original order.
Thus this postprocessing does not silently change the
enumerated scope.

Every remaining table with $H_1(-;\mathbb F_2)=0$ and
$\dim H_2(-;\mathbb F_2)=1$ is written to a literal table
stream.  The independent checkers reconstruct the normalized
bar complex, check all typed associativity equations, and
verify a nonboundary degree-two cocycle.  They check either
an explicit degree-three primitive for its square or positive
third mod-two homology.  The completed results are
\[
\begin{array}{c|r|r|r|r}
\text{objects}&\text{completed typed profiles}&\text{checked rows}
 &x^2=0&x^2\ne0,\ H_3\ne0\\\hline
3&588&3105576&3029149&76427\\
4&5336&1740083&1738853&1230\\
5&3839&30649&30642&7\\
6&724&725&725&0\\
7&22&5&5&0\\\hline
\text{total}&&4877038&4799374&77664.
\end{array}
\]
The profile column counts completed \texttt{START}/\texttt{DONE}
pairs before the four-object star filter.  In particular, the
three-object generator's $1183$ attempted distributions consist
of these $588$ completed profiles and $595$ impossible profiles
rejected earlier because a required composite hom-set is empty.
The two counts are not interchangeable.

The checks include $2055830258$ typed associativity equations.
For the two largest streams, the completed independent C++
certificates are linked to every source row.  The three-object
stream is covered by the disjoint half-open row intervals
\[
 [0,1666122),\qquad[1666122,2479174),\qquad
 [2479174,3105576).
\]
Their union is the entire completed stream, with no overlap
or gap.  The four-object certificate similarly covers the
entire filtered stream.  The earlier unfinished Python checker
prefixes are retained as historical artifacts and are not used
as completion evidence.  For five through seven objects the
independent Python checkers reached the exact final byte of
each completed stream.

The exact generators are
\begin{quote}\small
\path{code/c_categories15_linear_modular.cpp},\\
\path{code/c_categories15_modular4.cpp},\\
\path{code/c_categories15_modular5.cpp},\\
\path{code/c_categories15_high_objects.cpp}.
\end{quote}
The complete output stems in \path{code/outputs/} are
\path{c_categories_15_3_linear_modular_CP2},
\path{c_categories_15_4_modular_CP2} and its
\path{c_categories_15_4_star_reduced_CP2} substream,
\path{c_categories_15_5_modular_CP2}, and
\path{c_categories_15_6_reduced_CP2} through
\path{c_categories_15_7_reduced_CP2}.
The large-stream linkage records are
\path{b_category_tsv_15_3_complete_verified.json} and
\path{b_category_tsv_15_4_coverage_verified.json}; the
independent star-filter record is
\path{d_fifteen_four_star_filter_verified.json}.

Finally, \path{code/d_fifteen_category_coverage.py}
independently regenerates the $2002$ positive six-tuples and
their exact $588$ required profile orbits, reads every final
completion sentinel, verifies all row intervals and counts,
and checks the exact table-stream hashes and final byte
positions.  Its final record
\path{code/outputs/d_fifteen_category_coverage.json} has status
\texttt{PASS} and no pending scope.  Every possible projective
plane would have nonzero mod-two square and zero third
mod-two homology, so all rows are excluded.  No claim about
an uncomputed homological tail enters this conclusion.
\end{proof}

\begin{theorem}[The global lower bound of sixteen arrows]
\label{C:CP2-lower16}
Every finite category whose classifying space is homotopy
equivalent to $\mathbb {CP}^2$ has at least sixteen arrows,
identities counted.  Together with the explicit construction,
the minimum-arrow interval is $16$--$22$.
\end{theorem}
\begin{proof}
Corollary~\ref{C:fifteen-two-objects} excludes categories with
at most two objects and at most fifteen arrows.  Theorem
\ref{C:fifteen-many-objects} excludes every larger object
count.  The upper bound is Theorem~\ref{C:CP2-upper22}.
\end{proof}

\subsection{Act cones with one total row excess and the sixteen-arrow SCC case}

\begin{theorem}[One total row excess in a finite categorical base]
\label{C:one-row-excess-act-cups}
Let $S$ be a finite strongly connected category.  In its minimal
Rees ideal, let $r_c$ be the number of rows at object $c$.
Suppose
\[
 \sum_{c\in\operatorname{Ob}S}(r_c-1)=1.
\]
For every nonempty finite right $S$-act $F$, all rational
products of degree-two classes on the act cone
$B\mathcal C(S,F)$ vanish.  Here the act cone adds one sink
object, with hom-set $F(c)$ from $c$ to that sink.
\end{theorem}
\begin{proof}
Write $R=\mathbb QS$, and let $\mathbf Q$ denote the constant
module, with one copy of $\mathbb Q$ at every object.  Choose
an idempotent in the minimal ideal and average its finite
maximal group $H$:
\[
 p=|H|^{-1}\sum_{h\in H}h\in R.
\]
The module $P=Rp$ is projective.  At object $c$ it has one
basis vector for each of the $r_c$ minimal rows: multiplying
an arrow by $p$ averages its free $H$-coordinate and fixes
its minimal-ideal column.  Sending each of these basis vectors
to $1$ gives a natural surjection $P\to\mathbf Q$.
Its kernel $V$ is the objectwise row augmentation module.
The assumed total excess means that $V$ is one-dimensional
at a unique object $a$ and zero at all other objects.
Endomorphisms of $a$ act by a scalar character $\chi$.
Since its image is a finite multiplicative subset of
$\mathbb Q$, its values lie in $\{0,1,-1\}$.

Every set $F(c)$ is nonempty: a nonempty stalk propagates
along arrows in both directions in the strongly connected
base.  Let $L$ be the kernel of the objectwise augmentation
$\mathbb QF\to\mathbf Q$ of right modules, and let $D L$
denote its objectwise rational dual.  The act-cone bar and
vector-space duality give
\[
 H^2(B\mathcal C(S,F);\mathbb Q)
 \cong\operatorname{Ext}^1_R(\mathbf Q,D L).
\]
Applying $\operatorname{Hom}_R(-,D L)$ to
$0\to V\to P\to\mathbf Q\to0$ gives the natural
surjection
\[
 \operatorname{Hom}_R(V,D L)\longrightarrow
 \operatorname{Ext}^1_R(\mathbf Q,D L).
\]
Thus every degree-two class has a representative given by
a right-module map $\ell:L\to D V$.

Choose a base state $s_c\in F(c)$ and put
\[
 q_c(s)=\ell_c(e_s-e_{s_c}).
\]
This is zero off $a$ and is a rational-valued function at $a$.
For an arrow $f:c\to d$, write $s f$ for the right action
$F(d)\to F(c)$.  Naturality says that
\[
 q_c(s f)=D V(f)q_d(s)+q_c(s_d f).
\]
Consequently the image sets $T(c)=q_c(F(c))$ form an actual
finite quotient right act.  Off $a$ they are singletons.
Every cross-arrow map between them is constant, because
$D V(f)=0$ for distinct source and target objects.  At $a$
the maps have the affine form $x\mapsto\chi(f)x+b_f$.
If $T(a)$ is a singleton, the representing functional and
the class are zero.  Otherwise a slope-one map of the finite
set is the identity, a slope-minus-one map is its common
reflection, and a slope-zero map is constant.

The functional $\ell$ factors through the augmentation
module of $T$.  Naturality of the preceding connecting
surjection proves that the original class is pulled back
from the act cone on $T$ under the quotient functor
$\mathcal C(S,F)\to\mathcal C(S,T)$.

Extend $T(a)$ to its convex interval $I(a)$ and put $I(c)$
equal to a point for $c\ne a$.  Extend the maps affinely.
They satisfy the category composition equations on these
intervals because they agree on the spanning finite sets;
the singleton cases are immediate.  This is a strict
contravariant diagram of contractible spaces.  Its quotient
diagram $I(c)/T(c)$ has a wedge of circles at $a$ and a point
elsewhere.  Every cross map is a pointed constant map.

Taking degree two on every circle commutes with every
endomorphism at $a$: identity, common reflection, and
constants.  It also commutes with all cross maps.  Hence it
is a strict self-map of this pointed diagram.  The pointed
homotopy colimit of $I/T$ is the cofiber of
$\operatorname{hocolim}T\to\operatorname{hocolim}I$, and
$\operatorname{hocolim}I\to BS$ is an equivalence by
objectwise contractibility.  This is the required act cone.
The reduced cellular complex of the quotient diagram is
concentrated in degree one, so the induced map of its bar
complex is multiplication by two in every degree.  It acts
by two on every positive rational cohomology group of the
act cone.  Multiplicativity therefore gives
$2xy=4xy$, forcing every positive-degree rational product
on this quotient act cone to vanish.

Pulling back kills the square of the original arbitrary
degree-two class.  Polarization over $\mathbb Q$ proves
the assertion for every pair of degree-two classes on the
original cone.  This proof allows a nontrivial minimal group:
its only use was the projective averaging idempotent $p$.
\end{proof}

\begin{theorem}[Three-object SCC exclusion at sixteen arrows]
\label{C:sixteen-three-SCC}
A three-object finite category with at most sixteen arrows
and a nontrivial strongly connected component cannot have
classifying space homotopy equivalent to $\mathbb {CP}^2$.
\end{theorem}
\begin{proof}
The smaller totals have already been excluded.  If one object
is a retract of another, delete it; at least its identity
and two opposite cross arrows disappear, leaving at most
thirteen arrows.  We may assume there is no such retract.
In any nontrivial strongly connected block, this implies
that every object identity lies outside the minimal ideal.
Indeed, an identity in that ideal makes its corner a group;
composing any two opposite cross arrows and adjusting by
the inverse of their corner product exhibits a retract.

First suppose all three objects are strongly connected.
If the total number of rows or the total number of columns
is three, the corresponding averaged minimal representable
is the constant module, which is projective.  The category
is rationally acyclic and cannot have the proposed type.
Otherwise both totals are at least four.  The minimal ideal
then has at least sixteen arrows, in addition to the three
identities outside it, exceeding the budget.

There remains a strong two-object block $S$, a singleton
monoid $M$, and cross arrows in one direction.  Write $s,m,t$
for their sizes.  Here $s\geq6$ and $t\geq2$.
If $t=2$, the two cross-hom sets are singletons, so the
category is a categorical join.  Its positive rational
products vanish.  We may take $t\geq3$ and $s+m\leq13$,
and use opposition to make $S$ the source block.

If both total row and total column counts of $S$ are at
least three, its minimal ideal has at least nine arrows;
the two identities give $s\geq11$.  Hence $s\leq12$ and
$m\leq2$.  Its minimal group is trivial and both total
counts are exactly three.  Its total row excess is one.
The trivial rational modules of every monoid of order at
most two are projective.  The actual orbit-collapse functor
at the $M$ vertex therefore gives a rational cohomology-ring
equivalence to a finite right-$S$ act cone.  Apply
Theorem~\ref{C:one-row-excess-act-cups}.  This includes the
twelve-arrow strong block with an extra arrow outside its
minimal band; no deletion or classification of that arrow
is needed.

In the remaining case the total row or column count of
$S$ is two, and $BS$ is rationally acyclic by the projective
constant-module argument.  If $m\leq4$, both vertex blocks
are rationally acyclic.  Collapsing them in the correspondence
homotopy pushout gives a rational ring equivalence to a
suspension, so all products vanish.

Suppose $m\geq5$.  Then $s\leq8$.  The minimal group of
$S$ is trivial, and its total row and column counts are
either $(2,2)$, $(2,3)$, or $(3,2)$.  For $(2,2)$ the
constant modules on both sides are projective.  Collapsing
$S$ by actual cross-act coinvariants therefore gives a
rational ring equivalence to a finite act cone over $M$,
where $m\leq7$.  Its rational homology would be bounded
and its second cohomology one-dimensional; this contradicts
Lemma~\ref{C:nine-bounded-rank-one}.

For $(2,3)$ or $(3,2)$, the six-element band and the two
identities exhaust $s=8$.  Choose an object $b$ with one
minimal row, and another object $a$ in the block.  Any
minimal idempotent $e\in\operatorname{End}(b)$ fixes every
incoming cross arrow at $b$: all such arrows belong to the
minimal band and their target row is unique.  There are
no incoming arrows from outside this source block.  The
only unit of the corner is its identity.  Minimal-band
arrows $u:b\to a$ and $v:a\to b$ can be chosen with
$vu=e$.  Lemma~\ref{C:one-sided-corner-deletion} deletes
$b$, leaving a smaller category already excluded by
Theorem~\ref{C:CP2-lower16}.  This proves the last case.
\end{proof}

\subsection{Large corners with four objects and sixteen arrows}

\begin{lemma}\label{C:three-small-H2-zero}
A connected three-object category with at most seven arrows has
$H^2(BD;\mathbb Q)=0$. The same holds without connectedness at
at most six arrows.
\end{lemma}
\begin{proof}
In the connected case a corner has order at most three, since
there are two other identities and at least two cross arrows.
A retract deletion leaves at most four arrows, whose rational
$H^2$ is zero by the complete small-category census. Assume
no retract is present. A nontrivial strong pair without
such a retract requires two corners of order at least two,
two opposite cross arrows, and one further connecting cross
arrow, hence at least eight arrows. Thus the support is acyclic.
If a corner has order three, the remaining two cross arrows
are singletons and the other corners are trivial. There is
no composable pair of cross arrows, since that would require
a third cross hom-set. The nerve is a graph of vertex spaces
and singleton correspondence spaces. All vertex monoids have
rationally acyclic classifying spaces, by the complete order-four
classification; a singleton correspondence is their product.
The graph is a tree, so Mayer--Vietoris gives rational acyclicity.

Otherwise all corner algebras are rationally semisimple. If the
support has no three-object chain, its relative bar has length
at most one. If it has such a chain, write its three hom-set
sizes as $x,y,z$, with $z$ the long edge. If some corner has
order two, all three hom-sets are singletons, and one of the
trivial extreme corners is initial or terminal. If all corners
are trivial, $x+y+z\leq4$. The possibilities $x=2,y=z=1$
and $y=2,x=z=1$ have a terminal or initial object. In the other
possibilities $x=y=1$, the single two-simplex has a nonzero
boundary, so $H^2=0$.

For a disconnected three-object category with at most six
arrows, an isolated component monoid has order at most four.
A connected two-object component has at most five arrows. A nontrivial strong component either has a trivial
corner, which is a retract, or needs at least six arrows.
A triangular component with nontrivial corner and two cross
arrows has corner orders $(2,1)$; the order-two coinvariant
functor is exact and its orbit-collapse gives a graph.
All remaining triangular components have singleton cross hom
or both corners trivial. A singleton gives the join of the
vertex spaces, which is rationally acyclic here, and trivial
corners give a graph. These cases prove the assertion.
\end{proof}

\begin{proposition}\label{C:sixteen-four-large-corner}
A four-object category with at most sixteen arrows and an
endomorphism monoid of order at least eight cannot have type
$\mathbb{CP}^2$.
\end{proposition}
\begin{proof}
Assume the category is connected, and denote the large corner
by $v$, its order by $q\geq8$. A retract deletion leaves at
most thirteen arrows, already excluded. We may assume there
is no retract. Connectivity bounds $q\leq10$.

First exclude a nontrivial strong component. Every object in
one has a nontrivial corner, since otherwise it is a retract
of another object. The only possible budget is a strong pair
containing $v$, with corner orders $8,2$, two other trivial
corners, and exactly four cross arrows. Every nonempty cross
hom-set is a singleton. At the small corner the composite of
the opposite cross arrows is its nonidentity zero idempotent.
It fixes every incoming arrow. Lemma~\ref{C:one-sided-corner-deletion}
deletes that object and again reduces to at most thirteen arrows.
Thus the support can be assumed acyclic.

Suppose $v$ has both incoming and outgoing cross arrows. A
three-object chain through $v$ and its composite, together
with a connection to the fourth object, require at least
four cross arrows. If there are exactly four, the total
incoming and outgoing counts at $v$ must each be one:
a sum at least three requires at least five cross arrows,
whether the three incident arrows meet two or three neighbors.
The two adjacent corners are distinct, and at most one can
have order two; the other is trivial. Thus one of the two
singleton cofinal deletions deletes $v$.
If there are five cross arrows, $q=8$ and every other corner
is trivial. Both the incoming and outgoing total cannot be
at least two. With only two neighbors this would require
four incident arrows, a composite, and a further connecting
arrow; with three neighbors it would require four incident
arrows and at least two different composite hom-sets.
Either way it requires at least six arrows. Hence one of the
incoming and outgoing totals is one, and the adjacent corner
is trivial. The same cofinal deletion applies. Its result
has at most seven arrows. We may therefore assume that $v$
is extreme, and use opposition to make it a sink.

Let $D=C\setminus\{v\}$. If any incoming hom-set at $v$ is a
singleton, then $D$ has at most seven arrows. If it has seven,
the budget forces $q=8$ and exactly one cross arrow incident
to $v$, so connectivity forces $D$ to be connected. If it has
at most six, connectedness is unnecessary.
Lemma~\ref{C:three-small-H2-zero} gives $H^2(BD;\mathbb Q)=0$.
The singleton incoming arrow is fixed by every element of the
large monoid. Lemma~\ref{C:fixed-incoming-cups} therefore kills
all rational degree-two products.

It remains to suppose every nonempty incoming hom-set has at
least two elements. If there are at least two neighbors of
$v$, there are exactly two, each with two arrows to $v$;
$q=8$, all other corners are trivial, and one singleton cross
arrow connects the fourth object $w$ to one of these neighbors
$x$. It must be $x\to w$: the other orientation would force
a third incoming hom-set to $v$. Its sole incoming arrow
comes from the trivial corner $x$, so the singleton cofinal
deletion removes $w$, leaving at most fourteen arrows.

Finally suppose there is just one neighbor $w$, with
$s=|\operatorname{Hom}(w,v)|\geq2$. There are no cross arrows
into $w$ from $D\setminus\{w\}$, because their composites
would give another neighbor of $v$. Thus the nerve is the
union of $ND$ and the full two-object nerve $NY$ on $w,v$,
intersecting in $N\End(w)$; this is a homotopy pushout of
subcomplexes. The category $D$ is connected. Its two necessary
cross arrows give $s\leq3$ and $|\End(w)|\leq2$.
For $s=2$, $D$ has at most six arrows, so its rational $H^2$
is zero. The common corner is rationally acyclic. The
Mayer--Vietoris sequence consequently identifies $H^2(BC;\mathbb Q)$
with $H^2(BY;\mathbb Q)$ and injects $H^4(BC;\mathbb Q)$
into $H^4(BD;\mathbb Q)\oplus H^4(BY;\mathbb Q)$.
Lemma~\ref{C:two-object-cup-pruning} kills the degree-two
products on $BY$. Their restrictions on $BD$ are also zero,
so their products on $BC$ vanish.
For $s=3$, the entire budget forces $q=8$, every corner of
$D$ trivial, and exactly two singleton cross arrows in $D$.
Its support is a tree, so $BD$ is contractible, and the
pushout gives $BC\simeq BY$. This is a three-state act cone
with twelve arrows, excluded by
Theorem~\ref{C:three-state-cone-minimum22}.
Every case contradicts the primitive nonzero projective-plane
square or an already proved smaller-category exclusion.
\end{proof}

% Full independent literal audits by root and B passed; source note retained.
\subsection{Large corners with four objects and seventeen arrows}

\begin{lemma}[Fourth homology at excess three]
\label{C:arrow-excess-three-H4}
If a connected finite category has $k$ objects and at most
$2k+2$ arrows, then its fourth rational homology is zero.
\end{lemma}
\begin{proof}
Delete retract objects. Each deletion removes at least three
arrows, preserves homotopy type, and preserves the stated bound.
In the resulting category, a proper strong component on $t\geq2$
objects has at least $t^2+t$ internal arrows, since all its
identities lie outside its minimal ideal. Connectivity outside
that component gives a total of at least
$2k+t^2-1\geq2k+3$, a contradiction. If the whole category is
strongly connected, it has either one object and at most four
arrows, or two objects and exactly six arrows. The first case
is rationally acyclic by the small-monoid classification.
In the second case the minimal band has one row and one column
at each object and trivial maximal group. The minimal arrows
from a chosen object give a natural contraction, as in the
proof of Lemma~\ref{C:one-row-source-block-collapse}.

It remains to consider acyclic support. In the notation used
for the excess bounds, $D+P+R\leq3$. If the diagonal algebra
is semisimple, a strict chain has at most four objects, because
a five-object chain has support cycle rank six. The relative
bar therefore has length at most three and $H_4=0$.
Otherwise $D\geq2$. Every corner has order at most four and
is rationally acyclic. If $P=0$, the internal bars for a strict
chain have rational homology only in degree zero. Thus the
actual support functor is a rational homology equivalence;
the support has no four-simplex and its fourth homology is zero.
The sole remaining possibility is $D=2$, $P=1$, $R=0$.
There is one corner of order three, all other corners are
trivial, and the support is a tree with one doubled hom-set.
The trivial modules of every order-three monoid have rational
projective dimension at most one, by the explicit order-three
algebra list in Proposition~\ref{C:semisimple-ends-small-middle}:
the identity-adjoined two-element one-sided zero bands have
a projective one-term augmentation kernel on the nonprojective
side, and the other cases have projective trivial modules.
Consequently the only
possibly nontrivial edge collage, an ordinary two-state act
cone at that corner, has zero homology in degrees at least
three. The other edge collages are graphs or singleton joins.
Their corner intersections are rationally acyclic, so
Mayer--Vietoris gives zero fourth homology for their union.
\end{proof}

\begin{lemma}[A singleton arrow with a small outside budget]
\label{C:singleton-extreme-three-excess}
Let $C$ be connected and have a sink object $v$. Write $M$
for its corner, $D=C\setminus\{v\}$, and
$\ell=|\operatorname{Ob}D|\geq1$. Suppose that some incoming
hom-set at $v$ is a singleton and that
\[
 |D|+\sum_{x\in D}|C(x,v)|\leq2\ell+3.
\]
Then all rational degree-two products on $BC$ vanish.
The opposite assertion holds for a source object.
\end{lemma}
\begin{proof}
If $H^2(BD;\mathbb Q)=0$, the singleton fixed-arrow lemma
applies directly. Suppose instead this group is nonzero,
and let $D_1,\ldots,D_c$ be the connected components of $D$.
At least one component on $t$ objects has at least $2t+2$
arrows by Lemma~\ref{C:arrow-excess-two-H2}; each other
component has at least $2t-1$ arrows. Hence
\[
 |D|\geq2\ell-c+3.
\]
Connectivity requires at least one incoming arrow from each
component. The assumed budget forces equality throughout:
there is precisely one arrow from each component to $v$,
and the sum of the component excesses is exactly three.
Each component therefore satisfies the preceding lemma and
has zero fourth rational homology.

Let $w_i$ be the unique object of $D_i$ with an arrow to $v$.
It has no incoming cross arrow in $D_i$, since composing one
with $w_i\to v$ would give an additional arrow into $v$.
Its corner has order at most four: connectivity in $D_i$
and excess at most three leave at most three extra corner
arrows. In particular that corner is rationally acyclic.
Let $Z$ be the full star on $v,w_1,\ldots,w_c$. Its legs are
singletons. Each branch is the join of $BM$ with a rationally
acyclic connected corner space, so is rationally acyclic.
The branches meet along $BM$. Their Mayer--Vietoris sequence
gives $H^2(BZ;\mathbb Q)=0$, since $H^1(BM;\mathbb Q)=0$.

The nerve of $C$ is the actual union of this star and the
component nerves $ND_i$, meeting at their corner nerves
$N\operatorname{End}(w_i)$. Indeed, a chain using $v$ can
use no other object before $w_i$. These intersections are
rationally acyclic. Fourth-cohomology restriction to all
these pieces is injective. On $BZ$ a degree-two class is
zero, and on each $BD_i$ its square is zero because $H^4=0$.
The global square, and hence all mixed degree-two products
by polarization, is zero. Opposition proves the dual claim.
\end{proof}

\begin{proposition}\label{C:seventeen-four-large-corner}
A four-object category with at most seventeen arrows and an
endomorphism monoid of order at least eight cannot have type
$\mathbb CP^2$.
\end{proposition}
\begin{proof}
Assume it has that type and is connected. A retract deletion
leaves at most fourteen arrows, already excluded. Thus no
object is a retract of another. Let $v$ have corner $M$ of
order $q\geq8$.

First exclude a nontrivial strong component. One not containing
$v$ has at least six internal arrows. The other two corners
and the necessary interfaces give at least
$6+8+1+3=18$ arrows. A strong component of at least three
objects containing $v$ requires at least
$8+2+2+6+1+3=22$ arrows if proper, and still more if it is
the whole category. A strong pair containing $v$ has at least
$q+4$ internal arrows, and its two-object complement and
interfaces require at least five more. The only possible
case is $q=8$, internal size twelve, and total size seventeen.
It is extreme: a nonextreme pair needs two expanded interfaces
and their composite, giving at least seven outside arrows.
Its other corner is $\{1,e\}$, both internal cross hom-sets
are singletons, and their composite at that corner is $e$.
If the pair is a source, all incoming arrows at the small
corner are internal and fixed by $e$; if a sink, use outgoing
arrows. The one-sided corner-deletion lemma applies, leaving
a smaller forbidden category. We may assume acyclic support.

Let $d$ be the sum of the corner excesses away from $v$.
With $P,R$ denoting parallel excess and support cycle rank,
the arrow budget is
\[
 d+P+R\leq11-q\leq3.
\]
If $P=0$, every degree-two class pulls back from the support
poset, whose fourth homology is zero. We may assume $P\geq1$.

\paragraph{A nonextreme large corner.}
Suppose first that $v$ has both entrances and exits.
Then $R\geq1$, and hence $d\leq1$. If $d=1$, necessarily
$q=8$ and $P=R=1$. The support is one triangle through
$v$ and one leaf edge. The large corner is the middle of
the triangle. Its fourth object can attach only at one of
the two extremes and away from that extreme: any attachment
through the middle or toward the triangle would force another
comparison and a second cycle. Therefore the single order-two
corner, wherever it lies, is extreme in the whole category.
Its exact orbit quotient makes it trivial, preserving rational
cohomology and not increasing any hom-set. There is now at
most one doubled hom-set. One of the two adjacent hom-sets
at $v$ is a singleton, and its neighboring corner is trivial.
Singleton cofinal deletion removes $v$.
The remaining three corners are trivial, so their acyclic
category nerve has dimension at most two and zero fourth
homology. This contradicts the preserved rational square.

We may thus assume all other corners are trivial. If the
total number of incoming arrows at $v$ is one, or the total
number of outgoing arrows is one, the same cofinal deletion
applies, again leaving a nerve of dimension at most two.
Assume both totals are at least two. There are at
most six cross arrows in all.
If $v$ has only two neighbors, they lie on opposite sides.
The two incident hom-sets each have size two, their composite
hom-set is singleton, and the last object attaches by one
singleton edge. This uses all six arrows. Transitivity puts
that leaf edge at a trivial extreme, so the full triangle and
leaf edge form an actual nerve union along a point. The triangle
has a trivial extreme whose two stalks have sizes one and two;
Lemma~\ref{C:two-states-objectwise-act-cups} kills its products.
The leaf edge contributes no higher cohomology.
If $v$ has three neighbors, two are on one side and the third
on the other. Counting the two forced composite hom-sets
again uses exactly six arrows: the two-side hom-sets are
singletons, the other incident hom-set has size two, and
the composites are singletons. At the opposite trivial
extreme all stalks have size at most two. The same categorical
act-cone theorem kills all products. These exhaust the
nonextreme case.

\paragraph{An extreme large corner.}
Take opposites to make $v$ a sink. If it has a singleton
incoming hom-set, its outside budget is at most
$17-q\leq9=2\cdot3+3$, so the preceding lemma applies.
Assume henceforth every incoming hom-set has size at least two.
Let $r$ be the number of its neighbors.

For $r=3$, the entire budget forces $q=8$, all other corners
trivial, and three legs of size two. Identifying the trivial
leaves gives an ordinary act cone over $M$, adding only circles.
It inherits bounded rational homology and second-cohomology
rank one, contradicting Lemma~\ref{C:nine-bounded-rank-one}.

For $r=1$, denote its sole neighbor by $a$ and put
$s=|C(a,v)|$. The complement $D$ is connected, and no
cross arrow in $D$ enters $a$, by transitivity. The nerves
of $D$ and the two-object collage $Y$ on $a,v$ cover $C$
and meet in $B\operatorname{End}(a)$.
If $s=2$, then $|D|\leq7$ and its second rational cohomology
is zero by Lemma~\ref{C:three-small-H2-zero}. The shared
corner has order at most three and is rationally acyclic.
The universal two-cross-arrow theorem makes every degree-two
product on $BY$ zero. Fourth-cohomology restriction to $D,Y$
is injective, proving the contradiction.
If $s\geq3$, connectivity gives $|D|\geq5$, so $s\leq4$
and $q\leq9$. All corners of $D$ have order at most two.
Also $|D|\leq6$, so its second cohomology is zero by the
same three-object lemma. Its semisimple relative bar has
no high homology, and the
shared corner is rationally acyclic. Thus $Y$ inherits
bounded rational homology and second-cohomology rank one.
Exact coinvariants at its small corner give an ordinary
act over $M$ with at most four states. The same bounded
rank-one lemma excludes it.

It remains that $r=2$. Write $a,b$ for its neighbors and
$w$ for the fourth object. A cross arrow involving $w$
must point from $a$ or $b$ to $w$; an arrow in the opposite
direction would force $w\to v$. The two incoming hom-sets,
the remaining connection, and the corner identities force
$q\leq9$. If $q=9$ there is no remaining excess. If $q=8$
there is at most one further corner or cross arrow beyond
these minimum contributions.

If there is just one nonempty hom-set away from $v$, then
the category is the actual union of its central two-leg star
and the edge to $w$, along the corner at $a$ or $b$.
An order-two corner is the only possible extra corner, and
their common corner is rationally acyclic. A singleton edge
is a rationally acyclic join. If that edge has size two,
both endpoints are trivial and it contributes a circle,
contradicting simple connectivity. In the singleton case
the central star inherits bounded rational homology and
rank-one second cohomology. Its leaves have order at most
two, so exact leaf coinvariants followed by point identification
give an ordinary act over $M$, again excluded by the bounded
rank-one lemma.

Otherwise there are exactly two hom-sets away from $v$,
both singletons, $q=8$, all other corners are trivial,
and both incoming hom-sets at $v$ have size two.
If these two hom-sets are $a\to w$ and $b\to w$, the
support is an induced four-cycle. The support functor is
onto on fundamental groups, contradicting simple connectivity.
Otherwise, after swapping $a,b$, they are $a\to b$ and
$a\to w$. Delete the trivial leaf $w$. The remaining
triangle has hom-set sizes $(1,2,2)$, and composition is
an equivariant map between two-element sets. If constant,
its value is a fixed incoming arrow at $M$, so the fixed-arrow
lemma, Lemma~\ref{C:fixed-incoming-cups}, applies with $D$
the trivial two-object interval $a\to b$ and that constant
image as the $M$-fixed arrow into $v$. If it is
bijective, the actual inverse composition map deletes $b$
toward $a$, leaving an ordinary two-state act cone.
Both alternatives have zero degree-two products.
This exhausts all budgets and proves the proposition.
\end{proof}

% Root and B full literal audits PASS; these cuts contain no standalone table-census assertion.
\subsection{Additional profile cuts for the four-object search}

Three elementary reductions supplement the accepted profile lemmas in
the seventeen-arrow four-object search.  Each is applied to an actual
category; none assumes that an arbitrary quotient has finite homotopy type.

\paragraph{A full singleton cut.}
Suppose the object set is partitioned into two nonempty parts $A,B$,
there are no arrows from $B$ to $A$, and every hom-set from an object
of $A$ to an object of $B$ is a singleton.  The category is exactly
the categorical join of its two full subcategories.  Hence its nerve
realization is $BA*BB$.  Choosing a vertex in each nonempty space
and collapsing the union of the two cones on those vertices gives
the usual homotopy equivalence with the suspension of the smash
product.  The collapsed subspace is contractible and is a CW
subcomplex.  Thus all products of positive-degree reduced
cohomology classes vanish, without connectedness assumptions on
the two parts.

\paragraph{The six-arrow strong pair with a singleton side.}
Suppose two objects $a,b$ have corners of order two and exactly
one arrow in each direction.  In a minimal-arrow projective-plane
candidate there is no proper retract.  Consequently the composites
of those two cross arrows are the nonidentity idempotents of their
corners: if either composite were an identity, there would be a
proper retract.  In particular each corner is $\{1,e\}$ with $e$
a zero, and has no unit other than its identity.  If every incoming
hom-set at $b$ from another object is empty or a singleton, its
nonidentity idempotent fixes all incoming cross arrows, literally.
Lemma~\ref{C:one-sided-corner-deletion} therefore deletes $b$ by an
actual homotopy equivalence.  The same conclusion holds if every
outgoing cross hom-set at $b$ is empty or a singleton, by opposition.
At seventeen arrows the result has at most thirteen arrows, since
the two corner arrows and the two internal cross arrows disappear.
It is excluded by Theorem~\ref{C:CP2-lower17}.  The alternative
that the original category already has a retract is excluded by
the same lower bound.

\paragraph{Trees with low-dimensional internal corners.}
Let the undirected support of an acyclic category be a tree.
No two nonidentity cross arrows are composable, since transitivity
would supply a third support edge and hence a triangle.  The nerve
is therefore the union of its two-object edge subcategories, glued
along the corner nerves.  This is an iterated homotopy pushout of
CW subcomplexes along the support tree.

Assume every internal support vertex has a corner monoid of order
at most five.  The accepted classification through order five
gives $H_n(BM;\mathbb Q)=0$ for every such corner and every $n\geq3$.
In the Mayer--Vietoris sequence of the edge union, the intersections
that actually glue two pieces are precisely these internal corners.
It follows for every $n\geq4$ that
\[
 H_n(BC;\mathbb Q)\cong
 \bigoplus_{\{a,b\}\text{ a support edge}}
 H_n(BC|_{\{a,b\}};\mathbb Q).
\]
In particular bounded rational homology of $BC$ passes to every
edge piece, and restriction in fourth rational cohomology is
injective.  It suffices that every edge piece have zero products
of two degree-two classes.  Each of the following numerical
conditions supplies that conclusion:
\begin{enumerate}
\item its cross hom-set has at most two elements, by the universal
two-cross-arrow theorem (the singleton case is a join);
\item both endpoint corners have order at most four, so their
second rational cohomology vanishes and every degree-two class
is relative to the vertex spaces; the relative quotient is a
suspension;
\item one endpoint corner has order at most two and the other has
order at most eleven.  The first corner has projective trivial
left and right rational modules.  Its actual orbit collapse is a
rational cohomology-ring equivalence to an ordinary act cone over
the other corner.  The cone has bounded rational homology by the
preceding comparison, and Lemma~\ref{C:eleven-bounded-all-products}
annihilates all its degree-two products.
\end{enumerate}
All three conditions are statements about every degree-two product
on the edge piece.  In particular no second-cohomology rank-one
assumption is passed from the whole tree to its pieces.

% Root and B full literal audits PASS; complete all998 source/sentinel/independent-cup linkage PASS.
\subsection{The completed four-object seventeen-arrow exclusion}

\begin{theorem}\label{C:seventeen-four-objects}
No finite category with four objects and at most seventeen arrows
has classifying space homotopy equivalent to $\mathbb CP^2$.

\end{theorem}
\begin{proof}
Theorem~\ref{C:CP2-lower17} handles arrow totals at most sixteen.
A projective-plane category is connected.  Enumerate its undirected
support graph, every positive corner cardinality, and every ordered
pair of nonnegative cross cardinalities on each support edge, with
positive sum there and total seventeen.  The six support orbits
represent all $38$ connected labeled graphs on four vertices.
There are $8\,261\,184$ decorations on those representatives,
of which $830\,144$ satisfy the necessary transitivity condition
on nonempty hom-sets.

Apply the previously proved initial/terminal, singleton cofinal,
retract, singleton-sided corner, oriented zero-$H^2$ block, and
star reductions.  Every actual deletion used at this stage leaves
at most sixteen arrows.  Take orbits under the full automorphism
group of each support and opposition; exactly $5\,996$ cardinality
matrices remain.  The separate refinement applies
Theorem~\ref{C:seventeen-four-large-corner}, the accepted fixed-arrow,
parallel-hom, two-block, two-stalk and representable-cover cup
lemmas, and the three explicit cuts in
\path{sections/C_seventeen_four_profile_cuts.tex}.
The latter are the full singleton join, the six-arrow strong pair
with a singleton incoming or outgoing side, and the tree whose
bounded edge pieces have zero degree-two products.

The tree cut deserves one scope precision.  Internal corners have
order at most five, so their rational homology vanishes in every
degree at least three.  The Mayer--Vietoris sequence therefore
identifies fourth homology with the direct sum over edge pieces,
and transfers eventual homology vanishing to those pieces.  Each
used edge theorem annihilates all degree-two products.  No
rank-one second-cohomology hypothesis is passed to an edge.

Exactly $998$ profiles survive all these cuts.  Every one has
acyclic support, all corners have order at most seven, and the
largest number of normalized composable pairs is $71$.
The complete profile refinement, including every rejected cyclic
matrix, is independently checked in
\path{code/outputs/d_seventeen_four_scope_independently_verified.json}.

Fill each corner from the accepted complete ordinary monoid
catalog, retaining both multiplication orientations.  The literal
catalog counts through order seven are
$1,2,7,35,234,2325,32823$.  These are covers; no assertion of
distinctness under ordinary isomorphism is required here.
Every literal corner is checked for identities and associativity.
Internal corner permutations need not be enumerated again because
all remaining typed products are assigned exhaustively.

A corner choice is omitted only under an accepted sufficient
algebraic flag.  A nondegenerate regular trace form modulo $32003$
certifies rational semisimplicity.  The exact normalized bar test
$H^2(-;\mathbb F_2)=0$ certifies $H^2(-;\mathbb Q)=0$.
With semisimple diagonal, this four-object acyclic support has
relative bar length at most three and hence zero fourth rational
homology.  Otherwise the actual oriented two-block predicate
omits a choice only by Lemma~\ref{C:two-block-zero-H2-products}.

For every other choice, each unknown product ranges over its
entire target hom-set.  The finite-domain depth-first search uses
only forced associativity equations to propagate values and tries
every remaining value.  A literal proper retract is excluded by
the smaller-arrow lower bound.  At a complete table, the normalized
binary bar retains precisely the necessary homology profile
$H_1=0$, $\dim H_2=1$.

The search enumerates $215\,194\,501$ literal
corner-representative tables and retains $3\,911\,101$
records.  Permutations within cross hom-sets may remove repeats.
The resumed implementation uses the injective list of typed
nonidentity products as its orbit key and clears its cache above
$500\,000$ entries.  Clearing can only retain extra duplicates;
these are record counts, not claimed globally distinct orbits.
No new category can be lost by that operation.

Every retained record has a separately checked literal
normalized-bar certificate.  Exactly $3\,900\,881$ have
an explicit degree-three primitive for the square of an actual
nonboundary second cocycle.  Each of the other
$10\,220$ records has nonzero third homology over
$\mathbb F_2$.  Both alternatives contradict the homology ring
of $\mathbb CP^2$.  Thus all retained tables are excluded, while
every omitted table has one of the stated proved obstructions.

\end{proof}

\begin{remark}[Independent scope, source, and bar gates]
\label{C:seventeen-four-verification}
The source programs are
\path{code/d_seventeen_four_profiles.cpp},
\path{code/d_seventeen_four_refine.py}, and the two explicitly approved
generators
\path{code/d_categories17_four_profiles.cpp} and
\path{code/d_categories17_four_profiles_bounded.cpp}.
The independent topological-support program
\path{code/d_verify_seventeen_four_topological_scope.cpp}
enumerates all $18$ topological supports and independently
recovers the complete $4\,337$ acyclic baseline profile orbits.
The independent refinement checker tests fresh proof predicates
on all $5\,996$ baseline rows, including all $1\,659$ cyclic rows,
and matches the exact $998$-row input.

For each selected profile the generator emits one exact start
record, one completion record, and a scoped end sentinel.
The source TSV and the primary cup certificates are compressed
while streaming.  Their full literal row sequence is hashed in
both representations before a per-profile completion manifest is
published atomically.  Truncated output is never promoted.

The independent bar routine
\path{code/b_verify_seventeen_category_certificates.cpp} imports neither
the generator nor the primary cup solver.  It reconstructs every
typed table, normalized chain, cocycle, and supplied primitive.
The complete run checks $1\,987\,726\,399$ typed associativity
equations, $1\,128\,128\,424$ cocycle equations, the same number of boundary-square
equations, and $5\,789\,249\,390$ square-primitive equations.
The atomic records in \path{code/outputs/b_c17_four_literal_chunks} link every
literal source row and both gzip footers to its primary manifest.

Finally \path{code/d_verify_seventeen_four_completion.py} requires all
$998$ exact input indices, all scoped start/completion sentinels,
the two approved generator-source/binary pairs, unchanged profile
and checker hashes, and every independent source/cup linkage.
Its final-only record is
\path{code/outputs/d_seventeen_four_complete_independently_verified.json}.
The final full-run sentinel and the identity of the audited independent
checker used in every record are also matched in
\path{code/outputs/d_seventeen_four_final_sentinel_verified.json}.
The original completed profiles $0$ through $56$ remain unchanged;
every byte of the interrupted partial profile $57$ is preserved
separately.  The later disk-guard stop was at a clean profile
boundary.  Neither an interrupted partial nor a prefix is used as
evidence for the complete theorem.

\end{remark}

\subsection{Large corners in the sixteen-arrow chain search}

\begin{lemma}[Two blocks with zero second cohomology]
\label{C:two-block-zero-H2-products}
Let a finite category $C$ have its objects partitioned into two
nonempty full subcategories $A,D$, with no arrows from $D$ to $A$.
If $H^2(BA;\mathbb Q)=H^2(BD;\mathbb Q)=0$, every product of two
classes in $H^2(BC;\mathbb Q)$ is zero.
\end{lemma}
\begin{proof}
The arrows from $A$ to $D$ form a profunctor.  Its two-sided
bar realization $E$ has augmentation maps to $BA$ and $BD$.
A nerve simplex of $C$ either stays in one block or has exactly
one crossing from $A$ to $D$.  Splitting at that crossing gives
the correspondence homotopy pushout
\[
 BC\simeq BA\mathop{\cup}^{h}_{E}BD.
\]
This is the same normalized two-sided bar construction as for
monoid blocks; the identities and intermediate object indices
are retained in the two block bars.  Collapsing the two vertex
spaces separately to two points gives an actual map of these
homotopy pushouts $q:BC\to\Sigma E$, where the suspension is
unreduced.  The cohomology Mayer--Vietoris sequence and the two
vanishing hypotheses show that every degree-two class on $BC$
is in the image of the connecting map from $H^1(E;\mathbb Q)$.
Naturality identifies that connecting map with $q^*$ on
$H^2(\Sigma E;\mathbb Q)$.  It is therefore surjective.
Positive cohomology products on a suspension vanish; pulling
back proves the assertion.  If $E$ is disconnected, its
components introduce only the usual additional degree-one
suspension classes and do not change this argument.
\end{proof}

\begin{lemma}[Two objects and five arrows]
\label{C:two-object-five-H2-zero}
Every connected two-object category with at most five arrows
has zero second rational cohomology.
\end{lemma}
\begin{proof}
If arrows exist in both directions, one of the two corners
has order one: otherwise two corner identities and at least
two additional corner arrows, together with two cross arrows,
would already require six arrows.  That object is a retract
of the other.  Deleting it leaves a corner of order at most
two, whose rational cohomology is concentrated in degree zero.

Otherwise the category is triangular.  With one cross arrow
its classifying space is the join of its two connected corner
classifying spaces.  Such a join has zero second rational
cohomology.  With at least two cross arrows the corner orders
sum to at most three.  Both rational corner algebras are
semisimple, and the relative normalized bar has length one.
Its rational cohomology in degree two is consequently zero.
\end{proof}

\begin{proposition}[The only large-corner inputs for sixteen-arrow chains]
\label{C:sixteen-linear-large-corners}
After the singleton-join, singleton-cofinal, fixed-arrow and
triangle-compression reductions in the fifteen-arrow chain
proof, every remaining three-object linear sixteen-arrow
profile has corner orders at most eight.  A middle corner of
order eight is excluded structurally.  At an extreme corner
of order eight it suffices to enumerate the complete catalog
of monoids with possibly nonzero second rational cohomology.
No other eight-element monoid is needed.
\end{proposition}
\begin{proof}
Write the corner orders as $(m,n,p)$ and the three forward
hom-set sizes as $(x,y,z)$, in directions $01,12,02$.
These are positive integers with sum sixteen.  The profile
calculation retained in
\path{code/outputs/c_sixteen_linear_scope_preliminary.json}
forms all $\binom{15}{5}=3003$ positive six-tuples directly,
then applies exactly the stated inequalities.  Of these,
$506$ are singleton joins, $430$ singleton-cofinal cases,
$189$ fixed-arrow cases, and $23$ triangle-compression cases.
The remaining $1855$ tuples have corner orders at most eight.
One of them is $(1,8,1;2,2,2)$, excluded by
Corollary~\ref{C:middle-double-homs-rational}.
The other middle-eight tuples are
\[
 (1,8,2;2,1,2),\quad (1,8,2;2,2,1),\quad
 (2,8,2;1,1,2)
\]
and the opposites of the first two.

At an endpoint of order two, both constant rational modules
are projective.  Passing to actual cross-act orbits gives
a functor inducing a rational cohomology-ring equivalence,
by the exact coinvariant comparison already proved for
triangular blocks.  Collapse each such endpoint.  The two
endpoints are now trivial, and every forward hom-set has
size at most two.  If either adjacent hom-set is a singleton,
singleton-cofinal deletion leaves a two-object category of
two trivial corners, hence a graph.  If both adjacent sizes
are two and the long size is one, triangle compression gives
a monoid with a two-by-two minimal band.  Its degree-two
rational products vanish by Proposition~\ref{B:rank-one-square}, without
any homological boundedness assumption.  If all three sizes
are two, Corollary~\ref{C:middle-double-homs-rational} applies.
These alternatives exclude every large middle corner after
the rational orbit comparison.

Now suppose the eight-element corner is at the source.
Unless both outgoing hom-sets are singletons, a previously
removed join, they consume at least three arrows.  The
complementary two-object full subcategory consequently has
at most $16-8-3=5$ arrows.  It is connected, since its forward
hom-set is nonempty, and Lemma~\ref{C:two-object-five-H2-zero}
applies.  If the eight-element corner also has zero second
rational cohomology, Lemma~\ref{C:two-block-zero-H2-products}
excludes the whole category.  Opposition proves the sink case.
Thus only the already complete possible-positive-$H^2$
eight-element catalog is required.  The generator retains
all its literal tables and both orientations; it makes no
claim that this smaller catalog contains every monoid of
order eight.
\end{proof}

\subsection{Exact block cuts and sixteen-arrow star reductions}

\begin{corollary}[Height-one support]
\label{C:height-one-zero-H2-products}
Let $C$ be a finite category with no composable pair of arrows
between distinct objects.  If every corner has zero rational
second cohomology, all products in $H^2(BC;\mathbb Q)$ vanish.
\end{corollary}
\begin{proof}
Partition the objects into sources and sinks, assigning isolated
objects arbitrarily.  Each full block is a disjoint union of
corner monoids, so its second cohomology is zero.  No arrow
goes from the sink block to the source block.
Lemma~\ref{C:two-block-zero-H2-products} applies when both
blocks are nonempty.  If one is empty, the category itself is
a disjoint union of corners and has zero second cohomology.
\end{proof}

\begin{proposition}[Sufficient zero-cohomology flags for the search]
\label{C:exact-block-cut-flags}
The following are sufficient conditions for a full subcategory
$D$ to have $H^2(BD;\mathbb Q)=0$: each connected component
satisfies at least one of the following conditions.
\begin{enumerate}
\item it is a single corner whose normalized second homology
over $\mathbb F_2$ is zero, or whose rational algebra is semisimple;
\item it has two objects and at most five arrows, or exactly
one cross arrow, or is triangular with both corner algebras
semisimple over $\mathbb Q$;
\item it is connected, has three objects, and at most seven arrows.
\end{enumerate}
Thus an oriented partition of the objects into two blocks satisfying
these conditions excludes a projective-plane square, provided no
arrow goes from the second block to the first.
\end{proposition}
\begin{proof}
Cohomology of a finite disjoint union is the direct sum in degree
two.  For a corner of order $q$, the normalized bar groups in
degrees one through three have dimensions $q-1$, $(q-1)^2$,
and $(q-1)^3$.  Over $\mathbb F_2$ the two relevant differentials are
\[
 \partial[a|b]=[b]+[ab]+[a],\qquad
 \partial[a|b|c]=[b|c]+[ab|c]+[a|bc]+[a|b],
\]
with any identity-containing term omitted.
Hence the exact dimension is
$(q-1)^2-\operatorname{rank}\partial_2-
\operatorname{rank}\partial_3$.  An integral matrix has rational
rank at least its rank modulo two.  A zero value therefore forces
zero second rational homology, and then zero rational cohomology.
For a semisimple rational corner algebra all positive Ext groups
vanish.

Lemma~\ref{C:two-object-five-H2-zero} proves the five-arrow case.
A singleton cross hom-set gives the join of two connected corner
spaces, which has zero second rational cohomology.  In the
triangular semisimple case the relative bar over the product of
the two corner algebras has no term of length two, and exact
corner coinvariants leave no second cohomology.
Lemma~\ref{C:three-small-H2-zero} proves the three-object case.
The last assertion is precisely
Lemma~\ref{C:two-block-zero-H2-products}.

The implementation \path{code/c_cp2_block_prunes.hpp} computes
the two displayed ranks by exact binary elimination.  Its uses
have $q\leq8$, so the largest binary row has $(q-1)^2\leq49$
positions and fits in the stated 64-bit word.  It first splits
each proposed block into its actual connected components, and
tests every nonempty proper object subset for the absence of
reverse arrows.  A failed sufficient test never removes a case.
\end{proof}

\begin{proposition}[Four-object stars through sixteen arrows]
\label{C:sixteen-four-star-prunes}
The three exclusions in
Proposition~\ref{C:fifteen-four-star-prunes} remain valid for a
four-object inward or outward star with at most sixteen arrows.
\end{proposition}
\begin{proof}
The first exclusion uses only zero second cohomology at the
three leaves and a singleton incoming arrow fixed by the central
monoid.  Its proof by Lemma~\ref{C:fixed-incoming-cups} has no
fifteen-arrow restriction.

For a central monoid of order at most nine and leaves of orders
at most two, the exact leaf-coinvariant functors in the earlier
proof still give a rational cohomology-ring equivalence.  Identifying
the resulting three sink points gives the ordinary act cone with
two additional circle summands.  A hypothetical projective plane
would give bounded rational homology, degree-two rank one, and
a nonzero square, contrary to
Lemma~\ref{C:nine-bounded-rank-one}.  Its hypotheses explicitly
allow the additional first homology.

Finally two singleton legs give the same actual splitting
\eqref{C:four-star-two-singleton-splitting}.  The primitive
projective-plane square must lie entirely in one wedge summand.
The suspension summand has zero products, and consequently must
be integrally acyclic.  Its desuspension is the central monoid
space; Theorem~\ref{A:bounded} and the homological Whitehead
theorem make that space contractible.  The remaining two joins
are contractible, and the whole category reduces to its
center--third-leaf collage.  That collage has at most
\[
 16-(|A_1|+1)-(|A_2|+1)\leq12
\]
arrows, already excluded by Theorem~\ref{C:CP2-lower16}.
All three arguments are therefore valid at sixteen arrows.
\end{proof}

\subsection{The completed sixteen-arrow star calculation}

\begin{theorem}\label{C:sixteen-all-stars}
A three-object category with star support and at most sixteen
arrows cannot have classifying space homotopy equivalent to
$\mathbb {CP}^2$.
\end{theorem}
\begin{proof}
Theorem~\ref{C:CP2-lower16} treats smaller arrow counts.
Use opposition to make the star outward, and write
$(m,a,b,x,y)$ for the central corner order, the two leaf
corner orders, and the two cross hom-set sizes.  These are
positive integers of sum sixteen.  Order the leaves by
$(a,x)\leq(b,y)$.

The generic star exclusions used in
Lemma~\ref{C:fifteen-star-deletions} continue to apply.
The trivial-center wedge now has proper collages of at most
fourteen arrows, already excluded.  In its two coinvariant
reductions, replace the bound nine by ten, using
Lemma~\ref{C:ten-rational-act-cones} at order ten and
Lemma~\ref{C:nine-bounded-rank-one} below it.  In each use
the rational replacement of a hypothetical projective plane
has bounded homology and degree-two rank at most one; no
assertion about arbitrary mixed products of split row modules
is needed.

For reproducibility, the complete excluded cardinality
conditions are the following disjunction, together with the
condition that all three corner orders are at most four:
\[
\begin{gathered}
 m=1;\qquad (x=1,\ b\leq4);\qquad(y=1,\ a\leq4);
 \qquad x=y=1;\\
 m\leq10,\ a,b\leq2;\\
 m\leq4,\ b\leq4,\ x\leq2;\qquad
 m\leq4,\ a\leq4,\ y\leq2;\\
 m\leq2,\ b\leq4,\ a\leq10;\qquad
 m\leq2,\ a\leq4,\ b\leq10;\\
 m\leq4,\ x,y\leq2.
\end{gathered}
\]
The remaining list has exactly $74$ profiles.  It is obtained
both by five-integer enumeration and by independently cutting
sixteen into five positive parts.  The exact list is stored in
\path{d_sixteen_star_scope.json}.  Only two profiles have a
corner of order eight:
\[
 (3,1,8,1,3),\qquad (8,1,3,2,2).
\]
In each there is only one potentially positive-degree-two
corner, so the complete positive-$H^2$ order-eight catalog
suffices.

At least one vertex must have nonzero rational second
cohomology.  Otherwise every degree-two class is relative
to the vertex spaces, whose quotient is a wedge of suspended
edge bars; all such products vanish.  When just one corner
can have positive $H^2$, that corner ranges over the complete
positive catalog and every other corner ranges over its full
ordinary catalog.  When at least two corners can have
positive $H^2$, \emph{all} those corners range over their full
ordinary catalogs.  The budget and the two-singleton deletion
bound the orders of these full corners by seven.
Their exact positivity test is the minimal-ideal rational
augmentation pairing: matching row and column characters for
a $2\times2$ band, or an identity adjoined to a $2\times3$
or $3\times2$ band.  The smaller row/column and group cases
have zero rational $H^2$ in this order range.

There is one further literal corner-choice reduction.  A
singleton leg is a join of its two vertex spaces, so restriction
of every star degree-two class to both those vertices is zero.
If the opposite leaf also has zero $H^2$, every degree-two
class is relative to all vertices and has zero products.
Thus, when $x=1$, retain only choices with positive $H^2$ at
the second leaf; when $y=1$, retain only choices with positive
$H^2$ at the first leaf.  This removes actual corner choices,
not additional cardinality profiles.

For fixed corner tables, a leg is an arbitrary commuting
biset.  Enumerate every central right action up to exact
state conjugacy, then every commuting opposite-leaf action
in its actual relative position.  The two leg state sets
are disjoint, so their state conjugacies may be made
independently.  There is no composable pair between different
legs.  Taking every pair of the two resulting leg lists
therefore imposes all and only the associativity requirements
of the star.

The complete enumeration has $6\,384\,648$ cases.  In the
order of its tests the mutually exclusive outcomes are
\[
\begin{array}{l|r}
 H_1(-;\mathbb F_2)\ne0&1\,958\,746\\
 H_1=0,\quad\dim H_2(-;\mathbb F_2)\ne1&2\,852\,329\\
 \text{an explicit primitive for the degree-two square}&1\,564\,288\\
 \text{nonzero square and nonzero }H_3(-;\mathbb F_2)&9\,285.
\end{array}
\]
Every projective plane has zero first and third mod-two
homology, one-dimensional second cohomology, and a nonzero
generator square.  It satisfies none of the four outcomes.
This excludes outward stars, and opposition excludes inward
stars as well.
\end{proof}

\begin{remark}[Independent scope, table, and cochain gates]
\label{C:sixteen-stars-verification}
The exact inputs are generated by
\path{d_sixteen_star_inputs.py}; the primary producer is
\path{d_sixteen_stars.cpp}.  Its corner catalogs are the
previously independently verified ordinary and positive
catalogs, retaining both orientations where necessary.
The separate checker \path{d_verify_sixteen_star_catalogs.py}
checks all $35\,882$ distinct literal corner tables, including
$12\,047\,720$ associativity equations, and the full-catalog
$H^2$ flags.  It also canonically matches every flagged
positive class with the complete positive catalog.  Duplicate
ordinary catalog entries are permitted and are recorded;
they add cases rather than remove any.

The independent count producer
\path{d_sixteen_star_counts.cpp} imports neither the primary
pointwise-domain action solver nor any bar routine.
It instead chooses monoid generators, extends their
transformation images using all multiplication relations,
and uses exact stabilizer canonical augmentation for the
right action.  It enumerates the complete transformation
centralizer of that action and every left-action homomorphism
inside it.  This independently reproduces every one of the
$74$ profile counts, every corner-choice count, and every
structural-skip count.  The exact full-scope linkage is
\path{d_sixteen_star_count_scope_verified.json}; the additional
Python implementation is a control on its separately recorded
completed profiles, not a premise for completion of the C++
count gate.

The independent literal category and cochain checker is
\begin{center}\path{b_verify_category_certificates.cpp}.\end{center}
It does not import
the primary generic cup solver.  All $1\,573\,573$ retained
records are linked exactly to the primary compressed stream,
with complete profile counts, final metadata, gzip footer,
and whole-stream hash checked.  The checker reconstructs the
normalized bars, verifies that each degree-two cocycle is
nonzero in cohomology, checks every saved square primitive,
and independently recomputes positive third homology in the
remaining cases.  It checks $1\,006\,450\,110$ typed
associativity equations, $609\,451\,179$ cocycle and
boundary-squared equations, and $3\,427\,725\,756$
square-primitive equations.  The final record
\path{b_sixteen_stars_independently_verified.json} has status
\texttt{PASS}.
Completeness of the action search follows from the explicit
finite homomorphism enumeration above and its independent
regeneration; it is not inferred merely from the number of
saved cochains.
\end{remark}

% Complete gate and full root/B literal proof audits passed before input.
\subsection{The completed seventeen-arrow star calculation}

\begin{theorem}\label{C:seventeen-all-stars}
A three-object category with inward or outward star support and
at most seventeen arrows, identities counted, cannot have
classifying space homotopy equivalent to $\mathbb{CP}^2$.
\end{theorem}
\begin{proof}
The global lower bound treats at most sixteen arrows. For exactly
seventeen use opposition to make the star outward. Write
$(m,a,b,x,y)$ for its central corner order, its two leaf corner
orders, and the respective leg sizes. All five integers are
positive, have sum seventeen, and satisfy $(a,x)\leq(b,y)$.
There are exactly $924$ such canonical cardinality profiles:
one can enumerate the positive compositions, or independently
choose four cuts among the sixteen interior positions and then
order the leaves.

The following disjunction excludes $841$ entire profiles:
\[
\begin{gathered}
 \max(m,a,b)\leq4;\qquad m\leq4;\\
 (x=1,\ b\leq4);\qquad (y=1,\ a\leq4);
 \qquad x=y=1;\\
 m\leq11,\ a,b\leq2;\qquad
 (m,a,b,x,y)=(8,1,5,1,2).
\end{gathered}
\]
The first condition uses the relative suspension after restriction
to all three corner spaces: when their rational $H^2$ vanishes,
every degree-two product is zero. The second is
Corollary~\ref{C:star-small-center18}. A singleton leg gives
the join of its two corner spaces, and therefore kills restriction
of degree-two classes to both incident corners. This proves the
next three conditions by the same relative-suspension argument.
For leaves of order at most two, their trivial rational modules
are projective; Corollary~\ref{C:small-center-projective-leaves-starlimit}
proves the penultimate condition. The last is
Corollary~\ref{C:star-eight-one-five-singleton-two}.
The exact residual list consists of $83$ profiles and is saved in
\path{code/outputs/c_seventeen_star_scope.json}.

For each residual profile retain every relevant corner table.
At least one corner must have positive rational $H^2$.
If $x=1$, that corner must be the opposite leaf of order $b$;
if $y=1$, it must be the leaf of order $a$.
Without a singleton condition, a unique corner of order at
least five must have positive $H^2$. In precisely these forced
positions the complete possible-positive catalog suffices.
All other positions use the full ordinary monoid catalog,
including zero-$H^2$ corners. The largest ordinary corner has
order seven, and the largest forced-positive corner has order
nine. Any harmless extra entries in a possible-positive catalog
only enlarge the enumeration. No ordinary order-eight catalog
is substituted for its positive sublist.

There are also literal central-corner omissions. If its minimal
ideal has one row, its trivial left rational module is projective,
so Theorem~\ref{C:star-projective-center22} applies. If it has a
literal left zero, its classifying space is contractible, and
Corollary~\ref{C:star-contractible-center18} applies. Both flags
are checked on each actual multiplication table. In the first
case a uniform vector on a minimal left ideal is fixed by every
left multiplication; in the second the asserted left-zero
equations are checked literally. The remaining full-catalog
positivity flags are the small-corner Rees augmentation pairing.

For fixed corner tables, each leg is an arbitrary commuting
biset. Enumerate all right actions of the central monoid up to
state conjugacy and every commuting left action of its leaf
monoid in the actual relative position. The two leg state sets
are disjoint. Their relabelings can thus be chosen independently,
and taking every pair of the two complete leg lists gives all
stars. No cross arrow in one leg composes with a cross arrow
in the other. The multiplication laws inside each corner and
the two commuting biset laws are exactly the remaining category
axioms. This is a complete covering enumeration; it is not
asserted to be a census without repetitions of category
isomorphism classes.

There are $20\,193\,947$ cases. The mutually exclusive outcomes,
in the order in which the tests are applied, are
\[
\begin{array}{l|r}
 H_1(-;\mathbb F_2)\ne0&8\,320\,500\\
 H_1=0,\quad\dim H_2(-;\mathbb F_2)\ne1&7\,226\,817\\
 \text{an explicit primitive for the degree-two square}&4\,639\,637\\
 \text{nonzero square and nonzero }H_3(-;\mathbb F_2)&6\,993.
\end{array}
\]
There is no remaining case. Each outcome contradicts the
mod-two cohomology ring or homology of $\mathbb{CP}^2$.
Opposition proves the inward-star assertion.
\end{proof}

\begin{remark}[Independent finite scope and cochain verification]
\label{C:seventeen-stars-verification}
The producer is \path{code/c_seventeen_stars.cpp}, using
\path{code/c_seventeen_star_inputs.py}. The exact input catalog
is linked to the previously verified ordinary and positive
monoid catalogs. A separate composition-cut scope verifier
recovers all $924$ profiles and the $841/83$ partition. It checks
$44\,052$ distinct multiplication tables, $18\,003\,650$
monoid associativity equations, $203\,419$ full-catalog positivity
flags, and $168\,934$ literal central exclusion flags.
The plain text consumed by both compiled programs is matched
entry by entry to the hashed JSON input, including every table,
every flag, all counts, and end of file.

The independent generator-image enumerator
\path{code/c_seventeen_star_counts.cpp} uses no primary
pointwise-domain action solver and no bar routine. It chooses
generators, extends their transformation images through all
multiplication relations, and uses exact stabilizer augmentation
to enumerate right actions up to state conjugacy. It then
enumerates all left actions in the full transformation
centralizer. Every profile count, corner-choice count, and
structural-skip count matches the primary enumeration.

The separate literal checker
\path{code/c_verify_seventeen_star_certificates.cpp}
reconstructs the normalized bar complexes of all
$4\,646\,630$ retained categories, verifies the nonzero
degree-two cohomology generator, checks every recorded square
primitive, and independently verifies nonzero third homology
in the remaining cases. It imports no primary cup solver.
Its work comprises $3\,735\,095\,577$ typed associativity equations,
$2\,370\,755\,931$ cocycle equations, $2\,370\,755\,931$ boundary-square
checks, and $15\,112\,495\,568$ primitive equations.

The final gate also matches $314\,226\,734$ entries of retained
corner blocks to their literal source catalogs. Every retained
row is in exact order, every part is contiguous, every profile
count agrees, and the full stream hash, final completion record,
gzip footer, and end of file are checked. The complete result is
\begin{center}
\path{code/outputs/c_seventeen_stars_independently_verified.json}.
\end{center}
Its status is
\begin{center}\texttt{PASS\_COMPLETE\_SEVENTEEN\_ARROW\_STAR\_SCOPE}.
\end{center}
Neither a completed prefix nor the number of saved certificates
alone is used as evidence that the action scope is complete.
\end{remark}

\subsection{High-object reductions at sixteen arrows}

\begin{lemma}[A singleton tree with one nontrivial corner]
\label{C:tree-singleton-core-cups}
Suppose a connected finite category has at least two objects,
its undirected cross-arrow support is a tree, every nonempty
cross-hom set is a singleton, and at most one corner monoid
is nontrivial.  Every positive-degree rational product on its
classifying space is zero.
\end{lemma}
\begin{proof}
If an edge has arrows in both directions, any edge connecting
one of its vertices to a third vertex would force the other
side of a triangle by composition.  Connectedness of the tree
therefore leaves only those two objects.  At least one corner
is trivial, so that object is both initial and terminal and
the category is contractible.  We may exclude this case.

A directed path through three distinct objects would force
its composite cross-hom set and hence a triangle in the
undirected support.  No such path occurs.  Thus the nerve
is the graph of correspondence spaces indexed by that tree.
Remove a leaf whose corner and neighboring corner are both
trivial; its correspondence is a point and its mapping
cylinder is an interval, so this is a homotopy-equivalence
collapse.  Repeating leaves a star around the possible
nontrivial corner $M$, with $d\geq1$ trivial leaves.
Each edge correspondence is $BM$ and its map to the center
is the identity, because its unique cross arrow is fixed
by all actors.  The resulting space is the join of $BM$
with a discrete $d$-point set.  It is contractible for $d=1$
and a wedge of $d-1$ suspensions of $BM$ for $d>1$.
Its positive-degree products vanish.  If every corner is
trivial, the original tree itself is contractible.
\end{proof}

\begin{lemma}[Semisimple diagonal choices with many objects]
\label{C:sixteen-high-semisimple}
A finite category with sixteen arrows, at least five objects,
acyclic support, and semisimple rational corner algebras
cannot have classifying space $\mathbb CP^2$.
\end{lemma}
\begin{proof}
The relative normalized bar over the product of the corner
algebras is an exact rational cohomology model.  Its length
is at most the height of the support order.  A four-step
chain requires ten distinct cross-hom sets.  With $k$ objects,
connecting the other $k-5$ objects requires at least $k-5$
more cross arrows; together with identities this costs
$2k+5$.  For $k\geq6$ it exceeds sixteen, so the relative
bar has length at most three and fourth homology vanishes.

For $k=5$, a four-step chain forces all ten forward hom-sets.
If all corners are trivial, Theorem~\ref{C:all-acyclic21}
applies.  Otherwise exactly one corner can have order two,
all others are trivial, and every cross-hom set is a singleton.
At least one extreme object of the chain has trivial corner;
it is respectively initial or terminal.  The category is
contractible.  These cases exhaust the sixteen-arrow budget.
\end{proof}

\begin{theorem}[Seven or more objects at sixteen arrows]
\label{C:sixteen-seven-objects}
Every finite category with at least seven objects and at most
sixteen arrows fails to have classifying space $\mathbb CP^2$.
\end{theorem}
\begin{proof}
Smaller arrow totals are already excluded.  Assume connectedness.
The tree budget leaves only seven or eight objects.  For $k$
objects, write $e=\sum_c(|\operatorname{End}(c)|-1)$ for the
number of nonidentity corner arrows.  Relative to the minimum
$2k-1$ arrows, there are $17-2k$ extra arrows.

We first dispose of nontrivial strongly connected components.
A retract object can be deleted, reducing to the already
proved lower bound.  Without a retract, every corner in a
nontrivial strong component has order at least two.  A
component of $r$ objects also has all $r(r-1)$ directed cross
hom-sets nonempty.  Compared with a connecting tree on those
objects, these requirements cost at least
\[
 r+\bigl(r(r-1)-(r-1)\bigr)=r+(r-1)^2
\]
extra arrows.  This is at least seven for $r\geq3$ and three
for $r=2$.  At eight objects there is only one extra arrow,
so no such component remains.  At seven objects there are
three extras: a possible strong pair has both corner orders
two, one arrow each way, and every other corner and cross
hom-set is trivial or a singleton.  If an opposite composite
is a unit, it gives a retract.  Otherwise each corner is the
identity-adjoined zero monoid.  Its zero idempotent fixes every
incoming cross arrow, since each such hom-set is a singleton.
Lemma~\ref{C:one-sided-corner-deletion} removes a vertex of
the pair, again reducing the arrow count.  We may therefore
assume acyclic support.

For eight objects, $e\leq1$, so every corner has order at
most two and every corner algebra is semisimple over
$\mathbb Q$.  Lemma~\ref{C:sixteen-high-semisimple} applies.
For seven objects, $e\leq3$.  If all corner orders are at
most two, the same lemma applies.  Thus only a corner of
order three or four needs further attention.

If $e=3$, there are exactly six cross arrows.  Connectedness
makes their undirected support a tree and each cross-hom a
singleton.  Every corner has order at most four and hence
rationally acyclic classifying space, by the complete small
corner classification.  Every edge correspondence in the
graph-of-spaces description is the product of the two adjacent
corner classifying spaces, and is likewise rationally acyclic.
The Mayer--Vietoris sequence along the tree gives rational
acyclicity of the whole category.

The remaining case has $e=2$, one order-three corner, and
all other corners trivial.  There are seven cross arrows.
If they are in distinct hom-sets, projection to the support
poset is a rational homology equivalence.  Indeed, filter the
nerve by its strict object chain: over each such chain the
corner bars are products of rationally acyclic classifying
spaces, and the singleton cross arrows give their augmentation
maps.  The associated homology spectral sequence is precisely
the ordinary support-poset complex.  A four-object chain would
already require six cross arrows and three further connecting
arrows, so this support has height at most two.  Its fourth
homology is zero.

Otherwise there is exactly one two-element cross-hom set,
all others are singletons, and the undirected support is a
tree.  It has no directed path through three distinct objects.
All edge correspondences except possibly the doubled one are
rationally acyclic.  The doubled correspondence is either a
discrete two-point space or the action category of a two-state
act over the order-three corner.  Its rational homology is
zero in degrees at least two.  To check this last assertion,
the complete order-three corner list consists of group-type
monoids, whose trivial rational modules are projective, and
the identity-adjoined left- and right-zero bands of order two.
The latter rational algebras are triangular two-by-two matrix
algebras and have global dimension one.  The two-sided action
bar therefore has no homology in degrees at least two in
every case.  Mayer--Vietoris for the graph of correspondences
now gives zero fourth homology of the whole category.  This
exhausts $e=2,3$, and proves the theorem.
\end{proof}

\subsection{A uniform bound from the excess number of arrows}

\begin{theorem}\label{C:arrow-excess-four}
Let $C$ be a finite category with $k$ objects and $n$ arrows.
If $n\leq2k+3$, then $BC$ is not homotopy equivalent to
$\mathbb CP^2$.
In particular, every category with at least nine objects and
at most twenty-one arrows is excluded.
\end{theorem}
\begin{proof}
Only connected categories need consideration.  Their
undirected support has at least $k-1$ edges, and their
identities supply $k$ further arrows.  Define the nonnegative
excess
\[
 e=n-(2k-1)\leq4.
\]
Delete any object whose identity factors through another
object, using the already proved retract deletion.  This
preserves the classifying-space homotopy type and
connectedness.  At least its identity and two cross arrows
are removed; hence $k$ drops by one and $e$ drops by at
least one.  We may therefore assume that no object is a
retract of another, still with $e\leq4$.
If at most two objects remain, the arrow count is at most
seven, already excluded by Theorem~\ref{C:CP2-lower16}.

We first remove the possibility of a nontrivial strongly
connected component.  Let such a component have $t\geq2$
objects.  Every ordered cross hom-set inside it is nonempty.
Every corner inside it has at least two arrows: otherwise
the identity in that trivial corner factors through any
other object of the component.  Thus its full subcategory
contains at least $t^2+t$ arrows.  If it is the whole category,
then $t=k\geq3$ gives
\[
 e\geq k^2-k+1\geq7,
\]
which is impossible.

Suppose it is proper.  The other objects contribute at least
$k-t$ identities.  Contract the component in the undirected
support.  A connected graph on the remaining $k-t+1$
vertices needs at least $k-t$ edges.  At least one such edge
meets the contracted component, and any arrow across that
interface forces an arrow to or from every one of its $t$
objects, by composition inside the component.  Consequently
the arrows not internal to the component include at least
$k-1$ cross arrows.  In total
\[
 n\geq(t^2+t)+(k-t)+(k-1)=2k+t^2-1,
 \qquad e\geq t^2.
\]
Equality is therefore forced, with $t=2$ and $e=4$.
All inequalities used must be equalities.  There is exactly
one neighboring object $v$ outside the component; it has a
trivial corner, and the two cross hom-sets at this interface
are singletons, both oriented toward $v$ or both away from
$v$.  All other corners are trivial, and the contracted
support is a tree with singleton cross hom-sets.

The full subcategory on the component and $v$ has $v$ as a
terminal or initial object.  It meets the rest of the
category only in $v$; a nerve simplex cannot use an object
of the component and any other outside object.  Its
contractible nerve can therefore be removed by a homotopy
pushout along that point.  The remaining nerve is the tree
of singleton arrows on trivial corners, hence contractible.
This case cannot be a projective plane.

We are left with acyclic support.  Put
\[
 D=\sum_x(|\operatorname{End}(x)|-1),\qquad
 P=\sum_{x\ne y,\ C(x,y)\ne\varnothing}(|C(x,y)|-1),
 \qquad R=E-k+1,
\]
where $E$ is the number of edges of the undirected support.
There is at most one direction on each such edge, and
\begin{equation}\label{C:excess-decomposition}
 e=D+P+R\leq4.
\end{equation}
A strict chain of four arrows would supply all ten edges
on its five objects, and hence force $R\geq6$.
Thus there are at most three strict changes of object in
any simplex.

If all corner algebras are semisimple over $\mathbb Q$,
the relative bar over their product has no internal higher
Tor.  Its strict-object degree is at most three, and
$H_4(BC;\mathbb Q)=0$.

Every corner has at most five elements.  If one has five,
then $D=4$ and $P=R=0$.  There is only that one nontrivial
corner, the support is a tree, and every cross hom-set is
a singleton.  A valid acyclic category on a tree has no
composable pair of strict arrows, since their composite
would create a triangle.  The five-element corner is
incident to a singleton leg to a point.  The corresponding
two-object classifying space is the join of its corner
space with a point, so every global degree-two class
restricts to zero on that corner.  It also restricts to
zero on every other vertex.  The quotient by the union
of vertex spaces is a wedge of suspended edge bars, as
in Corollary~\ref{C:height-one-zero-H2-products}.
All the degree-two classes are relative and all their
products vanish.

Henceforth every corner has at most four elements.  By
Theorem~\ref{E:all-six-types} its classifying space is
rationally acyclic.  If $R=0$, support is again a tree of
height one, and Corollary~\ref{C:height-one-zero-H2-products}
applies.  If $P=0$, collapse every corner to its object.
This actual functor to the support poset is a rational
homology equivalence.  To see the comparison directly,
filter by the finite number of strict object changes:
because every cross hom-set is a singleton, each
associated-graded correspondence bar is the tensor
product of the ordinary bars of the corners in that
chain.  All those corner bars are rationally acyclic.
The comparison is therefore an isomorphism on the first
page and on total homology.  The support poset has
dimension at most three, so its fourth homology vanishes.

The only case still requiring consideration has a
nonsemisimple corner algebra and $P,R\geq1$.
An order-one or order-two monoid algebra is semisimple,
so $D\geq2$.  Equation~\eqref{C:excess-decomposition}
forces
\[
 D=2,\qquad P=1,\qquad R=1.
\]
There is exactly one corner of order three, all others
are trivial, the support is unicyclic, and exactly one
cross hom-set has two arrows; all other nonempty cross
hom-sets are singletons.

If the unique undirected cycle is not a triangle,
there can be no composable pair of strict arrows, and
the height-one corollary again applies.  Otherwise the
unique cycle is a transitive triangle.  The category
nerve is the union of its full triangle subcategory
and its full two-object bridge subcategories.  Any
simplex with two strict changes of object must lie in
that triangle; every other simplex lies in a bridge
piece or a corner.  These pieces form a tree of
subcomplexes meeting in their corner nerves.  The
intersection spaces are rationally acyclic.  The
Mayer--Vietoris sequence thus identifies fourth rational
homology of the whole nerve with the direct sum of
fourth homology of these pieces.

Each bridge piece has trivial corners except possibly
the single order-three corner.  Both trivial modules
of an order-three monoid algebra have projective
dimension at most one, as established in the proof
of Proposition~\ref{C:semisimple-ends-small-middle}.
Its edge bar has no homology in internal degree at
least two, and the bridge collage has zero fourth
homology.

If the triangle has only trivial corners, its nerve
has dimension at most two.  If its order-three corner
is an endpoint, the strict-object bar has internal
degree at most one, again by the trivial-module
projective-dimension bound at that endpoint; all other
corner algebras in the piece are $\mathbb Q$.
Total degree four is absent.  Finally, if the
order-three corner is in the middle, at least one
adjacent hom-set is a singleton, since there is only
one doubled hom-set in the entire category.  The
singleton-cofinal deletion removes that middle
object, leaving a graph.  Thus the triangle piece
also has zero fourth rational homology.
The direct-sum comparison proves $H_4(BC;\mathbb Q)=0$
in the last case.  Every alternative has now excluded
a projective-plane square or its fourth homology,
which proves the theorem.
\end{proof}

% Independently audited by root and B; source note retained.
\subsection{Acyclic support with at most six excess arrows}

\begin{lemma}[Singleton hom-sets and the support poset]
\label{C:singleton-support-H2}
Let $C$ be a finite connected category with at least two objects,
acyclic support, and exactly one arrow in each nonempty hom-set
between distinct objects.  Let $\Pi$ be its support poset.
The actual functor $q:C\to\Pi$ which collapses every corner
monoid to its identity induces an isomorphism
\[
 q^*:H^2(B\Pi;\mathbb Q)\xrightarrow{\ \cong\ }
      H^2(BC;\mathbb Q).
\]
No assumption on the positive rational homology of the corner
monoids is required.
\end{lemma}
\begin{proof}
The unique cross arrows make $q$ a functor: every composite
between different endpoints is the unique arrow with those
endpoints.  Filter normalized nerve chains by the number $p$
of strict changes of object.  For each chain of $p+1$ distinct
objects, the associated internal bar is the tensor product
of the ordinary bars of their corner monoids, because all
cross coefficients are singleton trivial modules.  Its
degree-zero homology is $\mathbb Q$, so the $q=0$ row is
exactly the simplicial chain complex of $B\Pi$.

Every finite monoid has zero first rational homology.  Hence
the entire internal-degree-one row is zero.  In internal
degree two, the bar for an edge $x<y$ has homology
\[
 H_2(B\operatorname{End}(x);\mathbb Q)
 \oplus H_2(B\operatorname{End}(y);\mathbb Q).
\]
Its first horizontal differential maps these two summands
individually, with the boundary signs, onto the corresponding
corner summands in horizontal degree zero.  Every object
has an incident edge because the support is connected and
has at least two vertices.  Thus $E^2_{0,2}=0$.
The only term of total degree two which remains is
$E^2_{2,0}=H_2(B\Pi;\mathbb Q)$.
No later differential enters it, and its only possible
outgoing differential lands in the already zero term
$E^2_{0,1}$.  Comparison under the actual functor $q$
therefore gives an isomorphism on second rational homology,
and duality gives the displayed cohomology isomorphism.
Naturality makes this a cup-product pullback, not merely
an abstract identification of vector spaces.
\end{proof}

\begin{lemma}[Small cycle rank of a support graph]
\label{C:small-cycle-support}
Let $\Pi$ be a finite connected poset, let $G$ be its
undirected comparability graph, and put
$R=|E(G)|-|V(G)|+1$.
If $R\leq6$, then $H_4(B\Pi;\mathbb Q)=0$.
If $R\leq3$ and $H_1(B\Pi;\mathbb F_2)=0$, then $G$ is
chordal.  In the latter case its maximal-clique decomposition
has the following possibilities after bridges and cut vertices
are separated:
\begin{itemize}
\item at $R=1$, a single triangle;
\item at $R=2$, triangles or two triangles sharing an edge;
\item at $R=3$, at most three triangles glued along edges or
vertices in a clique tree, or one complete four-vertex block.
\end{itemize}
The classifying-space nerve of a category with this support
is the corresponding union of its full clique subcategory
nerves, with their actual full intersection subcategories.
\end{lemma}
\begin{proof}
The order complex of a poset is the clique complex of its
comparability graph: a set of pairwise comparable objects
is a chain.  A four-simplex requires a complete graph on
five vertices, which already has cycle rank six.  If
$R\leq6$ and such a graph occurs, contracting that complete
graph leaves a tree; otherwise the cycle rank would exceed
six.  Consequently this is the only four-simplex, and
there is no five-simplex.  Its nonzero simplicial boundary
shows that fourth homology is zero.  If no such complete
graph occurs, there are no four-simplices at all.

Suppose now $R\leq3$ and $H_1(B\Pi;\mathbb F_2)=0$.
The cycle space of $G$ has dimension $R$ and is spanned
by triangle boundaries.  Choose a basis of at most three
such boundaries.  If an induced cycle of length at least
four existed, each triangle would contain at most one
edge of that cycle: two consecutive edges would give a
forbidden chord.  A sum of at most three basis triangles
therefore cannot contain all the edges of that cycle.
This contradicts the basis property, proving chordality.

The clique-tree description of a chordal graph follows
by removing a simplicial vertex successively.  In the
present small range it can also be read directly from
the cycle budget.  A complete four-vertex block consumes
all three cycles, so its remaining attachments are trees.
Otherwise every nontrivial maximal clique is a triangle;
adjoining a triangle along an edge or a vertex adds one
independent cycle.  There can be at most three triangles.
Two triangles sharing an edge are the diamond graph;
triangles joined only at a vertex belong to different
blocks.  These give exactly the alternatives listed.

Finally the distinct objects used by a nerve simplex
form a chain, hence a clique of $G$.  The simplex therefore
lies in a full maximal-clique subcategory.  Intersections
of these nerves are the nerves of the full intersections.
All inclusions are simplicial subcomplex inclusions, so
their realizations are the indicated homotopy pushouts.
\end{proof}

\begin{theorem}\label{C:acyclic-arrow-excess-six}
Let $C$ be a finite category with $k$ objects, $n$ arrows,
and acyclic support.  If $n\leq2k+5$, then
$BC\not\simeq\mathbb CP^2$.
In particular every category with acyclic support, at
least eight objects, and at most twenty-one arrows is
excluded.
\end{theorem}
\begin{proof}
Suppose to the contrary that $BC\simeq\mathbb CP^2$.
The category is connected.  If $k\leq2$, then $n\leq9$,
contrary to Theorem~\ref{C:CP2-lower16}.  Thus $k\geq3$.
Use the notation
\[
 \begin{split}
 D&=\sum_x(|\operatorname{End}(x)|-1),\\
 P&=\sum_{x<y}(|C(x,y)|-1),\qquad
 R=E-k+1,
 \end{split}
 \qquad D+P+R=n-(2k-1)\leq6,
\]
where the sum is over nonempty cross hom-sets and $E$
counts the edges of the comparability graph.

If $P=0$, Lemma~\ref{C:singleton-support-H2} pulls every
rational degree-two class back from the support poset.
Lemma~\ref{C:small-cycle-support} gives zero fourth
rational homology there.  Its degree-two products, and
therefore their pullbacks, are zero, a contradiction.

If all corner algebras are semisimple over $\mathbb Q$,
the relative bar over their product has no positive
internal Tor.  Since $P\geq1$ gives $R\leq5$, no chain
has five distinct objects.  The relative bar has strict
object degree at most three, so $H_4(BC;\mathbb Q)=0$.
We may suppose henceforth that a nonsemisimple corner
occurs.  An order-one or order-two monoid algebra is
semisimple, so $D\geq2$.

\paragraph{Tree support.}
Suppose $R=0$.  There is no composable pair of strict
arrows, since their composite would create a triangle.
If all corners have zero second rational cohomology,
Corollary~\ref{C:height-one-zero-H2-products} applies.
Otherwise there is exactly one corner $M$ which can have
nonzero second rational cohomology.  Indeed such a corner
has at least five elements by
Theorem~\ref{E:all-six-types}, and two would cost $D\geq8$.
Write $x$ for its object.  If any incident hom-set is a
singleton, its two-object collage is a join and has zero
second cohomology.  Every degree-two class of $BC$ then
restricts to zero on $BM$, and also on every other corner.
Collapsing the corner spaces in the height-one
correspondence pushout pulls all those classes back
from a wedge of suspensions.  Their products are zero.

Thus every edge incident with $x$ has at least two arrows.
If $d$ is its degree in the support tree, then
$d\leq P\leq6-D\leq2$.  When $d=2$, equality forces
$|M|=5$, $D=4$, and two incident hom-sets of size two;
all other corners and hom-sets are singletons.  The
five-element corner with nonzero $H^2$ is the
identity-adjoined $2\times2$ rectangular band.  The two
incident arrows have the same direction, since the
support has height one.  Identifying their two trivial
endpoint objects gives the ordinary four-state act
cone over $M$, or its opposite.  On realization this
is the identification of two points in a connected
CW space, so it induces an isomorphism on cohomology
in degrees at least two.  The act-cone products are
zero by Theorem~\ref{C:two-row-act-cups}; their pullbacks
give zero products on this two-leg star.

When $d=1$, the single incident hom-set has size at most
three.  Size two is covered by
Theorem~\ref{C:two-cross-arrows-universal}.  Size three
forces $|M|=5$, all other corners trivial, and is an
ordinary act cone over the same two-row band monoid;
Theorem~\ref{C:two-row-act-cups} again applies.
In both cases all remaining corners have order at most
two.  The other pieces of the tree have semisimple
diagonal and only one strict crossing, hence zero
rational homology in degrees at least two.  They attach
to the displayed collage or star only at corners of
order at most two.  These intersections are rationally
acyclic.  Mayer--Vietoris therefore makes restriction
in degree four injective and identifies degree-two
classes with their restrictions to the displayed piece.
The square vanishes on that piece and hence on all of
$BC$.  This excludes tree support.

We now have $P,R\geq1$.  Therefore
\[
 2\leq D\leq4,\qquad R\leq3,
\]
and every corner has at most five elements.  All such
corners have $H_3(-;\mathbb Q)=0$ by
Theorem~\ref{E:all-six-types}; those of order at most
four are rationally acyclic.
The actual functor $C\to\Pi$ is onto on fundamental
groups.  Indeed an edge loop in $B\Pi$ lifts by choosing
one actual arrow for each oriented edge, with formal
reversals when necessary.  Since $BC$ is simply
connected, $B\Pi$ is simply connected as well.
Lemma~\ref{C:small-cycle-support} now gives the stated
chordal decompositions.

We repeatedly use two elementary consequences of their
actual nerve pushouts.  A tree of pieces meeting at
rationally acyclic corners has positive rational
homology, in degrees at least two, equal to the direct
sum of the pieces' homology.  Thus each piece inherits
bounded rational homology and second-cohomology
dimension at most one from $BC$.  More generally, if
every intersection has zero third rational homology,
Mayer--Vietoris makes the map from the sum of fourth
homology of the pieces onto fourth homology of the
union surjective.  Equivalently, restriction in fourth
cohomology is injective.  Vanishing of squares on the
pieces then detects vanishing on the union.

\paragraph{One independent cycle.}
Let $R=1$.  The only non-bridge block is a triangle.
If $D\leq3$, all corner intersections are rationally
acyclic.  There is at most one corner of order at least
three, and it has order at most four; all other corners
have order at most two.
For a bridge piece, semisimple corners give zero $H^2$.
Otherwise exact coinvariants at its order-at-most-two
opposite corner give an actual rational cohomology-ring
equivalence to an ordinary act cone over a monoid of
order at most four.  To see the comparison directly,
the edge double bar computes the exact coinvariants
at that corner, and the two vertex bars have the same
rational augmentation.  The double-mapping-cylinder
comparison is consequently a rational equivalence.
The boundedness and rank-one condition inherited by
this act cone permit Lemma~\ref{C:nine-bounded-rank-one},
so its square is zero.

For the triangle, if its larger corner is an endpoint,
that corner has $H^2=0$ and the opposite two-object
block has semisimple diagonal and $H^2=0$.
Lemma~\ref{C:two-block-zero-H2-products} kills its
degree-two products.  If the larger corner is in the
middle, its endpoints are semisimple and its order is
at most four.  The inherited boundedness and
Corollary~\ref{C:semisimple-ends-middle-four} give
zero fourth homology.  A triangle with no larger corner
has semisimple diagonal and zero fourth homology directly.
Fourth-cohomology restriction now excludes a square
on the whole category.

If instead $D=4$, the budget forces $P=1$.  Every bridge
has one or two arrows, so its products vanish by the
join description or Theorem~\ref{C:two-cross-arrows-universal}.
The triangle has at most one doubled hom-set.  If all
its corners have order at most four, at least one
adjacent hom-set is a singleton.  That adjacent block
is a join with zero $H^2$, and the opposite corner
also has zero $H^2$.  The two-block lemma applies.
The only remaining possibility is a five-element corner,
with the other two corners trivial.  If it is the middle
corner, a singleton adjacent hom-set gives the usual
Quillen-A deletion of that middle object, leaving a
graph.  If it is an endpoint and its adjacent hom-set
is a singleton, use the same zero-$H^2$ two-block cut.
If its adjacent hom-set is doubled, the other endpoint
is terminal or initial, so the triangle is contractible.
All corner intersections still have zero $H_3$.
Thus restriction in degree four is injective and again
every global square is zero.

\paragraph{Two independent cycles.}
Let $R=2$.  Now $D\leq3$ and $D+P\leq4$, so all
corner intersections are rationally acyclic.  Bridge
and triangle blocks are handled exactly as in the
preceding $D\leq3$ argument.  The only new block is
a diamond: two triangles sharing an edge $u<v$, with
exclusive vertices $a,b$.  These vertices are
incomparable.  Transitivity forces them to occupy the
same position relative to $u<v$: both below, both
above, or both strictly between.

Call an edge potentially nonzero if its hom-set has
at least two arrows and one endpoint corner has at
least three elements.  Every other edge collage has
$H^2=0$, by the singleton join or semisimple diagonal
argument.  A potentially nonzero edge uses at least
three units of the combined $D+P$ budget.

Suppose first that $a,b<u<v$, the opposite case being
dual.  Cut the diamond into the full blocks $\{a,b\}$
and $\{u,v\}$.  The first has zero $H^2$.  If the shared
edge is not potentially nonzero, the second does too.
Otherwise the shared edge consumes at least three of
the at most four available units.  At least one of
$a,b$ consequently has a trivial corner and singleton
hom-sets to both $u,v$.  It is an initial object in
its full triangle with $u,v$.  Hence every global
degree-two class restricts to zero on the shared-edge
block, as well as on the first block.  The actual
two-block suspension quotient in the proof of
Lemma~\ref{C:two-block-zero-H2-products} pulls back
all these classes, so their products vanish.

Suppose instead $u<a,b<v$.  Consider the two stars
on $\{u,a,b\}$ and $\{a,b,v\}$.  A star has $H^2=0$
if neither edge is potentially nonzero: it is a union
of its two edge collages along a corner with zero
first and second rational homology.  If either star
has zero $H^2$, cut it from the opposite endpoint
corner and apply the two-block lemma.
Avoiding both cuts therefore requires a potentially
nonzero edge in each star.  There is at most one
corner of order at least three, so those edges must
meet at the same exclusive vertex, say $a$.
Their cost is at least two corner excesses and two
parallel-arrow excesses.  Equality with $D+P\leq4$
is forced: the corner at $a$ has order three,
$|C(u,a)|=|C(a,v)|=2$, and every other corner and
cross hom-set in the diamond is a singleton.

The full triangle on $u,b,v$ and its intersection
with the full triangle on $u,a,v$ are contractible;
the intersection is the singleton edge $u\to v$.
The nerve is their actual union, so the inclusion of
the latter triangle induces a homotopy equivalence.
It has semisimple endpoints and middle order three.
Its inherited bounded homology gives zero fourth
homology by Corollary~\ref{C:semisimple-ends-middle-four}.
This closes the diamond.  The corner-intersection
comparison then closes the whole $R=2$ case.

\paragraph{Three independent cycles.}
Finally $R=3$ forces $D=2$ and $P=1$.  There is one
order-three corner, every other corner is trivial,
and exactly one cross hom-set is doubled.  In a
bridge or clique intersection, the only possible
nontrivial corner is this order-three monoid.
Its trivial modules on both sides have projective
dimension at most one, as proved in
Proposition~\ref{C:semisimple-ends-small-middle}.
The edge bar therefore has no internal homology in
degree at least two, and its collage has no rational
homology in degrees at least three.  In particular
every clique intersection has zero third homology.

If the nontrivial cliques are triangles, each has
zero fourth homology.  With the order-three corner
at an endpoint, internal bar degree is at most one
and strict degree at most two.  With that corner
in the middle, at least one adjacent hom-set is a
singleton; the corresponding middle-object deletion
leaves a graph.  With all corners trivial the nerve
is two-dimensional.  Fourth-cohomology restriction
is injective across the clique tree, so this case
has zero fourth homology altogether.

It remains to consider the complete four-vertex
block, whose order is $0<1<2<3$.  A three-object
subchain with trivial corners and at most one doubled
hom-set has zero second homology: its triangle
boundary map is injective, since at least one
adjacent hom-set is a singleton, and the component
on the other adjacent edge distinguishes its
two-simplex labels.
If the order-three corner is outside this block, the
block has only trivial corners and its nerve has
dimension at most three, so its fourth homology is zero.
Otherwise argue as follows.
If the order-three corner is $0$ or $3$, separate
it from the other three objects.  Both full blocks
have zero second cohomology, so the two-block lemma
applies.
If that corner is $1$, first cut into $\{0,1\}$ and
$\{2,3\}$.  This is a zero-$H^2$ cut unless the
doubled hom-set is $C(0,1)$.  In that exception, cut
instead into $\{0\}$ and $\{1,2,3\}$: the second
block has all cross hom-sets singleton and hence
zero $H^2$ by Lemma~\ref{C:singleton-support-H2},
since its support is a simplex.  The case of corner
$2$ is the opposite argument.  Thus every square
vanishes on the four-vertex block.  Its remaining
attachments are bridges with zero third-homology
intersections, so restriction in degree four proves
that every global square vanishes as well.

All possible values of $P,R,D$ have now been covered.
Each contradicts the nonzero rational square of
$\mathbb CP^2$, proving the theorem.  For $k\geq8$
and $n\leq21$ one has $n-(2k-1)\leq6$, giving the
stated consequence.
\end{proof}

% D proof; root and B independent literal audits passed.
\subsection{Seven excess arrows with acyclic support}

\begin{lemma}[Two support-homology budgets]
\label{C:flag-second-cycle-budget}
Let $\Pi$ be a finite connected poset and let $R$ be the cycle
rank of its undirected comparability graph.  If $R\leq6$, then
$H_2(B\Pi;\mathbb Q)=0$.  If $R\leq8$, then
$H_4(B\Pi;\mathbb Q)=0$.
\end{lemma}
\begin{proof}
The order complex is the flag complex of the comparability graph.
Suppose it has nonzero second homology, and choose an induced
subcomplex with this property and the fewest vertices.  It is
connected.  Deleting any vertex kills its second homology.
The Mayer--Vietoris sequence for the deletion and the vertex
star therefore injects its nonzero second homology into the
first homology of the vertex link.  A flag complex with nonzero
first homology has at least four vertices: on at most three
vertices every graph cycle is a filled triangle.  Every vertex
of our induced graph consequently has degree at least four.
There are at least six vertices, since a five-vertex graph of
minimum degree four is complete and its flag complex is a simplex.
Writing $v,e$ for these vertex and edge counts gives
$e-v+1\geq v+1\geq7$.  Graph cycle rank cannot increase on
passing to a subgraph, proving the first assertion.

A four-simplex requires a complete five-vertex subgraph.
Two different such subgraphs have union of cycle rank at least
nine: the smallest value occurs when they share four vertices,
and then the union has six vertices and fourteen edges.
Disjoint complete subgraphs already contribute twelve cycles.
Thus $R\leq8$ permits at most one four-simplex and no
five-simplex.  The boundary of that one four-simplex, if present,
is nonzero, so fourth homology vanishes.
\end{proof}

\begin{lemma}[A cut through the parallel hom-sets]
\label{C:parallel-hom-cut}
Let $C$ be a finite category with acyclic support of cycle rank
at most six.  Suppose every corner has zero rational second
cohomology.  Call a nonempty cross hom-set parallel if it has
at least two arrows.  Suppose no two parallel hom-sets
$C(x,y)$ and $C(z,w)$ satisfy $y\leq z$ in the support order.
Then every product of two rational degree-two classes on $BC$
vanishes.
\end{lemma}
\begin{proof}
If there are no parallel hom-sets, apply
Lemma~\ref{C:singleton-support-H2} componentwise and
Lemma~\ref{C:flag-second-cycle-budget}.
Otherwise let $A$ be the downward closure of all sources of
parallel hom-sets, and let $D$ be its complement.  The stated
condition says precisely that every parallel hom-set crosses
from $A$ to $D$.  There are no arrows from $D$ to $A$.
Every connected component of either full block having at least
two objects has singleton cross hom-sets.  Its second
cohomology is that of its support by
Lemma~\ref{C:singleton-support-H2}, and is zero by the first
part of Lemma~\ref{C:flag-second-cycle-budget}.  A single-object
component has zero second cohomology by hypothesis.  Both
blocks therefore have zero second cohomology, and
Lemma~\ref{C:two-block-zero-H2-products} applies.

We shall also use the same proof when a corner has nonzero
second cohomology but the restriction of every global
degree-two class to that corner is zero.  Every global class
then restricts to zero on both full blocks, and the actual
two-block suspension quotient still pulls back all those
classes.  No vanishing assertion about that corner is needed.
\end{proof}

\begin{theorem}\label{C:acyclic-arrow-excess-seven}
Let $C$ have $k$ objects, $n$ arrows, and acyclic support.
If $n\leq2k+6$, then $BC\not\simeq\mathbb CP^2$.
In particular the acyclic seven-object case is excluded
through twenty arrows.
\end{theorem}
\begin{proof}
Suppose $BC\simeq\mathbb CP^2$.  We may assume that $C$ is
connected and $k\geq3$, since the cases $k\leq2$ have at
most ten arrows and are covered by the global lower bound.
Use the corner, parallel-arrow, and support excesses
\[
 D=\sum_x(|\operatorname{End}(x)|-1),\qquad
 P=\sum_{x<y}(|C(x,y)|-1),\qquad R=E-k+1.
\]
Thus $D+P+R=n-(2k-1)\leq7$.
The support functor is onto on fundamental groups, by lifting
edge loops to actual arrows and formal inverses.  Hence its
order complex has zero first homology over $\mathbb F_2$.

If $P=0$, every degree-two class is pulled back from the
support poset by Lemma~\ref{C:singleton-support-H2}.
Its fourth rational homology is zero by
Lemma~\ref{C:flag-second-cycle-budget}, so all these squares
vanish.

Suppose all corner algebras are semisimple.  If $R\leq5$,
the relative bar has strict-object length at most three,
so fourth homology vanishes.  The only remaining budget is
$D=0$, $P=1$, $R=6$.  There is at most one five-object
chain, and all corners are trivial.  A nondegenerate
four-simplex uses its four successive arrows.  If the unique
doubled hom-set is not one of these four, there is one such
simplex and its boundary is nonzero.  Otherwise there are
two such simplices.  Choose a face deleting a vertex which
is not an endpoint of the doubled successive hom-set.
That face retains its arrow label and distinguishes the two
simplices.  No other four-simplex can cancel either face,
since there is no other five-object chain.  In either case
the degree-four boundary is injective.  Thus $H_4=0$ here too.

We may now assume a nonsemisimple corner, so $D\geq2$.
If $P=1$, there is only one parallel hom-set, of size two,
and $R\leq4$.  At most one corner has nonzero second
rational cohomology, since two such corners would cost
$D\geq8$.  If none does, Lemma~\ref{C:parallel-hom-cut}
applies.  If the exceptional corner has a singleton incident
hom-set, that two-object subcategory is a join, so every
global degree-two class restricts to zero at the corner.
The final paragraph of that lemma again applies.
Otherwise the corner is a leaf on the unique doubled edge.
Its other endpoint has order at most four, hence is
rationally acyclic.  The remainder of the category has all
cross hom-sets singleton and support cycle rank at most six;
its second cohomology is zero.  The leaf edge has zero
degree-two products by
Theorem~\ref{C:two-cross-arrows-universal}.  The actual
nerve union meets at the rationally acyclic other corner,
so fourth-cohomology restriction is injective.  The global
square is therefore zero.

Henceforth $D,P\geq2$, and consequently $R\leq3$.
The support is chordal by Lemma~\ref{C:small-cycle-support}.
Its bridge and clique pieces are actual full subcategories,
whose nerves form the corresponding tree of unions.
We shall use rationally acyclic corner intersections to
inherit bounded homology and degree-two rank at most one
on a piece.  When an intersection merely has zero third
homology, fourth-cohomology restriction is still injective.
For boundedness of a piece it is enough that the intersections
have bounded rational homology.  All these statements follow
from the Mayer--Vietoris sequence of the actual nerve unions.

\paragraph{Tree support.}
If $R=0$ and every corner has zero second cohomology,
Corollary~\ref{C:height-one-zero-H2-products} applies.
There can otherwise be only one corner $M$ with nonzero
second cohomology.  A singleton incident edge kills every
global restriction to that corner, as before; the
height-one suspension quotient then kills all global products.
We may assume every incident edge has at least two arrows.
Since $P\geq2$ and $D+P\leq7$, this corner has order at
most six and every other corner has order at most two.
All its incident arrows have the same direction.

The other pieces of the support tree have semisimple
corners, one strict crossing, and no homology in degrees
at least two.  They meet the central star at rationally
acyclic corners.  The star thus inherits bounded rational
homology and degree-two rank at most one.  Exact coinvariants
at its order-at-most-two leaves give an actual rational
cohomology-ring equivalence to a star with trivial leaves.
Identify these leaf points.  The result is an ordinary
act cone over $M$; on realization the identification adds
only circle summands.  It preserves cohomology and products
in degrees at least two and preserves boundedness.
Lemma~\ref{C:nine-bounded-rank-one}, with $|M|\leq6$,
kills the square.  Fourth-cohomology restriction across
the tree then kills the global square.

\paragraph{One independent cycle.}
Let $R=1$.  Here $D+P\leq6$, so $D\leq4$ and
every corner has order at most five.  Such corners have
zero third rational homology.  The only non-bridge block
is a triangle, and fourth-cohomology restriction to it
and the bridges is injective.

A bridge whose endpoint corners have zero second
cohomology has zero degree-two products by the two-block
lemma, regardless of its hom-set size.  If an endpoint
has nonzero second cohomology, it is the five-element
identity-adjoined $2\times2$ band, every other corner is
trivial, and the bridge is an ordinary act cone over this
two-row monoid.  Its products vanish by
Theorem~\ref{C:two-row-act-cups}.

If $D\leq3$, a triangle has at most one corner of order
at least three, of order at most four.  At an endpoint,
cut this zero-$H^2$ corner from the opposite semisimple
two-object block.  In the middle, use inherited boundedness
and Corollary~\ref{C:semisimple-ends-middle-four}.
These give zero products or zero fourth homology.

Suppose $D=4$, so $P\leq2$.  If there is a five-element
corner, all others are trivial.  At the middle, a singleton
adjacent hom-set deletes that middle object; otherwise both
adjacent sets have size two and the long set is a singleton.
Triangle compression then gives a monoid with a $2\times2$
minimal band and zero degree-two products.
At an endpoint with zero $H^2$, the opposite two-object
block is semisimple and the two-block lemma applies.
At an endpoint with nonzero $H^2$, the monoid is the
five-element two-row band monoid.  If its adjacent hom-set
is a singleton, that adjacent join and the other endpoint
give a zero-$H^2$ cut.  If both adjacent sets are parallel,
they have size two and the long set is a singleton; the
trivial extreme object makes the triangle an objectwise
one- or two-state act cone, covered by
Lemma~\ref{C:two-states-objectwise-act-cups}.
In the other cases the opposite adjacent set is a singleton,
and its trivial endpoint permits the actual middle-object
deletion.  The resulting ordinary act cone has zero products
by Theorem~\ref{C:two-row-act-cups}.

It remains that all triangle corners have order at most
four.  If an adjacent hom-set is a singleton, its join
block and the opposite zero-$H^2$ corner give the cut.
Otherwise the two adjacent sets have size two and the
long set has size one.  A trivial endpoint makes the
triangle an objectwise one- or two-state categorical act
cone, so Lemma~\ref{C:two-states-objectwise-act-cups}
applies.  The same conclusion holds after exact rational
endpoint coinvariants if either endpoint has order two.
The only remaining corner orders are $(3,1,3)$.
Both trivial modules of an order-three monoid have
projective dimension at most one.  If either relevant
endpoint trivial module is projective, the strict-object
bar has total degree at most three and $H_4=0$.
Otherwise the source is the identity-adjoined two-element
left-zero band and the target its right-zero opposite.
On a two-element adjacent act the graph model in
Lemma~\ref{C:right-zero-action-graph} has nonzero first
homology only for the common constant retraction, when
it is the one-edge circle.  If either adjacent graph has
zero first homology, the triangle filtration has no total
degree-four class.  If both are circles, the category is
exactly the family of
Theorem~\ref{C:affine-chart-family-zero}.
Thus its products vanish in every case.

\paragraph{Two independent cycles.}
Let $R=2$.  Now $D+P\leq5$ and $D\leq3$.
Every corner is rationally acyclic.  Bridge and triangle
blocks are treated as above.  If $D=2$, the only nontrivial
corner has order three.  Every edge intersection then has
zero homology in degrees at least three, by the projective-
dimension-one bound for its trivial modules.  Restriction
in degree four to the triangle pieces is injective, and
each piece inherits bounded high homology.  An endpoint
order-three corner gives total bar degree at most three;
a middle order-three corner gives zero fourth homology by
Corollary~\ref{C:semisimple-ends-middle-four}.  This closes
the $D=2$ case, including its possible third parallel-arrow
excess.

We may therefore assume $D=3$, so $P\leq2$.
Consider a diamond with
shared edge $u<v$ and exclusive vertices $a,b$.
As in the excess-six proof, an edge can have nonzero
$H^2$ only if it is parallel and has an endpoint of
order at least three.  There is at most one such larger
corner in the whole category.

If $a,b<u<v$, cut $\{a,b\}$ from $\{u,v\}$.
Unless the shared edge is potentially nonzero, both
blocks have zero $H^2$.  If it is potentially nonzero,
it uses one of the at most two parallel-arrow excesses.
At least one of $a,b$ therefore has singleton hom-sets
to both $u,v$.  Its full triangle is the join of its
corner space with the shared-edge space.  This join has
zero $H^2$, even if that exclusive corner is not trivial.
Every global degree-two class hence restricts to zero
on the shared edge, and the same two-block suspension
argument applies.  The case $u<v<a,b$ is dual.

If $u<a,b<v$, a zero-$H^2$ star on either side gives
a cut from the opposite endpoint.  Avoiding both cuts
requires a potentially nonzero edge on each side.
They must meet at the unique larger exclusive corner,
say $a$.  Both its adjacent hom-sets have size two;
all other cross hom-sets in the diamond are singleton.
The corner at $a$ has order three or four, and at most
one other corner has order two.  In particular the
shared edge and the other triangle have contractible
classifying spaces: all their cross sets are singleton
and at least one of their endpoint corner spaces is a
point.  The full diamond is homotopy equivalent to the
triangle on $u,a,v$.  This triangle has semisimple
endpoints, middle order at most four, and inherited
bounded homology.  Its fourth homology is zero.

\paragraph{Three independent cycles.}
Finally $R=3$ forces $D=2$, $P=2$: one order-three
corner and all other corners trivial.  Every edge
intersection has zero homology in degrees at least
three, because the relevant trivial module of that
corner has projective dimension at most one.
Hence fourth-cohomology restriction to the clique
pieces is injective, and each piece inherits bounded
homology in sufficiently high degrees.

For triangle pieces an endpoint order-three corner
gives total bar degree at most three.  A middle
order-three corner gives zero fourth homology by
Corollary~\ref{C:semisimple-ends-middle-four} and the
inherited boundedness.  A triangle with trivial
corners is two-dimensional.
For a complete four-object block, if the only parallel
hom-set has size three, Lemma~\ref{C:parallel-hom-cut}
applies to that block.  Otherwise every cross hom-set
has at most two arrows.  At least one of its two extreme
corners is trivial.  It is therefore an objectwise
one- or two-state act cone over the other three objects,
or the opposite of such a cone.  Its products vanish
by Lemma~\ref{C:two-states-objectwise-act-cups}.
If the order-three corner lies outside the block, its
nerve has dimension at most three and $H_4=0$ directly.
Restriction now kills the global square.

Every budget case gives a zero rational square or zero
fourth rational homology, contradicting the projective
plane.  This proves the theorem.
\end{proof}

% Independently audited in full by root and B after the four precision repairs.
\subsection{Eight excess arrows with acyclic support}

\begin{lemma}[The rank-one form of the small-corner lemma]
\label{C:small-corner-two-stalks-rank-one}
Let a finite category $S$ be the source block of a collage
whose sink corner $N$ has at most four elements.  Suppose
every cross stalk is nonempty and has at most two elements.
If the collage has rational second-cohomology dimension at
most one, its degree-two square vanishes.  The opposite
statement holds for a small source corner.
\end{lemma}
\begin{proof}
This is the homological scope of the proof of
Lemma~\ref{C:small-corner-two-states-per-object}.
For a projective trivial right $\mathbb QN$-module, the
actual orbit quotient is a rational ring equivalence to
an objectwise one- or two-state categorical act cone.
All its positive products vanish by
Lemma~\ref{C:two-states-objectwise-act-cups}.
Otherwise the minimal ideal of $N$ is a right-zero band
of size two or three.  Its retractions coincide on each
one- or two-element stalk.  Their common image is an actual
biset retract on which $N$ acts trivially.  Since $BN$ is
contractible, that core collage is equivalent to an
objectwise one- or two-state categorical act cone and
has zero products.  The relative quotient is a suspension.
If the core has nonzero second cohomology, its retract
pullback spans the assumed rank-one space; otherwise every
degree-two class lifts from the suspension quotient.
In both cases the square vanishes.  These arguments use
only the stated rank bound, not bounded homology or finite
CW type of the collage.
\end{proof}

\begin{lemma}[Two parallel-arrow excesses with small corners]
\label{C:two-parallel-small-corners}
Let $C$ be a connected finite category with acyclic support,
all corner orders at most four, support cycle rank at most
six, and total parallel-arrow excess at most two.
If $\dim H^2(BC;\mathbb Q)\leq1$, its square vanishes.
\end{lemma}
\begin{proof}
All its corners are rationally acyclic by the small-monoid
classification.  If there is at most one parallel hom-set,
or if its two parallel hom-sets cannot occur serially,
Lemma~\ref{C:parallel-hom-cut} applies.
Thus there are exactly two doubled hom-sets
$C(a,b)$ and $C(c,d)$ with $b\leq c$.

If $b<c$, let $A$ be the full downset of $c$.
Every hom-set from $A\setminus\{c\}$ to $c$ is a singleton:
neither of the two doubled hom-sets ends at $c$.
Thus $BA$ is the join of the predecessor subcategory with
the corner at $c$, and has zero second rational cohomology.
The complementary full subcategory has only singleton
cross hom-sets and has zero second cohomology by
Lemmas~\ref{C:singleton-support-H2} and
\ref{C:flag-second-cycle-budget}, including its isolated
corners.  The two-block suspension lemma proves the claim.

It remains that the two doubled hom-sets are
$C(a,b)$ and $C(b,c)$.  Suppose $b$ has a neighbor
$x>b$ other than $c$.  Let $A$ be the full downset of $b$.
All hom-sets $A\to x$ are singleton, so the full
subcategory on $A\cup\{x\}$ is a join and has zero
second cohomology.  Every global degree-two class
therefore restricts to zero on $A$.  The complement
of $A$ has only singleton cross hom-sets and zero
second cohomology as above.  The actual two-block
quotient again kills the square.  A neighbor $x<b$
other than $a$ gives the opposite argument with the
upset of $b$.

Otherwise $b$ has precisely the two neighbors $a,c$.
Let $T$ be the full triangle on $a,b,c$ and let
$D=C\setminus\{b\}$.  Their actual nerve union is all
of $BC$, and their intersection is the singleton edge
$a\to c$.  Its classifying space is the join of two
rationally acyclic corner spaces, hence rationally acyclic.
The subcategory $D$ has only singleton cross hom-sets and
zero second cohomology.  Mayer--Vietoris identifies
$H^2(BT;\mathbb Q)$ with the global group, and makes
fourth-cohomology restriction injective.
The triangle has hom-set sizes $(2,2,1)$.  Its last
corner has order at most four and both incoming stalks
have size at most two.  The preceding lemma kills its
rank-one square; the restriction to $D$ is already zero.
This proves the assertion.
\end{proof}

\begin{theorem}\label{C:acyclic-arrow-excess-eight}
Let $C$ have $k$ objects, $n$ arrows, and acyclic support.
If $n\leq2k+7$, then $BC\not\simeq\mathbb CP^2$.
In particular every acyclic seven-object category with
at most twenty-one arrows is excluded.
\end{theorem}
\begin{proof}
Assume $BC\simeq\mathbb CP^2$, so $C$ is connected.
The cases $k\leq2$ have at most eleven arrows and are
already excluded.  Use the notation $D,P,R$ of
Theorem~\ref{C:acyclic-arrow-excess-seven}.
That theorem treats $D+P+R\leq7$, so it remains to
consider equality
\[
 D+P+R=8.
\]
As before the support functor is onto on fundamental
groups, and its support order complex has zero first
homology over $\mathbb F_2$.

If $P=0$, the actual support functor pulls back every
degree-two class and its target has zero fourth homology
by Lemma~\ref{C:flag-second-cycle-budget}.
If the corners are semisimple and $R\leq5$, the relative
bar has strict-object length at most three and $H_4=0$.
For semisimple corners and $R\geq6$, one has
$D+P\leq2$.  The support has at most one five-object
chain.  Choose a face of that chain deleting a vertex
which is not an endpoint of any parallel hom-set used
by its top simplex, nor a nontrivial corner.
Such a vertex exists: for $D=0$ at most two parallel
hom-sets exclude at most four vertices; for $D=1$ there
is at most one parallel hom-set and one nontrivial
corner, excluding at most three vertices.

The selected face retains all varying arrow labels.
Its remaining composition involves only singleton
hom-sets, except possibly a fixed singleton product
landing in a larger hom-set.  In that event the unique
source basis vector is sent to one target basis vector,
so this factor is still an injection.  The resulting
component of the top boundary is injective.  If the one nontrivial
semisimple corner occurs, it is retained by this face;
the component is the identity on the retained balanced
tensor factors, tensored with the just-described basis
injection.  Exact semisimple coinvariants preserve this
injection.  There is no other five-object
chain to contribute to the same boundary component.
Thus the degree-four relative-bar boundary is injective
and $H_4=0$ in all semisimple cases.

We may now assume $D\geq2$ and a nonsemisimple corner.
If $P=1$, the parallel-cut and exceptional-leaf argument
of the excess-seven proof applies unchanged.  There is
at most one corner with nonzero $H^2$, since $D\leq7$.
A singleton incident edge kills its global restriction.
Without such an edge it is a leaf on the sole doubled
hom-set, whose products vanish universally.  Its other
endpoint has order at most four and is rationally acyclic,
and the remaining singleton-hom category has zero $H^2$
since $R\leq5$.  Fourth-cohomology restriction finishes
this case.  Henceforth $D,P\geq2$, so $R\leq4$.

\paragraph{Tree support.}
When $R=0$, zero second cohomology at every corner gives
the height-one suspension obstruction.  There is otherwise
a unique corner $M$ with nonzero $H^2$, of order at most
seven.  A singleton incident edge again kills all global
vertex restrictions, so assume every incident edge is parallel.

If every other corner has order at most two, group all
edges incident with $M$ into its central star.  All remaining
pieces have semisimple diagonal and one strict crossing;
they have no homology in degrees at least two, and meet
the star at rationally acyclic corners.  The star inherits
bounded rational homology and second-cohomology rank at
most one.  Exact leaf coinvariants and identification of
the trivial leaf points give an ordinary act cone over
$M$, changing only first homology.  The bounded rank-one
act lemma for $|M|\leq9$ kills its square.

The only other budget is $D=6$, $P=2$, with the positive
corner the five-element sphere monoid and one other corner
of order three.  If the positive corner has degree two,
both incident hom-sets have size two.  Their individual
collages have zero products universally, and their common
five-element corner has $H_3(-;\mathbb Q)=0$.
Restriction in degree four is therefore injective on
their union.  Every remaining piece attaches at a
rationally acyclic corner and has zero products.
If the positive corner has degree one, its hom-set has
size two or three.  Size two is universal.  Size three
is either an ordinary act over the two-row sphere monoid,
or its three-state biset with the order-three corner.
For the latter, a projective trivial module at the small
opposite corner gives an ordinary act after exact
coinvariants.  The only nonprojective order-three corner
is the identity-adjoined two-element right-zero band,
and Corollary~\ref{C:P5-right-zero-three-all-states}
gives zero products universally.  No actual
projective-plane type of the piece is assumed.
This closes tree support.

\paragraph{One independent cycle.}
Let $R=1$.  The support consists of its one triangle
and bridge attachments.  If all corners have order at
most five, their third rational homology is zero, so
fourth-cohomology restriction to these pieces is injective.
When all corners have order at most four the intersections
are rationally acyclic: each piece inherits bounded homology
and $b_2\leq1$.

For $D\leq3$, bridge products vanish by the two-block
lemma, and triangle pieces are handled by the same
endpoint zero-$H^2$ cut or middle-order-at-most-four
bounded-Tor argument as in the excess-seven proof.
For $D=4$ or $5$, suppose first all corners have order
at most four.  Here respectively $P=3$ or $2$.
A singleton adjacent hom-set gives a join block and
a zero-$H^2$ cut.  Otherwise the adjacent sizes are
at least two.  The budget leaves either three hom-set
sizes at most two, or $(2,3,1)$ or $(3,2,1)$.
In each case one extreme corner has both incident
stalks of size at most two.  Lemma~\ref{C:small-corner-two-stalks-rank-one}
kills the inherited rank-one square.  Bridge products
vanish since both their corners have zero $H^2$.

If a five-element corner occurs, all other corners
have order at most two and $P\leq3$.
A bridge at a positive five-element corner is, after
exact opposite coinvariants, an ordinary act over the
two-row sphere monoid, so has zero products.  A bridge
whose corners have zero $H^2$ uses the two-block lemma.
At the middle of the triangle, collapse the semisimple
extremes rationally.  A singleton adjacent set deletes
the middle object.  Otherwise the adjacent sizes are
at least two; either all three sets have size two,
giving an objectwise two-state cone, or the long set
is singleton and one adjacent size is two, giving
a two-row or two-column band by triangle compression.
These all have zero products.

At an endpoint, a five-element corner with zero $H^2$
is cut from the opposite semisimple two-object block.
If its $H^2$ is nonzero, it is the sphere monoid.
A singleton incident hom-set is a fixed arrow; the
opposite semisimple block has zero $H^2$, so the
fixed-arrow lemma kills all products.  Without such
a singleton, both incident sizes are at least two.
The opposite endpoint is semisimple.  Except when
the long set has size three, its two incoming sets
have size at most two after exact endpoint orbits,
so the objectwise-two-state lemma applies.
The exceptional sizes are $(2,1,3)$, up to opposite;
they force $D=4$, so the two other corners are trivial.
Delete the middle object toward the other trivial
endpoint and apply the ordinary two-row act theorem.

The only corner of order six has $D=5$, $P=2$, with
all other corners trivial.  To avoid an unjustified
corner-intersection argument at its possibly nonzero
$H_3$, group every simplex using this corner into one
piece $X$: include all its incident bridge edges and,
if it lies on the support triangle, that triangle too.
The remaining bridge pieces and, if not already included,
the unique triangle attach to this grouped piece and
to one another only at trivial corner spaces.
Hence $X$ inherits bounded rational homology and
$b_2\leq1$, and fourth-cohomology restriction is
injective.  All other pieces have trivial corners
and dimension at most two, so have zero fourth homology.

If the order-six corner is not on the triangle, $X$
is a star, and identification of its trivial leaves
gives the ordinary bounded rank-one act obstruction.
If it is in the middle of the triangle, it has no
further neighbor: an extra comparable neighbor would
give a second support cycle.  A singleton adjacent
hom-set deletes the middle, while two parallel
adjacent sets must both have size two and the long
set is singleton.  The latter triangle is an
objectwise two-state cone and has zero products.
If the order-six corner is an extreme of the triangle,
all its additional neighbors are on that same side.
When the opposite adjacent hom-set is singleton,
delete the trivial middle object toward the other
trivial extreme, leaving a star and hence an ordinary
act after identifying leaves.  Otherwise that opposite
set is parallel.  The budget $P=2$ forces one of the
two hom-sets incident with the order-six corner in
the triangle to be singleton.  This is a fixed arrow,
and its full complement in $X$ has trivial corners
and height at most one, so has zero $H^2$.
The fixed-arrow lemma finishes this case.

\paragraph{Two parallel-arrow excesses and at least two cycles.}
Suppose $P=2$ and $R\geq2$.  Then $D\leq4$.
If all corners have order at most four,
Lemma~\ref{C:two-parallel-small-corners} applies.
Otherwise $D=4$, $R=2$, and the sole nontrivial corner
is a five-element monoid $M$.  If its $H^2$ is zero,
the parallel-cut proof applies whenever the parallel
hom-sets are not adjacent; in the adjacent degree-two
case its triangle has a trivial extreme and two-point
stalks, so the same proof applies.  Alternatively one
may use the following direct positive-corner argument
only when its $H^2$ is nonzero.

If $H^2(BM;\mathbb Q)\ne0$, then $M$ is the sphere
monoid.  A singleton incident hom-set makes every
global degree-two restriction to $BM$ zero; the same
parallel-cut proof then works, with the binary triangle
exception again an objectwise two-state cone.
Without such a singleton, $M$ has one or two neighbors.
Let $X$ be the full subcategory on $M$ and these neighbors,
and let $D'$ be its complement with only the $M$ object
removed.  Their nerves cover the category and meet in
a full subcategory on at most two trivial objects.
Their third intersection homology is zero.
The category $D'$ has trivial corners and support cycle
rank at most two, so its nerve has dimension at most two.
If $M$ has one neighbor, $X$ is an ordinary two- or
three-state act over the sphere monoid and has zero
products.  If it has two incomparable neighbors, $X$
is its two-leg, two-state star; identifying the two
trivial endpoints gives an ordinary four-state act
and preserves products in degrees at least two.
If its neighbors are comparable, $X$ is a binary
triangle with a trivial extreme and hence is an
objectwise two-state cone.  Its products vanish in
every case, and restriction in degree four finishes.

\paragraph{The remaining two-cycle budget.}
For $R=2$ and $P\geq3$, the possibilities are
$(D,P)=(2,4)$ or $(3,3)$.
For $D=2$ there is just one order-three corner.
Its trivial modules have projective dimension at
most one, so every edge intersection has zero
homology in degrees at least three.  Fourth-cohomology
restriction to the triangle pieces is injective,
and they inherit bounded high homology.  At an endpoint
the order-three corner gives total bar degree at
most three; in the middle the small-middle Tor
obstruction gives zero fourth homology.  This closes
the $D=2$ case.

For $D=3$ all corners are rationally acyclic and
there is at most one corner of order at least three.
Bridges and individual triangles are handled as in
the one-cycle $D\leq3$ case.  Only a diamond remains.
Write its shared edge as $u<v$ and its exclusive
vertices as $a,b$.
When $u<a,b<v$, avoiding a zero-$H^2$ star cut forces
both potentially nonzero edges to meet at the one
larger exclusive corner, say $a$.
The shared-edge endpoints are semisimple, so that
edge has no homology in degrees at least two.
Restriction in degree four to the two triangles is
injective, and bounded high homology passes to them.
The triangle through $a$ has middle order at most
four and semisimple endpoints; the other is semisimple.
Both have zero fourth homology.

Suppose instead $a,b<u<v$, the opposite case being
dual.  If the shared edge has zero $H^2$, use the
cut $\{a,b\}\mid\{u,v\}$.  Otherwise that edge is
parallel and touches the unique larger corner.
A clean exclusive vertex, with both hom-sets to
$u,v$ singleton, gives a join and kills every global
restriction to the shared edge, proving the same cut.
Avoiding such a clean vertex uses all three parallel
excesses: the shared edge is doubled, and one of
the two incident hom-sets at each of $a,b$ is doubled.
Thus every hom-set in the diamond has size at most two.

If the larger corner is $u$, the final corner at $v$
is semisimple.  Its exact orbit collapse makes it
trivial and preserves the two-point stalk bound;
the objectwise-two-state lemma kills all products.
Suppose the larger corner is $v$.
The other three corners are semisimple, and their
full height-one source block has zero $H^2$.
The corner at $v$ has order at most four and also
has zero second rational cohomology.  The actual
two-block suspension lemma for
$\{a,b,u\}\mid\{v\}$ therefore kills every
degree-two product directly.  This excludes the
diamond and completes the two-cycle case.

\paragraph{The last three-cycle budget.}
The only case not yet covered with $R=3$ is
$D=2$, $P=3$.  There is one order-three corner.
As above, edge intersections have zero homology in
degrees at least three, and triangle pieces have zero
fourth homology after inheriting boundedness.
For a complete four-object block, all hom-set sizes
at most two give the objectwise-two-state argument
at a trivial extreme.  A single parallel hom-set
of size four gives the parallel-hom cut.
The only other pattern consists of one hom-set of
size three and another of size two.
If they are not serial, use the parallel-hom cut.
If they are separated serially, cut at the downset
of the second source, exactly as in
Lemma~\ref{C:two-parallel-small-corners}.
If they are adjacent, their common vertex has an
additional neighbor in the four-object chain;
the join-cut argument in that lemma applies.
Thus all products vanish on the four-object block,
and fourth-cohomology restriction finishes the case.

The only four-cycle budget was $D=P=2$, already
covered by the two-parallel lemma.  These alternatives
exhaust $D,P\geq2$ and $D+P+R=8$.
Every one contradicts the projective-plane square,
proving the theorem and its seven-object consequence.
\end{proof}

% Independently audited by root and B; source note retained.
\subsection{Strong pairs in categories with many objects}

\begin{proposition}\label{C:seven-objects-strong-pair21}
Let $C$ be a connected finite category with at least seven
objects and at most twenty-one arrows.  Assume that no
object is a retract of another.  If $C$ has a proper
strongly connected component on two objects, then
$BC\not\simeq\mathbb CP^2$.
\end{proposition}
\begin{proof}
Let $S$ be such a component, write $s$ for its number of
arrows, and put $\ell=|\operatorname{Ob}C|-2\geq5$.
Its two corner identities lie outside its minimal ideal:
an identity in that ideal would make its object a retract
of the other, by the minimal-group inverse adjustment in
Lemma~\ref{C:finite-category-Rees}.  In particular $s\geq6$.

Contract $S$ to one vertex of the undirected support graph.
The resulting graph has $\ell+1$ vertices and is connected,
so it has at least $\ell$ edges.  At least one edge meets
the contracted vertex, and strong connectivity expands
each such interface to arrows at both objects of $S$.
Counting the $\ell$ outside identities gives
\[
 |C|\geq s+2\ell+1.
\]
Thus $s\leq10$.  Every such strong pair has
$H^2(BS;\mathbb Q)=0$.  Indeed the rational categorical
Rees comparison requires both total row and total column
counts to exceed two for nonzero second homology.  Its
minimal ideal would then have at least nine arrows, in
addition to the two identities outside it, giving
$s\geq11$.

\paragraph{An extreme component.}
First suppose $S$ is a source or a sink in the
condensation order.  Take opposites to make it a source,
and let $D=C\setminus S$ have $q$ connected components.
If $H^2(BD;\mathbb Q)\ne0$, one of its connected
components, on $t$ objects, has at least $2t+2$ arrows
by Lemma~\ref{C:arrow-excess-two-H2}.  Every other
component on $t'$ objects has at least $2t'-1$ arrows.
Therefore
\[
 |D|\geq2\ell-q+3.
\]
Each component of $D$ must meet $S$, and every such
interface supplies at least two cross arrows.
Consequently
\[
 |C|\geq s+(2\ell-q+3)+2q
       =s+2\ell+q+3\geq s+2\ell+4.
\]
For $s\geq8$ and $\ell\geq5$ this exceeds twenty-one.
Thus $H^2(BD;\mathbb Q)=0$ in those cases.  Since
$H^2(BS;\mathbb Q)=0$ as well, the actual two-block
suspension comparison of
Lemma~\ref{C:two-block-zero-H2-products} makes every
degree-two product on $BC$ zero.

If $s=6$ or $7$, the minimal group is trivial and the
minimal ideal has total row and column counts both two.
There is one minimal row and one minimal column at
each object.  Lemma~\ref{C:one-row-source-block-collapse}
gives an actual homotopy-equivalent category with $S$
replaced by a point.  This removes at least $s-1$
internal arrows and at least one cross arrow.  For the
latter assertion, take an outside object receiving
arrows from $S$.  The orbit set of its right $S$-act
is an image of either nonempty stalk and therefore
has size at most the smaller stalk size; replacing
the two stalks by it saves at least one arrow.
All other arrow counts do not increase.  The new
category has at most $21-s\leq15$ arrows, contrary
to Theorem~\ref{C:CP2-lower16}.

\paragraph{A nonextreme component.}
There are distinct outside objects $a,b$ with
$a\to S\to b$.  They are distinct because a path
leaving and reentering a strong component would put
its intermediate object in that component.
Composition gives $a\to b$.  The contracted support
graph therefore contains a triangle and has at least
$\ell+1$ edges; at least two of these edges meet $S$.

Let $P_{\mathrm{ext}}$ be the sum of $|C(d,i)|-1$
and $|C(i,d)|-1$ over nonempty hom-sets with $i\in S$
and $d\notin S$.  Each support interface already
contributes two cross arrows, and these parallel
excesses are additional.  Hence
\begin{equation}\label{C:nonextreme-pair-parallel-budget}
 |C|\geq s+2\ell+3+P_{\mathrm{ext}},
 \qquad
 P_{\mathrm{ext}}\leq8-s.
\end{equation}
In particular $s\leq8$.

Suppose first $s=6$ or $7$.  The minimal ideal is the
four-arrow band with a unique arrow $k_{ij}:i\to j$
for each pair of its two objects.  Put $e_i=k_{ii}$.
It is a two-sided zero of the corner monoid at $i$,
and it differs from that corner identity.
Call the incoming condition at $i$ defective if some
cross arrow $g:d\to i$ satisfies $e_i g\ne g$.
Call the outgoing condition defective if some
$h:i\to d$ satisfies $h e_i\ne h$.
There are four conditions in total, two at each object.

An internal minimal-ideal cross arrow satisfies both
relevant conditions.  Every extra internal cross arrow
can make at most two conditions defective, one outgoing
at its source and one incoming at its target.  There
are at most $s-6$ extra internal arrows, including any
extra endomorphisms, so these contribute at most
$2(s-6)$ defects.  An external singleton hom-set is
fixed by its endpoint idempotent.  Any external defect
therefore requires a parallel hom-set and affects only
the condition at its one endpoint in $S$.
There are at most $P_{\mathrm{ext}}$ such hom-sets.
By \eqref{C:nonextreme-pair-parallel-budget}, the total
number of defective conditions is at most
\[
 2(s-6)+P_{\mathrm{ext}}\leq s-4\leq3.
\]
At least one of the four conditions is consequently
not defective.

Choose its object $i$ and the two minimal arrows to
and from the other object.  Their composite at $i$
is $e_i$.  All incoming arrows are fixed by $e_i$,
or all outgoing arrows are fixed, as appropriate.
The unit hypothesis in the one-sided deletion lemma
is automatic because $e_i$ is a zero of its corner.
Lemma~\ref{C:one-sided-corner-deletion}, or its opposite,
therefore deletes $i$ by a homotopy equivalence.
At least six arrows are removed: its corner has at
least two, there is an internal arrow in each direction,
and nonextremality supplies an external entrance and
an external exit at $i$.  The resulting category has
at most fifteen arrows, again impossible.

Finally suppose $s=8$.  Equality in the budget forces
$P_{\mathrm{ext}}=0$, so every nonempty external
hom-set is a singleton.  The minimal group is trivial,
and its band has total row and column counts either
$(2,2)$, $(3,2)$, or $(2,3)$.  These are the only
products of positive totals at least two whose band
size, plus the two outside identities, is at most
eight.  Thus one of the two total counts is two.
If it is the row count, there is one minimal row at
each object and all external incoming arrows are
fixed by the chosen local minimal idempotents.
Lemma~\ref{C:one-row-fixed-entrance-collapse} replaces
the whole block by a point, preserving homotopy type.
If it is the column count, use the dual fixed-exit
statement.  In either case the eight internal arrows
become one and no other arrow count increases.
The resulting category has at most fourteen arrows,
which is impossible.  This exhausts all proper strong
pairs under the stated hypotheses.
\end{proof}

\begin{corollary}\label{C:eight-objects-through21}
No finite category with at least eight objects and at
most twenty-one arrows has classifying space
homotopy equivalent to $\mathbb CP^2$.
\end{corollary}
\begin{proof}
Suppose such a connected category exists, and delete
retract objects until none remain.  Two deletions
would leave at most fifteen arrows and are impossible
by Theorem~\ref{C:CP2-lower16}.  Thus the remaining
category has at least seven objects.  Its excess
relative to $2k-1$ is at most six, because each
deletion lowers that excess by at least one.

If the remaining category is strongly connected,
Corollary~\ref{C:strong-category-CP2-minimum23}
excludes it.  A proper component with at least three
objects is excluded by
Proposition~\ref{C:large-SCC-minimal-candidate21}.
A proper strong pair is excluded by the preceding
proposition.  Otherwise its support is acyclic,
and Theorem~\ref{C:acyclic-arrow-excess-six} applies.
Every alternative is impossible.
\end{proof}

\subsection{Five and six objects with sixteen arrows}

\begin{theorem}\label{C:sixteen-five-six-objects}
No finite category with five or six objects and at most sixteen
arrows has classifying space homotopy equivalent to
$\mathbb {CP}^2$.
\end{theorem}
\begin{proof}
Theorem~\ref{C:CP2-lower16} treats at most fifteen arrows.
We may assume that the category is connected and has exactly
sixteen arrows.  The following enumeration uses complete
ordinary corner-monoid catalogs, including both orientations;
it does not identify a monoid with its opposite when that
would omit an actual corner choice.

First enumerate the underlying simple undirected support:
an edge records that at least one of the two cross hom-sets
is nonempty.  For $k$ objects this is a connected graph with
at most $16-k$ edges.  Enumerate every graph on $k$ labeled
vertices, and retain one representative of each orbit under
all $k!$ vertex permutations.  The exact inventories are
\[
\begin{array}{c|r|r|r}
 k&\text{labeled supports}&\text{support orbits}&
 \text{raw cardinality decorations}\\ \hline
 5&728&21&9\,039\,898\\
 6&24\,763&94&7\,492\,422
\end{array}
\]
The raw decoration count is before any structural or symmetry
prune.  To explain its completeness, each vertex receives a
positive corner size, and an undirected edge receives an
ordered pair $(a,b)$ of nonnegative cross cardinalities with
$a+b>0$.  Subject only to total arrow count sixteen, all such
choices are tried.  For a fixed support with $e$ edges their
number is the coefficient of $x^{16-k-e}$ in
\[
 (2-x)^e(1-x)^{-k-2e}.
\]
Indeed, a corner's excess contributes $(1-x)^{-1}$ and an
edge's excess contributes
$\sum_{j\geq0}(j+2)x^j=(2-x)/(1-x)^2$.
The full automorphism group of the support, together with
opposition, removes duplicate decorations.  No additional
vertex permutation is needed: any isomorphism of decorated
categories must already be an automorphism of the underlying
support representative.

Apply the previously proved initial/terminal, retract,
singleton cofinal-deletion, and singleton-core prunes.
A five-object corner of order eight forces the remaining
cross support to be a singleton tree and is excluded by
Lemma~\ref{C:tree-singleton-core-cups}; a larger corner,
or an order-eight corner with six objects, cannot fit the
connected arrow budget.  Thus every remaining corner has
order at most seven.  Acyclic choices with semisimple
rational diagonal are excluded by
Lemma~\ref{C:sixteen-high-semisimple}.  The program uses a
nonzero regular-trace determinant modulo $32003$ only as a
sufficient certificate for rational semisimplicity; a zero
determinant never excludes a corner.

For every remaining decoration, insert every literal corner
multiplication table from the complete ordinary catalogs.
Every cross product is then a variable whose possible values
are exactly its target hom-set.  Associativity propagation
assigns a variable only when an already known associativity
equation forces it, and rejects only contradictory equations.
The depth-first search tries every remaining allowed value.
Consequently every category in the retained scope appears.
Permutations inside cross hom-sets remove duplicate finished
tables; omitting corner permutations at this stage can only
retain additional representatives.  An exact normalized-bar
calculation first retains the tables with mod-two homology
$H_1=0$, $H_2=\mathbb F_2$.

The completed searches and their independent literal bar gates
have the following counts:
\[
\begin{array}{c|r|r|r|r}
 k&\text{finished corner-representative tables}&
 \text{retained rows}&u^2=0&H_3(\mathbb F_2)\ne0\\ \hline
 5&8\,301\,658&411\,909&411\,734&175\\
 6&51\,458&4\,022&4\,022&0
\end{array}
\]
For every retained row the bar gate checks all typed category
axioms, computes a nonzero degree-two cocycle, and either
exhibits a degree-three primitive for its square or computes
positive third mod-two homology.  Neither alternative is
compatible with $\mathbb {CP}^2$.  This proves the theorem.
\end{proof}

\begin{remark}[Scope and independent verification]
\label{C:sixteen-five-six-verification}
The producer is \path{c_categories16_support_high.cpp}.
The completed source stems are
\path{c_categories_16_5_support_CP2} and
\path{c_categories_16_6_support_CP2}.
The independent program \path{d_sixteen_support_graph_audit.py}
constructs graph orbits by expanding every orbit, rather than
using the producer's minimum-representative test.  It verifies
all $21$ and $94$ support representatives, their full
automorphism groups, every formal decoration count, and the
complete start/finish support logs.  The earlier
\path{d_sixteen_generator_scope_audit.py} records the
structural-prune and code-difference audit.

The literal normalized-bar verifier
\path{c_category_tsv16_cups.cpp} imports neither the category
generator nor its homology filter.  Its mathematical routines
are the previously independently audited fifteen-arrow
verifier; the sixteen-arrow specialization changes only the
two input-size guards.  Its completed output stems are
\path{c_category_tsv_16_5_support_complete} and
\path{c_category_tsv_16_6_support_complete}.
It checks respectively $124\,599\,941$ and $633\,948$ typed
associativity equations, $63\,327\,254$ and $227\,919$
cocycle equations, and $298\,082\,532$ and $675\,597$
square-primitive equations.

The separate checker \path{d_verify_sixteen_high_linkage.py}
compares every literal source, target, and multiplication table
with its certificate, requires complete source and certificate
EOFs and final producer counts, and hashes the entire source.
It additionally reconstructs the category and normalized bars
from scratch in Python for all $4\,022$ six-object rows and
all $175$ five-object positive-$H_3$ rows, recomputing the
ranks, the nonzero cocycle, and every applicable primitive.
This additional Python check is not claimed for the remaining
five-object rows: their complete independent check is the
C++ bar verifier just described.  The combined exact-linkage
record is \path{d_sixteen_high_linkage_verified.json}.
No partial sixteen-arrow three- or four-object search is used
in this theorem.
\end{remark}

\subsection{The completed four-object sixteen-arrow exclusion}

\begin{theorem}\label{C:sixteen-four-objects}
No finite category with four objects and at most sixteen arrows
has classifying space homotopy equivalent to $\mathbb {CP}^2$.
\end{theorem}
\begin{proof}
Theorem~\ref{C:CP2-lower16} handles smaller arrow totals.
For a connected category with sixteen arrows, enumerate its
undirected support, then every positive corner size and every
ordered pair of nonnegative cross cardinalities on each edge,
with positive sum on that edge.  There are six support orbits,
representing all $38$ connected labeled graphs on four objects.
As in the support formula in
Theorem~\ref{C:sixteen-five-six-objects}, the complete number
of cardinality decorations on these six representatives is
$3\,440\,120$.

Nonempty composable cross hom-sets must have a nonempty composite
hom-set.  Exactly $360\,179$ decorations satisfy this necessary
transitivity condition.  Apply the proved initial/terminal,
retract, singleton cofinal-deletion, and singleton-core reductions.
The new cardinality block cuts use only the sufficient criteria of
Proposition~\ref{C:exact-block-cut-flags}; the star cuts are
Proposition~\ref{C:sixteen-four-star-prunes}.
Take decoration orbits under the full automorphism group of each
support and opposition.  Exactly $2\,074$ canonical profiles remain.
Nine have a corner of order at least eight and are excluded by
Theorem~\ref{C:sixteen-four-large-corner}.
The remaining $2\,065$ profiles have complete associative-table
searches.

Every corner of these profiles ranges over the complete ordinary
monoid catalog of its order, retaining both orientations.
For each actual corner choice, the regular-trace determinant test
can certify rational semisimplicity, and the literal mod-two bar
test can certify zero second rational cohomology.  These are
sufficient flags only.  An acyclic four-object support with
semisimple diagonal has relative normalized bar length at most
three, so its fourth rational homology is zero.  Otherwise a
corner choice is omitted only if an oriented two-block cut
satisfies Lemma~\ref{C:two-block-zero-H2-products}.

For every remaining choice, all unknown typed products range
over their complete target hom-sets.  Associativity propagation
forces only values required by a known associativity equation;
the depth-first search tries every remaining value.  A discovered
retract is discarded by the smaller-arrow lower bound.
At each completed table the normalized mod-two bar filter retains
exactly the cases with $H_1=0$ and $\dim H_2=1$.
Permutations inside cross hom-sets remove repetitions; corner
permutations are not used to discard additional tables.
The completed count is $25\,782\,736$ literal
corner-representative tables, and $1\,092\,077$ retained rows.

The independent cup computation finds exactly
\[
\begin{array}{l|r}
\text{explicit degree-three primitive for }u^2&1\,089\,656\\
u^2\ne0\text{ and }H_3(-;\mathbb F_2)\ne0&2\,421.
\end{array}
\]
Both alternatives contradict the integral projective-plane type.
This excludes every retained table, while all omitted tables have
one of the stated proved obstructions.  The theorem follows.
\end{proof}

\begin{remark}[Independent support, source, and bar gates]
\label{C:sixteen-four-verification}
The producer, its pruning header, and the full literal source are
\begin{center}\small
\path{c_categories16_support_four_v2.cpp},\\
\path{c_cp2_block_prunes.hpp},\\
\path{c_categories_16_4_support_v2_CP2.tsv}.
\end{center}
The independent program \path{d_four_object_v2_profiles.cpp}
constructs graph orbits by marking every image, allocates edge
sizes before corner sizes, and expresses the structural conditions
separately.  It reproduces all six support orbits, all $38$ labeled
supports, all $3\,440\,120$ decorations, and the exact $2\,074$
required profiles.  The separate checker
\path{d_verify_four_object_v2_scope.py} matches all completed
profile intervals and every one of the $1\,092\,077$ source rows,
including the nine explicit large-corner sentinels and the final
table counter.  Its record is
\path{d_four_object_v2_scope_verified.json}.

The primary normalized-bar certificates are produced by
\path{c_category_tsv16_cups.cpp}, independently of the table
generator, in
\path{c_category_tsv_16_4_support_v2_full.jsonl.gz}.
The further independent checker
\path{b_verify_category_certificates.cpp} reconstructs all literal
bars and imports neither that cup routine nor the generator.
It verifies $472\,994\,717$ typed associativity equations,
$263\,583\,036$ cocycle and boundary-squared equations, and
$1\,337\,626\,118$ square-primitive equations.
Its complete linkage record
\path{b_sixteen_four_v2_audit_complete_verified.json} checks every
literal source row, the whole-source hash, complete EOFs, the gzip
footer, and the final counts.  Its status is
\texttt{PASS\_COMPLETE\_INDEPENDENT\_BAR}.
The zero-row invocation with an accidentally explicit zero limit
and the earlier incomplete generator prefix are retained as work
records; neither is used as evidence of completion.
\end{remark}

\subsection{A forty-arrow category of type $\mathbb CP^2$}
\label{C:CP2-upper40-section}

\begin{proposition}\label{C:CP2-upper40}
There is a category with two objects and forty arrows whose
classifying space is homotopy equivalent to $\mathbb CP^2$.
Its complete composition table is supplied by
\path{code/d_cp2_category40.py} and
\path{code/outputs/d_cp2_category40_verified.json}.
\end{proposition}
\begin{proof}
We first specify the forty-arrow table directly. Let
$L_0=\{1,x_{00},x_{01},x_{10},x_{11}\}$, with
$x_{ij}x_{kl}=x_{il}$. Put
\[
\begin{split}
 O={}&\{(0,1,2),(0,0,0),(0,0,1),(0,0,2),
             (1,1,1),(1,1,2),(2,2,2)\},\\
 D={}&\{(0,0,0),(0,0,1),(1,1,1),(2,2,2)\}.
\end{split}
\]
These tuples are monotone self-maps of $\{0,1,2\}$.
Let $A$ be the 29-element monoid
\[
 A=(\{1,x_{00},x_{01}\}\times O)
       \cup(\{x_{10},x_{11}\}\times D),
 \qquad
 (l,\theta)(m,\eta)=(lm,\eta\circ\theta).
\]
The identity is $(1,(0,1,2))$. Define the ten-element set
\[
 S=\{(a,b,c):a\in\{0,1\},\ a\leq b\leq c\leq a+2\}
          \setminus\{(1,2,2),(1,3,3)\}.
\]
For an integer tuple $s$ write
$\nu(s)=s-2\lfloor s_0/2\rfloor(1,1,1)$, and define
$\lambda_1(t)=t$ and
$\lambda_{x_{ij}}(t)=2\lfloor(t+j)/2\rfloor+i$.
The left $A$-action on $S$ is
\[
 (l,\theta)\cdot s
   =\nu\bigl(\lambda_l(s_{\theta(0)}),
              \lambda_l(s_{\theta(1)}),
              \lambda_l(s_{\theta(2)})\bigr).
\]
The displayed sets are closed under these operations. The resulting
category $C$ has objects $a,*$, with
\[
 C(a,a)=A,\quad C(*,*)=\{1_*\},\quad
 C(*,a)=S,\quad C(a,*)=\varnothing.
\]
Composition with an arrow in $S$ is exactly the displayed action;
all other compositions are multiplication in $A$ or identities.
Thus the arrow count is $29+1+10=40$.
The independent table check verifies all 1142 typed pairs and
33\,100 associativity triples.

To prove the homotopy type, adjoin an intermediate object $b$ with
$\operatorname{End}(b)=\{1_b\}$, arrows
$i_k:b\to a$ for $k=0,1$, and arrows
$r_{j,t}:a\to b$ for $j=0,1$ and $t=0,1,2$. Specify
\[
 r_{j,t}i_k=1_b,\qquad
 i_kr_{j,t}=(x_{kj},\operatorname{const}_t).
\]
The old elements act by
\[
 (l,\theta)i_k=
 \begin{cases}i_k&l=1,\\i_i&l=x_{ij},\end{cases}
 \qquad
 r_{j,t}(l,\theta)=
 \begin{cases}
 r_{j,\theta(t)}&l=1,\\
 r_{v,\theta(t)}&l=x_{uv}.
 \end{cases}
\]
Extend $S$ to the covariant set-valued functor $P$ on the
two-object subcategory $K=\{a,b\}$ by
$P(a)=S$, $P(b)=\{s_0\}$,
$P(i_k)(s_0)=(k,k,k)$, and $P(r_{j,t})$ the unique map.
The collage $G^-$ of $P$ has fifty arrows: it adds
$*\to b$ to the original category and the nine other arrows
involving $b$. Its full table was independently checked on all
43\,302 composable triples.

The inclusion $C\hookrightarrow G^-$ is a homotopy equivalence.
For a target in $C$ its comma category has a terminal object.
For the target $b$, choose $i_0$ and $r_{0,0}$ with
$r_{0,0}i_0=1_b$. The constant functor to
$(a,r_{0,0})\in(C\downarrow b)$ receives a natural transformation
from the identity: at $(d,\alpha:d\to b)$ its component is
$i_0\alpha:d\to a$. This contracts the comma category.
Quillen's Theorem A proves the assertion.

Since $P(b)$ is a singleton and $b$ retracts from $a$,
Lemma~\ref{lem:B-framed-collage} identifies $BG^-$ with the
classifying space of the explicit monoid
\[
 T=A\amalg\bigl(S\times\{0,1\}\times\{0,1,2\}\bigr),
 \qquad |T|=29+60=89.
\]
The second summand is a rectangular band; the left action is the
one just displayed and the right action is
\[
 (s,j,t)(l,\theta)=
 \begin{cases}
 (s,j,\theta(t))&l=1,\\
 (s,v,\theta(t))&l=x_{uv}.
 \end{cases}
\]
Identify the twelve pairs
\[
 (x_{kj},\operatorname{const}_t)=((k,k,k),j,t).
\]
The resulting quotient is exactly the 77-element monoid of
Proposition~\ref{prop:B-CP77}; the program checks every equation
of this specified quotient map $T\to M_{77}$.

Here the homology-equivalence assertion is established by a complete
resolution and its actual comparison map. Number $T$ by listing
the old elements in the displayed $L_0$ order, each with the
displayed $O$ or $D$ order, followed by the band in lexicographic
row, $j$, $t$ order. Write $u_i$ for element $i$, starting at zero,
and $R=\Z T$. The finite projective resolution has terms
\[
\begin{array}{c|l|r}
 n&P_n&\operatorname{rank}_{\Z}P_n\\\hline
 0&Ru_{29}&10\\
 1&\bigl(R(u_0-u_3)\bigr)^3&108\\
 2&R(u_8-u_{29})\oplus Ru_{29}\oplus
        \bigl(R(u_7-u_{10})\bigr)^5\oplus
        \bigl(R(u_3-u_1)\bigr)^4&204\\
 3&\bigl(R(u_1-u_8)\bigr)^2\oplus
        \bigl(R(u_{10}-u_8)\bigr)^7&117\\
 4&R(u_8-u_{29})\oplus Ru_{29}&12.
\end{array}
\]
Each displayed difference is an idempotent formed from two nested
idempotents. The complete integral differential matrices are in
\path{code/outputs/b_resolution_D_CP50_framed89.json}.
The standalone verifier reconstructs their $T$-actions, verifies
equivariance for all 89 elements, computes saturated images by
integer Smith form, checks every kernel rank and zero composite,
and verifies that the final kernel is zero. After tensoring with
the trivial right module, only one copy of $\Z$ in degrees
zero, two and four remains.

The certificate
\path{code/outputs/b_resolution_map_D_CP50_framed89_to_CP77.json}
is an augmentation-preserving comparison for the specified quotient,
constructed by integral lifts on each idempotent generator. The
standalone comparison verifier checks every chain square and every
source action and recomputes the maps on trivial-coefficient chains.
They are $(1)$ in each of degrees zero, two and four. Hence
$BT\to BM_{77}$ is an integral homology equivalence. Every element
of $T$ has eventually constant powers, as also checked directly,
so its universal group is trivial. The target is likewise simply
connected. The homological Whitehead theorem now gives
\[
 BC\simeq BG^-\simeq BT\simeq BM_{77}\simeq\mathbb CP^2.
\]
This uses the literal comparison map in degree four; equal additive
homology alone, or an isomorphism only in degree two, would not
identify the homotopy type.
\end{proof}

\begin{proposition}[An all-order obstruction to literal small-target Hopf maps]
\label{C:aperiodic-Hopf-target}
\label{C:literal-Hopf-target}
Let $M$ be a finite monoid with $H_2(BM;\Q)=0$.
Every unital homomorphism $M\to L_0$ induces a nullhomotopic map
of classifying spaces. In particular no such homomorphism from
a finite monoid of type $S^3$ represents a Hopf map to the
five-element model of $S^2$.
\end{proposition}
\begin{proof}
First collapse each minimal $\mathcal H$-class of $M$, obtaining
the actual quotient $M\to\overline M$ of
Lemma~\ref{lem:B-collapse-minimal-group}. This is a rational
homology equivalence. Each minimal $\mathcal H$-class is a finite
group with its own idempotent identity, and a homomorphism into a
band is constant on that group. Therefore every map to $L_0$
factors through $\overline M$. We may replace $M$ by this quotient
and suppose its minimal ideal is a rectangular band.

Suppose that $\phi:M\twoheadrightarrow L_0$ is onto.
Its image sends the minimal ideal of $M$ onto the minimal ideal
of $L_0$. Write its rectangular minimal ideal as $K=I\times J$;
the target ideal is the
$2\times2$ band. Choose $e\in K$ and write $f=\phi(e)$.
The maps of principal ideals
\[
 Me\longrightarrow L_0f,\qquad eM\longrightarrow fL_0
\]
are surjections of the corresponding row and column sets.
They induce surjections of their rational augmentation kernels
$V_M\twoheadrightarrow V_{L_0}$ and
$W_M\twoheadrightarrow W_{L_0}$: lift each difference from a
chosen base row or column to a difference in the source.
There is also a surjective ring map
\[
 A_M=\Q M/\Q K\longrightarrow
 A_{L_0}=\Q L_0/\Q(K_{L_0}).
\]

The two connecting maps used in Lemma~\ref{lem:B-band-Tor}
are natural for this homomorphism. Indeed, apply $\phi$ to
the two augmentation sequences for $Me,eM$, and to their
two-sided bar complexes; this gives commuting maps on both
connecting homomorphisms. Under the resulting natural
identification of degree-two homology, the induced map is
therefore exactly
\[
 W_M\otimes_{A_M}V_M\longrightarrow
 W_{L_0}\otimes_{A_{L_0}}V_{L_0},
 \qquad w\otimes v\longmapsto\phi(w)\otimes\phi(v).
\]
It is onto, since every pure tensor has lifts in both factors.
Thus $H_2(BM;\Q)\to H_2(BL_0;\Q)=\Q$ is onto, contradicting
the hypothesis.

Consequently the image of $\phi$ is a proper unital submonoid
of $L_0$. If its band elements use two different rows and two
different columns, products fill the entire $2\times2$ band,
so the image would be all of $L_0$. Its nonidentity elements
therefore lie in a single row or a single column. If this set
is nonempty, any of its elements is respectively a right zero
or a left zero in the image monoid. Such a zero gives a natural
transformation between the identity functor and the constant
functor, so the image has contractible classifying space.
The image $\{1\}$ is also contractible. The map $B\phi$
factors through that contractible space, proving the claim.
\end{proof}

\begin{remark}[Scope of subsequent reductions]\label{C:CP40-reductions}
The program \texttt{c\_CP40\_all\_congruences.py} enumerates the
complete lattice of 28\,135 congruences of the forty-arrow table
that preserve its two objects and hom-set types. Among the proper
quotients, 25\,704 have zero rational degree-two homology at the
endomorphism-monoid vertex, which forces the degree-two square
of the act cone to vanish; the other 2430 have zero degree-two
homology over $\mathbb F_2$. Thus no proper quotient in this
specified sense is another $\mathbb CP^2$ model.

There are ten maximal proper wide subcategories of this table.
Their complete act-cone projective resolutions, reconstructed and
checked independently, all have different additive homology from
$\mathbb CP^2$. This does not exclude more deeply nested
subcategories. A separate search of all sixteen smaller full
subcategories retaining the cone object and at least one of its
proper Karoubi-object classes finds either the wrong second Betti
number or a zero degree-two square. These are explicitly scoped
reduction checks, not a proof that forty is the category minimum.
The records are
\path{code/outputs/c_CP40_all_congruences.json},
\path{code/outputs/d_cp2_subcategory_search.json}, and
\path{code/outputs/d_cp2_karoubi_search.json}.
\end{remark}

\subsection{A twenty-two-arrow category of type $\mathbb CP^2$}

\begin{theorem}\label{C:CP2-upper22}
There is a category with two objects and twenty-two arrows whose
classifying space is homotopy equivalent to $\mathbb CP^2$.
Its complete composition table and its comparison with the
twenty-three-element monoid in Theorem~\ref{B:CP23} are given in
\path{code/outputs/a_CP22_category.json}.
\end{theorem}
\begin{proof}
Use the explicit monoid $N$ from Theorem~\ref{B:CP23}.
Thus its minimal ideal is the band
\[
 K=I\times J,\qquad I=\{0,1,2\},\quad
 J=\{0,1,2,3,4\},
\]
and the remaining eight elements are $1,a,b,c,d,e,f,g$.
The row transformations of the seven nonidentity elements, in that
order, are
\[
 002,\quad002,\quad010,\quad202,\quad202,\quad202,\quad212.
\]
Every one preserves the two-element subset $I_0=\{0,2\}$.
Furthermore, a product of two elements outside $K$ that lies in
$K$ has its row in $I_0$.  Indeed none of the outside row
transformations is constant.  The only outside transformations
that take the value $1$ are $010$, $212$, and the identity, each
having singleton inverse image of $1$.  A composition of two
such nonconstant transformations cannot therefore be the constant
map with value $1$.

It follows that
\[
 P=(N\setminus K)\sqcup(I_0\times J)
\]
is an eighteen-element submonoid of $N$.  Products touching its
band stay in that band: left multiplication preserves $I_0$ and
right multiplication preserves the row.  Its minimal ideal is
exactly $I_0\times J$.  To see minimality directly, multiplication
on the two sides by this band reaches every one of its elements
from any fixed element of $P$.

Set $M=P^{\mathrm{op}}$ and let $S=I$.  Define a right $M$-action
by
\[
 i\cdot m=\sigma_m(i),
\]
where $\sigma_m$ is the original left row transformation in $N$.
The reversal of multiplication makes this a right action.
Every element of the minimal ideal of $M$ acts constantly on $S$.
The category $C=\mathcal C(M,S)$ is therefore explicitly specified
by
\[
 \End_C(a)=M,\qquad\End_C(b)=\{1_b\},\qquad
 \Hom_C(a,b)=I,\qquad\Hom_C(b,a)=\varnothing,
\]
with composition $i\circ m=\sigma_m(i)$.  Its number of arrows is
$18+1+3=22$.

Apply Theorem~\ref{C:constant-minimal-act-compression} to $(M,S)$.
The minimal ideal of $M$ has row set $J$ and column set $I_0$.
The compressed monoid has underlying set
\[
 (N\setminus K)\sqcup(J\times I).
\]
It is exactly $N^{\mathrm{op}}$: the identification is the identity
on $N\setminus K$ and sends the new element $(j,i)$ to the original
$x_{ij}$.  Products within the band, the two outside actions, and
outside products entering the band agree by their defining
formulas.  Consequently the actual comparison of classifying
spaces supplied by that theorem gives
\[
 BC\simeq B(N^{\mathrm{op}})\cong BN\simeq\mathbb CP^2.
\]
The last equivalence is the complete integral-resolution and
primitive-square result of Theorem~\ref{B:CP23}; thus this transfer
includes all homology degrees and the attaching-map information.

For a reproducible finite check, retain the source labels
$0,1,\ldots,5,11,12,\ldots,22$ in its numbering, take the opposite
table, and use the row action above.  The program
\path{code/a_CP22_category.py} checks all $972$ action equations,
all $6\,862$ composable associativity triples of the category,
and all $529$ product equations of the displayed compressed-monoid
isomorphism.  The same construction also works for the separately
certified source numbered $1651$; both literal models are retained
in the output file.  The separate verifier
\path{code/b_verify_CP22_category.py} imports no producer code and
independently reconstructs both submonoids, actions, categories,
and compressed multiplication tables; its results are recorded in
\path{code/outputs/b_CP22_category_independent.json}.
The additional independent reconstruction
\path{code/d_verify_CP22_category.py} also checks the comparison
functor on every composable pair.
\end{proof}

\subsection{The complete sixteen-arrow triangular exclusion}

\begin{lemma}[Right acts of the five-element sphere monoid]
\label{C:P5-right-act-normal-forms}
Let $P=\{1\}\sqcup\{x_{ij}:i,j\in\{0,1\}\}$, with
$x_{ij}x_{kl}=x_{il}$, and let $S$ be a finite nonempty right
$P$-set.  Its core $SK$ is a disjoint union of fixed singleton
orbits and two-point orbits whose points are labelled by the
column $j$.  Each point of $S\setminus SK$ is specified by an
ordered pair of core orbits, one for each row $i$; multiplication
by $x_{ij}$ sends it to the column-$j$ point of the chosen orbit.
Every such choice defines a right act.  Up to state isomorphism,
the seven- and eight-state right acts number $383$ and $1212$.
\end{lemma}
\begin{proof}
For $s\in SK$, its right $K$-orbit is a quotient of the two-point
column act.  It is therefore a singleton or a column-labelled
pair.  For an arbitrary $s$ and a fixed row $i$, the points
$s x_{i0},s x_{i1}$ lie in one such orbit, since
\[
 (s x_{ij})x_{kl}=s x_{il}.
\]
This determines the asserted ordered pair of orbits.  Conversely
that same equation verifies the action law for every choice.
Transient states with the same ordered pair are interchangeable.
Thus it suffices to enumerate numbers of singleton and two-point
core orbits and multisets of their ordered pairs.

The producer \path{code/b_p5_right_acts.hpp} constructs these acts
and takes their actual simultaneous state-conjugacy classes.
The independent checker
\path{code/b_verify_p5_right_normalforms.py} instead views the
transients as a directed multigraph on the core orbits.  It
canonicalizes its integer adjacency data under permutations
within the two sizes of core orbit.  It checks every literal
action equation and that the extracted multigraph signatures
are distinct and exhaust the independently generated signatures.
The two catalogs agree and have the stated sizes.  As an
additional control, direct comparison with the earlier pointwise
constraint solver gives exactly
\[
 1,3,6,17,44,130
\]
right-act classes on one through six states.
The complete seven- and eight-state records are
\path{code/outputs/b_p5_right_catalog_profile16_independently_verified.json}
and
\path{code/outputs/b_p5_right_catalog_profile17_independently_verified.json}.
\end{proof}

\begin{proposition}\label{C:sixteen-triangular}
No triangular finite category with two objects and at most
sixteen arrows has classifying space homotopy equivalent to
$\mathbb CP^2$.
\end{proposition}
\begin{proof}
The totals at most fifteen are already excluded by
Proposition~\ref{C:fifteen-triangular}.  Write $(m,n,s)$ for
the two corner orders and the cross-hom cardinality at total
sixteen.  A singleton cross-hom gives a join.  Two cross arrows
are excluded by Theorem~\ref{C:two-cross-arrows-universal}; three
are excluded by Corollary~\ref{C:three-cross-arrows-lower22}.
If both corner spaces have zero rational $H^2$, every degree-two
class of the collage is relative to the vertices and its products
vanish.  Otherwise choose a corner with nonzero rational $H^2$
as the first corner, taking the opposite category when necessary.
Then $m\geq5$.

The case $n=1$ is an ordinary act cone and is covered by
Corollary~\ref{C:act-cone-twenty-one-frontier}.
For $n=2$ and $s\geq4$ one has $m\leq10$; exact rational
coinvariants give an actual orbit-act functor and a rational
cohomology-ring equivalence.  Apply the bounded-rational
statements of Lemmas~\ref{C:nine-bounded-rank-one} and
\ref{C:ten-rational-act-cones}.  The profile $(9,3,4)$ is
excluded by Corollary~\ref{C:right-zero-four-state-bounds}.
Finally $(5,7,4)$ is covered by
Corollary~\ref{C:sixteen-P5-opposite-seven}: a second corner
with zero $H^2$ is excluded structurally, and a positive one is
covered by opposition into $(7,5,4)$.
The remaining profiles are precisely the thirteen rows below.

The first-corner catalogs are the complete positive-$H^2$
covers of orders five through eight already used in
Proposition~\ref{C:fifteen-triangular}; their sizes are
$1,5,43,516$.  The second corner ranges over every ordinary
monoid of the indicated order three through six, with catalog
sizes $7,35,228,2237$.  A second corner with projective trivial
right rational module is omitted only with an exact vector
$v$ satisfying $va=v$ for every corner element and
$\epsilon(v)=1$.  The resulting orbit-act comparison and
Lemma~\ref{C:nine-bounded-rank-one} exclude that case.
Thus no replacement of an arbitrary corner by a homotopy-equivalent
monoid, and no unproved restriction on a zero-$H^2$ second corner,
enters the catalog.

For every retained first corner, the solver enumerates right acts
up to the full state-permutation group.  For a fixed right act,
every commuting left act is enumerated, and only then quotiented
by the \emph{actual stabilizer of that right act}.
This retains every relative position of the two actions.
The complete stabilizer and the disjoint orbit weights are
stored and independently verified.

A proper equivariant idempotent of the biset gives the actual
smaller collage retract of Lemma~\ref{C:proper-biset-retract-CP2}.
If the original collage were a projective plane, so would be
that smaller retract, contradicting the complete fifteen-arrow
exclusion.  Trivial action sides give ordinary-act retracts
with at most fifteen arrows, and permutation actions on both
sides give the finite-cover transfer obstruction.  All these
omissions retain their literal witnesses.  Every remaining case
is checked by the normalized bar over $\mathbb F_2$: wrong
$b_1$ or $b_2$, an explicit square primitive, or positive $b_3$
excludes it.

For the last two profiles, use
Lemma~\ref{C:P5-right-act-normal-forms} before enumerating left
actions.  Exactly three entire right-action branches are
excluded in each: the trivial action, the action whose entire
core is a single fixed point, and the action whose entire core
is a two-point orbit, with every band actor constant.
The first uses the ordinary-act retract, the second has a
common point fixed by every commuting left action, and the third
uses Theorem~\ref{C:two-constant-image-bisets}.
Thus $380$ and $1209$ right-act classes remain.  These whole-branch
omissions are recorded separately and are not included in the
raw biset counts below.

The completed counts are
\[
\begin{array}{c|r|r}
(m,n,s)&\text{enumerated bisets}&\text{stabilizer representatives}\\\hline
(8,4,4)&5091042&2623517\\
(7,5,4)&4287807&1842296\\
(6,6,4)&5477024&1807300\\
(8,3,5)&5822749&2444567\\
(7,4,5)&5392764&1503278\\
(6,5,5)&7066364&1203505\\
(5,6,5)&19744699&2137018\\
(7,3,6)&4936972&1121518\\
(6,4,6)&10361883&816676\\
(5,5,6)&44714285&1410219\\
(6,3,7)&5992085&448933\\
(5,4,7)&7089764&762793\\
(5,3,8)&4173038&334064.
\end{array}
\]
Every one of the $130150476$ enumerated bisets is excluded.
The $18455684$ retained representatives have no surviving
projective-plane candidate.
\end{proof}

\begin{remark}[Complete independent verification]
\label{C:sixteen-triangular-verification}
The separate verifier \path{code/b_verify_core_biset_certificates.cpp}
rebuilds every action and literal category, checks the full
right-action stabilizer, checks disjoint left-action orbit
weights, and independently verifies every retraction, primitive
or positive-third-homology claim.  Its complete thirteen-profile
gate checks $4483769801$ action equations,
$288244240$ typed associativity equations,
$18570604$ cocycle equations and $115804412$ primitive equations.
There are $16922804$ checked proper-retraction witnesses.
The independent generator-image enumeration
\path{code/d_sixteen_biset_counts.cpp} additionally regenerates
all raw counts for profiles six through fourteen in the original
numbering, including the redundant $(5,7,4)$ control.  Its
$119988968$ bisets agree in every one of the $9651$ individual
corner-pair counts; this numerical gate is not asserted for the
other profiles.

The final linkage checker
\path{code/b_verify_sixteen_triangular_complete.py} verifies the
literal input catalogs, exact rational splitting vectors,
every active second-corner list for every right act, completed
profile sentinels, and all certificate-prefix hashes.  It also
links the last two profiles to the independent complete P5
right-act catalogs.  The combined result is
\path{code/outputs/b_sixteen_triangular_complete_verified.json}.
Only completed prefixes of the interrupted historical core stream
are used; unfinished tails have no role in this result.
\end{remark}

\begin{corollary}\label{C:sixteen-two-objects}
A finite category with at most two objects and classifying space
homotopy equivalent to $\mathbb CP^2$ has at least seventeen arrows.
\end{corollary}
\begin{proof}
The one-object case follows from
Theorem~\ref{B:CP2-minimum23}.  A connected two-object category
is strongly connected or triangular.  Apply
Theorem~\ref{C:twentytwo-strongly-connected} in the first case
and Proposition~\ref{C:sixteen-triangular} in the second.
\end{proof}

\subsection{The completed sixteen-arrow linear calculation}

\begin{theorem}\label{C:sixteen-three-linear}
No three-object finite category with linear acyclic support
and at most sixteen arrows has classifying space homotopy
equivalent to $\mathbb CP^2$.
\end{theorem}
\begin{proof}
Theorem~\ref{C:CP2-lower16} handles smaller arrow totals.
At total sixteen, write the corner orders as $(m,n,p)$
and the forward hom-set sizes as $(x,y,z)$ in directions
$01,12,02$.  All six integers are positive, and their sum
is sixteen.  There are $\binom{15}{5}=3003$ such ordered
six-tuples.  Opposition and object relabeling give exactly
$1547$ canonical Hom-cardinality profiles.

The structural omissions use the singleton joins and
singleton cofinal deletions, the fixed-arrow lemma, the
two-block zero-product lemma, and the triangle-band
compression already proved above.  Large middle corners
and the necessary extreme-corner catalog are treated in
Proposition~\ref{C:sixteen-linear-large-corners}.
The later exact middle-algebra reductions are
Corollary~\ref{C:semisimple-ends-middle-four} and
Corollary~\ref{C:semisimple-ends-binary-homs}.
The final endpoint reduction is
Theorem~\ref{C:semisimple-endpoints-sixteen}, with its
stated initial-left and final-right projectivity hypotheses.
No corresponding assertion for arbitrary larger arrow
totals is used.

The search was continued through successive completed-profile
prefixes as these reductions became available.  Across all
six retained streams, exactly $347$ distinct canonical
profiles have complete table searches.  The other $1200$
canonical profiles have one of the proved structural
obstructions.  A separate inventory enumerates the six
positive sizes directly, rather than the producer's
nine-bin distribution of arrows.  It reconstructs every
canonical profile and every structural reason, checks
that all completed profile sets are disjoint, and matches
the final required order of the remaining profiles.
Nontransitive assignments with two nonempty adjacent
hom-sets and an empty composite hom-set cannot be categories
and are not counted as valid linear profiles.

For each retained profile, all corner monoids of order at
most seven range over the complete ordinary catalogs,
including both orientations.  Only the proved possible-
positive-$H^2$ catalog is needed for an extreme corner of
order eight.  It contains $1032$ distinct literal table
orientations.  The number of literal corner tables across
the eight order catalogs is $36459$.
Each remaining cross product ranges over its entire typed
target hom-set.  Associativity propagation imposes only
equalities forced by the category axioms, and every
unforced value is tried.  Thus the search retains every
category in its stated scope.  Exact permutations inside
cross hom-sets may remove repeated tables; no relative
position of the corner actions is omitted.

The endpoint flags used in the final stream are separately
certified.  For every one of the $36459$ literal tables,
\path{code/b_verify_linear16_projective_endpoints.py}
solves the actual rational invariant-vector equations
on both sides, with augmentation equal to one.  It
compares these exact criteria with the producer's
minimal-ideal cardinality flags.  All $72918$ comparisons
agree, and all $12320569$ corner associativity equations
are checked.  The result is
\path{code/outputs/b_linear16_projective_endpoints_verified.json}.
The independent profile inventory also predicts the two
new producer counters exactly: $5311$ early endpoint
profile omissions and $838045$ projective-endpoint corner
choices.  Both equal the actual final footer.

For every completed table, an initial normalized-bar filter
retains the necessary mod-two conditions $H_1=0$ and
$\dim H_2=1$.  The independent normalized-bar computation
\path{code/c_category_tsv16_cups.cpp} then reconstructs
the literal category and all normalized chains through
degree four.  It checks the typed category axioms and
identifies a nonboundary degree-two cocycle.  Either it
supplies and evaluates a degree-three primitive for the
Alexander--Whitney square, or it independently computes
positive third mod-two homology.  The mutually exclusive
completed results are
\[
\begin{array}{c|r|r|r|r}
 \text{stream}&\text{complete profiles}&\text{retained rows}
   &u^2=0&H_3(\mathbb F_2)\ne0\\\hline
 \text{original}&65&1401244&1347792&53452\\
 v2&51&723515&709304&14211\\
 v3&59&865106&855934&9172\\
 v4&107&1592464&1585896&6568\\
 v5&10&415822&406100&9722\\
 v6&55&1080989&1069881&11108\\\hline
 \text{total}&347&6079140&5974907&104233.
\end{array}
\]
There is no surviving candidate.  Every retained row
therefore has an obstruction to the homology or cup ring
of $\mathbb CP^2$, proving the theorem.
\end{proof}

\begin{remark}[Completed prefixes and full scope]
\label{C:sixteen-linear-complete-verification}
The final producer is
\begin{quote}\small
\path{code/c_categories16_linear_resume_v6.cpp}.
\end{quote}
Its actual final footer records $97067056$ finished
corner-representative tables and $1080989$ retained rows
in that stream.  Its reported $292$ previously completed
keys agree with the disjoint prior sets of sizes
$65,51,59,107,10$; all $55$ newly required profiles finish.
The producer's larger distribution counter also includes
invalid supports rejected before a profile starts and
is not substituted for the completed-profile count.

Interrupted historical streams remain unchanged.
Only their exact completed-profile prefixes in the table
above are used.  The original source's storage-failure
tail and every later unfinished profile are excluded.
The bounded orbit cache uses an injective packed key
for the typed multiplication entries; clearing a full
cache can retain additional repetitions and cannot
discard a previously unseen category.

The separate completed-prefix gates match every literal
source, target and multiplication table with its saved
cochain record, require the verifier's final count and
the compressed certificate footer, and hash each exact
source prefix.  The final program
\path{code/b_finish_sixteen_linear_scope.py}
then re-reads all six prefixes.  It independently checks
every row's Hom profile, every profile boundary, all
six source hashes, all certificate-interval unions,
and the actual final source EOF.  It requires the full
inventory gate
\path{code/outputs/b_sixteen_linear_inventory_complete.json};
a stopped or merely growing stream cannot satisfy it.

The complete checks comprise $2767269079$ typed
associativity equations, $1483020541$ cocycle equations,
and $6463509055$ literal square-primitive equations.
Their final record is
\begin{quote}\small
\path{code/outputs/b_sixteen_linear_complete_verified.json},
\end{quote}
with status
\begin{quote}\small\ttfamily
PASS\_COMPLETE\_SIXTEEN\_ARROW\_LINEAR\_SCOPE.
\end{quote}
This records the full independent table-to-bar computation
and exact source linkage; no uncomputed homology tail is
used in the exclusion.
\end{remark}

\begin{theorem}[The global lower bound of seventeen arrows]
\label{C:CP2-lower17}
Every finite category whose classifying space is homotopy
equivalent to $\mathbb CP^2$ has at least seventeen arrows,
identities counted.  Together with the explicit construction,
the minimum-arrow interval is $17$--$22$.
\end{theorem}
\begin{proof}
Theorem~\ref{C:CP2-lower16} excludes all smaller totals,
so it remains to exclude sixteen arrows.  Connectedness
is necessary.  The one- and two-object cases are covered
by Corollary~\ref{C:sixteen-two-objects}.
For three objects, a nontrivial strongly connected
component is excluded by
Theorem~\ref{C:sixteen-three-SCC}.  An acyclic connected
support on three objects is an inward or outward star,
or a transitive linear support: a composable pair of
strict arrows forces its long composite hom-set to be
nonempty.  Apply Theorem~\ref{C:sixteen-all-stars} and
Theorem~\ref{C:sixteen-three-linear}, respectively.
Four objects are excluded by
Theorem~\ref{C:sixteen-four-objects}, five and six by
Theorem~\ref{C:sixteen-five-six-objects}, and every larger
object count by Theorem~\ref{C:sixteen-seven-objects}.
These cases exhaust all finite categories with sixteen
arrows.  The upper bound is Theorem~\ref{C:CP2-upper22}.
\end{proof}


% Root draft; independent literal audit requested before input.
\subsection{One-sided minimal rows and small categorical acts}

\begin{lemma}[At most two states at each object]
\label{C:two-states-objectwise-act-cups}
Let $S$ be a finite category and let
$F:S^{\mathrm{op}}\to\mathrm{Set}$ be a right act such that
$1\leq |F(c)|\leq2$ at every object. Adjoin one trivial sink
with outgoing hom-sets $F(c)$ from $S$, and call the resulting
act cone $X$. Every product of two positive-degree rational
cohomology classes on $BX$ is zero.
\end{lemma}
\begin{proof}
For every object $c$ take the convex simplex $I(c)$ on the
one- or two-element set $F(c)$, and extend each act map affinely.
These extensions compose strictly, since affine maps on a simplex
are determined by its vertices. Every $I(c)$ is contractible.
The pointed quotient $I(c)/F(c)$ is a point when $|F(c)|=1$
and a circle when $|F(c)|=2$. Every induced map between these
quotients is a pointed constant map, an identity, or a reflection,
after choosing an orientation for each circle.

The degree-two map on every such circle, and the identity of
each point, form a strict natural self-map of this pointed
diagram. Indeed, $t\mapsto2t$ on $\mathbb R/\mathbb Z$
commutes with $t\mapsto-t$, with identities and with pointed
constant maps. The pointed homotopy colimit of $I/F$ is the
cofiber of $\operatorname{hocolim}F\to\operatorname{hocolim}I$.
One can see this by first taking the ordinary homotopy colimit
of the objectwise quotient pushouts and then collapsing its
basepoint section $BS$. Since
$\operatorname{hocolim}I\to BS$ is an equivalence, this cofiber
is the classifying space of the stated act cone.

The reduced cellular chain complex of each quotient is
concentrated in degree one. The natural self-map therefore acts
as multiplication by two on every chain group of the reduced
pointed bar, and hence on all positive rational cohomology.
For positive-degree classes $x,y$, multiplicativity gives
$2xy=4xy$, so $xy=0$.
\end{proof}

\begin{lemma}[Collapsing a source block with one minimal row]
\label{C:one-row-source-block-collapse}
Let $C$ be a finite category partitioned into a full strongly
connected source block $S$ and its full complement $D$.
Assume the minimal ideal of $S$ has trivial maximal group and
exactly one minimal row at every object. No restriction is
imposed on its column counts. Replacing $S$ by one trivial
object and each incident right $S$-act by its orbit set gives
an actual quotient functor $Q:C\to\bar C$ inducing a homotopy
equivalence. Dually a sink block with trivial minimal group
and one minimal column at each object can be collapsed.
\end{lemma}
\begin{proof}
Choose an object $a\in S$ and a minimal idempotent
$e\in S(a,a)$. The minimal Rees description says that
$S(a,i)e$ has exactly one element for every $i$; call it
$u_i:a\to i$. At $i=a$ this element is $e$, and for every
$f:i\to j$ one has
\[
 u_i e=u_i,\qquad f u_i=u_j.
\]
The arrows $u_i$ give a natural transformation from the constant
functor at $a$ to the identity of $S$, proving that $BS$ is
contractible. A component of a natural transformation need not
be an identity, so $u_a=e$ causes no problem.

For a finite right $S$-act $F$, let its element category have
arrows $(i,x)\to(j,y)$ represented by $f:i\to j$ with $x=yf$.
Put $T=F(a)e$ and define
\[
 p(i,x)=x u_i\in T,\qquad s(t)=(a,t).
\]
The equation $f u_i=u_j$ proves that $p$ is constant on every
arrow, while $ps=\mathrm{id}_T$ because $te=t$.
The arrows $u_i$ give a natural transformation
$sp\Rightarrow\mathrm{id}$: its component at $(i,x)$ has
source $(a,xu_i)$ and target $(i,x)$, and its naturality is
exactly the same equation. Thus the element category is
homotopy equivalent to the discrete set $T$, which is its
orbit set. These functors and natural transformations commute
strictly with every morphism of right acts.

For each object $d\in D$ apply this construction to
$F_d(i)=C(i,d)$. Postcomposition by arrows of $D$ preserves
the whole construction. The two-block correspondence homotopy
pushout therefore compares $BC$ with $B\bar C$ by equivalences
on its source block and cross-arrow space, and by the identity
on $BD$. Its homotopy invariance proves the assertion for the
actual orbit-collapse functor. Since no arrow goes from $D$
back into $S$, quotient composition is exactly the induced
$D$-action on these orbit sets. Opposition gives the dual.
\end{proof}

\begin{proposition}\label{C:proper-strong-pair-at-least-seventeen}
Let $C$ be a connected finite category with at most twenty-one
arrows and a proper strongly connected component $S$ on two
objects. If $S$ has at least seventeen arrows, then
$BC\not\simeq\mathbb CP^2$.
Consequently a minimal-arrow projective-plane candidate through
twenty-one arrows has every proper two-object strong component
of size at most sixteen; moreover an extreme such component
with trivial minimal group cannot have one minimal row per
object at a source, or one minimal column per object at a sink.
\end{proposition}
\begin{proof}
Let $\ell$ be the number of objects outside $S$, so $\ell\geq1$.
Contracting $S$ in the connected undirected support leaves a
graph on $\ell+1$ vertices, which needs at least $\ell$ edges.
At least one such edge is incident with $S$. Strong connectivity
forces its incidence to include an arrow at both objects of $S$,
giving at least one additional arrow beyond this support count.
Together with the $\ell$ outside identities,
\[
 |C|\geq |S|+2\ell+1.
\]
For $|S|\geq17$ the bound twenty-one implies $\ell=1$.
Write $M$ for the remaining corner monoid and $t$ for the
number of cross arrows. All these arrows have one direction,
and $t\geq2$. If $t=2$, there is one arrow at each object of
$S$ and the category is the categorical join of its two blocks.
Its positive rational products vanish.

Otherwise $t\geq3$, and the budget forces
$|S|=17$, $|M|=1$, $t=3$. Take opposites if needed to make
$S$ the source. Its two outgoing stalks have sizes one and two,
so Lemma~\ref{C:two-states-objectwise-act-cups} excludes the
nonzero square of $\mathbb CP^2$.

The asserted size bound follows. The final extremal assertion
follows from Lemma~\ref{C:one-row-source-block-collapse}:
such a collapse preserves the homotopy type and strictly
decreases the arrow count, contradicting minimality. As always,
the minimality conclusion is about a category with the fewest
arrows, and does not assert that a nonminimal model cannot
contain the stated component.
\end{proof}

% Root draft, not input. Independent audit required.
\subsection{A one-row block with fixed external entrances}

\begin{lemma}\label{C:one-row-fixed-entrance-collapse}
Let $S$ be a full strongly connected component of a finite
category $C$, with at least two objects in $S$. Assume that its minimal group is trivial and
it has one minimal row at each object. Fix an object $a$,
a minimal idempotent $e\in S(a,a)$, the unique arrows
$u_i\in S(a,i)e$, and arrows $v_i\in eS(i,a)$.
Put $e_i=u_iv_i$; then $v_iu_i=e$.
Suppose $e_i p=p$ for every arrow $p:d\to i$ with
$d\notin S$. Replacing $S$ by one trivial object gives
a strictly smaller homotopy-equivalent finite category.
In particular the hypothesis holds when every nonempty
external incoming hom-set is a singleton. There is a dual
statement with one minimal column and fixed external exits.
\end{lemma}
\begin{proof}
Write $D=C\setminus S$. Define a category $\bar C$ with
objects $D\sqcup\{s\}$, retain all full $D$ hom-sets,
and put
\[
 \bar C(s,s)=\{1_s\},\qquad
 \bar C(d,s)=eC(d,a),\qquad
 \bar C(s,d)=C(a,d)e.
\]
Composition with arrows of $D$, and composition through
$s$ between two objects of $D$, are the original products
in $C$. They are well typed. A path leaving and then
reentering $S$ would put its intermediate objects in the
same strong component, so no such path occurs. Thus there
is no composite from $s$ through $D$ back to $s$ to define.
The displayed incoming and outgoing sets are fixed by $e$
on their respective sides. This verifies every associativity
equation involving $1_s$; the others are inherited from $C$.

Define $Q:C\to\bar C$ to fix $D$, send $S$ to $s$ and
every arrow of $S$ to $1_s$, and set
\[
 Q(p:d\to i)=v_i p,\qquad
 Q(q:i\to d)=q u_i.
\]
Since $S(a,i)e$ is a singleton, $f u_i=u_j$ for every
$f:i\to j$ in $S$. For an external entrance $p$,
$p=e_i p=u_i v_i p$. Therefore
\[
 v_j f p=v_j f u_i v_i p=v_j u_j v_i p=v_i p,
 \qquad q f u_i=q u_j.
\]
A composite through $S$ is also preserved, because
$(q u_i)(v_i p)=q e_i p=qp$.
These identities, together with associativity for retained
arrows of $D$, check every typed pair and prove that $Q$
is an identity-preserving functor.

There is a section $J:\bar C\to C$ choosing $a$ for $s$,
the literal displayed arrows for the two kinds of cross hom,
and $1_a$ for $1_s$. The definition of composition makes
this a functor and $QJ=\mathrm{id}_{\bar C}$.
The components $u_i:a\to i$ on $S$ and identities on $D$
give a natural transformation $JQ\Rightarrow\mathrm{id}_C$.
The internal equation is $f u_i=u_j$; the entrance equation
is $u_iv_ip=p$; the exit equation is the definition of $Q(q)$.
Hence $BQ$ and $BJ$ are homotopy inverses.
The quotient loses at least one object and never increases
an incident hom-set or the number of internal arrows, so it
strictly decreases the total arrow count. A singleton incoming
hom-set is necessarily fixed by the endomorphism $e_i$.
Taking opposites proves the dual statement.
\end{proof}

% Further root budget work, not yet a numbered assertion:
% For a no-retract proper two-object SCC of size s in a connected k-object
% category, n>=s+2(k-2)+1. If the SCC has both an incoming and outgoing
% component, the contracted support contains a triangle and two incidences
% with the SCC, giving n>=s+2(k-2)+3.
% At k=8,n<=21 this forces s<=8 at an extreme block and s=6 otherwise.
% An extreme G=1 strong pair of size<=8 is deletable: totals(2,2) collapse;
% totals(3,2)/(2,3) exhaust an eight-arrow minimal band plus identities,
% and a one-row (source) or one-column (sink) object is removed by the
% one-sided corner-deletion lemma. In the nonextreme s=6 equality case,
% every external nonempty hom is singleton, so the lemma above applies.
% Combined with D's audited acyclic excess-six theorem and the existing
% large-SCC exclusion, this will exclude every k>=8 candidate through21.
%
% At k=7, the same method appears to close SCC cases as well. An extreme
% block has s<=10. For s9/10, BS is rationally acyclic; the complement of
% at least5 objects cannot afford a component with nonzero H2 if the
% connected excess<=2 H2-vanishing lemma is supplied (D is checking it).
% Nonextreme blocks have s<=8. For band totals(2,2), s6 needs four
% entrance/exit obstructions to prevent deleting both two-element corners,
% but only two extra arrows are available; an internal extra cross arrow
% can cover two obstructions, while an external parallel extra covers one.
% The s7 cases still leave an unobstructed deletion; at s8 all external
% homs are singleton and the lemma above applies. The alternative s8
% band totals(3,2)/(2,3) permit a one-sided corner deletion directly.

% Uninput; independent literal audit requested.
\subsection{Two excess arrows cannot support second cohomology}

\begin{lemma}\label{C:arrow-excess-two-H2}
Let $C$ be a connected finite category with $k$ objects and
$n\leq2k+1$ arrows.  Then $H^2(BC;\mathbb Q)=0$.
\end{lemma}
\begin{proof}
Put $e=n-(2k-1)\leq2$.  Deleting an object which is a
retract of another preserves the classifying-space
homotopy type and connectedness, removes at least its
identity and two cross arrows, and lowers $e$ by at
least one.  Repeat until no such object remains.
If one object remains, its monoid has at most three
elements and has zero second rational cohomology by
Theorem~\ref{E:all-six-types}.

A nontrivial strongly connected component with $t$
objects has at least $t^2+t$ arrows: all cross hom-sets
are nonempty, and every corner has at least two
elements, since a trivial corner would give a retract.
If it is the whole category, its excess is at least
$t^2-t+1\geq3$.  If it is proper, the connected support
budget in the proof of
Theorem~\ref{C:arrow-excess-four} gives excess at least
$t^2\geq4$.  Both are impossible, so support is acyclic.

Write $D,P,R$ for the corner, parallel-arrow and support
cycle excesses as in
\eqref{C:excess-decomposition}.  Then
$D+P+R=e\leq2$.
If $P=0$, the actual support-poset functor induces an
isomorphism on $H^2(-;\mathbb Q)$ by
Lemma~\ref{C:singleton-support-H2}.  Its comparability
graph has cycle rank at most two.  It can contain at
most two triangles: the boundaries of any three
distinct triangles are linearly independent, because
each triangle has an edge absent from the other two.
The boundaries of one or two triangles are independent
as well.  Thus the degree-two simplicial boundary is
injective, and the support poset has zero second
homology.

Suppose $P\geq1$.  A nonsemisimple corner algebra would
require a monoid of order at least three and hence
$D\geq2$, contradicting the budget.  All corner
algebras are therefore semisimple.  If $R=0$, the
support tree has no composable strict pair, so the
relative bar has length at most one and has zero
second homology.  Otherwise equality is forced:
$D=0$, $P=1$, and $R=1$.  All corners are trivial and
there is only one doubled hom-set.  A unicyclic support
without a triangle has no two-simplices.  If its unique
cycle is a triangle, the nerve is the union of that
triangle category and its graph bridges, meeting at
points.  At least one adjacent hom-set in the triangle
is a singleton.  The component of its degree-two
boundary on the other adjacent edge is injective on
the two-simplex labels.  Thus its second homology is
zero, as is that of every bridge.  The union consequently
has zero second rational homology.
Duality completes the proof.
\end{proof}


% Root and B literal audits passed; grid endpoints explicitly repaired.
\subsection{Two minimal rows at each object of an act source}

\begin{lemma}[Finite affine diagrams on a common rational grid]
\label{C:categorical-grid-act-cups}
Let $D$ be a finite category and let $T:D^{\mathrm{op}}\to\mathrm{Set}$
be a finite right act. Suppose each $T_i$ is the set of all points
of $\delta\mathbb Z$ in a nonempty closed interval $I_i$ whose
endpoints lie in $\delta\mathbb Z$, where $\delta>0$ is fixed and rational. Suppose every act map is the
restriction of an affine map $I_j\to I_i$ of slope $0,1$, or $-1$.
Degenerate intervals are allowed. Then every positive-degree
rational cohomology product of the act cone is zero.
\end{lemma}
\begin{proof}
The affine maps give a diagram of contractible intervals.
For a nondegenerate source interval, their composition identities
follow from those on its two endpoints. For a degenerate source
interval there is only one point to check. Thus this is an actual
diagram, and its homotopy colimit maps equivalently to $BD$.
The pointed homotopy colimit of the objectwise quotients $I_i/T_i$
is therefore the cofiber of $\operatorname{hocolim}T\to BD$,
which is the classifying space of the act cone. As in
Lemma~\ref{C:two-states-objectwise-act-cups}, this identification
collapses the basepoint section of the ordinary category bar.

Each nondegenerate quotient is a wedge of circles indexed by
consecutive grid intervals of length $\delta$. An affine act map
of slope $1$ or $-1$ sends each such interval onto one interval
of the same length, preserving or reversing its orientation.
A slope-zero map becomes the pointed constant map. The
degree-two map on every circle, using its interval parameter
modulo endpoints, consequently commutes strictly with every act
map. It is also compatible with degenerate intervals.
The reduced cellular chain complex of each quotient is concentrated
in degree one, and this natural self-map is multiplication by two
on it. On the reduced pointed category bar it is multiplication
by two in every total degree. It thus acts by two on all positive
rational cohomology. Multiplicativity gives $2xy=4xy$, proving
$xy=0$ for any two positive-degree classes.
\end{proof}

\begin{theorem}[The categorical two-row act obstruction]
\label{C:categorical-two-row-act-cups}
Let $S$ be a finite strongly connected category. If every object
has at most two rows in its minimal Rees ideal, then every product
of rational degree-two classes on the cone of any nonempty finite
right $S$-act is zero. There is no restriction on the maximal
group of the minimal ideal or on the number of states. Dually,
at most two minimal columns at each object exclude such products
on every finite left-act cone.
\end{theorem}
\begin{proof}
Put $R=\mathbb QS$, and write $\mathbf1$ for the constant
left $S$-module. Choose a minimal idempotent at an object $a$,
let $G$ be its maximal group, and put
\[
 \bar e=\frac1{|G|}\sum_{g\in G}g,\qquad P=R\bar e.
\]
This is a projective left $R$-module. At an object $i$ its
basis consists of the $G$-averaged rows of $S(a,i)$ in the minimal
ideal. Every arrow sends these row basis vectors to row basis
vectors. Objectwise augmentation gives an exact sequence
\[
 0\longrightarrow V\longrightarrow P\longrightarrow\mathbf1
 \longrightarrow0.
\]
Here $V_i$ has dimension zero or one. Choose the usual difference
basis when there are two rows. For $f:i\to j$, its action on
the difference spaces is a scalar $\chi_f\in\{0,1,-1\}$,
with zero used whenever either space is zero. These are typed
module maps and satisfy the corresponding composition identities.

Let $F:S^{\mathrm{op}}\to\mathrm{Set}$ be the right act.
All its stalks are nonempty: strong connectivity propagates
nonemptiness along the contravariant maps. Put
$L_F=\ker(\mathbb QF\to\mathbf1_{\mathrm{right}})$.
The augmented category bar, followed by its cofiber, gives
\[
 H^2(B\mathcal C(S,F);\mathbb Q)
   \cong D\operatorname{Tor}^{R}_1(L_F,\mathbf1)
   \cong\operatorname{Ext}^{1}_R(\mathbf1,DL_F).
\]
These are the objectwise versions of
Lemma~\ref{C:act-cone-Tor}; explicitly the cofiber chain complex
is the suspension of $L_F\otimes_R^{\mathbf L}\mathbf1$.
Finite-dimensional vector-space duality gives the second equality.
Projectivity of $P$ gives a natural surjection
\[
 \operatorname{Hom}_R(V,DL_F)\longrightarrow
 H^2(B\mathcal C(S,F);\mathbb Q).
\]

Choose a lift of any class $x$ through this surjection. It is a
family of functionals $\ell_i:L_F(i)\to\mathbb Q$, with
$\ell_i=0$ if $V_i=0$, satisfying
\[
 \ell_i L_F(f)=\chi_f\ell_j\qquad(f:i\to j).
\]
Choose $s_i\in F(i)$ and set
$q_i(s)=\ell_i([s]-[s_i])$. There are rational numbers $\beta_f$
such that
\[
 q_i(F(f)s)=\chi_fq_j(s)+\beta_f.
\]
The equation follows by subtracting the value at $s_j$; one may
take $\beta_f=q_i(F(f)s_j)$. At a zero-dimensional $V_i$ the
set of values is the singleton $\{0\}$, which causes no exception.

Choose a positive rational $\delta$ for which all the finitely
many values $q_i(s)$ and $\beta_f$ lie in $\delta\mathbb Z$.
Let $I_i$ be the convex hull of $q_i(F(i))$, and enlarge that
finite set to $T_i=I_i\cap\delta\mathbb Z$. Each affine map
$t\mapsto\chi_ft+\beta_f$ carries $I_j$ into $I_i$ and the
whole grid $T_j$ into $T_i$. The typed composition identities
hold on these intervals: if $I_j$ is nondegenerate, equality
at its endpoints determines the affine map, and if it is
degenerate there is just one value to check. Thus the $T_i$
form a finite right $S$-act and $q:F\to T$ is an actual act map.

On $L_T(i)$ the coordinate functional
\[
 \ell_i^T\left(\sum_{t\in T_i}a_t[t]\right)=\sum_{t\in T_i}a_t t,
 \qquad \sum_{t\in T_i}a_t=0,
\]
is zero when $V_i=0$ and otherwise satisfies the same typed
equivariance equation. The translation disappears because the
coefficients sum to zero. Moreover $\ell_i=\ell_i^T L(q_i)$.
The natural connecting map therefore defines a class $y$ on
$B\mathcal C(S,T)$ with $x=(B\mathcal C(S,q))^*y$.
Lemma~\ref{C:categorical-grid-act-cups} gives $y^2=0$, hence
$x^2=0$. This applies to every $x$. Rational polarization proves
that every mixed degree-two product also vanishes. Taking
opposites proves the assertion about columns and left acts.
\end{proof}

\begin{corollary}\label{C:categorical-two-row-projective-sink}
The same degree-two product obstruction applies when the trivial
sink of the preceding theorem is replaced by a finite monoid $N$
whose trivial right rational module is projective, and the cross
sets form compatible $(N,S)$-actions. No finite-type hypothesis
is required for this assertion.
\end{corollary}
\begin{proof}
Proposition~\ref{C:finite-target-extreme-corner} gives an actual
quotient functor replacing $N$ by a point and every incident
left $N$-set by its orbit set. It is a rational cohomology-ring
equivalence. Apply Theorem~\ref{C:categorical-two-row-act-cups}
to the resulting right $S$-act.
\end{proof}

\begin{proposition}\label{C:strong-pair-fourteen-source-cone}
Let $S$ be a finite strongly connected category on two objects
with at most fourteen arrows. Adjoin a trivial sink with a
nonempty finite right $S$-act. If the resulting category has
at most twenty-one arrows, it cannot have type $\mathbb CP^2$.
The dual assertion holds for an adjoined trivial source.
\end{proposition}
\begin{proof}
If one object of $S$ is a retract of the other, deleting it
preserves the whole category's homotopy type and gives an
ordinary act cone with fewer than twenty-two arrows. This is
excluded by Corollary~\ref{C:act-cone-minimum-twenty-two}.
Assume henceforth that neither object is such a retract.
Both corner identities lie outside the minimal ideal of $S$.

If $H^2(BS;\mathbb Q)=0$, every degree-two class of the act cone
is relative to $BS$. Its relative quotient is a suspension, so
all degree-two products vanish. Otherwise the categorical Rees
comparison requires both total minimal row and column counts,
say $r,c$, to be at least three. If the minimal group had order
at least two, its ideal and the two outside identities would
already require at least $2\cdot3\cdot3+2=20$ arrows.
The minimal group is therefore trivial, and $rc+2\leq14$.
The only possibilities are $(r,c)=(3,3),(3,4),(4,3)$.

If each object has at most two rows, apply
Theorem~\ref{C:categorical-two-row-act-cups}.
The only remaining case has row counts $(3,1)$, total column
count three, and exactly fourteen arrows. Every arrow other
than the two identities is in the minimal band. Let $b$ be
the one-row object and choose a minimal idempotent $e_b$ there.
Every incoming arrow from the other object lies in the minimal
band and is fixed by left multiplication by $e_b$. The corner
at $b$ has no unit other than its identity. Minimal cross arrows
in the two directions have composite $e_b$.
There is no external entrance at $b$, because the added object
is a sink. Lemma~\ref{C:one-sided-corner-deletion} consequently
deletes $b$ from the whole category by a homotopy equivalence.

The remaining source monoid consists of the identity and a
$3\times c_a$ band, where $c_a=1$ or $2$. If $c_a=1$, its
classifying space is contractible because the band elements
are left zeros, so the preceding relative-suspension argument
applies. If $c_a=2$, it has only one element outside its
two-column band, and
Lemma~\ref{C:three-outside-all-two-column-act-cups} makes all
degree-two products of every act cone zero. Both alternatives
contradict the nonzero projective-plane square. Opposition
proves the dual assertion.
\end{proof}


\begin{corollary}\label{C:two-valued-representable-act-cups}
Let $S$ be a finite category with an object $a$ for which
$1\leq|S(a,i)|\leq2$ at every object $i$. For every finite right
$S$-act whose stalks are all nonempty, every rational degree-two
product of its act cone is zero.
\end{corollary}
\begin{proof}
Repeat the proof of Theorem~\ref{C:categorical-two-row-act-cups}
with the projective representable module $P=\mathbb QS(a,-)$.
Its row basis at each object has one or two elements, and
composition sends basis elements to basis elements. Thus its
augmentation kernel has the same objectwise zero- or
one-dimensional form. The Ext lift, finite rational grid, actual
act map, and scalar-degree argument apply verbatim. Strong
connectivity was used there only to construct this projective
cover and to ensure nonempty stalks, which are now explicit
hypotheses.
\end{proof}

\begin{corollary}\label{C:two-orbit-projective-cover-act-cups}
Let $S$ be a finite category, $a$ an object, $e\in\operatorname{End}_S(a)$
an idempotent, and $G$ a finite subgroup of the unit group of
$e\operatorname{End}_S(a)e$. Suppose that the right $G$-orbit
set $S(a,i)e/G$ has one or two elements at every object $i$.
For every finite right $S$-act with all stalks nonempty, all
rational degree-two products of its act cone vanish.
The conclusion remains true after replacing the trivial sink by
a monoid with projective trivial right rational module.
\end{corollary}
\begin{proof}
The averaging idempotent $\bar e=|G|^{-1}\sum_{g\in G}g$ gives
the projective module $P=\mathbb QS\bar e$. At $i$, a basis
consists of the normalized averages of the right $G$-orbits
in $S(a,i)e$. The action need not be free. Left composition
maps one orbit equivariantly onto another, and sends its normalized
average to the normalized average of that target orbit.
Consequently the augmentation $P\to\mathbf1$ is precisely
the projective basis cover used in the proof of
Theorem~\ref{C:categorical-two-row-act-cups}, with one or two
basis elements at each object. That proof applies unchanged.
For the last assertion, the actual extreme-corner orbit functor
is a rational cohomology-ring equivalence and does not increase
any cross stalk. Apply the first assertion to its act cone.
\end{proof}


% Root continuation; not input pending independent literal audit.
\subsection{Strong pairs with an adjoined trivial endpoint}

\begin{lemma}\label{C:active-corner-fixed-state}
Let $S$ be a finite strongly connected two-object category with
rectangular-band minimal ideal. Suppose that its row counts are
$(r,1)$, with the larger row set at $a$. If a finite right
$S$-act $F$ has an $\operatorname{End}_S(a)$-fixed state in $F(a)$,
every degree-two rational product of its act cone is zero.
\end{lemma}
\begin{proof}
In the natural categorical balanced-tensor presentation of
$H_2(BS;\mathbb Q)$, the row augmentation at the other object
is zero. Thus every class is represented by a band-grid cycle
at $a$: inclusion of the corner monoid at $a$ surjects onto
$H_2(BS;\mathbb Q)$. The chosen fixed state gives an actual
functor from that corner monoid to the element category of $F$.
Its composite with the projection to $S$ is that corner inclusion.
Consequently $H_2$ of the element category surjects onto $H_2(BS)$.
The cofiber sequence makes every degree-two class of the act cone
relative to $BS$. The relative quotient is a suspension, so its
positive-degree cup products, and their absolute images, vanish.
\end{proof}

\begin{theorem}\label{C:strong-pair-trivial-endpoint21}
A category obtained by adjoining a trivial sink or source to a
strongly connected two-object category cannot have type
$\mathbb CP^2$ with at most twenty-one arrows.
\end{theorem}
\begin{proof}
Take opposites to make the new object a sink, and call the strong
source $S$, its arrow count $s$, and its right act $F$. The two
stalks of $F$ are nonempty. A retract object of $S$ can be deleted
from the whole category, giving an ordinary act cone with fewer
than twenty-two arrows, which is impossible. Assume no such
retract exists.

Proposition~\ref{C:proper-strong-pair-at-least-seventeen} treats
$s\geq17$, and Proposition~\ref{C:strong-pair-fourteen-source-cone}
treats $s\leq14$. Hence $s=15$ or $16$.
If $H^2(BS;\mathbb Q)=0$, the relative-suspension argument
already gives zero products. Otherwise both total minimal row
and column counts are at least three. A nontrivial minimal group
would require at least twenty arrows including the two identities.
Thus the group is trivial. If every object has at most two rows,
Theorem~\ref{C:categorical-two-row-act-cups} applies.

The only remaining shape is a rectangular band with total row
count four and column count three, with row counts $(3,1)$.
Indeed a larger total row count would require at least
$5\cdot3+2=17$ arrows, while total column count at least four
would require at least $4\cdot4+2=18$. Write $a$ for the
three-row object and $b$ for the one-row object. Nonzero second
homology forces the column counts to be $(2,1)$ at $(a,b)$:
in the other distribution each objectwise row-column augmentation
tensor is zero. The minimal band and identities use fourteen
arrows. There are at most two extra internal arrows.

The corner at $b$ has a minimal two-sided zero $e_b$.
If every incoming arrow $a\to b$ belongs to the minimal ideal,
it is fixed by $e_b$, and
Lemma~\ref{C:one-sided-corner-deletion} deletes $b$ from the whole
category. The result is an ordinary act cone with fewer than
twenty-two arrows, again impossible.
We may therefore assume that at least one extra arrow goes
from $a$ to $b$.

If there is no extra endomorphism at $a$, its endomorphism monoid
is precisely the identity-adjoined $3\times2$ band. The left row
augmentation module $V$ in the projective sequence
$0\to V\to P\to\mathbf1\to0$ is supported only at $a$.
It is two-dimensional there, every nonidentity corner element
acts by zero, and every cross map is zero. Consequently
$V=L\oplus L$ for a one-dimensional rational $S$-module $L$.
Using the natural act-cone Tor formula and dimension shifting,
\[
 H_4(B\mathcal C(S,F);\mathbb Q)
   \cong\operatorname{Tor}^{\mathbb QS}_3(L_F,\mathbf1)
   \cong\operatorname{Tor}^{\mathbb QS}_2(L_F,V)
   \cong\operatorname{Tor}^{\mathbb QS}_2(L_F,L)^{\oplus2}.
\]
Its dimension is even, contradicting the fourth Betti number one
of a projective plane. This covers every remaining $s=15$ case.

Finally suppose an extra endomorphism at $a$ exists. Along with
the extra incoming arrow at $b$ it uses both extra arrows, so
$s=16$. Every internal arrow $b\to a$ is minimal and is fixed
on the right by $e_b$, because the column set at $b$ is a singleton.
The total state count satisfies $|F(a)|+|F(b)|\leq4$.
If $|F(b)|=1$, its external outgoing arrow is also fixed by
$e_b$, and the outgoing version of the one-sided deletion lemma
deletes $b$. If $|F(a)|=1$, apply
Lemma~\ref{C:active-corner-fixed-state}.
The only remaining positive pair is $(|F(a)|,|F(b)|)=(2,2)$,
excluded by Lemma~\ref{C:two-states-objectwise-act-cups}.
Every case is impossible.
\end{proof}

% Root and D completed independent literal audits before input.
\subsection{Projective opposite corners of a strong pair}

\begin{proposition}\label{C:finite-target-strong-pair-twenty}
Let $S$ be a finite strongly connected two-object category and
$F:S^{\mathrm{op}}\to\mathrm{Set}$ a finite right act with both
stalks nonempty. Suppose its act category has at most twenty
arrows. Its classifying space cannot satisfy all four conditions
of Theorem~\ref{C:finite-target-small-act-twenty}: simple
connectivity and rational projective-plane homology, an integral
degree-two generator with nonzero primitive square modulo
torsion, nullity of that generator on every finite group, and
Betti numbers $(1,0,1)$ in degrees two through four over
$\mathbb F_{65521}$.
\end{proposition}
\begin{proof}
Suppose all four conditions hold. Write $s=|S|$ and
$t=|F(a)|+|F(b)|$, so $s+t+1\leq20$ and $t\geq2$.
If one object of $S$ is a retract of the other, delete it
from the entire category. This is an actual homotopy equivalence
to an ordinary act cone with fewer arrows. All four conditions
survive that equivalence, contradicting
Theorem~\ref{C:finite-target-small-act-twenty}.
Assume no such retract exists.

If $s\geq16$, the two stalks have total size at most three.
Lemma~\ref{C:two-states-objectwise-act-cups} gives zero rational
products, a contradiction. We may assume $s\leq15$.
If $H^2(BS;\mathbb Q)=0$, all degree-two classes of the act
cone are relative to $BS$, and the suspension quotient gives
the same contradiction. Otherwise total minimal row and column
counts are both at least three. The two corner identities are
outside the minimal ideal. A nontrivial minimal group would
therefore require at least twenty arrows in $S$, so that group
is trivial.

If every object has at most two minimal rows, apply
Theorem~\ref{C:categorical-two-row-act-cups}.
The only remaining possibility has four total rows, three total
columns, and row counts $(3,1)$ at objects $(a,b)$.
Indeed five rows would already cost seventeen arrows, and
four columns together with four rows would cost eighteen.
Nonzero second homology forces column counts $(2,1)$ at
$(a,b)$, by the objectwise balanced-tensor presentation.
The band and its two outside identities cost fourteen arrows;
there is at most one extra arrow in $S$.

If every arrow $a\to b$ is minimal, the minimal idempotent
$e_b$ fixes every incoming arrow. There is no external entrance
at $b$. Lemma~\ref{C:one-sided-corner-deletion} deletes $b$
by an actual homotopy equivalence. The resulting ordinary act
cone still has at most twenty arrows and still satisfies all
four conditions, contradicting
Theorem~\ref{C:finite-target-small-act-twenty}.

Otherwise the sole extra internal arrow goes from $a$ to $b$.
The corner at $a$ is exactly the identity-adjoined $3\times2$
band. In the projective augmentation sequence
$0\to V\to P\to\mathbf1\to0$ over $\mathbb QS$, the row
augmentation module has dimension two at $a$ and zero at $b$.
Every nonidentity corner element acts by zero on it, and every
cross map is zero. Hence $V=L\oplus L$ for a one-dimensional
module $L$. Put $L_F=\ker(\mathbb QF\to\mathbf1_{\rm right})$.
The natural categorical act-cone Tor formula and dimension
shifting give
\[
 H_4(B\mathcal C(S,F);\mathbb Q)
 \cong\operatorname{Tor}^{\mathbb QS}_3(L_F,\mathbf1)
 \cong\operatorname{Tor}^{\mathbb QS}_2(L_F,L)^{\oplus2}.
\]
Its dimension is even, contradicting the fourth Betti number
one. This exhausts the possibilities. Every use of the ordinary
act theorem followed an actual homotopy equivalence; no finite
CW type was assumed for an intermediate space.
\end{proof}

\begin{corollary}\label{C:strong-pair-projective-endpoint21}
Let $C$ have three objects, with a strongly connected two-object
source component $S$ and a remaining sink corner $N$. Suppose
the trivial right $\mathbb QN$-module is projective.
If $|C|\leq21$, then $BC\not\simeq\mathbb CP^2$.
Dually, a projective trivial left module at a source corner
gives the same conclusion when the strong pair is the sink.
\end{corollary}
\begin{proof}
If $|N|=1$, apply
Theorem~\ref{C:strong-pair-trivial-endpoint21}.
Otherwise suppose $BC\simeq\mathbb CP^2$ and use the actual
extreme-corner quotient of
Proposition~\ref{C:finite-target-extreme-corner}.
It replaces $N$ by a point and each incoming left $N$-set by
its orbit set, producing a right $S$-act cone $Y$ with
\[
 |\bar C|\leq |C|-|N|+1\leq20.
\]
The quotient is onto on fundamental groups and is a rational
cohomology-ring equivalence. It has a homotopy left inverse
$h:Y\to BC$. If $u$ generates $H^2(BC;\mathbb Z)$, put
$v=h^*u$. The integral argument in
Corollary~\ref{C:finite-target-corner-CP2} applies unchanged:
$Y$ is simply connected, $H^2(Y;\mathbb Z)\cong\mathbb Z$,
$v$ is its generator, and $v^2$ is primitive modulo torsion.
The identities $Q^*v=u$ and $Q^*(v^2)=u^2$ prove the last
two assertions by the split surjections in degrees two and four.
For every finite group $H$ and map $z:BH\to Y$, Miller's
theorem makes $hz$ null-homotopic, so $z^*v=0$ integrally.
Finally the minimal local group of $N$ has order at most
$|N|\leq21$, so $65521$ is a good prime for the same quotient.
Thus $Y$ satisfies all four conditions of the preceding
proposition, a contradiction. Taking opposites proves the dual.
\end{proof}

% Root draft; uninput pending independent literal audit.
\subsection{The seventeen-arrow three-object strong-component cases}

\begin{lemma}\label{C:small-corner-two-states-per-object}
Let $S$ be a finite category and let $N$ be a monoid of order
at most four. Form a collage with full source block $S$, sink
corner $N$, and compatible cross sets $F(i)$ satisfying
$1\leq|F(i)|\leq2$ for every object of $S$. Its classifying
space cannot have type $\mathbb CP^2$.
\end{lemma}
\begin{proof}
If the trivial right rational $N$-module is projective, the
actual orbit-collapse functor of
Proposition~\ref{C:finite-target-extreme-corner} is a rational
cohomology-ring equivalence to an ordinary categorical act cone
with at most two states at each object. Apply
Lemma~\ref{C:two-states-objectwise-act-cups}.

Otherwise the small-corner Rees argument in
Theorem~\ref{C:small-opposite-corner-three-states} says that the
minimal ideal of $N$ is a right-zero band of size two or three.
Its left actors are retractions onto the same subset $U(i)$ of
each $F(i)$. On a set of at most two elements these retractions
are identical: if the common image has one element they are
the same constant map, and if it has two they are identities.
Choose any minimal element $p$ and write $P_i$ for this common
map. Every $N$-actor fixes $U(i)$ pointwise. The identities
$np=p$ and $pn$ minimal therefore show
$nP_i=P_i=P_in$. The family $P_i$ also commutes with all
right $S$-act maps. It is consequently an actual biset retraction
onto the categorical act $U$ with trivial left $N$-action.

Let $X$ be the original collage space and $Y$ the core collage
on $U$. This gives an inclusion and retraction $Y\leftrightarrows X$.
The monoid $N$ has a right zero, so $BN$ is contractible. The
actual two-block pushout identifies $Y$, up to homotopy, with
the ordinary categorical act cone on $U$: its edge correspondence
is the product of the $S$-element space with $BN$, and projection
is an equivalence on every space of the pushout diagram.
Every positive rational product of $Y$ is therefore zero.

Both block spaces are contained in $Y$. The mixed-nerve
description makes the relative quotient $X/Y$ a suspension of
the quotient of their edge correspondence spaces. Suppose
$X\simeq\mathbb CP^2$. If $H^2(Y;\mathbb Q)=0$, all its
degree-two classes lift from that suspension and square to zero.
If $H^2(Y;\mathbb Q)\ne0$, the split injection from the
retraction identifies it with the one-dimensional $H^2(X;\mathbb Q)$,
whose generator again has zero square by pullback from $Y$.
Both cases contradict the projective-plane square.
\end{proof}

\begin{theorem}\label{C:seventeen-three-strong-component}
No three-object category with at most seventeen arrows and a
nontrivial strongly connected component can have classifying
space of type $\mathbb CP^2$.
\end{theorem}
\begin{proof}
A disconnected category is impossible. If an object is a retract
of another, deleting it removes at least three arrows and gives
a category with at most fourteen, contrary to
Theorem~\ref{C:CP2-lower17}. A strongly connected whole category
is excluded by Corollary~\ref{C:strong-category-CP2-minimum23}.
Thus there is a proper strong pair $S$ and one remaining monoid
corner $N$. Take opposites so that $S$ is the source block.
Write $s=|S|$, $m=|N|$ and $t=|F(a)|+|F(b)|$ for its two
nonempty cross stalks. Then $s\geq6$, $t\geq2$, and
$s+m+t\leq17$.

If $m=1$, apply Theorem~\ref{C:strong-pair-trivial-endpoint21}.
If both block spaces have zero second rational cohomology,
Lemma~\ref{C:two-block-zero-H2-products} excludes their square.

Suppose first that $H^2(BS;\mathbb Q)\ne0$.
The categorical Rees count gives $s\geq11$, so $m+t\leq6$.
For $m\geq3$, one has $m\leq4$ and $t\leq3$, and the
preceding lemma applies to the two stalks. If $m=2$, its
trivial right rational module is projective and $s\leq13$.
Nonzero second homology then forces a trivial-group $3\times3$
minimal band in $S$. Each object has at most two rows, and
Corollary~\ref{C:categorical-two-row-projective-sink} applies.

It remains that $H^2(BS;\mathbb Q)=0$ and
$H^2(BN;\mathbb Q)\ne0$. The latter requires $m\geq5$ by
the small-monoid classification, hence $s\leq10$.
If $t\leq3$, at least one cross stalk is a singleton.
Its state is fixed by $N$ and gives an actual section
$BN\to E$ of the edge-space projection $E\to BN$.
Since the other block has zero second cohomology, the whole
vertex degree-two group injects into that of $E$.
All degree-two classes of the pushout are therefore relative
to the vertex spaces and pull back from a suspension; their
products vanish.
We may consequently assume $t\geq4$, which forces $s\leq8$.

The minimal ideal of $S$ now has trivial group. If its total
row count is two, there is one row at each object and
Lemma~\ref{C:one-row-source-block-collapse} collapses $S$
to one trivial object. This is an actual homotopy-equivalent
ordinary left-act cone over $N$, with fewer than twenty-two
arrows, which is impossible.
The only remaining shape has total row and column counts
$(3,2)$ and $s=8$. Its row distribution is $(2,1)$ and its
column distribution is $(1,1)$; every nonidentity arrow of
$S$ is minimal. At its one-row object, every incoming internal
cross arrow is fixed by the minimal idempotent, and there are
no external incoming arrows. The one-sided deletion lemma
deletes that object from the entire collage. It removes at
least two corner arrows, three internal cross arrows, and
one external cross arrow. The resulting category has at most
eleven arrows, contrary to the global lower bound.
This exhausts all cases.
\end{proof}

% Root proof; independent full literal audits by D and B passed.
\subsection{The eighteen-arrow three-object strong-component cases}

\begin{lemma}[A singleton and a three-state right-zero stalk]
\label{C:right-zero-singleton-three-stalks}
Let $S$ be a finite two-object category, let
$R=\{1,p,q\}$ be the identity-adjoined right-zero two-band,
and let compatible left $R$-sets and right $S$-actions form
the cross sets of a collage with source block $S$ and sink
corner $R$. Suppose the two stalk sizes are one and three.
If all degree-two rational products vanish on every ordinary
finite right $S$-act cone, then a one-dimensional second
rational cohomology group of this collage has zero square.
\end{lemma}
\begin{proof}
Use the natural right-$S$ diagram of left-$R$ action graphs,
as in Lemma~\ref{C:right-zero-action-graph}. The singleton
stalk gives a point. At the three-state stalk, either the graph
is a forest, or it is a point with two loops, or it is a loop
together with an isolated vertex. Indeed a common core of
three elements gives a forest with no edges. A core of two
elements gives one edge, producing either a forest or a loop
and an isolated vertex. A core of one element gives two loops.
The outside edges join the two images under the minimal
retractions.

If that graph is a forest, the actual natural graph-to-component
map is an objectwise homotopy equivalence. The category bar
and its cofiber therefore identify the collage with the
ordinary act cone on the finite component diagram. Its
degree-two products vanish by hypothesis.

In either loop case, double every loop and fix every vertex.
This is a strict natural self-map of the whole $S$-diagram.
An endomorphism at the three-state stalk sends a loop to a
loop or a vertex; maps to or from the singleton stalk fix the
selected vertices and collapse loops. There is no other
nontrivial edge whose parametrization would require a choice.
Thus the degree-two circle maps commute with every actual
diagram map.

Let $q:X\to Y$ be the actual cofiber comparison to the ordinary
component-act cone. The graph-chain kernel has homology only
in degree one, naturally equal to the graph first-homology
module. The self-map is multiplication by two on that module
and the identity on the component diagram. Apply the free
category bar and its homology exact sequence. Exactly as in
Lemma~\ref{C:graph-component-eigenvalues}, on each rational
cohomology group its pullback satisfies
\[
 (f^*-1)(f^*-2)=0,\qquad
 \ker(f^*-1)=\operatorname{im}(q^*).
\]
This argument uses the objectwise natural graph-chain maps,
so it applies to a finite category in place of a monoid.
On the one-dimensional second cohomology the eigenvalue is
one or two. In the first case the class pulls back from $Y$
and has zero square. In the second its square has eigenvalue
four, forbidden by the displayed polynomial. This proves
the assertion.
\end{proof}

\begin{lemma}[Two-state actions of the five-element sphere band]
\label{C:sphere-band-two-state-strong-source}
Let $N$ be the identity-adjoined $2\times2$ rectangular band,
and let $S$ be a finite strongly connected category with
contractible classifying space. Suppose every cross stalk
of an $(N,S)$-collage has two elements, and $H^2(BS;\mathbb Q)=0$.
Then the collage cannot have the rational cohomology ring
of $\mathbb CP^2$.
\end{lemma}
\begin{proof}
If one stalk has an $N$-fixed state, its action category
contains an actual copy of $BN$ over $BN$. The cross-arrow
space in the mixed-nerve pushout therefore surjects onto
$H_2(BN;\mathbb Q)$. Since $H^2(BS;\mathbb Q)=0$, every
degree-two class of the collage is relative to its two
vertex blocks. Their quotient is a suspension, so its
degree-two products vanish.

Otherwise no stalk has an $N$-fixed state. Every minimal
band element acts idempotently on a two-element set, hence
as the identity or as a constant. If $A_{i_0j_0}=1$ for
the images $A_{ij}$ of the minimal actors, the equations
$A_{i_0j_0}A_{kl}=A_{i_0l}$ first make every image depend
only on its column. Then
$A_{ij}A_{i_0j_0}=A_{ij_0}=1$ gives $A_{ij}=1$ for
every pair. This gives fixed states, contrary to the
present case. Thus every minimal actor is constant.
The equations then force its constant value to depend
only on its row. The two rows must give the two distinct
states, or a common fixed state would again result.

Consequently every stalk is canonically the same row
$N$-set, and every map between stalks is the unique
row-preserving bijection. The whole right $S$-action on
this row set is therefore trivial after these canonical
identifications. The cross-arrow space is
$BS\times B(N\ltimes\{0,1\})$. Contractibility of $BS$
and homotopy invariance of the mixed-nerve pushout identify
the collage with the ordinary left $N$-act cone on this
two-element set. Its degree-two products vanish by the
two-state act theorem.
\end{proof}

\begin{theorem}\label{C:eighteen-three-strong-component}
Every finite category with three objects, at most eighteen
arrows, and a nontrivial strongly connected component is
excluded from type $\mathbb CP^2$.
\end{theorem}
\begin{proof}
Assume such a category is a projective plane. It is connected.
A retract object can be deleted, giving a two-object category
with at most fifteen arrows, contrary to the global lower
bound. A wholly strongly connected category is excluded by
the exact strongly connected minimum twenty-three.
Take opposites so that the proper strong pair $S$ is the
source and the remaining corner $N$ the sink. Put
$s=|S|$, $m=|N|$, and let $t$ be the total cross-set size.
Then $s\geq6$, $t\geq2$, and $s+m+t\leq18$.
Projective trivial right modules at $N$ are excluded by
Corollary~\ref{C:strong-pair-projective-endpoint21}.
Thus $m\geq3$ and that module is nonprojective.

If both vertex blocks have zero second rational cohomology,
Lemma~\ref{C:two-block-zero-H2-products} applies.
Suppose first $H^2(BS;\mathbb Q)\ne0$. The categorical Rees
budget gives $s\geq11$. If $t=2$, every stalk is a singleton
and the collage is a join, so its positive products vanish.
For $t\geq3$ the budget leaves $m\leq4$. If both stalks
have size at most two, use
Lemma~\ref{C:small-corner-two-states-per-object}.
The only remaining possibility is
\[
 (s,m,t)=(11,3,4),\qquad
 (|F(a)|,|F(b)|)=(1,3)\text{ or }(3,1).
\]
The nonprojective three-element corner is $R=\{1,p,q\}$
with a right-zero minimal two-band. The strong pair has
trivial minimal group, total row and column counts three,
and only its two identities outside the minimal band.
It therefore has at most two minimal rows at each object.
Theorem~\ref{C:categorical-two-row-act-cups} and the first
lemma exclude this last case.

It remains that $H^2(BS;\mathbb Q)=0$ and
$H^2(BN;\mathbb Q)\ne0$, so $m\geq5$.
A singleton cross stalk supplies an $N$-fixed state and
an actual section over $BN$ in the cross-arrow space.
As in the preceding lemma, every degree-two class is
relative to the two blocks and has zero square. Thus both
stalks have size at least two, $t\geq4$, and $s\leq9$.
The minimal group of the nonretract strong pair is trivial.
If it has one minimal row at each object, the actual
source-block collapse of
Lemma~\ref{C:one-row-source-block-collapse} gives an ordinary
left $N$-act cone with fewer than twenty-two arrows,
contrary to its exact minimum.

The only remaining Rees shape has total row count three,
total column count two, and row counts $(2,1)$ at $(a,b)$.
Its two column counts are both one. The minimal band and
identities cost eight arrows. If every incoming arrow
$a\to b$ is minimal, delete $b$ using
Lemma~\ref{C:one-sided-corner-deletion}; the result has fewer
than seventeen arrows, a contradiction. Otherwise the sole
extra arrow is a nonminimal $a\to b$, and the budget forces
\[
 s=9,\qquad m=5,\qquad |F(a)|=|F(b)|=2.
\]
The five-element corner with positive second rational
cohomology is the identity-adjoined $2\times2$ band.
The strong source has one minimal column per object and
trivial minimal group, so its classifying space is
contractible: the dual one-column collapse applied to $S$
alone supplies the actual equivalence to a point.
The preceding two-state sphere-band lemma now gives the
contradiction. This exhausts every budget case.
\end{proof}

% Root proof; D independent complete literal audit passed; acyclic excess-eight root+B gate passed.
\subsection{An excess-eight bound without a support restriction}

\begin{proposition}\label{C:strong-pair-excess-eight-through21}
Let $C$ be a connected finite category with $k$ objects and $n$ arrows,
including identities, such that
\[
 n\leq\min\{21,2k+7\}.
\]
Suppose that no object is a retract of another. If $C$ has a proper
strongly connected component on two objects, then
$BC\not\simeq\mathbb CP^2$.
\end{proposition}
\begin{proof}
Write $S$ for such a component, $s=|S|$, and $\ell=k-2\geq1$.
Both identities of $S$ lie outside its minimal ideal by the
minimal-group inverse adjustment in Lemma~\ref{C:finite-category-Rees}.
Consequently $s\geq6$. Contracting $S$ in the undirected support graph
leaves $\ell+1$ vertices. Its at least $\ell$ support edges include
an interface with $S$, which supplies arrows at both objects of $S$.
Together with the $\ell$ outside identities this gives
\[
 n\geq s+2\ell+1,\qquad n\leq2\ell+11,\qquad s\leq10.
\]
The rational categorical Rees comparison then gives
$H^2(BS;\mathbb Q)=0$: nonzero second homology requires both total
row and total column counts at least three, and therefore at least
nine minimal-ideal arrows and the two additional identities.

\paragraph{An extreme component.}
Take opposites to make an extreme $S$ a source. Let
$D=C\setminus S$ have $q$ connected components. If
$H^2(BD;\mathbb Q)\ne0$, Lemma~\ref{C:arrow-excess-two-H2}
forces at least $2t+2$ arrows in one connected component with $t$
objects. Each other component with $t'$ objects has at least
$2t'-1$ arrows. Thus $|D|\geq2\ell-q+3$.
Every component of $D$ meets $S$ in at least two cross arrows,
so
\[
 n\geq s+2\ell+q+3\geq s+2\ell+4.
\]
For $s\geq8$ this contradicts $n\leq2\ell+11$.
Both vertex blocks consequently have zero rational second
cohomology in that case, and
Lemma~\ref{C:two-block-zero-H2-products} makes every degree-two
product zero.

For $s=6$ or $7$, the minimal group is trivial and there is one
minimal row and one minimal column at each object.
Lemma~\ref{C:one-row-source-block-collapse} replaces $S$ by a point
through an actual homotopy equivalence. It saves $s-1$ internal
arrows and at least one cross arrow: for any nonempty receiving
stalk family, its orbit set is an image of either of the two
nonempty stalks, so its size is at most their smaller size.
No other hom-set increases. The resulting category has at most
$n-s\leq15$ arrows, contrary to Theorem~\ref{C:CP2-lower17}.

\paragraph{A nonextreme component.}
There are outside objects $a,b$ with $a\to S\to b$.
They are distinct, and composition supplies $a\to b$.
Thus the contracted support has at least $\ell+1$ edges,
at least two of which meet $S$. Let $P_{\rm ext}$ be the sum
of $|C(d,i)|-1$ and $|C(i,d)|-1$ over nonempty hom-sets with
$i\in S$ and $d\notin S$. Counting as above gives
\[
 n\geq s+2\ell+3+P_{\rm ext},\qquad
 P_{\rm ext}\leq8-s,\qquad s\leq8.
\]
For $s=6$ or $7$, let $e_i$ be the unique minimal-ideal
endomorphism at each of the two objects $i$. It is a zero
of its corner. Consider the four conditions asserting that
$e_i$ fixes every incoming cross arrow, or every outgoing
cross arrow, at either object. An extra internal arrow can
violate at most two conditions, and there are at most $s-6$
such arrows. An external violation requires a nonsingleton
hom-set and affects only one of the four conditions, so
the total number of violated conditions is at most
\[
 2(s-6)+P_{\rm ext}\leq s-4\leq3.
\]
Choose a condition that holds. The minimal arrows to and from
the other object have composite $e_i$ at the chosen object.
The unit hypothesis of Lemma~\ref{C:one-sided-corner-deletion}
holds because $e_i$ is a zero: for a unit $h$ and an incoming
$g$ fixed by $e_i$, one has $hg=he_i g=e_i g=g$; the outgoing
version is dual. That lemma deletes the object by a homotopy
equivalence. Its corner contributes at least two deleted
arrows, the other object contributes an arrow in each direction,
and nonextremality contributes an external entrance and an
external exit. At least six arrows disappear, leaving at most
fifteen, again impossible.

For $s=8$ the budget forces $P_{\rm ext}=0$.
The minimal group is trivial, and the possible total row and
column counts are $(2,2)$, $(3,2)$, or $(2,3)$.
There is therefore one minimal row per object or one minimal
column per object. All external hom-sets are singletons, so
all external entrances, respectively exits, are fixed by the
local minimal idempotents. Lemma~\ref{C:one-row-fixed-entrance-collapse}
or its dual replaces $S$ by a point through an actual
homotopy equivalence. Its eight internal arrows become one,
and no other hom-set increases. The result has at most
fourteen arrows. This exhausts the possibilities.
\end{proof}

\begin{theorem}\label{C:all-category-excess-eight-through21}
Every finite category with $k$ objects and
$n\leq\min\{21,2k+7\}$ arrows has classifying space not
homotopy equivalent to $\mathbb CP^2$.
\end{theorem}
\begin{proof}
Assume otherwise; the category is connected. Delete retract
objects successively. Each deletion is an actual homotopy
equivalence and removes at least three arrows (the deleted
identity and the two retract arrows). If $k'$ and $n'$ are
the remaining counts, this preserves
$n'\leq21$ and $n'\leq2k'+7$.

Suppose the remaining category has a proper strong component
on $t\geq3$ objects. Every identity of this component lies
outside its minimal ideal, so its internal arrow count is
at least $t^2+t$. There are $\ell=k'-t\geq1$ outside objects.
The connected support after contracting the component has
at least $\ell$ edges, including an interface which expands
to arrows at all $t$ objects of the component. With the
$\ell$ outside identities, the total is at least
\[
 n'\geq t^2+t+2\ell+(t-1)
       =2k'+t^2-1\geq2k'+8,
\]
a contradiction. A wholly strongly connected category is
excluded by the exact minimum twenty-three in
Corollary~\ref{C:strong-category-CP2-minimum23}.
A proper strong pair is excluded by the preceding proposition.
Otherwise every strong component is a single object and the
support is acyclic. Theorem~\ref{C:acyclic-arrow-excess-eight}
then excludes the remaining category. All cases are impossible.
\end{proof}

\begin{corollary}\label{C:CP2-object-bounds-through21}
A category of type $\mathbb CP^2$ with seventeen arrows has
at most four objects. With eighteen or nineteen arrows it
has at most five objects; with twenty or twenty-one arrows
it has at most six objects.
\end{corollary}
\begin{proof}
Apply the preceding theorem with, respectively,
$k\geq5$, $k\geq6$, and $k\geq7$.
\end{proof}

% Root and B full literal audits passed. No numerical table premise is used.
\subsection{A ninth excess with a proper strong pair}

\begin{proposition}\label{C:strong-pair-excess-nine-through21}
Let $C$ be a connected finite category with $k$ objects and $n$ arrows.
Suppose no object is a retract of another, $C$ has a proper strongly
connected component on two objects, and
\[
 n\leq\min\{21,2k+8\}.
\]
Then $BC$ is not homotopy equivalent to $\mathbb CP^2$.
\end{proposition}
\begin{proof}
The cases with $n\leq16$ follow from Theorem~\ref{C:CP2-lower17}.
Write $S$ for the strong pair, $s=|S|$, and $\ell=k-2\geq1$.
The no-retract hypothesis puts both identities outside its minimal
ideal, so $s\geq6$. Contracting $S$ in the undirected support gives
at least $\ell$ edges, and at least one interface expands into
arrows incident with both objects of $S$. Counting the outside
identities therefore gives
\[
 n\geq s+2\ell+1,\qquad n\leq2\ell+12,\qquad s\leq11.
\]
For $s\leq10$, the categorical Rees comparison gives
$H^2(BS;\mathbb Q)=0$: nonzero second homology requires total
row and column counts at least three, hence at least nine ideal
arrows and the two additional identities.

\paragraph{An extreme pair.}
By opposition suppose $S$ is a source. Let the complement $D$ have
$q$ connected components, and let $T$ be the total number of arrows
from $S$ to $D$. Every component receives at least two such arrows.

For $s=6,7$, the minimal group is trivial with one row and one
column at each object. The actual source-block collapse of
Lemma~\ref{C:one-row-source-block-collapse} saves $s-1$ internal
arrows and at least one interface arrow. Its result has at most
$n-s\leq15$ arrows, which is impossible.

Suppose $8\leq s\leq10$. If $H^2(BD;\mathbb Q)=0$, the
two-block cup lemma finishes. Otherwise the excess-two lemma gives
\[
 |D|\geq2\ell-q+3,\qquad
 n\geq s+2\ell+q+3\geq s+2\ell+4.
\]
The upper bound forces $s=8$, $q=1$, $|D|=2\ell+2$ and $T=2$.
There is exactly one external arrow from each object of $S$, and
all other external hom-sets at $S$ are empty. The minimal group of
$S$ is trivial, and its total row/column counts are $(2,2)$,
$(3,2)$ or $(2,3)$. It has one minimal row per object or one
minimal column per object. All external entrances and exits are
singleton or empty, and hence fixed by its minimal idempotents.
Lemma~\ref{C:one-row-fixed-entrance-collapse}, or its opposite,
replaces $S$ by a point through an actual homotopy equivalence,
without increasing any other hom-set. At least seven arrows
disappear, leaving at most fourteen.

For $s=11$, the elementary connectivity bound
\[
 n\geq11+(2\ell-q)+2q=2\ell+11+q
\]
forces $q=1$, $|D|=2\ell-1$ and $T=2$.
Thus $D$ has trivial corners and singleton hom-sets on an
undirected tree; its classifying space is contractible.
Both interface arrows end at a single object $w$ of $D$,
because an interface meets both objects of $S$. There is no
cross arrow out of $w$ in $D$, since composing it with the two
interface arrows would contradict $T=2$. The full subcategory
on $S\cup\{w\}$ has terminal object $w$.
The nerve of $C$ is exactly the union of that contractible
subcategory nerve and $ND$, meeting along the point $w$.
This is a union of CW subcomplexes and is contractible.
This excludes every extreme case.

\paragraph{A nonextreme pair.}
Choose outside objects $a,b$ with $a\to S\to b$.
Composition supplies $a\to b$. The contracted undirected
support therefore has at least $\ell+1$ edges, including two
interfaces. Let $P_{\rm ext}$ be the parallel excess summed
over all individual hom-sets between $S$ and its complement.
The count is
\[
 n\geq s+2\ell+3+P_{\rm ext},\qquad
 P_{\rm ext}\leq9-s,\qquad s\leq9.
\]

For $s=6,7$, the unique minimal-ideal endomorphism $e_i$ at
each object is a zero of its corner. There are four possible
one-sided fixed-arrow conditions: all entrances or all exits
at either object are fixed by its $e_i$. Internal extra arrows
can violate at most $2(s-6)$ of these conditions, and external
parallel excess can violate at most $P_{\rm ext}$ more.
This is the same literal four-condition count proved in
Proposition~\ref{C:strong-pair-excess-eight-through21}.
Consequently a condition holds whenever
\[
 2(s-6)+P_{\rm ext}<4.
\]
The one-sided deletion lemma applies. Its unit hypothesis holds
because a unit fixes every arrow fixed by a zero idempotent.
The deleted object removes at least its two corner arrows,
the two internal cross arrows, and an external entrance and
exit. The result has at most $n-6\leq15$ arrows.
This treats $s=6$ and $s=7$ with $P_{\rm ext}\leq1$.

For $s=8,9$ and $P_{\rm ext}=0$, the minimal group is trivial,
and the minimal band has one row per object or one column per
object, as its size is at most seven. All external hom-sets
are singleton or empty. The fixed-entrance block collapse or
its dual applies as above, saving at least seven arrows.

The only remaining numerical possibilities are
\[
 (s,P_{\rm ext})=(7,2),\qquad(8,1).
\]
In either case equality holds throughout the lower-bound count:
$n=2\ell+12$, all outside corners are trivial, every outside
cross hom-set is a singleton, the contracted support has exactly
$\ell+1$ edges, and exactly two of those edges meet $S$.
Its unique cycle is the triangle $a,S,b$. All other contracted
edges form trees attached to that triangle. They cannot attach
directly to $S$, since there are only two interfaces.

Every nonempty such tree has a leaf outside $\{a,S,b\}$.
It is an object with trivial corner and a single singleton
hom-set incident with another outside object of trivial corner.
The singleton cofinal deletion removes it by an actual homotopy
equivalence. Repeating deletes all $\ell-2$ extra outside objects,
two arrows each time. The four-object category that remains has
\[
 n-2(\ell-2)=16
\]
arrows, and is excluded by Theorem~\ref{C:CP2-lower17}.
If $\ell=2$ there was no tree to delete and the same final bound
holds already. This completes the nonextreme cases.
\end{proof}

\begin{corollary}\label{C:eighteen-five-acyclic-scope}
A category of type $\mathbb CP^2$ with eighteen
arrows and at least four objects has either four objects, or
five objects and acyclic support. In the five-object case its
corner, parallel and support-cycle excesses satisfy
$D+P+R=9$.
\end{corollary}
\begin{proof}
The accepted object bound gives at most five objects. A proper
retract deletion would leave at most fifteen arrows, so no
retract exists. Proper strong components on at least three
objects are excluded through twenty-one arrows by
Proposition~\ref{C:large-SCC-minimal-candidate21}, and a wholly
strongly connected category needs at least twenty-three by
Corollary~\ref{C:strong-category-CP2-minimum23}.
For five objects the preceding proposition excludes every
proper strong pair, because $18=2\cdot5+8$. Acyclicity follows.
The excess identity is $18-(2\cdot5-1)=9$.
\end{proof}

% Root proof; B independent complete literal audit passed for both statements.
\subsection{A projective central corner in a star}

\begin{theorem}\label{C:star-projective-center22}
Let $C$ have outward star support with at least two leaves and
at most twenty-two arrows. Suppose the trivial left module of
its central endomorphism monoid is projective over its rational
monoid algebra. Then $BC\not\simeq\mathbb CP^2$.
For an inward star the same conclusion holds with the trivial
right module of the central monoid in place of the left module.
\end{theorem}
\begin{proof}
Assume $BC\simeq\mathbb CP^2$. Let $M$ be the central monoid,
let $N_i$ be the leaf monoids, and let ${}_{N_i}S_i{}_M$ be
the nonempty cross bisets. Apply the initial-object version
of Proposition~\ref{C:finite-target-extreme-corner} at the center.
It gives an actual functor $Q:C\to\bar C$ replacing $M$ by
the trivial monoid and $S_i$ by the left $N_i$-set
$T_i=S_i/M$. Write $Y=B\bar C$.
The map $BQ$ is onto on fundamental groups, induces rational
and $\mathbb F_{65521}$ cohomology-ring isomorphisms, and has
a homotopy left inverse $h:Y\to BC$. The prime is good
because every relevant finite group has order at most $22$.

There is no nonidentity arrow between different branches.
Thus $Y$ is the actual wedge of the ordinary left-act cones
\[
 Y_i=B\mathcal C^{\rm left}(N_i,T_i)
\]
at their common trivial central vertex. Each $Y_i$ is connected.
Since $Y$ is simply connected, the free-product formula implies
that every $Y_i$ is simply connected. The finite category nerve
has finitely many simplices in each degree, so all its homology
groups in a fixed degree are finitely generated. In particular,
simple connectivity and rational second rank one give
$H^2(Y;\mathbb Z)\cong\mathbb Z$.

Let $u$ generate $H^2(BC;\mathbb Z)$ and put $v=h^*u$.
The equation $h\,BQ\simeq1$ shows that $v$ is an integral
generator, and that $v^2$ is nonzero and primitive modulo
torsion in $H^4(Y;\mathbb Z)$. The direct sum decomposition
of second cohomology for this finite wedge puts $v$ in one
unique summand $Y_j$; all other integral second-cohomology
groups are zero. Write $v_j$ for its restriction there.
The rational cohomology of $Y$ is that of the projective plane,
and $v_j^2\ne0$, so $Y_j$ has exactly that rational homology.
The square $v_j^2$ is primitive modulo torsion: divisibility
in this wedge summand would give the same divisibility of
$v^2$ in the direct sum.

The map $h$ restricted to $Y_j$, followed by a chosen equivalence
$BC\simeq\mathbb CP^2$, pulls the projective-plane generator
back to $v_j$. Miller's theorem therefore makes the restriction
of $v_j$ zero along every map from a finite group. Over
$\mathbb F_{65521}$, the class $v^2$ remains nonzero because
$(BQ)^*v^2=u^2$. It lies in the same wedge summand $Y_j$.
As the total Betti numbers in degrees two, three and four
are $(1,0,1)$, those of $Y_j$ must also be $(1,0,1)$.

This ordinary act cone has at most twenty arrows. Indeed,
there is at least one other nonempty branch, whose leaf
identity and one cross arrow contribute two arrows; hence
\[
 |N_j|+|T_j|+1
 \leq |C|-|M|-2+1\leq20.
\]
All four conditions of
Theorem~\ref{C:finite-target-small-act-twenty} now hold for
$Y_j$, after opposition to turn its left act into a right act.
This is impossible. Taking opposites proves the inward version.
No finite-CW-type assertion about $Y$ or its wedge summands
has been used.
\end{proof}

\begin{corollary}\label{C:star-contractible-center18}
A finite inward or outward star with at least two leaves,
at most eighteen arrows, and contractible central monoid
classifying space cannot have type $\mathbb CP^2$.
In particular this applies if the central monoid has a
left zero or a right zero.
\end{corollary}
\begin{proof}
The full branch nerves meet in the contractible central
nerve, so their actual union is homotopy equivalent to
their wedge. The free-product fundamental-group formula
and the integral cohomology ring of $\mathbb CP^2$ then
force one branch to have type $\mathbb CP^2$: precisely one
summand supports both the degree-two generator and its
nonzero square, and every other simply connected summand
is acyclic and contractible. Each branch has at most
sixteen arrows, since a different branch contributes at
least two. Theorem~\ref{C:CP2-lower17} excludes it.
A left or right zero gives a natural transformation
between the identity and a constant functor on the
one-object category, proving the stated sufficient condition.
\end{proof}

\begin{corollary}\label{C:star-small-center18}
An inward or outward star category with at least two leaves,
at most eighteen arrows, and central endomorphism monoid of
order at most four cannot have type $\mathbb CP^2$.
\end{corollary}
\begin{proof}
Take opposites to make the star outward and write $M$ for
its center. If the trivial left rational $M$-module is
projective, the preceding theorem applies. Otherwise the
minimal ideal has at least two rows. Its identity lies
outside that ideal. Since $|M|\leq4$, its minimal Rees
data must have one column and trivial maximal group.
A minimal idempotent $e$ then satisfies $em=e$ for every
$m\in M$. It is a left zero, so $BM$ is contractible.

The nerve of the star is the union of its branch nerves,
intersecting in the full central nerve $BM$. These are
simplicial cofibrations. Their homotopy pushout over the
contractible common subspace is therefore a wedge of the
individual two-object branch spaces. If this wedge has
type $\mathbb CP^2$, one branch has that type. Indeed each
branch is simply connected by the free-product formula;
reduced integral homology splits over the wedge, and the
primitive nonzero square lies in its unique degree-two
summand. That same summand carries degree four. All other
summands are simply connected and acyclic, hence contractible,
and the chosen branch is homotopy equivalent to the wedge.

Every branch has at most sixteen arrows, since another
branch contributes at least its leaf identity and one
cross arrow. Theorem~\ref{C:CP2-lower17} excludes it.
\end{proof}

% Root proof; D independent complete literal audit passed.
\subsection{Bounded act cones through eleven monoid elements}

\begin{lemma}\label{C:eleven-bounded-all-products}
Let $M$ be a finite monoid with $|M|\leq11$ and let $S$ be
any nonempty finite right $M$-set. If its ordinary act cone
has bounded rational homology, then every product of two
rational degree-two classes is zero. Neither a first-homology
condition nor a bound on the second-cohomology rank is required.
\end{lemma}
\begin{proof}
If $H^2(BM;\mathbb Q)=0$, every degree-two class is relative
to the monoid vertex, and the suspension quotient kills products.
Otherwise the minimal Rees ideal has at least two rows and
two columns. At most two rows give the all-order conclusion
of Theorem~\ref{C:two-row-act-cups}. We may assume at least
three rows and at least two columns. The identity lies outside
the ideal, and a nontrivial maximal group would require at
least twelve ideal elements. Thus the ideal is a rectangular
band. The only possible shapes and outside-element bounds are
\[
\begin{array}{c|c}
 \text{minimal band}&q=|M\setminus K|\\\hline
 3\times2&q\leq5\\
 4\times2&q\leq3\\
 5\times2&q\leq1\\
 3\times3&q\leq2.
\end{array}
\]
The first row is Proposition~\ref{C:three-row-five-outside-all-acts}.
The next two follow from
Lemma~\ref{C:three-outside-all-two-column-act-cups}.
For the last row, Lemma~\ref{C:two-outside-act-obstruction}
and the act-cone Tor formula make fourth rational homology
zero under the stated boundedness hypothesis. This excludes
all degree-two products in the last case as well.
\end{proof}

\begin{corollary}\label{C:small-center-projective-leaves-starlimit}
Let $C$ have outward star support with any positive number
of leaves. Suppose its central monoid $M$ has at most eleven
elements and the trivial right rational module at every leaf
monoid is projective. If $BC$ has bounded rational homology,
then every rational degree-two product on $BC$ is zero.
The inward-star assertion is dual, with left modules at the leaves.
There is no bound on the numbers of cross arrows.
\end{corollary}
\begin{proof}
Apply Proposition~\ref{C:finite-target-extreme-corner} at each
leaf. Its actual orbit quotient gives a rational cohomology-ring
equivalence to a star with the same central monoid and trivial
leaves, with nonempty finite right $M$-acts $S_i$ as its legs.
Now identify those leaf vertices. On nerves this gives exactly
the ordinary act cone on $S=\coprod_i S_i$: no nonidentity
arrow exits a leaf. On realizations, identifying these finitely
many points in a connected CW space adds only a wedge of
circles up to homotopy. This preserves bounded rational
homology, cohomology in degrees at least two, and all its
products. The preceding lemma applies to that act cone.
Pulling back the vanishing through these actual quotient maps
proves the assertion for the original star. Opposition gives
the inward version.
\end{proof}

% Root proof; D independent complete literal audit passed.
\subsection{A singleton cone beside a two-state sphere-band leg}

\begin{lemma}\label{C:star-singleton-sphere-two-states}
Let $C$ be an outward three-object star whose center is any
finite monoid $M$. Suppose one leaf is trivial and its leg
has one arrow, while the other leaf is the five-element
identity-adjoined $2\times2$ rectangular band $P_5$ and its
leg has two arrows. Then every rational degree-two product
on $BC$ is zero.
\end{lemma}
\begin{proof}
Write $S$ for the two-element $(P_5,M)$-biset and $E_S$ for
its mixed action bar. The singleton trivial-leaf branch is
the cone on $BM$. Associativity of the homotopy pushout
therefore gives the actual comparison
\[
 BC\simeq\operatorname{hocofib}(E_S\longrightarrow BP_5).
\]
Classify the left $P_5$-action on $S$. Every band actor is
an idempotent transformation of two points, hence an identity
or a constant. If one actor is the identity, the rectangular
relations force all four actors to be identities. Otherwise
all are constant and their values depend only on the row.
This leaves exactly three cases: all identities, one common
constant, or the two distinct constants associated to the rows.

In the identity case the left action category is the disjoint
union of two copies of $BP_5$, and the commuting right
$M$-action changes only the state. Thus
$E_S=B(S\mathbin{/\mkern-6mu/}M)\times BP_5$, and its map
to $BP_5$ is projection. It has a section, so restriction
of every degree-two cofiber class to $BP_5$ is zero.
Those classes pull back from the suspension obtained by
collapsing $BP_5$ in the cofiber; their products vanish.

In the common-constant case commutation forces its value
$s_0$ to be fixed by $M$. The full action subcategory at
$s_0$ supplies a copy of $BM\times BP_5$ inside $E_S$,
whose projection to $BP_5$ is the second projection.
Choosing a point of $BM$ again gives a section of
$E_S\to BP_5$, and the same relative-suspension argument
kills every product.

In the remaining case $S$ is the left row set of $P_5$.
The only $P_5$-equivariant self-map of this set is the
identity, since it must fix each row constant. Therefore
the right $M$-action is trivial and
\[
 E_S=BM\times B(P_5\ltimes S).
\]
The left set $S$ is $P_5e$ for a minimal idempotent $e$,
so its action category has contractible classifying space.
For an explicit contraction, every object is $x\in P_5e$;
the arrow labelled $x$ from $e$ to $x$ is natural, since
an action arrow labelled $n$ has $nx=y$ and the composite
label is $nx=y$. These arrows define a natural transformation
from the constant functor at $e$ to the identity functor.
The map $E_S\to BP_5$ factors through this contractible
action-category space and is nullhomotopic. Its cofiber
is consequently $BP_5\vee\Sigma E_S$.
As $BP_5\simeq S^2$, both wedge summands have zero
positive-degree products. This proves the assertion.
\end{proof}

\begin{corollary}\label{C:star-eight-one-five-singleton-two}
Every outward star with corner and leg sizes
$(m,a,b,x,y)=(8,1,5,1,2)$ is excluded from type
$\mathbb CP^2$; the same holds after opposition.
\end{corollary}
\begin{proof}
The singleton branch kills restriction of every global
degree-two class to the central monoid. If the five-element
leaf has zero rational second cohomology, all such classes
are relative to all vertex spaces. The quotient is a wedge
of suspended edge bars, so products vanish. The only
five-element monoid with nonzero rational second cohomology
is $P_5$, by the independently checked small-monoid census.
The preceding lemma handles that case. Opposition preserves
classifying-space homotopy type and the cup obstruction.
\end{proof}

